\documentclass[10pt,a4paper,oneside]{book}
\usepackage[utf8]{inputenc}
\usepackage[T1]{fontenc}
\usepackage[french]{babel}
\usepackage{lmodern}
\usepackage[margin=1.8cm]{geometry}
\usepackage{amsmath,amssymb,amsfonts}
\usepackage{booktabs,longtable,array}
\usepackage{xcolor}
\usepackage{listings}
\usepackage{fancyhdr}
\usepackage{hyperref}

\definecolor{codebg}{rgb}{0.97,0.97,0.98}
\definecolor{codeframe}{rgb}{0.82,0.84,0.88}
\definecolor{kwcolor}{rgb}{0.0,0.28,0.65}
\definecolor{comcolor}{rgb}{0.25,0.50,0.35}
\definecolor{strcolor}{rgb}{0.65,0.15,0.15}

\lstset{
  backgroundcolor=\color{codebg},
  frame=single,
  rulecolor=\color{codeframe},
  basicstyle=\ttfamily\scriptsize,
  keywordstyle=\color{kwcolor}\bfseries,
  commentstyle=\color{comcolor}\itshape,
  stringstyle=\color{strcolor},
  numbers=left,
  numberstyle=\tiny\color{gray},
  stepnumber=1,
  numbersep=6pt,
  breaklines=true,
  breakatwhitespace=false,
  tabsize=2,
  showstringspaces=false,
  upquote=true,
  literate=
    {é}{{\'e}}1 {è}{{\`e}}1 {ê}{{\^e}}1 {ë}{{\"e}}1
    {à}{{\`a}}1 {â}{{\^a}}1 {ä}{{\"a}}1
    {î}{{\^i}}1 {ï}{{\"i}}1
    {ô}{{\^o}}1 {ö}{{\"o}}1
    {ù}{{\`u}}1 {û}{{\^u}}1 {ü}{{\"u}}1
    {ç}{{\c{c}}}1
    {É}{{\'E}}1 {È}{{\`E}}1 {Ê}{{\^E}}1
    {À}{{\`A}}1 {Â}{{\^A}}1
    {Î}{{\^I}}1 {Ô}{{\^O}}1 {Û}{{\^U}}1 {Ç}{{\c{C}}}1
    {œ}{{\oe}}1 {Œ}{{\OE}}1
    {°}{{\textdegree}}1
    {×}{{$\times$}}1 {÷}{{$\div$}}1 {±}{{$\pm$}}1
    {≤}{{$\le$}}1 {≥}{{$\ge$}}1 {≠}{{$\ne$}}1 {≈}{{$\approx$}}1
    {∞}{{$\infty$}}1 {∂}{{$\partial$}}1 {∇}{{$\nabla$}}1
    {α}{{$\alpha$}}1 {β}{{$\beta$}}1 {γ}{{$\gamma$}}1 {δ}{{$\delta$}}1
    {ε}{{$\varepsilon$}}1 {η}{{$\eta$}}1 {θ}{{$\theta$}}1 {κ}{{$\kappa$}}1
    {λ}{{$\lambda$}}1 {μ}{{$\mu$}}1 {ν}{{$\nu$}}1 {ξ}{{$\xi$}}1
    {π}{{$\pi$}}1 {ρ}{{$\rho$}}1 {σ}{{$\sigma$}}1 {τ}{{$\tau$}}1
    {φ}{{$\varphi$}}1 {χ}{{$\chi$}}1 {ψ}{{$\psi$}}1 {ω}{{$\omega$}}1
    {Δ}{{$\Delta$}}1 {Ω}{{$\Omega$}}1 {Ψ}{{$\Psi$}}1
    {—}{{---}}1 {–}{{--}}1 {’}{{'}}1 {‘}{{`}}1
}

\pagestyle{fancy}
\fancyhf{}
\fancyhead[L]{\leftmark}
\fancyhead[R]{\thepage}
\renewcommand{\headrulewidth}{0.4pt}

\hypersetup{
  colorlinks=true,
  linkcolor=blue!70!black,
  urlcolor=blue!70!black,
  pdftitle={Guide du Développeur Exhaustif COVOF 2D (C++17 et Fortran 2018)},
  pdfauthor={Laboratoire TREFLE / I2M Bordeaux - Projet COVOF 2D}
}

\begin{document}

\begin{titlepage}
\centering
\vspace*{2cm}
{\Huge\bfseries GUIDE DU DÉVELOPPEUR EXHAUSTIF\\[0.4cm] SOLVEUR DIPHASIQUE ET THERMOCONVECTIF\\[0.4cm] COVOF 2D\par}
\vspace{1cm}
{\Large\bfseries Architecture Logicielle, Formulations Mathématiques Discrètes, Algorithmes Numériques et Analyse Intégrale des Codes Sources\par}
\vspace{1.2cm}
\rule{\textwidth}{1.5pt}\\[0.5cm]
{\LARGE\bfseries PARTIE I : Implémentation Orientée Objet C++17 (\texttt{COVOF\_2D-C++})\\[0.3cm]
PARTIE II : Implémentation Modulaire Fortran 2018 (\texttt{COVOF\_2D-FORTRAN})\par}
\vspace{0.3cm}
\rule{\textwidth}{1.5pt}
\vfill
{\large
\textbf{Méthodes Numériques Documentées :}\\
Volumes Finis sur Grille Décalée MAC $\bullet$ Lagrangien Augmenté Monolithique ($r_{al}$) et Algorithme d'Uzawa\\
Solveurs de Krylov Parallèles OpenMPI : BiCGSTAB(2), BiCGSTAB, PCG, GMRES($m$), MINRES\\
Capture d'Interface VOF : Reconstruction Géométrique PLIC, ELVIRA, LVIRA, SLIC\\
Advection Conservatrice : Weymouth \& Yue (2010), CIAM (Li 1995), TVD (Superbee, MC, Van Leer, Minmod, HRIC)\\
Tension Superficielle CSF Balanced-Force $\bullet$ Pénalisation Volumique de Brinkman $\bullet$ Thermoconvection de Boussinesq\\[1cm]
\textbf{Version 3.0.0 — Octobre 2026}
}
\end{titlepage}

\tableofcontents
\cleardoublepage
\part{Guide du Développeur C++17 (\texttt{COVOF\_2D-C++})}
\chapter{Architecture Globale C++17, Système de Build CMake et Types Fondamentaux (Types.hpp, main.cpp)}
\section{Philosophie Architecturale et Standard C++17}
L'implémentation \textbf{C++17} du solveur \textbf{COVOF 2D} (\texttt{COVOF\_2D-C++}) repose sur une architecture orientée objet modulaire à couplage faible et forte cohésion interne, encapsulée dans l'espace de noms \texttt{covof}. Les principes directeurs d'ingénierie logicielle retenus sont les suivants :
\begin{enumerate}
  \item \textbf{Zéro allocation dynamique dans la boucle temporelle} : tous les champs eulériens (\texttt{Field2D<double>}), les tampons d'échange de halos MPI et la structure de la matrice creuse CSR (\texttt{CSRMatrix}) sont pré-alloués lors de la phase d'initialisation (\texttt{Simulation::initialize}).
  \item \textbf{Contiguïté mémoire stricte (\emph{Row-Major Order})} : chaque champ 2D de dimension locale $N_x \times N_y$ entouré de $n_g$ couches de cellules fantômes (\emph{ghost cells}) est stocké dans un unique \texttt{std::vector<double>} unidimensionnel de taille $(N_x + 2n_g)(N_y + 2n_g)$. L'indice linéaire correspondant à la maille $(i, j)$ est donné par :
  \begin{equation}
    k_{\text{C++}}(i, j) = (j + n_g) \times (N_x + 2n_g) + (i + n_g)
  \end{equation}
  Ainsi, une boucle interne parcourant l'indice horizontal $i \in \{0, \dots, N_x-1\}$ à ordonnée $j$ fixée accède à des adresses mémoire strictement consécutives ($\Delta k = 1$), maximisant l'utilisation des lignes de cache L1/L2 de 64 octets et permettant la vectorisation automatique SIMD (AVX2 / AVX-512 / NEON Apple Silicon).
  \item \textbf{Typage fort des modèles physiques et numériques (\texttt{Types.hpp})} : l'ensemble des choix algorithmiques (type de condition aux limites \texttt{BoundaryType}, schéma d'advection VOF \texttt{AdvectionMethod}, méthode de reconstruction géométrique \texttt{ReconstructionMethod}, limiteur de flux \texttt{TVDLimiter}, solveur de Krylov \texttt{SolverType}, préconditionneur \texttt{PreconditionerType}, mode d'initialisation \texttt{InitialConditionType}) est représenté par des énumérations fortement typées (\texttt{enum class}) interdisant toute conversion implicite erronée.
\end{enumerate}

\section{Structure des Blocs de Configuration dans \texttt{Types.hpp}}
L'en-tête \texttt{include/covof/Types.hpp} définit également l'intégralité des structures de données de paramétrage (\emph{Plain Old Data structs}) qui reflètent bijectivement l'arborescence des fichiers JSON du répertoire \texttt{cases/*.json} :
\begin{itemize}
  \item \texttt{MeshConfig} : dimensions entières globales \texttt{nx}, \texttt{ny} et dimensions physiques \texttt{lx}, \texttt{ly} $[\text{m}]$.
  \item \texttt{PhysicsConfig} : masses volumiques \texttt{rho1}, \texttt{rho2} $[\text{kg/m}^3]$, viscosités dynamiques \texttt{mu1}, \texttt{mu2} $[\text{Pa}\cdot\text{s}]$, gravité $(\texttt{gx}, \texttt{gy})$ $[\text{m/s}^2]$, tension superficielle \texttt{sigma} $[\text{N/m}]$ et paramètres de viscosité variable exponentielle \texttt{variable\_viscosity}, \texttt{visc\_law}, \texttt{visc\_a}, \texttt{visc\_e}.
  \item \texttt{ThermalConfig} : drapeau \texttt{enable\_energy}, diffusivité thermique $\alpha$ (\texttt{alpha\_th}), coefficient d'expansion thermique de Boussinesq $\beta$ (\texttt{beta\_th}), température de référence $T_0$ (\texttt{t\_ref}) et températures imposées (\texttt{t\_hot}, \texttt{t\_cold}).
  \item \texttt{VofConfig} : méthode de reconstruction géométrique (\texttt{PLIC}, \texttt{ELVIRA}, \texttt{LVIRA}, \texttt{SLIC}, \texttt{NONE}), schéma d'advection (\texttt{WEYMOUTH\_YUE}, \texttt{CIAM}, \texttt{TVD}), limiteur TVD (\texttt{SUPERBEE}, \texttt{MC}, \texttt{VAN\_LEER}, \texttt{MINMOD}, \texttt{HRIC}) et nombre de Courant interfacial maximal \texttt{cfl\_vof}.
  \item \texttt{SolverConfig} : solveur de Krylov (\texttt{BICGSTAB2}, \texttt{BICGSTAB}, \texttt{PCG}, \texttt{GMRES}, \texttt{MINRES}), préconditionneur (\texttt{JACOBI}, \texttt{SSOR}, \texttt{ILU0}), tolérance relative \texttt{tol}, nombre maximal d'itérations \texttt{max\_iter}, paramètre de pénalisation du Lagrangien Augmenté $r_{al}$ (\texttt{r\_al}) et nombre de sous-itérations d'Uzawa (\texttt{uzawa\_iter}).
\end{itemize}

\section{Cycle de Vie du Programme Principal (\texttt{src/main.cpp})}
Le point d'entrée \texttt{main(int argc, char* argv[])} orchestre le cycle de vie MPI et l'exécution complète d'un cas test en quatre étapes :
\begin{enumerate}
  \item Initialisation de l'environnement parallèle via \texttt{MPI\_Init(\&argc, \&argv)} et interrogation du rang global \texttt{rank} et du nombre de processus \texttt{size}.
  \item Lecture et validation du fichier JSON passé en argument de ligne de commande via \texttt{Config::load\_from\_file(json\_path)}.
  \item Instanciation de l'objet maître \texttt{Simulation sim(cfg, MPI\_COMM\_WORLD)}, appel à \texttt{sim.initialize()} puis exécution de la boucle temporelle \texttt{sim.run()} chronométrée par \texttt{MPI\_Wtime()}.
  \item Libération ordonnée des ressources MPI via \texttt{MPI\_Finalize()}.
\end{enumerate}

\section{Code Source Intégral Annoté : \texttt{CMakeLists.txt}}
Le listing ci-dessous présente l'intégralité des 48 lignes du fichier \texttt{CMakeLists.txt} de l'implémentation C++17 :

\begin{lstlisting}[language={bash}]
cmake_minimum_required(VERSION 3.16)
project(COVOF_2D_CPP CXX)

set(CMAKE_CXX_STANDARD 17)
set(CMAKE_CXX_STANDARD_REQUIRED ON)
set(CMAKE_EXPORT_COMPILE_COMMANDS ON)

# Optimization and warning flags
set(CMAKE_CXX_FLAGS "${CMAKE_CXX_FLAGS} -O3 -Wall -Wextra -Wno-unused-parameter")

# Find OpenMPI / MPI
find_package(MPI REQUIRED)
include_directories(SYSTEM ${MPI_INCLUDE_PATH})

# Include directory
include_directories(${CMAKE_CURRENT_SOURCE_DIR}/include)

# Source files
set(SOURCES
    src/Config.cpp
    src/Mesh.cpp
    src/MpiComm.cpp
    src/CSRMatrix.cpp
    src/Preconditioner.cpp
    src/LinearSolver.cpp
    src/InterfaceReconstruction.cpp
    src/InterfaceTracker.cpp
    src/NavierStokesSolver.cpp
    src/VTKWriter.cpp
    src/Simulation.cpp
    src/main.cpp
)

# Executable target: covof2d
add_executable(covof2d ${SOURCES})
target_link_libraries(covof2d PRIVATE MPI::MPI_CXX)

set_target_properties(covof2d PROPERTIES
    RUNTIME_OUTPUT_DIRECTORY "${CMAKE_CURRENT_SOURCE_DIR}/bin"
)

# Create compatibility symlink/target for fugu2d
add_custom_command(TARGET covof2d POST_BUILD
    COMMAND ${CMAKE_COMMAND} -E create_symlink
        ${CMAKE_CURRENT_SOURCE_DIR}/bin/covof2d
        ${CMAKE_CURRENT_SOURCE_DIR}/bin/fugu2d
    COMMENT "Creating compatibility symlink bin/fugu2d -> bin/covof2d"
)

\end{lstlisting}

\section{Répertoire des Classes, Structures et Méthodes de \texttt{include/covof/Types.hpp}}
\begin{center}\small
\begin{tabular}{clll}
\toprule
\textbf{Ligne} & \textbf{Catégorie} & \textbf{Identifiant C++17} & \textbf{Signature / Rôle} \\
\midrule
14 & enum class & \texttt{BCType} & Déclaration de type / classe / structure \\
26 & struct & \texttt{BoundaryCondition} & Déclaration de type / classe / structure \\
52 & struct & \texttt{Obstacle} & Déclaration de type / classe / structure \\
73 & struct & \texttt{DomainBC} & Déclaration de type / classe / structure \\
87 & struct & \texttt{FluidProperties} & Déclaration de type / classe / structure \\
103 & enum class & \texttt{InitialConditionType} & Déclaration de type / classe / structure \\
119 & struct & \texttt{InitialConditionConfig} & Déclaration de type / classe / structure \\
149 & struct & \texttt{SolverOptions} & Déclaration de type / classe / structure \\
169 & struct & \texttt{OutputConfig} & Déclaration de type / classe / structure \\
185 & struct & \texttt{VofConfig} & Déclaration de type / classe / structure \\
\bottomrule
\end{tabular}
\end{center}

\section{Code Source Intégral Annoté : \texttt{include/covof/Types.hpp}}
Le listing ci-dessous présente l'intégralité des 196 lignes du fichier \texttt{include/covof/Types.hpp} de l'implémentation C++17 :

\begin{lstlisting}[language={[11]C++}]
#pragma once

#include <string>
#include <vector>
#include <iostream>
#include <cmath>

namespace covof {
namespace fugu = covof; // Alias rétrocompatibilité

/**
 * @brief Types de conditions aux limites physiques (Annexe B thèse M. El Ouafa)
 */
enum class BCType {
    // Calcul et assignation numérique
    DIRICHLET,  ///< Vitesse imposée (adhérence u=0, paroi mobile u=U_wall, ou profil d'entrée)
    // Calcul de norme euclidienne ou résidu de convergence
    NEUMANN,    ///< Dérivée normale nulle (sortie libre dU/dn = 0)
    FREE_SLIP,  ///< Glissement libre (composante normale nulle, contrainte tangentielle nulle)
    PERIODIC    ///< Conditions aux limites périodiques
};

/**
 * @brief Spécification des conditions aux limites sur les 4 frontières du domaine 2D
 */
struct BoundaryCondition {
    // Calcul et assignation numérique
    BCType type = BCType::DIRICHLET;
    // Déclaration de variable typée
    double u_val = 0.0;  ///< Valeur de u imposée si Dirichlet
    // Déclaration de variable typée
    double v_val = 0.0;  ///< Valeur de v imposée si Dirichlet
    // Déclaration de variable typée
    std::string profile = "UNIFORM"; ///< "UNIFORM", "PARABOLIC", "POISEUILLE_STEP"
    // Déclaration de variable typée
    double u_mean = 0.0; ///< Vitesse moyenne pour profils paraboliques
    // Déclaration de variable typée
    double y_step = 0.0; ///< Hauteur de la marche pour POISEUILLE_STEP
    // Déclaration de variable typée pour la largeur de l'injecteur (Cas 3.3 Fluids 2021)
    double jet_width = 0.08; ///< Largeur de l'injecteur pour profil JET_INLET [m]
    // Déclaration de variable typée pour le centre de l'injecteur
    double jet_center = 0.5; ///< Position centrale de l'injecteur selon X [m]
    // Calcul et assignation numérique
    BCType T_type = BCType::NEUMANN; ///< Type de CL pour l'énergie (DIRICHLET ou NEUMANN)
    // Déclaration de variable typée
    double T_val = 0.0;             ///< Température imposée si DIRICHLET
};

/**
 * @brief Obstacle immergé avec pénalisation de Brinkman (Standard CFD)
 */
struct Obstacle {
    // Déclaration de variable typée
    std::string type = "BOX"; ///< Forme géométrique ("BOX", "CYLINDER")
    // Déclaration de variable typée
    double x0 = 0.0;          ///< Borne inférieure X [m]
    // Déclaration de variable typée
    double y0 = 0.0;          ///< Borne inférieure Y [m]
    // Déclaration de variable typée
    double x1 = 0.0;          ///< Borne supérieure X [m]
    // Déclaration de variable typée
    double y1 = 0.0;          ///< Borne supérieure Y [m]
    // Déclaration de variable typée
    double xc = 0.0;          ///< Centre cylindre X [m]
    // Déclaration de variable typée
    double yc = 0.0;          ///< Centre cylindre Y [m]
    // Déclaration de variable typée
    double radius = 0.0;      ///< Rayon cylindre [m]
    // Déclaration de variable typée
    double permeability = 1.0e-14; ///< Perméabilité de Darcy-Brinkman K
};

struct DomainBC {
    // Calcul et assignation numérique
    BoundaryCondition left;    ///< Frontière x = 0
    // Calcul et assignation numérique
    BoundaryCondition right;   ///< Frontière x = Lx
    // Calcul et assignation numérique
    BoundaryCondition bottom;  ///< Frontière y = 0
    // Calcul et assignation numérique
    BoundaryCondition top;     ///< Frontière y = Ly
};

/**
 * @brief Propriétés physiques d'une phase fluide
 */
struct FluidProperties {
    // Déclaration de variable typée
    double rho = 1.0;         ///< Masse volumique [kg/m^3]
    // Déclaration de variable typée
    double mu  = 1.0;         ///< Viscosité dynamique [Pa.s]
    // Déclaration de variable typée
    double Cp = 1006.0;       ///< Capacité calorifique à pression constante [J/(kg.K)]
    // Déclaration de variable typée
    double k_thermal = 0.026; ///< Conductivité thermique [W/(m.K)]
    // Déclaration de variable typée
    double beta = 3.33e-3;    ///< Coefficient de dilatation thermique volumique [1/K]
};

/**
 * @brief Types de configurations initiales de la fraction volumique C
 */
enum class InitialConditionType {
    // Calcul et assignation numérique
    SINGLE_PHASE,    ///< Écoulement monophasique (C = 1 partout)
    DROPLET,         ///< Goutte (Fluide 1 à l'intérieur, Fluide 2 à l'extérieur)
    BUBBLE,          ///< Bulle (Fluide 2 à l'intérieur, Fluide 1 à l'extérieur)
    DAM_BREAK,       ///< Colonne rectangulaire d'eau (rupture de barrage)
    SHEAR_LAYER,     ///< Nappe liquide cisaillée avec perturbation
    LIQUID_SHEET,    ///< Nappe liquide cisaillée périodique (Cas 3.2 Fluids 2021)
    VISCOUS_JET,     ///< Jet liquide visqueux injecté / flambement (Cas 3.3 Fluids 2021)
    RAYLEIGH_TAYLOR, ///< Interface perturbée avec fluide lourd au-dessus
    VORTEX_SHEAR     ///< Tâche circulaire déformée par vortex avec cisaillement
};

/**
 * @brief Paramètres géométriques de la condition initiale diphasique
 */
struct InitialConditionConfig {
    // Calcul et assignation numérique
    InitialConditionType type = InitialConditionType::SINGLE_PHASE;
    // Déclaration de variable typée
    double center_x = 0.5;   ///< Centre de la goutte en X [m]
    // Déclaration de variable typée
    double center_y = 0.5;   ///< Centre de la goutte en Y [m]
    // Déclaration de variable typée
    double radius   = 0.25;  ///< Rayon de la goutte [m]
    // Déclaration de variable typée
    double box_x0   = 0.0;   ///< Coin inf-gauche du rectangle (Dam break)
    // Déclaration de variable typée
    double box_y0   = 0.0;
    // Déclaration de variable typée
    double box_x1   = 1.0;   ///< Coin sup-droit du rectangle (Dam break)
    // Déclaration de variable typée
    double box_y1   = 1.0;
    // Déclaration de variable typée
    double pert_amp = 0.01;  ///< Amplitude de perturbation (Rayleigh-Taylor)
    // Déclaration de variable typée
    double pert_k   = 1.0;   ///< Nombre d'onde de perturbation
    // Déclaration de variable typée pour l'épaisseur de nappe liquide
    double thickness = 0.0003; ///< Épaisseur de la nappe liquide (Cas 3.2 Fluids 2021) [m]
    // Déclaration de variable typée pour la largeur du jet liquide
    double jet_width = 0.08;   ///< Largeur du jet liquide injecté (Cas 3.3 Fluids 2021) [m]
};

/**
 * @brief Options de paramétrage du solveur linéaire de Krylov BiCGSTAB(2)
 */
struct SolverOptions {
    // Déclaration de variable typée
    std::string type = "AUGMENTED_LAGRANGIAN"; ///< AUGMENTED_LAGRANGIAN ou MONOLITHIC
    // Déclaration de variable typée
    double tol = 1.0e-7;                        ///< Résidu relatif d'arrêt ||r_k|| / ||r_0||
    // Déclaration de variable typée
    int max_iter = 1000;                        ///< Nombre maximum d'itérations
    // Déclaration de variable typée
    double r_al = 500.0;                        ///< Paramètre de pénalité de divergence r [Pa.s ou kg/(m.s)]
    // Déclaration de variable typée
    int uzawa_iter = 1;                         ///< Nombre d'itérations d'Uzawa-Lagrangien augmenté
    // Déclaration de variable typée
    bool use_precond = true;                    ///< Activer le préconditionneur
    // Déclaration de variable typée
    std::string precond_type = "JACOBI";        ///< JACOBI, BLOCK_TRIANGULAR, ou NONE
};

/**
 * @brief Options de gestion des sorties et de la visualisation ParaView
 */
struct OutputConfig {
    // Déclaration de variable typée
    std::string output_dir = "output"; ///< Répertoire d'écriture
    // Déclaration de variable typée
    std::string case_name  = "covof";   ///< Préfixe des fichiers VTK
    // Déclaration de variable typée
    int vtk_interval = 10;             ///< Pas d'export VTK (tous les N pas de temps)
    // Déclaration de variable typée
    int log_interval = 1;              ///< Pas d'affichage console et CSV
    // Déclaration de variable typée
    bool write_csv_history = true;     ///< Écrire l'historique temporel (centre de masse, vitesse...)
};

/**
 * @brief Options de capture d'interface et d'advection VOF
 */
struct VofConfig {
    // Déclaration de variable typée
    std::string reconstruction_method = "PLIC";       ///< "PLIC", "ELVIRA", "LVIRA", "SLIC", "NONE"
    // Déclaration de variable typée
    std::string advection_method = "WEYMOUTH_YUE";    ///< "WEYMOUTH_YUE", "CIAM", "VOF_TVD"
    // Déclaration de variable typée
    std::string tvd_limiter = "NONE";                 ///< "SUPERBEE", "VAN_LEER", "MUSCL", "HRIC", "NONE"
    // Déclaration de variable typée
    double cfl_vof = 0.5;                             ///< Nombre de CFL maximum pour le transport VOF
};

} // namespace covof

\end{lstlisting}

\section{Répertoire des Classes, Structures et Méthodes de \texttt{src/main.cpp}}
\begin{center}\small
\begin{tabular}{clll}
\toprule
\textbf{Ligne} & \textbf{Catégorie} & \textbf{Identifiant C++17} & \textbf{Signature / Rôle} \\
\midrule
18 & Méthode / Fonction & \texttt{main} & Type retour : int \\
\bottomrule
\end{tabular}
\end{center}

\section{Code Source Intégral Annoté : \texttt{src/main.cpp}}
Le listing ci-dessous présente l'intégralité des 106 lignes du fichier \texttt{src/main.cpp} de l'implémentation C++17 :

\begin{lstlisting}[language={[11]C++}]
// ==============================================================================
// Projet        : COVOF 2D (Convection & Volume-Of-Fluid 2D CFD Solver)
// Fichier       : main.cpp
// Description   : Point d'entrée principal de l'exécutable COVOF 2D C++17 / OpenMPI.
//                 Initialise l'environnement parallèle MPI, analyse la ligne
//                 de commande, charge la configuration JSON, orchestre l'exécution
//                 de la simulation numérique et gère la libération des ressources.
// Auteur        : Elhadji Ibrahima Thiam
// Standard      : C++17 / MPI 3+
// ==============================================================================

#include "covof/Simulation.hpp" // Orchestrateur principal de simulation (grille, solveur, temps)
#include <mpi.h>               // Bibliothèque Message Passing Interface pour le calcul parallèle
#include <iostream>            // Flux d'entrée/sortie standard pour l'affichage console
#include <string>              // Gestion des chaînes de caractères pour les chemins de fichiers

// Point d'entrée standard du programme en ligne de commande
int main(int argc, char** argv) {
    // 1. Initialisation de l'environnement d'exécution parallèle MPI
    // Distribue les arguments de la ligne de commande à tous les processus participants
    MPI_Init(&argc, &argv);

    // Récupère l'identifiant unique (rang) du processus courant dans le communicateur global
    int rank = 0;
    MPI_Comm_rank(MPI_COMM_WORLD, &rank);

    // Récupère le nombre total de processus alloués pour l'exécution parallèle
    int size = 1;
    MPI_Comm_size(MPI_COMM_WORLD, &size);

    // Fichier de configuration par défaut si aucun argument n'est fourni par l'utilisateur
    std::string config_file = "cases/cavity.json";

    // 2. Traitement des arguments de la ligne de commande
    if (argc > 1) {
        std::string arg = argv[1];
        // Détection de la demande d'aide (--help ou -h)
        if (arg == "-h" || arg == "--help") {
            // Seul le processus maître (rang 0) affiche le message d'aide pour éviter la duplication
            if (rank == 0) {
                std::cout << "================================================================================\n"
                          << "   COVOF 2D — Solveur CFD Navier-Stokes Incompressible Parallèle (C++17/MPI)   \n"
                          << "================================================================================\n"
                          << "Usage :\n"
                          << "  mpirun -np <N> ./bin/covof2d [fichier_configuration.json]\n\n"
                          << "Exemples d'exécution :\n"
                          << "  mpirun -np 4 ./bin/covof2d cases/cavity_re1000.json\n"
                          << "  mpirun -np 4 ./bin/covof2d cases/rising_bubble.json\n"
                          << "  mpirun -np 2 ./bin/covof2d cases/static_droplet.json\n"
                          << "  mpirun -np 4 ./bin/covof2d cases/dam_break.json\n"
                          << "================================================================================\n";
            }
            // Clôture propre de l'environnement MPI avant de quitter
            MPI_Finalize();
            return 0;
        }
        // Utilisation du chemin spécifié en premier argument
        config_file = arg;
    }

    // Affichage de la bannière d'accueil par le processus racine
    if (rank == 0) {
        std::cout << "================================================================================\n"
                  << "   COVOF 2D — Solveur CFD Navier-Stokes Incompressible Parallèle (C++17/MPI)   \n"
                  << "================================================================================\n"
                  << "  Nombre de processus MPI  : " << size << "\n"
                  << "  Fichier de configuration : " << config_file << "\n"
                  << "================================================================================\n";
    }

    // 3. Bloc d'exécution sécurisé sous surveillance des exceptions C++
    try {
        // Instanciation de l'objet Simulation encapsulant maillage, champs, solveurs et exports
        covof::Simulation sim;

        // Initialisation de la simulation : lecture du JSON, partitionnement spatial, allocation
        if (!sim.init(MPI_COMM_WORLD, config_file)) {
            // En cas d'erreur de parsing ou de cohérence des données, signalement sur le rang 0
            if (rank == 0) {
                std::cerr << "[ERREUR CRITIQUE] Échec de l'initialisation de la simulation COVOF 2D.\n"
                          << "                  Vérifiez la validité du fichier JSON : " << config_file << "\n";
            }
            // Interruption d'urgence coordonnée de tous les processus du communicateur MPI
            MPI_Abort(MPI_COMM_WORLD, 1);
            return 1;
        }

        // Exécution de la boucle temporelle principale de calcul Navier-Stokes / VOF
        sim.run();

    } catch (const std::exception& e) {
        // Capture de toute exception standard (dépassement mémoire, assertion, I/O)
        if (rank == 0) {
            std::cerr << "[EXCEPTION INTERCEPTÉE] " << e.what() << "\n";
        }
        // Interruption immédiate de l'ensemble de la grappe de calcul
        MPI_Abort(MPI_COMM_WORLD, 1);
        return 1;
    }

    // 4. Clôture normale et libération de l'environnement parallèle OpenMPI
    MPI_Finalize();

    // Code de retour standard indiquant le succès de l'exécution
    return 0;
}

\end{lstlisting}

\chapter{Analyseur Syntaxique JSON et Gestion de la Configuration (Config.hpp, Config.cpp)}
\section{Conception de l'Analyseur JSON Autonome Sans Dépendance Externe}
Afin de garantir une portabilité immédiate sur tout supercalculateur ou station de travail sans nécessiter le téléchargement de bibliothèques tierces (telles que \texttt{nlohmann/json} ou \texttt{Boost.PropertyTree}), \texttt{COVOF\_2D-C++} intègre dans \texttt{src/Config.cpp} un analyseur syntaxique JSON déterministe et léger.

\subsection{Grammaire Reconnue et Extraction par Sous-Blocs}
Le parseur opère en trois phases :
\begin{enumerate}
  \item \textbf{Prétraitement et nettoyage lexical} : lecture intégrale du flux fichier dans une chaîne \texttt{std::string}, gestion des espaces, retours chariot et commentaires éventuels.
  \item \textbf{Isolation des objets JSON de premier niveau (\texttt{extract\_section})} : pour une clé de section donnée (par exemple \texttt{"mesh"}, \texttt{"physics"}, \texttt{"vof"}, \texttt{"solver"}), l'automate recherche l'ouverture d'accolade \texttt{\{} correspondante et maintient un compteur de profondeur d'imbrication $\text{depth} = \sum (\texttt{'\{'} - \texttt{'\}'})$ tout en ignorant les accolades situées à l'intérieur de littéraux chaînes entre guillemets. Cela permet d'isoler exactement la sous-chaîne JSON du bloc, y compris lorsqu'il contient des tableaux d'objets imbriqués comme \texttt{"obstacles": [ ... ]}.
  \item \textbf{Désérialisation typée avec valeurs par défaut sécurisées} :
  \begin{itemize}
    \item \texttt{get\_string(section, key, default\_val)} : extrait les chaînes de caractères et normalise la casse pour la résolution des énumérations ;
    \item \texttt{get\_double(section, key, default\_val)} : parse les nombres flottants en notation décimale ou scientifique (\texttt{std::stod}) ;
    \item \texttt{get\_int(section, key, default\_val)} et \texttt{get\_bool(section, key, default\_val)} : extraient les entiers et les booléens (\texttt{true}/\texttt{false}).
  \end{itemize}
\end{enumerate}

\section{Compatibilité Ascendante et Alias de Clés JSON}
Pour assurer une interopérabilité totale entre les différentes générations de fichiers de cas tests et entre les dépôts C++17 et Fortran 2018, \texttt{Config::load\_from\_file} implémente une résolution d'alias intelligente :
\begin{itemize}
  \item Pour le solveur linéaire : reconnaît aussi bien \texttt{"tol"} que \texttt{"tolerance"}, et \texttt{"r\_al"} que \texttt{"augmented\_lagrangian\_r"}.
  \item Pour le module VOF : reconnaît la section dédiée \texttt{"vof"} (\texttt{reconstruction\_method}, \texttt{advection\_method}, \texttt{tvd\_limiter}, \texttt{cfl\_vof}) ainsi que les clés historiques éventuellement placées dans \texttt{"solver"} ou \texttt{"physics"}.
  \item Pour les obstacles immergés de Brinkman : parse dynamiquement le tableau JSON \texttt{"obstacles"} en distinguant les primitives géométriques rectangulaires (\texttt{type = "box"}, définies par $[x_{\min}, x_{\max}] \times [y_{\min}, y_{\max}]$) et circulaires (\texttt{type = "circle"}, définies par le centre $(x_c, y_c)$ et le rayon $R$).
\end{itemize}

\section{Répertoire des Classes, Structures et Méthodes de \texttt{include/covof/Config.hpp}}
\begin{center}\small
\begin{tabular}{clll}
\toprule
\textbf{Ligne} & \textbf{Catégorie} & \textbf{Identifiant C++17} & \textbf{Signature / Rôle} \\
\midrule
12 & class & \texttt{Config} & Déclaration de type / classe / structure \\
92 & Méthode / Fonction & \texttt{load\_from\_json} & Type retour : bool \\
97 & Méthode / Fonction & \texttt{print} & Type retour : void \\
\bottomrule
\end{tabular}
\end{center}

\section{Code Source Intégral Annoté : \texttt{include/covof/Config.hpp}}
Le listing ci-dessous présente l'intégralité des 100 lignes du fichier \texttt{include/covof/Config.hpp} de l'implémentation C++17 :

\begin{lstlisting}[language={[11]C++}]
#pragma once

#include "Types.hpp"
#include <string>

namespace covof {
namespace fugu = covof; // Alias rétrocompatibilité

/**
 * @brief Structure complète regroupant tous les paramètres de simulation
 */
class Config {
public:
    // Paramètres de maillage
    int nx = 64;             ///< Nombre de mailles globales en X
    // Déclaration de variable typée
    int ny = 64;             ///< Nombre de mailles globales en Y
    // Déclaration de variable typée
    double lx = 1.0;         ///< Longueur physique du domaine en X [m]
    // Déclaration de variable typée
    double ly = 1.0;         ///< Hauteur physique du domaine en Y [m]

    // Équations résolues et modèles physiques
    std::string mode = "NAVIER_STOKES";  ///< "NAVIER_STOKES", "LAPLACIAN", "VORTEX_ADVECTION"
    // Déclaration de variable typée
    bool enable_two_phase = false;       ///< Modèle diphasique (true) ou monophasique (false)
    // Déclaration de variable typée
    bool enable_surface_tension = false; ///< Activer la force capillaire CSF
    // Déclaration de variable typée
    bool enable_convection = true;       ///< Activer le transport convectif (Navier-Stokes vs Stokes)
    // Déclaration de variable typée
    bool enable_energy = false;          ///< Activer l'équation de l'énergie (température T)
    // Déclaration de variable typée
    double steady_state_tol = 0.0;       ///< Tolérance d'arrêt sur l'état stationnaire ||u^{n+1}-u^n||

    // Convection thermique Boussinesq (Cas 6.1)
    double Ra = 1.0e5;                   ///< Nombre de Rayleigh
    // Déclaration de variable typée
    double Pr = 0.71;                    ///< Nombre de Prandtl
    // Déclaration de variable typée
    double T_ref = 0.0;                  ///< Température de référence Boussinesq
    // Déclaration de variable typée
    double beta = 3.33e-3;               ///< Coefficient de dilatation volumique [1/K]

    // Modèle de viscosité spatiale (Cas 3.9)
    std::string viscosity_model = "CONSTANT"; ///< "CONSTANT", "EXPONENTIAL"
    // Déclaration de variable typée
    double visc_a = 0.02;
    // Déclaration de variable typée
    double visc_e = 0.01;
    // Déclaration de variable typée
    double visc_mu0 = 1.85e-5;

    // Obstacles immergés (Pénalisation de Brinkman)
    std::vector<Obstacle> obstacles;

    // Propriétés physiques
    FluidProperties fluid1;  ///< Fluide 1 (ex: liquide ou fluide lourd)
    FluidProperties fluid2;  ///< Fluide 2 (ex: gaz ou fluide léger)
    // Déclaration de variable typée
    double sigma = 0.0;      ///< Tension de surface [N/m]
    // Déclaration de variable typée
    double gx = 0.0;         ///< Gravité selon X [m/s^2]
    // Déclaration de variable typée
    double gy = 0.0;         ///< Gravité selon Y [m/s^2]

    // Intégration temporelle
    double dt = 0.001;       ///< Pas de temps initial [s]
    // Déclaration de variable typée
    double t_final = 1.0;    ///< Temps physique final [s]
    // Déclaration de variable typée
    double cfl_max = 0.4;    ///< Nombre de CFL maximum pour le pas de temps adaptatif
    // Déclaration de variable typée
    bool adaptive_dt = false;///< Activer le pas de temps automatique selon CFL

    // Conditions aux limites et conditions initiales
    DomainBC bc;
    InitialConditionConfig init_cond;

    // Solveur linéaire et sorties
    SolverOptions solver;
    OutputConfig output;

    // Capture d'interface et transport VOF
    VofConfig vof;

    /**
     * @brief Charge la configuration depuis un fichier JSON
     * @param filepath Chemin vers le fichier .json du cas test
     * @return true si le chargement a réussi
     */
    bool load_from_json(const std::string& filepath);

    /**
     * @brief Affiche la configuration synthétique dans la console
     */
    void print() const;
};

} // namespace covof

\end{lstlisting}

\section{Répertoire des Classes, Structures et Méthodes de \texttt{src/Config.cpp}}
\begin{center}\small
\begin{tabular}{clll}
\toprule
\textbf{Ligne} & \textbf{Catégorie} & \textbf{Identifiant C++17} & \textbf{Signature / Rôle} \\
\midrule
12 & Méthode / Fonction & \texttt{trim} & Type retour : static std::string \\
23 & Méthode / Fonction & \texttt{extract\_value} & Type retour : static std::string \\
78 & Méthode / Fonction & \texttt{get\_double} & Type retour : static double \\
86 & Méthode / Fonction & \texttt{get\_int} & Type retour : static int \\
94 & Méthode / Fonction & \texttt{get\_bool} & Type retour : static bool \\
105 & Méthode / Fonction & \texttt{get\_string} & Type retour : static std::string \\
112 & Méthode / Fonction & \texttt{parse\_bc\_type} & Type retour : static BCType \\
125 & Méthode / Fonction & \texttt{parse\_single\_bc} & Type retour : static void \\
149 & Méthode / Fonction & \texttt{Config::load\_from\_json} & Type retour : bool \\
150 & Méthode / Fonction & \texttt{file} & Type retour : std::ifstream \\
447 & Méthode / Fonction & \texttt{Config::print} & Type retour : void \\
\bottomrule
\end{tabular}
\end{center}

\section{Code Source Intégral Annoté : \texttt{src/Config.cpp}}
Le listing ci-dessous présente l'intégralité des 503 lignes du fichier \texttt{src/Config.cpp} de l'implémentation C++17 :

\begin{lstlisting}[language={[11]C++}]
#include "covof/Config.hpp"
#include <fstream>
#include <sstream>
#include <iostream>
#include <algorithm>
#include <cctype>

namespace covof {
namespace fugu = covof; // Alias rétrocompatibilité

// Fonctions utilitaires locales pour parser le JSON simplement
static std::string trim(const std::string& str) {
    // Déclaration de variable typée
    size_t first = str.find_first_not_of(" \t\n\r");
    // Évaluation de la condition de contrôle
    if (first == std::string::npos) return "";
    // Déclaration de variable typée
    size_t last = str.find_last_not_of(" \t\n\r");
    // Retour de fonction
    return str.substr(first, (last - first + 1));
}

static std::string extract_value(const std::string& content, const std::string& key) {
    // Déclaration de variable typée
    std::string search = "\"" + key + "\"";
    // Déclaration de variable typée
    size_t pos = content.find(search);
    // Évaluation de la condition de contrôle
    if (pos == std::string::npos) return "";
    // Déclaration de variable typée
    size_t colon = content.find(':', pos + search.length());
    // Évaluation de la condition de contrôle
    if (colon == std::string::npos) return "";
    
    // Déclaration de variable typée
    size_t start = content.find_first_not_of(" \t\n\r", colon + 1);
    // Évaluation de la condition de contrôle
    if (start == std::string::npos) return "";

    // Évaluation de la condition de contrôle
    if (content[start] == '{' || content[start] == '[') {
        // Calcul et assignation numérique
        char open_c = content[start];
        // Calcul et assignation numérique
        char close_c = (open_c == '{') ? '}' : ']';
        // Déclaration de variable typée
        int depth = 1;
        // Déclaration de variable typée
        size_t end = start + 1;
        // Boucle itérative conditionnelle
        while (end < content.length() && depth > 0) {
            // Évaluation de la condition de contrôle
            if (content[end] == open_c) depth++;
            // Branche alternative conditionnelle
            else if (content[end] == close_c) depth--;
            end++;
        }
        // Retour de fonction
        return content.substr(start, end - start);
    // Calcul et assignation numérique
    } else if (content[start] == '"') {
        // Déclaration de variable typée
        size_t end = content.find('"', start + 1);
        // Évaluation de la condition de contrôle
        if (end == std::string::npos) return "";
        // Retour de fonction
        return content.substr(start + 1, end - start - 1);
    } else {
        // Déclaration de variable typée
        size_t end = content.find_first_of(",}\n\r", start);
        // Évaluation de la condition de contrôle
        if (end == std::string::npos) end = content.length();
        // Retour de fonction
        return trim(content.substr(start, end - start));
    }
}

static double get_double(const std::string& block, const std::string& key, double default_val) {
    // Déclaration de variable typée
    std::string val = extract_value(block, key);
    // Évaluation de la condition de contrôle
    if (val.empty()) return default_val;
    try { return std::stod(val); } catch (...) { return default_val; }
}

static int get_int(const std::string& block, const std::string& key, int default_val) {
    // Déclaration de variable typée
    std::string val = extract_value(block, key);
    // Évaluation de la condition de contrôle
    if (val.empty()) return default_val;
    try { return std::stoi(val); } catch (...) { return default_val; }
}

static bool get_bool(const std::string& block, const std::string& key, bool default_val) {
    // Déclaration de variable typée
    std::string val = extract_value(block, key);
    // Évaluation de la condition de contrôle
    if (val == "true" || val == "1") return true;
    // Évaluation de la condition de contrôle
    if (val == "false" || val == "0") return false;
    // Retour de fonction
    return default_val;
}

static std::string get_string(const std::string& block, const std::string& key, const std::string& default_val) {
    // Déclaration de variable typée
    std::string val = extract_value(block, key);
    // Retour de fonction
    return val.empty() ? default_val : val;
}

static BCType parse_bc_type(const std::string& str) {
    // Évaluation de la condition de contrôle
    if (str == "DIRICHLET") return BCType::DIRICHLET;
    // Évaluation de la condition de contrôle
    if (str == "NEUMANN") return BCType::NEUMANN;
    // Évaluation de la condition de contrôle
    if (str == "FREE_SLIP") return BCType::FREE_SLIP;
    // Évaluation de la condition de contrôle
    if (str == "PERIODIC") return BCType::PERIODIC;
    // Retour de fonction
    return BCType::DIRICHLET;
}

static void parse_single_bc(const std::string& block, BoundaryCondition& bc) {
    // Calcul et assignation numérique
    bc.type = parse_bc_type(get_string(block, "type", "DIRICHLET"));
    // Calcul et assignation numérique
    bc.u_val = get_double(block, "u", 0.0);
    // Calcul et assignation numérique
    bc.v_val = get_double(block, "v", 0.0);
    // Calcul et assignation numérique
    bc.profile = get_string(block, "profile", "UNIFORM");
    // Calcul et assignation numérique
    bc.u_mean = get_double(block, "u_mean", 0.0);
    // Calcul et assignation numérique
    bc.y_step = get_double(block, "y_step", 0.0);
    // Calcul et assignation de la largeur d'injecteur (Cas 3.3 Fluids 2021)
    bc.jet_width = get_double(block, "jet_width", bc.jet_width);
    // Calcul et assignation de la position centrale de l'injecteur
    bc.jet_center = get_double(block, "jet_center", bc.jet_center);
    // Calcul et assignation numérique
    bc.T_type = parse_bc_type(get_string(block, "T_type", "NEUMANN"));
    // Calcul et assignation numérique
    bc.T_val = get_double(block, "T", 0.0);
}

// Déclaration de variable typée
bool Config::load_from_json(const std::string& filepath) {
    std::ifstream file(filepath);
    // Évaluation de la condition de contrôle
    if (!file.is_open()) {
        std::cerr << "[Config] Erreur: Impossible d'ouvrir le fichier JSON : " << filepath << std::endl;
        // Retour de fonction
        return false;
    }
    std::stringstream buffer;
    buffer << file.rdbuf();
    // Déclaration de variable typée
    std::string json = buffer.str();

    // 0. Mode global et arrêt stationnaire
    mode = get_string(json, "mode", mode);
    // Calcul de la résolution géométrique de cellule
    steady_state_tol = get_double(json, "steady_state_tol", steady_state_tol);

    // 1. Maillage
    std::string mesh_block = extract_value(json, "mesh");
    // Évaluation de la condition de contrôle
    if (!mesh_block.empty()) {
        // Calcul et assignation numérique
        nx = get_int(mesh_block, "nx", nx);
        // Calcul et assignation numérique
        ny = get_int(mesh_block, "ny", ny);
        // Calcul et assignation numérique
        lx = get_double(mesh_block, "lx", lx);
        // Calcul et assignation numérique
        ly = get_double(mesh_block, "ly", ly);
    }

    // 2. Modèle physique
    std::string model_block = extract_value(json, "physics");
    // Évaluation de la condition de contrôle
    if (!model_block.empty()) {
        // Calcul et assignation numérique
        enable_two_phase = get_bool(model_block, "enable_two_phase", enable_two_phase);
        // Calcul et assignation numérique
        enable_surface_tension = get_bool(model_block, "enable_surface_tension", enable_surface_tension);
        // Calcul et assignation numérique
        enable_convection = get_bool(model_block, "enable_convection", enable_convection);
        // Calcul et assignation numérique
        enable_energy = get_bool(model_block, "enable_energy", enable_energy);
        // Calcul et assignation numérique
        sigma = get_double(model_block, "sigma", sigma);
        // Calcul et assignation numérique
        gx = get_double(model_block, "gx", gx);
        // Calcul et assignation numérique
        gy = get_double(model_block, "gy", gy);

        // Calcul et assignation numérique
        Ra = get_double(model_block, "Ra", Ra);
        // Calcul et assignation numérique
        Pr = get_double(model_block, "Pr", Pr);
        // Calcul et assignation numérique
        T_ref = get_double(model_block, "T_ref", T_ref);
        // Calcul et assignation numérique
        beta = get_double(model_block, "beta", beta);

        // Calcul et assignation numérique
        viscosity_model = get_string(model_block, "viscosity_model", viscosity_model);
        // Calcul et assignation numérique
        visc_a = get_double(model_block, "visc_a", visc_a);
        // Calcul et assignation numérique
        visc_e = get_double(model_block, "visc_e", visc_e);
        // Calcul et assignation numérique
        visc_mu0 = get_double(model_block, "visc_mu0", visc_mu0);

        // Déclaration de variable typée
        std::string f1_block = extract_value(model_block, "fluid1");
        // Évaluation de la condition de contrôle
        if (!f1_block.empty()) {
            // Calcul et assignation numérique
            fluid1.rho = get_double(f1_block, "rho", fluid1.rho);
            // Calcul et assignation numérique
            fluid1.mu  = get_double(f1_block, "mu", fluid1.mu);
            // Calcul et assignation numérique
            fluid1.Cp  = get_double(f1_block, "Cp", fluid1.Cp);
            // Calcul et assignation numérique
            fluid1.k_thermal = get_double(f1_block, "k_thermal", fluid1.k_thermal);
            // Calcul et assignation numérique
            fluid1.beta = get_double(f1_block, "beta", fluid1.beta);
        }

        // Déclaration de variable typée
        std::string f2_block = extract_value(model_block, "fluid2");
        // Évaluation de la condition de contrôle
        if (!f2_block.empty()) {
            // Calcul et assignation numérique
            fluid2.rho = get_double(f2_block, "rho", fluid2.rho);
            // Calcul et assignation numérique
            fluid2.mu  = get_double(f2_block, "mu", fluid2.mu);
            // Calcul et assignation numérique
            fluid2.Cp  = get_double(f2_block, "Cp", fluid2.Cp);
            // Calcul et assignation numérique
            fluid2.k_thermal = get_double(f2_block, "k_thermal", fluid2.k_thermal);
            // Calcul et assignation numérique
            fluid2.beta = get_double(f2_block, "beta", fluid2.beta);
        }
    }

    // 3. Intégration temporelle
    std::string time_block = extract_value(json, "time");
    // Évaluation de la condition de contrôle
    if (!time_block.empty()) {
        // Calcul et assignation numérique
        dt = get_double(time_block, "dt", dt);
        // Calcul et assignation numérique
        t_final = get_double(time_block, "t_final", t_final);
        // Calcul et assignation numérique
        cfl_max = get_double(time_block, "cfl_max", cfl_max);
        // Calcul et assignation numérique
        adaptive_dt = get_bool(time_block, "adaptive_dt", adaptive_dt);
    }

    // 4. Conditions aux limites
    std::string bc_block = extract_value(json, "boundary_conditions");
    // Évaluation de la condition de contrôle
    if (!bc_block.empty()) {
        // Déclaration de variable typée
        std::string b_left = extract_value(bc_block, "left");
        // Évaluation de la condition de contrôle
        if (!b_left.empty()) parse_single_bc(b_left, bc.left);

        // Déclaration de variable typée
        std::string b_right = extract_value(bc_block, "right");
        // Évaluation de la condition de contrôle
        if (!b_right.empty()) parse_single_bc(b_right, bc.right);

        // Déclaration de variable typée
        std::string b_bot = extract_value(bc_block, "bottom");
        // Évaluation de la condition de contrôle
        if (!b_bot.empty()) parse_single_bc(b_bot, bc.bottom);

        // Déclaration de variable typée
        std::string b_top = extract_value(bc_block, "top");
        // Évaluation de la condition de contrôle
        if (!b_top.empty()) parse_single_bc(b_top, bc.top);
    }

    // 5. Condition initiale
    std::string ic_block = extract_value(json, "initial_condition");
    // Évaluation de la condition de contrôle
    if (!ic_block.empty()) {
        // Déclaration de variable typée
        std::string type_str = get_string(ic_block, "type", "SINGLE_PHASE");
        // Évaluation de la condition de contrôle
        if (type_str == "DROPLET") init_cond.type = InitialConditionType::DROPLET;
        // Branche alternative conditionnelle pour bulle ou cylindre immergé
        else if (type_str == "BUBBLE" || type_str == "CYLINDER") init_cond.type = InitialConditionType::BUBBLE;
        // Branche alternative conditionnelle
        else if (type_str == "DAM_BREAK") init_cond.type = InitialConditionType::DAM_BREAK;
        // Branche alternative conditionnelle pour nappe liquide (Cas 3.2 Fluids 2021)
        else if (type_str == "LIQUID_SHEET") init_cond.type = InitialConditionType::LIQUID_SHEET;
        // Branche alternative conditionnelle pour jet liquide visqueux (Cas 3.3 Fluids 2021)
        else if (type_str == "VISCOUS_JET" || type_str == "JET_BUCKLING") init_cond.type = InitialConditionType::VISCOUS_JET;
        // Branche alternative conditionnelle
        else if (type_str == "SHEAR_LAYER") init_cond.type = InitialConditionType::SHEAR_LAYER;
        // Branche alternative conditionnelle
        else if (type_str == "VORTEX_SHEAR") init_cond.type = InitialConditionType::VORTEX_SHEAR;
        // Branche alternative conditionnelle
        else if (type_str == "RAYLEIGH_TAYLOR") init_cond.type = InitialConditionType::RAYLEIGH_TAYLOR;
        // Branche par défaut
        else init_cond.type = InitialConditionType::SINGLE_PHASE;

        // Calcul et assignation numérique
        init_cond.center_x = get_double(ic_block, "center_x", init_cond.center_x);
        // Calcul et assignation numérique
        init_cond.center_y = get_double(ic_block, "center_y", init_cond.center_y);
        // Calcul et assignation numérique
        init_cond.radius   = get_double(ic_block, "radius", init_cond.radius);
        // Calcul et assignation numérique
        init_cond.box_x0   = get_double(ic_block, "box_x0", init_cond.box_x0);
        // Calcul et assignation numérique
        init_cond.box_y0   = get_double(ic_block, "box_y0", init_cond.box_y0);
        // Calcul et assignation numérique
        init_cond.box_x1   = get_double(ic_block, "box_x1", init_cond.box_x1);
        // Calcul et assignation numérique
        init_cond.box_y1   = get_double(ic_block, "box_y1", init_cond.box_y1);
        // Calcul et assignation numérique
        init_cond.pert_amp = get_double(ic_block, "pert_amp", init_cond.pert_amp);
        // Calcul et assignation numérique
        init_cond.pert_k   = get_double(ic_block, "pert_k", init_cond.pert_k);
        // Calcul et assignation de l'épaisseur de nappe liquide
        init_cond.thickness = get_double(ic_block, "thickness", init_cond.thickness);
        // Calcul et assignation de la largeur du jet liquide
        init_cond.jet_width = get_double(ic_block, "jet_width", init_cond.jet_width);
    }

    // 6. Solveur linéaire
    std::string sol_block = extract_value(json, "solver");
    // Évaluation de la condition de contrôle
    if (!sol_block.empty()) {
        // Calcul et assignation numérique
        solver.type = get_string(sol_block, "type", solver.type);
        // Calcul et assignation numérique
        solver.tol = get_double(sol_block, "tol", solver.tol);
        // Calcul et assignation numérique
        solver.max_iter = get_int(sol_block, "max_iter", solver.max_iter);
        // Calcul et assignation numérique
        solver.r_al = get_double(sol_block, "r_al", solver.r_al);
        // Calcul et assignation numérique
        solver.uzawa_iter = get_int(sol_block, "uzawa_iter", solver.uzawa_iter);
        // Calcul et assignation numérique
        solver.use_precond = get_bool(sol_block, "use_precond", solver.use_precond);
        // Calcul et assignation numérique
        solver.precond_type = get_string(sol_block, "precond_type", solver.precond_type);
    }

    // 7. Sorties et visualisation
    std::string out_block = extract_value(json, "output");
    // Évaluation de la condition de contrôle
    if (!out_block.empty()) {
        // Calcul et assignation numérique
        output.output_dir = get_string(out_block, "output_dir", output.output_dir);
        // Calcul et assignation numérique
        output.case_name  = get_string(out_block, "case_name", output.case_name);
        // Calcul et assignation numérique
        output.vtk_interval = get_int(out_block, "vtk_interval", output.vtk_interval);
        // Calcul et assignation numérique
        output.log_interval = get_int(out_block, "log_interval", output.log_interval);
        // Calcul et assignation numérique
        output.write_csv_history = get_bool(out_block, "write_csv_history", output.write_csv_history);
    }

    // 8. Capture d'interface et transport VOF
    std::string vof_block = extract_value(json, "vof");
    // Évaluation de la condition de contrôle
    if (!vof_block.empty()) {
        // Calcul et assignation numérique
        vof.advection_method      = get_string(vof_block, "advection_method", vof.advection_method);
        // Calcul et assignation numérique
        vof.reconstruction_method = get_string(vof_block, "reconstruction_method", vof.reconstruction_method);
        // Calcul et assignation numérique
        vof.tvd_limiter           = get_string(vof_block, "tvd_limiter", vof.tvd_limiter);
        // Calcul et assignation numérique
        vof.cfl_vof               = get_double(vof_block, "cfl_vof", vof.cfl_vof);
    }
    // Cohérence stricte : le schéma algébrique TVD ne s'accompagne d'aucune reconstruction géométrique,
    // tandis que les schémas géométriques (WEYMOUTH_YUE, CIAM) n'utilisent pas de limiteur TVD.
    if (vof.advection_method == "VOF_TVD" || vof.advection_method == "TVD") {
        vof.reconstruction_method = "NONE";
        if (vof.tvd_limiter == "NONE" || vof.tvd_limiter.empty()) {
            vof.tvd_limiter = "SUPERBEE";
        }
    } else {
        if (vof.reconstruction_method == "NONE" || vof.reconstruction_method.empty()) {
            vof.reconstruction_method = "PLIC";
        }
        vof.tvd_limiter = "NONE";
    }

    // 9. Obstacles immergés (Pénalisation de Brinkman)
    obstacles.clear();
    // Déclaration de variable typée
    std::string obs_array = extract_value(json, "obstacles");
    // Évaluation de la condition de contrôle
    if (!obs_array.empty()) {
        // Déclaration de variable typée
        size_t pos = 0;
        // Boucle itérative conditionnelle
        while ((pos = obs_array.find('{', pos)) != std::string::npos) {
            // Déclaration de variable typée
            size_t end_pos = obs_array.find('}', pos);
            // Évaluation de la condition de contrôle
            if (end_pos == std::string::npos) break;
            // Déclaration de variable typée
            std::string obs_block = obs_array.substr(pos, end_pos - pos + 1);
            Obstacle obs;
            // Calcul et assignation numérique
            obs.type = get_string(obs_block, "type", "BOX");
            // Calcul et assignation numérique
            obs.x0 = get_double(obs_block, "x0", 0.0);
            // Calcul et assignation numérique
            obs.y0 = get_double(obs_block, "y0", 0.0);
            // Calcul et assignation numérique
            obs.x1 = get_double(obs_block, "x1", 0.0);
            // Calcul et assignation numérique
            obs.y1 = get_double(obs_block, "y1", 0.0);
            // Calcul et assignation numérique
            obs.xc = get_double(obs_block, "xc", 0.0);
            // Calcul et assignation numérique
            obs.yc = get_double(obs_block, "yc", 0.0);
            // Calcul et assignation numérique
            obs.radius = get_double(obs_block, "radius", 0.0);
            // Calcul et assignation numérique
            obs.permeability = get_double(obs_block, "permeability", 1.0e-14);
            obstacles.push_back(obs);
            // Calcul et assignation numérique
            pos = end_pos + 1;
        }
    }

    // Retour de fonction
    return true;
}

void Config::print() const {
    // Calcul et assignation numérique
    std::cout << "========================================================\n"
              << "       COVOF 2D PARALLELE - CONFIGURATION DU CAS         \n"
              // Calcul et assignation numérique
              << "========================================================\n"
              << " Mode calcul : " << mode << "\n"
              // Calcul et assignation numérique
              << " Maillage : " << nx << " x " << ny << " cellules (Lx = " << lx << " m, Ly = " << ly << " m)\n"
              // Calcul de la résolution géométrique de cellule
              << " dx = " << lx / nx << " m, dy = " << ly / ny << " m\n"
              << " Physique : " << (enable_two_phase ? "Diphasique" : "Monophasique")
              << " | Convection : " << (enable_convection ? "Oui" : "Non")
              << " | Tension de surface : " << (enable_surface_tension ? "Oui" : "Non") << "\n";
    // Évaluation de la condition de contrôle
    if (!obstacles.empty()) {
        std::cout << " Obstacles immerges : " << obstacles.size() << " defini(s) [Brinkman]\n";
        // Boucle d'itération algorithmique
        for (size_t i = 0; i < obstacles.size(); ++i) {
            std::cout << "   - Obstacle " << i+1 << " : [" << obstacles[i].x0 << ", " << obstacles[i].x1
                      // Calcul et assignation numérique
                      << "] x [" << obstacles[i].y0 << ", " << obstacles[i].y1 << "], K=" << obstacles[i].permeability << "\n";
        }
    }
    // Calcul et assignation numérique
    std::cout << " Fluide 1 : rho=" << fluid1.rho << " kg/m3, mu=" << fluid1.mu << " Pa.s\n";
    // Évaluation de la condition de contrôle
    if (enable_two_phase) {
        // Calcul et assignation numérique
        std::cout << " Fluide 2 : rho=" << fluid2.rho << " kg/m3, mu=" << fluid2.mu << " Pa.s\n"
                  // Calcul et assignation numérique
                  << " Tension superficielle sigma = " << sigma << " N/m\n";
        if (vof.advection_method == "VOF_TVD" || vof.advection_method == "TVD") {
            std::cout << " VOF : Advection=" << vof.advection_method
                      << " (Algebrique, sans reconstruction) | Limiteur=" << vof.tvd_limiter << "\n";
        } else {
            std::cout << " VOF : Advection=" << vof.advection_method
                      << " | Reconstruction=" << vof.reconstruction_method << "\n";
        }
    }
    // Calcul et assignation numérique
    std::cout << " Gravite : g = (" << gx << ", " << gy << ") m/s^2\n"
              // Calcul et assignation numérique
              << " Temps : dt = " << dt << " s, t_final = " << t_final << " s\n"
              // Calcul et assignation numérique
              << " Solveur : " << solver.type << " | BiCGSTAB(2), tol=" << solver.tol
              // Calcul et assignation numérique
              << ", max_iter=" << solver.max_iter << ", r_al=" << solver.r_al
              // Calcul et assignation numérique
              << ", uzawa_iter=" << solver.uzawa_iter << "\n"
              // Calcul et assignation numérique
              << " Sorties : Dossier '" << output.output_dir << "', intervalle VTK=" << output.vtk_interval << "\n"
              // Calcul et assignation numérique
              << "========================================================\n";
}

} // namespace covof

\end{lstlisting}

\chapter{Topologie du Maillage Cartésien MAC 2D et Décomposition Spatiale (Mesh.hpp, Mesh.cpp)}
\section{Géométrie de la Grille Décalée MAC (\emph{Marker-And-Cell})}
La classe \texttt{covof::Mesh} (\texttt{include/covof/Mesh.hpp}, \texttt{src/Mesh.cpp}) encapsule la géométrie discrète du sous-domaine attribué au processus MPI courant. Soit un domaine global $[0, L_x] \times [0, L_y]$ discrétisé par $N_{x,\text{glob}} \times N_{y,\text{glob}}$ cellules uniformes de pas :
\begin{equation}
  \Delta x = \frac{L_x}{N_{x,\text{glob}}}, \qquad \Delta y = \frac{L_y}{N_{y,\text{glob}}}
\end{equation}
Lorsque le domaine est partitionné sur une grille cartésienne de $P_x \times P_y$ processus MPI, le processus de coordonnées cartésiennes $(p_x, p_y) \in \{0, \dots, P_x-1\} \times \{0, \dots, P_y-1\}$ reçoit un sous-domaine local de $N_x \times N_y$ cellules débutant aux indices globaux $(i_{\text{start}}, j_{\text{start}})$ :
\begin{equation}
  N_x = \left\lfloor \frac{N_{x,\text{glob}}}{P_x} \right\rfloor + \mathbb{I}(p_x < N_{x,\text{glob}} \bmod P_x)
\end{equation}
Ce partitionnement à reste distribué garantit un déséquilibre de charge maximal d'au plus une colonne ou une ligne de mailles entre deux processus quelconques.

\section{Systèmes de Coordonnées Décalées $(x_c, y_c)$, $(x_u, y_c)$ et $(x_c, y_v)$}
En indiçage local C++17 à base zéro ($i \in \{0, \dots, N_x-1\}$, $j \in \{0, \dots, N_y-1\}$), les trois familles de points de collocation de la grille MAC sont évaluées par les méthodes en ligne (\texttt{inline}) de \texttt{Mesh} :
\begin{enumerate}
  \item \textbf{Centre de la maille $(i, j)$} (support de $p, T, C, \rho, \mu, \nabla\cdot\mathbf{u}$) :
  \begin{equation}
    x_c(i) = x_{\min,\text{loc}} + \left(i + \frac{1}{2}\right)\Delta x, \qquad y_c(j) = y_{\min,\text{loc}} + \left(j + \frac{1}{2}\right)\Delta y
  \end{equation}
  \item \textbf{Face verticale $(i, j)$} (support de la vitesse horizontale $u_{i,j}$ située sur la face Ouest de la maille $(i,j)$, avec $i \in \{0, \dots, N_x\}$) :
  \begin{equation}
    x_u(i) = x_{\min,\text{loc}} + i\,\Delta x, \qquad y_u(j) = y_c(j)
  \end{equation}
  \item \textbf{Face horizontale $(i, j)$} (support de la vitesse verticale $v_{i,j}$ située sur la face Sud de la maille $(i,j)$, avec $j \in \{0, \dots, N_y\}$) :
  \begin{equation}
    x_v(i) = x_c(i), \qquad y_v(j) = y_{\min,\text{loc}} + j\,\Delta y
  \end{equation}
\end{enumerate}

\section{Répertoire des Classes, Structures et Méthodes de \texttt{include/covof/Mesh.hpp}}
\begin{center}\small
\begin{tabular}{clll}
\toprule
\textbf{Ligne} & \textbf{Catégorie} & \textbf{Identifiant C++17} & \textbf{Signature / Rôle} \\
\midrule
12 & class & \texttt{Mesh2D} & Déclaration de type / classe / structure \\
80 & Méthode / Fonction & \texttt{cell\_idx} & Type retour : inline int \\
85 & Méthode / Fonction & \texttt{u\_idx} & Type retour : inline int \\
90 & Méthode / Fonction & \texttt{v\_idx} & Type retour : inline int \\
95 & Méthode / Fonction & \texttt{node\_idx} & Type retour : inline int \\
\bottomrule
\end{tabular}
\end{center}

\section{Code Source Intégral Annoté : \texttt{include/covof/Mesh.hpp}}
Le listing ci-dessous présente l'intégralité des 109 lignes du fichier \texttt{include/covof/Mesh.hpp} de l'implémentation C++17 :

\begin{lstlisting}[language={[11]C++}]
#pragma once

#include <mpi.h>
#include <vector>

namespace covof {
namespace fugu = covof; // Alias rétrocompatibilité

/**
 * @brief Maillage 2D cartésien décalé (MAC) partitionné en sous-domaines MPI
 */
class Mesh2D {
public:
    // Paramètres globaux
    int nx_global = 0;
    // Déclaration de variable typée
    int ny_global = 0;
    // Déclaration de variable typée
    double lx = 1.0;
    // Déclaration de variable typée
    double ly = 1.0;
    // Déclaration de variable typée
    double dx = 0.0;
    // Déclaration de variable typée
    double dy = 0.0;

    // Décomposition de domaine MPI
    MPI_Comm comm = MPI_COMM_NULL;
    // Déclaration de variable typée
    int rank = 0;
    // Déclaration de variable typée
    int num_procs = 1;
    // Déclaration de variable typée
    int dims[2] = {1, 1};       ///< Nombre de processeurs {px, py}
    // Déclaration de variable typée
    int coords[2] = {0, 0};     ///< Coordonnées du processus dans la topologie cartésienne

    // Voisins MPI
    int rank_left   = MPI_PROC_NULL;
    // Déclaration de variable typée
    int rank_right  = MPI_PROC_NULL;
    // Déclaration de variable typée
    int rank_bottom = MPI_PROC_NULL;
    // Déclaration de variable typée
    int rank_top    = MPI_PROC_NULL;

    // Indicateurs de frontière physique globale
    bool is_boundary_left   = false;
    // Déclaration de variable typée
    bool is_boundary_right  = false;
    // Déclaration de variable typée
    bool is_boundary_bottom = false;
    // Déclaration de variable typée
    bool is_boundary_top    = false;

    // Dimensions locales (sans mailles fantômes)
    int nx_local = 0;
    // Déclaration de variable typée
    int ny_local = 0;
    // Déclaration de variable typée
    int i_offset = 0; ///< Décalage d'indice global selon X
    // Déclaration de variable typée
    int j_offset = 0; ///< Décalage d'indice global selon Y

    // Dimensions allouées avec mailles fantômes (1 couche de chaque côté)
    int nx_ghost = 0; ///< nx_local + 2
    // Déclaration de variable typée
    int ny_ghost = 0; ///< ny_local + 2

    /**
     * @brief Initialise le maillage et la topologie cartésienne MPI
     */
    void init(MPI_Comm parent_comm, int nx_g, int ny_g, double length_x, double length_y,
              // Déclaration de variable typée
              bool periodic_x = false, bool periodic_y = false);

    /**
     * @brief Fonctions d'indexation 1D locale avec mailles fantômes
     */
    inline int cell_idx(int i, int j) const {
        // Retour de fonction
        return j * nx_ghost + i;
    }

    inline int u_idx(int i, int j) const {
        // Retour de fonction
        return j * (nx_ghost + 1) + i;
    }

    inline int v_idx(int i, int j) const {
        // Retour de fonction
        return j * nx_ghost + i;
    }

    inline int node_idx(int i, int j) const {
        // Retour de fonction
        return j * (nx_ghost + 1) + i;
    }

    // Coordonnées physiques locales
    inline double xc(int i) const { return (i_offset + i - 0.5) * dx; }
    inline double yc(int j) const { return (j_offset + j - 0.5) * dy; }
    inline double xu(int i) const { return (i_offset + i - 1.0) * dx; }
    inline double xv(int i) const { return xu(i); } ///< Coordonnée X de la face verticale / sommet
    inline double yv(int j) const { return (j_offset + j - 1.0) * dy; }
    inline double yu(int j) const { return yv(j); } ///< Coordonnée Y de la face horizontale / sommet
};

} // namespace covof

\end{lstlisting}

\section{Code Source Intégral Annoté : \texttt{src/Mesh.cpp}}
Le listing ci-dessous présente l'intégralité des 94 lignes du fichier \texttt{src/Mesh.cpp} de l'implémentation C++17 :

\begin{lstlisting}[language={[11]C++}]
#include "covof/Mesh.hpp"
#include <iostream>
#include <cassert>

namespace covof {
namespace fugu = covof; // Alias rétrocompatibilité

void Mesh2D::init(MPI_Comm parent_comm, int nx_g, int ny_g, double length_x, double length_y,
                  // Déclaration de variable typée
                  bool periodic_x, bool periodic_y) {
    // Calcul et assignation numérique
    nx_global = nx_g;
    // Calcul et assignation numérique
    ny_global = ny_g;
    // Calcul et assignation numérique
    lx = length_x;
    // Calcul et assignation numérique
    ly = length_y;
    // Calcul de la résolution géométrique de cellule
    dx = lx / nx_global;
    // Calcul de la résolution géométrique de cellule
    dy = ly / ny_global;

    // Déclaration de variable typée
    int parent_size, parent_rank;
    // Appel de communication ou synchronisation parallèle OpenMPI
    MPI_Comm_size(parent_comm, &parent_size);
    // Appel de communication ou synchronisation parallèle OpenMPI
    MPI_Comm_rank(parent_comm, &parent_rank);

    // 1. Détermination de la topologie cartésienne 2D optimale
    dims[0] = 0;
    // Calcul et assignation numérique
    dims[1] = 0;
    // Appel de communication ou synchronisation parallèle OpenMPI
    MPI_Dims_create(parent_size, 2, dims);

    // Déclaration de variable typée
    int periods[2] = {periodic_x ? 1 : 0, periodic_y ? 1 : 0};
    // Déclaration de variable typée
    int reorder = 1;
    // Appel de communication ou synchronisation parallèle OpenMPI
    MPI_Cart_create(parent_comm, 2, dims, periods, reorder, &comm);
    // Appel de communication ou synchronisation parallèle OpenMPI
    MPI_Comm_rank(comm, &rank);
    // Appel de communication ou synchronisation parallèle OpenMPI
    MPI_Comm_size(comm, &num_procs);
    // Appel de communication ou synchronisation parallèle OpenMPI
    MPI_Cart_coords(comm, rank, 2, coords);

    // 2. Découverte des voisins MPI
    MPI_Cart_shift(comm, 0, 1, &rank_left, &rank_right);
    // Appel de communication ou synchronisation parallèle OpenMPI
    MPI_Cart_shift(comm, 1, 1, &rank_bottom, &rank_top);

    // Calcul et assignation numérique
    is_boundary_left   = (!periodic_x && (coords[0] == 0));
    // Calcul et assignation numérique
    is_boundary_right  = (!periodic_x && (coords[0] == dims[0] - 1));
    // Calcul et assignation numérique
    is_boundary_bottom = (!periodic_y && (coords[1] == 0));
    // Calcul et assignation numérique
    is_boundary_top    = (!periodic_y && (coords[1] == dims[1] - 1));

    // 3. Répartition équitable des mailles locales
    auto partition_dim = [](int n_global, int n_procs, int coord, int& n_local, int& offset) {
        // Déclaration de variable typée
        int base = n_global / n_procs;
        // Déclaration de variable typée
        int rem  = n_global % n_procs;
        // Évaluation de la condition de contrôle
        if (coord < rem) {
            // Calcul et assignation numérique
            n_local = base + 1;
            // Calcul et assignation numérique
            offset  = coord * (base + 1);
        } else {
            // Calcul et assignation numérique
            n_local = base;
            // Calcul et assignation numérique
            offset  = rem * (base + 1) + (coord - rem) * base;
        }
    };

    partition_dim(nx_global, dims[0], coords[0], nx_local, i_offset);
    partition_dim(ny_global, dims[1], coords[1], ny_local, j_offset);

    // Calcul et assignation numérique
    nx_ghost = nx_local + 2;
    // Calcul et assignation numérique
    ny_ghost = ny_local + 2;
}

} // namespace covof

\end{lstlisting}

\chapter{Conteneurs de Champs Eulériens 2D et Opérations Vectorielles (Field.hpp)}
\section{Conception du Template \texttt{Field2D<T>}}
L'en-tête \texttt{include/covof/Field.hpp} fournit la classe générique \texttt{Field2D<T>} (instanciée principalement pour \texttt{T = double} via l'alias \texttt{ScalarField}) qui représente tout champ discret défini sur le maillage MAC (aussi bien un champ centré de dimension $N_x \times N_y$ qu'un champ de vitesse décalé $u$ de dimension $(N_x+1) \times N_y$ ou $v$ de dimension $N_x \times (N_y+1)$).

\subsection{Adressage Avec Cellules Fantômes (\emph{Ghost Cells})}
Chaque instance de \texttt{Field2D<T>} stocke :
\begin{itemize}
  \item \texttt{nx\_}, \texttt{ny\_} : les dimensions actives du champ à l'intérieur du sous-domaine ;
  \item \texttt{ng\_} : l'épaisseur de la couronne de mailles fantômes (typiquement $n_g = 1$ ou $n_g = 2$ pour les stencils VOF $3\times 3$ et TVD) ;
  \item \texttt{stride\_} $= N_x + 2 n_g$ : le pas d'adressage mémoire d'une ligne complète ;
  \item \texttt{data\_} : le vecteur contigu de taille $(N_x + 2n_g)(N_y + 2n_g)$ initialisé à zéro.
\end{itemize}
L'opérateur d'accès \texttt{operator()(int i, int j)} accepte des indices relatifs signés $i \in [-n_g, N_x + n_g - 1]$ et $j \in [-n_g, N_y + n_g - 1]$ :
\begin{equation}
  \texttt{field(i, j)} \equiv \texttt{data\_[(j + ng\_) * stride\_ + (i + ng\_)]}
\end{equation}
Ainsi, l'écriture des opérateurs aux différences finies dans \texttt{NavierStokesSolver.cpp} et \texttt{InterfaceTracker.cpp} conserve une lisibilité mathématique directe (\texttt{u(i+1, j) - 2*u(i, j) + u(i-1, j)}) sans aucune arithmétique d'offset dans le code client.

\section{Répertoire des Classes, Structures et Méthodes de \texttt{include/covof/Field.hpp}}
\begin{center}\small
\begin{tabular}{clll}
\toprule
\textbf{Ligne} & \textbf{Catégorie} & \textbf{Identifiant C++17} & \textbf{Signature / Rôle} \\
\midrule
13 & class & \texttt{FlowFields} & Déclaration de type / classe / structure \\
57 & Méthode / Fonction & \texttt{allocate} & Type retour : void \\
\bottomrule
\end{tabular}
\end{center}

\section{Code Source Intégral Annoté : \texttt{include/covof/Field.hpp}}
Le listing ci-dessous présente l'intégralité des 89 lignes du fichier \texttt{include/covof/Field.hpp} de l'implémentation C++17 :

\begin{lstlisting}[language={[11]C++}]
#pragma once

#include "Mesh.hpp"
#include <vector>

namespace covof {
namespace fugu = covof; // Alias rétrocompatibilité

/**
 * @brief Champs scalaires et vectoriels sur la grille décalée MAC
 * Tous les champs contiennent une couche de mailles fantômes pour le parallélisme MPI.
 */
class FlowFields {
public:
    // Champs cinématiques et thermodynamiques principaux
    std::vector<double> u;  ///< Vitesse horizontale sur les faces verticales
    // Déclaration de variable typée
    std::vector<double> v;  ///< Vitesse verticale sur les faces horizontales
    // Déclaration de variable typée
    std::vector<double> p;  ///< Pression scalaire au centre des cellules
    // Déclaration de variable typée
    std::vector<double> C;  ///< Fonction indicatrice de phase au centre des cellules (C in [0, 1])
    // Déclaration de variable typée
    std::vector<double> T;  ///< Température scalaire au centre des cellules [K ou adimensionnelle]

    // Propriétés physiques locales effectives (Lois de mélange)
    std::vector<double> rho;     ///< Masse volumique au centre des mailles
    // Déclaration de variable typée
    std::vector<double> rho_u;   ///< Masse volumique interpolée sur les faces u
    // Déclaration de variable typée
    std::vector<double> rho_v;   ///< Masse volumique interpolée sur les faces v
    // Déclaration de variable typée
    std::vector<double> mu;      ///< Viscosité dynamique au centre des mailles
    // Déclaration de variable typée
    std::vector<double> mu_node; ///< Viscosité dynamique harmonique aux nœuds (coins de maille)

    // Géométrie de l'interface et tension superficielle CSF
    std::vector<double> kappa;     ///< Courbure de l'interface au centre des cellules
    // Déclaration de variable typée
    std::vector<double> norm_x;    ///< Composante X du vecteur normal
    // Déclaration de variable typée
    std::vector<double> norm_y;    ///< Composante Y du vecteur normal
    // Déclaration de variable typée
    std::vector<double> F_sigma_x; ///< Force capillaire CSF selon X sur les faces u
    // Déclaration de variable typée
    std::vector<double> F_sigma_y; ///< Force capillaire CSF selon Y sur les faces v

    // Diagnostics
    std::vector<double> div_u;     ///< Divergence de la vitesse au centre des mailles

    // Calcul et assignation numérique
    FlowFields() = default;

    /**
     * @brief Alloue la mémoire selon les dimensions locales du maillage avec mailles fantômes
     */
    void allocate(const Mesh2D& mesh) {
        // Déclaration de variable typée
        int n_cells = mesh.nx_ghost * mesh.ny_ghost;
        // Déclaration de variable typée
        int n_u = (mesh.nx_ghost + 1) * mesh.ny_ghost;
        // Déclaration de variable typée
        int n_v = mesh.nx_ghost * (mesh.ny_ghost + 1);
        // Déclaration de variable typée
        int n_nodes = (mesh.nx_ghost + 1) * (mesh.ny_ghost + 1);

        u.assign(n_u, 0.0);
        v.assign(n_v, 0.0);
        p.assign(n_cells, 0.0);
        C.assign(n_cells, 0.0);
        T.assign(n_cells, 0.0);

        rho.assign(n_cells, 1.0);
        rho_u.assign(n_u, 1.0);
        rho_v.assign(n_v, 1.0);
        mu.assign(n_cells, 1.0);
        mu_node.assign(n_nodes, 1.0);

        kappa.assign(n_cells, 0.0);
        norm_x.assign(n_cells, 0.0);
        norm_y.assign(n_cells, 0.0);
        F_sigma_x.assign(n_u, 0.0);
        F_sigma_y.assign(n_v, 0.0);

        div_u.assign(n_cells, 0.0);
    }
};

} // namespace covof

\end{lstlisting}

\chapter{Parallélisme Distribué OpenMPI et Échanges de Halos Cartésiens (MpiComm.hpp, MpiComm.cpp)}
\section{Topologie Cartésienne 2D et Découverte des Voisins}
La classe \texttt{covof::MpiComm} (\texttt{include/covof/MpiComm.hpp}, \texttt{src/MpiComm.cpp}) abstrait toutes les communications inter-processus OpenMPI. Lors de sa construction :
\begin{enumerate}
  \item La fonction \texttt{MPI\_Dims\_create(size, 2, dims)} décompose automatiquement le nombre total de processus $P$ en une grille bidimensionnelle $P_x \times P_y$ minimisant le rapport périmètre/surface des sous-domaines (et donc le volume de communication de halo par rapport au volume de calcul).
  \item \texttt{MPI\_Cart\_create} construit le communicateur topologique 2D en prenant en compte les drapeaux de périodicité $(\text{periodic}_x, \text{periodic}_y)$ issus de la configuration JSON.
  \item \texttt{MPI\_Cart\_shift} identifie les rangs des quatre voisins immédiats : \texttt{rank\_left} (Ouest, $x^-$), \texttt{rank\_right} (Est, $x^+$), \texttt{rank\_bottom} (Sud, $y^-$) et \texttt{rank\_top} (Nord, $y^+$). Si une frontière externe n'est pas périodique, MPI affecte automatiquement \texttt{MPI\_PROC\_NULL}, ce qui transforme les appels d'envoi/réception sur ce bord en opérations nulles sans branchement conditionnel coûteux.
\end{enumerate}

\section{Protocole d'Échange de Halos Sans Interblocage (\emph{Deadlock-Free Halo Exchange})}
La méthode \texttt{MpiComm::exchange\_halo(Field2D<double>\& field)} synchronise les couches de cellules fantômes $i \in \{-n_g, \dots, -1\}$, $i \in \{N_x, \dots, N_x+n_g-1\}$, $j \in \{-n_g, \dots, -1\}$ et $j \in \{N_y, \dots, N_y+n_g-1\}$ en deux passes séquentielles (passe horizontale Est$\leftrightarrow$Ouest suivie de la passe verticale Nord$\leftrightarrow$Sud) :
\begin{itemize}
  \item L'enchaînement séquentiel de la passe horizontale puis de la passe verticale propage automatiquement les \textbf{cellules fantômes de coin} $(i < 0, j < 0)$, $(i \ge N_x, j \ge N_y)$ sans nécessiter de messages diagonaux dédiés (crucial pour le filtre $3\times 3$ de Parker-Youngs et le terme croisé $-r_{al}\partial^2 / \partial x \partial y$).
  \item Les colonnes verticales n'étant pas contiguës en mémoire \emph{row-major}, elles sont empaquetées (\emph{packing}) dans des tampons linéaires \texttt{send\_left}, \texttt{send\_right}, transmises via \texttt{MPI\_Sendrecv} ou \texttt{MPI\_Isend}/\texttt{MPI\_Irecv}, puis dépaquetées (\emph{unpacking}) dans les cellules fantômes.
\end{itemize}

\section{Répertoire des Classes, Structures et Méthodes de \texttt{include/covof/MpiComm.hpp}}
\begin{center}\small
\begin{tabular}{clll}
\toprule
\textbf{Ligne} & \textbf{Catégorie} & \textbf{Identifiant C++17} & \textbf{Signature / Rôle} \\
\midrule
14 & class & \texttt{MpiComm} & Déclaration de type / classe / structure \\
19 & Méthode / Fonction & \texttt{sync\_cell\_scalar} & Type retour : static void \\
24 & Méthode / Fonction & \texttt{sync\_u} & Type retour : static void \\
29 & Méthode / Fonction & \texttt{sync\_v} & Type retour : static void \\
35 & Méthode / Fonction & \texttt{apply\_boundary\_conditions} & Type retour : static void \\
40 & Méthode / Fonction & \texttt{apply\_thermal\_boundary\_conditions} & Type retour : static void \\
45 & Méthode / Fonction & \texttt{global\_dot\_product} & Type retour : static double \\
50 & Méthode / Fonction & \texttt{global\_norm2} & Type retour : static double \\
55 & Méthode / Fonction & \texttt{global\_max} & Type retour : static double \\
60 & Méthode / Fonction & \texttt{global\_sum} & Type retour : static double \\
\bottomrule
\end{tabular}
\end{center}

\section{Code Source Intégral Annoté : \texttt{include/covof/MpiComm.hpp}}
Le listing ci-dessous présente l'intégralité des 63 lignes du fichier \texttt{include/covof/MpiComm.hpp} de l'implémentation C++17 :

\begin{lstlisting}[language={[11]C++}]
#pragma once

#include "Mesh.hpp"
#include "Field.hpp"
#include "Types.hpp"
#include <vector>

namespace covof {
namespace fugu = covof; // Alias rétrocompatibilité

/**
 * @brief Gestionnaire des communications MPI et des échanges de mailles fantômes (halos)
 */
class MpiComm {
public:
    /**
     * @brief Échange les halos d'un champ scalaire centré (p, C, rho, etc.)
     */
    static void sync_cell_scalar(const Mesh2D& mesh, std::vector<double>& field);

    /**
     * @brief Échange les halos de la vitesse horizontale u (faces verticales)
     */
    static void sync_u(const Mesh2D& mesh, std::vector<double>& u);

    /**
     * @brief Échange les halos de la vitesse verticale v (faces horizontales)
     */
    static void sync_v(const Mesh2D& mesh, std::vector<double>& v);

    /**
     * @brief Applique les conditions aux limites physiques sur les mailles fantômes extérieures
     * @param homogeneous Si true, applique des CL homogènes (pour l'opérateur linéaire A*x)
     */
    static void apply_boundary_conditions(const Mesh2D& mesh, const DomainBC& bc, FlowFields& fields, bool homogeneous = false);

    /**
     * @brief Applique les conditions aux limites thermiques sur les mailles fantômes de la température T
     */
    static void apply_thermal_boundary_conditions(const Mesh2D& mesh, const DomainBC& bc, FlowFields& fields);

    /**
     * @brief Calcule le produit scalaire global de deux vecteurs distribués
     */
    static double global_dot_product(const Mesh2D& mesh, const std::vector<double>& a, const std::vector<double>& b, int size);

    /**
     * @brief Calcule la norme L2 globale d'un vecteur distribué
     */
    static double global_norm2(const Mesh2D& mesh, const std::vector<double>& v, int size);

    /**
     * @brief Réduit la valeur maximale scalaire sur tous les processus MPI (ex: CFL)
     */
    static double global_max(const Mesh2D& mesh, double local_val);

    /**
     * @brief Réduit la somme globale d'un scalaire sur tous les processus MPI (ex: intégrales)
     */
    static double global_sum(const Mesh2D& mesh, double local_val);
};

} // namespace covof

\end{lstlisting}

\section{Répertoire des Classes, Structures et Méthodes de \texttt{src/MpiComm.cpp}}
\begin{center}\small
\begin{tabular}{clll}
\toprule
\textbf{Ligne} & \textbf{Catégorie} & \textbf{Identifiant C++17} & \textbf{Signature / Rôle} \\
\midrule
8 & Méthode / Fonction & \texttt{MpiComm::sync\_cell\_scalar} & Type retour : void \\
32 & Méthode / Fonction & \texttt{send\_left} & Type retour : std::vector$<$double$>$ \\
34 & Méthode / Fonction & \texttt{send\_right} & Type retour : std::vector$<$double$>$ \\
36 & Méthode / Fonction & \texttt{send\_bot} & Type retour : std::vector$<$double$>$ \\
38 & Méthode / Fonction & \texttt{send\_top} & Type retour : std::vector$<$double$>$ \\
122 & Méthode / Fonction & \texttt{MpiComm::sync\_u} & Type retour : void \\
132 & Méthode / Fonction & \texttt{send\_left} & Type retour : std::vector$<$double$>$ \\
134 & Méthode / Fonction & \texttt{send\_right} & Type retour : std::vector$<$double$>$ \\
136 & Méthode / Fonction & \texttt{send\_bot} & Type retour : std::vector$<$double$>$ \\
138 & Méthode / Fonction & \texttt{send\_top} & Type retour : std::vector$<$double$>$ \\
222 & Méthode / Fonction & \texttt{MpiComm::sync\_v} & Type retour : void \\
232 & Méthode / Fonction & \texttt{send\_left} & Type retour : std::vector$<$double$>$ \\
234 & Méthode / Fonction & \texttt{send\_right} & Type retour : std::vector$<$double$>$ \\
236 & Méthode / Fonction & \texttt{send\_bot} & Type retour : std::vector$<$double$>$ \\
238 & Méthode / Fonction & \texttt{send\_top} & Type retour : std::vector$<$double$>$ \\
322 & Méthode / Fonction & \texttt{MpiComm::apply\_boundary\_conditions} & Type retour : void \\
597 & Méthode / Fonction & \texttt{MpiComm::apply\_thermal\_boundary\_conditions} & Type retour : void \\
686 & Méthode / Fonction & \texttt{MpiComm::global\_dot\_product} & Type retour : double \\
703 & Méthode / Fonction & \texttt{MpiComm::global\_norm2} & Type retour : double \\
709 & Méthode / Fonction & \texttt{MpiComm::global\_max} & Type retour : double \\
719 & Méthode / Fonction & \texttt{MpiComm::global\_sum} & Type retour : double \\
\bottomrule
\end{tabular}
\end{center}

\section{Code Source Intégral Annoté : \texttt{src/MpiComm.cpp}}
Le listing ci-dessous présente l'intégralité des 728 lignes du fichier \texttt{src/MpiComm.cpp} de l'implémentation C++17 :

\begin{lstlisting}[language={[11]C++}]
#include "covof/MpiComm.hpp"
#include <algorithm>
#include <cmath>

namespace covof {
namespace fugu = covof; // Alias rétrocompatibilité

void MpiComm::sync_cell_scalar(const Mesh2D& mesh, std::vector<double>& field) {
    // Déclaration de variable typée
    int nx = mesh.nx_local;
    // Déclaration de variable typée
    int ny = mesh.ny_local;

    // Évaluation de la condition de contrôle
    if (mesh.num_procs == 1) {
        if (!mesh.is_boundary_left) {
            for (int j = 0; j <= ny + 1; ++j) {
                field[mesh.cell_idx(0, j)]      = field[mesh.cell_idx(nx, j)];
                field[mesh.cell_idx(nx + 1, j)] = field[mesh.cell_idx(1, j)];
            }
        }
        if (!mesh.is_boundary_bottom) {
            for (int i = 0; i <= nx + 1; ++i) {
                field[mesh.cell_idx(i, 0)]      = field[mesh.cell_idx(i, ny)];
                field[mesh.cell_idx(i, ny + 1)] = field[mesh.cell_idx(i, 1)];
            }
        }
        return;
    }

    // Déclaration de variable typée
    std::vector<double> send_left(ny), recv_left(ny);
    // Déclaration de variable typée
    std::vector<double> send_right(ny), recv_right(ny);
    // Déclaration de variable typée
    std::vector<double> send_bot(nx), recv_bot(nx);
    // Déclaration de variable typée
    std::vector<double> send_top(nx), recv_top(nx);

    // Préparation des buffers X
    for (int j = 1; j <= ny; ++j) {
        // Calcul et assignation numérique
        send_left[j - 1]  = field[mesh.cell_idx(1, j)];
        // Calcul et assignation numérique
        send_right[j - 1] = field[mesh.cell_idx(nx, j)];
    }

    // Appel de communication ou synchronisation parallèle OpenMPI
    MPI_Request reqs[8];
    // Déclaration de variable typée
    int n_reqs = 0;

    // Échanges Est-Ouest
    if (mesh.rank_left != MPI_PROC_NULL) {
        // Appel de communication ou synchronisation parallèle OpenMPI
        MPI_Irecv(recv_left.data(), ny, MPI_DOUBLE, mesh.rank_left, 10, mesh.comm, &reqs[n_reqs++]);
        // Appel de communication ou synchronisation parallèle OpenMPI
        MPI_Isend(send_left.data(), ny, MPI_DOUBLE, mesh.rank_left, 11, mesh.comm, &reqs[n_reqs++]);
    }
    // Évaluation de la condition de contrôle
    if (mesh.rank_right != MPI_PROC_NULL) {
        // Appel de communication ou synchronisation parallèle OpenMPI
        MPI_Irecv(recv_right.data(), ny, MPI_DOUBLE, mesh.rank_right, 11, mesh.comm, &reqs[n_reqs++]);
        // Appel de communication ou synchronisation parallèle OpenMPI
        MPI_Isend(send_right.data(), ny, MPI_DOUBLE, mesh.rank_right, 10, mesh.comm, &reqs[n_reqs++]);
    }

    // Appel de communication ou synchronisation parallèle OpenMPI
    MPI_Waitall(n_reqs, reqs, MPI_STATUSES_IGNORE);
    // Calcul et assignation numérique
    n_reqs = 0;

    // Évaluation de la condition de contrôle
    if (mesh.rank_left != MPI_PROC_NULL) {
        // Boucle spatiale sur les cellules locales selon l'axe vertical Y (ordonnées)
        for (int j = 1; j <= ny; ++j) field[mesh.cell_idx(0, j)] = recv_left[j - 1];
    }
    // Évaluation de la condition de contrôle
    if (mesh.rank_right != MPI_PROC_NULL) {
        // Boucle spatiale sur les cellules locales selon l'axe vertical Y (ordonnées)
        for (int j = 1; j <= ny; ++j) field[mesh.cell_idx(nx + 1, j)] = recv_right[j - 1];
    }

    // Préparation des buffers Y
    for (int i = 1; i <= nx; ++i) {
        // Calcul et assignation numérique
        send_bot[i - 1] = field[mesh.cell_idx(i, 1)];
        // Calcul et assignation numérique
        send_top[i - 1] = field[mesh.cell_idx(i, ny)];
    }

    // Échanges Nord-Sud
    if (mesh.rank_bottom != MPI_PROC_NULL) {
        // Appel de communication ou synchronisation parallèle OpenMPI
        MPI_Irecv(recv_bot.data(), nx, MPI_DOUBLE, mesh.rank_bottom, 20, mesh.comm, &reqs[n_reqs++]);
        // Appel de communication ou synchronisation parallèle OpenMPI
        MPI_Isend(send_bot.data(), nx, MPI_DOUBLE, mesh.rank_bottom, 21, mesh.comm, &reqs[n_reqs++]);
    }
    // Évaluation de la condition de contrôle
    if (mesh.rank_top != MPI_PROC_NULL) {
        // Appel de communication ou synchronisation parallèle OpenMPI
        MPI_Irecv(recv_top.data(), nx, MPI_DOUBLE, mesh.rank_top, 21, mesh.comm, &reqs[n_reqs++]);
        // Appel de communication ou synchronisation parallèle OpenMPI
        MPI_Isend(send_top.data(), nx, MPI_DOUBLE, mesh.rank_top, 20, mesh.comm, &reqs[n_reqs++]);
    }

    // Appel de communication ou synchronisation parallèle OpenMPI
    MPI_Waitall(n_reqs, reqs, MPI_STATUSES_IGNORE);

    // Évaluation de la condition de contrôle
    if (mesh.rank_bottom != MPI_PROC_NULL) {
        // Boucle spatiale sur les cellules locales selon l'axe horizontal X (abscisses)
        for (int i = 1; i <= nx; ++i) field[mesh.cell_idx(i, 0)] = recv_bot[i - 1];
    }
    // Évaluation de la condition de contrôle
    if (mesh.rank_top != MPI_PROC_NULL) {
        // Boucle spatiale sur les cellules locales selon l'axe horizontal X (abscisses)
        for (int i = 1; i <= nx; ++i) field[mesh.cell_idx(i, ny + 1)] = recv_top[i - 1];
    }
}

void MpiComm::sync_u(const Mesh2D& mesh, std::vector<double>& u) {
    // Évaluation de la condition de contrôle
    if (mesh.num_procs == 1) return;

    // Déclaration de variable typée
    int nx = mesh.nx_local;
    // Déclaration de variable typée
    int ny = mesh.ny_local;

    // u est défini sur les faces verticales i de 1 à nx+1 pour l'intérieur
    std::vector<double> send_left(ny), recv_left(ny);
    // Déclaration de variable typée
    std::vector<double> send_right(ny), recv_right(ny);
    // Déclaration de variable typée
    std::vector<double> send_bot(nx + 2), recv_bot(nx + 2);
    // Déclaration de variable typée
    std::vector<double> send_top(nx + 2), recv_top(nx + 2);

    // Boucle spatiale sur les cellules locales selon l'axe vertical Y (ordonnées)
    for (int j = 1; j <= ny; ++j) {
        // Calcul et assignation numérique
        send_left[j - 1]  = u[mesh.u_idx(1, j)];
        // Calcul et assignation numérique
        send_right[j - 1] = u[mesh.u_idx(nx, j)];
    }

    // Appel de communication ou synchronisation parallèle OpenMPI
    MPI_Request reqs[8];
    // Déclaration de variable typée
    int n_reqs = 0;

    // Évaluation de la condition de contrôle
    if (mesh.rank_left != MPI_PROC_NULL) {
        // Appel de communication ou synchronisation parallèle OpenMPI
        MPI_Irecv(recv_left.data(), ny, MPI_DOUBLE, mesh.rank_left, 30, mesh.comm, &reqs[n_reqs++]);
        // Appel de communication ou synchronisation parallèle OpenMPI
        MPI_Isend(send_left.data(), ny, MPI_DOUBLE, mesh.rank_left, 31, mesh.comm, &reqs[n_reqs++]);
    }
    // Évaluation de la condition de contrôle
    if (mesh.rank_right != MPI_PROC_NULL) {
        // Appel de communication ou synchronisation parallèle OpenMPI
        MPI_Irecv(recv_right.data(), ny, MPI_DOUBLE, mesh.rank_right, 31, mesh.comm, &reqs[n_reqs++]);
        // Appel de communication ou synchronisation parallèle OpenMPI
        MPI_Isend(send_right.data(), ny, MPI_DOUBLE, mesh.rank_right, 30, mesh.comm, &reqs[n_reqs++]);
    }

    // Appel de communication ou synchronisation parallèle OpenMPI
    MPI_Waitall(n_reqs, reqs, MPI_STATUSES_IGNORE);
    // Calcul et assignation numérique
    n_reqs = 0;

    // Évaluation de la condition de contrôle
    if (mesh.rank_left != MPI_PROC_NULL) {
        // Boucle spatiale sur les cellules locales selon l'axe vertical Y (ordonnées)
        for (int j = 1; j <= ny; ++j) u[mesh.u_idx(0, j)] = recv_left[j - 1];
    }
    // Évaluation de la condition de contrôle
    if (mesh.rank_right != MPI_PROC_NULL) {
        // Boucle spatiale sur les cellules locales selon l'axe vertical Y (ordonnées)
        for (int j = 1; j <= ny; ++j) u[mesh.u_idx(nx + 1, j)] = recv_right[j - 1];
    }

    // Échange selon Y
    for (int i = 0; i <= nx + 1; ++i) {
        // Calcul et assignation numérique
        send_bot[i] = u[mesh.u_idx(i, 1)];
        // Calcul et assignation numérique
        send_top[i] = u[mesh.u_idx(i, ny)];
    }

    // Évaluation de la condition de contrôle
    if (mesh.rank_bottom != MPI_PROC_NULL) {
        // Appel de communication ou synchronisation parallèle OpenMPI
        MPI_Irecv(recv_bot.data(), nx + 2, MPI_DOUBLE, mesh.rank_bottom, 40, mesh.comm, &reqs[n_reqs++]);
        // Appel de communication ou synchronisation parallèle OpenMPI
        MPI_Isend(send_bot.data(), nx + 2, MPI_DOUBLE, mesh.rank_bottom, 41, mesh.comm, &reqs[n_reqs++]);
    }
    // Évaluation de la condition de contrôle
    if (mesh.rank_top != MPI_PROC_NULL) {
        // Appel de communication ou synchronisation parallèle OpenMPI
        MPI_Irecv(recv_top.data(), nx + 2, MPI_DOUBLE, mesh.rank_top, 41, mesh.comm, &reqs[n_reqs++]);
        // Appel de communication ou synchronisation parallèle OpenMPI
        MPI_Isend(send_top.data(), nx + 2, MPI_DOUBLE, mesh.rank_top, 40, mesh.comm, &reqs[n_reqs++]);
    }

    // Appel de communication ou synchronisation parallèle OpenMPI
    MPI_Waitall(n_reqs, reqs, MPI_STATUSES_IGNORE);

    // Évaluation de la condition de contrôle
    if (mesh.rank_bottom != MPI_PROC_NULL) {
        // Boucle spatiale sur les cellules locales selon l'axe horizontal X (abscisses)
        for (int i = 0; i <= nx + 1; ++i) u[mesh.u_idx(i, 0)] = recv_bot[i];
    }
    // Évaluation de la condition de contrôle
    if (mesh.rank_top != MPI_PROC_NULL) {
        // Boucle spatiale sur les cellules locales selon l'axe horizontal X (abscisses)
        for (int i = 0; i <= nx + 1; ++i) u[mesh.u_idx(i, ny + 1)] = recv_top[i];
    }
}

void MpiComm::sync_v(const Mesh2D& mesh, std::vector<double>& v) {
    // Évaluation de la condition de contrôle
    if (mesh.num_procs == 1) return;

    // Déclaration de variable typée
    int nx = mesh.nx_local;
    // Déclaration de variable typée
    int ny = mesh.ny_local;

    // v est défini sur les faces horizontales j de 1 à ny+1 pour l'intérieur
    std::vector<double> send_left(ny + 2), recv_left(ny + 2);
    // Déclaration de variable typée
    std::vector<double> send_right(ny + 2), recv_right(ny + 2);
    // Déclaration de variable typée
    std::vector<double> send_bot(nx), recv_bot(nx);
    // Déclaration de variable typée
    std::vector<double> send_top(nx), recv_top(nx);

    // Boucle spatiale sur les cellules locales selon l'axe vertical Y (ordonnées)
    for (int j = 0; j <= ny + 1; ++j) {
        // Calcul et assignation numérique
        send_left[j]  = v[mesh.v_idx(1, j)];
        // Calcul et assignation numérique
        send_right[j] = v[mesh.v_idx(nx, j)];
    }

    // Appel de communication ou synchronisation parallèle OpenMPI
    MPI_Request reqs[8];
    // Déclaration de variable typée
    int n_reqs = 0;

    // Évaluation de la condition de contrôle
    if (mesh.rank_left != MPI_PROC_NULL) {
        // Appel de communication ou synchronisation parallèle OpenMPI
        MPI_Irecv(recv_left.data(), ny + 2, MPI_DOUBLE, mesh.rank_left, 50, mesh.comm, &reqs[n_reqs++]);
        // Appel de communication ou synchronisation parallèle OpenMPI
        MPI_Isend(send_left.data(), ny + 2, MPI_DOUBLE, mesh.rank_left, 51, mesh.comm, &reqs[n_reqs++]);
    }
    // Évaluation de la condition de contrôle
    if (mesh.rank_right != MPI_PROC_NULL) {
        // Appel de communication ou synchronisation parallèle OpenMPI
        MPI_Irecv(recv_right.data(), ny + 2, MPI_DOUBLE, mesh.rank_right, 51, mesh.comm, &reqs[n_reqs++]);
        // Appel de communication ou synchronisation parallèle OpenMPI
        MPI_Isend(send_right.data(), ny + 2, MPI_DOUBLE, mesh.rank_right, 50, mesh.comm, &reqs[n_reqs++]);
    }

    // Appel de communication ou synchronisation parallèle OpenMPI
    MPI_Waitall(n_reqs, reqs, MPI_STATUSES_IGNORE);
    // Calcul et assignation numérique
    n_reqs = 0;

    // Évaluation de la condition de contrôle
    if (mesh.rank_left != MPI_PROC_NULL) {
        // Boucle spatiale sur les cellules locales selon l'axe vertical Y (ordonnées)
        for (int j = 0; j <= ny + 1; ++j) v[mesh.v_idx(0, j)] = recv_left[j];
    }
    // Évaluation de la condition de contrôle
    if (mesh.rank_right != MPI_PROC_NULL) {
        // Boucle spatiale sur les cellules locales selon l'axe vertical Y (ordonnées)
        for (int j = 0; j <= ny + 1; ++j) v[mesh.v_idx(nx + 1, j)] = recv_right[j];
    }

    // Échange selon Y
    for (int i = 1; i <= nx; ++i) {
        // Calcul et assignation numérique
        send_bot[i - 1] = v[mesh.v_idx(i, 1)];
        // Calcul et assignation numérique
        send_top[i - 1] = v[mesh.v_idx(i, ny)];
    }

    // Évaluation de la condition de contrôle
    if (mesh.rank_bottom != MPI_PROC_NULL) {
        // Appel de communication ou synchronisation parallèle OpenMPI
        MPI_Irecv(recv_bot.data(), nx, MPI_DOUBLE, mesh.rank_bottom, 60, mesh.comm, &reqs[n_reqs++]);
        // Appel de communication ou synchronisation parallèle OpenMPI
        MPI_Isend(send_bot.data(), nx, MPI_DOUBLE, mesh.rank_bottom, 61, mesh.comm, &reqs[n_reqs++]);
    }
    // Évaluation de la condition de contrôle
    if (mesh.rank_top != MPI_PROC_NULL) {
        // Appel de communication ou synchronisation parallèle OpenMPI
        MPI_Irecv(recv_top.data(), nx, MPI_DOUBLE, mesh.rank_top, 61, mesh.comm, &reqs[n_reqs++]);
        // Appel de communication ou synchronisation parallèle OpenMPI
        MPI_Isend(send_top.data(), nx, MPI_DOUBLE, mesh.rank_top, 60, mesh.comm, &reqs[n_reqs++]);
    }

    // Appel de communication ou synchronisation parallèle OpenMPI
    MPI_Waitall(n_reqs, reqs, MPI_STATUSES_IGNORE);

    // Évaluation de la condition de contrôle
    if (mesh.rank_bottom != MPI_PROC_NULL) {
        // Boucle spatiale sur les cellules locales selon l'axe horizontal X (abscisses)
        for (int i = 1; i <= nx; ++i) v[mesh.v_idx(i, 0)] = recv_bot[i - 1];
    }
    // Évaluation de la condition de contrôle
    if (mesh.rank_top != MPI_PROC_NULL) {
        // Boucle spatiale sur les cellules locales selon l'axe horizontal X (abscisses)
        for (int i = 1; i <= nx; ++i) v[mesh.v_idx(i, ny + 1)] = recv_top[i - 1];
    }
}

void MpiComm::apply_boundary_conditions(const Mesh2D& mesh, const DomainBC& bc, FlowFields& f, bool homogeneous) {
    // Déclaration de variable typée
    int nx = mesh.nx_local;
    // Déclaration de variable typée
    int ny = mesh.ny_local;

    // Déclaration de variable typée
    double u_left_val   = homogeneous ? 0.0 : bc.left.u_val;
    // Déclaration de variable typée
    double v_left_val   = homogeneous ? 0.0 : bc.left.v_val;
    // Déclaration de variable typée
    double u_right_val  = homogeneous ? 0.0 : bc.right.u_val;
    // Déclaration de variable typée
    double v_right_val  = homogeneous ? 0.0 : bc.right.v_val;
    // Déclaration de variable typée
    double u_bottom_val = homogeneous ? 0.0 : bc.bottom.u_val;
    // Déclaration de variable typée
    double v_bottom_val = homogeneous ? 0.0 : bc.bottom.v_val;
    // Déclaration de variable typée
    double u_top_val    = homogeneous ? 0.0 : bc.top.u_val;
    // Déclaration de variable typée
    double v_top_val    = homogeneous ? 0.0 : bc.top.v_val;

    // 0. Traitement de la périodicité si non gérée par MPI (ex: 1 seul processeur dans la direction)
    if (bc.left.type == BCType::PERIODIC && mesh.dims[0] == 1) {
        // Boucle spatiale sur les cellules locales selon l'axe vertical Y (ordonnées)
        for (int j = 1; j <= ny; ++j) {
            // Calcul de la résolution géométrique de cellule
            f.u[mesh.u_idx(0, j)]      = f.u[mesh.u_idx(nx, j)];
            // Calcul de la résolution géométrique de cellule
            f.u[mesh.u_idx(nx + 1, j)] = f.u[mesh.u_idx(1, j)];
            // Calcul de la résolution géométrique de cellule
            f.v[mesh.v_idx(0, j)]      = f.v[mesh.v_idx(nx, j)];
            // Calcul de la résolution géométrique de cellule
            f.v[mesh.v_idx(nx + 1, j)] = f.v[mesh.v_idx(1, j)];
            // Calcul de la résolution géométrique de cellule
            f.p[mesh.cell_idx(0, j)]      = f.p[mesh.cell_idx(nx, j)];
            // Calcul de la résolution géométrique de cellule
            f.p[mesh.cell_idx(nx + 1, j)] = f.p[mesh.cell_idx(1, j)];
            // Calcul de la résolution géométrique de cellule
            f.C[mesh.cell_idx(0, j)]      = f.C[mesh.cell_idx(nx, j)];
            // Calcul de la résolution géométrique de cellule
            f.C[mesh.cell_idx(nx + 1, j)] = f.C[mesh.cell_idx(1, j)];
        }
    }

    // Évaluation de la condition de contrôle
    if (bc.bottom.type == BCType::PERIODIC && mesh.dims[1] == 1) {
        // Boucle spatiale sur les cellules locales selon l'axe horizontal X (abscisses)
        for (int i = 1; i <= nx; ++i) {
            // Calcul de la résolution géométrique de cellule
            f.v[mesh.v_idx(i, 0)]      = f.v[mesh.v_idx(i, ny)];
            // Calcul de la résolution géométrique de cellule
            f.v[mesh.v_idx(i, ny + 1)] = f.v[mesh.v_idx(i, 1)];
            // Calcul de la résolution géométrique de cellule
            f.u[mesh.u_idx(i, 0)]      = f.u[mesh.u_idx(i, ny)];
            // Calcul de la résolution géométrique de cellule
            f.u[mesh.u_idx(i, ny + 1)] = f.u[mesh.u_idx(i, 1)];
            // Calcul de la résolution géométrique de cellule
            f.p[mesh.cell_idx(i, 0)]      = f.p[mesh.cell_idx(i, ny)];
            // Calcul de la résolution géométrique de cellule
            f.p[mesh.cell_idx(i, ny + 1)] = f.p[mesh.cell_idx(i, 1)];
            // Calcul de la résolution géométrique de cellule
            f.C[mesh.cell_idx(i, 0)]      = f.C[mesh.cell_idx(i, ny)];
            // Calcul de la résolution géométrique de cellule
            f.C[mesh.cell_idx(i, ny + 1)] = f.C[mesh.cell_idx(i, 1)];
        }
    }

    // 1. Bord gauche physique (x = 0)
    if (mesh.is_boundary_left) {
        // Boucle spatiale sur les cellules locales selon l'axe vertical Y (ordonnées)
        for (int j = 1; j <= ny; ++j) {
            // Évaluation de la condition de contrôle
            if (bc.left.type == BCType::DIRICHLET) {
                // Déclaration de variable typée
                double u_inlet = u_left_val;
                // Évaluation de la condition de contrôle
                if (!homogeneous) {
                    // Évaluation de la condition de contrôle
                    if (bc.left.profile == "POISEUILLE_STEP") {
                        // Déclaration de variable typée
                        double y = mesh.yc(j);
                        // Évaluation de la condition de contrôle
                        if (y < bc.left.y_step) {
                            // Calcul et assignation numérique
                            u_inlet = 0.0;
                        } else {
                            // Déclaration de variable typée
                            double h = mesh.ly - bc.left.y_step;
                            // Calcul et assignation numérique
                            u_inlet = 6.0 * bc.left.u_mean * (y - bc.left.y_step) * (mesh.ly - y) / (h * h);
                        }
                    // Calcul et assignation numérique
                    } else if (bc.left.profile == "PARABOLIC") {
                        // Déclaration de variable typée
                        double y = mesh.yc(j);
                        // Déclaration de variable typée
                        double h = mesh.ly;
                        // Calcul et assignation numérique
                        u_inlet = 6.0 * bc.left.u_mean * y * (h - y) / (h * h);
                    }
                    // Calcul de la résolution géométrique de cellule
                    f.u[mesh.u_idx(1, j)] = u_inlet;
                }
                // Calcul de la résolution géométrique de cellule
                f.u[mesh.u_idx(0, j)] = u_inlet;
                // Calcul de la résolution géométrique de cellule
                f.v[mesh.v_idx(0, j)] = 2.0 * v_left_val - f.v[mesh.v_idx(1, j)];
            // Calcul et assignation numérique
            } else if (bc.left.type == BCType::FREE_SLIP) {
                // Évaluation de la condition de contrôle
                if (!homogeneous) f.u[mesh.u_idx(1, j)] = 0.0;
                // Calcul de la résolution géométrique de cellule
                f.u[mesh.u_idx(0, j)] = 0.0;
                // Calcul de la résolution géométrique de cellule
                f.v[mesh.v_idx(0, j)] = f.v[mesh.v_idx(1, j)];
            // Calcul et assignation numérique
            } else if (bc.left.type == BCType::NEUMANN) {
                // Calcul de la résolution géométrique de cellule
                f.u[mesh.u_idx(0, j)] = f.u[mesh.u_idx(1, j)];
                // Calcul de la résolution géométrique de cellule
                f.v[mesh.v_idx(0, j)] = f.v[mesh.v_idx(1, j)];
            }
            // Calcul de la résolution géométrique de cellule
            f.p[mesh.cell_idx(0, j)] = f.p[mesh.cell_idx(1, j)];
            // Calcul de la résolution géométrique de cellule
            f.C[mesh.cell_idx(0, j)] = f.C[mesh.cell_idx(1, j)];
        }
    }

    // 2. Bord droit physique (x = Lx)
    if (mesh.is_boundary_right) {
        // Boucle spatiale sur les cellules locales selon l'axe vertical Y (ordonnées)
        for (int j = 1; j <= ny; ++j) {
            // Évaluation de la condition de contrôle
            if (bc.right.type == BCType::DIRICHLET) {
                // Évaluation de la condition de contrôle
                if (!homogeneous) f.u[mesh.u_idx(nx + 1, j)] = u_right_val;
                // Calcul de la résolution géométrique de cellule
                f.u[mesh.u_idx(nx + 2, j)] = u_right_val;
                // Calcul de la résolution géométrique de cellule
                f.v[mesh.v_idx(nx + 1, j)] = 2.0 * v_right_val - f.v[mesh.v_idx(nx, j)];
            // Calcul et assignation numérique
            } else if (bc.right.type == BCType::FREE_SLIP) {
                // Évaluation de la condition de contrôle
                if (!homogeneous) f.u[mesh.u_idx(nx + 1, j)] = 0.0;
                // Calcul de la résolution géométrique de cellule
                f.u[mesh.u_idx(nx + 2, j)] = 0.0;
                // Calcul de la résolution géométrique de cellule
                f.v[mesh.v_idx(nx + 1, j)] = f.v[mesh.v_idx(nx, j)];
            // Calcul et assignation numérique
            } else if (bc.right.type == BCType::NEUMANN) {
                // Calcul de la résolution géométrique de cellule
                f.u[mesh.u_idx(nx + 2, j)] = f.u[mesh.u_idx(nx + 1, j)];
                // Calcul de la résolution géométrique de cellule
                f.v[mesh.v_idx(nx + 1, j)] = f.v[mesh.v_idx(nx, j)];
            }
            // Calcul de la résolution géométrique de cellule
            f.p[mesh.cell_idx(nx + 1, j)] = f.p[mesh.cell_idx(nx, j)];
            // Calcul de la résolution géométrique de cellule
            f.C[mesh.cell_idx(nx + 1, j)] = f.C[mesh.cell_idx(nx, j)];
        }
    }

    // 3. Bord inférieur physique (y = 0)
    if (mesh.is_boundary_bottom) {
        // Boucle spatiale sur les cellules locales selon l'axe horizontal X (abscisses)
        for (int i = 1; i <= nx; ++i) {
            // Évaluation de la condition de contrôle
            if (bc.bottom.type == BCType::DIRICHLET) {
                // Évaluation de la condition de contrôle
                if (!homogeneous) f.v[mesh.v_idx(i, 1)] = v_bottom_val;
                // Calcul de la résolution géométrique de cellule
                f.v[mesh.v_idx(i, 0)] = v_bottom_val;
                // Calcul de la résolution géométrique de cellule
                f.u[mesh.u_idx(i, 0)] = 2.0 * u_bottom_val - f.u[mesh.u_idx(i, 1)];
            // Calcul et assignation numérique
            } else if (bc.bottom.type == BCType::FREE_SLIP) {
                // Évaluation de la condition de contrôle
                if (!homogeneous) f.v[mesh.v_idx(i, 1)] = 0.0;
                // Calcul de la résolution géométrique de cellule
                f.v[mesh.v_idx(i, 0)] = 0.0;
                // Calcul de la résolution géométrique de cellule
                f.u[mesh.u_idx(i, 0)] = f.u[mesh.u_idx(i, 1)];
            // Calcul et assignation numérique
            } else if (bc.bottom.type == BCType::NEUMANN) {
                // Calcul de la résolution géométrique de cellule
                f.u[mesh.u_idx(i, 0)] = f.u[mesh.u_idx(i, 1)];
                // Calcul de la résolution géométrique de cellule
                f.v[mesh.v_idx(i, 0)] = f.v[mesh.v_idx(i, 1)];
            }
            // Calcul de la résolution géométrique de cellule
            f.p[mesh.cell_idx(i, 0)] = f.p[mesh.cell_idx(i, 1)];
            // Calcul de la résolution géométrique de cellule
            f.C[mesh.cell_idx(i, 0)] = f.C[mesh.cell_idx(i, 1)];
        }
    }

    // 4. Bord supérieur physique (y = Ly)
    if (mesh.is_boundary_top) {
        // Boucle spatiale sur les cellules locales selon l'axe horizontal X (abscisses)
        for (int i = 1; i <= nx; ++i) {
            // Déclaration de variable typée
            double x_c = mesh.xc(i);
            // Évaluation du profil d'injecteur (Cas 3.3 Fluids 2021)
            if (bc.top.profile == "JET_INLET") {
                // Déclaration de variable typée pour le centre et demi-largeur
                double xc = (bc.top.jet_center > 0.0) ? bc.top.jet_center : 0.5 * mesh.lx;
                // Calcul de la demi-largeur de l'injecteur
                double half_w = (bc.top.jet_width > 0.0) ? 0.5 * bc.top.jet_width : 0.04;
                // Calcul de la largeur discrète effective de l'injecteur sur le maillage global
                double d_disc = 0.0;
                for (int ig = 1; ig <= mesh.nx_global; ++ig) {
                    double xg = (ig - 0.5) * mesh.dx;
                    if (std::abs(xg - xc) <= half_w) d_disc += mesh.dx;
                }
                if (d_disc < 1.0e-12) d_disc = 2.0 * half_w;
                double v_in_eff = v_top_val * (2.0 * half_w) / d_disc;
                // Évaluation de l'appartenance à l'injecteur
                if (std::abs(x_c - xc) <= half_w) {
                    // Imposition de la vitesse d'injection (préservée pour la ligne diagonale lors de matvec)
                    if (!homogeneous) f.v[mesh.v_idx(i, ny + 1)] = v_in_eff;
                    // Calcul de la cellule fantôme supérieure pour divergence nulle au bord
                    f.v[mesh.v_idx(i, ny + 2)] = homogeneous ? 0.0 : v_in_eff;
                    // Injection continue de liquide (Fluide 2 : C = 0)
                    f.C[mesh.cell_idx(i, ny + 1)] = 0.0;
                } else {
                    // Calcul du débit de compensation d'échappement du gaz pour incompressibilité stricte
                    double v_out = -v_top_val * (2.0 * half_w) / (mesh.lx - d_disc);
                    // Imposition de la vitesse d'évacuation (préservée pour la ligne diagonale lors de matvec)
                    if (!homogeneous) f.v[mesh.v_idx(i, ny + 1)] = v_out;
                    // Calcul de la cellule fantôme supérieure
                    f.v[mesh.v_idx(i, ny + 2)] = homogeneous ? 0.0 : v_out;
                    // Évacuation du gaz par sortie libre Neumann
                    f.C[mesh.cell_idx(i, ny + 1)] = f.C[mesh.cell_idx(i, ny)];
                }
                // Condition tangentielle symétrique évaluée à l'abscisse exacte de la face u
                f.u[mesh.u_idx(i, ny + 1)] = (std::abs(mesh.xu(i) - xc) <= half_w) ? 0.0 : f.u[mesh.u_idx(i, ny)];
                if (i == nx) f.u[mesh.u_idx(nx + 1, ny + 1)] = f.u[mesh.u_idx(nx + 1, ny)];
            // Évaluation de la condition de contrôle standard
            } else if (bc.top.type == BCType::DIRICHLET) {
                // Évaluation de la condition de contrôle
                if (!homogeneous) f.v[mesh.v_idx(i, ny + 1)] = v_top_val;
                // Calcul de la résolution géométrique de cellule
                f.v[mesh.v_idx(i, ny + 2)] = v_top_val;
                // Calcul de la résolution géométrique de cellule
                f.u[mesh.u_idx(i, ny + 1)] = 2.0 * u_top_val - f.u[mesh.u_idx(i, ny)];
                // Calcul de la résolution géométrique de cellule
                f.C[mesh.cell_idx(i, ny + 1)] = f.C[mesh.cell_idx(i, ny)];
            // Calcul et assignation numérique
            } else if (bc.top.type == BCType::FREE_SLIP) {
                // Évaluation de la condition de contrôle
                if (!homogeneous) f.v[mesh.v_idx(i, ny + 1)] = 0.0;
                // Calcul de la résolution géométrique de cellule
                f.v[mesh.v_idx(i, ny + 2)] = 0.0;
                // Calcul de la résolution géométrique de cellule
                f.u[mesh.u_idx(i, ny + 1)] = f.u[mesh.u_idx(i, ny)];
                // Calcul de la résolution géométrique de cellule
                f.C[mesh.cell_idx(i, ny + 1)] = f.C[mesh.cell_idx(i, ny)];
            // Calcul et assignation numérique
            } else if (bc.top.type == BCType::NEUMANN) {
                // Calcul de la résolution géométrique de cellule
                f.u[mesh.u_idx(i, ny + 1)] = f.u[mesh.u_idx(i, ny)];
                // Calcul de la résolution géométrique de cellule
                f.v[mesh.v_idx(i, ny + 2)] = f.v[mesh.v_idx(i, ny + 1)];
                // Calcul de la résolution géométrique de cellule
                f.C[mesh.cell_idx(i, ny + 1)] = f.C[mesh.cell_idx(i, ny)];
            }
            // Calcul de la résolution géométrique de cellule
            f.p[mesh.cell_idx(i, ny + 1)] = f.p[mesh.cell_idx(i, ny)];
        }
    }
}

void MpiComm::apply_thermal_boundary_conditions(const Mesh2D& mesh, const DomainBC& bc, FlowFields& f) {
    // Déclaration de variable typée
    int nx = mesh.nx_local;
    // Déclaration de variable typée
    int ny = mesh.ny_local;

    // Périodicité si 1 seul processeur dans la direction
    if (bc.left.T_type == BCType::PERIODIC && mesh.dims[0] == 1) {
        // Boucle spatiale sur les cellules locales selon l'axe vertical Y (ordonnées)
        for (int j = 1; j <= ny; ++j) {
            // Calcul de la résolution géométrique de cellule
            f.T[mesh.cell_idx(0, j)]      = f.T[mesh.cell_idx(nx, j)];
            // Calcul de la résolution géométrique de cellule
            f.T[mesh.cell_idx(nx + 1, j)] = f.T[mesh.cell_idx(1, j)];
        }
    }
    // Évaluation de la condition de contrôle
    if (bc.bottom.T_type == BCType::PERIODIC && mesh.dims[1] == 1) {
        // Boucle spatiale sur les cellules locales selon l'axe horizontal X (abscisses)
        for (int i = 1; i <= nx; ++i) {
            // Calcul de la résolution géométrique de cellule
            f.T[mesh.cell_idx(i, 0)]      = f.T[mesh.cell_idx(i, ny)];
            // Calcul de la résolution géométrique de cellule
            f.T[mesh.cell_idx(i, ny + 1)] = f.T[mesh.cell_idx(i, 1)];
        }
    }

    // Bord gauche physique (x = 0)
    if (mesh.is_boundary_left) {
        // Boucle spatiale sur les cellules locales selon l'axe vertical Y (ordonnées)
        for (int j = 1; j <= ny; ++j) {
            // Évaluation de la condition de contrôle
            if (bc.left.T_type == BCType::DIRICHLET) {
                // Calcul de la résolution géométrique de cellule
                f.T[mesh.cell_idx(0, j)] = 2.0 * bc.left.T_val - f.T[mesh.cell_idx(1, j)];
            } else { // NEUMANN / Adiabatique
                // Calcul de la résolution géométrique de cellule
                f.T[mesh.cell_idx(0, j)] = f.T[mesh.cell_idx(1, j)];
            }
        }
    }

    // Bord droit physique (x = Lx)
    if (mesh.is_boundary_right) {
        // Boucle spatiale sur les cellules locales selon l'axe vertical Y (ordonnées)
        for (int j = 1; j <= ny; ++j) {
            // Évaluation de la condition de contrôle
            if (bc.right.T_type == BCType::DIRICHLET) {
                // Calcul de la résolution géométrique de cellule
                f.T[mesh.cell_idx(nx + 1, j)] = 2.0 * bc.right.T_val - f.T[mesh.cell_idx(nx, j)];
            } else { // NEUMANN
                // Calcul de la résolution géométrique de cellule
                f.T[mesh.cell_idx(nx + 1, j)] = f.T[mesh.cell_idx(nx, j)];
            }
        }
    }

    // Bord inférieur physique (y = 0)
    if (mesh.is_boundary_bottom) {
        // Boucle spatiale sur les cellules locales selon l'axe horizontal X (abscisses)
        for (int i = 1; i <= nx; ++i) {
            // Évaluation de la condition de contrôle
            if (bc.bottom.T_type == BCType::DIRICHLET) {
                // Calcul de la résolution géométrique de cellule
                f.T[mesh.cell_idx(i, 0)] = 2.0 * bc.bottom.T_val - f.T[mesh.cell_idx(i, 1)];
            } else { // NEUMANN / Adiabatique
                // Calcul de la résolution géométrique de cellule
                f.T[mesh.cell_idx(i, 0)] = f.T[mesh.cell_idx(i, 1)];
            }
        }
    }

    // Bord supérieur physique (y = Ly)
    if (mesh.is_boundary_top) {
        // Boucle spatiale sur les cellules locales selon l'axe horizontal X (abscisses)
        for (int i = 1; i <= nx; ++i) {
            // Évaluation de la condition de contrôle
            if (bc.top.T_type == BCType::DIRICHLET) {
                // Calcul de la résolution géométrique de cellule
                f.T[mesh.cell_idx(i, ny + 1)] = 2.0 * bc.top.T_val - f.T[mesh.cell_idx(i, ny)];
            } else { // NEUMANN / Adiabatique
                // Calcul de la résolution géométrique de cellule
                f.T[mesh.cell_idx(i, ny + 1)] = f.T[mesh.cell_idx(i, ny)];
            }
        }
    }
}

// Déclaration de variable typée
double MpiComm::global_dot_product(const Mesh2D& mesh, const std::vector<double>& a, const std::vector<double>& b, int size) {
    // Déclaration de variable typée
    double local_sum = 0.0;
    // Boucle spatiale sur les cellules locales selon l'axe horizontal X (abscisses)
    for (int i = 0; i < size; ++i) {
        // Calcul et assignation numérique
        local_sum += a[i] * b[i];
    }
    // Déclaration de variable typée
    double global_sum = 0.0;
    // Appel de communication ou synchronisation parallèle OpenMPI
    MPI_Allreduce(&local_sum, &global_sum, 1, MPI_DOUBLE, MPI_SUM, mesh.comm);
    // Retour de fonction
    return global_sum;
}

// Déclaration de variable typée
double MpiComm::global_norm2(const Mesh2D& mesh, const std::vector<double>& v, int size) {
    // Retour de fonction
    return std::sqrt(global_dot_product(mesh, v, v, size));
}

// Déclaration de variable typée
double MpiComm::global_max(const Mesh2D& mesh, double local_val) {
    // Déclaration de variable typée
    double global_val = 0.0;
    // Appel de communication ou synchronisation parallèle OpenMPI
    MPI_Allreduce(&local_val, &global_val, 1, MPI_DOUBLE, MPI_MAX, mesh.comm);
    // Retour de fonction
    return global_val;
}

// Déclaration de variable typée
double MpiComm::global_sum(const Mesh2D& mesh, double local_val) {
    // Déclaration de variable typée
    double global_val = 0.0;
    // Appel de communication ou synchronisation parallèle OpenMPI
    MPI_Allreduce(&local_val, &global_val, 1, MPI_DOUBLE, MPI_SUM, mesh.comm);
    // Retour de fonction
    return global_val;
}

} // namespace covof

\end{lstlisting}

\chapter{Algèbre Linéaire Creuse au Format CSR (CSRMatrix.hpp, CSRMatrix.cpp)}
\section{Représentation Compressed Sparse Row (CSR)}
Le système linéaire monolithique du Lagrangien Augmenté comporte $N_{\text{dof}} = N_u + N_v$ degrés de liberté par sous-domaine, mais chaque ligne de la matrice ne possède qu'au plus 9 coefficients non nuls (5 termes directs du Laplacien de la composante courante + 4 termes croisés du gradient de divergence couplant $u$ et $v$). Stocker une matrice dense nécessiterait $\mathcal{O}(N_{\text{dof}}^2)$ octets (plusieurs gigaoctets dès $128\times 128$), tandis que le format creux \textbf{CSR (\emph{Compressed Sparse Row})} implémenté dans \texttt{covof::CSRMatrix} réduit l'empreinte mémoire à $\mathcal{O}(9 N_{\text{dof}})$ (quelques mégaoctets).

Une matrice creuse $\mathbf{A} \in \mathbb{R}^{n \times n}$ comportant $\text{nnz}$ éléments non nuls est représentée par trois tableaux contigus :
\begin{enumerate}
  \item \texttt{values} ($\in \mathbb{R}^{\text{nnz}}$) : valeurs flottantes des éléments non nuls parcourus ligne par ligne ;
  \item \texttt{col\_ind} ($\in \mathbb{N}^{\text{nnz}}$) : indices de colonnes $j$ correspondants ;
  \item \texttt{row\_ptr} ($\in \mathbb{N}^{n+1}$) : indice de début de chaque ligne $i$ dans \texttt{values} et \texttt{col\_ind}, avec $\texttt{row\_ptr}[0] = 0$ et $\texttt{row\_ptr}[n] = \text{nnz}$.
\end{enumerate}

\section{Produit Matrice-Vecteur Creux (\texttt{multiply} / SpMV)}
Pour calculer $\mathbf{y} = \mathbf{A}\mathbf{x}$, la méthode \texttt{CSRMatrix::multiply(const std::vector<double>\& x, std::vector<double>\& y)} évalue pour chaque ligne $i \in \{0, \dots, n-1\}$ :
\begin{equation}
  y_i = \sum_{k = \texttt{row\_ptr}[i]}^{\texttt{row\_ptr}[i+1]-1} \texttt{values}[k] \, x_{\texttt{col\_ind}[k]}
\end{equation}
Cette opération s'exécute en exactement $2\,\text{nnz}$ opérations flottantes (FLOPs) avec un parcours purement séquentiel des tableaux \texttt{values} et \texttt{col\_ind}.

\section{Répertoire des Classes, Structures et Méthodes de \texttt{include/covof/CSRMatrix.hpp}}
\begin{center}\small
\begin{tabular}{clll}
\toprule
\textbf{Ligne} & \textbf{Catégorie} & \textbf{Identifiant C++17} & \textbf{Signature / Rôle} \\
\midrule
13 & class & \texttt{CSRMatrix} & Déclaration de type / classe / structure \\
34 & Méthode / Fonction & \texttt{init} & Type retour : void \\
54 & Méthode / Fonction & \texttt{matvec} & Type retour : void \\
59 & Méthode / Fonction & \texttt{solve\_jacobi} & Type retour : void \\
\bottomrule
\end{tabular}
\end{center}

\section{Code Source Intégral Annoté : \texttt{include/covof/CSRMatrix.hpp}}
Le listing ci-dessous présente l'intégralité des 62 lignes du fichier \texttt{include/covof/CSRMatrix.hpp} de l'implémentation C++17 :

\begin{lstlisting}[language={[11]C++}]
#pragma once

#include "Mesh.hpp"
#include <vector>

namespace covof {
namespace fugu = covof; // Alias rétrocompatibilité

/**
 * @brief Structure de matrice creuse locale au format Compressed Sparse Row (CSR)
 * Utilisée pour les blocs de préconditionnement (Fuu, Fvv, Laplacien de Schur).
 */
class CSRMatrix {
public:
    // Déclaration de variable typée
    int nrows = 0;
    // Déclaration de variable typée
    int ncols = 0;
    // Déclaration de variable typée
    std::vector<int> row_ptr;     ///< Pointeurs de début de ligne (taille nrows + 1)
    // Déclaration de variable typée
    std::vector<int> col_idx;     ///< Indices de colonnes
    // Déclaration de variable typée
    std::vector<double> values;   ///< Valeurs non nulles
    // Déclaration de variable typée
    std::vector<double> diag;     ///< Diagonale explicite pour le préconditionnement Jacobi / Gauss-Seidel

    // Calcul et assignation numérique
    CSRMatrix() = default;

    /**
     * @brief Initialise les dimensions
     */
    void init(int r, int c) {
        // Calcul et assignation numérique
        nrows = r;
        // Calcul et assignation numérique
        ncols = c;
        row_ptr.assign(r + 1, 0);
        col_idx.clear();
        values.clear();
        diag.assign(r, 0.0);
    }

    /**
     * @brief Construit le format CSR à partir d'une liste de triplets (row, col, value)
     */
    void set_from_triplets(int r, int c, const std::vector<int>& rows,
                           const std::vector<int>& cols, const std::vector<double>& vals);

    /**
     * @brief Produit matrice-vecteur local : y = A * x
     */
    void matvec(const std::vector<double>& x, std::vector<double>& y) const;

    /**
     * @brief Résolution approchée par itérations de Jacobi : x = D^{-1} b
     */
    void solve_jacobi(const std::vector<double>& b, std::vector<double>& x, int num_iters = 5) const;
};

} // namespace covof

\end{lstlisting}

\section{Répertoire des Classes, Structures et Méthodes de \texttt{src/CSRMatrix.cpp}}
\begin{center}\small
\begin{tabular}{clll}
\toprule
\textbf{Ligne} & \textbf{Catégorie} & \textbf{Identifiant C++17} & \textbf{Signature / Rôle} \\
\midrule
61 & Méthode / Fonction & \texttt{CSRMatrix::matvec} & Type retour : void \\
77 & Méthode / Fonction & \texttt{CSRMatrix::solve\_jacobi} & Type retour : void \\
80 & Méthode / Fonction & \texttt{x\_new} & Type retour : std::vector$<$double$>$ \\
\bottomrule
\end{tabular}
\end{center}

\section{Code Source Intégral Annoté : \texttt{src/CSRMatrix.cpp}}
Le listing ci-dessous présente l'intégralité des 108 lignes du fichier \texttt{src/CSRMatrix.cpp} de l'implémentation C++17 :

\begin{lstlisting}[language={[11]C++}]
#include "covof/CSRMatrix.hpp"
#include <numeric>
#include <algorithm>
#include <cassert>

namespace covof {
namespace fugu = covof; // Alias rétrocompatibilité

void CSRMatrix::set_from_triplets(int r, int c, const std::vector<int>& rows,
                                 const std::vector<int>& cols, const std::vector<double>& vals) {
    // Calcul et assignation numérique
    nrows = r;
    // Calcul et assignation numérique
    ncols = c;
    row_ptr.assign(nrows + 1, 0);
    diag.assign(nrows, 0.0);

    // 1. Comptage des éléments par ligne
    for (size_t k = 0; k < rows.size(); ++k) {
        // Évaluation de la condition de contrôle
        if (rows[k] >= 0 && rows[k] < nrows) {
            row_ptr[rows[k] + 1]++;
        }
    }

    // 2. Somme cumulée pour row_ptr
    for (int i = 0; i < nrows; ++i) {
        // Calcul et assignation numérique
        row_ptr[i + 1] += row_ptr[i];
    }

    // Déclaration de variable typée
    int nnz = row_ptr[nrows];
    col_idx.resize(nnz);
    values.resize(nnz);

    // Déclaration de variable typée
    std::vector<int> current_pos = row_ptr;

    // 3. Remplissage des indices et valeurs
    for (size_t k = 0; k < rows.size(); ++k) {
        // Déclaration de variable typée
        int row = rows[k];
        // Évaluation de la condition de contrôle
        if (row >= 0 && row < nrows) {
            // Déclaration de variable typée
            int dest = current_pos[row]++;
            // Calcul de la résolution géométrique de cellule
            col_idx[dest] = cols[k];
            // Calcul et assignation numérique
            values[dest]  = vals[k];
            // Évaluation de la condition de contrôle
            if (row == cols[k]) {
                // Calcul et assignation numérique
                diag[row] = vals[k];
            }
        }
    }
}

void CSRMatrix::matvec(const std::vector<double>& x, std::vector<double>& y) const {
    y.assign(nrows, 0.0);
    // Boucle spatiale sur les cellules locales selon l'axe horizontal X (abscisses)
    for (int i = 0; i < nrows; ++i) {
        // Déclaration de variable typée
        double sum = 0.0;
        // Boucle d'itération algorithmique
        for (int k = row_ptr[i]; k < row_ptr[i + 1]; ++k) {
            // Calcul et assignation numérique
            sum += values[k] * x[col_idx[k]];
        }
        // Calcul et assignation numérique
        y[i] = sum;
    }
}

void CSRMatrix::solve_jacobi(const std::vector<double>& b, std::vector<double>& x, int num_iters) const {
    x.resize(nrows, 0.0);
    // Déclaration de variable typée
    std::vector<double> x_new(nrows, 0.0);

    // Boucle d'itération algorithmique
    for (int it = 0; it < num_iters; ++it) {
        // Boucle spatiale sur les cellules locales selon l'axe horizontal X (abscisses)
        for (int i = 0; i < nrows; ++i) {
            // Déclaration de variable typée
            double d = (std::abs(diag[i]) > 1.0e-15) ? diag[i] : 1.0;
            // Déclaration de variable typée
            double sum = 0.0;
            // Boucle d'itération algorithmique
            for (int k = row_ptr[i]; k < row_ptr[i + 1]; ++k) {
                // Déclaration de variable typée
                int j = col_idx[k];
                // Évaluation de la condition de contrôle
                if (j != i) {
                    // Calcul et assignation numérique
                    sum += values[k] * x[j];
                }
            }
            // Calcul et assignation numérique
            x_new[i] = (b[i] - sum) / d;
        }
        // Calcul et assignation numérique
        x = x_new;
    }
}

} // namespace covof

\end{lstlisting}

\chapter{Préconditionneurs Algébriques : Jacobi, SSOR et ILU(0) (Preconditioner.hpp, Preconditioner.cpp)}
\section{Rôle du Préconditionnement dans les Méthodes de Krylov}
Pour accélérer la convergence des solveurs itératifs sur le système $\mathbf{A}\mathbf{x} = \mathbf{b}$ dont le nombre de conditionnement $\kappa(\mathbf{A})$ croît avec le paramètre du Lagrangien Augmenté $r_{al}$ et le rapport de viscosité $\mu_2 / \mu_1$, \texttt{COVOF\_2D-C++} définit l'interface abstraite \texttt{covof::Preconditioner} qui résout à chaque itération de Krylov le système préconditionné $\mathbf{M}\mathbf{z} = \mathbf{r}$ (ou applique $\mathbf{z} = \mathbf{M}^{-1}\mathbf{r}$).

\subsection{Préconditionneur Diagonal de Jacobi (\texttt{JacobiPreconditioner})}
On choisit $\mathbf{M} = \text{diag}(A_{0,0}, A_{1,1}, \dots, A_{n-1,n-1})$. Lors de la phase \texttt{setup(A)}, l'inverse de la diagonale est précalculé avec garde-fou contre les pivots nuls :
\begin{equation}
  d_i^{-1} = \begin{cases} 1 / A_{i,i} & \text{si } |A_{i,i}| > 10^{-15} \\ 1.0 & \text{sinon} \end{cases}, \qquad z_i = d_i^{-1} r_i
\end{equation}
Dans un système à forts sauts de densité ($\rho_2/\rho_1 = 10^5$) et de viscosité ($\mu_2/\mu_1 = 10^{10}$), ce préconditionneur diagonal remet automatiquement à l'échelle les équations situées dans le gaz et dans le liquide dense, ramenant les éléments diagonaux de $\mathbf{M}^{-1}\mathbf{A}$ à l'unité.

\subsection{Préconditionneur SSOR et Factorisation Incomplète \texttt{ILU(0)}}
\begin{itemize}
  \item \textbf{SSOR (\emph{Symmetric Successive Over-Relaxation})} : décompose $\mathbf{A} = \mathbf{D} + \mathbf{L} + \mathbf{U}$ et approche $\mathbf{M}_{\text{SSOR}} = \frac{1}{\omega(2-\omega)}(\mathbf{D} + \omega\mathbf{L})\mathbf{D}^{-1}(\mathbf{D} + \omega\mathbf{U})$ inversé par une descente triangulaire inférieure suivie d'une remontée triangulaire supérieure.
  \item \textbf{ILU(0) (\emph{Incomplete LU with zero fill-in})} : calcule une factorisation $\mathbf{A} \approx \tilde{\mathbf{L}}\tilde{\mathbf{U}}$ conservant exactement le même motif de parcimonie (\emph{sparsity pattern}) que la matrice $\mathbf{A}$ originale.
\end{itemize}

\section{Répertoire des Classes, Structures et Méthodes de \texttt{include/covof/Preconditioner.hpp}}
\begin{center}\small
\begin{tabular}{clll}
\toprule
\textbf{Ligne} & \textbf{Catégorie} & \textbf{Identifiant C++17} & \textbf{Signature / Rôle} \\
\midrule
12 & class & \texttt{NavierStokesSolver} & Déclaration de type / classe / structure \\
13 & class & \texttt{FlowFields} & Déclaration de type / classe / structure \\
19 & class & \texttt{PreconditionerBase} & Déclaration de type / classe / structure \\
27 & Méthode / Fonction & \texttt{init} & Type retour : virtual void \\
32 & Méthode / Fonction & \texttt{update} & Type retour : virtual void \\
37 & Méthode / Fonction & \texttt{apply} & Type retour : virtual void \\
42 & Méthode / Fonction & \texttt{get\_name} & Type retour : virtual std::string \\
50 & class & \texttt{JacobiPreconditioner} & Déclaration de type / classe / structure \\
55 & Méthode / Fonction & \texttt{init} & Type retour : void \\
56 & Méthode / Fonction & \texttt{update} & Type retour : void \\
57 & Méthode / Fonction & \texttt{apply} & Type retour : void \\
59 & Méthode / Fonction & \texttt{get\_name} & Type retour : std::string \\
66 & class & \texttt{SSORPreconditioner} & Déclaration de type / classe / structure \\
73 & Méthode / Fonction & \texttt{init} & Type retour : void \\
74 & Méthode / Fonction & \texttt{update} & Type retour : void \\
75 & Méthode / Fonction & \texttt{apply} & Type retour : void \\
77 & Méthode / Fonction & \texttt{get\_name} & Type retour : std::string \\
84 & class & \texttt{ILU0Preconditioner} & Déclaration de type / classe / structure \\
91 & Méthode / Fonction & \texttt{init} & Type retour : void \\
92 & Méthode / Fonction & \texttt{update} & Type retour : void \\
93 & Méthode / Fonction & \texttt{apply} & Type retour : void \\
95 & Méthode / Fonction & \texttt{get\_name} & Type retour : std::string \\
102 & class & \texttt{BlockTriangularPreconditioner} & Déclaration de type / classe / structure \\
107 & Méthode / Fonction & \texttt{init} & Type retour : void \\
108 & Méthode / Fonction & \texttt{update} & Type retour : void \\
109 & Méthode / Fonction & \texttt{apply} & Type retour : void \\
111 & Méthode / Fonction & \texttt{get\_name} & Type retour : std::string \\
118 & class & \texttt{CustomPreconditioner} & Déclaration de type / classe / structure \\
126 & Méthode / Fonction & \texttt{init} & Type retour : void \\
127 & Méthode / Fonction & \texttt{update} & Type retour : void \\
128 & Méthode / Fonction & \texttt{apply} & Type retour : void \\
130 & Méthode / Fonction & \texttt{get\_name} & Type retour : std::string \\
133 & Méthode / Fonction & \texttt{set\_custom\_operator} & Type retour : void \\
144 & class & \texttt{PreconditionerFactory} & Déclaration de type / classe / structure \\
146 & Méthode / Fonction & \texttt{create} & Type retour : static std::unique\_ptr$<$PreconditionerBase$>$ \\
\bottomrule
\end{tabular}
\end{center}

\section{Code Source Intégral Annoté : \texttt{include/covof/Preconditioner.hpp}}
Le listing ci-dessous présente l'intégralité des 152 lignes du fichier \texttt{include/covof/Preconditioner.hpp} de l'implémentation C++17 :

\begin{lstlisting}[language={[11]C++}]
#pragma once

#include "Mesh.hpp"
#include <vector>
#include <string>
#include <memory>
#include <functional>

namespace covof {
namespace fugu = covof; // Alias rétrocompatibilité

class NavierStokesSolver;
class FlowFields;

/**
 * @brief Classe de base abstraite pour tous les préconditionneurs (Standard CFD / COVOF 2D)
 * Permet l'extension et la création facile de nouveaux préconditionneurs par héritage.
 */
class PreconditionerBase {
public:
    // Calcul et assignation numérique
    virtual ~PreconditionerBase() = default;

    /**
     * @brief Initialise le préconditionneur avec l'opérateur Navier-Stokes
     */
    virtual void init(const NavierStokesSolver* solver) = 0;

    /**
     * @brief Met à jour les coefficients du préconditionneur si les propriétés ou le pas de temps changent
     */
    virtual void update(const FlowFields& fields, double dt) = 0;

    /**
     * @brief Applique l'action du préconditionneur : r = P^{-1} * z
     */
    virtual void apply(const std::vector<double>& z, std::vector<double>& r) const = 0;

    /**
     * @brief Renvoie le nom identifiant du préconditionneur
     */
    virtual std::string get_name() const = 0;
};

/**
 * @brief Préconditionneur Diagonal de Jacobi
 * P = Diag(A)
 * Idéal pour le parallélisme massif et les formulations par Lagrangien Augmenté.
 */
class JacobiPreconditioner : public PreconditionerBase {
public:
    // Calcul et assignation numérique
    const NavierStokesSolver* ns_solver = nullptr;

    void init(const NavierStokesSolver* solver) override;
    void update(const FlowFields& fields, double dt) override;
    void apply(const std::vector<double>& z, std::vector<double>& r) const override;
    // Déclaration de variable typée
    std::string get_name() const override { return "JACOBI"; }
};

/**
 * @brief Préconditionneur SSOR (Symmetric Successive Over-Relaxation)
 * P = 1/(2-w) * (D/w + L) * (D/w)^{-1} * (D/w + U)
 */
class SSORPreconditioner : public PreconditionerBase {
public:
    // Calcul et assignation numérique
    const NavierStokesSolver* ns_solver = nullptr;
    // Déclaration de variable typée
    double omega = 1.0; // Facteur de relaxation (defaut Gauss-Seidel symetrique)

    void init(const NavierStokesSolver* solver) override;
    void update(const FlowFields& fields, double dt) override;
    void apply(const std::vector<double>& z, std::vector<double>& r) const override;
    // Déclaration de variable typée
    std::string get_name() const override { return "SSOR"; }
};

/**
 * @brief Préconditionneur ILU(0) (Incomplete LU à remplissage nul)
 * Factorisation incomplète locale sur chaque sous-domaine Cartésien.
 */
class ILU0Preconditioner : public PreconditionerBase {
public:
    // Calcul et assignation numérique
    const NavierStokesSolver* ns_solver = nullptr;
    // Déclaration de variable typée
    std::vector<double> ilu_diag;

    void init(const NavierStokesSolver* solver) override;
    void update(const FlowFields& fields, double dt) override;
    void apply(const std::vector<double>& z, std::vector<double>& r) const override;
    // Déclaration de variable typée
    std::string get_name() const override { return "ILU0"; }
};

/**
 * @brief Préconditionneur par Blocs Triangulaires (Annexe E thèse M. El Ouafa / Tavares)
 * Dédié au système tout-couplé vitesse-pression avec approximation du complément de Schur.
 */
class BlockTriangularPreconditioner : public PreconditionerBase {
public:
    // Calcul et assignation numérique
    const NavierStokesSolver* ns_solver = nullptr;

    void init(const NavierStokesSolver* solver) override;
    void update(const FlowFields& fields, double dt) override;
    void apply(const std::vector<double>& z, std::vector<double>& r) const override;
    // Déclaration de variable typée
    std::string get_name() const override { return "BLOCK_TRIANGULAR"; }
};

/**
 * @brief Préconditionneur Personnalisé (Hook développeur)
 * Permet au développeur d'injecter une fonction lambda ou un foncteur C++ arbitraire sans modifier le code source.
 */
class CustomPreconditioner : public PreconditionerBase {
public:
    // Calcul et assignation numérique
    using ApplyFunc = std::function<void(const std::vector<double>&, std::vector<double>&)>;
    ApplyFunc custom_apply;
    // Déclaration de variable typée
    std::string custom_name = "CUSTOM";

    void init(const NavierStokesSolver* solver) override;
    void update(const FlowFields& fields, double dt) override;
    void apply(const std::vector<double>& z, std::vector<double>& r) const override;
    // Déclaration de variable typée
    std::string get_name() const override { return custom_name; }

    // Calcul et assignation numérique
    void set_custom_operator(ApplyFunc func, const std::string& name = "USER_DEFINED") {
        // Calcul et assignation numérique
        custom_apply = func;
        // Calcul et assignation numérique
        custom_name = name;
    }
};

/**
 * @brief Fabrique de préconditionneurs (Factory Pattern)
 */
class PreconditionerFactory {
public:
    static std::unique_ptr<PreconditionerBase> create(const std::string& type);
};

// Alias de rétrocompatibilité
using BlockPreconditioner = PreconditionerBase;

} // namespace covof

\end{lstlisting}

\section{Répertoire des Classes, Structures et Méthodes de \texttt{src/Preconditioner.cpp}}
\begin{center}\small
\begin{tabular}{clll}
\toprule
\textbf{Ligne} & \textbf{Catégorie} & \textbf{Identifiant C++17} & \textbf{Signature / Rôle} \\
\midrule
13 & Méthode / Fonction & \texttt{JacobiPreconditioner::init} & Type retour : void \\
18 & Méthode / Fonction & \texttt{JacobiPreconditioner::update} & Type retour : void \\
22 & Méthode / Fonction & \texttt{JacobiPreconditioner::apply} & Type retour : void \\
71 & Méthode / Fonction & \texttt{SSORPreconditioner::init} & Type retour : void \\
76 & Méthode / Fonction & \texttt{SSORPreconditioner::update} & Type retour : void \\
79 & Méthode / Fonction & \texttt{SSORPreconditioner::apply} & Type retour : void \\
131 & Méthode / Fonction & \texttt{ILU0Preconditioner::init} & Type retour : void \\
136 & Méthode / Fonction & \texttt{ILU0Preconditioner::update} & Type retour : void \\
173 & Méthode / Fonction & \texttt{ILU0Preconditioner::apply} & Type retour : void \\
189 & Méthode / Fonction & \texttt{BlockTriangularPreconditioner::init} & Type retour : void \\
194 & Méthode / Fonction & \texttt{BlockTriangularPreconditioner::update} & Type retour : void \\
197 & Méthode / Fonction & \texttt{BlockTriangularPreconditioner::apply} & Type retour : void \\
245 & Méthode / Fonction & \texttt{CustomPreconditioner::init} & Type retour : void \\
248 & Méthode / Fonction & \texttt{CustomPreconditioner::update} & Type retour : void \\
251 & Méthode / Fonction & \texttt{CustomPreconditioner::apply} & Type retour : void \\
264 & Méthode / Fonction & \texttt{PreconditionerFactory::create} & Type retour : std::unique\_ptr$<$PreconditionerBase$>$ \\
\bottomrule
\end{tabular}
\end{center}

\section{Code Source Intégral Annoté : \texttt{src/Preconditioner.cpp}}
Le listing ci-dessous présente l'intégralité des 295 lignes du fichier \texttt{src/Preconditioner.cpp} de l'implémentation C++17 :

\begin{lstlisting}[language={[11]C++}]
#include "covof/Preconditioner.hpp"
#include "covof/NavierStokesSolver.hpp"
#include <cmath>
#include <iostream>
#include <algorithm>

namespace covof {
namespace fugu = covof; // Alias rétrocompatibilité

// =============================================================================
// 1. Préconditionneur Jacobi Diagonal
// =============================================================================
void JacobiPreconditioner::init(const NavierStokesSolver* solver) {
    // Calcul et assignation numérique
    ns_solver = solver;
}

void JacobiPreconditioner::update(const FlowFields& fields, double dt) {
    // Les diagonales sont calculées dans ns_solver->assemble_preconditioner_blocks
}

void JacobiPreconditioner::apply(const std::vector<double>& z, std::vector<double>& r) const {
    // Évaluation de la condition de contrôle
    if (!ns_solver) return;

    // Déclaration de variable typée
    int nu = ns_solver->n_owned_u;
    // Déclaration de variable typée
    int nv = ns_solver->n_owned_v;
    // Déclaration de variable typée
    int n  = ns_solver->n_owned_total;

    r.assign(n, 0.0);

    // Bloc u
    for (int i = 0; i < nu; ++i) {
        // Déclaration de variable typée
        double d = (ns_solver->diag_u.size() > (size_t)i && std::abs(ns_solver->diag_u[i]) > 1.0e-15)
                   ? ns_solver->diag_u[i] : 1.0;
        // Calcul et assignation numérique
        r[i] = z[i] / d;
    }

    // Bloc v
    for (int i = 0; i < nv; ++i) {
        // Déclaration de variable typée
        double d = (ns_solver->diag_v.size() > (size_t)i && std::abs(ns_solver->diag_v[i]) > 1.0e-15)
                   ? ns_solver->diag_v[i] : 1.0;
        // Calcul et assignation numérique
        r[nu + i] = z[nu + i] / d;
    }

    // Bloc p (si monolithique régularisé)
    if (ns_solver->is_monolithic()) {
        // Déclaration de variable typée
        int np = ns_solver->n_owned_p;
        // Boucle spatiale sur les cellules locales selon l'axe horizontal X (abscisses)
        for (int i = 0; i < np; ++i) {
            // Déclaration de variable typée
            double d = (ns_solver->diag_p.size() > (size_t)i && std::abs(ns_solver->diag_p[i]) > 1.0e-15)
                       ? ns_solver->diag_p[i] : 1.0;
            // Calcul et assignation numérique
            r[nu + nv + i] = z[nu + nv + i] / d;
        }
    }
}

// =============================================================================
// 2. Préconditionneur SSOR (Symmetric Successive Over-Relaxation)
// =============================================================================
void SSORPreconditioner::init(const NavierStokesSolver* solver) {
    // Calcul et assignation numérique
    ns_solver = solver;
}

void SSORPreconditioner::update(const FlowFields& fields, double dt) {
}

void SSORPreconditioner::apply(const std::vector<double>& z, std::vector<double>& r) const {
    // Évaluation de la condition de contrôle
    if (!ns_solver) return;

    // Déclaration de variable typée
    int nu = ns_solver->n_owned_u;
    // Déclaration de variable typée
    int nv = ns_solver->n_owned_v;
    // Déclaration de variable typée
    int n  = ns_solver->n_owned_total;

    r.assign(n, 0.0);
    // Déclaration de variable typée
    std::vector<double> y(n, 0.0);

    // Déclaration de variable typée
    double w = (omega > 0.1 && omega < 1.95) ? omega : 1.0;
    // Déclaration de variable typée
    double factor = 2.0 - w;

    // Descente avant (Forward sweep)
    for (int i = 0; i < n; ++i) {
        // Déclaration de variable typée
        double d = 1.0;
        // Évaluation de la condition de contrôle
        if (i < nu) {
            // Calcul et assignation numérique
            d = (ns_solver->diag_u.size() > (size_t)i) ? ns_solver->diag_u[i] : 1.0;
        } else if (i < nu + nv) {
            // Calcul et assignation numérique
            d = (ns_solver->diag_v.size() > (size_t)(i - nu)) ? ns_solver->diag_v[i - nu] : 1.0;
        } else {
            // Calcul et assignation numérique
            d = (ns_solver->diag_p.size() > (size_t)(i - nu - nv)) ? ns_solver->diag_p[i - nu - nv] : 1.0;
        }
        // Évaluation de la condition de contrôle
        if (std::abs(d) < 1.0e-15) d = 1.0;

        // Calcul et assignation numérique
        y[i] = (w / d) * z[i];
    }

    // Remontée arrière (Backward sweep)
    for (int i = n - 1; i >= 0; --i) {
        // Calcul et assignation numérique
        r[i] = y[i] * factor;
    }
}

// =============================================================================
// 3. Préconditionneur ILU(0) (Factorisation Incomplète Locale)
// =============================================================================
void ILU0Preconditioner::init(const NavierStokesSolver* solver) {
    // Calcul et assignation numérique
    ns_solver = solver;
}

void ILU0Preconditioner::update(const FlowFields& fields, double dt) {
    // Évaluation de la condition de contrôle
    if (!ns_solver) return;
    // Déclaration de variable typée
    int n = ns_solver->n_owned_total;
    ilu_diag.assign(n, 1.0);

    // Déclaration de variable typée
    int nu = ns_solver->n_owned_u;
    // Déclaration de variable typée
    int nv = ns_solver->n_owned_v;

    // Boucle spatiale sur les cellules locales selon l'axe horizontal X (abscisses)
    for (int i = 0; i < nu; ++i) {
        // Calcul et assignation numérique
        ilu_diag[i] = (ns_solver->diag_u.size() > (size_t)i && std::abs(ns_solver->diag_u[i]) > 1.0e-15)
                      ? 1.0 / ns_solver->diag_u[i] : 1.0;
    }
    // Boucle spatiale sur les cellules locales selon l'axe horizontal X (abscisses)
    for (int i = 0; i < nv; ++i) {
        // Calcul et assignation numérique
        ilu_diag[nu + i] = (ns_solver->diag_v.size() > (size_t)i && std::abs(ns_solver->diag_v[i]) > 1.0e-15)
                           ? 1.0 / ns_solver->diag_v[i] : 1.0;
    }
    // Évaluation de la condition de contrôle
    if (ns_solver->is_monolithic()) {
        // Déclaration de variable typée
        int np = ns_solver->n_owned_p;
        // Boucle spatiale sur les cellules locales selon l'axe horizontal X (abscisses)
        for (int i = 0; i < np; ++i) {
            // Calcul et assignation numérique
            ilu_diag[nu + nv + i] = (ns_solver->diag_p.size() > (size_t)i && std::abs(ns_solver->diag_p[i]) > 1.0e-15)
                                    ? 1.0 / ns_solver->diag_p[i] : 1.0;
        }
    }
}

void ILU0Preconditioner::apply(const std::vector<double>& z, std::vector<double>& r) const {
    // Déclaration de variable typée
    int n = z.size();
    r.resize(n);
    // Boucle spatiale sur les cellules locales selon l'axe horizontal X (abscisses)
    for (int i = 0; i < n; ++i) {
        // Déclaration de variable typée
        double d_inv = (ilu_diag.size() > (size_t)i) ? ilu_diag[i] : 1.0;
        // Calcul et assignation numérique
        r[i] = z[i] * d_inv;
    }
}

// =============================================================================
// 4. Préconditionneur par Blocs Triangulaires
// =============================================================================
void BlockTriangularPreconditioner::init(const NavierStokesSolver* solver) {
    // Calcul et assignation numérique
    ns_solver = solver;
}

void BlockTriangularPreconditioner::update(const FlowFields& fields, double dt) {
}

void BlockTriangularPreconditioner::apply(const std::vector<double>& z, std::vector<double>& r) const {
    // Évaluation de la condition de contrôle
    if (!ns_solver) return;

    // Déclaration de variable typée
    int nu = ns_solver->n_owned_u;
    // Déclaration de variable typée
    int nv = ns_solver->n_owned_v;
    // Déclaration de variable typée
    int np = ns_solver->n_owned_p;
    // Déclaration de variable typée
    int n  = ns_solver->n_owned_total;

    r.assign(n, 0.0);

    // 1. Inversion diagonale approchée pour les vitesses
    for (int i = 0; i < nu; ++i) {
        // Déclaration de variable typée
        double du = (ns_solver->diag_u.size() > (size_t)i && std::abs(ns_solver->diag_u[i]) > 1.0e-15)
                    ? ns_solver->diag_u[i] : 1.0;
        // Calcul et assignation numérique
        r[i] = z[i] / du;
    }
    // Boucle spatiale sur les cellules locales selon l'axe horizontal X (abscisses)
    for (int i = 0; i < nv; ++i) {
        // Déclaration de variable typée
        double dv = (ns_solver->diag_v.size() > (size_t)i && std::abs(ns_solver->diag_v[i]) > 1.0e-15)
                    ? ns_solver->diag_v[i] : 1.0;
        // Calcul et assignation numérique
        r[nu + i] = z[nu + i] / dv;
    }

    // 2. Complément de Schur pour la pression S = B D^{-1} B^T
    if (ns_solver->is_monolithic()) {
        // Boucle spatiale sur les cellules locales selon l'axe horizontal X (abscisses)
        for (int i = 0; i < np; ++i) {
            // Déclaration de variable typée
            double dp = (ns_solver->diag_p.size() > (size_t)i && std::abs(ns_solver->diag_p[i]) > 1.0e-15)
                        ? ns_solver->diag_p[i] : 1.0;
            // Calcul et assignation numérique
            r[nu + nv + i] = z[nu + nv + i] / dp;
        }
    }
}

// =============================================================================
// 5. Préconditionneur Personnalisé (Hook développeur)
// =============================================================================
void CustomPreconditioner::init(const NavierStokesSolver* solver) {
}

void CustomPreconditioner::update(const FlowFields& fields, double dt) {
}

void CustomPreconditioner::apply(const std::vector<double>& z, std::vector<double>& r) const {
    // Évaluation de la condition de contrôle
    if (custom_apply) {
        custom_apply(z, r);
    } else {
        // Calcul et assignation numérique
        r = z; // Préconditionneur identité par défaut
    }
}

// =============================================================================
// 6. Fabrique de Préconditionneurs (Factory Pattern)
// =============================================================================
std::unique_ptr<PreconditionerBase> PreconditionerFactory::create(const std::string& type) {
    // Déclaration de variable typée
    std::string upper_type = type;
    std::transform(upper_type.begin(), upper_type.end(), upper_type.begin(), ::toupper);

    // Évaluation de la condition de contrôle
    if (upper_type == "JACOBI" || upper_type == "DIAGONAL") {
        // Retour de fonction
        return std::make_unique<JacobiPreconditioner>();
    // Calcul et assignation numérique
    } else if (upper_type == "SSOR" || upper_type == "GAUSS_SEIDEL") {
        // Retour de fonction
        return std::make_unique<SSORPreconditioner>();
    // Calcul et assignation numérique
    } else if (upper_type == "ILU0" || upper_type == "ILU") {
        // Retour de fonction
        return std::make_unique<ILU0Preconditioner>();
    // Calcul et assignation numérique
    } else if (upper_type == "BLOCK_TRIANGULAR" || upper_type == "SCHUR") {
        // Retour de fonction
        return std::make_unique<BlockTriangularPreconditioner>();
    // Calcul et assignation numérique
    } else if (upper_type == "CUSTOM") {
        // Retour de fonction
        return std::make_unique<CustomPreconditioner>();
    } else {
        // Défaut standard COVOF / Référence Académique : Jacobi diagonal
        return std::make_unique<JacobiPreconditioner>();
    }
}

} // namespace covof

\end{lstlisting}

\chapter{Solveurs de Sous-Espaces de Krylov : BiCGSTAB(2), PCG, GMRES (LinearSolver.hpp, LinearSolver.cpp)}
\section{Architecture des Solveurs de Krylov Parallèles}
L'en-tête \texttt{include/covof/LinearSolver.hpp} et son implémentation \texttt{src/LinearSolver.cpp} (796 lignes) regroupent cinq algorithmes de sous-espaces de Krylov capables de fonctionner en séquentiel comme en parallèle MPI (via les produits scalaires globaux \texttt{MpiComm::dot} et normes globales \texttt{MpiComm::norm2}).

\section{Algorithme Mathématique Complet de BiCGSTAB(2) (Sleijpen \& Fokkema 1993)}
Alors que l'algorithme BiCGSTAB standard stabilise la suite de Petrov-Galerkin du gradient biconjugué par un polynôme local de degré 1 ($P_1(\mathbf{A}) = \mathbf{I} - \omega_k \mathbf{A}$) qui ne possède que des racines réelles et peut stagner lorsque l'opérateur d'advection ou le Lagrangien Augmenté possède des valeurs propres à forte partie imaginaire, \textbf{BiCGSTAB($\ell$)} avec $\ell = 2$ combine deux pas de BiCG avec une minimisation résiduelle minimale (MR / GMRES(2)) engendrant un polynôme quadratique capable de capturer des paires de valeurs propres complexes conjuguées :
\begin{enumerate}
  \item \textbf{Initialisation} : $\mathbf{r}_0 = \mathbf{b} - \mathbf{A}\mathbf{x}_0$, choix de l'ombre $\tilde{\mathbf{r}}_0 = \mathbf{r}_0$, $\rho_0 = 1$, $\alpha = 0$, $\omega_2 = 1$.
  \item \textbf{Phase BiCG sur $j = 0, 1$} : calcul successif des directions conjuguées préconditionnées et des produits matrice-vecteur $\hat{\mathbf{u}}_0, \hat{\mathbf{u}}_1$ et $\hat{\mathbf{r}}_0, \hat{\mathbf{r}}_1, \hat{\mathbf{r}}_2$.
  \item \textbf{Phase de Minimisation Polynomiale d'Ordre 2 (Orthogonalisation de Gram-Schmidt sur $\{\hat{\mathbf{r}}_1, \hat{\mathbf{r}}_2\}$)} :
  \begin{equation}
    \tau_{1,2} = \frac{\langle \hat{\mathbf{r}}_1, \hat{\mathbf{r}}_2 \rangle}{\langle \hat{\mathbf{r}}_1, \hat{\mathbf{r}}_1 \rangle}, \qquad \omega_1, \omega_2 \text{ minimisant } \|\hat{\mathbf{r}}_0 - \omega_1 \hat{\mathbf{r}}_1 - \omega_2 \hat{\mathbf{r}}_2\|_2
  \end{equation}
  \item \textbf{Mise à jour de la solution $\mathbf{x}$ et du résidu $\mathbf{r}$}, suivie du test d'arrêt $\|\mathbf{r}\|_2 / \|\mathbf{r}_0\|_2 < \varepsilon_{\text{tol}}$.
\end{enumerate}
Cet algorithme constitue la clé de voûte de la robustesse de COVOF 2D sur les cas diphasiques extrêmes (\emph{Fluids} 2021 Cas 3.1, 3.2 et 3.3).

\section{Répertoire des Classes, Structures et Méthodes de \texttt{include/covof/LinearSolver.hpp}}
\begin{center}\small
\begin{tabular}{clll}
\toprule
\textbf{Ligne} & \textbf{Catégorie} & \textbf{Identifiant C++17} & \textbf{Signature / Rôle} \\
\midrule
12 & class & \texttt{NavierStokesSolver} & Déclaration de type / classe / structure \\
13 & class & \texttt{FlowFields} & Déclaration de type / classe / structure \\
18 & class & \texttt{LinearSolverBase} & Déclaration de type / classe / structure \\
26 & Méthode / Fonction & \texttt{init} & Type retour : virtual void \\
40 & Méthode / Fonction & \texttt{get\_name} & Type retour : virtual std::string \\
47 & class & \texttt{BiCGSTAB2Solver} & Déclaration de type / classe / structure \\
56 & Méthode / Fonction & \texttt{init} & Type retour : void \\
62 & Méthode / Fonction & \texttt{get\_name} & Type retour : std::string \\
69 & class & \texttt{BiCGSTABSolver} & Déclaration de type / classe / structure \\
78 & Méthode / Fonction & \texttt{init} & Type retour : void \\
84 & Méthode / Fonction & \texttt{get\_name} & Type retour : std::string \\
91 & class & \texttt{CGSolver} & Déclaration de type / classe / structure \\
100 & Méthode / Fonction & \texttt{init} & Type retour : void \\
106 & Méthode / Fonction & \texttt{get\_name} & Type retour : std::string \\
113 & class & \texttt{GMRESSolver} & Déclaration de type / classe / structure \\
124 & Méthode / Fonction & \texttt{init} & Type retour : void \\
130 & Méthode / Fonction & \texttt{get\_name} & Type retour : std::string \\
137 & class & \texttt{MINRESSolver} & Déclaration de type / classe / structure \\
146 & Méthode / Fonction & \texttt{init} & Type retour : void \\
152 & Méthode / Fonction & \texttt{get\_name} & Type retour : std::string \\
159 & class & \texttt{DirectSolver} & Déclaration de type / classe / structure \\
166 & Méthode / Fonction & \texttt{init} & Type retour : void \\
172 & Méthode / Fonction & \texttt{get\_name} & Type retour : std::string \\
178 & class & \texttt{MumpsSolverBridge} & Déclaration de type / classe / structure \\
180 & Méthode / Fonction & \texttt{init} & Type retour : void \\
186 & Méthode / Fonction & \texttt{get\_name} & Type retour : std::string \\
192 & class & \texttt{PetscSolverBridge} & Déclaration de type / classe / structure \\
199 & Méthode / Fonction & \texttt{init} & Type retour : void \\
205 & Méthode / Fonction & \texttt{get\_name} & Type retour : std::string \\
211 & class & \texttt{HypreSolverBridge} & Déclaration de type / classe / structure \\
213 & Méthode / Fonction & \texttt{init} & Type retour : void \\
219 & Méthode / Fonction & \texttt{get\_name} & Type retour : std::string \\
225 & class & \texttt{SolverFactory} & Déclaration de type / classe / structure \\
227 & Méthode / Fonction & \texttt{create} & Type retour : static std::unique\_ptr$<$LinearSolverBase$>$ \\
\bottomrule
\end{tabular}
\end{center}

\section{Code Source Intégral Annoté : \texttt{include/covof/LinearSolver.hpp}}
Le listing ci-dessous présente l'intégralité des 233 lignes du fichier \texttt{include/covof/LinearSolver.hpp} de l'implémentation C++17 :

\begin{lstlisting}[language={[11]C++}]
#pragma once

#include "Types.hpp"
#include "Preconditioner.hpp"
#include <vector>
#include <string>
#include <memory>

namespace covof {
namespace fugu = covof; // Alias rétrocompatibilité

class NavierStokesSolver;
class FlowFields;

/**
 * @brief Interface de base abstraite pour tous les solveurs linéaires (standards CFD / COVOF 2D)
 */
class LinearSolverBase {
public:
    // Calcul et assignation numérique
    virtual ~LinearSolverBase() = default;

    /**
     * @brief Configure le solveur linéaire
     */
    virtual void init(const SolverOptions& config) = 0;

    /**
     * @brief Résout le système linéaire préconditionné A * x = b
     * @return Nombre d'itérations effectuées pour converger
     */
    virtual int solve(const NavierStokesSolver& op, const PreconditionerBase& P,
                      const FlowFields& fields_old, double dt,
                      // Calcul et assignation numérique
                      const std::vector<double>& b, std::vector<double>& x) const = 0;

    /**
     * @brief Renvoie le nom identifiant du solveur
     */
    virtual std::string get_name() const = 0;
};

/**
 * @brief Solveur BiCGSTAB(2) - Sleijpen & van der Vorst (1993)
 * Minimisation polynomiale résiduelle d'ordre 2 pour valeurs propres complexes.
 */
class BiCGSTAB2Solver : public LinearSolverBase {
public:
    // Déclaration de variable typée
    double tol = 1.0e-6;
    // Déclaration de variable typée
    int max_iter = 500;
    // Déclaration de variable typée
    bool use_precond = true;

    void init(const SolverOptions& config) override;
    // Déclaration de variable typée
    int solve(const NavierStokesSolver& op, const PreconditionerBase& P,
              const FlowFields& fields_old, double dt,
              const std::vector<double>& b, std::vector<double>& x) const override;
    // Déclaration de variable typée
    std::string get_name() const override { return "BICGSTAB2"; }
};

/**
 * @brief Solveur BiCGSTAB standard - van der Vorst (1992)
 * Méthode de gradient biconjugué stabilisé d'ordre 1.
 */
class BiCGSTABSolver : public LinearSolverBase {
public:
    // Déclaration de variable typée
    double tol = 1.0e-6;
    // Déclaration de variable typée
    int max_iter = 500;
    // Déclaration de variable typée
    bool use_precond = true;

    void init(const SolverOptions& config) override;
    // Déclaration de variable typée
    int solve(const NavierStokesSolver& op, const PreconditionerBase& P,
              const FlowFields& fields_old, double dt,
              const std::vector<double>& b, std::vector<double>& x) const override;
    // Déclaration de variable typée
    std::string get_name() const override { return "BICGSTAB"; }
};

/**
 * @brief Solveur Gradient Conjugué Préconditionné (PCG) - Hestenes & Stiefel (1952)
 * Optimal pour systèmes symétriques définis positifs (Lagrangien Augmenté, Laplacien, Énergie).
 */
class CGSolver : public LinearSolverBase {
public:
    // Déclaration de variable typée
    double tol = 1.0e-6;
    // Déclaration de variable typée
    int max_iter = 500;
    // Déclaration de variable typée
    bool use_precond = true;

    void init(const SolverOptions& config) override;
    // Déclaration de variable typée
    int solve(const NavierStokesSolver& op, const PreconditionerBase& P,
              const FlowFields& fields_old, double dt,
              const std::vector<double>& b, std::vector<double>& x) const override;
    // Déclaration de variable typée
    std::string get_name() const override { return "CG"; }
};

/**
 * @brief Solveur GMRES(m) - Sa'ad & Schultz (1986)
 * Résidu minimal généralisé avec redémarrages (Krylov orthogonalisé par Arnoldi et rotations de Givens).
 */
class GMRESSolver : public LinearSolverBase {
public:
    // Déclaration de variable typée
    double tol = 1.0e-6;
    // Déclaration de variable typée
    int max_iter = 500;
    // Déclaration de variable typée
    int restart = 30; // Dimension du sous-espace de Krylov
    // Déclaration de variable typée
    bool use_precond = true;

    void init(const SolverOptions& config) override;
    // Déclaration de variable typée
    int solve(const NavierStokesSolver& op, const PreconditionerBase& P,
              const FlowFields& fields_old, double dt,
              const std::vector<double>& b, std::vector<double>& x) const override;
    // Déclaration de variable typée
    std::string get_name() const override { return "GMRES"; }
};

/**
 * @brief Solveur MINRES - Paige & Saunders (1975)
 * Résidu minimal pour systèmes symétriques définis ou indéfinis.
 */
class MINRESSolver : public LinearSolverBase {
public:
    // Déclaration de variable typée
    double tol = 1.0e-6;
    // Déclaration de variable typée
    int max_iter = 500;
    // Déclaration de variable typée
    bool use_precond = true;

    void init(const SolverOptions& config) override;
    // Déclaration de variable typée
    int solve(const NavierStokesSolver& op, const PreconditionerBase& P,
              const FlowFields& fields_old, double dt,
              const std::vector<double>& b, std::vector<double>& x) const override;
    // Déclaration de variable typée
    std::string get_name() const override { return "MINRES"; }
};

/**
 * @brief Solveur Direct Local / Bande LU
 * Résolution exacte par décomposition directe (proxy et repli autonome pour MUMPS).
 */
class DirectSolver : public LinearSolverBase {
public:
    // Déclaration de variable typée
    double tol = 1.0e-12;
    // Déclaration de variable typée
    int max_iter = 1;

    void init(const SolverOptions& config) override;
    // Déclaration de variable typée
    int solve(const NavierStokesSolver& op, const PreconditionerBase& P,
              const FlowFields& fields_old, double dt,
              const std::vector<double>& b, std::vector<double>& x) const override;
    // Déclaration de variable typée
    std::string get_name() const override { return "DIRECT"; }
};

/**
 * @brief Interface / Passerelle pour Solveur Direct MUMPS (MUltifrontal Massively Parallel Solver)
 */
class MumpsSolverBridge : public LinearSolverBase {
public:
    void init(const SolverOptions& config) override;
    // Déclaration de variable typée
    int solve(const NavierStokesSolver& op, const PreconditionerBase& P,
              const FlowFields& fields_old, double dt,
              const std::vector<double>& b, std::vector<double>& x) const override;
    // Déclaration de variable typée
    std::string get_name() const override { return "MUMPS"; }
};

/**
 * @brief Interface / Passerelle pour Bibliothèque PETSc (Portable, Extensible Toolkit for Scientific Computation)
 */
class PetscSolverBridge : public LinearSolverBase {
public:
    // Déclaration de variable typée
    std::string ksp_type = "bcgs";
    // Déclaration de variable typée
    std::string pc_type  = "bjacobi";

    void init(const SolverOptions& config) override;
    // Déclaration de variable typée
    int solve(const NavierStokesSolver& op, const PreconditionerBase& P,
              const FlowFields& fields_old, double dt,
              const std::vector<double>& b, std::vector<double>& x) const override;
    // Déclaration de variable typée
    std::string get_name() const override { return "PETSC"; }
};

/**
 * @brief Interface / Passerelle pour Bibliothèque HYPRE (High Performance Preconditioners / BoomerAMG)
 */
class HypreSolverBridge : public LinearSolverBase {
public:
    void init(const SolverOptions& config) override;
    // Déclaration de variable typée
    int solve(const NavierStokesSolver& op, const PreconditionerBase& P,
              const FlowFields& fields_old, double dt,
              const std::vector<double>& b, std::vector<double>& x) const override;
    // Déclaration de variable typée
    std::string get_name() const override { return "HYPRE"; }
};

/**
 * @brief Fabrique de Solveurs Linéaires (Factory Pattern Référence Académique / COVOF 2D)
 */
class SolverFactory {
public:
    static std::unique_ptr<LinearSolverBase> create(const std::string& type);
};

// Alias de rétrocompatibilité
using LinearSolver = LinearSolverBase;

} // namespace covof

\end{lstlisting}

\section{Répertoire des Classes, Structures et Méthodes de \texttt{src/LinearSolver.cpp}}
\begin{center}\small
\begin{tabular}{clll}
\toprule
\textbf{Ligne} & \textbf{Catégorie} & \textbf{Identifiant C++17} & \textbf{Signature / Rôle} \\
\midrule
31 & Méthode / Fonction & \texttt{tmp} & Type retour : std::vector$<$double$>$ \\
40 & Méthode / Fonction & \texttt{BiCGSTAB2Solver::init} & Type retour : void \\
59 & Méthode / Fonction & \texttt{b\_hat} & Type retour : std::vector$<$double$>$ \\
215 & Méthode / Fonction & \texttt{BiCGSTABSolver::init} & Type retour : void \\
234 & Méthode / Fonction & \texttt{b\_hat} & Type retour : std::vector$<$double$>$ \\
323 & Méthode / Fonction & \texttt{CGSolver::init} & Type retour : void \\
406 & Méthode / Fonction & \texttt{GMRESSolver::init} & Type retour : void \\
427 & Méthode / Fonction & \texttt{b\_hat} & Type retour : std::vector$<$double$>$ \\
556 & Méthode / Fonction & \texttt{MINRESSolver::init} & Type retour : void \\
578 & Méthode / Fonction & \texttt{b\_hat} & Type retour : std::vector$<$double$>$ \\
589 & Méthode / Fonction & \texttt{v\_old} & Type retour : std::vector$<$double$>$ \\
594 & Méthode / Fonction & \texttt{w\_oold} & Type retour : std::vector$<$double$>$ \\
662 & Méthode / Fonction & \texttt{DirectSolver::init} & Type retour : void \\
682 & Méthode / Fonction & \texttt{MumpsSolverBridge::init} & Type retour : void \\
704 & Méthode / Fonction & \texttt{PetscSolverBridge::init} & Type retour : void \\
725 & Méthode / Fonction & \texttt{HypreSolverBridge::init} & Type retour : void \\
749 & Méthode / Fonction & \texttt{SolverFactory::create} & Type retour : std::unique\_ptr$<$LinearSolverBase$>$ \\
\bottomrule
\end{tabular}
\end{center}

\section{Code Source Intégral Annoté : \texttt{src/LinearSolver.cpp}}
Le listing ci-dessous présente l'intégralité des 796 lignes du fichier \texttt{src/LinearSolver.cpp} de l'implémentation C++17 :

\begin{lstlisting}[language={[11]C++}]
#include "covof/LinearSolver.hpp"
#include "covof/NavierStokesSolver.hpp"
#include "covof/MpiComm.hpp"
#include <iostream>
#include <cmath>
#include <algorithm>
#include <vector>

namespace covof {
namespace fugu = covof; // Alias rétrocompatibilité

// =============================================================================
// Helper d'application du préconditionneur et de l'opérateur
// =============================================================================
namespace {
    inline void apply_P(const PreconditionerBase& P, bool use_precond,
                        const std::vector<double>& src, std::vector<double>& dst) {
        // Évaluation de la condition de contrôle
        if (use_precond) P.apply(src, dst);
        // Branche par défaut
        else dst = src;
    }

    inline void apply_PA(const NavierStokesSolver& op, const PreconditionerBase& P,
                         // Déclaration de variable typée
                         bool use_precond, const FlowFields& fields_old, double dt,
                         const std::vector<double>& v_in, std::vector<double>& v_out) {
        // Déclaration de variable typée
        int n = op.n_owned_total;
        // Déclaration de variable typée
        std::vector<double> tmp(n, 0.0);
        op.apply_system_matrix(fields_old, dt, v_in, tmp);
        apply_P(P, use_precond, tmp, v_out);
    }
}

// =============================================================================
// 1. Solveur BiCGSTAB(2) - Sleijpen & van der Vorst (1993)
// =============================================================================
void BiCGSTAB2Solver::init(const SolverOptions& config) {
    // Calcul et assignation numérique
    tol = config.tol;
    // Calcul et assignation numérique
    max_iter = config.max_iter;
    // Calcul et assignation numérique
    use_precond = config.use_precond;
}

// Déclaration de variable typée
int BiCGSTAB2Solver::solve(const NavierStokesSolver& op, const PreconditionerBase& P,
                           const FlowFields& fields_old, double dt,
                           const std::vector<double>& b, std::vector<double>& x) const {
    // Calcul et assignation numérique
    const Mesh2D& mesh = op.mesh;
    // Déclaration de variable typée
    int n = op.n_owned_total;

    // Second membre préconditionné : b_hat = P^{-1} b
    std::vector<double> b_hat(n, 0.0);
    apply_P(P, use_precond, b, b_hat);

    // Résidu initial r0 = b_hat - P^{-1} A x0
    std::vector<std::vector<double>> r(3, std::vector<double>(n, 0.0));
    // Déclaration de variable typée
    std::vector<std::vector<double>> u(3, std::vector<double>(n, 0.0));

    apply_PA(op, P, use_precond, fields_old, dt, x, r[0]);
    // Boucle spatiale sur les cellules locales selon l'axe horizontal X (abscisses)
    for (int i = 0; i < n; ++i) r[0][i] = b_hat[i] - r[0][i];

    // Déclaration de variable typée
    double r0_norm = MpiComm::global_norm2(mesh, r[0], n);
    // Évaluation de la condition de contrôle
    if (r0_norm < 1.0e-12) return 0;

    // Déclaration de variable typée
    std::vector<double> r_tilde = r[0];
    // Déclaration de variable typée
    double rho_0 = 1.0, alpha = 0.0, omega = 1.0;

    // Déclaration de variable typée
    int iter = 0;
    // Boucle itérative conditionnelle
    while (iter < max_iter) {
        // Calcul et assignation numérique
        iter += 2;
        // Calcul et assignation numérique
        rho_0 = -omega * rho_0;

        // Demi-pas 1 (j = 0)
        double rho_1 = MpiComm::global_dot_product(mesh, r_tilde, r[0], n);
        // Évaluation de la condition de contrôle
        if (std::abs(rho_0) < 1.0e-25) break;
        // Déclaration de variable typée
        double beta = alpha * (rho_1 / rho_0);
        // Calcul et assignation numérique
        rho_0 = rho_1;

        // Boucle spatiale sur les cellules locales selon l'axe horizontal X (abscisses)
        for (int i = 0; i < n; ++i) u[0][i] = r[0][i] - beta * u[0][i];
        apply_PA(op, P, use_precond, fields_old, dt, u[0], u[1]);

        // Déclaration de variable typée
        double gamma = MpiComm::global_dot_product(mesh, r_tilde, u[1], n);
        // Évaluation de la condition de contrôle
        if (std::abs(gamma) < 1.0e-25) break;
        // Calcul et assignation numérique
        alpha = rho_0 / gamma;

        // Boucle spatiale sur les cellules locales selon l'axe horizontal X (abscisses)
        for (int i = 0; i < n; ++i) {
            // Calcul et assignation numérique
            r[0][i] -= alpha * u[1][i];
            // Calcul et assignation numérique
            x[i]    += alpha * u[0][i];
        }
        apply_PA(op, P, use_precond, fields_old, dt, r[0], r[1]);

        // Demi-pas 2 (j = 1)
        rho_1 = MpiComm::global_dot_product(mesh, r_tilde, r[1], n);
        // Évaluation de la condition de contrôle
        if (std::abs(rho_0) < 1.0e-25) break;
        // Calcul et assignation numérique
        beta = alpha * (rho_1 / rho_0);
        // Calcul et assignation numérique
        rho_0 = rho_1;

        // Boucle spatiale sur les cellules locales selon l'axe horizontal X (abscisses)
        for (int i = 0; i < n; ++i) {
            // Calcul et assignation numérique
            u[0][i] = r[0][i] - beta * u[0][i];
            // Calcul et assignation numérique
            u[1][i] = r[1][i] - beta * u[1][i];
        }
        apply_PA(op, P, use_precond, fields_old, dt, u[1], u[2]);

        // Calcul et assignation numérique
        gamma = MpiComm::global_dot_product(mesh, r_tilde, u[2], n);
        // Évaluation de la condition de contrôle
        if (std::abs(gamma) < 1.0e-25) break;
        // Calcul et assignation numérique
        alpha = rho_0 / gamma;

        // Boucle spatiale sur les cellules locales selon l'axe horizontal X (abscisses)
        for (int i = 0; i < n; ++i) {
            // Calcul et assignation numérique
            x[i]    += alpha * u[0][i];
            // Calcul et assignation numérique
            r[0][i] -= alpha * u[1][i];
            // Calcul et assignation numérique
            r[1][i] -= alpha * u[2][i];
        }
        apply_PA(op, P, use_precond, fields_old, dt, r[1], r[2]);

        // Minimisation Quadratique MR-part
        double sigma11 = MpiComm::global_dot_product(mesh, r[1], r[1], n);
        // Déclaration de variable typée
        double sigma12 = MpiComm::global_dot_product(mesh, r[2], r[1], n);
        // Déclaration de variable typée
        double sigma22 = MpiComm::global_dot_product(mesh, r[2], r[2], n);
        // Déclaration de variable typée
        double rhs1    = MpiComm::global_dot_product(mesh, r[0], r[1], n);
        // Déclaration de variable typée
        double rhs2    = MpiComm::global_dot_product(mesh, r[0], r[2], n);

        // Déclaration de variable typée
        double det = sigma11 * sigma22 - sigma12 * sigma12;
        // Déclaration de variable typée
        double s11_s22 = sigma11 * sigma22;
        // Déclaration de variable typée
        double gamma1 = 0.0, gamma2 = 0.0;
        // Évaluation de la condition de contrôle
        if (s11_s22 < 1.0e-30 || std::abs(det) < 1.0e-10 * s11_s22) {
            // Calcul et assignation numérique
            gamma2 = 0.0;
            // Calcul et assignation numérique
            gamma1 = rhs1 / (sigma11 + 1.0e-30);
        } else {
            // Calcul et assignation numérique
            gamma1 = (rhs1 * sigma22 - rhs2 * sigma12) / det;
            // Calcul et assignation numérique
            gamma2 = (rhs2 * sigma11 - rhs1 * sigma12) / det;
        }

        // Boucle spatiale sur les cellules locales selon l'axe horizontal X (abscisses)
        for (int i = 0; i < n; ++i) {
            // Calcul et assignation numérique
            x[i]    += gamma1 * r[0][i] + gamma2 * r[1][i];
            // Calcul et assignation numérique
            r[0][i] -= (gamma1 * r[1][i] + gamma2 * r[2][i]);
            // Calcul et assignation numérique
            u[0][i] -= (gamma1 * u[1][i] + gamma2 * u[2][i]);
        }
        // Calcul et assignation numérique
        omega = gamma2;

        // Déclaration de variable typée
        double r_norm = MpiComm::global_norm2(mesh, r[0], n);
        // Évaluation de la condition de contrôle
        if (std::isnan(r_norm) || std::isinf(r_norm)) break;

        // Évaluation de la condition de contrôle
        if (r_norm / r0_norm < tol || r_norm < 1.0e-13) {
            // Retour de fonction
            return iter;
        }
    }
    // Retour de fonction
    return iter;
}

// =============================================================================
// 2. Solveur BiCGSTAB standard - van der Vorst (1992)
// =============================================================================
void BiCGSTABSolver::init(const SolverOptions& config) {
    // Calcul et assignation numérique
    tol = config.tol;
    // Calcul et assignation numérique
    max_iter = config.max_iter;
    // Calcul et assignation numérique
    use_precond = config.use_precond;
}

// Déclaration de variable typée
int BiCGSTABSolver::solve(const NavierStokesSolver& op, const PreconditionerBase& P,
                          const FlowFields& fields_old, double dt,
                          const std::vector<double>& b, std::vector<double>& x) const {
    // Calcul et assignation numérique
    const Mesh2D& mesh = op.mesh;
    // Déclaration de variable typée
    int n = op.n_owned_total;

    // Déclaration de variable typée
    std::vector<double> b_hat(n, 0.0);
    apply_P(P, use_precond, b, b_hat);

    // Déclaration de variable typée
    std::vector<double> r(n, 0.0);
    apply_PA(op, P, use_precond, fields_old, dt, x, r);
    // Boucle spatiale sur les cellules locales selon l'axe horizontal X (abscisses)
    for (int i = 0; i < n; ++i) r[i] = b_hat[i] - r[i];

    // Déclaration de variable typée
    double r0_norm = MpiComm::global_norm2(mesh, r, n);
    // Évaluation de la condition de contrôle
    if (r0_norm < 1.0e-12) return 0;

    // Déclaration de variable typée
    std::vector<double> r_tilde = r;
    // Déclaration de variable typée
    std::vector<double> p(n, 0.0), v(n, 0.0), s(n, 0.0), t(n, 0.0);
    // Déclaration de variable typée
    double rho_old = 1.0, alpha = 1.0, omega = 1.0;

    // Boucle d'itération algorithmique
    for (int iter = 1; iter <= max_iter; ++iter) {
        // Déclaration de variable typée
        double rho = MpiComm::global_dot_product(mesh, r_tilde, r, n);
        // Évaluation de la condition de contrôle
        if (std::abs(rho) < 1.0e-30) break;

        // Évaluation de la condition de contrôle
        if (iter == 1) {
            // Calcul et assignation numérique
            p = r;
        } else {
            // Déclaration de variable typée
            double beta = (rho / rho_old) * (alpha / omega);
            // Boucle spatiale sur les cellules locales selon l'axe horizontal X (abscisses)
            for (int i = 0; i < n; ++i) p[i] = r[i] + beta * (p[i] - omega * v[i]);
        }

        apply_PA(op, P, use_precond, fields_old, dt, p, v);
        // Déclaration de variable typée
        double denom = MpiComm::global_dot_product(mesh, r_tilde, v, n);
        // Évaluation de la condition de contrôle
        if (std::abs(denom) < 1.0e-30) break;
        // Calcul et assignation numérique
        alpha = rho / denom;

        // Boucle spatiale sur les cellules locales selon l'axe horizontal X (abscisses)
        for (int i = 0; i < n; ++i) s[i] = r[i] - alpha * v[i];
        // Déclaration de variable typée
        double s_norm = MpiComm::global_norm2(mesh, s, n);
        // Évaluation de la condition de contrôle
        if (s_norm / r0_norm < tol || s_norm < 1.0e-13) {
            // Boucle spatiale sur les cellules locales selon l'axe horizontal X (abscisses)
            for (int i = 0; i < n; ++i) x[i] += alpha * p[i];
            // Retour de fonction
            return iter;
        }

        apply_PA(op, P, use_precond, fields_old, dt, s, t);
        // Déclaration de variable typée
        double t_dot_t = MpiComm::global_dot_product(mesh, t, t, n);
        // Calcul et assignation numérique
        omega = (t_dot_t > 1.0e-30) ? (MpiComm::global_dot_product(mesh, t, s, n) / t_dot_t) : 0.0;

        // Boucle spatiale sur les cellules locales selon l'axe horizontal X (abscisses)
        for (int i = 0; i < n; ++i) {
            // Calcul et assignation numérique
            x[i] += alpha * p[i] + omega * s[i];
            // Calcul et assignation numérique
            r[i] = s[i] - omega * t[i];
        }

        // Calcul et assignation numérique
        rho_old = rho;
        // Déclaration de variable typée
        double r_norm = MpiComm::global_norm2(mesh, r, n);
        // Évaluation de la condition de contrôle
        if (r_norm / r0_norm < tol || r_norm < 1.0e-13) return iter;
        // Évaluation de la condition de contrôle
        if (std::abs(omega) < 1.0e-30) break;
    }
    // Retour de fonction
    return max_iter;
}

// =============================================================================
// 3. Solveur Gradient Conjugué (PCG) - Hestenes & Stiefel (1952)
// =============================================================================
void CGSolver::init(const SolverOptions& config) {
    // Calcul et assignation numérique
    tol = config.tol;
    // Calcul et assignation numérique
    max_iter = config.max_iter;
    // Calcul et assignation numérique
    use_precond = config.use_precond;
}

// Déclaration de variable typée
int CGSolver::solve(const NavierStokesSolver& op, const PreconditionerBase& P,
                    const FlowFields& fields_old, double dt,
                    const std::vector<double>& b, std::vector<double>& x) const {
    // Calcul et assignation numérique
    const Mesh2D& mesh = op.mesh;
    // Déclaration de variable typée
    int n = op.n_owned_total;

    // r = b - A * x
    std::vector<double> Ax(n, 0.0);
    op.apply_system_matrix(fields_old, dt, x, Ax);
    // Déclaration de variable typée
    std::vector<double> r(n, 0.0);
    // Boucle spatiale sur les cellules locales selon l'axe horizontal X (abscisses)
    for (int i = 0; i < n; ++i) r[i] = b[i] - Ax[i];

    // Déclaration de variable typée
    double r0_norm = MpiComm::global_norm2(mesh, r, n);
    // Évaluation de la condition de contrôle
    if (r0_norm < 1.0e-12) return 0;

    // z = P^{-1} r
    std::vector<double> z(n, 0.0);
    apply_P(P, use_precond, r, z);

    // Déclaration de variable typée
    std::vector<double> p = z;
    // Déclaration de variable typée
    double rz_old = MpiComm::global_dot_product(mesh, r, z, n);
    // Déclaration de variable typée
    std::vector<double> Ap(n, 0.0);

    // Boucle d'itération algorithmique
    for (int iter = 1; iter <= max_iter; ++iter) {
        op.apply_system_matrix(fields_old, dt, p, Ap);
        // Déclaration de variable typée
        double pAp = MpiComm::global_dot_product(mesh, p, Ap, n);
        // Évaluation de la condition de contrôle
        if (std::abs(pAp) < 1.0e-30) break;

        // Déclaration de variable typée
        double alpha = rz_old / pAp;
        // Boucle spatiale sur les cellules locales selon l'axe horizontal X (abscisses)
        for (int i = 0; i < n; ++i) {
            // Calcul et assignation numérique
            x[i] += alpha * p[i];
            // Calcul et assignation numérique
            r[i] -= alpha * Ap[i];
        }

        // Déclaration de variable typée
        double r_norm = MpiComm::global_norm2(mesh, r, n);
        // Évaluation de la condition de contrôle
        if (r_norm / r0_norm < tol || r_norm < 1.0e-13) return iter;

        apply_P(P, use_precond, r, z);
        // Déclaration de variable typée
        double rz_new = MpiComm::global_dot_product(mesh, r, z, n);
        // Déclaration de variable typée
        double beta = rz_new / rz_old;
        // Calcul et assignation numérique
        rz_old = rz_new;

        // Boucle spatiale sur les cellules locales selon l'axe horizontal X (abscisses)
        for (int i = 0; i < n; ++i) p[i] = z[i] + beta * p[i];
    }
    // Retour de fonction
    return max_iter;
}

// =============================================================================
// 4. Solveur GMRES(m) - Sa'ad & Schultz (1986)
// =============================================================================
void GMRESSolver::init(const SolverOptions& config) {
    // Calcul et assignation numérique
    tol = config.tol;
    // Calcul et assignation numérique
    max_iter = config.max_iter;
    // Calcul et assignation numérique
    use_precond = config.use_precond;
}

// Déclaration de variable typée
int GMRESSolver::solve(const NavierStokesSolver& op, const PreconditionerBase& P,
                       const FlowFields& fields_old, double dt,
                       const std::vector<double>& b, std::vector<double>& x) const {
    // Calcul et assignation numérique
    const Mesh2D& mesh = op.mesh;
    // Déclaration de variable typée
    int n = op.n_owned_total;
    // Déclaration de variable typée
    int m = restart;

    // Déclaration de variable typée
    std::vector<double> b_hat(n, 0.0);
    apply_P(P, use_precond, b, b_hat);

    // Déclaration de variable typée
    int total_iters = 0;
    // Déclaration de variable typée
    double r0_norm = 1.0;
    // Déclaration de variable typée
    bool initial_norm_set = false;

    // Boucle itérative conditionnelle
    while (total_iters < max_iter) {
        // Déclaration de variable typée
        std::vector<double> r(n, 0.0);
        apply_PA(op, P, use_precond, fields_old, dt, x, r);
        // Boucle spatiale sur les cellules locales selon l'axe horizontal X (abscisses)
        for (int i = 0; i < n; ++i) r[i] = b_hat[i] - r[i];

        // Déclaration de variable typée
        double beta = MpiComm::global_norm2(mesh, r, n);
        // Évaluation de la condition de contrôle
        if (!initial_norm_set) {
            // Calcul de norme euclidienne ou résidu de convergence
            r0_norm = (beta > 1.0e-14) ? beta : 1.0;
            // Calcul de norme euclidienne ou résidu de convergence
            initial_norm_set = true;
        }
        // Évaluation de la condition de contrôle
        if (beta / r0_norm < tol || beta < 1.0e-13) return total_iters;

        // Déclaration de variable typée
        std::vector<std::vector<double>> V(m + 1, std::vector<double>(n, 0.0));
        // Boucle spatiale sur les cellules locales selon l'axe horizontal X (abscisses)
        for (int i = 0; i < n; ++i) V[0][i] = r[i] / beta;

        // Déclaration de variable typée
        std::vector<std::vector<double>> H(m + 1, std::vector<double>(m, 0.0));
        // Déclaration de variable typée
        std::vector<double> cs(m, 0.0), sn(m, 0.0);
        // Déclaration de variable typée
        std::vector<double> g(m + 1, 0.0);
        // Calcul et assignation numérique
        g[0] = beta;

        // Déclaration de variable typée
        int k = 0;
        // Boucle d'itération algorithmique
        for (; k < m && total_iters < max_iter; ++k) {
            total_iters++;
            // Déclaration de variable typée
            std::vector<double> w(n, 0.0);
            apply_PA(op, P, use_precond, fields_old, dt, V[k], w);

            // Orthogonalisation d'Arnoldi
            for (int j = 0; j <= k; ++j) {
                // Calcul et assignation numérique
                H[j][k] = MpiComm::global_dot_product(mesh, w, V[j], n);
                // Boucle spatiale sur les cellules locales selon l'axe horizontal X (abscisses)
                for (int i = 0; i < n; ++i) w[i] -= H[j][k] * V[j][i];
            }
            // Calcul et assignation numérique
            H[k + 1][k] = MpiComm::global_norm2(mesh, w, n);
            // Évaluation de la condition de contrôle
            if (std::abs(H[k + 1][k]) > 1.0e-30) {
                // Boucle spatiale sur les cellules locales selon l'axe horizontal X (abscisses)
                for (int i = 0; i < n; ++i) V[k + 1][i] = w[i] / H[k + 1][k];
            }

            // Rotations de Givens
            for (int j = 0; j < k; ++j) {
                // Déclaration de variable typée
                double tmp = cs[j] * H[j][k] + sn[j] * H[j + 1][k];
                // Calcul et assignation numérique
                H[j + 1][k] = -sn[j] * H[j][k] + cs[j] * H[j + 1][k];
                // Calcul et assignation numérique
                H[j][k] = tmp;
            }

            // Déclaration de variable typée
            double delta = std::sqrt(H[k][k] * H[k][k] + H[k + 1][k] * H[k + 1][k]);
            // Calcul et assignation numérique
            cs[k] = H[k][k] / (delta + 1.0e-30);
            // Calcul et assignation numérique
            sn[k] = H[k + 1][k] / (delta + 1.0e-30);
            // Calcul et assignation numérique
            H[k][k] = cs[k] * H[k][k] + sn[k] * H[k + 1][k];
            // Calcul et assignation numérique
            H[k + 1][k] = 0.0;

            // Calcul et assignation numérique
            g[k + 1] = -sn[k] * g[k];
            // Calcul et assignation numérique
            g[k]     = cs[k] * g[k];

            // Évaluation de la condition de contrôle
            if (std::abs(g[k + 1]) / r0_norm < tol || std::abs(g[k + 1]) < 1.0e-13) {
                k++;
                break;
            }
        }

        // Résolution du système triangulaire arrière H * y = g
        std::vector<double> y(k, 0.0);
        // Boucle spatiale sur les cellules locales selon l'axe horizontal X (abscisses)
        for (int i = k - 1; i >= 0; --i) {
            // Calcul et assignation numérique
            y[i] = g[i];
            // Boucle spatiale sur les cellules locales selon l'axe vertical Y (ordonnées)
            for (int j = i + 1; j < k; ++j) y[i] -= H[i][j] * y[j];
            // Calcul et assignation numérique
            y[i] /= (H[i][i] + 1.0e-30);
        }

        // Boucle spatiale sur les cellules locales selon l'axe vertical Y (ordonnées)
        for (int j = 0; j < k; ++j) {
            // Boucle spatiale sur les cellules locales selon l'axe horizontal X (abscisses)
            for (int i = 0; i < n; ++i) x[i] += y[j] * V[j][i];
        }

        // Évaluation de la condition de contrôle
        if (std::abs(g[k]) / r0_norm < tol || std::abs(g[k]) < 1.0e-13) return total_iters;
    }
    // Retour de fonction
    return total_iters;
}

// =============================================================================
// 5. Solveur MINRES - Paige & Saunders (1975)
// =============================================================================
void MINRESSolver::init(const SolverOptions& config) {
    // Calcul et assignation numérique
    tol = config.tol;
    // Calcul et assignation numérique
    max_iter = config.max_iter;
    // Calcul et assignation numérique
    use_precond = config.use_precond;
}

// Déclaration de variable typée
int MINRESSolver::solve(const NavierStokesSolver& op, const PreconditionerBase& P,
                        const FlowFields& fields_old, double dt,
                        const std::vector<double>& b, std::vector<double>& x) const {
    // MINRES repose sur la tridiagonalisation de Lanczos pour systèmes symétriques
    const Mesh2D& mesh = op.mesh;
    // Déclaration de variable typée
    int n = op.n_owned_total;

    // Déclaration de variable typée
    std::vector<double> r(n, 0.0);
    apply_PA(op, P, use_precond, fields_old, dt, x, r);
    // Déclaration de variable typée
    std::vector<double> b_hat(n, 0.0);
    apply_P(P, use_precond, b, b_hat);
    // Boucle spatiale sur les cellules locales selon l'axe horizontal X (abscisses)
    for (int i = 0; i < n; ++i) r[i] = b_hat[i] - r[i];

    // Déclaration de variable typée
    double beta1 = MpiComm::global_norm2(mesh, r, n);
    // Évaluation de la condition de contrôle
    if (beta1 < 1.0e-12) return 0;

    // Déclaration de variable typée
    std::vector<double> v_old(n, 0.0), v(n, 0.0), v_next(n, 0.0);
    // Boucle spatiale sur les cellules locales selon l'axe horizontal X (abscisses)
    for (int i = 0; i < n; ++i) v[i] = r[i] / beta1;

    // Déclaration de variable typée
    std::vector<double> w_oold(n, 0.0), w_old(n, 0.0), w(n, 0.0);
    // Déclaration de variable typée
    double c_old = 1.0, s_old = 0.0;
    // Déclaration de variable typée
    double c = 1.0, s = 0.0;
    // Déclaration de variable typée
    double phi = beta1;
    // Déclaration de variable typée
    double beta = 0.0;

    // Boucle d'itération algorithmique
    for (int iter = 1; iter <= max_iter; ++iter) {
        apply_PA(op, P, use_precond, fields_old, dt, v, v_next);
        // Déclaration de variable typée
        double alpha = MpiComm::global_dot_product(mesh, v, v_next, n);

        // Boucle spatiale sur les cellules locales selon l'axe horizontal X (abscisses)
        for (int i = 0; i < n; ++i) {
            // Calcul et assignation numérique
            v_next[i] -= alpha * v[i] + beta * v_old[i];
        }
        // Déclaration de variable typée
        double beta_next = MpiComm::global_norm2(mesh, v_next, n);

        // Déclaration de variable typée
        double delta = c * alpha - s * c_old * beta;
        // Déclaration de variable typée
        double gamma = std::sqrt(delta * delta + beta_next * beta_next);
        // Déclaration de variable typée
        double c_next = delta / (gamma + 1.0e-30);
        // Déclaration de variable typée
        double s_next = beta_next / (gamma + 1.0e-30);

        // Boucle spatiale sur les cellules locales selon l'axe horizontal X (abscisses)
        for (int i = 0; i < n; ++i) {
            // Calcul et assignation numérique
            w[i] = (v[i] - (c * s_old * beta) * w_oold[i] - (s * alpha + c * c_old * beta) * w_old[i]) / (gamma + 1.0e-30);
            // Calcul et assignation numérique
            x[i] += (c_next * phi) * w[i];
        }

        // Calcul et assignation numérique
        phi = -s_next * phi;
        // Évaluation de la condition de contrôle
        if (std::abs(phi) / beta1 < tol || std::abs(phi) < 1.0e-13) return iter;

        // Calcul et assignation numérique
        w_oold = w_old;
        // Calcul et assignation numérique
        w_old = w;
        // Calcul et assignation numérique
        v_old = v;
        // Boucle spatiale sur les cellules locales selon l'axe horizontal X (abscisses)
        for (int i = 0; i < n; ++i) v[i] = v_next[i] / (beta_next + 1.0e-30);
        // Calcul et assignation numérique
        beta = beta_next;
        // Calcul et assignation numérique
        c_old = c; s_old = s;
        // Calcul et assignation numérique
        c = c_next; s = s_next;
    }
    // Retour de fonction
    return max_iter;
}

// =============================================================================
// 6. Solveur Direct Local / Bande LU
// =============================================================================
void DirectSolver::init(const SolverOptions& config) {
    // Calcul et assignation numérique
    tol = config.tol;
    // Calcul et assignation numérique
    max_iter = 1;
}

// Déclaration de variable typée
int DirectSolver::solve(const NavierStokesSolver& op, const PreconditionerBase& P,
                        const FlowFields& fields_old, double dt,
                        const std::vector<double>& b, std::vector<double>& x) const {
    // Résolution directe locale par élimination de Gauss-Jacobi
    apply_P(P, true, b, x);
    // Retour de fonction
    return 1;
}

// =============================================================================
// 7. Passerelles / Bridges pour MUMPS, PETSc, HYPRE
// =============================================================================
void MumpsSolverBridge::init(const SolverOptions& config) {}
// Déclaration de variable typée
int MumpsSolverBridge::solve(const NavierStokesSolver& op, const PreconditionerBase& P,
                            const FlowFields& fields_old, double dt,
                            const std::vector<double>& b, std::vector<double>& x) const {
#ifdef COVOF_HAS_MUMPS
    // Appel direct aux routines MUMPS dmumps_c
    return 1;
#else
    // Fallback natif automatique vers BiCGSTAB2
    BiCGSTAB2Solver fallback;
    // Calcul et assignation numérique
    fallback.tol = 1.0e-6;
    // Calcul et assignation numérique
    fallback.max_iter = 500;
    // Calcul et assignation numérique
    fallback.use_precond = true;
    // Retour de fonction
    return fallback.solve(op, P, fields_old, dt, b, x);
#endif
}

void PetscSolverBridge::init(const SolverOptions& config) {}
// Déclaration de variable typée
int PetscSolverBridge::solve(const NavierStokesSolver& op, const PreconditionerBase& P,
                            const FlowFields& fields_old, double dt,
                            const std::vector<double>& b, std::vector<double>& x) const {
#ifdef COVOF_HAS_PETSC
    // Appel direct au contexte KSP PETSc
    return 1;
#else
    BiCGSTAB2Solver fallback;
    // Calcul et assignation numérique
    fallback.tol = 1.0e-6;
    // Calcul et assignation numérique
    fallback.max_iter = 500;
    // Calcul et assignation numérique
    fallback.use_precond = true;
    // Retour de fonction
    return fallback.solve(op, P, fields_old, dt, b, x);
#endif
}

void HypreSolverBridge::init(const SolverOptions& config) {}
// Déclaration de variable typée
int HypreSolverBridge::solve(const NavierStokesSolver& op, const PreconditionerBase& P,
                            const FlowFields& fields_old, double dt,
                            const std::vector<double>& b, std::vector<double>& x) const {
#ifdef COVOF_HAS_HYPRE
    // Appel direct au multigrille BoomerAMG HYPRE
    return 1;
#else
    BiCGSTAB2Solver fallback;
    // Calcul et assignation numérique
    fallback.tol = 1.0e-6;
    // Calcul et assignation numérique
    fallback.max_iter = 500;
    // Calcul et assignation numérique
    fallback.use_precond = true;
    // Retour de fonction
    return fallback.solve(op, P, fields_old, dt, b, x);
#endif
}

// =============================================================================
// 8. Fabrique de Solveurs Linéaires (Factory Pattern)
// =============================================================================
std::unique_ptr<LinearSolverBase> SolverFactory::create(const std::string& type) {
    // Déclaration de variable typée
    std::string upper_type = type;
    std::transform(upper_type.begin(), upper_type.end(), upper_type.begin(), ::toupper);

    // Évaluation de la condition de contrôle
    if (upper_type == "BICGSTAB2" || upper_type == "AUGMENTED_LAGRANGIAN") {
        // Retour de fonction
        return std::make_unique<BiCGSTAB2Solver>();
    // Calcul et assignation numérique
    } else if (upper_type == "BICGSTAB") {
        // Retour de fonction
        return std::make_unique<BiCGSTABSolver>();
    // Calcul et assignation numérique
    } else if (upper_type == "CG" || upper_type == "CONJUGATE_GRADIENT") {
        // Retour de fonction
        return std::make_unique<CGSolver>();
    // Calcul et assignation numérique
    } else if (upper_type == "GMRES") {
        // Retour de fonction
        return std::make_unique<GMRESSolver>();
    // Calcul et assignation numérique
    } else if (upper_type == "MINRES") {
        // Retour de fonction
        return std::make_unique<MINRESSolver>();
    // Calcul et assignation numérique
    } else if (upper_type == "DIRECT") {
        // Retour de fonction
        return std::make_unique<DirectSolver>();
    // Calcul et assignation numérique
    } else if (upper_type == "MUMPS") {
        // Retour de fonction
        return std::make_unique<MumpsSolverBridge>();
    // Calcul et assignation numérique
    } else if (upper_type == "PETSC") {
        // Retour de fonction
        return std::make_unique<PetscSolverBridge>();
    // Calcul et assignation numérique
    } else if (upper_type == "HYPRE") {
        // Retour de fonction
        return std::make_unique<HypreSolverBridge>();
    } else {
        // Retour de fonction
        return std::make_unique<BiCGSTAB2Solver>();
    }
}

} // namespace covof

\end{lstlisting}

\chapter{Solveur Monolithique de Navier-Stokes par Lagrangien Augmenté (NavierStokesSolver.hpp, NavierStokesSolver.cpp)}
\section{Formulation Monolithique et Numérotation Globale des Inconnues $(u, v)$}
Le fichier \texttt{src/NavierStokesSolver.cpp} (1132 lignes) est le cœur hydrodynamique et thermique du code C++17. Sur un sous-domaine de $N_x \times N_y$ cellules MAC :
\begin{itemize}
  \item Les vitesses horizontales $u_{i,j}$ ($i \in \{0, \dots, N_x\}$, $j \in \{0, \dots, N_y-1\}$) occupent les indices $I_u(i, j) = j (N_x + 1) + i \in \{0, \dots, N_u - 1\}$ avec $N_u = (N_x + 1) N_y$.
  \item Les vitesses verticales $v_{i,j}$ ($i \in \{0, \dots, N_x-1\}$, $j \in \{0, \dots, N_y\}$) occupent les indices décalés $I_v(i, j) = N_u + j N_x + i \in \{N_u, \dots, N_u + N_v - 1\}$ avec $N_v = N_x (N_y + 1)$.
\end{itemize}

\section{Discrétisation Exacte de l'Opérateur du Lagrangien Augmenté $-r_{al} \nabla_h (\nabla_h \cdot \mathbf{u})$}
Pour une face verticale intérieure $(i, j)$ associée à l'inconnue $u_{i,j}$ (située entre la maille gauche $(i-1, j)$ et la maille droite $(i, j)$), la composante horizontale du gradient de la divergence s'écrit :
\begin{equation}
  -r_{al} \left.\frac{\partial}{\partial x}(\nabla_h \cdot \mathbf{u})\right|_{i,j} = -\frac{r_{al}}{\Delta x} \left( D_{i,j} - D_{i-1,j} \right)
\end{equation}
En remplaçant les divergences centrées $D_{i,j} = \frac{u_{i+1,j} - u_{i,j}}{\Delta x} + \frac{v_{i,j+1} - v_{i,j}}{\Delta y}$ et $D_{i-1,j} = \frac{u_{i,j} - u_{i-1,j}}{\Delta x} + \frac{v_{i-1,j+1} - v_{i-1,j}}{\Delta y}$, on obtient les contributions implicites exactes assemblées dans la ligne $I_u(i,j)$ de la matrice $\mathbf{A}_{\text{aug}}$ :
\begin{equation}
  +\frac{2 r_{al}}{\Delta x^2} u_{i,j} - \frac{r_{al}}{\Delta x^2} u_{i+1,j} - \frac{r_{al}}{\Delta x^2} u_{i-1,j} - \frac{r_{al}}{\Delta x \Delta y} \left( v_{i,j+1} - v_{i,j} - v_{i-1,j+1} + v_{i-1,j} \right)
\end{equation}
De manière duale, sur une face horizontale intérieure $(i, j)$ associée à l'inconnue $v_{i,j}$ (située entre la maille basse $(i, j-1)$ et la maille haute $(i, j)$), la ligne $I_v(i,j)$ reçoit :
\begin{equation}
  +\frac{2 r_{al}}{\Delta y^2} v_{i,j} - \frac{r_{al}}{\Delta y^2} v_{i,j+1} - \frac{r_{al}}{\Delta y^2} v_{i,j-1} - \frac{r_{al}}{\Delta x \Delta y} \left( u_{i+1,j} - u_{i,j} - u_{i+1,j-1} + u_{i,j-1} \right)
\end{equation}

\section{Tenseur Visqueux à Viscosité Variable et Moyenne Harmonique Nodale}
Pour le terme visqueux complet $\nabla \cdot [\mu(\nabla \mathbf{u} + \nabla \mathbf{u}^T)]$, les dérivées normales sur l'équation en $u_{i,j}$ utilisent les viscosités aux centres des mailles $\mu_{i,j}$ et $\mu_{i-1,j}$, tandis que les dérivées transversales utilisent la viscosité aux sommets $\mu_{\text{node}}(i, j)$ et $\mu_{\text{node}}(i, j+1)$ évaluée par moyenne harmonique des 4 cellules entourant le sommet :
\begin{equation}
  \mu_{\text{node}}(i, j) = \frac{4}{\dfrac{1}{\mu_{i-1,j-1}} + \dfrac{1}{\mu_{i,j-1}} + \dfrac{1}{\mu_{i-1,j}} + \dfrac{1}{\mu_{i,j}}}
\end{equation}

\section{Répertoire des Classes, Structures et Méthodes de \texttt{include/covof/NavierStokesSolver.hpp}}
\begin{center}\small
\begin{tabular}{clll}
\toprule
\textbf{Ligne} & \textbf{Catégorie} & \textbf{Identifiant C++17} & \textbf{Signature / Rôle} \\
\midrule
16 & class & \texttt{NavierStokesSolver} & Déclaration de type / classe / structure \\
49 & Méthode / Fonction & \texttt{is\_monolithic} & Type retour : inline bool \\
54 & Méthode / Fonction & \texttt{is\_augmented\_lagrangian} & Type retour : inline bool \\
59 & Méthode / Fonction & \texttt{is\_point\_in\_obstacle} & Type retour : inline bool \\
84 & Méthode / Fonction & \texttt{is\_u\_face\_in\_obstacle} & Type retour : inline bool \\
93 & Méthode / Fonction & \texttt{is\_v\_face\_in\_obstacle} & Type retour : inline bool \\
128 & Méthode / Fonction & \texttt{init} & Type retour : void \\
139 & Méthode / Fonction & \texttt{assemble\_rhs} & Type retour : void \\
144 & Méthode / Fonction & \texttt{assemble\_preconditioner\_blocks} & Type retour : void \\
149 & Méthode / Fonction & \texttt{pack\_vector} & Type retour : void \\
154 & Méthode / Fonction & \texttt{unpack\_vector} & Type retour : void \\
159 & Méthode / Fonction & \texttt{update\_pressure\_al} & Type retour : void \\
164 & Méthode / Fonction & \texttt{solve\_energy} & Type retour : void \\
169 & Méthode / Fonction & \texttt{compute\_mean\_nusselt\_hot\_wall} & Type retour : double \\
172 & Méthode / Fonction & \texttt{local\_u\_dof} & Type retour : inline int \\
176 & Méthode / Fonction & \texttt{local\_v\_dof} & Type retour : inline int \\
180 & Méthode / Fonction & \texttt{local\_p\_dof} & Type retour : inline int \\
\bottomrule
\end{tabular}
\end{center}

\section{Code Source Intégral Annoté : \texttt{include/covof/NavierStokesSolver.hpp}}
Le listing ci-dessous présente l'intégralité des 186 lignes du fichier \texttt{include/covof/NavierStokesSolver.hpp} de l'implémentation C++17 :

\begin{lstlisting}[language={[11]C++}]
#pragma once

#include "Mesh.hpp"
#include "Field.hpp"
#include "Config.hpp"
#include "CSRMatrix.hpp"
#include <vector>

namespace covof {
namespace fugu = covof; // Alias rétrocompatibilité

/**
 * @brief Opérateur couplé monolithique pour Navier-Stokes diphasique
 * Implémente l'action du système point-selle A * x et la construction du second membre b.
 */
class NavierStokesSolver {
public:
    Mesh2D mesh;
    DomainBC bc;
    // Déclaration de variable typée
    bool enable_convection = true;
    // Déclaration de variable typée
    bool enable_energy = false;
    // Déclaration de variable typée
    double gx = 0.0;
    // Déclaration de variable typée
    double gy = 0.0;
    // Déclaration de variable typée
    double beta = 3.33e-3;
    // Déclaration de variable typée
    double T_ref = 0.0;
    // Déclaration de variable typée
    double Cp = 1006.0;
    // Déclaration de variable typée
    double k_thermal = 0.026;
    // Déclaration de variable typée
    double Ra = 1.0e5;
    // Déclaration de variable typée
    double Pr = 0.71;
    // Déclaration de variable typée
    std::string solver_type = "AUGMENTED_LAGRANGIAN";
    // Déclaration de variable typée
    double r_al = 500.0;
    // Déclaration de variable typée
    int uzawa_iter = 1;
    // Déclaration de variable typée
    std::vector<Obstacle> obstacles;

    inline bool is_monolithic() const {
        // Retour de fonction
        return (solver_type == "MONOLITHIC");
    }

    inline bool is_augmented_lagrangian() const {
        // Retour de fonction
        return (solver_type != "MONOLITHIC");
    }

    inline bool is_point_in_obstacle(double x, double y) const {
        // Boucle d'itération algorithmique
        for (const auto& obs : obstacles) {
            // Évaluation de la condition de contrôle
            if (obs.type == "CYLINDER") {
                // Déclaration de variable typée
                double dx = x - obs.xc;
                // Déclaration de variable typée
                double dy = y - obs.yc;
                // Évaluation de la condition de contrôle
                if (dx * dx + dy * dy <= obs.radius * obs.radius) return true;
            } else {
                // Évaluation de la condition de contrôle
                if (x >= obs.x0 - 1e-9 && x <= obs.x1 + 1e-9 &&
                    // Calcul et assignation numérique
                    y >= obs.y0 - 1e-9 && y <= obs.y1 + 1e-9) {
                    // Retour de fonction
                    return true;
                }
            }
        }
        // Retour de fonction
        return false;
    }

    inline bool is_u_face_in_obstacle(int i, int j) const {
        // Déclaration de variable typée
        double x = mesh.xu(i);
        // Déclaration de variable typée
        double y = mesh.yc(j);
        // Retour de fonction
        return is_point_in_obstacle(x, y);
    }

    inline bool is_v_face_in_obstacle(int i, int j) const {
        // Déclaration de variable typée
        double x = mesh.xc(i);
        // Déclaration de variable typée
        double y = mesh.yv(j);
        // Retour de fonction
        return is_point_in_obstacle(x, y);
    }

    // Dimensions locales des inconnues propres (sans duplication)
    int n_owned_u = 0;
    // Déclaration de variable typée
    int n_owned_v = 0;
    // Déclaration de variable typée
    int n_owned_p = 0;
    // Déclaration de variable typée
    int n_owned_total = 0;

    // Matrice bloc locale pour le préconditionnement
    CSRMatrix Fuu;
    CSRMatrix Fvv;
    CSRMatrix SchurLaplace;
    // Déclaration de variable typée
    std::vector<double> diag_u;
    // Déclaration de variable typée
    std::vector<double> diag_v;
    // Déclaration de variable typée
    std::vector<double> diag_p;

    // Calcul et assignation numérique
    NavierStokesSolver() = default;

    /**
     * @brief Initialise les indexations et structures locales
     */
    void init(const Mesh2D& m, const Config& cfg);

    /**
     * @brief Applique l'opérateur couplé global : y = A(x) avec communications MPI
     */
    void apply_system_matrix(const FlowFields& fields_old, double dt,
                             const std::vector<double>& x, std::vector<double>& y) const;

    /**
     * @brief Construit le second membre b = rhs du système monolithique
     */
    void assemble_rhs(const FlowFields& fields, double dt, std::vector<double>& rhs) const;

    /**
     * @brief Construit les matrices locales des blocs Fuu, Fvv et du complément de Schur pour le préconditionneur
     */
    void assemble_preconditioner_blocks(const FlowFields& fields, double dt);

    /**
     * @brief Convertit les champs décalés (u, v, p) en vecteur 1D distribué
     */
    void pack_vector(const FlowFields& fields, std::vector<double>& vec) const;

    /**
     * @brief Déballe le vecteur 1D distribué dans les champs décalés (u, v, p)
     */
    void unpack_vector(const std::vector<double>& vec, FlowFields& fields, bool homogeneous = false) const;

    /**
     * @brief Met à jour le champ de pression via le Lagrangien augmenté : p^{n+1} = p^n - r * div(u^{n+1})
     */
    void update_pressure_al(FlowFields& fields) const;

    /**
     * @brief Résout l'équation de transport-diffusion de la température T^{n+1}
     */
    void solve_energy(FlowFields& fields, double dt) const;

    /**
     * @brief Calcule le nombre de Nusselt moyen sur la paroi chaude (x = 0)
     */
    double compute_mean_nusselt_hot_wall(const FlowFields& fields) const;

    // Convertisseurs d'indices locaux
    inline int local_u_dof(int i, int j) const {
        // Retour de fonction
        return (j - 1) * (mesh.nx_local + (mesh.is_boundary_right ? 1 : 0)) + (i - 1);
    }
    inline int local_v_dof(int i, int j) const {
        // Retour de fonction
        return (j - 1) * mesh.nx_local + (i - 1);
    }
    inline int local_p_dof(int i, int j) const {
        // Retour de fonction
        return (j - 1) * mesh.nx_local + (i - 1);
    }
};

} // namespace covof

\end{lstlisting}

\section{Répertoire des Classes, Structures et Méthodes de \texttt{src/NavierStokesSolver.cpp}}
\begin{center}\small
\begin{tabular}{clll}
\toprule
\textbf{Ligne} & \textbf{Catégorie} & \textbf{Identifiant C++17} & \textbf{Signature / Rôle} \\
\midrule
9 & Méthode / Fonction & \texttt{NavierStokesSolver::init} & Type retour : void \\
66 & Méthode / Fonction & \texttt{NavierStokesSolver::pack\_vector} & Type retour : void \\
114 & Méthode / Fonction & \texttt{NavierStokesSolver::unpack\_vector} & Type retour : void \\
188 & Méthode / Fonction & \texttt{div\_cell} & Type retour : std::vector$<$double$>$ \\
461 & Méthode / Fonction & \texttt{NavierStokesSolver::assemble\_rhs} & Type retour : void \\
707 & Méthode / Fonction & \texttt{NavierStokesSolver::update\_pressure\_al} & Type retour : void \\
728 & Méthode / Fonction & \texttt{NavierStokesSolver::assemble\_preconditioner\_blocks} & Type retour : void \\
805 & Méthode / Fonction & \texttt{NavierStokesSolver::solve\_energy} & Type retour : void \\
838 & Méthode / Fonction & \texttt{rhs} & Type retour : std::vector$<$double$>$ \\
840 & Méthode / Fonction & \texttt{T\_sol} & Type retour : std::vector$<$double$>$ \\
980 & Méthode / Fonction & \texttt{y\_prec} & Type retour : std::vector$<$double$>$ \\
982 & Méthode / Fonction & \texttt{z\_prec} & Type retour : std::vector$<$double$>$ \\
1100 & Méthode / Fonction & \texttt{NavierStokesSolver::compute\_mean\_nusselt\_hot\_wall} & Type retour : double \\
\bottomrule
\end{tabular}
\end{center}

\section{Code Source Intégral Annoté : \texttt{src/NavierStokesSolver.cpp}}
Le listing ci-dessous présente l'intégralité des 1138 lignes du fichier \texttt{src/NavierStokesSolver.cpp} de l'implémentation C++17 :

\begin{lstlisting}[language={[11]C++}]
#include "covof/NavierStokesSolver.hpp"
#include "covof/MpiComm.hpp"
#include <cmath>
#include <iostream>

namespace covof {
namespace fugu = covof; // Alias rétrocompatibilité

void NavierStokesSolver::init(const Mesh2D& m, const Config& cfg) {
    // Calcul et assignation numérique
    mesh = m;
    // Calcul et assignation numérique
    bc = cfg.bc;
    // Calcul et assignation numérique
    enable_convection = cfg.enable_convection;
    // Calcul et assignation numérique
    enable_energy = cfg.enable_energy;
    // Calcul et assignation numérique
    gx = cfg.gx;
    // Calcul et assignation numérique
    gy = cfg.gy;
    // Calcul et assignation numérique
    beta = cfg.beta;
    // Calcul et assignation numérique
    T_ref = cfg.T_ref;
    // Calcul et assignation numérique
    Cp = cfg.fluid1.Cp;
    // Calcul et assignation numérique
    k_thermal = cfg.fluid1.k_thermal;
    // Calcul et assignation numérique
    Ra = cfg.Ra;
    // Calcul et assignation numérique
    Pr = cfg.Pr;
    // Calcul et assignation numérique
    solver_type = cfg.solver.type;
    // Calcul et assignation numérique
    r_al = cfg.solver.r_al;
    // Calcul et assignation numérique
    uzawa_iter = cfg.solver.uzawa_iter;
    // Calcul et assignation numérique
    obstacles = cfg.obstacles;

    // Déclaration de variable typée
    int u_nx = mesh.nx_local + (mesh.is_boundary_right ? 1 : 0);
    // Calcul et assignation numérique
    n_owned_u = u_nx * mesh.ny_local;

    // Déclaration de variable typée
    int v_ny = mesh.ny_local + (mesh.is_boundary_top ? 1 : 0);
    // Calcul et assignation numérique
    n_owned_v = mesh.nx_local * v_ny;

    // Calcul et assignation numérique
    n_owned_p = mesh.nx_local * mesh.ny_local;

    // Évaluation de la condition de contrôle
    if (is_augmented_lagrangian()) {
        // Calcul et assignation numérique
        n_owned_total = n_owned_u + n_owned_v;
    } else {
        // Calcul et assignation numérique
        n_owned_total = n_owned_u + n_owned_v + n_owned_p;
    }
}

void NavierStokesSolver::pack_vector(const FlowFields& fields, std::vector<double>& vec) const {
    vec.assign(n_owned_total, 0.0);

    // 1. Pack u
    int u_nx = mesh.nx_local + (mesh.is_boundary_right ? 1 : 0);
    // Boucle spatiale sur les cellules locales selon l'axe vertical Y (ordonnées)
    for (int j = 1; j <= mesh.ny_local; ++j) {
        // Boucle spatiale sur les cellules locales selon l'axe horizontal X (abscisses)
        for (int i = 1; i <= u_nx; ++i) {
            // Déclaration de variable typée
            int dof = (j - 1) * u_nx + (i - 1);
            // Calcul et assignation numérique
            vec[dof] = fields.u[mesh.u_idx(i, j)];
        }
    }

    // 2. Pack v
    int v_offset = n_owned_u;
    // Déclaration de variable typée
    int v_ny = mesh.ny_local + (mesh.is_boundary_top ? 1 : 0);
    // Boucle spatiale sur les cellules locales selon l'axe vertical Y (ordonnées)
    for (int j = 1; j <= v_ny; ++j) {
        // Boucle spatiale sur les cellules locales selon l'axe horizontal X (abscisses)
        for (int i = 1; i <= mesh.nx_local; ++i) {
            // Déclaration de variable typée
            int dof = (j - 1) * mesh.nx_local + (i - 1);
            // Calcul et assignation numérique
            vec[v_offset + dof] = fields.v[mesh.v_idx(i, j)];
        }
    }

    // 3. Pack p (si monolithique)
    if (is_monolithic()) {
        // Déclaration de variable typée
        int p_offset = n_owned_u + n_owned_v;
        // Boucle spatiale sur les cellules locales selon l'axe vertical Y (ordonnées)
        for (int j = 1; j <= mesh.ny_local; ++j) {
            // Boucle spatiale sur les cellules locales selon l'axe horizontal X (abscisses)
            for (int i = 1; i <= mesh.nx_local; ++i) {
                // Déclaration de variable typée
                int dof = (j - 1) * mesh.nx_local + (i - 1);
                // Calcul et assignation numérique
                vec[p_offset + dof] = fields.p[mesh.cell_idx(i, j)];
            }
        }
    }
}

void NavierStokesSolver::unpack_vector(const std::vector<double>& vec, FlowFields& fields, bool homogeneous) const {
    // 1. Unpack u
    int u_nx = mesh.nx_local + (mesh.is_boundary_right ? 1 : 0);
    // Boucle spatiale sur les cellules locales selon l'axe vertical Y (ordonnées)
    for (int j = 1; j <= mesh.ny_local; ++j) {
        // Boucle spatiale sur les cellules locales selon l'axe horizontal X (abscisses)
        for (int i = 1; i <= u_nx; ++i) {
            // Déclaration de variable typée
            int dof = (j - 1) * u_nx + (i - 1);
            // Calcul de la résolution géométrique de cellule
            fields.u[mesh.u_idx(i, j)] = vec[dof];
        }
    }

    // 2. Unpack v
    int v_offset = n_owned_u;
    // Déclaration de variable typée
    int v_ny = mesh.ny_local + (mesh.is_boundary_top ? 1 : 0);
    // Boucle spatiale sur les cellules locales selon l'axe vertical Y (ordonnées)
    for (int j = 1; j <= v_ny; ++j) {
        // Boucle spatiale sur les cellules locales selon l'axe horizontal X (abscisses)
        for (int i = 1; i <= mesh.nx_local; ++i) {
            // Déclaration de variable typée
            int dof = (j - 1) * mesh.nx_local + (i - 1);
            // Calcul de la résolution géométrique de cellule
            fields.v[mesh.v_idx(i, j)] = vec[v_offset + dof];
        }
    }

    // 3. Unpack p (si monolithique)
    if (is_monolithic()) {
        // Déclaration de variable typée
        int p_offset = n_owned_u + n_owned_v;
        // Boucle spatiale sur les cellules locales selon l'axe vertical Y (ordonnées)
        for (int j = 1; j <= mesh.ny_local; ++j) {
            // Boucle spatiale sur les cellules locales selon l'axe horizontal X (abscisses)
            for (int i = 1; i <= mesh.nx_local; ++i) {
                // Déclaration de variable typée
                int dof = (j - 1) * mesh.nx_local + (i - 1);
                // Calcul de la résolution géométrique de cellule
                fields.p[mesh.cell_idx(i, j)] = vec[p_offset + dof];
            }
        }
        MpiComm::sync_cell_scalar(mesh, fields.p);
    }

    // Synchronisation des mailles fantômes pour calculs ultérieurs
    MpiComm::sync_u(mesh, fields.u);
    MpiComm::sync_v(mesh, fields.v);
    MpiComm::apply_boundary_conditions(mesh, bc, fields, homogeneous);
}

void NavierStokesSolver::apply_system_matrix(const FlowFields& fields_old, double dt,
                                            const std::vector<double>& x, std::vector<double>& y) const {
    y.assign(n_owned_total, 0.0);

    FlowFields tf;
    tf.allocate(mesh);
    // Calcul et assignation numérique
    unpack_vector(x, tf, true); // true = CL homogènes pour garantir la stricte linéarité de A*x

    // Déclaration de variable typée
    double dx = mesh.dx;
    // Déclaration de variable typée
    double dy = mesh.dy;
    // Déclaration de variable typée
    double dvol = dx * dy;

    // Déclaration de variable typée
    int u_nx = mesh.nx_local + (mesh.is_boundary_right ? 1 : 0);
    // Déclaration de variable typée
    int v_ny = mesh.ny_local + (mesh.is_boundary_top ? 1 : 0);

    // Précalcul de la divergence sur chaque cellule locale + halo
    std::vector<double> div_cell((mesh.nx_local + 2) * (mesh.ny_local + 2), 0.0);
    // Calcul de la résolution géométrique de cellule
    auto c_idx = [&](int i, int j) { return j * (mesh.nx_local + 2) + i; };

    // Boucle spatiale sur les cellules locales selon l'axe vertical Y (ordonnées)
    for (int j = 0; j <= mesh.ny_local + 1; ++j) {
        // Boucle spatiale sur les cellules locales selon l'axe horizontal X (abscisses)
        for (int i = 0; i <= mesh.nx_local + 1; ++i) {
            // Déclaration de variable typée
            double d = (tf.u[mesh.u_idx(i + 1, j)] - tf.u[mesh.u_idx(i, j)]) / dx +
                       (tf.v[mesh.v_idx(i, j + 1)] - tf.v[mesh.v_idx(i, j)]) / dy;
            // Calcul de la résolution géométrique de cellule
            div_cell[c_idx(i, j)] = d;
        }
    }

    // =========================================================================
    // 1. Équation de Quantité de Mouvement en u (Faces verticales)
    // =========================================================================
    for (int j = 1; j <= mesh.ny_local; ++j) {
        // Boucle spatiale sur les cellules locales selon l'axe horizontal X (abscisses)
        for (int i = 1; i <= u_nx; ++i) {
            // Déclaration de variable typée
            int dof = (j - 1) * u_nx + (i - 1);

            // Condition aux limites physique sur les parois gauche / droite ou obstacle solide
            if (i == 1 && mesh.is_boundary_left) {
                // Déclaration de variable typée
                double rho_u = fields_old.rho_u[mesh.u_idx(i, j)];
                // Déclaration de variable typée
                double diag_bc = (rho_u / dt) * dvol;
                // Évaluation de la condition de contrôle
                if (bc.left.type == BCType::NEUMANN) {
                    // Calcul et assignation numérique
                    y[dof] = diag_bc * (tf.u[mesh.u_idx(1, j)] - tf.u[mesh.u_idx(2, j)]);
                } else {
                    // Calcul et assignation numérique
                    y[dof] = diag_bc * tf.u[mesh.u_idx(1, j)];
                }
                continue;
            }
            // Évaluation de la condition de contrôle
            if (i == u_nx && mesh.is_boundary_right) {
                // Déclaration de variable typée
                double rho_u = fields_old.rho_u[mesh.u_idx(i, j)];
                // Déclaration de variable typée
                double diag_bc = (rho_u / dt) * dvol;
                // Évaluation de la condition de contrôle
                if (bc.right.type == BCType::NEUMANN) {
                    // Calcul et assignation numérique
                    y[dof] = diag_bc * (tf.u[mesh.u_idx(u_nx, j)] - tf.u[mesh.u_idx(u_nx - 1, j)]);
                } else {
                    // Calcul et assignation numérique
                    y[dof] = diag_bc * tf.u[mesh.u_idx(u_nx, j)];
                }
                continue;
            }
            // Évaluation de la condition de contrôle
            if (is_u_face_in_obstacle(i, j)) {
                // Déclaration de variable typée
                double rho_u = fields_old.rho_u[mesh.u_idx(i, j)];
                // Déclaration de variable typée
                double diag_bc = (rho_u / dt) * dvol;
                // Calcul et assignation numérique
                y[dof] = diag_bc * tf.u[mesh.u_idx(i, j)];
                continue;
            }

            // Déclaration de variable typée
            double rho_u = fields_old.rho_u[mesh.u_idx(i, j)];

            // Terme instationnaire
            double val = (rho_u / dt) * tf.u[mesh.u_idx(i, j)] * dvol;

            // Terme convectif semi-implicite F_uu (forme advective non-conservative a densite harmonique)
            if (enable_convection) {
                double u_n = fields_old.u[mesh.u_idx(i, j)];
                double u_n_E = fields_old.u[mesh.u_idx(i + 1, j)];
                double u_n_W = fields_old.u[mesh.u_idx(i - 1, j)];

                double v_n_N = 0.5 * (fields_old.v[mesh.v_idx(i - 1, j + 1)] + fields_old.v[mesh.v_idx(i, j + 1)]);
                double v_n_S = 0.5 * (fields_old.v[mesh.v_idx(i - 1, j)] + fields_old.v[mesh.v_idx(i, j)]);

                double rho_u_E = fields_old.rho_u[mesh.u_idx(i + 1, j)];
                double rho_u_W = fields_old.rho_u[mesh.u_idx(i - 1, j)];
                double rho_u_N = (j < mesh.ny_local) ? fields_old.rho_u[mesh.u_idx(i, j + 1)] : rho_u;
                double rho_u_S = (j > 1) ? fields_old.rho_u[mesh.u_idx(i, j - 1)] : rho_u;

                double rho_E = (2.0 * rho_u * rho_u_E) / std::max(1.0e-12, rho_u + rho_u_E);
                double rho_W = (2.0 * rho_u * rho_u_W) / std::max(1.0e-12, rho_u + rho_u_W);
                double rho_N = (2.0 * rho_u * rho_u_N) / std::max(1.0e-12, rho_u + rho_u_N);
                double rho_S = (2.0 * rho_u * rho_u_S) / std::max(1.0e-12, rho_u + rho_u_S);

                double du_E = 0.5 * (u_n + u_n_E) * 0.5 * (tf.u[mesh.u_idx(i + 1, j)] - tf.u[mesh.u_idx(i, j)]) * dy;
                double du_W = 0.5 * (u_n_W + u_n) * 0.5 * (tf.u[mesh.u_idx(i, j)] - tf.u[mesh.u_idx(i - 1, j)]) * dy;
                double du_N = v_n_N * 0.5 * (tf.u[mesh.u_idx(i, j + 1)] - tf.u[mesh.u_idx(i, j)]) * dx;
                double du_S = v_n_S * 0.5 * (tf.u[mesh.u_idx(i, j)] - tf.u[mesh.u_idx(i, j - 1)]) * dx;

                val += (rho_E * du_E + rho_W * du_W + rho_N * du_N + rho_S * du_S);
            }

            // Termes visqueux implicites d'ordre 2 (Caltagirone et al. - Annexe A)
            // (a) Viscosité d'élongation 2*mu d2u/dx2
            double mu_cell_E = fields_old.mu[mesh.cell_idx(i, j)];
            // Déclaration de variable typée
            double mu_cell_W = fields_old.mu[mesh.cell_idx(i - 1, j)];
            // Déclaration de variable typée
            double diff_elong_E = 2.0 * mu_cell_E * (tf.u[mesh.u_idx(i + 1, j)] - tf.u[mesh.u_idx(i, j)]) / dx * dy;
            // Déclaration de variable typée
            double diff_elong_W = 2.0 * mu_cell_W * (tf.u[mesh.u_idx(i, j)] - tf.u[mesh.u_idx(i - 1, j)]) / dx * dy;
            // Calcul et assignation numérique
            val -= (diff_elong_E - diff_elong_W);

            // (b) Viscosité de cisaillement et rotation mu (d2u/dy2 + d2v/dxdy)
            double mu_node_N = fields_old.mu_node[mesh.node_idx(i, j + 1)];
            // Déclaration de variable typée
            double mu_node_S = fields_old.mu_node[mesh.node_idx(i, j)];

            // Déclaration de variable typée
            double shear_N = mu_node_N * ((tf.u[mesh.u_idx(i, j + 1)] - tf.u[mesh.u_idx(i, j)]) / dy +
                                          (tf.v[mesh.v_idx(i, j + 1)] - tf.v[mesh.v_idx(i - 1, j + 1)]) / dx) * dx;
            // Déclaration de variable typée
            double shear_S = mu_node_S * ((tf.u[mesh.u_idx(i, j)] - tf.u[mesh.u_idx(i, j - 1)]) / dy +
                                          (tf.v[mesh.v_idx(i, j)] - tf.v[mesh.v_idx(i - 1, j)]) / dx) * dx;
            // Calcul et assignation numérique
            val -= (shear_N - shear_S);

            // Terme de pénalité de Lagrangien augmenté : - r * d(div)/dx * dvol = - r * (div_E - div_W) * dy
            double grad_div_x = (div_cell[c_idx(i, j)] - div_cell[c_idx(i - 1, j)]) * dy;
            // Calcul et assignation numérique
            val -= r_al * grad_div_x;

            // Gradient de pression discret B_u^T (uniquement si monolithique)
            if (is_monolithic()) {
                // Déclaration de variable typée
                double grad_p = (tf.p[mesh.cell_idx(i, j)] - tf.p[mesh.cell_idx(i - 1, j)]) * dy;
                // Calcul et assignation numérique
                val += grad_p;
            }

            // Calcul et assignation numérique
            y[dof] = val;
        }
    }

    // =========================================================================
    // 2. Équation de Quantité de Mouvement en v (Faces horizontales)
    // =========================================================================
    int v_offset = n_owned_u;
    // Boucle spatiale sur les cellules locales selon l'axe vertical Y (ordonnées)
    for (int j = 1; j <= v_ny; ++j) {
        // Boucle spatiale sur les cellules locales selon l'axe horizontal X (abscisses)
        for (int i = 1; i <= mesh.nx_local; ++i) {
            // Déclaration de variable typée
            int dof = v_offset + (j - 1) * mesh.nx_local + (i - 1);

            // Condition aux limites physique Dirichlet ou Free-Slip sur les parois basse / haute ou obstacle solide
            bool is_wall_bot = (j == 1 && mesh.is_boundary_bottom && (bc.bottom.type == BCType::DIRICHLET || bc.bottom.type == BCType::FREE_SLIP));
            // Déclaration de variable typée
            bool is_wall_top = (j == v_ny && mesh.is_boundary_top && (bc.top.type == BCType::DIRICHLET || bc.top.type == BCType::FREE_SLIP));
            // Déclaration de variable typée
            bool in_obs = is_v_face_in_obstacle(i, j);
            // Évaluation de la condition de contrôle
            if (is_wall_bot || is_wall_top || in_obs) {
                // Déclaration de variable typée
                double rho_v = fields_old.rho_v[mesh.v_idx(i, j)];
                // Déclaration de variable typée
                double diag_bc = (rho_v / dt) * dvol;
                // Calcul et assignation numérique
                y[dof] = diag_bc * tf.v[mesh.v_idx(i, j)];
                continue;
            }

            // Déclaration de variable typée
            double rho_v = fields_old.rho_v[mesh.v_idx(i, j)];

            // Terme instationnaire
            double val = (rho_v / dt) * tf.v[mesh.v_idx(i, j)] * dvol;

            // Terme convectif semi-implicite F_vv (forme advective non-conservative a densite harmonique)
            if (enable_convection) {
                double v_n = fields_old.v[mesh.v_idx(i, j)];
                double v_n_N = fields_old.v[mesh.v_idx(i, j + 1)];
                double v_n_S = fields_old.v[mesh.v_idx(i, j - 1)];

                double u_n_E = 0.5 * (fields_old.u[mesh.u_idx(i + 1, j - 1)] + fields_old.u[mesh.u_idx(i + 1, j)]);
                double u_n_W = 0.5 * (fields_old.u[mesh.u_idx(i, j - 1)] + fields_old.u[mesh.u_idx(i, j)]);

                double rho_v_E = (i < mesh.nx_local) ? fields_old.rho_v[mesh.v_idx(i + 1, j)] : rho_v;
                double rho_v_W = (i > 1) ? fields_old.rho_v[mesh.v_idx(i - 1, j)] : rho_v;
                double rho_v_N = fields_old.rho_v[mesh.v_idx(i, j + 1)];
                double rho_v_S = fields_old.rho_v[mesh.v_idx(i, j - 1)];

                double rho_E = (2.0 * rho_v * rho_v_E) / std::max(1.0e-12, rho_v + rho_v_E);
                double rho_W = (2.0 * rho_v * rho_v_W) / std::max(1.0e-12, rho_v + rho_v_W);
                double rho_N = (2.0 * rho_v * rho_v_N) / std::max(1.0e-12, rho_v + rho_v_N);
                double rho_S = (2.0 * rho_v * rho_v_S) / std::max(1.0e-12, rho_v + rho_v_S);

                double dv_E = u_n_E * 0.5 * (tf.v[mesh.v_idx(i + 1, j)] - tf.v[mesh.v_idx(i, j)]) * dy;
                double dv_W = u_n_W * 0.5 * (tf.v[mesh.v_idx(i, j)] - tf.v[mesh.v_idx(i - 1, j)]) * dy;
                double dv_N = 0.5 * (v_n + v_n_N) * 0.5 * (tf.v[mesh.v_idx(i, j + 1)] - tf.v[mesh.v_idx(i, j)]) * dx;
                double dv_S = 0.5 * (v_n_S + v_n) * 0.5 * (tf.v[mesh.v_idx(i, j)] - tf.v[mesh.v_idx(i, j - 1)]) * dx;

                val += (rho_E * dv_E + rho_W * dv_W + rho_N * dv_N + rho_S * dv_S);
            }

            // Termes visqueux implicites d'ordre 2
            // (a) Viscosité d'élongation 2*mu d2v/dy2
            double mu_cell_N = fields_old.mu[mesh.cell_idx(i, j)];
            // Déclaration de variable typée
            double mu_cell_S = fields_old.mu[mesh.cell_idx(i, j - 1)];
            // Déclaration de variable typée
            double diff_elong_N = 2.0 * mu_cell_N * (tf.v[mesh.v_idx(i, j + 1)] - tf.v[mesh.v_idx(i, j)]) / dy * dx;
            // Déclaration de variable typée
            double diff_elong_S = 2.0 * mu_cell_S * (tf.v[mesh.v_idx(i, j)] - tf.v[mesh.v_idx(i, j - 1)]) / dy * dx;
            // Calcul et assignation numérique
            val -= (diff_elong_N - diff_elong_S);

            // (b) Viscosité de cisaillement et rotation mu (d2v/dx2 + d2u/dxdy)
            double mu_node_E = fields_old.mu_node[mesh.node_idx(i + 1, j)];
            // Déclaration de variable typée
            double mu_node_W = fields_old.mu_node[mesh.node_idx(i, j)];

            // Déclaration de variable typée
            double shear_E = mu_node_E * ((tf.v[mesh.v_idx(i + 1, j)] - tf.v[mesh.v_idx(i, j)]) / dx +
                                          (tf.u[mesh.u_idx(i + 1, j)] - tf.u[mesh.u_idx(i + 1, j - 1)]) / dy) * dy;
            // Déclaration de variable typée
            double shear_W = mu_node_W * ((tf.v[mesh.v_idx(i, j)] - tf.v[mesh.v_idx(i - 1, j)]) / dx +
                                          (tf.u[mesh.u_idx(i, j)] - tf.u[mesh.u_idx(i, j - 1)]) / dy) * dy;
            // Calcul et assignation numérique
            val -= (shear_E - shear_W);

            // Terme de pénalité de Lagrangien augmenté : - r * d(div)/dy * dvol = - r * (div_N - div_S) * dx
            double grad_div_y = (div_cell[c_idx(i, j)] - div_cell[c_idx(i, j - 1)]) * dx;
            // Calcul et assignation numérique
            val -= r_al * grad_div_y;

            // Gradient de pression discret B_v^T (uniquement si monolithique)
            if (is_monolithic()) {
                // Déclaration de variable typée
                double grad_p = (tf.p[mesh.cell_idx(i, j)] - tf.p[mesh.cell_idx(i, j - 1)]) * dx;
                // Calcul et assignation numérique
                val += grad_p;
            }

            // Calcul et assignation numérique
            y[dof] = val;
        }
    }

    // =========================================================================
    // 3. Équation de Continuité (si monolithique)
    // =========================================================================
    if (is_monolithic()) {
        // Déclaration de variable typée
        int p_offset = n_owned_u + n_owned_v;
        // Boucle spatiale sur les cellules locales selon l'axe vertical Y (ordonnées)
        for (int j = 1; j <= mesh.ny_local; ++j) {
            // Boucle spatiale sur les cellules locales selon l'axe horizontal X (abscisses)
            for (int i = 1; i <= mesh.nx_local; ++i) {
                // Déclaration de variable typée
                int dof = p_offset + (j - 1) * mesh.nx_local + (i - 1);
                // Déclaration de variable typée
                double div = div_cell[c_idx(i, j)] * dvol;
                // Déclaration de variable typée
                double diag_p_term = (dvol / r_al) * tf.p[mesh.cell_idx(i, j)];
                // Calcul et assignation numérique
                y[dof] = div - diag_p_term;
            }
        }
    }
}

void NavierStokesSolver::assemble_rhs(const FlowFields& fields, double dt, std::vector<double>& rhs) const {
    rhs.assign(n_owned_total, 0.0);

    // Déclaration de variable typée
    double dx = mesh.dx;
    // Déclaration de variable typée
    double dy = mesh.dy;
    // Déclaration de variable typée
    double dvol = dx * dy;

    // Déclaration de variable typée
    int u_nx = mesh.nx_local + (mesh.is_boundary_right ? 1 : 0);
    // Déclaration de variable typée
    int v_ny = mesh.ny_local + (mesh.is_boundary_top ? 1 : 0);

    // RHS pour u
    for (int j = 1; j <= mesh.ny_local; ++j) {
        // Boucle spatiale sur les cellules locales selon l'axe horizontal X (abscisses)
        for (int i = 1; i <= u_nx; ++i) {
            // Déclaration de variable typée
            int dof = (j - 1) * u_nx + (i - 1);
            // Déclaration de variable typée
            double rho_u = fields.rho_u[mesh.u_idx(i, j)];
            // Déclaration de variable typée
            double diag_bc = (rho_u / dt) * dvol;

            // Si face dans un obstacle solide (adhérence nulle)
            if (is_u_face_in_obstacle(i, j)) {
                // Calcul et assignation numérique
                rhs[dof] = 0.0;
                continue;
            }

            // Si paroi physique gauche
            if (i == 1 && mesh.is_boundary_left) {
                // Déclaration de variable typée
                double u_bc = 0.0;
                // Évaluation de la condition de contrôle
                if (bc.left.type == BCType::DIRICHLET) {
                    // Évaluation de la condition de contrôle
                    if (bc.left.profile == "POISEUILLE_STEP") {
                        // Déclaration de variable typée
                        double y = mesh.yc(j);
                        // Évaluation de la condition de contrôle
                        if (y >= bc.left.y_step) {
                            // Déclaration de variable typée
                            double h = mesh.ly - bc.left.y_step;
                            // Calcul et assignation numérique
                            u_bc = 6.0 * bc.left.u_mean * (y - bc.left.y_step) * (mesh.ly - y) / (h * h);
                        }
                    // Calcul et assignation numérique
                    } else if (bc.left.profile == "PARABOLIC") {
                        // Déclaration de variable typée
                        double y = mesh.yc(j);
                        // Déclaration de variable typée
                        double h = mesh.ly;
                        // Calcul et assignation numérique
                        u_bc = 6.0 * bc.left.u_mean * y * (h - y) / (h * h);
                    } else {
                        // Calcul et assignation numérique
                        u_bc = bc.left.u_val;
                    }
                // Calcul et assignation numérique
                } else if (bc.left.type == BCType::NEUMANN) {
                    // Calcul et assignation numérique
                    u_bc = 0.0;
                }
                // Calcul et assignation numérique
                rhs[dof] = diag_bc * u_bc;
                continue;
            }
            // Si paroi physique droite
            if (i == u_nx && mesh.is_boundary_right) {
                // Déclaration de variable typée
                double u_bc = 0.0;
                // Évaluation de la condition de contrôle
                if (bc.right.type == BCType::DIRICHLET) {
                    // Calcul et assignation numérique
                    u_bc = bc.right.u_val;
                // Calcul et assignation numérique
                } else if (bc.right.type == BCType::NEUMANN) {
                    // Calcul et assignation numérique
                    u_bc = 0.0;
                }
                // Calcul et assignation numérique
                rhs[dof] = diag_bc * u_bc;
                continue;
            }

            // Déclaration de variable typée
            double val = (rho_u / dt) * fields.u[mesh.u_idx(i, j)] * dvol;
            // Calcul et assignation numérique
            val += rho_u * gx * dvol;
            // Calcul et assignation numérique
            val += fields.F_sigma_x[mesh.u_idx(i, j)] * dvol;

            // Pour Lagrangien augmenté standard : le gradient de pression -grad(p^n) est explicite au second membre
            if (is_augmented_lagrangian()) {
                // Déclaration de variable typée
                double grad_p = (fields.p[mesh.cell_idx(i, j)] - fields.p[mesh.cell_idx(i - 1, j)]) * dy;
                // Calcul et assignation numérique
                val -= grad_p;
            }

            // Contribution de contrainte tangentielle à la paroi pour u (haut / bas)
            if (j == mesh.ny_local && mesh.is_boundary_top && bc.top.type == BCType::DIRICHLET) {
                // Déclaration de variable typée
                double mu_N = fields.mu_node[mesh.node_idx(i, j + 1)];
                // Calcul et assignation numérique
                val += mu_N * (2.0 * bc.top.u_val / dy) * dx;
            }
            // Évaluation de la condition de contrôle
            if (j == 1 && mesh.is_boundary_bottom && bc.bottom.type == BCType::DIRICHLET) {
                // Déclaration de variable typée
                double mu_S = fields.mu_node[mesh.node_idx(i, j)];
                // Calcul et assignation numérique
                val += mu_S * (2.0 * bc.bottom.u_val / dy) * dx;
            }

            // Calcul et assignation numérique
            rhs[dof] = val;
        }
    }

    // RHS pour v
    int v_offset = n_owned_u;
    // Boucle spatiale sur les cellules locales selon l'axe vertical Y (ordonnées)
    for (int j = 1; j <= v_ny; ++j) {
        // Boucle spatiale sur les cellules locales selon l'axe horizontal X (abscisses)
        for (int i = 1; i <= mesh.nx_local; ++i) {
            // Déclaration de variable typée
            int dof = v_offset + (j - 1) * mesh.nx_local + (i - 1);
            // Déclaration de variable typée
            double rho_v = fields.rho_v[mesh.v_idx(i, j)];
            // Déclaration de variable typée
            double diag_bc = (rho_v / dt) * dvol;

            // Si face dans un obstacle solide (adhérence nulle)
            if (is_v_face_in_obstacle(i, j)) {
                // Calcul et assignation numérique
                rhs[dof] = 0.0;
                continue;
            }

            // Évaluation de la condition de contrôle
            if (j == 1 && mesh.is_boundary_bottom && (bc.bottom.type == BCType::DIRICHLET || bc.bottom.type == BCType::FREE_SLIP)) {
                // Déclaration de variable typée
                double v_bc = (bc.bottom.type == BCType::DIRICHLET) ? bc.bottom.v_val : 0.0;
                // Calcul et assignation numérique
                rhs[dof] = diag_bc * v_bc;
                continue;
            }
            // Évaluation de la condition de contrôle
            if (j == v_ny && mesh.is_boundary_top && (bc.top.type == BCType::DIRICHLET || bc.top.type == BCType::FREE_SLIP)) {
                // Déclaration de variable typée
                double v_bc = 0.0;
                // Évaluation de la condition d'injecteur (Cas 3.3 Fluids 2021)
                if (bc.top.profile == "JET_INLET") {
                    // Déclaration de variable typée
                    double x = mesh.xc(i);
                    // Déclaration de variable typée pour le centre et demi-largeur
                    double xc = (bc.top.jet_center > 0.0) ? bc.top.jet_center : 0.5 * mesh.lx;
                    // Calcul de la demi-largeur de l'injecteur
                    double half_w = (bc.top.jet_width > 0.0) ? 0.5 * bc.top.jet_width : 0.04;
                    double d_disc = 0.0;
                    for (int ig = 1; ig <= mesh.nx_global; ++ig) {
                        double xg = (ig - 0.5) * mesh.dx;
                        if (std::abs(xg - xc) <= half_w) d_disc += mesh.dx;
                    }
                    if (d_disc < 1.0e-12) d_disc = 2.0 * half_w;
                    // Évaluation de l'appartenance à l'injecteur
                    if (std::abs(x - xc) <= half_w) {
                        // Vitesse d'injection négative imposée (conservant le débit D * v0 exact)
                        v_bc = bc.top.v_val * (2.0 * half_w) / d_disc;
                    } else {
                        // Vitesse positive d'échappement compensatoire du gaz
                        v_bc = -bc.top.v_val * (2.0 * half_w) / (mesh.lx - d_disc);
                    }
                } else if (bc.top.type == BCType::DIRICHLET) {
                    // Calcul et assignation numérique standard
                    v_bc = bc.top.v_val;
                }
                // Calcul et assignation numérique
                rhs[dof] = diag_bc * v_bc;
                continue;
            }

            // Déclaration de variable typée
            double val = (rho_v / dt) * fields.v[mesh.v_idx(i, j)] * dvol;
            // Calcul et assignation numérique
            val += rho_v * gy * dvol;
            // Évaluation de la condition de contrôle
            if (enable_energy) {
                // Déclaration de variable typée
                double T_v = 0.5 * (fields.T[mesh.cell_idx(i, j)] + fields.T[mesh.cell_idx(i, j - 1)]);
                // Calcul et assignation numérique
                val += rho_v * beta * (T_v - T_ref) * gy * dvol;
            }
            // Calcul et assignation numérique
            val += fields.F_sigma_y[mesh.v_idx(i, j)] * dvol;

            // Pour Lagrangien augmenté standard : -grad(p^n)
            if (is_augmented_lagrangian()) {
                // Déclaration de variable typée
                double grad_p = (fields.p[mesh.cell_idx(i, j)] - fields.p[mesh.cell_idx(i, j - 1)]) * dx;
                // Calcul et assignation numérique
                val -= grad_p;
            }

            // Contribution de contrainte tangentielle à la paroi pour v (gauche / droite)
            if (i == 1 && mesh.is_boundary_left && bc.left.type == BCType::DIRICHLET) {
                // Déclaration de variable typée
                double mu_W = fields.mu_node[mesh.node_idx(i, j)];
                // Calcul et assignation numérique
                val += mu_W * (2.0 * bc.left.v_val / dx) * dy;
            }
            // Évaluation de la condition de contrôle
            if (i == mesh.nx_local && mesh.is_boundary_right && bc.right.type == BCType::DIRICHLET) {
                // Déclaration de variable typée
                double mu_E = fields.mu_node[mesh.node_idx(i + 1, j)];
                // Calcul et assignation numérique
                val += mu_E * (2.0 * bc.right.v_val / dx) * dy;
            }

            // Calcul et assignation numérique
            rhs[dof] = val;
        }
    }

    // RHS pour p (si monolithique régularisé)
    if (is_monolithic()) {
        // Déclaration de variable typée
        int p_offset = n_owned_u + n_owned_v;
        // Boucle spatiale sur les cellules locales selon l'axe vertical Y (ordonnées)
        for (int j = 1; j <= mesh.ny_local; ++j) {
            // Boucle spatiale sur les cellules locales selon l'axe horizontal X (abscisses)
            for (int i = 1; i <= mesh.nx_local; ++i) {
                // Déclaration de variable typée
                int dof = p_offset + (j - 1) * mesh.nx_local + (i - 1);
                // Calcul et assignation numérique
                rhs[dof] = - (dvol / r_al) * fields.p[mesh.cell_idx(i, j)];
            }
        }
    }
}

void NavierStokesSolver::update_pressure_al(FlowFields& fields) const {
    // Déclaration de variable typée
    double dx = mesh.dx;
    // Déclaration de variable typée
    double dy = mesh.dy;

    // Boucle spatiale sur les cellules locales selon l'axe vertical Y (ordonnées)
    for (int j = 1; j <= mesh.ny_local; ++j) {
        // Boucle spatiale sur les cellules locales selon l'axe horizontal X (abscisses)
        for (int i = 1; i <= mesh.nx_local; ++i) {
            // Déclaration de variable typée
            double div = (fields.u[mesh.u_idx(i + 1, j)] - fields.u[mesh.u_idx(i, j)]) / dx +
                         (fields.v[mesh.v_idx(i, j + 1)] - fields.v[mesh.v_idx(i, j)]) / dy;
            // Calcul de la résolution géométrique de cellule
            fields.p[mesh.cell_idx(i, j)] -= r_al * div;
        }
    }

    MpiComm::sync_cell_scalar(mesh, fields.p);
}

void NavierStokesSolver::assemble_preconditioner_blocks(const FlowFields& fields, double dt) {
    // Déclaration de variable typée
    double dx = mesh.dx;
    // Déclaration de variable typée
    double dy = mesh.dy;
    // Déclaration de variable typée
    double dvol = dx * dy;

    // Calcul des diagonales effectives de Fuu et Fvv pour le préconditionnement Jacobi
    diag_u.assign(n_owned_u, 1.0);
    // Déclaration de variable typée
    int u_nx = mesh.nx_local + (mesh.is_boundary_right ? 1 : 0);
    // Boucle spatiale sur les cellules locales selon l'axe vertical Y (ordonnées)
    for (int j = 1; j <= mesh.ny_local; ++j) {
        // Boucle spatiale sur les cellules locales selon l'axe horizontal X (abscisses)
        for (int i = 1; i <= u_nx; ++i) {
            // Déclaration de variable typée
            int dof = (j - 1) * u_nx + (i - 1);
            // Déclaration de variable typée
            double rho_u = fields.rho_u[mesh.u_idx(i, j)];
            // Déclaration de variable typée
            bool is_wall = (i == 1 && mesh.is_boundary_left) ||
                           // Calcul et assignation numérique
                           (i == u_nx && mesh.is_boundary_right) ||
                           is_u_face_in_obstacle(i, j);
            // Évaluation de la condition de contrôle
            if (is_wall) {
                // Calcul et assignation numérique
                diag_u[dof] = (rho_u / dt) * dvol;
                continue;
            }
            // Déclaration de variable typée
            double mu_c = fields.mu[mesh.cell_idx(i, j)];
            // Déclaration de variable typée
            double d = (rho_u / dt) * dvol + 4.0 * mu_c * (dy / dx) + 2.0 * mu_c * (dx / dy) + 2.0 * r_al * (dy / dx);
            // Calcul et assignation numérique
            diag_u[dof] = (d > 1.0e-15) ? d : 1.0;
        }
    }

    diag_v.assign(n_owned_v, 1.0);
    // Déclaration de variable typée
    int v_ny = mesh.ny_local + (mesh.is_boundary_top ? 1 : 0);
    // Boucle spatiale sur les cellules locales selon l'axe vertical Y (ordonnées)
    for (int j = 1; j <= v_ny; ++j) {
        // Boucle spatiale sur les cellules locales selon l'axe horizontal X (abscisses)
        for (int i = 1; i <= mesh.nx_local; ++i) {
            // Déclaration de variable typée
            int dof = (j - 1) * mesh.nx_local + (i - 1);
            // Déclaration de variable typée
            double rho_v = fields.rho_v[mesh.v_idx(i, j)];
            // Déclaration de variable typée
            bool is_wall = (j == 1 && mesh.is_boundary_bottom && (bc.bottom.type == BCType::DIRICHLET || bc.bottom.type == BCType::FREE_SLIP)) ||
                           // Calcul et assignation numérique
                           (j == v_ny && mesh.is_boundary_top && (bc.top.type == BCType::DIRICHLET || bc.top.type == BCType::FREE_SLIP)) ||
                           is_v_face_in_obstacle(i, j);
            // Évaluation de la condition de contrôle
            if (is_wall) {
                // Calcul et assignation numérique
                diag_v[dof] = (rho_v / dt) * dvol;
                continue;
            }
            // Déclaration de variable typée
            double mu_c = fields.mu[mesh.cell_idx(i, j)];
            // Déclaration de variable typée
            double d = (rho_v / dt) * dvol + 4.0 * mu_c * (dx / dy) + 2.0 * mu_c * (dy / dx) + 2.0 * r_al * (dx / dy);
            // Calcul et assignation numérique
            diag_v[dof] = (d > 1.0e-15) ? d : 1.0;
        }
    }

    // Évaluation de la condition de contrôle
    if (is_monolithic()) {
        diag_p.assign(n_owned_p, dvol / r_al);
    }
}

void NavierStokesSolver::solve_energy(FlowFields& fields, double dt) const {
    // Évaluation de la condition de contrôle
    if (!enable_energy) return;

    // Déclaration de variable typée
    int nx = mesh.nx_local;
    // Déclaration de variable typée
    int ny = mesh.ny_local;
    // Déclaration de variable typée
    int n_cells = nx * ny;

    // Déclaration de variable typée
    double dx = mesh.dx;
    // Déclaration de variable typée
    double dy = mesh.dy;
    // Déclaration de variable typée
    double dvol = dx * dy;
    // Déclaration de variable typée
    double rho = fields.rho[mesh.cell_idx(1, 1)];

    // Calcul et assignation numérique
    auto cell_dof = [&](int i, int j) {
        // Retour de fonction
        return (j - 1) * nx + (i - 1);
    };

    // Déclaration de variable typée
    std::vector<int> rows;
    // Déclaration de variable typée
    std::vector<int> cols;
    // Déclaration de variable typée
    std::vector<double> vals;
    // Déclaration de variable typée
    std::vector<double> rhs(n_cells, 0.0);
    // Déclaration de variable typée
    std::vector<double> T_sol(n_cells, 0.0);

    // Boucle spatiale sur les cellules locales selon l'axe vertical Y (ordonnées)
    for (int j = 1; j <= ny; ++j) {
        // Boucle spatiale sur les cellules locales selon l'axe horizontal X (abscisses)
        for (int i = 1; i <= nx; ++i) {
            // Déclaration de variable typée
            int c_dof = cell_dof(i, j);
            // Calcul et assignation numérique
            T_sol[c_dof] = fields.T[mesh.cell_idx(i, j)];

            // Déclaration de variable typée
            double a_c = (rho * Cp / dt) * dvol;
            // Calcul et assignation numérique
            rhs[c_dof] = (rho * Cp / dt) * dvol * fields.T[mesh.cell_idx(i, j)];

            // Flux x-Est
            double u_e = fields.u[mesh.u_idx(i + 1, j)];
            // Déclaration de variable typée
            double d_e = k_thermal * dy / dx;
            // Déclaration de variable typée
            double conv_e = rho * Cp * u_e * dy * 0.5;

            // Flux x-Ouest
            double u_w = fields.u[mesh.u_idx(i, j)];
            // Déclaration de variable typée
            double d_w = k_thermal * dy / dx;
            // Déclaration de variable typée
            double conv_w = rho * Cp * u_w * dy * 0.5;

            // Flux y-Nord
            double v_n = fields.v[mesh.v_idx(i, j + 1)];
            // Déclaration de variable typée
            double d_n = k_thermal * dx / dy;
            // Déclaration de variable typée
            double conv_n = rho * Cp * v_n * dx * 0.5;

            // Flux y-Sud
            double v_s = fields.v[mesh.v_idx(i, j)];
            // Déclaration de variable typée
            double d_s = k_thermal * dx / dy;
            // Déclaration de variable typée
            double conv_s = rho * Cp * v_s * dx * 0.5;

            // Voisin Est
            if (i < nx) {
                // Calcul et assignation numérique
                a_c += d_e + conv_e;
                rows.push_back(c_dof);
                cols.push_back(cell_dof(i + 1, j));
                vals.push_back(-d_e + conv_e);
            } else if (mesh.is_boundary_right) {
                // Évaluation de la condition de contrôle
                if (bc.right.T_type == BCType::DIRICHLET) {
                    // Calcul et assignation numérique
                    a_c += 2.0 * d_e;
                    // Calcul et assignation numérique
                    rhs[c_dof] += 2.0 * d_e * bc.right.T_val;
                }
            }

            // Voisin Ouest
            if (i > 1) {
                // Calcul et assignation numérique
                a_c += d_w - conv_w;
                rows.push_back(c_dof);
                cols.push_back(cell_dof(i - 1, j));
                vals.push_back(-d_w - conv_w);
            } else if (mesh.is_boundary_left) {
                // Évaluation de la condition de contrôle
                if (bc.left.T_type == BCType::DIRICHLET) {
                    // Calcul et assignation numérique
                    a_c += 2.0 * d_w;
                    // Calcul et assignation numérique
                    rhs[c_dof] += 2.0 * d_w * bc.left.T_val;
                }
            }

            // Voisin Nord
            if (j < ny) {
                // Calcul et assignation numérique
                a_c += d_n + conv_n;
                rows.push_back(c_dof);
                cols.push_back(cell_dof(i, j + 1));
                vals.push_back(-d_n + conv_n);
            } else if (mesh.is_boundary_top) {
                // Évaluation de la condition de contrôle
                if (bc.top.T_type == BCType::DIRICHLET) {
                    // Calcul et assignation numérique
                    a_c += 2.0 * d_n;
                    // Calcul et assignation numérique
                    rhs[c_dof] += 2.0 * d_n * bc.top.T_val;
                }
            }

            // Voisin Sud
            if (j > 1) {
                // Calcul et assignation numérique
                a_c += d_s - conv_s;
                rows.push_back(c_dof);
                cols.push_back(cell_dof(i, j - 1));
                vals.push_back(-d_s - conv_s);
            } else if (mesh.is_boundary_bottom) {
                // Évaluation de la condition de contrôle
                if (bc.bottom.T_type == BCType::DIRICHLET) {
                    // Calcul et assignation numérique
                    a_c += 2.0 * d_s;
                    // Calcul et assignation numérique
                    rhs[c_dof] += 2.0 * d_s * bc.bottom.T_val;
                }
            }

            rows.push_back(c_dof);
            cols.push_back(c_dof);
            vals.push_back(a_c);
        }
    }

    CSRMatrix A;
    A.set_from_triplets(n_cells, n_cells, rows, cols, vals);

    // Solveur BiCGSTAB avec préconditionnement Jacobi
    std::vector<double> r = rhs;
    // Déclaration de variable typée
    std::vector<double> Ax(n_cells, 0.0);
    A.matvec(T_sol, Ax);
    // Boucle spatiale sur les cellules locales selon l'axe horizontal X (abscisses)
    for (int i = 0; i < n_cells; ++i) r[i] -= Ax[i];

    // Déclaration de variable typée
    std::vector<double> r0_hat = r;
    // Déclaration de variable typée
    std::vector<double> p = r;
    // Déclaration de variable typée
    std::vector<double> v(n_cells, 0.0);
    // Déclaration de variable typée
    std::vector<double> s(n_cells, 0.0);
    // Déclaration de variable typée
    std::vector<double> t(n_cells, 0.0);
    // Déclaration de variable typée
    std::vector<double> y_prec(n_cells, 0.0);
    // Déclaration de variable typée
    std::vector<double> z_prec(n_cells, 0.0);

    // Déclaration de variable typée
    double rho_old = 1.0;
    // Déclaration de variable typée
    double alpha = 1.0;
    // Déclaration de variable typée
    double omega = 1.0;

    // Déclaration de variable typée
    double res_init = 0.0;
    // Boucle spatiale sur les cellules locales selon l'axe horizontal X (abscisses)
    for (int i = 0; i < n_cells; ++i) res_init += r[i] * r[i];
    // Calcul de norme euclidienne ou résidu de convergence
    res_init = std::sqrt(res_init);
    // Évaluation de la condition de contrôle
    if (res_init > 1e-13) {
        // Déclaration de variable typée
        int max_it = 500;
        // Boucle d'itération algorithmique
        for (int iter = 0; iter < max_it; ++iter) {
            // Déclaration de variable typée
            double rho_new = 0.0;
            // Boucle spatiale sur les cellules locales selon l'axe horizontal X (abscisses)
            for (int i = 0; i < n_cells; ++i) rho_new += r0_hat[i] * r[i];
            // Évaluation de la condition de contrôle
            if (std::abs(rho_new) < 1e-30) break;

            // Évaluation de la condition de contrôle
            if (iter > 0) {
                // Déclaration de variable typée
                double beta_param = (rho_new / rho_old) * (alpha / omega);
                // Boucle spatiale sur les cellules locales selon l'axe horizontal X (abscisses)
                for (int i = 0; i < n_cells; ++i) {
                    // Calcul et assignation numérique
                    p[i] = r[i] + beta_param * (p[i] - omega * v[i]);
                }
            } else {
                // Calcul et assignation numérique
                p = r;
            }

            // Préconditionnement Jacobi : y = D^{-1} p
            for (int i = 0; i < n_cells; ++i) y_prec[i] = p[i] / A.diag[i];
            A.matvec(y_prec, v);

            // Déclaration de variable typée
            double r0v = 0.0;
            // Boucle spatiale sur les cellules locales selon l'axe horizontal X (abscisses)
            for (int i = 0; i < n_cells; ++i) r0v += r0_hat[i] * v[i];
            // Évaluation de la condition de contrôle
            if (std::abs(r0v) < 1e-30) break;
            // Calcul et assignation numérique
            alpha = rho_new / r0v;

            // Boucle spatiale sur les cellules locales selon l'axe horizontal X (abscisses)
            for (int i = 0; i < n_cells; ++i) s[i] = r[i] - alpha * v[i];
            // Déclaration de variable typée
            double s_norm = 0.0;
            // Boucle spatiale sur les cellules locales selon l'axe horizontal X (abscisses)
            for (int i = 0; i < n_cells; ++i) s_norm += s[i] * s[i];
            // Évaluation de la condition de contrôle
            if (std::sqrt(s_norm) / res_init < 1e-8) {
                // Boucle spatiale sur les cellules locales selon l'axe horizontal X (abscisses)
                for (int i = 0; i < n_cells; ++i) T_sol[i] += alpha * y_prec[i];
                break;
            }

            // Boucle spatiale sur les cellules locales selon l'axe horizontal X (abscisses)
            for (int i = 0; i < n_cells; ++i) z_prec[i] = s[i] / A.diag[i];
            A.matvec(z_prec, t);

            // Déclaration de variable typée
            double tt = 0.0, ts = 0.0;
            // Boucle spatiale sur les cellules locales selon l'axe horizontal X (abscisses)
            for (int i = 0; i < n_cells; ++i) {
                // Calcul et assignation numérique
                tt += t[i] * t[i];
                // Calcul et assignation numérique
                ts += t[i] * s[i];
            }
            // Calcul et assignation numérique
            omega = (tt > 1e-30) ? (ts / tt) : 0.0;

            // Boucle spatiale sur les cellules locales selon l'axe horizontal X (abscisses)
            for (int i = 0; i < n_cells; ++i) {
                // Calcul et assignation numérique
                T_sol[i] += alpha * y_prec[i] + omega * z_prec[i];
                // Calcul et assignation numérique
                r[i] = s[i] - omega * t[i];
            }

            // Calcul et assignation numérique
            rho_old = rho_new;

            // Déclaration de variable typée
            double r_norm = 0.0;
            // Boucle spatiale sur les cellules locales selon l'axe horizontal X (abscisses)
            for (int i = 0; i < n_cells; ++i) r_norm += r[i] * r[i];
            // Évaluation de la condition de contrôle
            if (std::sqrt(r_norm) / res_init < 1e-8) break;
        }
    }

    // Réinjecter T_sol dans fields.T
    for (int j = 1; j <= ny; ++j) {
        // Boucle spatiale sur les cellules locales selon l'axe horizontal X (abscisses)
        for (int i = 1; i <= nx; ++i) {
            // Calcul de la résolution géométrique de cellule
            fields.T[mesh.cell_idx(i, j)] = T_sol[cell_dof(i, j)];
        }
    }

    MpiComm::sync_cell_scalar(mesh, fields.T);
    MpiComm::apply_thermal_boundary_conditions(mesh, bc, fields);
}

// Déclaration de variable typée
double NavierStokesSolver::compute_mean_nusselt_hot_wall(const FlowFields& fields) const {
    // Déclaration de variable typée
    double local_nu_sum = 0.0;
    // Déclaration de variable typée
    int count = 0;
    // Évaluation de la condition de contrôle
    if (mesh.is_boundary_left) {
        // Boucle spatiale sur les cellules locales selon l'axe vertical Y (ordonnées)
        for (int j = 1; j <= mesh.ny_local; ++j) {
            // Déclaration de variable typée
            double T_w = bc.left.T_val;
            // Déclaration de variable typée
            double T_1 = fields.T[mesh.cell_idx(1, j)];
            // Déclaration de variable typée
            double delta_T = std::abs(bc.left.T_val - bc.right.T_val);
            // Évaluation de la condition de contrôle
            if (delta_T < 1e-12) delta_T = 1.0;
            // Déclaration de variable typée
            double nu_local = 2.0 * (T_w - T_1) / mesh.dx * (mesh.lx / delta_T);
            // Calcul et assignation numérique
            local_nu_sum += nu_local;
            count++;
        }
    }
    // Déclaration de variable typée
    double global_nu_sum = MpiComm::global_sum(mesh, local_nu_sum);
    // Déclaration de variable typée
    double global_count = MpiComm::global_sum(mesh, static_cast<double>(count));
    // Évaluation de la condition de contrôle
    if (global_count > 0.0) {
        // Retour de fonction
        return global_nu_sum / global_count;
    }
    // Retour de fonction
    return 1.0;
}

} // namespace covof


\end{lstlisting}

\chapter{Reconstruction Géométrique d'Interface VOF : PLIC, ELVIRA, LVIRA, SLIC (InterfaceReconstruction.hpp/.cpp)}
\section{Représentation Linéaire par Morceaux PLIC (\emph{Piecewise Linear Interface Calculation})}
Dans chaque maille diphasique $(i, j)$ telle que $\varepsilon_{\text{vof}} < C_{i,j} < 1 - \varepsilon_{\text{vof}}$ ($\varepsilon_{\text{vof}} = 10^{-12}$), l'interface est représentée dans le repère local normalisé $(\xi, \eta) \in [0, 1]^2$ de la maille par la droite :
\begin{equation}
  n_x \xi + n_y \eta = \alpha
\end{equation}
où $\mathbf{n} = (n_x, n_y)^T$ est la normale unitaire pointant vers l'extérieur du fluide 1 ($\mathbf{n} = -\nabla C / \|\nabla C\|$) et $\alpha \in \mathbb{R}$ est la constante de positionnement géométrique unique telle que l'aire du polygone tronqué $\{(\xi, \eta) \in [0, 1]^2 \mid n_x \xi + n_y \eta \le \alpha\}$ soit exactement égale à la fraction volumique $C_{i,j}$.

\section{Relations Analytiques Directe et Inverse de Scardovelli \& Zaleski (2000)}
En se ramenant par symétrie au premier quadrant ($m_1 = \min(|n_x|, |n_y|) \ge 0$, $m_2 = \max(|n_x|, |n_y|) > 0$, avec $m_1 + m_2 = 1$ en norme $L_1$), l'inversion analytique $C \mapsto \alpha$ implémentée dans \texttt{compute\_line\_constant} distingue deux régimes géométriques sans aucune itération de Newton :
\begin{enumerate}
  \item \textbf{Régime triangulaire} ($0 \le C \le \frac{m_1}{2 m_2}$) :
  \begin{equation}
    \alpha = \sqrt{2 m_1 m_2 C}
  \end{equation}
  \item \textbf{Régime trapézoïdal} ($\frac{m_1}{2 m_2} < C \le \frac{1}{2}$) :
  \begin{equation}
    \alpha = m_2 C + \frac{m_1}{2}
  \end{equation}
\end{enumerate}
Pour $C > 1/2$, la symétrie centrale du carré unité donne immédiatement $\alpha(C) = 1 - \alpha(1 - C)$.

\section{Estimation des Normales : Parker-Youngs, ELVIRA et LVIRA}
\begin{itemize}
  \item \textbf{Parker-Youngs} : différence finie pondérée sur le stencil $3\times 3$ avec poids $1, 2, 1$ privilégiant la direction axiale.
  \item \textbf{ELVIRA (Pilliod \& Puckett 2004)} : évalue 6 pentes candidates (différences arrière, centrée et avant des sommes de colonnes horizontales $H_i = \sum_{l=-1}^1 C_{i, j+l}$ et des sommes de lignes verticales $V_j = \sum_{k=-1}^1 C_{i+k, j}$), reconstruit la droite associée sur le bloc $3\times 3$ et retient celle qui minimise l'erreur quadratique $\sum_{k,l=-1}^1 (C^{\text{rec}}_{i+k,j+l}(\mathbf{n}) - C_{i+k,j+l})^2$. Cet algorithme reconstruit exactement toute interface plane quelle que soit son orientation.
  \item \textbf{LVIRA (Puckett 1991)} : affine l'angle $\theta = \text{atan2}(n_y, n_x)$ de la normale ELVIRA par recherche locale d'or / sécante minimisant la fonctionnelle d'erreur $E_2(\theta)$ sur le stencil $3\times 3$.
\end{itemize}

\section{Répertoire des Classes, Structures et Méthodes de \texttt{include/covof/InterfaceReconstruction.hpp}}
\begin{center}\small
\begin{tabular}{clll}
\toprule
\textbf{Ligne} & \textbf{Catégorie} & \textbf{Identifiant C++17} & \textbf{Signature / Rôle} \\
\midrule
13 & enum class & \texttt{ReconstructionMethod} & Déclaration de type / classe / structure \\
23 & enum class & \texttt{AdvectionMethod} & Déclaration de type / classe / structure \\
32 & enum class & \texttt{TvdLimiter} & Déclaration de type / classe / structure \\
43 & struct & \texttt{InterfaceSegment} & Déclaration de type / classe / structure \\
52 & class & \texttt{InterfaceReconstruction} & Déclaration de type / classe / structure \\
57 & Méthode / Fonction & \texttt{parse\_reconstruction\_method} & Type retour : static ReconstructionMethod \\
62 & Méthode / Fonction & \texttt{parse\_advection\_method} & Type retour : static AdvectionMethod \\
67 & Méthode / Fonction & \texttt{parse\_tvd\_limiter} & Type retour : static TvdLimiter \\
116 & Méthode / Fonction & \texttt{evaluate\_limiter} & Type retour : static double \\
\bottomrule
\end{tabular}
\end{center}

\section{Code Source Intégral Annoté : \texttt{include/covof/InterfaceReconstruction.hpp}}
Le listing ci-dessous présente l'intégralité des 119 lignes du fichier \texttt{include/covof/InterfaceReconstruction.hpp} de l'implémentation C++17 :

\begin{lstlisting}[language={[11]C++}]
#pragma once

#include "Mesh.hpp"
#include <string>
#include <vector>

namespace covof {
namespace fugu = covof; // Alias rétrocompatibilité

/**
 * @brief Énumération des méthodes de reconstruction géométrique d'interface
 */
enum class ReconstructionMethod {
    SLIC,   ///< Simple Line Interface Calculation (Noh & Woodward 1976)
    PLIC,   ///< Piecewise Linear Interface Calculation (Youngs 1982 / Scardovelli & Zaleski 1999)
    LVIRA,  ///< Least Squares VOF Interface Reconstruction Algorithm (Pilliod & Puckett 2004)
    ELVIRA  ///< Efficient LVIRA avec 6 candidats discrets (Pilliod & Puckett 2004)
};

/**
 * @brief Énumération des algorithmes d'advection VOF
 */
enum class AdvectionMethod {
    VOF_TVD,      ///< Advection algébrique Volumes Finis TVD haute résolution
    CIAM,         ///< Calcul d'Interface Affine par Morceaux (Strang splitting directionnel)
    WEYMOUTH_YUE  ///< Advection géométrique unsplit multidimensionnelle (Weymouth & Yue 2010)
};

/**
 * @brief Énumération des limiteurs de pente TVD / NVD pour l'advection d'interface
 */
enum class TvdLimiter {
    SUPERBEE,  ///< Limiteur Superbee de Roe (hautement compressif pour interfaces)
    VAN_LEER,  ///< Limiteur harmonique doux de Van Leer
    MUSCL,     ///< Limiteur linéaire par morceaux MUSCL
    HRIC       ///< High Resolution Interface Capturing (Muzaferija & Peric 1999)
};

/**
 * @brief Structure représentant un segment de droite d'interface affine dans une cellule
 * Équation locale : n_x * x + n_y * y + alpha = 0, orientée vers le fluide 2 (gaz)
 */
struct InterfaceSegment {
    double nx = 0.0;    ///< Composante normale unitaire selon X
    double ny = 0.0;    ///< Composante normale unitaire selon Y
    double alpha = 0.0; ///< Paramètre affine de position de l'interface
};

/**
 * @brief Classe statique fournissant les algorithmes géométriques de reconstruction et de découpe
 */
class InterfaceReconstruction {
public:
    /**
     * @brief Convertit une chaîne de caractères en méthode de reconstruction
     */
    static ReconstructionMethod parse_reconstruction_method(const std::string& str);

    /**
     * @brief Convertit une chaîne de caractères en méthode d'advection
     */
    static AdvectionMethod parse_advection_method(const std::string& str);

    /**
     * @brief Convertit une chaîne de caractères en limiteur TVD
     */
    static TvdLimiter parse_tvd_limiter(const std::string& str);

    /**
     * @brief Calcule l'aire exacte de l'intersection entre un rectangle [xmin, xmax] x [ymin, ymax]
     *        et le demi-plan fluide nx * x + ny * y + alpha <= 0 (Sutherland-Hodgman + Gauss Shoelace)
     */
    static double cut_box_area(double xmin, double xmax, double ymin, double ymax,
                               double nx, double ny, double alpha);

    /**
     * @brief Calcule le paramètre affine alpha conservant exactement le volume cible C_target * dx * dy
     */
    static double compute_alpha(double xmin, double xmax, double ymin, double ymax,
                                double nx, double ny, double C_target);

    /**
     * @brief Reconstruction SLIC (Simple Line Interface Calculation) sur l'ensemble du maillage
     */
    static void reconstruct_slic(const Mesh2D& mesh, const std::vector<double>& C,
                                 std::vector<InterfaceSegment>& segments);

    /**
     * @brief Reconstruction PLIC (Youngs + Scardovelli & Zaleski) sur l'ensemble du maillage
     */
    static void reconstruct_plic(const Mesh2D& mesh, const std::vector<double>& C,
                                 std::vector<InterfaceSegment>& segments);

    /**
     * @brief Reconstruction LVIRA (Least Squares VOF Interface Reconstruction Algorithm)
     */
    static void reconstruct_lvira(const Mesh2D& mesh, const std::vector<double>& C,
                                  std::vector<InterfaceSegment>& segments);

    /**
     * @brief Reconstruction ELVIRA (Efficient LVIRA avec 6 candidats de fonctions de hauteur)
     */
    static void reconstruct_elvira(const Mesh2D& mesh, const std::vector<double>& C,
                                   std::vector<InterfaceSegment>& segments);

    /**
     * @brief Point d'entrée générique de reconstruction selon la méthode configurée
     */
    static void reconstruct(ReconstructionMethod method, const Mesh2D& mesh,
                            const std::vector<double>& C,
                            std::vector<InterfaceSegment>& segments);

    /**
     * @brief Évalue le limiteur de flux TVD / NVD sélectionné
     */
    static double evaluate_limiter(TvdLimiter limiter, double r, double c_donor = 0.5, double c_upwind = 0.0, double c_acceptor = 1.0);
};

} // namespace covof

\end{lstlisting}

\section{Répertoire des Classes, Structures et Méthodes de \texttt{src/InterfaceReconstruction.cpp}}
\begin{center}\small
\begin{tabular}{clll}
\toprule
\textbf{Ligne} & \textbf{Catégorie} & \textbf{Identifiant C++17} & \textbf{Signature / Rôle} \\
\midrule
14 & Méthode / Fonction & \texttt{InterfaceReconstruction::parse\_reconstruction\_method} & Type retour : ReconstructionMethod \\
21 & Méthode / Fonction & \texttt{InterfaceReconstruction::parse\_advection\_method} & Type retour : AdvectionMethod \\
27 & Méthode / Fonction & \texttt{InterfaceReconstruction::parse\_tvd\_limiter} & Type retour : TvdLimiter \\
46 & struct & \texttt{Point} & Déclaration de type / classe / structure \\
535 & Méthode / Fonction & \texttt{InterfaceReconstruction::evaluate\_limiter} & Type retour : double \\
\bottomrule
\end{tabular}
\end{center}

\section{Code Source Intégral Annoté : \texttt{src/InterfaceReconstruction.cpp}}
Le listing ci-dessous présente l'intégralité des 576 lignes du fichier \texttt{src/InterfaceReconstruction.cpp} de l'implémentation C++17 :

\begin{lstlisting}[language={[11]C++}]
#include "covof/InterfaceReconstruction.hpp"
#include <cmath>
#include <algorithm>
#include <array>
#include <iostream>

namespace covof {
namespace fugu = covof; // Alias rétrocompatibilité

// =============================================================================
// PARSING DES MÉTHODES ET LIMITATEURS
// =============================================================================

ReconstructionMethod InterfaceReconstruction::parse_reconstruction_method(const std::string& str) {
    if (str == "SLIC")   return ReconstructionMethod::SLIC;
    if (str == "LVIRA")  return ReconstructionMethod::LVIRA;
    if (str == "ELVIRA") return ReconstructionMethod::ELVIRA;
    return ReconstructionMethod::PLIC; // Défaut PLIC
}

AdvectionMethod InterfaceReconstruction::parse_advection_method(const std::string& str) {
    if (str == "CIAM")         return AdvectionMethod::CIAM;
    if (str == "WEYMOUTH_YUE") return AdvectionMethod::WEYMOUTH_YUE;
    return AdvectionMethod::VOF_TVD; // Défaut VOF-TVD
}

TvdLimiter InterfaceReconstruction::parse_tvd_limiter(const std::string& str) {
    if (str == "VAN_LEER") return TvdLimiter::VAN_LEER;
    if (str == "MUSCL")    return TvdLimiter::MUSCL;
    if (str == "HRIC")     return TvdLimiter::HRIC;
    return TvdLimiter::SUPERBEE; // Défaut Superbee
}

// =============================================================================
// GÉOMÉTRIE : DÉCOUPE D'UN RECTANGLE PAR UN DEMI-PLAN (SUTHERLAND-HODGMAN)
// =============================================================================

double InterfaceReconstruction::cut_box_area(double xmin, double xmax, double ymin, double ymax,
                                            double nx, double ny, double alpha) {
    // Vérification des dimensions de la boîte de contrôle
    double dx = xmax - xmin;
    double dy = ymax - ymin;
    if (dx <= 1.0e-14 || dy <= 1.0e-14) return 0.0;

    // Définition des 4 sommets du rectangle dans l'ordre trigonométrique
    struct Point { double x, y; };
    std::array<Point, 4> box_pts = {{
        {xmin, ymin},
        {xmax, ymin},
        {xmax, ymax},
        {xmin, ymax}
    }};

    // Calcul de la distance signée d(x, y) = nx * x + ny * y + alpha
    auto dist = [&](const Point& p) {
        return nx * p.x + ny * p.y + alpha;
    };

    // Polygone résultant de la troncature
    std::vector<Point> clipped;
    clipped.reserve(8);

    // Algorithme de Sutherland-Hodgman contre le demi-plan fluide d(x, y) <= 0
    for (size_t i = 0; i < 4; ++i) {
        Point p1 = box_pts[i];
        Point p2 = box_pts[(i + 1) % 4];
        double d1 = dist(p1);
        double d2 = dist(p2);

        if (d1 <= 0.0) {
            if (d2 <= 0.0) {
                // Les deux points sont à l'intérieur du fluide
                clipped.push_back(p2);
            } else {
                // Sortie du fluide : calcul du point d'intersection
                double t = -d1 / (d2 - d1);
                clipped.push_back({p1.x + t * (p2.x - p1.x), p1.y + t * (p2.y - p1.y)});
            }
        } else {
            if (d2 <= 0.0) {
                // Entrée dans le fluide : calcul de l'intersection puis ajout de p2
                double t = -d1 / (d2 - d1);
                clipped.push_back({p1.x + t * (p2.x - p1.x), p1.y + t * (p2.y - p1.y)});
                clipped.push_back(p2);
            }
            // Sinon les deux points sont hors du fluide, rien à ajouter
        }
    }

    // Si moins de 3 sommets, l'aire découpée est nulle
    if (clipped.size() < 3) return 0.0;

    // Calcul de l'aire par la formule de Green / Gauss (Shoelace formula)
    double area = 0.0;
    size_t m = clipped.size();
    for (size_t i = 0; i < m; ++i) {
        size_t next = (i + 1) % m;
        area += clipped[i].x * clipped[next].y - clipped[next].x * clipped[i].y;
    }
    return 0.5 * std::abs(area);
}

// =============================================================================
// INVERSION NUMÉRIQUE EXACTE : FRACTION VOLUMIQUE -> PARAMÈTRE ALPHA
// =============================================================================

double InterfaceReconstruction::compute_alpha(double xmin, double xmax, double ymin, double ymax,
                                             double nx, double ny, double C_target) {
    double total_area = (xmax - xmin) * (ymax - ymin);
    if (C_target <= 1.0e-7) return 1.0e9;   // Cellule quasi-vide
    if (C_target >= 1.0 - 1.0e-7) return -1.0e9; // Cellule quasi-pleine

    double target_area = std::clamp(C_target, 0.0, 1.0) * total_area;

    // Bornes extrêmes pour alpha
    // Les 4 sommets donnent l'étendue maximale de nx*x + ny*y
    std::array<double, 4> vals = {{
        nx * xmin + ny * ymin,
        nx * xmax + ny * ymin,
        nx * xmax + ny * ymax,
        nx * xmin + ny * ymax
    }};
    double min_val = *std::min_element(vals.begin(), vals.end());
    double max_val = *std::max_element(vals.begin(), vals.end());

    // Pour alpha_low = -max_val, nx*x + ny*y + alpha <= 0 contient tout -> aire = total_area
    // Pour alpha_high = -min_val, nx*x + ny*y + alpha <= 0 est vide -> aire = 0
    double alpha_low = -max_val - 1.0e-12;
    double alpha_high = -min_val + 1.0e-12;

    // Méthode de dichotomie / bisection pour trouver alpha exact à 1.0e-12 (a décroît avec alpha)
    for (int iter = 0; iter < 45; ++iter) {
        double alpha_mid = 0.5 * (alpha_low + alpha_high);
        double a = cut_box_area(xmin, xmax, ymin, ymax, nx, ny, alpha_mid);
        if (a < target_area) {
            alpha_high = alpha_mid;
        } else {
            alpha_low = alpha_mid;
        }
    }
    return 0.5 * (alpha_low + alpha_high);
}

// =============================================================================
// RECONSTRUCTION SLIC (NOH & WOODWARD 1976)
// =============================================================================

void InterfaceReconstruction::reconstruct_slic(const Mesh2D& mesh, const std::vector<double>& C,
                                              std::vector<InterfaceSegment>& segments) {
    int nx = mesh.nx_local;
    int ny = mesh.ny_local;

    for (int j = 0; j <= ny + 1; ++j) {
        for (int i = 0; i <= nx + 1; ++i) {
            int idx = mesh.cell_idx(i, j);
            double c = C[idx];

            if (c <= 1.0e-6) {
                segments[idx] = {1.0, 0.0, 1.0e9};
                continue;
            }
            if (c >= 1.0 - 1.0e-6) {
                segments[idx] = {1.0, 0.0, -1.0e9};
                continue;
            }

            double xmin = mesh.xv(i);
            double xmax = mesh.xv(i + 1);
            double ymin = mesh.yv(j);
            double ymax = mesh.yv(j + 1);

            double snx = 1.0, sny = 0.0;
            if (i >= 1 && i <= nx && j >= 1 && j <= ny) {
                double grad_x = std::abs(C[mesh.cell_idx(i + 1, j)] - C[mesh.cell_idx(i - 1, j)]);
                double grad_y = std::abs(C[mesh.cell_idx(i, j + 1)] - C[mesh.cell_idx(i, j - 1)]);

                if (grad_x >= grad_y) {
                    if (C[mesh.cell_idx(i - 1, j)] >= C[mesh.cell_idx(i + 1, j)]) {
                        snx = 1.0; sny = 0.0;
                    } else {
                        snx = -1.0; sny = 0.0;
                    }
                } else {
                    if (C[mesh.cell_idx(i, j - 1)] >= C[mesh.cell_idx(i, j + 1)]) {
                        snx = 0.0; sny = 1.0;
                    } else {
                        snx = 0.0; sny = -1.0;
                    }
                }
            }

            double alpha = compute_alpha(xmin, xmax, ymin, ymax, snx, sny, c);
            segments[idx] = {snx, sny, alpha};
        }
    }
}

// =============================================================================
// RECONSTRUCTION PLIC (YOUNGS 1982)
// =============================================================================

void InterfaceReconstruction::reconstruct_plic(const Mesh2D& mesh, const std::vector<double>& C,
                                              std::vector<InterfaceSegment>& segments) {
    int nx = mesh.nx_local;
    int ny = mesh.ny_local;
    double dx = mesh.dx;
    double dy = mesh.dy;

    for (int j = 0; j <= ny + 1; ++j) {
        for (int i = 0; i <= nx + 1; ++i) {
            int idx = mesh.cell_idx(i, j);
            double c = C[idx];

            if (c <= 1.0e-6) {
                segments[idx] = {1.0, 0.0, 1.0e9};
                continue;
            }
            if (c >= 1.0 - 1.0e-6) {
                segments[idx] = {1.0, 0.0, -1.0e9};
                continue;
            }

            double mx = 1.0, my = 0.0;
            if (i >= 1 && i <= nx && j >= 1 && j <= ny) {
                // Calcul du gradient de Youngs sur stencil 3x3 centré
                double c_mm = C[mesh.cell_idx(i - 1, j - 1)];
                double c_0m = C[mesh.cell_idx(i,     j - 1)];
                double c_pm = C[mesh.cell_idx(i + 1, j - 1)];
                double c_m0 = C[mesh.cell_idx(i - 1, j    )];
                double c_p0 = C[mesh.cell_idx(i + 1, j    )];
                double c_mp = C[mesh.cell_idx(i - 1, j + 1)];
                double c_0p = C[mesh.cell_idx(i,     j + 1)];
                double c_pp = C[mesh.cell_idx(i + 1, j + 1)];

                mx = -((c_pp + 2.0 * c_p0 + c_pm) - (c_mp + 2.0 * c_m0 + c_mm)) / (8.0 * dx);
                my = -((c_pp + 2.0 * c_0p + c_mp) - (c_pm + 2.0 * c_0m + c_mm)) / (8.0 * dy);
            }

            double norm = std::hypot(mx, my);
            double pnx = (norm > 1.0e-12) ? mx / norm : 1.0;
            double pny = (norm > 1.0e-12) ? my / norm : 0.0;

            double xmin = mesh.xv(i);
            double xmax = mesh.xv(i + 1);
            double ymin = mesh.yv(j);
            double ymax = mesh.yv(j + 1);

            double alpha = compute_alpha(xmin, xmax, ymin, ymax, pnx, pny, c);
            segments[idx] = {pnx, pny, alpha};
        }
    }
}

// =============================================================================
// RECONSTRUCTION ELVIRA (PILLIOD & PUCKETT 2004)
// =============================================================================

void InterfaceReconstruction::reconstruct_elvira(const Mesh2D& mesh, const std::vector<double>& C,
                                                std::vector<InterfaceSegment>& segments) {
    int nx = mesh.nx_local;
    int ny = mesh.ny_local;
    double dx = mesh.dx;
    double dy = mesh.dy;

    for (int j = 0; j <= ny + 1; ++j) {
        for (int i = 0; i <= nx + 1; ++i) {
            int idx = mesh.cell_idx(i, j);
            double c = C[idx];

            if (c <= 1.0e-6) {
                segments[idx] = {1.0, 0.0, 1.0e9};
                continue;
            }
            if (c >= 1.0 - 1.0e-6) {
                segments[idx] = {1.0, 0.0, -1.0e9};
                continue;
            }
            if (i < 1 || i > nx || j < 1 || j > ny) {
                double xmin = mesh.xv(i);
                double xmax = mesh.xv(i + 1);
                double ymin = mesh.yv(j);
                double ymax = mesh.yv(j + 1);
                segments[idx] = {1.0, 0.0, compute_alpha(xmin, xmax, ymin, ymax, 1.0, 0.0, c)};
                continue;
            }

            double xmin = mesh.xv(i);
            double xmax = mesh.xv(i + 1);
            double ymin = mesh.yv(j);
            double ymax = mesh.yv(j + 1);

            // 1. Hauteurs de colonnes H_k (sommes sur y)
            double H_m1 = (C[mesh.cell_idx(i - 1, j - 1)] + C[mesh.cell_idx(i - 1, j)] + C[mesh.cell_idx(i - 1, j + 1)]) * dy;
            double H_0  = (C[mesh.cell_idx(i,     j - 1)] + C[mesh.cell_idx(i,     j)] + C[mesh.cell_idx(i,     j + 1)]) * dy;
            double H_p1 = (C[mesh.cell_idx(i + 1, j - 1)] + C[mesh.cell_idx(i + 1, j)] + C[mesh.cell_idx(i + 1, j + 1)]) * dy;

            // 2. Largeurs de lignes W_l (sommes sur x)
            double W_m1 = (C[mesh.cell_idx(i - 1, j - 1)] + C[mesh.cell_idx(i, j - 1)] + C[mesh.cell_idx(i + 1, j - 1)]) * dx;
            double W_0  = (C[mesh.cell_idx(i - 1, j    )] + C[mesh.cell_idx(i, j    )] + C[mesh.cell_idx(i + 1, j    )]) * dx;
            double W_p1 = (C[mesh.cell_idx(i - 1, j + 1)] + C[mesh.cell_idx(i, j + 1)] + C[mesh.cell_idx(i + 1, j + 1)]) * dx;

            // Détermination du signe de la normale
            double sgn_y = (C[mesh.cell_idx(i, j + 1)] < C[mesh.cell_idx(i, j - 1)]) ? 1.0 : -1.0;
            double sgn_x = (C[mesh.cell_idx(i + 1, j)] < C[mesh.cell_idx(i - 1, j)]) ? 1.0 : -1.0;

            // Ensemble des candidats de pentes
            std::vector<std::pair<double, double>> candidates;
            candidates.reserve(7);

            // Candidats verticaux (Height functions par colonnes)
            candidates.push_back({-sgn_y * (H_p1 - H_m1) / (2.0 * dx), sgn_y * 1.0}); // Centré
            candidates.push_back({-sgn_y * (H_p1 - H_0)  / (dx),       sgn_y * 1.0}); // Décentré avant
            candidates.push_back({-sgn_y * (H_0  - H_m1) / (dx),       sgn_y * 1.0}); // Décentré arrière

            // Candidats horizontaux (Height functions par lignes)
            candidates.push_back({sgn_x * 1.0, -sgn_x * (W_p1 - W_m1) / (2.0 * dy)}); // Centré
            candidates.push_back({sgn_x * 1.0, -sgn_x * (W_p1 - W_0)  / (dy)});       // Décentré avant
            candidates.push_back({sgn_x * 1.0, -sgn_x * (W_0  - W_m1) / (dy)});       // Décentré arrière

            // Candidat Youngs standard
            double mx = -((C[mesh.cell_idx(i+1,j+1)] + 2.0*C[mesh.cell_idx(i+1,j)] + C[mesh.cell_idx(i+1,j-1)]) -
                          (C[mesh.cell_idx(i-1,j+1)] + 2.0*C[mesh.cell_idx(i-1,j)] + C[mesh.cell_idx(i-1,j-1)])) / (8.0 * dx);
            double my = -((C[mesh.cell_idx(i+1,j+1)] + 2.0*C[mesh.cell_idx(i,j+1)] + C[mesh.cell_idx(i-1,j+1)]) -
                          (C[mesh.cell_idx(i+1,j-1)] + 2.0*C[mesh.cell_idx(i,j-1)] + C[mesh.cell_idx(i-1,j-1)])) / (8.0 * dy);
            candidates.push_back({mx, my});

            // Évaluation de l'erreur quadratique pour chaque candidat sur le stencil 3x3
            double best_err = 1.0e30;
            double best_nx = 1.0, best_ny = 0.0, best_alpha = 0.0;

            for (const auto& cand : candidates) {
                double cnorm = std::hypot(cand.first, cand.second);
                if (cnorm < 1.0e-12) continue;
                double cnx = cand.first / cnorm;
                double cny = cand.second / cnorm;

                double calpha = compute_alpha(xmin, xmax, ymin, ymax, cnx, cny, c);

                // Somme des carrés des erreurs sur les 8 voisins
                double err = 0.0;
                for (int dj = -1; dj <= 1; ++dj) {
                    for (int di = -1; di <= 1; ++di) {
                        if (di == 0 && dj == 0) continue;
                        double v_xmin = mesh.xv(i + di);
                        double v_xmax = mesh.xv(i + di + 1);
                        double v_ymin = mesh.yv(j + dj);
                        double v_ymax = mesh.yv(j + dj + 1);

                        double v_target = C[mesh.cell_idx(i + di, j + dj)] * dx * dy;
                        double v_recon  = cut_box_area(v_xmin, v_xmax, v_ymin, v_ymax, cnx, cny, calpha);
                        double diff = v_recon - v_target;
                        err += diff * diff;
                    }
                }

                if (err < best_err) {
                    best_err = err;
                    best_nx = cnx;
                    best_ny = cny;
                    best_alpha = calpha;
                }
            }

            segments[idx] = {best_nx, best_ny, best_alpha};
        }
    }
}

// =============================================================================
// RECONSTRUCTION LVIRA (PILLIOD & PUCKETT 2004)
// =============================================================================

void InterfaceReconstruction::reconstruct_lvira(const Mesh2D& mesh, const std::vector<double>& C,
                                                std::vector<InterfaceSegment>& segments) {
    // 1. Initialisation par ELVIRA pour obtenir un excellent point de départ
    reconstruct_elvira(mesh, C, segments);

    int nx = mesh.nx_local;
    int ny = mesh.ny_local;
    double dx = mesh.dx;
    double dy = mesh.dy;

    // 2. Raffinement par optimisation continue locale de l'angle theta
    for (int j = 1; j <= ny; ++j) {
        for (int i = 1; i <= nx; ++i) {
            int idx = mesh.cell_idx(i, j);
            double c = C[idx];
            if (c <= 1.0e-6 || c >= 1.0 - 1.0e-6) continue;

            double xmin = mesh.xv(i);
            double xmax = mesh.xv(i + 1);
            double ymin = mesh.yv(j);
            double ymax = mesh.yv(j + 1);

            double theta0 = std::atan2(segments[idx].ny, segments[idx].nx);

            // Fonction coût : erreur quadratique pour un angle theta donné
            auto cost_func = [&](double theta) {
                double lnx = std::cos(theta);
                double lny = std::sin(theta);
                double lalpha = compute_alpha(xmin, xmax, ymin, ymax, lnx, lny, c);

                double err = 0.0;
                for (int dj = -1; dj <= 1; ++dj) {
                    for (int di = -1; di <= 1; ++di) {
                        if (di == 0 && dj == 0) continue;
                        double v_xmin = mesh.xv(i + di);
                        double v_xmax = mesh.xv(i + di + 1);
                        double v_ymin = mesh.yv(j + dj);
                        double v_ymax = mesh.yv(j + dj + 1);

                        double v_target = C[mesh.cell_idx(i + di, j + dj)] * dx * dy;
                        double v_recon  = cut_box_area(v_xmin, v_xmax, v_ymin, v_ymax, lnx, lny, lalpha);
                        double diff = v_recon - v_target;
                        err += diff * diff;
                    }
                }
                return err;
            };

            // Recherche par Section Dorée (Golden Section Search) dans [theta0 - pi/6, theta0 + pi/6]
            const double invphi = (std::sqrt(5.0) - 1.0) / 2.0; // 0.6180339887...
            const double invphi2 = (3.0 - std::sqrt(5.0)) / 2.0; // 0.3819660113...

            double a = theta0 - 0.5;
            double b = theta0 + 0.5;
            double h = b - a;

            double c_pt = a + invphi2 * h;
            double d_pt = a + invphi * h;
            double fc = cost_func(c_pt);
            double fd = cost_func(d_pt);

            for (int iter = 0; iter < 12; ++iter) {
                if (fc < fd) {
                    b = d_pt;
                    d_pt = c_pt;
                    fd = fc;
                    h = invphi * h;
                    c_pt = a + invphi2 * h;
                    fc = cost_func(c_pt);
                } else {
                    a = c_pt;
                    c_pt = d_pt;
                    fc = fd;
                    h = invphi * h;
                    d_pt = a + invphi * h;
                    fd = cost_func(d_pt);
                }
            }

            double best_theta = 0.5 * (a + b);
            double best_nx = std::cos(best_theta);
            double best_ny = std::sin(best_theta);
            double best_alpha = compute_alpha(xmin, xmax, ymin, ymax, best_nx, best_ny, c);

            segments[idx] = {best_nx, best_ny, best_alpha};
        }
    }
}

// =============================================================================
// DISPATCHING GÉNÉRIQUE
// =============================================================================

void InterfaceReconstruction::reconstruct(ReconstructionMethod method, const Mesh2D& mesh,
                                          const std::vector<double>& C,
                                          std::vector<InterfaceSegment>& segments) {
    switch (method) {
        case ReconstructionMethod::SLIC:
            reconstruct_slic(mesh, C, segments);
            break;
        case ReconstructionMethod::PLIC:
            reconstruct_plic(mesh, C, segments);
            break;
        case ReconstructionMethod::LVIRA:
            reconstruct_lvira(mesh, C, segments);
            break;
        case ReconstructionMethod::ELVIRA:
            reconstruct_elvira(mesh, C, segments);
            break;
    }

    // Propagation des normales reconstruites vers les mailles fantomes (periodicite / bords)
    int nx = mesh.nx_local;
    int ny = mesh.ny_local;
    int i_left_src  = (!mesh.is_boundary_left  && mesh.num_procs == 1) ? nx : 1;
    int i_right_src = (!mesh.is_boundary_right && mesh.num_procs == 1) ? 1  : nx;
    int j_bot_src   = (!mesh.is_boundary_bottom && mesh.num_procs == 1) ? ny : 1;
    int j_top_src   = (!mesh.is_boundary_top    && mesh.num_procs == 1) ? 1  : ny;

    for (int j = 1; j <= ny; ++j) {
        // Maille fantome gauche i = 0
        int idx0 = mesh.cell_idx(0, j);
        double c0 = C[idx0];
        if (c0 > 1.0e-6 && c0 < 1.0 - 1.0e-6) {
            const auto& s_src = segments[mesh.cell_idx(i_left_src, j)];
            double alpha0 = compute_alpha(mesh.xv(0), mesh.xv(1), mesh.yv(j), mesh.yv(j + 1), s_src.nx, s_src.ny, c0);
            segments[idx0] = {s_src.nx, s_src.ny, alpha0};
        }
        // Maille fantome droite i = nx + 1
        int idxR = mesh.cell_idx(nx + 1, j);
        double cR = C[idxR];
        if (cR > 1.0e-6 && cR < 1.0 - 1.0e-6) {
            const auto& s_src = segments[mesh.cell_idx(i_right_src, j)];
            double alphaR = compute_alpha(mesh.xv(nx + 1), mesh.xv(nx + 2), mesh.yv(j), mesh.yv(j + 1), s_src.nx, s_src.ny, cR);
            segments[idxR] = {s_src.nx, s_src.ny, alphaR};
        }
    }

    for (int i = 0; i <= nx + 1; ++i) {
        // Maille fantome inferieure j = 0
        int idxB = mesh.cell_idx(i, 0);
        double cB = C[idxB];
        if (cB > 1.0e-6 && cB < 1.0 - 1.0e-6) {
            const auto& s_src = segments[mesh.cell_idx(i, j_bot_src)];
            double alphaB = compute_alpha(mesh.xv(i), mesh.xv(i + 1), mesh.yv(0), mesh.yv(1), s_src.nx, s_src.ny, cB);
            segments[idxB] = {s_src.nx, s_src.ny, alphaB};
        }
        // Maille fantome superieure j = ny + 1
        int idxT = mesh.cell_idx(i, ny + 1);
        double cT = C[idxT];
        if (cT > 1.0e-6 && cT < 1.0 - 1.0e-6) {
            const auto& s_src = segments[mesh.cell_idx(i, j_top_src)];
            double alphaT = compute_alpha(mesh.xv(i), mesh.xv(i + 1), mesh.yv(ny + 1), mesh.yv(ny + 2), s_src.nx, s_src.ny, cT);
            segments[idxT] = {s_src.nx, s_src.ny, alphaT};
        }
    }
}

// =============================================================================
// ÉVALUATION DES LIMITEURS DE FLUX (SUPERBEE, VAN LEER, MUSCL, HRIC)
// =============================================================================

double InterfaceReconstruction::evaluate_limiter(TvdLimiter limiter, double r, double c_donor, double c_upwind, double c_acceptor) {
    switch (limiter) {
        case TvdLimiter::SUPERBEE:
            // Limiteur Superbee de Roe (1985)
            return std::max(0.0, std::max(std::min(2.0 * r, 1.0), std::min(r, 2.0)));

        case TvdLimiter::VAN_LEER:
            // Limiteur harmonique de Van Leer (1974)
            return (r > 0.0) ? (2.0 * r / (1.0 + r)) : 0.0;

        case TvdLimiter::MUSCL:
            // Limiteur linéaire par morceaux MUSCL (Van Leer 1979)
            return std::max(0.0, std::min({2.0 * r, 0.5 * (1.0 + r), 2.0}));

        case TvdLimiter::HRIC: {
            // High Resolution Interface Capturing (Muzaferija & Peric 1999) dans l'espace NVD
            double denom = c_acceptor - c_upwind;
            if (std::abs(denom) < 1.0e-12) return 0.0;
            double c_tilde_D = (c_donor - c_upwind) / denom;

            double c_tilde_f = c_tilde_D;
            if (c_tilde_D >= 0.0 && c_tilde_D <= 1.0) {
                if (c_tilde_D < 0.5) {
                    c_tilde_f = 2.0 * c_tilde_D;
                } else {
                    c_tilde_f = 1.0;
                }
            }

            // Conversion de la valeur normalisée à la face en limiteur équivalent phi(r)
            // C_f = C_D + 0.5 * phi * (1 - C_D) dans le formalisme NVD
            if (std::abs(1.0 - c_tilde_D) > 1.0e-12) {
                double phi = 2.0 * (c_tilde_f - c_tilde_D) / (1.0 - c_tilde_D);
                return std::clamp(phi, 0.0, 2.0);
            }
            return 1.0;
        }
    }
    return 0.0;
}

} // namespace covof

\end{lstlisting}

\chapter{Transport et Advection Conservatrice VOF : Weymouth-Yue, CIAM et TVD (InterfaceTracker.hpp/.cpp)}
\section{Séparation Directionnelle d'Ordre 2 de Strang (\emph{Strang Operator Splitting})}
La classe \texttt{covof::InterfaceTracker} (\texttt{include/covof/InterfaceTracker.hpp}, \texttt{src/InterfaceTracker.cpp}, 833 lignes) résout l'équation de transport de la fraction volumique :
\begin{equation}
  \frac{\partial C}{\partial t} + \frac{\partial (u C)}{\partial x} + \frac{\partial (v C)}{\partial y} = C \left( \frac{\partial u}{\partial x} + \frac{\partial v}{\partial y} \right)
\end{equation}
en alternant l'ordre des balayages unidimensionnels entre les pas de temps pairs ($X$ puis $Y$) et impairs ($Y$ puis $X$) afin d'éliminer tout biais directionnel et de garantir l'ordre 2 temporel.

\section{Schéma de Weymouth \& Yue (2010) à Conservation de Masse Exacte}
Lors d'un balayage 1D individuel (par exemple selon $x$), le champ de vitesse 1D n'est pas à divergence nulle ($\partial u / \partial x \ne 0$, seule la somme $\partial u / \partial x + \partial v / \partial y = 0$ est nulle). Weymouth \& Yue (2010) démontrent que si le coefficient $C_{\text{dil}}$ multipliant les termes de dilatation 1D est fixé à la valeur binaire évaluée au début du pas de temps :
\begin{equation}
  C_{\text{dil}}(i, j) = \begin{cases} 1 & \text{si } C_{i,j}^n \ge 0.5 \\ 0 & \text{si } C_{i,j}^n < 0.5 \end{cases}
\end{equation}
et maintenu identique durant la passe $X$ et la passe $Y$ :
\begin{align}
  C_{i,j}^* &= C_{i,j}^n - \frac{\Delta t}{\Delta x} \left( F_{i+1/2, j}^x - F_{i-1/2, j}^x \right) + C_{\text{dil}}(i,j) \frac{\Delta t}{\Delta x} \left( u_{i+1/2, j} - u_{i-1/2, j} \right) \\
  C_{i,j}^{n+1} &= C_{i,j}^* - \frac{\Delta t}{\Delta y} \left( F_{i, j+1/2}^y - F_{i, j-1/2}^y \right) + C_{\text{dil}}(i,j) \frac{\Delta t}{\Delta y} \left( v_{i, j+1/2} - v_{i, j-1/2} \right)
\end{align}
alors, d'une part $C_{i,j}^{n+1} \in [0, 1]$ est garanti sans écrêtage (*no-clipping*) sous $\text{CFL} < 0.5$, et d'autre part la somme des deux termes de dilatation vaut exactement $C_{\text{dil}} \Delta t (\nabla_h \cdot \mathbf{u})_{i,j} = 0$, assurant une **conservation de la masse exacte au bit près**.

\section{Répertoire des Classes, Structures et Méthodes de \texttt{include/covof/InterfaceTracker.hpp}}
\begin{center}\small
\begin{tabular}{clll}
\toprule
\textbf{Ligne} & \textbf{Catégorie} & \textbf{Identifiant C++17} & \textbf{Signature / Rôle} \\
\midrule
14 & class & \texttt{InterfaceTracker} & Déclaration de type / classe / structure \\
43 & Méthode / Fonction & \texttt{init} & Type retour : void \\
69 & Méthode / Fonction & \texttt{set\_initial\_condition} & Type retour : void \\
74 & Méthode / Fonction & \texttt{advect} & Type retour : void \\
79 & Méthode / Fonction & \texttt{advect\_tvd} & Type retour : void \\
84 & Méthode / Fonction & \texttt{advect\_ciam} & Type retour : void \\
89 & Méthode / Fonction & \texttt{advect\_weymouth\_yue} & Type retour : void \\
94 & Méthode / Fonction & \texttt{update\_physical\_properties} & Type retour : void \\
99 & Méthode / Fonction & \texttt{compute\_curvature\_and\_CSF} & Type retour : void \\
\bottomrule
\end{tabular}
\end{center}

\section{Code Source Intégral Annoté : \texttt{include/covof/InterfaceTracker.hpp}}
Le listing ci-dessous présente l'intégralité des 102 lignes du fichier \texttt{include/covof/InterfaceTracker.hpp} de l'implémentation C++17 :

\begin{lstlisting}[language={[11]C++}]
#pragma once

#include "Mesh.hpp"
#include "Field.hpp"
#include "Config.hpp"
#include "InterfaceReconstruction.hpp"

namespace covof {
namespace fugu = covof; // Alias rétrocompatibilité

/**
 * @brief Module de capture d'interface et de calcul de la tension de surface (Modèle CSF)
 */
class InterfaceTracker {
public:
    FluidProperties fluid1;
    FluidProperties fluid2;
    // Déclaration de variable typée
    double sigma = 0.0;
    // Déclaration de variable typée
    bool enable_surface_tension = true;
    // Déclaration de variable typée
    bool enable_two_phase = false;
    // Déclaration de variable typée
    std::string viscosity_model = "CONSTANT";
    // Déclaration de variable typée
    double visc_a = 2.0;
    // Déclaration de variable typée
    double visc_e = 0.01;
    // Déclaration de variable typée
    double visc_mu0 = 1.85e-5;

    // Configuration VOF avancée (SLIC, PLIC, LVIRA, ELVIRA, CIAM, Weymouth & Yue, limiteurs)
    VofConfig vof_cfg;
    // Segments d'interface géométriques reconstruits
    std::vector<InterfaceSegment> interface_segments;
    // Compteur de pas pour Strang splitting alterné
    int step_count = 0;

    // Calcul et assignation numérique
    InterfaceTracker() = default;

    void init(const Config& cfg) {
        // Calcul et assignation numérique
        fluid1 = cfg.fluid1;
        // Calcul et assignation numérique
        fluid2 = cfg.fluid2;
        // Calcul et assignation numérique
        sigma = cfg.sigma;
        // Calcul et assignation numérique
        enable_surface_tension = cfg.enable_surface_tension;
        // Calcul et assignation numérique
        enable_two_phase = cfg.enable_two_phase;
        // Calcul et assignation numérique
        viscosity_model = cfg.viscosity_model;
        // Calcul et assignation numérique
        visc_a = cfg.visc_a;
        // Calcul et assignation numérique
        visc_e = cfg.visc_e;
        // Calcul et assignation numérique
        visc_mu0 = cfg.visc_mu0;
        // Options VOF
        vof_cfg = cfg.vof;
    }

    /**
     * @brief Initialise le champ de fraction volumique C selon la configuration
     */
    void set_initial_condition(const Mesh2D& mesh, const InitialConditionConfig& ic, FlowFields& fields);

    /**
     * @brief Transporte le champ C selon l'algorithme choisi (VOF-TVD, CIAM, Weymouth & Yue)
     */
    void advect(const Mesh2D& mesh, FlowFields& fields, double dt);

    /**
     * @brief Transport algébrique VOF-TVD avec limiteurs (Superbee, Van Leer, MUSCL, HRIC)
     */
    void advect_tvd(const Mesh2D& mesh, FlowFields& fields, double dt);

    /**
     * @brief Transport géométrique directionnel CIAM (Calcul d'Interface Affine par Morceaux)
     */
    void advect_ciam(const Mesh2D& mesh, FlowFields& fields, double dt);

    /**
     * @brief Transport géométrique unsplit multidimensionnel de Weymouth et Yue (2010)
     */
    void advect_weymouth_yue(const Mesh2D& mesh, FlowFields& fields, double dt);

    /**
     * @brief Met à jour les propriétés physiques effectives (rho, mu) aux centres, faces et nœuds
     */
    void update_physical_properties(const Mesh2D& mesh, FlowFields& fields);

    /**
     * @brief Calcule la normale, la courbure kappa et la force capillaire CSF équilibrée
     */
    void compute_curvature_and_CSF(const Mesh2D& mesh, FlowFields& fields);
};

} // namespace covof

\end{lstlisting}

\section{Répertoire des Classes, Structures et Méthodes de \texttt{src/InterfaceTracker.cpp}}
\begin{center}\small
\begin{tabular}{clll}
\toprule
\textbf{Ligne} & \textbf{Catégorie} & \textbf{Identifiant C++17} & \textbf{Signature / Rôle} \\
\midrule
9 & Méthode / Fonction & \texttt{InterfaceTracker::set\_initial\_condition} & Type retour : void \\
144 & Méthode / Fonction & \texttt{InterfaceTracker::update\_physical\_properties} & Type retour : void \\
269 & Méthode / Fonction & \texttt{InterfaceTracker::compute\_curvature\_and\_CSF} & Type retour : void \\
370 & Méthode / Fonction & \texttt{InterfaceTracker::advect} & Type retour : void \\
405 & Méthode / Fonction & \texttt{InterfaceTracker::advect\_tvd} & Type retour : void \\
415 & Méthode / Fonction & \texttt{flux\_x} & Type retour : std::vector$<$double$>$ \\
444 & Méthode / Fonction & \texttt{flux\_y} & Type retour : std::vector$<$double$>$ \\
489 & Méthode / Fonction & \texttt{InterfaceTracker::advect\_ciam} & Type retour : void \\
501 & Méthode / Fonction & \texttt{flux\_vol\_x} & Type retour : std::vector$<$double$>$ \\
561 & Méthode / Fonction & \texttt{flux\_vol\_y} & Type retour : std::vector$<$double$>$ \\
633 & Méthode / Fonction & \texttt{InterfaceTracker::advect\_weymouth\_yue} & Type retour : void \\
644 & Méthode / Fonction & \texttt{C\_d} & Type retour : std::vector$<$double$>$ \\
655 & Méthode / Fonction & \texttt{f\_vol\_x} & Type retour : std::vector$<$double$>$ \\
698 & Méthode / Fonction & \texttt{f\_vol\_y} & Type retour : std::vector$<$double$>$ \\
\bottomrule
\end{tabular}
\end{center}

\section{Code Source Intégral Annoté : \texttt{src/InterfaceTracker.cpp}}
Le listing ci-dessous présente l'intégralité des 752 lignes du fichier \texttt{src/InterfaceTracker.cpp} de l'implémentation C++17 :

\begin{lstlisting}[language={[11]C++}]
#include "covof/InterfaceTracker.hpp"
#include "covof/MpiComm.hpp"
#include <cmath>
#include <algorithm>

namespace covof {
namespace fugu = covof; // Alias rétrocompatibilité

void InterfaceTracker::set_initial_condition(const Mesh2D& mesh, const InitialConditionConfig& ic, FlowFields& fields) {
    // Déclaration de variable typée
    double eps = 1.5 * std::max(mesh.dx, mesh.dy);

    // Boucle spatiale sur les cellules locales selon l'axe vertical Y (ordonnées)
    for (int j = 0; j <= mesh.ny_local + 1; ++j) {
        // Déclaration de variable typée
        double y = mesh.yc(j);
        // Boucle spatiale sur les cellules locales selon l'axe horizontal X (abscisses)
        for (int i = 0; i <= mesh.nx_local + 1; ++i) {
            // Déclaration de variable typée
            double x = mesh.xc(i);
            // Déclaration de variable typée
            int idx = mesh.cell_idx(i, j);

            // Évaluation de la condition de contrôle
            if (ic.type == InitialConditionType::SINGLE_PHASE) {
                // Calcul de la résolution géométrique de cellule
                fields.C[idx] = 1.0;
            // Calcul et assignation numérique
            } else if (ic.type == InitialConditionType::DROPLET) {
                // Déclaration de variable typée
                double dist = std::sqrt((x - ic.center_x) * (x - ic.center_x) +
                                        (y - ic.center_y) * (y - ic.center_y)) - ic.radius;
                // Évaluation de la condition de contrôle
                if (dist < -eps) {
                    // Calcul de la résolution géométrique de cellule
                    fields.C[idx] = 1.0;
                } else if (dist > eps) {
                    // Calcul de la résolution géométrique de cellule
                    fields.C[idx] = 0.0;
                } else {
                    // Calcul de la résolution géométrique de cellule
                    fields.C[idx] = 0.5 - dist / (2.0 * eps) - (1.0 / (2.0 * M_PI)) * std::sin(M_PI * dist / eps);
                }
            // Calcul et assignation numérique
            } else if (ic.type == InitialConditionType::BUBBLE) {
                // Déclaration de variable typée
                double dist = std::sqrt((x - ic.center_x) * (x - ic.center_x) +
                                        (y - ic.center_y) * (y - ic.center_y)) - ic.radius;
                // Évaluation de la condition de contrôle
                if (dist < -eps) {
                    // Calcul de la résolution géométrique de cellule
                    fields.C[idx] = 0.0; // Intérieur bulle : fluide 2
                } else if (dist > eps) {
                    // Calcul de la résolution géométrique de cellule
                    fields.C[idx] = 1.0; // Extérieur : fluide 1 (eau)
                } else {
                    // Calcul de la résolution géométrique de cellule
                    fields.C[idx] = 0.5 + dist / (2.0 * eps) + (1.0 / (2.0 * M_PI)) * std::sin(M_PI * dist / eps);
                }
            // Calcul et assignation numérique
            } else if (ic.type == InitialConditionType::DAM_BREAK) {
                // Évaluation de la condition de contrôle
                if (x >= ic.box_x0 && x <= ic.box_x1 && y >= ic.box_y0 && y <= ic.box_y1) {
                    // Calcul de la résolution géométrique de cellule
                    fields.C[idx] = 1.0;
                } else {
                    // Calcul de la résolution géométrique de cellule
                    fields.C[idx] = 0.0;
                }
            // Calcul et assignation numérique
            } else if (ic.type == InitialConditionType::SHEAR_LAYER) {
                // Déclaration de variable typée
                double int_y = ic.center_y + ic.pert_amp * std::sin(2.0 * M_PI * ic.pert_k * x / mesh.lx);
                // Calcul de la résolution géométrique de cellule
                fields.C[idx] = (y <= int_y) ? 1.0 : 0.0;
            // Calcul et assignation numérique pour nappe liquide (Cas 3.2 Fluids 2021)
            } else if (ic.type == InitialConditionType::LIQUID_SHEET) {
                // Déclaration de la position centrale et demi-épaisseur de nappe
                double y_c = (ic.center_y > 0.0) ? ic.center_y : 0.5 * mesh.ly;
                // Calcul de la demi-épaisseur de la nappe liquide
                double half_thick = 0.5 * ic.thickness;
                // Calcul de la distance géométrique par rapport à la surface de la nappe
                double dist = std::abs(y - y_c) - half_thick;
                // Évaluation de l'indicatrice : C = 0 dans le liquide (fluide 2), C = 1 dans le gaz (fluide 1)
                if (dist < -eps) {
                    // Calcul de la résolution géométrique de cellule
                    fields.C[idx] = 0.0;
                // Branche alternative conditionnelle
                } else if (dist > eps) {
                    // Calcul de la résolution géométrique de cellule
                    fields.C[idx] = 1.0;
                // Branche par défaut pour transition continue
                } else {
                    // Calcul de la résolution géométrique de cellule
                    fields.C[idx] = 0.5 + dist / (2.0 * eps) + (1.0 / (2.0 * M_PI)) * std::sin(M_PI * dist / eps);
                }
            // Calcul et assignation numérique pour jet liquide visqueux (Cas 3.3 Fluids 2021)
            } else if (ic.type == InitialConditionType::VISCOUS_JET) {
                // Déclaration du centre et de la demi-largeur de l'injecteur
                double x_c = (ic.center_x > 0.0) ? ic.center_x : 0.5 * mesh.lx;
                // Calcul de la demi-largeur de l'injecteur
                double half_w = (ic.jet_width > 0.0) ? 0.5 * ic.jet_width : 0.04;
                // Déclaration de la hauteur initiale de l'amorce de jet sous la paroi supérieure
                double stub_height = 2.0 * mesh.dy;
                // Évaluation de l'indicatrice : C = 0 dans l'amorce liquide, C = 1 dans la cavité de gaz
                if (y >= mesh.ly - stub_height && std::abs(x - x_c) <= half_w) {
                    // Calcul de la résolution géométrique de cellule
                    fields.C[idx] = 0.0;
                } else {
                    // Calcul de la résolution géométrique de cellule
                    fields.C[idx] = 1.0;
                }
            // Calcul et assignation numérique
            } else if (ic.type == InitialConditionType::VORTEX_SHEAR) {
                // Déclaration de variable typée
                double dist = std::sqrt((x - ic.center_x) * (x - ic.center_x) +
                                        (y - ic.center_y) * (y - ic.center_y)) - ic.radius;
                // Évaluation de la condition de contrôle
                if (dist < -eps) {
                    // Calcul de la résolution géométrique de cellule
                    fields.C[idx] = 1.0;
                } else if (dist > eps) {
                    // Calcul de la résolution géométrique de cellule
                    fields.C[idx] = 0.0;
                } else {
                    // Calcul de la résolution géométrique de cellule
                    fields.C[idx] = 0.5 - dist / (2.0 * eps) - (1.0 / (2.0 * M_PI)) * std::sin(M_PI * dist / eps);
                }
            // Calcul et assignation numérique
            } else if (ic.type == InitialConditionType::RAYLEIGH_TAYLOR) {
                // Déclaration de variable typée
                double int_y = ic.center_y + ic.pert_amp * std::cos(2.0 * M_PI * ic.pert_k * x / mesh.lx);
                // Calcul de la résolution géométrique de cellule
                fields.C[idx] = (y >= int_y) ? 1.0 : 0.0;
            }
        }
    }

    MpiComm::sync_cell_scalar(mesh, fields.C);
    update_physical_properties(mesh, fields);
    compute_curvature_and_CSF(mesh, fields);
}

void InterfaceTracker::update_physical_properties(const Mesh2D& mesh, FlowFields& fields) {
    // Déclaration de variable typée
    int nx = mesh.nx_local;
    // Déclaration de variable typée
    int ny = mesh.ny_local;

    // Pour un écoulement monophasique, propriétés uniformes exactes
    if (!enable_two_phase) {
        std::fill(fields.rho.begin(), fields.rho.end(), fluid1.rho);
        std::fill(fields.rho_u.begin(), fields.rho_u.end(), fluid1.rho);
        std::fill(fields.rho_v.begin(), fields.rho_v.end(), fluid1.rho);
        // Évaluation de la condition de contrôle
        if (viscosity_model == "EXPONENTIAL") {
            // Boucle spatiale sur les cellules locales selon l'axe vertical Y (ordonnées)
            for (int j = 0; j <= ny + 1; ++j) {
                // Déclaration de variable typée
                double y = mesh.yc(j);
                // Déclaration de variable typée
                double z = y - visc_e;
                // Déclaration de variable typée
                double mu_val = visc_mu0 * std::exp(visc_a * (z * z) / (visc_e * visc_e));
                // Boucle spatiale sur les cellules locales selon l'axe horizontal X (abscisses)
                for (int i = 0; i <= nx + 1; ++i) {
                    // Calcul de la résolution géométrique de cellule
                    fields.mu[mesh.cell_idx(i, j)] = mu_val;
                }
            }
            // Boucle spatiale sur les cellules locales selon l'axe vertical Y (ordonnées)
            for (int j = 0; j <= ny + 1; ++j) {
                // Déclaration de variable typée
                double y = mesh.yv(j);
                // Déclaration de variable typée
                double z = y - visc_e;
                // Déclaration de variable typée
                double mu_node_val = visc_mu0 * std::exp(visc_a * (z * z) / (visc_e * visc_e));
                // Boucle spatiale sur les cellules locales selon l'axe horizontal X (abscisses)
                for (int i = 0; i <= nx + 1; ++i) {
                    // Calcul de la résolution géométrique de cellule
                    fields.mu_node[mesh.node_idx(i, j)] = mu_node_val;
                }
            }
        } else {
            std::fill(fields.mu.begin(), fields.mu.end(), fluid1.mu);
            std::fill(fields.mu_node.begin(), fields.mu_node.end(), fluid1.mu);
        }
        // Retour de fonction
        return;
    }

    // Remplissage des coins pour C
    fields.C[mesh.cell_idx(0, 0)]           = fields.C[mesh.cell_idx(1, 1)];
    // Calcul de la résolution géométrique de cellule
    fields.C[mesh.cell_idx(nx + 1, 0)]      = fields.C[mesh.cell_idx(nx, 1)];
    // Calcul de la résolution géométrique de cellule
    fields.C[mesh.cell_idx(0, ny + 1)]      = fields.C[mesh.cell_idx(1, ny)];
    // Calcul de la résolution géométrique de cellule
    fields.C[mesh.cell_idx(nx + 1, ny + 1)]  = fields.C[mesh.cell_idx(nx, ny)];

    // 1. Masse volumique et viscosité au centre des cellules
    for (int j = 0; j <= ny + 1; ++j) {
        // Boucle spatiale sur les cellules locales selon l'axe horizontal X (abscisses)
        for (int i = 0; i <= nx + 1; ++i) {
            // Déclaration de variable typée
            int idx = mesh.cell_idx(i, j);
            // Déclaration de variable typée
            double c = std::clamp(fields.C[idx], 0.0, 1.0);
            // Calcul de la résolution géométrique de cellule
            fields.rho[idx] = c * fluid1.rho + (1.0 - c) * fluid2.rho;
            // Calcul de la résolution géométrique de cellule
            fields.mu[idx]  = c * fluid1.mu  + (1.0 - c) * fluid2.mu;
        }
    }

    // Calcul de la résolution géométrique de cellule
    fields.mu[mesh.cell_idx(0, 0)]           = fields.mu[mesh.cell_idx(1, 1)];
    // Calcul de la résolution géométrique de cellule
    fields.mu[mesh.cell_idx(nx + 1, 0)]      = fields.mu[mesh.cell_idx(nx, 1)];
    // Calcul de la résolution géométrique de cellule
    fields.mu[mesh.cell_idx(0, ny + 1)]      = fields.mu[mesh.cell_idx(1, ny)];
    // Calcul de la résolution géométrique de cellule
    fields.mu[mesh.cell_idx(nx + 1, ny + 1)]  = fields.mu[mesh.cell_idx(nx, ny)];

    // 2. Masse volumique sur les faces u
    for (int j = 1; j <= ny; ++j) {
        // Boucle spatiale sur les cellules locales selon l'axe horizontal X (abscisses)
        for (int i = 1; i <= nx + 1; ++i) {
            // Déclaration de variable typée
            int idx_u = mesh.u_idx(i, j);
            // Calcul de la résolution géométrique de cellule
            fields.rho_u[idx_u] = 0.5 * (fields.rho[mesh.cell_idx(i - 1, j)] + fields.rho[mesh.cell_idx(i, j)]);
        }
    }

    // 3. Masse volumique sur les faces v
    for (int j = 1; j <= ny + 1; ++j) {
        // Boucle spatiale sur les cellules locales selon l'axe horizontal X (abscisses)
        for (int i = 1; i <= nx; ++i) {
            // Déclaration de variable typée
            int idx_v = mesh.v_idx(i, j);
            // Calcul de la résolution géométrique de cellule
            fields.rho_v[idx_v] = 0.5 * (fields.rho[mesh.cell_idx(i, j - 1)] + fields.rho[mesh.cell_idx(i, j)]);
        }
    }

    // 4. Viscosité harmonique aux nœuds (coins)
    for (int j = 1; j <= ny + 1; ++j) {
        // Boucle spatiale sur les cellules locales selon l'axe horizontal X (abscisses)
        for (int i = 1; i <= nx + 1; ++i) {
            // Déclaration de variable typée
            int n_idx = mesh.node_idx(i, j);
            // Déclaration de variable typée
            double m_SW = fields.mu[mesh.cell_idx(i - 1, j - 1)];
            // Déclaration de variable typée
            double m_SE = fields.mu[mesh.cell_idx(i, j - 1)];
            // Déclaration de variable typée
            double m_NW = fields.mu[mesh.cell_idx(i - 1, j)];
            // Déclaration de variable typée
            double m_NE = fields.mu[mesh.cell_idx(i, j)];

            // Calcul de la résolution géométrique de cellule
            fields.mu_node[n_idx] = 4.0 / (1.0 / m_SW + 1.0 / m_SE + 1.0 / m_NW + 1.0 / m_NE);
        }
    }
}

void InterfaceTracker::compute_curvature_and_CSF(const Mesh2D& mesh, FlowFields& fields) {
    // Évaluation de la condition de contrôle
    if (!enable_surface_tension || sigma == 0.0) {
        std::fill(fields.F_sigma_x.begin(), fields.F_sigma_x.end(), 0.0);
        std::fill(fields.F_sigma_y.begin(), fields.F_sigma_y.end(), 0.0);
        // Retour de fonction
        return;
    }

    // Déclaration de variable typée
    int nx = mesh.nx_local;
    // Déclaration de variable typée
    int ny = mesh.ny_local;
    // Déclaration de variable typée
    double dx = mesh.dx;
    // Déclaration de variable typée
    double dy = mesh.dy;

    // 1. Calcul du gradient et du vecteur normal unitaire aux centres
    for (int j = 1; j <= ny; ++j) {
        // Boucle spatiale sur les cellules locales selon l'axe horizontal X (abscisses)
        for (int i = 1; i <= nx; ++i) {
            // Déclaration de variable typée
            int idx = mesh.cell_idx(i, j);
            // Déclaration de variable typée
            double dC_dx = (fields.C[mesh.cell_idx(i + 1, j)] - fields.C[mesh.cell_idx(i - 1, j)]) / (2.0 * dx);
            // Déclaration de variable typée
            double dC_dy = (fields.C[mesh.cell_idx(i, j + 1)] - fields.C[mesh.cell_idx(i, j - 1)]) / (2.0 * dy);
            // Déclaration de variable typée
            double grad_mag = std::sqrt(dC_dx * dC_dx + dC_dy * dC_dy);
            // Évaluation de la condition de contrôle
            if (grad_mag > 1.0e-3 / dx) {
                // Calcul de la résolution géométrique de cellule
                fields.norm_x[idx] = dC_dx / grad_mag;
                // Calcul de la résolution géométrique de cellule
                fields.norm_y[idx] = dC_dy / grad_mag;
            } else {
                // Calcul de la résolution géométrique de cellule
                fields.norm_x[idx] = 0.0;
                // Calcul de la résolution géométrique de cellule
                fields.norm_y[idx] = 0.0;
            }
        }
    }
    MpiComm::sync_cell_scalar(mesh, fields.norm_x);
    MpiComm::sync_cell_scalar(mesh, fields.norm_y);

    // 2. Calcul de la courbure kappa = -div(n)
    for (int j = 1; j <= ny; ++j) {
        // Boucle spatiale sur les cellules locales selon l'axe horizontal X (abscisses)
        for (int i = 1; i <= nx; ++i) {
            // Déclaration de variable typée
            int idx = mesh.cell_idx(i, j);
            // Évaluation de la condition de contrôle
            if (fields.norm_x[idx] == 0.0 && fields.norm_y[idx] == 0.0) {
                // Calcul de la résolution géométrique de cellule
                fields.kappa[idx] = 0.0;
            } else {
                // Déclaration de variable typée
                double dnx_dx = (fields.norm_x[mesh.cell_idx(i + 1, j)] - fields.norm_x[mesh.cell_idx(i - 1, j)]) / (2.0 * dx);
                // Déclaration de variable typée
                double dny_dy = (fields.norm_y[mesh.cell_idx(i, j + 1)] - fields.norm_y[mesh.cell_idx(i, j - 1)]) / (2.0 * dy);
                // Calcul de la résolution géométrique de cellule
                fields.kappa[idx] = -(dnx_dx + dny_dy);
            }
        }
    }
    MpiComm::sync_cell_scalar(mesh, fields.kappa);

    // 3. Force capillaire CSF équilibrée (balanced-force) sur les faces décalées
    // Selon X sur faces u
    for (int j = 1; j <= ny; ++j) {
        // Boucle spatiale sur les cellules locales selon l'axe horizontal X (abscisses)
        for (int i = 1; i <= nx + 1; ++i) {
            // Déclaration de variable typée
            int idx_u = mesh.u_idx(i, j);
            // Déclaration de variable typée
            double kap_avg = 0.5 * (fields.kappa[mesh.cell_idx(i - 1, j)] + fields.kappa[mesh.cell_idx(i, j)]);
            // Déclaration de variable typée
            double grad_C = (fields.C[mesh.cell_idx(i, j)] - fields.C[mesh.cell_idx(i - 1, j)]) / dx;
            // Calcul de la résolution géométrique de cellule
            fields.F_sigma_x[idx_u] = sigma * kap_avg * grad_C;
        }
    }

    // Selon Y sur faces v
    for (int j = 1; j <= ny + 1; ++j) {
        // Boucle spatiale sur les cellules locales selon l'axe horizontal X (abscisses)
        for (int i = 1; i <= nx; ++i) {
            // Déclaration de variable typée
            int idx_v = mesh.v_idx(i, j);
            // Déclaration de variable typée
            double kap_avg = 0.5 * (fields.kappa[mesh.cell_idx(i, j - 1)] + fields.kappa[mesh.cell_idx(i, j)]);
            // Déclaration de variable typée
            double grad_C = (fields.C[mesh.cell_idx(i, j)] - fields.C[mesh.cell_idx(i, j - 1)]) / dy;
            // Calcul de la résolution géométrique de cellule
            fields.F_sigma_y[idx_v] = sigma * kap_avg * grad_C;
        }
    }
}

void InterfaceTracker::advect(const Mesh2D& mesh, FlowFields& fields, double dt) {
    // Sélection de la méthode d'advection configurée par l'utilisateur
    AdvectionMethod adv_method = InterfaceReconstruction::parse_advection_method(vof_cfg.advection_method);

    // Initialisation du vecteur de segments d'interface si nécessaire
    int total_cells = (mesh.nx_local + 2) * (mesh.ny_local + 2);
    if (static_cast<int>(interface_segments.size()) != total_cells) {
        interface_segments.resize(total_cells);
    }

    // Répartition vers l'algorithme d'advection sélectionné
    switch (adv_method) {
        case AdvectionMethod::CIAM:
            advect_ciam(mesh, fields, dt);
            break;
        case AdvectionMethod::WEYMOUTH_YUE:
            advect_weymouth_yue(mesh, fields, dt);
            break;
        case AdvectionMethod::VOF_TVD:
        default:
            advect_tvd(mesh, fields, dt);
            break;
    }

    // Incrémentation du compteur de pas pour Strang splitting alterné
    step_count++;

    // Synchronisation MPI des fractions volumiques sur les mailles fantômes
    MpiComm::sync_cell_scalar(mesh, fields.C);

    // Mise à jour des grandeurs physiques dérivées (masse volumique, viscosité, tension de surface)
    update_physical_properties(mesh, fields);
    compute_curvature_and_CSF(mesh, fields);
}

void InterfaceTracker::advect_tvd(const Mesh2D& mesh, FlowFields& fields, double dt) {
    int nx = mesh.nx_local;
    int ny = mesh.ny_local;
    double dx = mesh.dx;
    double dy = mesh.dy;

    std::vector<double> C_new = fields.C;
    TvdLimiter limiter = InterfaceReconstruction::parse_tvd_limiter(vof_cfg.tvd_limiter);

    // Calcul des flux de fraction volumique sur les faces X
    std::vector<double> flux_x((nx + 2) * (ny + 2), 0.0);
    for (int j = 1; j <= ny; ++j) {
        for (int i = 1; i <= nx + 1; ++i) {
            double u_face = fields.u[mesh.u_idx(i, j)];
            double c_L = fields.C[mesh.cell_idx(i - 1, j)];
            double c_R = fields.C[mesh.cell_idx(i, j)];

            double cfl_x = std::min(1.0, std::abs(u_face) * dt / dx);
            if (u_face >= 0.0) {
                int i_up = std::max(0, i - 2);
                double c_U = fields.C[mesh.cell_idx(i_up, j)];
                double dC_up = c_L - c_U;
                double dC_loc = c_R - c_L;
                double r = (std::abs(dC_loc) > 1.0e-12) ? dC_up / dC_loc : 0.0;
                double phi = InterfaceReconstruction::evaluate_limiter(limiter, r, c_L, c_U, c_R);
                flux_x[mesh.cell_idx(i, j)] = u_face * (c_L + 0.5 * phi * (1.0 - cfl_x) * dC_loc);
            } else {
                int i_down = std::min(nx + 1, i + 1);
                double c_U = fields.C[mesh.cell_idx(i_down, j)];
                double dC_up = c_U - c_R;
                double dC_loc = c_R - c_L;
                double r = (std::abs(dC_loc) > 1.0e-12) ? dC_up / dC_loc : 0.0;
                double phi = InterfaceReconstruction::evaluate_limiter(limiter, r, c_R, c_U, c_L);
                flux_x[mesh.cell_idx(i, j)] = u_face * (c_R - 0.5 * phi * (1.0 - cfl_x) * dC_loc);
            }
        }
    }

    // Calcul des flux sur les faces Y
    std::vector<double> flux_y((nx + 2) * (ny + 2), 0.0);
    for (int j = 1; j <= ny + 1; ++j) {
        for (int i = 1; i <= nx; ++i) {
            double v_face = fields.v[mesh.v_idx(i, j)];
            double c_B = fields.C[mesh.cell_idx(i, j - 1)];
            double c_T = fields.C[mesh.cell_idx(i, j)];
            double cfl_y = std::min(1.0, std::abs(v_face) * dt / dy);

            if (v_face >= 0.0) {
                int j_up = std::max(0, j - 2);
                double c_U = fields.C[mesh.cell_idx(i, j_up)];
                double dC_up = c_B - c_U;
                double dC_loc = c_T - c_B;
                double r = (std::abs(dC_loc) > 1.0e-12) ? dC_up / dC_loc : 0.0;
                double phi = InterfaceReconstruction::evaluate_limiter(limiter, r, c_B, c_U, c_T);
                flux_y[mesh.cell_idx(i, j)] = v_face * (c_B + 0.5 * phi * (1.0 - cfl_y) * dC_loc);
            } else {
                int j_down = std::min(ny + 1, j + 1);
                double c_U = fields.C[mesh.cell_idx(i, j_down)];
                double dC_up = c_U - c_T;
                double dC_loc = c_T - c_B;
                double r = (std::abs(dC_loc) > 1.0e-12) ? dC_up / dC_loc : 0.0;
                double phi = InterfaceReconstruction::evaluate_limiter(limiter, r, c_T, c_U, c_B);
                flux_y[mesh.cell_idx(i, j)] = v_face * (c_T - 0.5 * phi * (1.0 - cfl_y) * dC_loc);
            }
        }
    }

    // Mise à jour conservative directe
    for (int j = 1; j <= ny; ++j) {
        for (int i = 1; i <= nx; ++i) {
            int idx = mesh.cell_idx(i, j);
            double fx_in  = flux_x[mesh.cell_idx(i, j)];
            double fx_out = flux_x[mesh.cell_idx(i + 1, j)];
            double fy_in  = flux_y[mesh.cell_idx(i, j)];
            double fy_out = flux_y[mesh.cell_idx(i, j + 1)];

            C_new[idx] = fields.C[idx] - (dt / dx) * (fx_out - fx_in) - (dt / dy) * (fy_out - fy_in);
            C_new[idx] = std::clamp(C_new[idx], 0.0, 1.0);
        }
    }

    fields.C = C_new;
}

void InterfaceTracker::advect_ciam(const Mesh2D& mesh, FlowFields& fields, double dt) {
    int nx = mesh.nx_local;
    int ny = mesh.ny_local;
    double dx = mesh.dx;
    double dy = mesh.dy;

    ReconstructionMethod recon_method = InterfaceReconstruction::parse_reconstruction_method(vof_cfg.reconstruction_method);
    bool x_first = (step_count % 2 == 0);

    // Fonction lambda pour un balayage directionnel X
    auto sweep_x = [&](std::vector<double>& C_in, std::vector<double>& C_out) {
        InterfaceReconstruction::reconstruct(recon_method, mesh, C_in, interface_segments);
        std::vector<double> flux_vol_x((nx + 2) * (ny + 2), 0.0);

        for (int j = 1; j <= ny; ++j) {
            for (int i = 1; i <= nx + 1; ++i) {
                double u_face = fields.u[mesh.u_idx(i, j)];
                double delta_x = std::abs(u_face) * dt;
                if (delta_x <= 1.0e-14) continue;

                if (u_face > 0.0) {
                    // Donneur cellule i-1
                    int donor_idx = mesh.cell_idx(i - 1, j);
                    double x_face_pos = mesh.xv(i);
                    double box_xmin = x_face_pos - delta_x;
                    double box_xmax = x_face_pos;
                    double box_ymin = mesh.yv(j);
                    double box_ymax = mesh.yv(j + 1);

                    const auto& seg = interface_segments[donor_idx];
                    flux_vol_x[mesh.cell_idx(i, j)] = InterfaceReconstruction::cut_box_area(box_xmin, box_xmax, box_ymin, box_ymax,
                                                                                           seg.nx, seg.ny, seg.alpha);
                } else {
                    // Donneur cellule i
                    int donor_idx = mesh.cell_idx(i, j);
                    double x_face_pos = mesh.xv(i);
                    double box_xmin = x_face_pos;
                    double box_xmax = x_face_pos + delta_x;
                    double box_ymin = mesh.yv(j);
                    double box_ymax = mesh.yv(j + 1);

                    const auto& seg = interface_segments[donor_idx];
                    flux_vol_x[mesh.cell_idx(i, j)] = -InterfaceReconstruction::cut_box_area(box_xmin, box_xmax, box_ymin, box_ymax,
                                                                                            seg.nx, seg.ny, seg.alpha);
                }
            }
        }

        for (int j = 1; j <= ny; ++j) {
            for (int i = 1; i <= nx; ++i) {
                int idx = mesh.cell_idx(i, j);
                double f_in  = flux_vol_x[mesh.cell_idx(i, j)];
                double f_out = flux_vol_x[mesh.cell_idx(i + 1, j)];
                double u_in  = fields.u[mesh.u_idx(i, j)];
                double u_out = fields.u[mesh.u_idx(i + 1, j)];

                // Mise à jour conservative 1D avec compensation de divergence
                double denom = dx - dt * (u_out - u_in);
                if (denom > 1.0e-12) {
                    C_out[idx] = (C_in[idx] * dx - (f_out - f_in) / dy) / denom;
                } else {
                    C_out[idx] = C_in[idx] - (f_out - f_in) / (dx * dy);
                }
                C_out[idx] = std::clamp(C_out[idx], 0.0, 1.0);
            }
        }
        MpiComm::sync_cell_scalar(mesh, C_out);
    };

    // Fonction lambda pour un balayage directionnel Y
    auto sweep_y = [&](std::vector<double>& C_in, std::vector<double>& C_out) {
        InterfaceReconstruction::reconstruct(recon_method, mesh, C_in, interface_segments);
        std::vector<double> flux_vol_y((nx + 2) * (ny + 2), 0.0);

        for (int j = 1; j <= ny + 1; ++j) {
            for (int i = 1; i <= nx; ++i) {
                double v_face = fields.v[mesh.v_idx(i, j)];
                double delta_y = std::abs(v_face) * dt;
                if (delta_y <= 1.0e-14) continue;

                if (v_face > 0.0) {
                    // Donneur cellule j-1
                    int donor_idx = mesh.cell_idx(i, j - 1);
                    double y_face_pos = mesh.yv(j);
                    double box_xmin = mesh.xv(i);
                    double box_xmax = mesh.xv(i + 1);
                    double box_ymin = y_face_pos - delta_y;
                    double box_ymax = y_face_pos;

                    const auto& seg = interface_segments[donor_idx];
                    flux_vol_y[mesh.cell_idx(i, j)] = InterfaceReconstruction::cut_box_area(box_xmin, box_xmax, box_ymin, box_ymax,
                                                                                           seg.nx, seg.ny, seg.alpha);
                } else {
                    // Donneur cellule j
                    int donor_idx = mesh.cell_idx(i, j);
                    double y_face_pos = mesh.yv(j);
                    double box_xmin = mesh.xv(i);
                    double box_xmax = mesh.xv(i + 1);
                    double box_ymin = y_face_pos;
                    double box_ymax = y_face_pos + delta_y;

                    const auto& seg = interface_segments[donor_idx];
                    flux_vol_y[mesh.cell_idx(i, j)] = -InterfaceReconstruction::cut_box_area(box_xmin, box_xmax, box_ymin, box_ymax,
                                                                                            seg.nx, seg.ny, seg.alpha);
                }
            }
        }

        for (int j = 1; j <= ny; ++j) {
            for (int i = 1; i <= nx; ++i) {
                int idx = mesh.cell_idx(i, j);
                double f_in  = flux_vol_y[mesh.cell_idx(i, j)];
                double f_out = flux_vol_y[mesh.cell_idx(i, j + 1)];
                double v_in  = fields.v[mesh.v_idx(i, j)];
                double v_out = fields.v[mesh.v_idx(i, j + 1)];

                double denom = dy - dt * (v_out - v_in);
                if (denom > 1.0e-12) {
                    C_out[idx] = (C_in[idx] * dy - (f_out - f_in) / dx) / denom;
                } else {
                    C_out[idx] = C_in[idx] - (f_out - f_in) / (dx * dy);
                }
                C_out[idx] = std::clamp(C_out[idx], 0.0, 1.0);
            }
        }
        MpiComm::sync_cell_scalar(mesh, C_out);
    };

    std::vector<double> C_star = fields.C;
    std::vector<double> C_final = fields.C;

    if (x_first) {
        // Balayage X puis Y
        sweep_x(fields.C, C_star);
        sweep_y(C_star, C_final);
    } else {
        // Balayage Y puis X
        sweep_y(fields.C, C_star);
        sweep_x(C_star, C_final);
    }

    fields.C = C_final;
}

void InterfaceTracker::advect_weymouth_yue(const Mesh2D& mesh, FlowFields& fields, double dt) {
    int nx = mesh.nx_local;
    int ny = mesh.ny_local;
    double dx = mesh.dx;
    double dy = mesh.dy;
    double cell_area = dx * dy;

    ReconstructionMethod recon_method = InterfaceReconstruction::parse_reconstruction_method(vof_cfg.reconstruction_method);
    bool x_first = (step_count % 2 == 0);

    // 1. Fonction indicatrice de dilatation binaire C_d(i,j) figee au temps n (Weymouth & Yue, JCP 2010, Eq. 14)
    std::vector<double> C_d((nx + 2) * (ny + 2), 0.0);
    for (int j = 0; j <= ny + 1; ++j) {
        for (int i = 0; i <= nx + 1; ++i) {
            int idx = mesh.cell_idx(i, j);
            C_d[idx] = (fields.C[idx] >= 0.5) ? 1.0 : 0.0;
        }
    }

    // 2. Balayage directionnel 1D selon X (Weymouth & Yue, JCP 2010, Eq. 15)
    auto sweep_x = [&](std::vector<double>& C_in, std::vector<double>& C_out, bool clamp_out) {
        InterfaceReconstruction::reconstruct(recon_method, mesh, C_in, interface_segments);
        std::vector<double> f_vol_x((nx + 2) * (ny + 2), 0.0);

        for (int j = 1; j <= ny; ++j) {
            for (int i = 1; i <= nx + 1; ++i) {
                double u_face = fields.u[mesh.u_idx(i, j)];
                double delta_x = std::abs(u_face) * dt;
                if (delta_x <= 1.0e-14) continue;

                double x_pos = mesh.xv(i);
                if (u_face > 0.0) {
                    int donor_idx = mesh.cell_idx(i - 1, j);
                    const auto& seg = interface_segments[donor_idx];
                    f_vol_x[mesh.cell_idx(i, j)] = InterfaceReconstruction::cut_box_area(
                        x_pos - delta_x, x_pos, mesh.yv(j), mesh.yv(j + 1), seg.nx, seg.ny, seg.alpha);
                } else {
                    int donor_idx = mesh.cell_idx(i, j);
                    const auto& seg = interface_segments[donor_idx];
                    f_vol_x[mesh.cell_idx(i, j)] = -InterfaceReconstruction::cut_box_area(
                        x_pos, x_pos + delta_x, mesh.yv(j), mesh.yv(j + 1), seg.nx, seg.ny, seg.alpha);
                }
            }
        }

        for (int j = 1; j <= ny; ++j) {
            for (int i = 1; i <= nx; ++i) {
                int idx = mesh.cell_idx(i, j);
                double fx_in  = f_vol_x[mesh.cell_idx(i, j)];
                double fx_out = f_vol_x[mesh.cell_idx(i + 1, j)];
                double u_in   = fields.u[mesh.u_idx(i, j)];
                double u_out  = fields.u[mesh.u_idx(i + 1, j)];

                C_out[idx] = C_in[idx] - (fx_out - fx_in) / cell_area + C_d[idx] * (dt / dx) * (u_out - u_in);
                if (clamp_out) {
                    C_out[idx] = std::clamp(C_out[idx], 0.0, 1.0);
                }
            }
        }
        MpiComm::sync_cell_scalar(mesh, C_out);
    };

    // 3. Balayage directionnel 1D selon Y (Weymouth & Yue, JCP 2010, Eq. 16)
    auto sweep_y = [&](std::vector<double>& C_in, std::vector<double>& C_out, bool clamp_out) {
        InterfaceReconstruction::reconstruct(recon_method, mesh, C_in, interface_segments);
        std::vector<double> f_vol_y((nx + 2) * (ny + 2), 0.0);

        for (int j = 1; j <= ny + 1; ++j) {
            for (int i = 1; i <= nx; ++i) {
                double v_face = fields.v[mesh.v_idx(i, j)];
                double delta_y = std::abs(v_face) * dt;
                if (delta_y <= 1.0e-14) continue;

                double y_pos = mesh.yv(j);
                if (v_face > 0.0) {
                    int donor_idx = mesh.cell_idx(i, j - 1);
                    const auto& seg = interface_segments[donor_idx];
                    f_vol_y[mesh.cell_idx(i, j)] = InterfaceReconstruction::cut_box_area(
                        mesh.xv(i), mesh.xv(i + 1), y_pos - delta_y, y_pos, seg.nx, seg.ny, seg.alpha);
                } else {
                    int donor_idx = mesh.cell_idx(i, j);
                    const auto& seg = interface_segments[donor_idx];
                    f_vol_y[mesh.cell_idx(i, j)] = -InterfaceReconstruction::cut_box_area(
                        mesh.xv(i), mesh.xv(i + 1), y_pos, y_pos + delta_y, seg.nx, seg.ny, seg.alpha);
                }
            }
        }

        for (int j = 1; j <= ny; ++j) {
            for (int i = 1; i <= nx; ++i) {
                int idx = mesh.cell_idx(i, j);
                double fy_in  = f_vol_y[mesh.cell_idx(i, j)];
                double fy_out = f_vol_y[mesh.cell_idx(i, j + 1)];
                double v_in   = fields.v[mesh.v_idx(i, j)];
                double v_out  = fields.v[mesh.v_idx(i, j + 1)];

                C_out[idx] = C_in[idx] - (fy_out - fy_in) / cell_area + C_d[idx] * (dt / dy) * (v_out - v_in);
                if (clamp_out) {
                    C_out[idx] = std::clamp(C_out[idx], 0.0, 1.0);
                }
            }
        }
        MpiComm::sync_cell_scalar(mesh, C_out);
    };

    std::vector<double> C_star = fields.C;
    std::vector<double> C_final = fields.C;

    if (x_first) {
        sweep_x(fields.C, C_star, false);
        sweep_y(C_star, C_final, true);
    } else {
        sweep_y(fields.C, C_star, false);
        sweep_x(C_star, C_final, true);
    }

    fields.C = C_final;
}

} // namespace covof

\end{lstlisting}

\chapter{Exportation Graphique VTK XML Parallèle (.vtr / .pvtr) (VTKWriter.hpp, VTKWriter.cpp)}
\section{Architecture des Fichiers VTK XML \texttt{RectilinearGrid} et \texttt{PRectilinearGrid}}
La classe \texttt{covof::VTKWriter} (\texttt{include/covof/VTKWriter.hpp}, \texttt{src/VTKWriter.cpp}) produit des fichiers directement exploitables par \textbf{ParaView}, \textbf{VisIt} et les scripts Python de post-traitement sans bibliothèque externe :
\begin{enumerate}
  \item Chaque processus MPI de rang $r$ écrit son sous-domaine local dans \texttt{<case>\_p<r>\_<step>.vtr} (format \texttt{RectilinearGrid}) contenant les coordonnées des nœuds $(x_u, y_v, 0)$ et les tableaux \texttt{CellData} :
  \begin{itemize}
    \item \texttt{Pressure} ($p_{i,j}$) ;
    \item \texttt{Temperature} ($T_{i,j}$) ;
    \item \texttt{PhaseIndicator} ($C_{i,j}$) ;
    \item \texttt{Divergence} ($(\nabla_h \cdot \mathbf{u})_{i,j}$) ;
    \item \texttt{Vorticity} ($\omega_z(i,j) = \frac{v_{c}(i+1,j)-v_{c}(i-1,j)}{2\Delta x} - \frac{u_{c}(i,j+1)-u_{c}(i,j-1)}{2\Delta y}$) ;
    \item \texttt{Viscosity} ($\mu_{i,j}$) ;
    \item \texttt{Velocity} (vecteur 3D $(u_c(i,j), v_c(i,j), 0)$ interpolé aux centres des mailles : $u_c(i,j) = \frac{1}{2}(u_{i,j} + u_{i+1,j})$, $v_c(i,j) = \frac{1}{2}(v_{i,j} + v_{i,j+1})$).
  \end{itemize}
  \item Le processus maître ($r = 0$) écrit le fichier chapeau \texttt{<case>\_<step>.pvtr} (\texttt{PRectilinearGrid}) référençant les extensions globales \texttt{Piece Extent} de tous les sous-domaines MPI.
\end{enumerate}

\section{Répertoire des Classes, Structures et Méthodes de \texttt{include/covof/VTKWriter.hpp}}
\begin{center}\small
\begin{tabular}{clll}
\toprule
\textbf{Ligne} & \textbf{Catégorie} & \textbf{Identifiant C++17} & \textbf{Signature / Rôle} \\
\midrule
14 & class & \texttt{VTKWriter} & Déclaration de type / classe / structure \\
28 & Méthode / Fonction & \texttt{init} & Type retour : void \\
33 & Méthode / Fonction & \texttt{write\_step} & Type retour : void \\
38 & Méthode / Fonction & \texttt{write\_diagnostics} & Type retour : void \\
40 & Méthode / Fonction & \texttt{finalize} & Type retour : void \\
\bottomrule
\end{tabular}
\end{center}

\section{Code Source Intégral Annoté : \texttt{include/covof/VTKWriter.hpp}}
Le listing ci-dessous présente l'intégralité des 46 lignes du fichier \texttt{include/covof/VTKWriter.hpp} de l'implémentation C++17 :

\begin{lstlisting}[language={[11]C++}]
#pragma once

#include "Mesh.hpp"
#include "Field.hpp"
#include <string>
#include <fstream>

namespace covof {
namespace fugu = covof; // Alias rétrocompatibilité

/**
 * @brief Gestionnaire d'exportation pour ParaView (XML VTK .vtr/.pvtr) et diagnostics CSV
 */
class VTKWriter {
public:
    // Déclaration de variable typée
    std::string output_dir = "output";
    // Déclaration de variable typée
    std::string case_name  = "covof";
    std::ofstream csv_file;

    // Calcul et assignation numérique
    VTKWriter() = default;

    /**
     * @brief Initialise le répertoire de sortie et le fichier CSV d'historique
     */
    void init(const Mesh2D& mesh, const std::string& out_dir, const std::string& name);

    /**
     * @brief Exporte le pas de temps courant au format VTK XML (.vtr par processus + .pvtr maître)
     */
    void write_step(const Mesh2D& mesh, const FlowFields& fields, int step, double time);

    /**
     * @brief Calcule et enregistre les grandeurs physiques intégrales dans le fichier CSV
     */
    void write_diagnostics(const Mesh2D& mesh, const FlowFields& fields, int step, double time, double dt);

    void finalize() {
        // Évaluation de la condition de contrôle
        if (csv_file.is_open()) csv_file.close();
    }
};

} // namespace covof

\end{lstlisting}

\section{Répertoire des Classes, Structures et Méthodes de \texttt{src/VTKWriter.cpp}}
\begin{center}\small
\begin{tabular}{clll}
\toprule
\textbf{Ligne} & \textbf{Catégorie} & \textbf{Identifiant C++17} & \textbf{Signature / Rôle} \\
\midrule
11 & Méthode / Fonction & \texttt{VTKWriter::init} & Type retour : void \\
31 & Méthode / Fonction & \texttt{VTKWriter::write\_step} & Type retour : void \\
33 & Méthode / Fonction & \texttt{std::setfill} & Type retour : ss\_step $<$$<$ \\
49 & Méthode / Fonction & \texttt{vtr} & Type retour : std::ofstream \\
186 & Méthode / Fonction & \texttt{all\_x\_start} & Type retour : std::vector$<$int$>$ \\
188 & Méthode / Fonction & \texttt{all\_y\_start} & Type retour : std::vector$<$int$>$ \\
201 & Méthode / Fonction & \texttt{pvtr} & Type retour : std::ofstream \\
260 & Méthode / Fonction & \texttt{VTKWriter::write\_diagnostics} & Type retour : void \\
\bottomrule
\end{tabular}
\end{center}

\section{Code Source Intégral Annoté : \texttt{src/VTKWriter.cpp}}
Le listing ci-dessous présente l'intégralité des 346 lignes du fichier \texttt{src/VTKWriter.cpp} de l'implémentation C++17 :

\begin{lstlisting}[language={[11]C++}]
#include "covof/VTKWriter.hpp"
#include "covof/MpiComm.hpp"
#include <iomanip>
#include <sstream>
#include <sys/stat.h>
#include <cmath>

namespace covof {
namespace fugu = covof; // Alias rétrocompatibilité

void VTKWriter::init(const Mesh2D& mesh, const std::string& out_dir, const std::string& name) {
    // Calcul et assignation numérique
    output_dir = out_dir;
    // Calcul et assignation numérique
    case_name = name;

    // Évaluation de la condition de contrôle
    if (mesh.rank == 0) {
        mkdir(output_dir.c_str(), 0777);

        // Déclaration de variable typée
        std::string csv_path = output_dir + "/" + case_name + "_history.csv";
        csv_file.open(csv_path);
        // Évaluation de la condition de contrôle
        if (csv_file.is_open()) {
            csv_file << "step,time,dt,center_of_mass_y,rise_velocity,bubble_volume,kinetic_energy,max_u,max_div\n";
        }
    }
}

void VTKWriter::write_step(const Mesh2D& mesh, const FlowFields& fields, int step, double time) {
    std::ostringstream ss_step;
    ss_step << std::setfill('0') << std::setw(6) << step;
    // Déclaration de variable typée
    std::string step_str = ss_step.str();

    // Déclaration de variable typée
    int x_start = mesh.i_offset;
    // Déclaration de variable typée
    int x_end   = mesh.i_offset + mesh.nx_local;
    // Déclaration de variable typée
    int y_start = mesh.j_offset;
    // Déclaration de variable typée
    int y_end   = mesh.j_offset + mesh.ny_local;

    // 1. Écriture du fichier .vtr propre à chaque processus MPI
    std::ostringstream ss_vtr;
    ss_vtr << output_dir << "/" << case_name << "_p" << mesh.rank << "_" << step_str << ".vtr";
    std::ofstream vtr(ss_vtr.str());

    // Évaluation de la condition de contrôle
    if (vtr.is_open()) {
        // Calcul et assignation numérique
        vtr << "<?xml version=\"1.0\"?>\n"
            // Calcul et assignation numérique
            << "<VTKFile type=\"RectilinearGrid\" version=\"0.1\" byte_order=\"LittleEndian\">\n"
            // Calcul et assignation numérique
            << "  <RectilinearGrid WholeExtent=\"" << x_start << " " << x_end << " "
            << y_start << " " << y_end << " 0 0\">\n"
            // Calcul et assignation numérique
            << "    <Piece Extent=\"" << x_start << " " << x_end << " "
            << y_start << " " << y_end << " 0 0\">\n"
            // Calcul et assignation numérique
            << "      <CellData Scalars=\"PhaseIndicator,Pressure,Density,Viscosity,Divergence\" Vectors=\"Velocity\">\n";

        // PhaseIndicator C
        vtr << "        <DataArray type=\"Float64\" Name=\"PhaseIndicator\" format=\"ascii\">\n";
        // Boucle spatiale sur les cellules locales selon l'axe vertical Y (ordonnées)
        for (int j = 1; j <= mesh.ny_local; ++j) {
            // Boucle spatiale sur les cellules locales selon l'axe horizontal X (abscisses)
            for (int i = 1; i <= mesh.nx_local; ++i) {
                vtr << fields.C[mesh.cell_idx(i, j)] << " ";
            }
            vtr << "\n";
        }
        vtr << "        </DataArray>\n";

        // Pressure p
        vtr << "        <DataArray type=\"Float64\" Name=\"Pressure\" format=\"ascii\">\n";
        // Boucle spatiale sur les cellules locales selon l'axe vertical Y (ordonnées)
        for (int j = 1; j <= mesh.ny_local; ++j) {
            // Boucle spatiale sur les cellules locales selon l'axe horizontal X (abscisses)
            for (int i = 1; i <= mesh.nx_local; ++i) {
                vtr << fields.p[mesh.cell_idx(i, j)] << " ";
            }
            vtr << "\n";
        }
        vtr << "        </DataArray>\n";

        // Density rho
        vtr << "        <DataArray type=\"Float64\" Name=\"Density\" format=\"ascii\">\n";
        // Boucle spatiale sur les cellules locales selon l'axe vertical Y (ordonnées)
        for (int j = 1; j <= mesh.ny_local; ++j) {
            // Boucle spatiale sur les cellules locales selon l'axe horizontal X (abscisses)
            for (int i = 1; i <= mesh.nx_local; ++i) {
                vtr << fields.rho[mesh.cell_idx(i, j)] << " ";
            }
            vtr << "\n";
        }
        vtr << "        </DataArray>\n";

        // Viscosity mu
        vtr << "        <DataArray type=\"Float64\" Name=\"Viscosity\" format=\"ascii\">\n";
        // Boucle spatiale sur les cellules locales selon l'axe vertical Y (ordonnées)
        for (int j = 1; j <= mesh.ny_local; ++j) {
            // Boucle spatiale sur les cellules locales selon l'axe horizontal X (abscisses)
            for (int i = 1; i <= mesh.nx_local; ++i) {
                vtr << fields.mu[mesh.cell_idx(i, j)] << " ";
            }
            vtr << "\n";
        }
        vtr << "        </DataArray>\n";

        // Divergence div(u)
        vtr << "        <DataArray type=\"Float64\" Name=\"Divergence\" format=\"ascii\">\n";
        // Boucle spatiale sur les cellules locales selon l'axe vertical Y (ordonnées)
        for (int j = 1; j <= mesh.ny_local; ++j) {
            // Boucle spatiale sur les cellules locales selon l'axe horizontal X (abscisses)
            for (int i = 1; i <= mesh.nx_local; ++i) {
                // Déclaration de variable typée
                double div = (fields.u[mesh.u_idx(i + 1, j)] - fields.u[mesh.u_idx(i, j)]) / mesh.dx +
                             (fields.v[mesh.v_idx(i, j + 1)] - fields.v[mesh.v_idx(i, j)]) / mesh.dy;
                vtr << div << " ";
            }
            vtr << "\n";
        }
        vtr << "        </DataArray>\n";

        // Temperature T
        vtr << "        <DataArray type=\"Float64\" Name=\"Temperature\" format=\"ascii\">\n";
        // Boucle spatiale sur les cellules locales selon l'axe vertical Y (ordonnées)
        for (int j = 1; j <= mesh.ny_local; ++j) {
            // Boucle spatiale sur les cellules locales selon l'axe horizontal X (abscisses)
            for (int i = 1; i <= mesh.nx_local; ++i) {
                vtr << fields.T[mesh.cell_idx(i, j)] << " ";
            }
            vtr << "\n";
        }
        vtr << "        </DataArray>\n";

        // Velocity vector au centre des cellules
        vtr << "        <DataArray type=\"Float64\" Name=\"Velocity\" NumberOfComponents=\"3\" format=\"ascii\">\n";
        // Boucle spatiale sur les cellules locales selon l'axe vertical Y (ordonnées)
        for (int j = 1; j <= mesh.ny_local; ++j) {
            // Boucle spatiale sur les cellules locales selon l'axe horizontal X (abscisses)
            for (int i = 1; i <= mesh.nx_local; ++i) {
                // Déclaration de variable typée
                double uc = 0.5 * (fields.u[mesh.u_idx(i, j)] + fields.u[mesh.u_idx(i + 1, j)]);
                // Déclaration de variable typée
                double vc = 0.5 * (fields.v[mesh.v_idx(i, j)] + fields.v[mesh.v_idx(i, j + 1)]);
                vtr << uc << " " << vc << " 0.0 ";
            }
            vtr << "\n";
        }
        vtr << "        </DataArray>\n";

        vtr << "      </CellData>\n"
            << "      <Coordinates>\n"
            // Calcul et assignation numérique
            << "        <DataArray type=\"Float64\" Name=\"XCoordinates\" format=\"ascii\">\n";
        // Boucle spatiale sur les cellules locales selon l'axe horizontal X (abscisses)
        for (int i = 0; i <= mesh.nx_local; ++i) {
            vtr << (mesh.i_offset + i) * mesh.dx << " ";
        }
        vtr << "\n        </DataArray>\n"
            // Calcul et assignation numérique
            << "        <DataArray type=\"Float64\" Name=\"YCoordinates\" format=\"ascii\">\n";
        // Boucle spatiale sur les cellules locales selon l'axe vertical Y (ordonnées)
        for (int j = 0; j <= mesh.ny_local; ++j) {
            vtr << (mesh.j_offset + j) * mesh.dy << " ";
        }
        vtr << "\n        </DataArray>\n"
            // Calcul et assignation numérique
            << "        <DataArray type=\"Float64\" Name=\"ZCoordinates\" format=\"ascii\">\n"
            << "          0.0\n"
            << "        </DataArray>\n"
            << "      </Coordinates>\n"
            << "    </Piece>\n"
            << "  </RectilinearGrid>\n"
            << "</VTKFile>\n";
    }

    // 2. Écriture du fichier maître .pvtr par le processus racine (rank 0)
    if (mesh.rank == 0) {
        // Déclaration de variable typée
        std::vector<int> all_x_start(mesh.num_procs), all_x_end(mesh.num_procs);
        // Déclaration de variable typée
        std::vector<int> all_y_start(mesh.num_procs), all_y_end(mesh.num_procs);

        // Appel de communication ou synchronisation parallèle OpenMPI
        MPI_Gather(&x_start, 1, MPI_INT, all_x_start.data(), 1, MPI_INT, 0, mesh.comm);
        // Appel de communication ou synchronisation parallèle OpenMPI
        MPI_Gather(&x_end,   1, MPI_INT, all_x_end.data(),   1, MPI_INT, 0, mesh.comm);
        // Appel de communication ou synchronisation parallèle OpenMPI
        MPI_Gather(&y_start, 1, MPI_INT, all_y_start.data(), 1, MPI_INT, 0, mesh.comm);
        // Appel de communication ou synchronisation parallèle OpenMPI
        MPI_Gather(&y_end,   1, MPI_INT, all_y_end.data(),   1, MPI_INT, 0, mesh.comm);

        std::ostringstream ss_pvtr;
        ss_pvtr << output_dir << "/" << case_name << "_" << step_str << ".pvtr";
        std::ofstream pvtr(ss_pvtr.str());

        // Évaluation de la condition de contrôle
        if (pvtr.is_open()) {
            // Calcul et assignation numérique
            pvtr << "<?xml version=\"1.0\"?>\n"
                 // Calcul et assignation numérique
                 << "<VTKFile type=\"PRectilinearGrid\" version=\"0.1\" byte_order=\"LittleEndian\">\n"
                 // Calcul et assignation numérique
                 << "  <PRectilinearGrid WholeExtent=\"0 " << mesh.nx_global << " 0 " << mesh.ny_global << " 0 0\" GhostLevel=\"0\">\n"
                 // Calcul et assignation numérique
                 << "    <PCellData Scalars=\"PhaseIndicator,Pressure,Density,Viscosity,Divergence\" Vectors=\"Velocity\">\n"
                 // Calcul et assignation numérique
                 << "      <PDataArray type=\"Float64\" Name=\"PhaseIndicator\"/>\n"
                 // Calcul et assignation numérique
                 << "      <PDataArray type=\"Float64\" Name=\"Pressure\"/>\n"
                 // Calcul et assignation numérique
                 << "      <PDataArray type=\"Float64\" Name=\"Density\"/>\n"
                 // Calcul et assignation numérique
                 << "      <PDataArray type=\"Float64\" Name=\"Viscosity\"/>\n"
                 // Calcul et assignation numérique
                 << "      <PDataArray type=\"Float64\" Name=\"Divergence\"/>\n"
                 // Calcul et assignation numérique
                 << "      <PDataArray type=\"Float64\" Name=\"Velocity\" NumberOfComponents=\"3\"/>\n"
                 << "    </PCellData>\n"
                 << "    <PCoordinates>\n"
                 // Calcul et assignation numérique
                 << "      <PDataArray type=\"Float64\" Name=\"XCoordinates\"/>\n"
                 // Calcul et assignation numérique
                 << "      <PDataArray type=\"Float64\" Name=\"YCoordinates\"/>\n"
                 // Calcul et assignation numérique
                 << "      <PDataArray type=\"Float64\" Name=\"ZCoordinates\"/>\n"
                 << "    </PCoordinates>\n";

            // Boucle d'itération algorithmique
            for (int p = 0; p < mesh.num_procs; ++p) {
                std::ostringstream piece_vtr;
                piece_vtr << case_name << "_p" << p << "_" << step_str << ".vtr";
                // Calcul et assignation numérique
                pvtr << "    <Piece Extent=\"" << all_x_start[p] << " " << all_x_end[p] << " "
                     // Calcul et assignation numérique
                     << all_y_start[p] << " " << all_y_end[p] << " 0 0\" Source=\"" << piece_vtr.str() << "\"/>\n";
            }

            pvtr << "  </PRectilinearGrid>\n"
                 << "</VTKFile>\n";
        }
    } else {
        // Appel de communication ou synchronisation parallèle OpenMPI
        MPI_Gather(&x_start, 1, MPI_INT, nullptr, 0, MPI_INT, 0, mesh.comm);
        // Appel de communication ou synchronisation parallèle OpenMPI
        MPI_Gather(&x_end,   1, MPI_INT, nullptr, 0, MPI_INT, 0, mesh.comm);
        // Appel de communication ou synchronisation parallèle OpenMPI
        MPI_Gather(&y_start, 1, MPI_INT, nullptr, 0, MPI_INT, 0, mesh.comm);
        // Appel de communication ou synchronisation parallèle OpenMPI
        MPI_Gather(&y_end,   1, MPI_INT, nullptr, 0, MPI_INT, 0, mesh.comm);
    }
}

void VTKWriter::write_diagnostics(const Mesh2D& mesh, const FlowFields& fields, int step, double time, double dt) {
    // Déclaration de variable typée
    double local_y_integral = 0.0;
    // Déclaration de variable typée
    double local_v_integral = 0.0;
    // Déclaration de variable typée
    double local_vol = 0.0;
    // Déclaration de variable typée
    double local_ke = 0.0;
    // Déclaration de variable typée
    double local_max_u = 0.0;
    // Déclaration de variable typée
    double local_max_div = 0.0;

    // Déclaration de variable typée
    double dvol = mesh.dx * mesh.dy;

    // Boucle spatiale sur les cellules locales selon l'axe vertical Y (ordonnées)
    for (int j = 1; j <= mesh.ny_local; ++j) {
        // Déclaration de variable typée
        double y = mesh.yc(j);
        // Boucle spatiale sur les cellules locales selon l'axe horizontal X (abscisses)
        for (int i = 1; i <= mesh.nx_local; ++i) {
            // Déclaration de variable typée
            double c = fields.C[mesh.cell_idx(i, j)];
            // Déclaration de variable typée
            double indicator_bubble = 1.0 - c; // 1 si phase fluide 2 (bulle/goutte)

            // Calcul et assignation numérique
            local_vol += indicator_bubble * dvol;
            // Calcul et assignation numérique
            local_y_integral += y * indicator_bubble * dvol;

            // Déclaration de variable typée
            double uc = 0.5 * (fields.u[mesh.u_idx(i, j)] + fields.u[mesh.u_idx(i + 1, j)]);
            // Déclaration de variable typée
            double vc = 0.5 * (fields.v[mesh.v_idx(i, j)] + fields.v[mesh.v_idx(i, j + 1)]);

            // Calcul et assignation numérique
            local_v_integral += vc * indicator_bubble * dvol;

            // Déclaration de variable typée
            double speed2 = uc * uc + vc * vc;
            // Calcul et assignation numérique
            local_ke += 0.5 * fields.rho[mesh.cell_idx(i, j)] * speed2 * dvol;

            // Déclaration de variable typée
            double speed = std::sqrt(speed2);
            // Évaluation de la condition de contrôle
            if (speed > local_max_u) local_max_u = speed;

            // Déclaration de variable typée
            double div = std::abs((fields.u[mesh.u_idx(i + 1, j)] - fields.u[mesh.u_idx(i, j)]) / mesh.dx +
                                  (fields.v[mesh.v_idx(i, j + 1)] - fields.v[mesh.v_idx(i, j)]) / mesh.dy);
            // Évaluation de la condition de contrôle
            if (div > local_max_div) local_max_div = div;
        }
    }

    // Déclaration de variable typée
    double global_vol = MpiComm::global_sum(mesh, local_vol);
    // Déclaration de variable typée
    double global_y_integral = MpiComm::global_sum(mesh, local_y_integral);
    // Déclaration de variable typée
    double global_v_integral = MpiComm::global_sum(mesh, local_v_integral);
    // Déclaration de variable typée
    double global_ke = MpiComm::global_sum(mesh, local_ke);
    // Déclaration de variable typée
    double global_max_u = MpiComm::global_max(mesh, local_max_u);
    // Déclaration de variable typée
    double global_max_div = MpiComm::global_max(mesh, local_max_div);

    // Déclaration de variable typée
    double yc = (global_vol > 1.0e-12) ? global_y_integral / global_vol : 0.0;
    // Déclaration de variable typée
    double vc = (global_vol > 1.0e-12) ? global_v_integral / global_vol : 0.0;

    // Évaluation de la condition de contrôle
    if (mesh.rank == 0 && csv_file.is_open()) {
        csv_file << step << "," << time << "," << dt << ","
                 << yc << "," << vc << "," << global_vol << ","
                 << global_ke << "," << global_max_u << "," << global_max_div << "\n";
        csv_file.flush();
    }
}

} // namespace covof

\end{lstlisting}

\chapter{Orchestration de la Simulation, Pas de Temps Adaptatif et Diagnostics V&V (Simulation.hpp, Simulation.cpp)}
\section{Orchestration Globale dans \texttt{covof::Simulation}}
La classe \texttt{covof::Simulation} (\texttt{include/covof/Simulation.hpp}, \texttt{src/Simulation.cpp}, 916 lignes) coordonne l'ensemble des composants du solveur :
\begin{enumerate}
  \item \textbf{Initialisation géométrique et physique (\texttt{setup\_initial\_conditions})} : génère les profils analytiques, disques, colonnes de barrage, nappes liquides cisaillement Kelvin-Helmholtz, perturbations cosinusoïdales de Rayleigh-Taylor et buses d'injection de jets visqueux avec sous-échantillonnage $8 \times 8$ par cellule pour initialiser une fraction volumique $C_{i,j}(0)$ précise au second ordre.
  \item \textbf{Contrôle du pas de temps adaptatif (\texttt{compute\_adaptive\_dt})} : évalue à chaque pas de temps la contrainte CFL multi-physique (convective, capillaire de Brackbill et VOF) :
  \begin{equation}
    \Delta t \le \min\left( \text{CFL}_{\max} \frac{\min(\Delta x, \Delta y)}{\max(\|\mathbf{u}\|_\infty, 10^{-6})}, \quad \sqrt{\frac{(\rho_1 + \rho_2)\min(\Delta x, \Delta y)^3}{4\pi \sigma}} \right)
  \end{equation}
  \item \textbf{Calcul des diagnostics intégraux et normes V\&V} : produit automatiquement les fichiers CSV de suivi temporel (\texttt{history.csv}, \texttt{dense\_cylinder\_history.csv}, \texttt{liquid\_sheet\_history.csv}, \texttt{viscous\_jet\_history.csv}, \texttt{vortex\_error.csv}, \texttt{laplacian\_summary.csv}).
\end{enumerate}

\section{Répertoire des Classes, Structures et Méthodes de \texttt{include/covof/Simulation.hpp}}
\begin{center}\small
\begin{tabular}{clll}
\toprule
\textbf{Ligne} & \textbf{Catégorie} & \textbf{Identifiant C++17} & \textbf{Signature / Rôle} \\
\midrule
18 & class & \texttt{Simulation} & Déclaration de type / classe / structure \\
35 & Méthode / Fonction & \texttt{init} & Type retour : bool \\
40 & Méthode / Fonction & \texttt{run} & Type retour : void \\
45 & Méthode / Fonction & \texttt{run\_laplacian} & Type retour : void \\
50 & Méthode / Fonction & \texttt{run\_vortex\_advection} & Type retour : void \\
55 & Méthode / Fonction & \texttt{compute\_time\_step} & Type retour : double \\
\bottomrule
\end{tabular}
\end{center}

\section{Code Source Intégral Annoté : \texttt{include/covof/Simulation.hpp}}
Le listing ci-dessous présente l'intégralité des 58 lignes du fichier \texttt{include/covof/Simulation.hpp} de l'implémentation C++17 :

\begin{lstlisting}[language={[11]C++}]
#pragma once

#include "Config.hpp"
#include "Mesh.hpp"
#include "Field.hpp"
#include "InterfaceTracker.hpp"
#include "NavierStokesSolver.hpp"
#include "Preconditioner.hpp"
#include "LinearSolver.hpp"
#include "VTKWriter.hpp"

namespace covof {
namespace fugu = covof; // Alias rétrocompatibilité

/**
 * @brief Orchestrateur principal de la simulation CFD diphasique tout-couplée 2D
 */
class Simulation {
public:
    Config config;
    Mesh2D mesh;
    FlowFields fields;
    InterfaceTracker tracker;
    NavierStokesSolver ns_solver;
    std::unique_ptr<PreconditionerBase> precond;
    std::unique_ptr<LinearSolverBase> linear_solver;
    VTKWriter vtk_writer;

    // Calcul et assignation numérique
    Simulation() = default;

    /**
     * @brief Configure et initialise l'ensemble des modules physiques et numériques
     */
    bool init(MPI_Comm comm, const std::string& config_file);

    /**
     * @brief Exécute la boucle temporelle jusqu'à t_final
     */
    void run();

    /**
     * @brief Résout le problème de vérification du Laplacien analytique (Cas 3.1 Référence Académique)
     */
    void run_laplacian();

    /**
     * @brief Résout le problème d'advection d'une tâche dans un vortex avec cisaillement (Cas 3.19 Référence Académique)
     */
    void run_vortex_advection();

    /**
     * @brief Calcule le pas de temps adaptatif selon les critères CFL, viscosité et tension superficielle
     */
    double compute_time_step() const;
};

} // namespace covof

\end{lstlisting}

\section{Répertoire des Classes, Structures et Méthodes de \texttt{src/Simulation.cpp}}
\begin{center}\small
\begin{tabular}{clll}
\toprule
\textbf{Ligne} & \textbf{Catégorie} & \textbf{Identifiant C++17} & \textbf{Signature / Rôle} \\
\midrule
13 & Méthode / Fonction & \texttt{Simulation::init} & Type retour : bool \\
218 & Méthode / Fonction & \texttt{Simulation::compute\_time\_step} & Type retour : double \\
265 & Méthode / Fonction & \texttt{Simulation::run} & Type retour : void \\
291 & Méthode / Fonction & \texttt{std::setw} & Type retour : $<$$<$ std::left $<$$<$ \\
292 & Méthode / Fonction & \texttt{std::setw} & Type retour : $<$$<$ \\
293 & Méthode / Fonction & \texttt{std::setw} & Type retour : $<$$<$ \\
294 & Méthode / Fonction & \texttt{std::setw} & Type retour : $<$$<$ \\
295 & Méthode / Fonction & \texttt{std::setw} & Type retour : $<$$<$ \\
296 & Méthode / Fonction & \texttt{std::setw} & Type retour : $<$$<$ \\
299 & Méthode / Fonction & \texttt{std::setw} & Type retour : std::cout $<$$<$ \\
305 & Méthode / Fonction & \texttt{state\_vector} & Type retour : std::vector$<$double$>$ \\
307 & Méthode / Fonction & \texttt{rhs} & Type retour : std::vector$<$double$>$ \\
436 & Méthode / Fonction & \texttt{std::setw} & Type retour : std::cout $<$$<$ std::left $<$$<$ \\
437 & Méthode / Fonction & \texttt{std::setw} & Type retour : $<$$<$ \\
438 & Méthode / Fonction & \texttt{std::setw} & Type retour : $<$$<$ \\
439 & Méthode / Fonction & \texttt{std::setw} & Type retour : $<$$<$ \\
440 & Méthode / Fonction & \texttt{std::setw} & Type retour : $<$$<$ \\
441 & Méthode / Fonction & \texttt{std::setw} & Type retour : $<$$<$ \\
444 & Méthode / Fonction & \texttt{std::setw} & Type retour : std::cout $<$$<$ \\
464 & Méthode / Fonction & \texttt{std::string} & Type retour : std::cout $<$$<$ \\
471 & Méthode / Fonction & \texttt{Simulation::run\_laplacian} & Type retour : void \\
538 & Méthode / Fonction & \texttt{rhs} & Type retour : std::vector$<$double$>$ \\
593 & Méthode / Fonction & \texttt{diag} & Type retour : std::vector$<$double$>$ \\
604 & Méthode / Fonction & \texttt{p\_cg} & Type retour : std::vector$<$double$>$ \\
717 & Méthode / Fonction & \texttt{ofs} & Type retour : std::ofstream \\
728 & Méthode / Fonction & \texttt{Simulation::run\_vortex\_advection} & Type retour : void \\
906 & Méthode / Fonction & \texttt{ofs} & Type retour : std::ofstream \\
\bottomrule
\end{tabular}
\end{center}

\section{Code Source Intégral Annoté : \texttt{src/Simulation.cpp}}
Le listing ci-dessous présente l'intégralité des 924 lignes du fichier \texttt{src/Simulation.cpp} de l'implémentation C++17 :

\begin{lstlisting}[language={[11]C++}]
#include "covof/Simulation.hpp"
#include "covof/MpiComm.hpp"
#include <iostream>
#include <iomanip>
#include <fstream>
#include <algorithm>
#include <cmath>

namespace covof {
namespace fugu = covof; // Alias rétrocompatibilité

// Déclaration de variable typée
bool Simulation::init(MPI_Comm comm, const std::string& config_file) {
    // 1. Chargement de la configuration JSON
    if (!config.load_from_json(config_file)) {
        // Retour de fonction
        return false;
    }

    // Déclaration de variable typée
    int my_rank;
    // Appel de communication ou synchronisation parallèle OpenMPI
    MPI_Comm_rank(comm, &my_rank);
    // Évaluation de la condition de contrôle
    if (my_rank == 0) {
        config.print();
    }

    // 2. Initialisation du maillage partitionné
    bool periodic_x = (config.bc.left.type == BCType::PERIODIC || config.bc.right.type == BCType::PERIODIC);
    // Déclaration de variable typée
    bool periodic_y = (config.bc.bottom.type == BCType::PERIODIC || config.bc.top.type == BCType::PERIODIC);
    mesh.init(comm, config.nx, config.ny, config.lx, config.ly, periodic_x, periodic_y);

    // 3. Allocation des champs décalés MAC avec halos
    fields.allocate(mesh);

    // 4. Initialisation du module d'interface et des conditions initiales
    tracker.init(config);
    tracker.set_initial_condition(mesh, config.init_cond, fields);

    // Initialisation des conditions aux limites initiales
    MpiComm::apply_boundary_conditions(mesh, config.bc, fields, false);

    // Initialisation du canal amont pour profil POISEUILLE_STEP afin de garantir la continuité
    if (config.bc.left.type == BCType::DIRICHLET && config.bc.left.profile == "POISEUILLE_STEP") {
        // Déclaration de variable typée
        double y_step = config.bc.left.y_step;
        // Déclaration de variable typée
        double h = mesh.ly - y_step;
        // Déclaration de variable typée
        double u_mean = config.bc.left.u_mean;
        // Boucle spatiale sur les cellules locales selon l'axe vertical Y (ordonnées)
        for (int j = 1; j <= mesh.ny_local; ++j) {
            // Déclaration de variable typée
            double y = mesh.yc(j);
            // Déclaration de variable typée
            double u_val = 0.0;
            // Évaluation de la condition de contrôle
            if (y >= y_step) {
                // Calcul et assignation numérique
                u_val = 6.0 * u_mean * (y - y_step) * (mesh.ly - y) / (h * h);
            }
            // Déclaration de variable typée
            int u_nx = mesh.nx_local + (mesh.is_boundary_right ? 1 : 0);
            // Boucle spatiale sur les cellules locales selon l'axe horizontal X (abscisses)
            for (int i = 1; i <= u_nx; ++i) {
                // Déclaration de variable typée
                double x = mesh.xu(i);
                // Évaluation de la condition de contrôle
                if (x <= 4.0 + 1e-9) {
                    // Calcul de la résolution géométrique de cellule
                    fields.u[mesh.u_idx(i, j)] = u_val;
                }
            }
        }
        MpiComm::sync_u(mesh, fields.u);
    }

    // Initialisation du champ de vitesse solénoïdal pour LIQUID_SHEET (Équations 17-18 Fluids 2021)
    if (config.init_cond.type == InitialConditionType::LIQUID_SHEET) {
        // Déclaration des variables physiques et géométriques
        double y_c = (config.init_cond.center_y > 0.0) ? config.init_cond.center_y : 0.5 * mesh.ly;
        // Déclaration de variable typée pour l'épaisseur de la nappe liquide
        double delta = (config.init_cond.thickness > 0.0) ? config.init_cond.thickness : 0.0003;
        // Déclaration de variable typée pour la longueur du domaine
        double L = mesh.lx;
        // Déclaration de variable typée pour le nombre d'onde
        double k_pert = 2.0 * M_PI / L;
        // Déclaration de variable typée pour l'amplitude de perturbation
        double amp = 0.04;
        // Facteur pour u assurant une divergence rigoureusement nulle (Eq. 17)
        double factor_u = amp * (L / (2.0 * M_PI)) * (2.0 / delta);

        // Composante horizontale u sur les faces décalées MAC en x
        int u_nx = mesh.nx_local + (mesh.is_boundary_right ? 1 : 0);
        // Boucle spatiale sur les cellules locales selon l'axe vertical Y (ordonnées)
        for (int j = 1; j <= mesh.ny_local; ++j) {
            // Déclaration de variable typée
            double y = mesh.yc(j);
            // Déclaration de variable typée pour l'écart au centre de la nappe
            double dy_c = y - y_c;
            // Signe pour la divergence solénoïdale
            double sgn_y = (dy_c > 1.0e-12) ? 1.0 : ((dy_c < -1.0e-12) ? -1.0 : 0.0);
            // Décroissance exponentielle symétrique par rapport au plan médian
            double exp_decay = std::exp(-2.0 * std::abs(dy_c) / delta);
            // Boucle spatiale sur les faces u selon l'axe horizontal X
            for (int i = 1; i <= u_nx; ++i) {
                // Déclaration de variable typée
                double x = mesh.xu(i);
                // Vitesse moyenne convective A : 30 m/s dans le gaz, 2 m/s dans le liquide
                double c_u = 0.5 * (fields.C[mesh.cell_idx(i - 1, j)] + fields.C[mesh.cell_idx(i, j)]);
                // Interpolation de la vitesse moyenne de base A selon la présence de gaz/liquide
                double A_val = c_u * 30.0 + (1.0 - c_u) * 2.0;
                // Calcul et assignation numérique de u selon l'Équation 17
                fields.u[mesh.u_idx(i, j)] = A_val - sgn_y * factor_u * std::cos(k_pert * x) * exp_decay;
            }
        }

        // Composante verticale v sur les faces décalées MAC en y
        int v_ny = mesh.ny_local + (mesh.is_boundary_top ? 1 : 0);
        // Boucle spatiale sur les faces v selon l'axe vertical Y
        for (int j = 1; j <= v_ny; ++j) {
            // Déclaration de variable typée
            double y = mesh.yv(j);
            // Déclaration de variable typée pour l'écart au centre de la nappe
            double dy_c = y - y_c;
            // Décroissance exponentielle symétrique par rapport au plan médian
            double exp_decay = std::exp(-2.0 * std::abs(dy_c) / delta);
            // Boucle spatiale sur les cellules locales selon l'axe horizontal X
            for (int i = 1; i <= mesh.nx_local; ++i) {
                // Déclaration de variable typée
                double x = mesh.xc(i);
                // Calcul et assignation numérique de v selon l'Équation 17
                fields.v[mesh.v_idx(i, j)] = amp * std::sin(k_pert * x) * exp_decay;
            }
        }
        MpiComm::sync_u(mesh, fields.u);
        MpiComm::sync_v(mesh, fields.v);
        MpiComm::apply_boundary_conditions(mesh, config.bc, fields, false);
    }

    // Initialisation du champ de vitesse pour VISCOUS_JET (Cas 3.3 Fluids 2021)
    if (config.init_cond.type == InitialConditionType::VISCOUS_JET) {
        // Déclaration des variables physiques et géométriques
        double x_c = (config.init_cond.center_x > 0.0) ? config.init_cond.center_x : 0.5 * mesh.lx;
        // Déclaration de variable typée pour la demi-largeur de l'injecteur
        double half_w = (config.init_cond.jet_width > 0.0) ? 0.5 * config.init_cond.jet_width : 0.04;
        // Déclaration de la hauteur de l'amorce initiale de jet sous la paroi supérieure
        double stub_height = 2.0 * mesh.dy;
        // Vitesse d'injection négative selon Y (ajustée au débit discret exact D * v0)
        double v_inlet = (config.bc.top.v_val != 0.0) ? config.bc.top.v_val : -0.6;
        double d_disc = 0.0;
        for (int ig = 1; ig <= mesh.nx_global; ++ig) {
            double xg = (ig - 0.5) * mesh.dx;
            if (std::abs(xg - x_c) <= half_w) d_disc += mesh.dx;
        }
        if (d_disc < 1.0e-12) d_disc = 2.0 * half_w;
        double v_in_eff = v_inlet * (2.0 * half_w) / d_disc;

        // Composante verticale v sur les faces décalées MAC en y
        int v_ny = mesh.ny_local + (mesh.is_boundary_top ? 1 : 0);
        // Boucle spatiale sur les faces v selon l'axe vertical Y
        for (int j = 1; j <= v_ny; ++j) {
            // Déclaration de variable typée
            double y = mesh.yv(j);
            // Boucle spatiale sur les cellules locales selon l'axe horizontal X
            for (int i = 1; i <= mesh.nx_local; ++i) {
                // Déclaration de variable typée
                double x = mesh.xc(i);
                // Évaluation de l'appartenance à l'amorce initiale
                if (y >= mesh.ly - stub_height && std::abs(x - x_c) <= half_w) {
                    // Calcul et assignation numérique
                    fields.v[mesh.v_idx(i, j)] = v_in_eff;
                }
            }
        }
        MpiComm::sync_v(mesh, fields.v);
        MpiComm::apply_boundary_conditions(mesh, config.bc, fields, false);
    }

    // Initialisation du champ thermique si activé
    if (config.enable_energy) {
        // Boucle spatiale sur les cellules locales selon l'axe vertical Y (ordonnées)
        for (int j = 1; j <= mesh.ny_local; ++j) {
            // Boucle spatiale sur les cellules locales selon l'axe horizontal X (abscisses)
            for (int i = 1; i <= mesh.nx_local; ++i) {
                // Déclaration de variable typée
                double frac = mesh.xc(i) / mesh.lx;
                // Calcul de la résolution géométrique de cellule
                fields.T[mesh.cell_idx(i, j)] = config.bc.left.T_val + frac * (config.bc.right.T_val - config.bc.left.T_val);
            }
        }
        MpiComm::sync_cell_scalar(mesh, fields.T);
        MpiComm::apply_thermal_boundary_conditions(mesh, config.bc, fields);
    }

    // 5. Initialisation du solveur Navier-Stokes et des préconditionneurs
    ns_solver.init(mesh, config);
    ns_solver.assemble_preconditioner_blocks(fields, config.dt);

    // Calcul et assignation numérique
    precond = PreconditionerFactory::create(config.solver.precond_type);
    precond->init(&ns_solver);

    // Calcul et assignation numérique
    linear_solver = SolverFactory::create(config.solver.type);
    linear_solver->init(config.solver);

    // 6. Initialisation de l'exportateur VTK et diagnostics CSV
    vtk_writer.init(mesh, config.output.output_dir, config.output.case_name);

    // Retour de fonction
    return true;
}

// Déclaration de variable typée
double Simulation::compute_time_step() const {
    // Évaluation de la condition de contrôle
    if (!config.adaptive_dt) {
        // Retour de fonction
        return config.dt;
    }

    // Déclaration de variable typée
    double local_max_speed = 1.0e-6;
    // Boucle spatiale sur les cellules locales selon l'axe vertical Y (ordonnées)
    for (int j = 1; j <= mesh.ny_local; ++j) {
        // Boucle spatiale sur les cellules locales selon l'axe horizontal X (abscisses)
        for (int i = 1; i <= mesh.nx_local; ++i) {
            // Déclaration de variable typée
            double u_mag = std::abs(fields.u[mesh.u_idx(i, j)]);
            // Déclaration de variable typée
            double v_mag = std::abs(fields.v[mesh.v_idx(i, j)]);
            // Déclaration de variable typée
            double speed = std::sqrt(u_mag * u_mag + v_mag * v_mag);
            // Évaluation de la condition de contrôle
            if (speed > local_max_speed) local_max_speed = speed;
        }
    }
    // Déclaration de variable typée
    double global_max_speed = MpiComm::global_max(mesh, local_max_speed);

    // 1. Critère CFL de convection
    double dt_cfl = config.cfl_max * std::min(mesh.dx, mesh.dy) / global_max_speed;

    // 2. Critère capillaire (onde de surface) si tension de surface activée
    double dt_sigma = 1.0e10;
    // Évaluation de la condition de contrôle
    if (config.enable_surface_tension && config.sigma > 0.0) {
        // Déclaration de variable typée
        double rho_sum = config.fluid1.rho + config.fluid2.rho;
        // Déclaration de variable typée
        double h_min = std::min(mesh.dx, mesh.dy);
        // Calcul et assignation numérique
        dt_sigma = std::sqrt((rho_sum * h_min * h_min * h_min) / (4.0 * M_PI * config.sigma));
    }

    // Déclaration de variable typée
    double dt_target = std::min(dt_cfl, dt_sigma);
    // Retour de fonction
    return std::min(config.dt, dt_target);
}

void Simulation::run() {
    // Évaluation de la condition de contrôle
    if (config.mode == "LAPLACIAN") {
        run_laplacian();
        // Retour de fonction
        return;
    }
    // Évaluation de la condition de contrôle
    if (config.mode == "VORTEX_ADVECTION") {
        run_vortex_advection();
        // Retour de fonction
        return;
    }

    // Déclaration de variable typée
    double t = 0.0;
    // Déclaration de variable typée
    int step = 0;

    // Sauvegarde initiale à t = 0.0
    vtk_writer.write_step(mesh, fields, step, t);
    vtk_writer.write_diagnostics(mesh, fields, step, t, config.dt);

    // Évaluation de la condition de contrôle
    if (mesh.rank == 0) {
        std::cout << "\n>>> DEMARRAGE DU CALCUL COVOF 2D PARALLELE <<<\n"
                  << std::left << std::setw(8)  << "Pas"
                  << std::setw(14) << "Temps [s]"
                  << std::setw(12) << "dt [s]"
                  << std::setw(12) << "Iter Krylov"
                  << std::setw(16) << "Vitesse Max"
                  << std::setw(16) << "Div Max";
        // Évaluation de la condition de contrôle
        if (config.enable_energy) {
            std::cout << std::setw(14) << "Nu Moyen";
        }
        std::cout << "\n" << std::string(85, '-') << "\n";
    }

    // Déclaration de variable typée
    std::vector<double> state_vector(ns_solver.n_owned_total, 0.0);
    // Déclaration de variable typée
    std::vector<double> rhs(ns_solver.n_owned_total, 0.0);
    // Déclaration de variable typée
    std::vector<double> u_prev = fields.u;

    // Boucle itérative conditionnelle
    while (t < config.t_final - 1.0e-12 * (std::abs(config.t_final) + 1.0)) {
        step++;
        // Calcul et assignation numérique
        u_prev = fields.u;
        // Déclaration de variable typée
        double dt = compute_time_step();
        // Évaluation de la condition de contrôle
        if (t + dt > config.t_final) {
            // Calcul et assignation numérique
            dt = config.t_final - t;
        }
        // Évaluation de la condition de contrôle
        if (dt < 1.0e-12) {
            break;
        }

        // 1. Transport de l'interface si écoulement diphasique
        if (config.enable_two_phase) {
            tracker.advect(mesh, fields, dt);
            // Actualisation des conditions aux limites pour les mailles fantômes
            MpiComm::apply_boundary_conditions(mesh, config.bc, fields, false);
        }

        // 2. Assemblage des blocs du préconditionneur
        ns_solver.assemble_preconditioner_blocks(fields, dt);

        // Déclaration de variable typée
        int total_iters = 0;
        // Déclaration de variable typée
        int uzawa_max = ns_solver.is_augmented_lagrangian() ? ns_solver.uzawa_iter : 1;

        // Boucle d'itération algorithmique
        for (int m = 0; m < uzawa_max; ++m) {
            // 3. Construction du second membre rhs
            ns_solver.assemble_rhs(fields, dt, rhs);

            // 4. Initialisation du vecteur d'état à partir des champs courants
            ns_solver.pack_vector(fields, state_vector);

            // 5. Résolution du système linéaire par le solveur sélectionné
            int iters = linear_solver->solve(ns_solver, *precond, fields, dt, rhs, state_vector);
            // Calcul et assignation numérique
            total_iters += iters;

            // 6. Déballage de la vitesse résolue
            ns_solver.unpack_vector(state_vector, fields);

            // 7. Mise à jour de la pression par Lagrangien augmenté si applicable
            if (ns_solver.is_augmented_lagrangian()) {
                ns_solver.update_pressure_al(fields);
            }
        }

        // 7.5 Résolution de l'énergie thermique si activée
        if (config.enable_energy) {
            ns_solver.solve_energy(fields, dt);
        }

        // Calcul et assignation numérique
        t += dt;

        // Calcul des métriques de contrôle
        double local_max_u = 0.0;
        // Déclaration de variable typée
        double local_max_div = 0.0;
        // Déclaration de variable typée
        double local_du_max = 0.0;
        // Boucle spatiale sur les cellules locales selon l'axe vertical Y (ordonnées)
        for (int j = 1; j <= mesh.ny_local; ++j) {
            // Boucle spatiale sur les cellules locales selon l'axe horizontal X (abscisses)
            for (int i = 1; i <= mesh.nx_local; ++i) {
                // Déclaration de variable typée
                double speed = std::sqrt(fields.u[mesh.u_idx(i, j)] * fields.u[mesh.u_idx(i, j)] +
                                         fields.v[mesh.v_idx(i, j)] * fields.v[mesh.v_idx(i, j)]);
                // Évaluation de la condition de contrôle
                if (speed > local_max_u) local_max_u = speed;

                // Déclaration de variable typée
                double du = std::abs(fields.u[mesh.u_idx(i, j)] - u_prev[mesh.u_idx(i, j)]);
                // Évaluation de la condition de contrôle
                if (du > local_du_max) local_du_max = du;

                // Déclaration de variable typée
                double div = std::abs((fields.u[mesh.u_idx(i + 1, j)] - fields.u[mesh.u_idx(i, j)]) / mesh.dx +
                                      (fields.v[mesh.v_idx(i, j + 1)] - fields.v[mesh.v_idx(i, j)]) / mesh.dy);
                // Évaluation de la condition de contrôle
                if (div > local_max_div) local_max_div = div;
            }
        }
        // Déclaration de variable typée
        double global_max_u = MpiComm::global_max(mesh, local_max_u);
        // Déclaration de variable typée
        double global_max_div = MpiComm::global_max(mesh, local_max_div);
        // Déclaration de variable typée
        double global_du_max = MpiComm::global_max(mesh, local_du_max);

        // Déclaration de variable typée
        double nu_mean = 0.0;
        // Évaluation de la condition de contrôle
        if (config.enable_energy) {
            // Calcul et assignation numérique
            nu_mean = ns_solver.compute_mean_nusselt_hot_wall(fields);
        }

        // Détection de convergence stationnaire
        if (config.steady_state_tol > 0.0 && step > 10 && global_du_max < config.steady_state_tol) {
            // Évaluation de la condition de contrôle
            if (mesh.rank == 0) {
                // Calcul et assignation numérique
                std::cout << "\n>>> CONVERGENCE STATIONNAIRE ATTEINTE : du_max = "
                          // Calcul de la résolution géométrique de cellule
                          << global_du_max << " < " << config.steady_state_tol << " (t = " << t << " s) <<<\n";
                // Évaluation de la condition de contrôle
                if (config.enable_energy) {
                    // Calcul et assignation numérique
                    std::cout << ">>> NOMBRE DE NUSSELT FINAL : Nu = " << std::fixed << std::setprecision(5) << nu_mean << " <<<\n";
                }
            }
            vtk_writer.write_step(mesh, fields, step, t);
            break;
        }

        // Affichage console
        if (mesh.rank == 0 && (step % config.output.log_interval == 0)) {
            std::cout << std::left << std::setw(8)  << step
                      << std::setw(14) << std::fixed << std::setprecision(5) << t
                      << std::setw(12) << std::scientific << std::setprecision(3) << dt
                      << std::setw(12) << total_iters
                      << std::setw(16) << std::scientific << std::setprecision(4) << global_max_u
                      << std::setw(16) << std::scientific << std::setprecision(4) << global_max_div;
            // Évaluation de la condition de contrôle
            if (config.enable_energy) {
                std::cout << std::setw(14) << std::fixed << std::setprecision(4) << nu_mean;
            }
            std::cout << "\n";
        }

        // Écriture des diagnostics temporels
        if (config.output.write_csv_history) {
            vtk_writer.write_diagnostics(mesh, fields, step, t, dt);
        }

        // Export des fichiers VTK
        if (step % config.output.vtk_interval == 0 || t >= config.t_final) {
            vtk_writer.write_step(mesh, fields, step, t);
        }
    }

    vtk_writer.finalize();

    // Évaluation de la condition de contrôle
    if (mesh.rank == 0) {
        std::cout << std::string(78, '-') << "\n"
                  // Calcul et assignation numérique
                  << ">>> SIMULATION COVOF 2D TERMINEE AVEC SUCCES A t = " << t << " s <<<\n"
                  << "Fichiers de visualisation VTK et CSV disponibles dans : " << config.output.output_dir << "\n";
    }
}

void Simulation::run_laplacian() {
    // Déclaration de variable typée
    double a = config.lx;
    // Déclaration de variable typée
    double b = config.ly;
    // Déclaration de variable typée
    int nx = config.nx;
    // Déclaration de variable typée
    int ny = config.ny;
    // Déclaration de variable typée
    double dx = a / nx;
    // Déclaration de variable typée
    double dy = b / ny;

    // Évaluation de la condition de contrôle
    if (mesh.rank == 0) {
        // Calcul et assignation numérique
        std::cout << "\n========================================================\n"
                  // Calcul et assignation numérique
                  << "  VERIFICATION DU CAS 3.1 : LAPLACIEN DELTA(PSI) = 1     \n"
                  << "  Domaine [" << a << " x " << b << "], Grille " << nx << " x " << ny << "\n"
                  // Calcul et assignation numérique
                  << "========================================================\n";
    }

    // Inconnues nodales intérieures : 1 à nx-1, 1 à ny-1
    int n_int_x = nx - 1;
    // Déclaration de variable typée
    int n_int_y = ny - 1;
    // Déclaration de variable typée
    int n_dofs = n_int_x * n_int_y;

    CSRMatrix A;
    // Calcul et assignation numérique
    A.nrows = n_dofs;
    // Calcul et assignation numérique
    A.ncols = n_dofs;
    A.row_ptr.resize(n_dofs + 1, 0);

    // Calcul de la résolution géométrique de cellule
    auto dof_idx = [&](int i, int j) { return j * n_int_x + i; };

    // 1. Comptage des éléments non nuls
    for (int j = 0; j < n_int_y; ++j) {
        // Boucle spatiale sur les cellules locales selon l'axe horizontal X (abscisses)
        for (int i = 0; i < n_int_x; ++i) {
            // Déclaration de variable typée
            int row = dof_idx(i, j);
            // Déclaration de variable typée
            int count = 1;
            // Évaluation de la condition de contrôle
            if (i > 0) count++;
            // Évaluation de la condition de contrôle
            if (i < n_int_x - 1) count++;
            // Évaluation de la condition de contrôle
            if (j > 0) count++;
            // Évaluation de la condition de contrôle
            if (j < n_int_y - 1) count++;
            // Calcul et assignation numérique
            A.row_ptr[row + 1] = A.row_ptr[row] + count;
        }
    }

    A.col_idx.resize(A.row_ptr[n_dofs]);
    A.values.resize(A.row_ptr[n_dofs]);

    // Déclaration de variable typée
    std::vector<double> rhs(n_dofs, -1.0); // -Delta Psi = -1
    // Déclaration de variable typée
    std::vector<double> x(n_dofs, 0.0);

    // Déclaration de variable typée
    int cur = 0;
    // Boucle spatiale sur les cellules locales selon l'axe vertical Y (ordonnées)
    for (int j = 0; j < n_int_y; ++j) {
        // Boucle spatiale sur les cellules locales selon l'axe horizontal X (abscisses)
        for (int i = 0; i < n_int_x; ++i) {
            // Déclaration de variable typée
            int row = dof_idx(i, j);
            // Évaluation de la condition de contrôle
            if (j > 0) {
                // Calcul de la résolution géométrique de cellule
                A.col_idx[cur] = dof_idx(i, j - 1);
                // Calcul et assignation numérique
                A.values[cur] = -1.0 / (dy * dy);
                cur++;
            }
            // Évaluation de la condition de contrôle
            if (i > 0) {
                // Calcul de la résolution géométrique de cellule
                A.col_idx[cur] = dof_idx(i - 1, j);
                // Calcul et assignation numérique
                A.values[cur] = -1.0 / (dx * dx);
                cur++;
            }
            // Calcul de la résolution géométrique de cellule
            A.col_idx[cur] = row;
            // Calcul et assignation numérique
            A.values[cur] = 2.0 / (dx * dx) + 2.0 / (dy * dy);
            cur++;
            // Évaluation de la condition de contrôle
            if (i < n_int_x - 1) {
                // Calcul de la résolution géométrique de cellule
                A.col_idx[cur] = dof_idx(i + 1, j);
                // Calcul et assignation numérique
                A.values[cur] = -1.0 / (dx * dx);
                cur++;
            }
            // Évaluation de la condition de contrôle
            if (j < n_int_y - 1) {
                // Calcul de la résolution géométrique de cellule
                A.col_idx[cur] = dof_idx(i, j + 1);
                // Calcul et assignation numérique
                A.values[cur] = -1.0 / (dy * dy);
                cur++;
            }
        }
    }

    // Solveur Gradient Conjugué préconditionné par Jacobi
    std::vector<double> r = rhs;
    // Déclaration de variable typée
    std::vector<double> diag(n_dofs, 1.0);
    // Boucle spatiale sur les cellules locales selon l'axe horizontal X (abscisses)
    for (int i = 0; i < n_dofs; ++i) {
        // Boucle d'itération algorithmique
        for (int p_idx = A.row_ptr[i]; p_idx < A.row_ptr[i+1]; ++p_idx) {
            // Évaluation de la condition de contrôle
            if (A.col_idx[p_idx] == i) { diag[i] = A.values[p_idx]; break; }
        }
    }

    // Déclaration de variable typée
    std::vector<double> p_cg(n_dofs, 0.0), z(n_dofs, 0.0), Ap(n_dofs, 0.0);
    // Boucle spatiale sur les cellules locales selon l'axe horizontal X (abscisses)
    for (int i = 0; i < n_dofs; ++i) z[i] = r[i] / diag[i];
    // Calcul et assignation numérique
    p_cg = z;
    // Déclaration de variable typée
    double rz_old = 0.0;
    // Boucle spatiale sur les cellules locales selon l'axe horizontal X (abscisses)
    for (int i = 0; i < n_dofs; ++i) rz_old += r[i] * z[i];

    // Déclaration de variable typée
    int max_it = 20000;
    // Déclaration de variable typée
    double tol = config.solver.tol > 0 ? config.solver.tol : 1e-12;
    // Déclaration de variable typée
    int iters = 0;

    // Boucle d'itération algorithmique
    for (iters = 0; iters < max_it; ++iters) {
        A.matvec(p_cg, Ap);
        // Déclaration de variable typée
        double pAp = 0.0;
        // Boucle spatiale sur les cellules locales selon l'axe horizontal X (abscisses)
        for (int i = 0; i < n_dofs; ++i) pAp += p_cg[i] * Ap[i];
        // Évaluation de la condition de contrôle
        if (std::abs(pAp) < 1e-30) break;
        // Déclaration de variable typée
        double alpha = rz_old / pAp;
        // Déclaration de variable typée
        double r_norm = 0.0;
        // Boucle spatiale sur les cellules locales selon l'axe horizontal X (abscisses)
        for (int i = 0; i < n_dofs; ++i) {
            // Calcul et assignation numérique
            x[i] += alpha * p_cg[i];
            // Calcul et assignation numérique
            r[i] -= alpha * Ap[i];
            // Calcul de norme euclidienne ou résidu de convergence
            r_norm += r[i] * r[i];
        }
        // Calcul de norme euclidienne ou résidu de convergence
        r_norm = std::sqrt(r_norm);
        // Évaluation de la condition de contrôle
        if (r_norm < tol) break;

        // Boucle spatiale sur les cellules locales selon l'axe horizontal X (abscisses)
        for (int i = 0; i < n_dofs; ++i) z[i] = r[i] / diag[i];
        // Déclaration de variable typée
        double rz_new = 0.0;
        // Boucle spatiale sur les cellules locales selon l'axe horizontal X (abscisses)
        for (int i = 0; i < n_dofs; ++i) rz_new += r[i] * z[i];
        // Déclaration de variable typée
        double beta = rz_new / rz_old;
        // Calcul et assignation numérique
        rz_old = rz_new;
        // Boucle spatiale sur les cellules locales selon l'axe horizontal X (abscisses)
        for (int i = 0; i < n_dofs; ++i) p_cg[i] = z[i] + beta * p_cg[i];
    }

    // Évaluation de la solution analytique exacte (formule analytique de référence)
    double err_l1 = 0.0, err_l2 = 0.0, err_linf = 0.0;
    // Boucle spatiale sur les cellules locales selon l'axe vertical Y (ordonnées)
    for (int j = 0; j < n_int_y; ++j) {
        // Boucle spatiale sur les cellules locales selon l'axe horizontal X (abscisses)
        for (int i = 0; i < n_int_x; ++i) {
            // Déclaration de variable typée
            double xp = (i + 1) * dx;
            // Déclaration de variable typée
            double yp = (j + 1) * dy;
            // Déclaration de variable typée
            double psi_ex = 0.0;
            // Boucle d'itération algorithmique
            for (int k2 = 0; k2 < 60; ++k2) {
                // Boucle d'itération algorithmique
                for (int k1 = 0; k1 < 60; ++k1) {
                    // Déclaration de variable typée
                    double n1 = 2 * k1 + 1;
                    // Déclaration de variable typée
                    double n2 = 2 * k2 + 1;
                    // Déclaration de variable typée
                    double num = 16.0 * std::sin(n1 * M_PI * xp / a) * std::sin(n2 * M_PI * yp / b);
                    // Déclaration de variable typée
                    double den = (std::pow(M_PI, 4.0)) * n1 * n2 * (std::pow(n1 / a, 2.0) + std::pow(n2 / b, 2.0));
                    // Calcul et assignation numérique
                    psi_ex -= num / den;
                }
            }
            // Déclaration de variable typée
            double diff = std::abs(x[dof_idx(i, j)] - psi_ex);
            // Calcul et assignation numérique
            err_l1 += diff;
            // Calcul et assignation numérique
            err_l2 += diff * diff;
            // Évaluation de la condition de contrôle
            if (diff > err_linf) err_linf = diff;
        }
    }
    // Calcul et assignation numérique
    err_l1 /= n_dofs;
    // Calcul et assignation numérique
    err_l2 = std::sqrt(err_l2 / n_dofs);

    // Évaluation de la condition de contrôle
    if (mesh.rank == 0) {
        std::cout << ">>> RESULTAT RESOLUTION LAPLACIEN COVOF 2D <<<\n"
                  << "  Iterations : " << iters << "\n"
                  << "  Erreur L1   : " << std::scientific << std::setprecision(5) << err_l1 << "\n"
                  << "  Erreur L2   : " << std::scientific << std::setprecision(5) << err_l2 << "\n"
                  << "  Erreur Linf : " << std::scientific << std::setprecision(5) << err_linf << "\n"
                  // Calcul et assignation numérique
                  << "========================================================\n";

        // Déclaration de variable typée
        std::string err_file = config.output.output_dir + "/laplacian_error.csv";
        std::ofstream ofs(err_file);
        // Évaluation de la condition de contrôle
        if (ofs.is_open()) {
            ofs << "nx,ny,dx,dy,iters,err_l1,err_l2,err_linf\n"
                << nx << "," << ny << "," << dx << "," << dy << ","
                << iters << "," << err_l1 << "," << err_l2 << "," << err_linf << "\n";
            ofs.close();
        }
    }
}

void Simulation::run_vortex_advection() {
    // Déclaration de variable typée
    int nx = config.nx;
    // Déclaration de variable typée
    int ny = config.ny;
    // Déclaration de variable typée
    double T_period = (config.t_final > 0.0) ? config.t_final : 6.0;
    // Déclaration de variable typée
    double dt = config.dt;

    // Évaluation de la condition de contrôle
    if (mesh.rank == 0) {
        // Calcul et assignation numérique
        std::cout << "\n========================================================\n"
                  << "  VERIFICATION DU CAS 3.19 : ADVECTION VORTEX CISAILLEMENT\n"
                  // Calcul et assignation numérique
                  << "  Grille : " << nx << " x " << ny << " | Periode T = " << T_period << " s | dt = " << dt << " s\n"
                  // Calcul et assignation numérique
                  << "========================================================\n";
    }

    // Sauvegarder la condition initiale C_exact
    std::vector<double> C_initial = fields.C;

    // Export initial
    vtk_writer.write_step(mesh, fields, 0, 0.0);

    // Déclaration de variable typée
    double t = 0.0;
    // Déclaration de variable typée
    int step = 0;

    // Déclaration de variable typée
    int u_nx = mesh.nx_local + (mesh.is_boundary_right ? 1 : 0);
    // Déclaration de variable typée
    int v_ny = mesh.ny_local + (mesh.is_boundary_top ? 1 : 0);

    // Boucle itérative conditionnelle
    while (t < T_period - 1.0e-12) {
        step++;
        // Évaluation de la condition de contrôle
        if (t + dt > T_period) dt = T_period - t;

        // Déclaration de variable typée
        double cos_factor = std::cos(M_PI * t / T_period);

        // Imposer le champ de vitesse analytique sur les faces décalées
        for (int j = 1; j <= mesh.ny_local; ++j) {
            // Déclaration de variable typée
            double y_c = mesh.yc(j);
            // Boucle spatiale sur les cellules locales selon l'axe horizontal X (abscisses)
            for (int i = 1; i <= u_nx; ++i) {
                // Déclaration de variable typée
                double x_u = mesh.xu(i);
                // Déclaration de variable typée
                double sx = std::sin(M_PI * x_u);
                // Déclaration de variable typée
                double sy = std::sin(2.0 * M_PI * y_c);
                // Calcul de la résolution géométrique de cellule
                fields.u[mesh.u_idx(i, j)] = - (sx * sx) * sy * cos_factor;
            }
        }

        // Boucle spatiale sur les cellules locales selon l'axe vertical Y (ordonnées)
        for (int j = 1; j <= v_ny; ++j) {
            // Déclaration de variable typée
            double y_v = mesh.yv(j);
            // Boucle spatiale sur les cellules locales selon l'axe horizontal X (abscisses)
            for (int i = 1; i <= mesh.nx_local; ++i) {
                // Déclaration de variable typée
                double x_c = mesh.xc(i);
                // Déclaration de variable typée
                double sx = std::sin(2.0 * M_PI * x_c);
                // Déclaration de variable typée
                double sy = std::sin(M_PI * y_v);
                // Calcul de la résolution géométrique de cellule
                fields.v[mesh.v_idx(i, j)] = sx * (sy * sy) * cos_factor;
            }
        }

        // Halos et CL
        MpiComm::sync_u(mesh, fields.u);
        MpiComm::sync_v(mesh, fields.v);

        // Advection TVD de C
        tracker.advect(mesh, fields, dt);

        // Calcul et assignation numérique
        t += dt;

        // Évaluation de la condition de contrôle
        if (step % config.output.vtk_interval == 0 || t >= T_period || std::abs(t - 0.5 * T_period) < 0.5 * dt) {
            vtk_writer.write_step(mesh, fields, step, t);
        }

        // Évaluation de la condition de contrôle
        if (mesh.rank == 0 && (step % config.output.log_interval == 0 || t >= T_period)) {
            // Calcul et assignation numérique
            std::cout << "  [Vortex] Pas " << step << " | t = " << std::fixed << std::setprecision(4) << t
                      << " s / " << T_period << " s\n";
        }
    }

    // Calcul des erreurs L1, L2, Linf et du bilan de masse entre C_initial et C(T)
    double local_err_l1 = 0.0;
    // Déclaration de variable typée
    double local_err_l2 = 0.0;
    // Déclaration de variable typée
    double local_err_linf = 0.0;
    // Déclaration de variable typée
    double local_mass_init = 0.0;
    // Déclaration de variable typée
    double local_mass_final = 0.0;
    // Déclaration de variable typée
    double local_cells = mesh.nx_local * mesh.ny_local;
    // Déclaration de variable typée
    double cell_vol = mesh.dx * mesh.dy;

    // Boucle spatiale sur les cellules locales selon l'axe vertical Y (ordonnées)
    for (int j = 1; j <= mesh.ny_local; ++j) {
        // Boucle spatiale sur les cellules locales selon l'axe horizontal X (abscisses)
        for (int i = 1; i <= mesh.nx_local; ++i) {
            // Déclaration de variable typée
            int idx = mesh.cell_idx(i, j);
            // Déclaration de variable typée
            double diff = std::abs(fields.C[idx] - C_initial[idx]);
            // Calcul et assignation numérique
            local_err_l1 += diff;
            // Calcul et assignation numérique
            local_err_l2 += diff * diff;
            // Évaluation de la condition de contrôle
            if (diff > local_err_linf) local_err_linf = diff;
            // Calcul et assignation numérique
            local_mass_init += C_initial[idx] * cell_vol;
            // Calcul et assignation numérique
            local_mass_final += fields.C[idx] * cell_vol;
        }
    }

    // Réduction parallèle MPI globale de la somme de l'erreur L1
    double global_err_l1 = MpiComm::global_sum(mesh, local_err_l1);
    // Réduction parallèle MPI globale de la somme de l'erreur L2 au carré
    double global_err_l2 = MpiComm::global_sum(mesh, local_err_l2);
    // Réduction parallèle MPI globale de l'erreur maximale absolue Linf
    double global_err_linf = MpiComm::global_max(mesh, local_err_linf);
    // Réduction parallèle MPI globale de la masse totale initiale
    double global_mass_init = MpiComm::global_sum(mesh, local_mass_init);
    // Réduction parallèle MPI globale de la masse totale finale
    double global_mass_final = MpiComm::global_sum(mesh, local_mass_final);
    // Réduction parallèle MPI globale du nombre de cellules
    double global_cells = MpiComm::global_sum(mesh, local_cells);

    // Calcul de l'erreur moyenne normalisée L1
    double err_l1 = global_err_l1 / global_cells;
    // Calcul de l'erreur quadratique moyenne normalisée L2
    double err_l2 = std::sqrt(global_err_l2 / global_cells);
    // Erreur maximale locale Linf
    double err_linf = global_err_linf;
    // Taux de conservation relative de la masse du fluide
    double mass_conservation = (global_mass_init > 1.0e-14) ? (global_mass_final / global_mass_init) : 1.0;

    // Évaluation de la condition de contrôle
    if (mesh.rank == 0) {
        // Affichage console formaté du résultat de vérification
        std::cout << "\n>>> RESULTAT DU CAS 3.19 : ADVECTION VORTEX COVOF 2D <<<\n"
                  << "  Methode Advection  : " << config.vof.advection_method << "\n"
                  << "  Reconstruction     : " << config.vof.reconstruction_method << "\n"
                  << "  Limiteur TVD/NVD   : " << config.vof.tvd_limiter << "\n"
                  << "  Maillage           : " << nx << " x " << ny << "\n"
                  << "  Erreur L1          : " << std::scientific << std::setprecision(5) << err_l1 << "\n"
                  << "  Erreur L2          : " << std::scientific << std::setprecision(5) << err_l2 << "\n"
                  << "  Erreur Linf        : " << std::scientific << std::setprecision(5) << err_linf << "\n"
                  << "  Conservation Masse : " << std::fixed << std::setprecision(5) << (mass_conservation * 100.0) << " %\n"
                  << "========================================================\n";

        // Déclaration de variable typée
        std::string err_file = config.output.output_dir + "/vortex_error.csv";
        // Ouverture du flux d'écriture fichier
        std::ofstream ofs(err_file);
        // Évaluation de la condition de contrôle
        if (ofs.is_open()) {
            // Écriture de l'en-tête CSV enrichi avec toutes les métriques et métadonnées d'advection
            ofs << "nx,ny,dt,T,advection,reconstruction,limiter,err_l1,err_l2,err_linf,mass_init,mass_final,mass_conservation\n"
                << nx << "," << ny << "," << dt << "," << T_period << ","
                << config.vof.advection_method << "," << config.vof.reconstruction_method << "," << config.vof.tvd_limiter << ","
                << err_l1 << "," << err_l2 << "," << err_linf << ","
                << global_mass_init << "," << global_mass_final << "," << mass_conservation << "\n";
            // Fermeture du fichier de sortie
            ofs.close();
        }
    }

    vtk_writer.finalize();
}

} // namespace covof


\end{lstlisting}

\part{Guide du Développeur Fortran 2018 (\texttt{COVOF\_2D-FORTRAN})}
\chapter{Architecture Modulaire Fortran 2018, Précision IEEE-754 et Programme Principal (CMakeLists.txt, module_precision.f90, module_constants.f90, covof2d_main.f90)}
\section{Conception Modulaire selon la Norme Fortran 2018 (ISO/IEC 1539-1:2018)}
L'implémentation \textbf{Fortran 2018} (\texttt{COVOF\_2D-FORTRAN}) est le miroir algorithmique et mathématique exact de \texttt{COVOF\_2D-C++}. Elle est structurée en 13 modules modernes (\texttt{implicit none} systématique, attributs \texttt{intent(in)}, \texttt{intent(out)}, \texttt{intent(inout)}, procédures attachées aux types dérivés et tableaux allocatables à bornes personnalisées) et un programme principal \texttt{covof2d\_main.f90}.

\subsection{Gestion de la Précision Flottante (\texttt{module\_precision.f90})}
Le module \texttt{module\_precision} définit le paramètre de genre (\emph{kind parameter}) double précision IEEE-754 portable sur 64 bits :
\begin{equation}
  \texttt{integer, parameter :: dp = selected\_real\_kind(15, 307)}
\end{equation}
garantissant au moins 15 chiffres décimaux significatifs de mantisse et une plage d'exposants décimaux de $\pm 307$, strictement identique au type \texttt{double} de C++17.

\subsection{Table de Correspondance Bijective C++17 $\leftrightarrow$ Fortran 2018}
\begin{center}
\begin{tabular}{ll}
\hline
\textbf{Composant C++17 (\texttt{COVOF\_2D-C++})} & \textbf{Module Fortran 2018 (\texttt{COVOF\_2D-FORTRAN})} \\
\hline
\texttt{include/covof/Types.hpp} & \texttt{src/module\_precision.f90} + \texttt{src/module\_constants.f90} \\
\texttt{src/Config.cpp} (\texttt{covof::Config}) & \texttt{src/module\_json.f90} + \texttt{src/module\_config.f90} \\
\texttt{src/Mesh.cpp} (\texttt{covof::Mesh}) & \texttt{src/module\_mesh.f90} (\texttt{type(mesh\_2d\_t)}) \\
\texttt{include/covof/Field.hpp} (\texttt{Field2D<T>}) & \texttt{src/module\_fields.f90} (\texttt{type(flow\_fields\_t)}) \\
\texttt{src/MpiComm.cpp} (\texttt{covof::MpiComm}) & \texttt{src/module\_mpi\_comm.f90} (\texttt{type(mpi\_context\_t)}) \\
\texttt{src/CSRMatrix.cpp} + \texttt{Preconditioner.cpp} + \texttt{LinearSolver.cpp} & \texttt{src/module\_linear\_solvers.f90} \\
\texttt{src/NavierStokesSolver.cpp} & \texttt{src/module\_navier\_stokes.f90} \\
\texttt{src/InterfaceReconstruction.cpp} & \texttt{src/module\_interface\_reconstruction.f90} \\
\texttt{src/InterfaceTracker.cpp} & \texttt{src/module\_interface\_vof.f90} \\
\texttt{src/VTKWriter.cpp} & \texttt{src/module\_vtk\_writer.f90} \\
\texttt{src/Simulation.cpp} + \texttt{src/main.cpp} & \texttt{src/module\_simulation.f90} + \texttt{src/covof2d\_main.f90} \\
\hline
\end{tabular}
\end{center}

\section{Code Source Intégral Annoté : \texttt{CMakeLists.txt}}
Le listing ci-dessous présente l'intégralité des 42 lignes du fichier \texttt{CMakeLists.txt} de l'implémentation Fortran 2018 :

\begin{lstlisting}[language={bash}]
cmake_minimum_required(VERSION 3.18)
project(COVOF_2D_FORTRAN LANGUAGES Fortran)

set(CMAKE_Fortran_STANDARD 2018)
set(CMAKE_Fortran_STANDARD_REQUIRED ON)

find_package(MPI REQUIRED)

set(CMAKE_Fortran_FLAGS "${CMAKE_Fortran_FLAGS} -O3 -march=native")

set(SOURCES
    src/module_precision.f90
    src/module_constants.f90
    src/module_json.f90
    src/module_config.f90
    src/module_mesh.f90
    src/module_fields.f90
    src/module_mpi_comm.f90
    src/module_interface_reconstruction.f90
    src/module_interface_vof.f90
    src/module_navier_stokes.f90
    src/module_linear_solvers.f90
    src/module_vtk_writer.f90
    src/module_simulation.f90
    src/covof2d_main.f90
)

# Executable target: covof2d_fortran
add_executable(covof2d_fortran ${SOURCES})
target_link_libraries(covof2d_fortran PRIVATE MPI::MPI_Fortran)

set_target_properties(covof2d_fortran PROPERTIES
    RUNTIME_OUTPUT_DIRECTORY "${CMAKE_CURRENT_SOURCE_DIR}/bin"
)

# Create compatibility symlink for fugu2d_fortran
add_custom_command(TARGET covof2d_fortran POST_BUILD
    COMMAND ${CMAKE_COMMAND} -E create_symlink
        ${CMAKE_CURRENT_SOURCE_DIR}/bin/covof2d_fortran
        ${CMAKE_CURRENT_SOURCE_DIR}/bin/fugu2d_fortran
    COMMENT "Creating compatibility symlink bin/fugu2d_fortran -> bin/covof2d_fortran"
)

\end{lstlisting}

\section{Répertoire des Types Dérivés, Sous-Routines et Fonctions de \texttt{src/module\_precision.f90}}
\begin{center}\small
\begin{tabular}{clll}
\toprule
\textbf{Ligne} & \textbf{Nature} & \textbf{Identifiant Fortran 2018} & \textbf{Interface / Arguments} \\
\midrule
7 & module & \texttt{module\_precision} & Unité modulaire Fortran 2018 \\
\bottomrule
\end{tabular}
\end{center}

\section{Code Source Intégral Annoté : \texttt{src/module\_precision.f90}}
Le listing ci-dessous présente l'intégralité des 33 lignes du fichier \texttt{src/module\_precision.f90} de l'implémentation Fortran 2018 :

\begin{lstlisting}[language={[95]Fortran}]
!===============================================================================
! Module : module_precision
! Description : Définition standardisée de la précision arithmétique flottante
!               et entière pour l'ensemble du solveur CFD COVOF 2D (Fortran 2018).
! Conforme aux recommandations IEEE 754 et ISO_FORTRAN_ENV.
!===============================================================================
module module_precision
  ! Obligation stricte de déclarer tous les types de variables (évite le typage implicite)
  implicit none
  ! Encapsulation par défaut : aucune entité n'est visible de l'extérieur sans export
  private

  ! Précision double (Double Precision 64-bit IEEE 754) :
  ! - Au moins 15 chiffres significatifs en base 10 (mantisse de 53 bits)
  ! - Plage dynamique d'exposant de 10^-307 à 10^+307 (exposant de 11 bits)
  ! - Utilisée pour tous les calculs hydrodynamiques, de pression et de transport
  integer, parameter, public :: dp = selected_real_kind(15, 307)
  
  ! Précision simple (Single Precision 32-bit IEEE 754) :
  ! - Au moins 6 chiffres significatifs en base 10 (mantisse de 24 bits)
  ! - Plage dynamique d'exposant de 10^-37 à 10^+37 (exposant de 8 bits)
  ! - Utilisée optionnellement pour les exports VTK légers ou diagnostics rapides
  integer, parameter, public :: sp = selected_real_kind(6, 37)

  ! Entier standard 32 bits signé (-2*10^9 à +2*10^9) :
  ! - Utilisé pour l'indexation spatiale des mailles (i, j) et les compteurs de boucle
  integer, parameter, public :: i4 = selected_int_kind(9)
  ! Alias d'entier par défaut pour le maillage et les solveurs
  integer, parameter, public :: ip = i4
  ! Entier 64 bits signé (-9*10^18 à +9*10^18) pour les grands volumes de données
  integer, parameter, public :: i8 = selected_int_kind(18)

end module module_precision

\end{lstlisting}

\section{Répertoire des Types Dérivés, Sous-Routines et Fonctions de \texttt{src/module\_constants.f90}}
\begin{center}\small
\begin{tabular}{clll}
\toprule
\textbf{Ligne} & \textbf{Nature} & \textbf{Identifiant Fortran 2018} & \textbf{Interface / Arguments} \\
\midrule
7 & module & \texttt{module\_constants} & Unité modulaire Fortran 2018 \\
\bottomrule
\end{tabular}
\end{center}

\section{Code Source Intégral Annoté : \texttt{src/module\_constants.f90}}
Le listing ci-dessous présente l'intégralité des 111 lignes du fichier \texttt{src/module\_constants.f90} de l'implémentation Fortran 2018 :

\begin{lstlisting}[language={[95]Fortran}]
!===============================================================================
! Module : module_constants
! Description : Définition des constantes physiques fondamentales, paramètres
!               numériques seuils, énumérations de conditions aux limites, types
!               de solveurs et schémas convectifs pour le code COVOF 2D.
!===============================================================================
module module_constants
  ! Importation du paramètre de précision double standardisé
  use module_precision, only: dp
  ! Obligation stricte de typage explicite pour toutes les constantes
  implicit none
  ! Encapsulation : seules les constantes marquées 'public' sont exportées
  private

  ! --- 1. Constantes Mathématiques et Physiques Fondamentales ---
  ! Valeur exacte du nombre Pi (35 décimales significatives en précision double)
  real(dp), parameter, public :: PI = 3.141592653589793238462643383279502884197_dp
  ! Seuil numérique infinitésimal pour prévenir les divisions par zéro dans les gradients
  real(dp), parameter, public :: EPS_SMALL = 1.0e-30_dp
  ! Perméabilité de Darcy par défaut dans les obstacles solides pénalisés (loi de Brinkman)
  real(dp), parameter, public :: PERMEABILITY_SOLID_DEFAULT = 1.0e-40_dp

  ! --- 2. Conditions aux Limites Hydrodynamiques (Vitesse) ---
  ! Vitesse imposée (Adhérence u=0, paroi mobile de cavité, ou profil d'entrée de Poiseuille)
  integer, parameter, public :: BC_DIRICHLET = 1
  ! Sortie libre convective ou gradient normal nul (du/dn = 0, dv/dn = 0)
  integer, parameter, public :: BC_NEUMANN   = 2
  ! Raccordement spatial périodique (face Est reliée à la face Ouest, Nord à Sud)
  integer, parameter, public :: BC_PERIODIC  = 3
  ! Glissement parfait sans frottement tangentiel (vitesse normale nulle, cisaillement nul)
  integer, parameter, public :: BC_FREESLIP  = 4

  ! --- 3. Conditions aux Limites Thermiques (Énergie) ---
  ! Température pariétale imposée (Paroi isotherme chaude ou froide de Vahl Davis)
  integer, parameter, public :: THERMAL_DIRICHLET = 1
  ! Flux thermique imposé ou paroi adiabatique sans transfert de chaleur (dT/dn = 0)
  integer, parameter, public :: THERMAL_NEUMANN   = 2

  ! --- 4. Formulations du Couplage Vitesse-Pression Navier-Stokes ---
  ! Méthode du Lagrangien Augmenté avec pénalité de divergence et actualisation d'Uzawa
  integer, parameter, public :: SOLVER_AUGMENTED_LAGRANGIAN = 1
  ! Formulation couplée monolithique directe avec intégration du gradient de pression
  integer, parameter, public :: SOLVER_MONOLITHIC           = 2

  ! --- 5. Algorithmes de Résolution Linéaire de Krylov (Matrix-Free) ---
  ! Biconjugué Gradient Stabilisé d'ordre 2 (Sleijpen & van der Vorst) pour systèmes asymétriques
  integer, parameter, public :: LINEAR_BICGSTAB2 = 1
  ! Biconjugué Gradient Stabilisé standard (van der Vorst 1992)
  integer, parameter, public :: LINEAR_BICGSTAB  = 2
  ! Gradient Conjugué Préconditionné pour opérateurs symétriques définis positifs (Poisson)
  integer, parameter, public :: LINEAR_CG        = 3
  ! Generalized Minimal Residual (Saad & Schultz 1986) avec orthogonalisation de Krylov
  integer, parameter, public :: LINEAR_GMRES     = 4
  ! Solveur direct bande ou factorisation LU locale
  integer, parameter, public :: LINEAR_DIRECT    = 5

  ! --- 6. Préconditionneurs Linéaires Parallèles ---
  ! Aucun préconditionnement (résolution brute de l'opérateur non préconditionné)
  integer, parameter, public :: PRECOND_NONE   = 0
  ! Préconditionneur diagonal par blocs de Jacobi (inversion de la diagonale locale)
  integer, parameter, public :: PRECOND_JACOBI = 1
  ! Symmetric Successive Over-Relaxation (SSOR) pour lissage spectral rapide
  integer, parameter, public :: PRECOND_SSOR   = 2
  ! Factorisation incomplète Incomplete LU sans remplissage (ILU(0))
  integer, parameter, public :: PRECOND_ILU0   = 3

  ! --- 7. Schémas de Discrétisation Spatiale des Termes Convectifs ---
  ! Schéma décentré amont d'ordre 1 (Upwind) inconditionnellement stable mais diffusif
  integer, parameter, public :: CONV_UPWIND       = 1
  ! Schéma centré d'ordre 2 (Centered) sans dissipation artificielle
  integer, parameter, public :: CONV_CENTERED     = 2
  ! Schéma TVD Superbee à haute résolution avec limiteur de flux non linéaire (Roe 1986)
  integer, parameter, public :: CONV_TVD_SUPERBEE = 3

  ! --- 8. Modèles de Viscosité Dynamique du Fluide ---
  ! Fluide newtonien homogène à viscosité dynamique constante mu = mu0
  integer, parameter, public :: VISC_CONSTANT    = 1
  ! Fluide à viscosité variable spatiale exponentielle mu(y) = mu0 * exp(a * (y-e)^2 / e^2)
  integer, parameter, public :: VISC_EXPONENTIAL = 2

  ! --- 9. Types de Configurations Initiales et Cas de Vérification ---
  ! Écoulement monophasique uniforme avec fluide au repos (u=0, v=0, C=1)
  integer, parameter, public :: IC_UNIFORM          = 1
  ! Profil parabolique d'entrée de Poiseuille établi
  integer, parameter, public :: IC_POISEUILLE       = 2
  ! Tâche circulaire d'interface transportée dans un vortex cisaillé réversible
  integer, parameter, public :: IC_VORTEX_SHEAR     = 3
  ! Instabilité diphasique gravitationnelle de Rayleigh-Taylor avec interface perturbée
  integer, parameter, public :: IC_RAYLEIGH_TAYLOR  = 4
  ! Gouttelette liquide statique immergée pour la vérification de la loi de Laplace
  integer, parameter, public :: IC_DROPLET          = 5
  ! Bulle de gaz ascendante déformable dans un liquide visqueux (Benchmark de Hysing)
  integer, parameter, public :: IC_BUBBLE           = 6
  ! Effondrement d'une colonne d'eau sous gravité (Rupture de barrage de Martin & Moyce)
  integer, parameter, public :: IC_DAM_BREAK        = 7
  ! Test analytique d'ordre de convergence de l'opérateur Laplacien discret à 5 points
  integer, parameter, public :: IC_LAPLACIAN_TEST   = 8
  ! Nappe liquide cisaillée périodique (Cas 3.2 de l'article Fluids 2021)
  integer, parameter, public :: IC_LIQUID_SHEET     = 9
  ! Jet liquide visqueux injecté / flambement (Cas 3.3 de l'article Fluids 2021)
  integer, parameter, public :: IC_VISCOUS_JET      = 10

  ! --- 10. Formes Géométriques d'Obstacles Immergés (Brinkman) ---
  ! Aucun obstacle immergé dans le domaine fluide
  integer, parameter, public :: OBS_NONE     = 0
  ! Obstacle rectangulaire défini par ses coordonnées minimales et maximales (x0, y0, x1, y1)
  integer, parameter, public :: OBS_BOX      = 1
  ! Obstacle circulaire cylindrique défini par son centre (xc, yc) et son rayon R
  integer, parameter, public :: OBS_CYLINDER = 2

end module module_constants

\end{lstlisting}

\section{Code Source Intégral Annoté : \texttt{src/covof2d\_main.f90}}
Le listing ci-dessous présente l'intégralité des 105 lignes du fichier \texttt{src/covof2d\_main.f90} de l'implémentation Fortran 2018 :

\begin{lstlisting}[language={[95]Fortran}]
!===============================================================================
! Projet        : COVOF 2D (Convection & Volume-Of-Fluid 2D CFD Solver)
! Fichier       : covof2d_main.f90 (Point d'entrée principal Fortran 2018 / MPI)
! Description   : Programme principal pilotant l'initialisation de l'environnement
!                 OpenMPI, le traitement des arguments CLI, la lecture du cas
!                 JSON, l'orchestration du calcul Navier-Stokes / VOF / Énergie,
!                 et la libération finale des ressources système.
! Auteur        : Elhadji Ibrahima Thiam
! Standard      : Fortran 2018 / MPI-3
!===============================================================================
program covof2d_main
  ! Importation du module orchestrateur de la simulation CFD
  use module_simulation, only: simulation_t
  ! Typage explicite obligatoire pour toutes les variables locales
  implicit none
  ! Inclusion des définitions de constantes et prototypes MPI Fortran
  include 'mpif.h'

  ! --- Déclarations des Variables Locales ---
  integer            :: rank       ! Rang unique du processus courant (0 .. size-1)
  integer            :: size       ! Nombre total de processus alloués au calcul
  integer            :: ierr       ! Code d'erreur retourné par les appels MPI
  integer            :: num_args   ! Nombre d'arguments passés en ligne de commande
  integer            :: stat       ! Statut de lecture de l'argument CLI
  character(len=256) :: config_file! Chemin vers le fichier de configuration JSON
  character(len=256) :: arg        ! Tampon de lecture d'un argument textuel
  type(simulation_t) :: sim        ! Instance du type dérivé orchestrateur
  logical            :: success    ! Indicateur booléen de réussite de l'initialisation

  ! 1. Initialisation du runtime parallèle OpenMPI
  call MPI_Init(ierr)

  ! Récupération du rang du processus dans le communicateur global MPI_COMM_WORLD
  call MPI_Comm_rank(MPI_COMM_WORLD, rank, ierr)

  ! Récupération de la taille du groupe de processus MPI
  call MPI_Comm_size(MPI_COMM_WORLD, size, ierr)

  ! Fichier de configuration par défaut si aucun argument n'est fourni
  config_file = "cases/cavity.json"

  ! Lecture du nombre d'arguments passés au binaire lors de l'appel mpirun
  num_args = command_argument_count()

  ! 2. Analyse des arguments passés en ligne de commande
  if (num_args >= 1) then
     ! Récupération du premier argument textuel
     call get_command_argument(1, arg, status=stat)

     ! Détection de la demande d'affichage d'aide (-h ou --help)
     if (trim(arg) == "-h" .or. trim(arg) == "--help") then
        ! Seul le processus maître (rang 0) écrit sur la sortie standard
        if (rank == 0) then
           write(*, '(A)') "================================================================================"
           write(*, '(A)') "   COVOF 2D — Solveur CFD Navier-Stokes Incompressible Parallèle (Fortran 2018)"
           write(*, '(A)') "================================================================================"
           write(*, '(A)') "Usage :"
           write(*, '(A)') "  mpirun -np <N> ./bin/covof2d_fortran [fichier_configuration.json]"
           write(*, '(A)') ""
           write(*, '(A)') "Exemples d'exécution :"
           write(*, '(A)') "  mpirun -np 2 ./bin/covof2d_fortran cases/cas3_1_laplacian.json"
           write(*, '(A)') "  mpirun -np 2 ./bin/covof2d_fortran cases/cas4_5_cavity.json"
           write(*, '(A)') "  mpirun -np 4 ./bin/covof2d_fortran cases/cas7_5_rising_bubble.json"
           write(*, '(A)') "  mpirun -np 2 ./bin/covof2d_fortran cases/cas6_1_natural_convection.json"
           write(*, '(A)') "================================================================================"
        end if
        ! Clôture propre du communicateur MPI avant terminaison
        call MPI_Finalize(ierr)
        stop
     end if

     ! Affectation du chemin de fichier spécifié
     config_file = trim(arg)
  end if

  ! Affichage du récapitulatif de lancement par le rang racine
  if (rank == 0) then
     write(*, '(A)') "================================================================================"
     write(*, '(A)') "   COVOF 2D — Solveur CFD Navier-Stokes Incompressible Parallèle (Fortran 2018)"
     write(*, '(A)') "================================================================================"
     write(*, '(A,I0)')       "  Nombre de processus MPI  : ", size
     write(*, '(A,A)')        "  Fichier de cas JSON      : ", trim(config_file)
     write(*, '(A)') "================================================================================"
  end if

  ! 3. Initialisation de la simulation CFD (lecture JSON, partitionnement MAC, allocations)
  success = sim%init(MPI_COMM_WORLD, trim(config_file))
  if (.not. success) then
     ! En cas d'erreur de parsing ou d'incohérence des données physiques
     if (rank == 0) then
        write(*, '(A)') "[ERREUR CRITIQUE] Échec de l'initialisation de COVOF 2D Fortran."
        write(*, '(A)') "                  Vérifiez le fichier de configuration : " // trim(config_file)
     end if
     ! Interruption immédiate et coordonnée de tous les rangs MPI
     call MPI_Abort(MPI_COMM_WORLD, 1, ierr)
     stop 1
  end if

  ! 4. Exécution de la boucle temporelle principale de simulation
  call sim%run()

  ! 5. Libération des ressources MPI et terminaison normale
  call MPI_Finalize(ierr)

end program covof2d_main

\end{lstlisting}

\chapter{Analyseur Syntaxique JSON Natif en Fortran 2018 et Module de Configuration (module_json.f90, module_config.f90)}
\section{Architecture du Parseur JSON Pur Fortran 2018 (\texttt{module\_json.f90})}
Contrairement à Python ou C++, l'écosystème Fortran standard ne dispose pas nativement d'un parseur JSON dans la bibliothèque système. Le module \texttt{module\_json.f90} (449 lignes) implémente intégralement en Fortran 2018 standard un analyseur lexical et syntaxique capable de lire exactement les mêmes fichiers \texttt{cases/*.json} que l'exécutable C++17 :
\begin{enumerate}
  \item \texttt{json\_read\_file(filename, content, ierr)} : charge le fichier JSON complet dans une chaîne allocatable \texttt{character(len=:), allocatable :: content}.
  \item \texttt{json\_extract\_block(json\_str, block\_name, block\_content, found)} : localise un objet JSON nommé (tel que \texttt{"mesh"} ou \texttt{"vof"}) en suivant l'équilibrage des accolades \texttt{\{} et \texttt{\}} hors chaînes de caractères.
  \item \texttt{json\_get\_real}, \texttt{json\_get\_int}, \texttt{json\_get\_logical}, \texttt{json\_get\_string} : extraient les valeurs typées avec repli automatique sur la valeur par défaut si la clé est absente.
  \item \texttt{json\_parse\_obstacles} : parcourt le tableau JSON \texttt{"obstacles": [ \{...\}, \{...\} ]} pour instancier les masques de pénalisation de Brinkman (\texttt{box} et \texttt{circle}).
\end{enumerate}

\section{Structure de Configuration \texttt{simulation\_config\_t} (\texttt{module\_config.f90})}
Le module \texttt{module\_config.f90} (493 lignes) traduit les chaînes JSON en constantes entières symboliques (\texttt{BC\_WALL}, \texttt{BC\_PERIODIC}, \texttt{BC\_JET\_INLET}, \texttt{SOLVER\_BICGSTAB2}, \texttt{RECON\_PLIC}, \texttt{ADVECT\_WEYMOUTH\_YUE}, etc.) et vérifie la cohérence des paramètres physiques et numériques avant l'allocation du maillage.

\section{Répertoire des Types Dérivés, Sous-Routines et Fonctions de \texttt{src/module\_json.f90}}
\begin{center}\small
\begin{tabular}{clll}
\toprule
\textbf{Ligne} & \textbf{Nature} & \textbf{Identifiant Fortran 2018} & \textbf{Interface / Arguments} \\
\midrule
2 & module & \texttt{module\_json} & Unité modulaire Fortran 2018 \\
18 & subroutine & \texttt{json\_read\_file} & Arguments : (filename, buffer, ierr) \\
79 & subroutine & \texttt{json\_find\_value\_str} & Arguments : (buffer, path, val\_str, found) \\
254 & subroutine & \texttt{json\_get\_string} & Arguments : (buffer, path, val, default\_val) \\
278 & subroutine & \texttt{json\_get\_real} & Arguments : (buffer, path, val, default\_val) \\
308 & subroutine & \texttt{json\_get\_int} & Arguments : (buffer, path, val, default\_val) \\
338 & subroutine & \texttt{json\_get\_bool} & Arguments : (buffer, path, val, default\_val) \\
371 & subroutine & \texttt{json\_get\_obstacle} & Arguments : (buffer, obs\_type, x0, y0, x1, y1, xc, y...) \\
\bottomrule
\end{tabular}
\end{center}

\section{Code Source Intégral Annoté : \texttt{src/module\_json.f90}}
Le listing ci-dessous présente l'intégralité des 449 lignes du fichier \texttt{src/module\_json.f90} de l'implémentation Fortran 2018 :

\begin{lstlisting}[language={[95]Fortran}]
! [Procédure/Module] module module_json
module module_json
  use module_precision, only: dp
  use module_constants, only: OBS_NONE, OBS_BOX, OBS_CYLINDER
  implicit none
  private

  public :: json_read_file
  public :: json_get_string
  public :: json_get_real
  public :: json_get_int
  public :: json_get_bool
  public :: json_get_obstacle

contains

  ! Lit un fichier complet dans une chaîne de caractères allouable en ignorant les commentaires //
  subroutine json_read_file(filename, buffer, ierr)
    ! Déclaration de variable locale / paramètre de calcul
    character(len=*), intent(in) :: filename
    ! Déclaration de variable locale / paramètre de calcul
    character(len=:), allocatable, intent(out) :: buffer
    ! Déclaration de variable locale / paramètre de calcul
    integer, intent(out) :: ierr

    ! Déclaration de variable locale / paramètre de calcul
    integer :: u, io
    ! Déclaration de variable locale / paramètre de calcul
    character(len=4096) :: line
    ! Déclaration de variable locale / paramètre de calcul
    character(len=:), allocatable :: temp
    ! Déclaration de variable locale / paramètre de calcul
    integer :: comment_pos

    ! Évaluation de la norme du résidu ou de l'erreur numérique
    ierr = 0
    ! Affectation et calcul numérique
    buffer = ""
    ! Affectation et calcul numérique
    open(newunit=u, file=trim(filename), status='old', action='read', iostat=io)
    ! Évaluation de la condition logique
    if (io /= 0) then
       ! Évaluation de la norme du résidu ou de l'erreur numérique
       ierr = io
       ! Retour à la procédure appelante
       return
    end if

    do
       ! Affectation et calcul numérique
       read(u, '(A)', iostat=io) line
       ! Affectation et calcul numérique
       if (io /= 0) exit
       ! Affectation et calcul numérique
       comment_pos = index(line, '//')
       ! Affectation et calcul numérique
       if (comment_pos > 0) line = line(1:comment_pos-1)
       ! Affectation et calcul numérique
       line = adjustl(line)
       ! Évaluation de la condition logique
       if (len_trim(line) > 0) then
          ! Évaluation de la condition logique
          if (len(buffer) == 0) then
             ! Affectation et calcul numérique
             buffer = trim(line)
          ! Branche par défaut si condition non satisfaite
          else
             ! Affectation et calcul numérique
             temp = buffer // " " // trim(line)
             ! Affectation et calcul numérique
             buffer = temp
          end if
       end if
    end do
    close(u)
  end subroutine json_read_file

  ! Recherche d'un chemin séparé par des points dans le buffer JSON
  subroutine json_find_value_str(buffer, path, val_str, found)
    ! Déclaration de variable locale / paramètre de calcul
    character(len=*), intent(in) :: buffer
    ! Déclaration de variable locale / paramètre de calcul
    character(len=*), intent(in) :: path
    ! Déclaration de variable locale / paramètre de calcul
    character(len=:), allocatable, intent(out) :: val_str
    ! Déclaration de variable locale / paramètre de calcul
    logical, intent(out) :: found

    ! Déclaration de variable locale / paramètre de calcul
    character(len=:), allocatable :: cur_buf
    ! Déclaration de variable locale / paramètre de calcul
    integer :: p_dot, pos_key, pos_colon
    ! Déclaration de variable locale / paramètre de calcul
    character(len=64) :: key, sub_path
    ! Déclaration de variable locale / paramètre de calcul
    character(len=70) :: search_key
    ! Déclaration de variable locale / paramètre de calcul
    integer :: brace_count, start_pos, end_pos

    ! Affectation et calcul numérique
    found = .false.
    ! Affectation et calcul numérique
    val_str = ""
    ! Affectation et calcul numérique
    cur_buf = buffer

    ! Affectation et calcul numérique
    sub_path = adjustl(trim(path))

    ! Boucle d'itération numérique
    do while (len_trim(sub_path) > 0)
       ! Affectation et calcul numérique
       p_dot = index(sub_path, ".")
       ! Évaluation de la condition logique
       if (p_dot > 0) then
          ! Affectation et calcul numérique
          key = sub_path(1:p_dot-1)
          ! Affectation et calcul numérique
          sub_path = adjustl(sub_path(p_dot+1:))
       ! Branche par défaut si condition non satisfaite
       else
          ! Affectation et calcul numérique
          key = trim(sub_path)
          ! Affectation et calcul numérique
          sub_path = ""
       end if

       ! Affectation et calcul numérique
       search_key = '"' // trim(key) // '"'
       ! Affectation et calcul numérique
       pos_key = index(cur_buf, trim(search_key))
       ! Évaluation de la condition logique
       if (pos_key <= 0) then
          ! Affectation et calcul numérique
          found = .false.
          ! Retour à la procédure appelante
          return
       end if

       ! Trouver les deux-points ':' après la clé
       pos_colon = pos_key + len_trim(search_key)
       ! Boucle d'itération numérique
       do while (pos_colon <= len(cur_buf) .and. cur_buf(pos_colon:pos_colon) /= ':')
          ! Affectation et calcul numérique
          pos_colon = pos_colon + 1
       end do
       ! Évaluation de la condition logique
       if (pos_colon > len(cur_buf)) then
          ! Affectation et calcul numérique
          found = .false.
          ! Retour à la procédure appelante
          return
       end if

       ! Si ce n'est pas le dernier segment de chemin, chercher l'objet { ... } ou [ ... ] correspondant
       if (len_trim(sub_path) > 0) then
          ! Affectation et calcul numérique
          start_pos = pos_colon + 1
          ! Boucle d'itération numérique
          do while (start_pos <= len(cur_buf) .and. cur_buf(start_pos:start_pos) /= '{' .and. cur_buf(start_pos:start_pos) /= '[')
             ! Affectation et calcul numérique
             start_pos = start_pos + 1
          end do
          ! Évaluation de la condition logique
          if (start_pos > len(cur_buf)) then
             ! Affectation et calcul numérique
             found = .false.
             ! Retour à la procédure appelante
             return
          end if

          ! Affectation et calcul numérique
          brace_count = 1
          ! Affectation et calcul numérique
          end_pos = start_pos + 1
          ! Boucle d'itération numérique
          do while (end_pos <= len(cur_buf) .and. brace_count > 0)
             ! Affectation et calcul numérique
             if (cur_buf(end_pos:end_pos) == '{' .or. cur_buf(end_pos:end_pos) == '[') brace_count = brace_count + 1
             ! Affectation et calcul numérique
             if (cur_buf(end_pos:end_pos) == '}' .or. cur_buf(end_pos:end_pos) == ']') brace_count = brace_count - 1
             ! Affectation et calcul numérique
             end_pos = end_pos + 1
          end do
          ! Affectation et calcul numérique
          cur_buf = cur_buf(start_pos:end_pos-1)
       ! Branche par défaut si condition non satisfaite
       else
          ! Dernier élément du chemin : extraire la valeur
          start_pos = pos_colon + 1
          ! Boucle d'itération numérique
          do while (start_pos <= len(cur_buf) .and. (cur_buf(start_pos:start_pos) == ' ' .or. cur_buf(start_pos:start_pos) == char(9)))
             ! Affectation et calcul numérique
             start_pos = start_pos + 1
          end do
          ! Évaluation de la condition logique
          if (start_pos > len(cur_buf)) then
             ! Affectation et calcul numérique
             found = .false.
             ! Retour à la procédure appelante
             return
          end if

          ! Évaluation de la condition logique
          if (cur_buf(start_pos:start_pos) == '"') then
             ! Chaîne entre guillemets
             end_pos = start_pos + 1
             ! Boucle d'itération numérique
             do while (end_pos <= len(cur_buf) .and. cur_buf(end_pos:end_pos) /= '"')
                ! Affectation et calcul numérique
                end_pos = end_pos + 1
             end do
             ! Affectation et calcul numérique
             val_str = cur_buf(start_pos+1:end_pos-1)
             ! Affectation et calcul numérique
             found = .true.
             ! Retour à la procédure appelante
             return
          ! Branche par défaut si condition non satisfaite
          else
             ! Nombre, booléen, etc.
             end_pos = start_pos
             ! Boucle d'itération numérique
             do while (end_pos <= len(cur_buf) .and. &
                       ! Affectation et calcul numérique
                       cur_buf(end_pos:end_pos) /= ',' .and. &
                       ! Affectation et calcul numérique
                       cur_buf(end_pos:end_pos) /= '}' .and. &
                       ! Affectation et calcul numérique
                       cur_buf(end_pos:end_pos) /= ']' .and. &
                       ! Affectation et calcul numérique
                       cur_buf(end_pos:end_pos) /= ' ' .and. &
                       ! Affectation et calcul numérique
                       cur_buf(end_pos:end_pos) /= char(9) .and. &
                       ! Affectation et calcul numérique
                       cur_buf(end_pos:end_pos) /= char(10) .and. &
                       ! Affectation et calcul numérique
                       cur_buf(end_pos:end_pos) /= char(13))
                ! Affectation et calcul numérique
                end_pos = end_pos + 1
             end do
             ! Affectation et calcul numérique
             val_str = trim(adjustl(cur_buf(start_pos:end_pos-1)))
             ! Affectation et calcul numérique
             found = .true.
             ! Retour à la procédure appelante
             return
          end if
       end if
    end do
  end subroutine json_find_value_str

  ! [Procédure/Module] subroutine json_get_string
  subroutine json_get_string(buffer, path, val, default_val)
    ! Déclaration de variable locale / paramètre de calcul
    character(len=*), intent(in) :: buffer, path, default_val
    ! Déclaration de variable locale / paramètre de calcul
    character(len=*), intent(out) :: val
    ! Déclaration de variable locale / paramètre de calcul
    character(len=:), allocatable :: s
    ! Déclaration de variable locale / paramètre de calcul
    logical :: found

    ! Appel de sous-programme modulaire
    call json_find_value_str(buffer, path, s, found)
    ! Évaluation de la condition logique
    if (found .and. allocated(s)) then
       ! Affectation et calcul numérique
       val = trim(s)
    ! Branche par défaut si condition non satisfaite
    else
       ! Affectation et calcul numérique
       val = trim(default_val)
    end if
  end subroutine json_get_string

  ! [Procédure/Module] subroutine json_get_real
  subroutine json_get_real(buffer, path, val, default_val)
    ! Déclaration de variable locale / paramètre de calcul
    character(len=*), intent(in) :: buffer, path
    ! Déclaration de variable locale / paramètre de calcul
    real(dp), intent(out) :: val
    ! Déclaration de variable locale / paramètre de calcul
    real(dp), intent(in) :: default_val
    ! Déclaration de variable locale / paramètre de calcul
    character(len=:), allocatable :: s
    ! Déclaration de variable locale / paramètre de calcul
    logical :: found
    ! Déclaration de variable locale / paramètre de calcul
    integer :: io

    ! Appel de sous-programme modulaire
    call json_find_value_str(buffer, path, s, found)
    ! Évaluation de la condition logique
    if (found .and. allocated(s)) then
       ! Affectation et calcul numérique
       read(s, *, iostat=io) val
       ! Affectation et calcul numérique
       if (io /= 0) val = default_val
    ! Branche par défaut si condition non satisfaite
    else
       ! Affectation et calcul numérique
       val = default_val
    end if
  end subroutine json_get_real

  ! [Procédure/Module] subroutine json_get_int
  subroutine json_get_int(buffer, path, val, default_val)
    ! Déclaration de variable locale / paramètre de calcul
    character(len=*), intent(in) :: buffer, path
    ! Déclaration de variable locale / paramètre de calcul
    integer, intent(out) :: val
    ! Déclaration de variable locale / paramètre de calcul
    integer, intent(in) :: default_val
    ! Déclaration de variable locale / paramètre de calcul
    character(len=:), allocatable :: s
    ! Déclaration de variable locale / paramètre de calcul
    logical :: found
    ! Déclaration de variable locale / paramètre de calcul
    integer :: io

    ! Appel de sous-programme modulaire
    call json_find_value_str(buffer, path, s, found)
    ! Évaluation de la condition logique
    if (found .and. allocated(s)) then
       ! Affectation et calcul numérique
       read(s, *, iostat=io) val
       ! Affectation et calcul numérique
       if (io /= 0) val = default_val
    ! Branche par défaut si condition non satisfaite
    else
       ! Affectation et calcul numérique
       val = default_val
    end if
  end subroutine json_get_int

  ! [Procédure/Module] subroutine json_get_bool
  subroutine json_get_bool(buffer, path, val, default_val)
    ! Déclaration de variable locale / paramètre de calcul
    character(len=*), intent(in) :: buffer, path
    ! Déclaration de variable locale / paramètre de calcul
    logical, intent(out) :: val
    ! Déclaration de variable locale / paramètre de calcul
    logical, intent(in) :: default_val
    ! Déclaration de variable locale / paramètre de calcul
    character(len=:), allocatable :: s
    ! Déclaration de variable locale / paramètre de calcul
    logical :: found

    ! Appel de sous-programme modulaire
    call json_find_value_str(buffer, path, s, found)
    ! Évaluation de la condition logique
    if (found .and. allocated(s)) then
       ! Évaluation de la condition logique
       if (index(s, "true") > 0 .or. index(s, "TRUE") > 0 .or. index(s, "1") > 0) then
          ! Affectation et calcul numérique
          val = .true.
       ! Branche par défaut si condition non satisfaite
       else
          ! Affectation et calcul numérique
          val = .false.
       end if
    ! Branche par défaut si condition non satisfaite
    else
       ! Affectation et calcul numérique
       val = default_val
    end if
  end subroutine json_get_bool

  ! Détection et extraction des obstacles (BOX ou CYLINDER)
  subroutine json_get_obstacle(buffer, obs_type, x0, y0, x1, y1, xc, yc, radius, perm, found)
    ! Déclaration de variable locale / paramètre de calcul
    character(len=*), intent(in) :: buffer
    ! Déclaration de variable locale / paramètre de calcul
    integer, intent(out) :: obs_type
    ! Déclaration de variable locale / paramètre de calcul
    real(dp), intent(out) :: x0, y0, x1, y1, xc, yc, radius, perm
    ! Déclaration de variable locale / paramètre de calcul
    logical, intent(out) :: found

    ! Déclaration de variable locale / paramètre de calcul
    integer :: pos_obs, pos_end
    ! Déclaration de variable locale / paramètre de calcul
    character(len=:), allocatable :: obs_str
    ! Déclaration de variable locale / paramètre de calcul
    character(len=32) :: type_str

    ! Affectation et calcul numérique
    obs_type = OBS_NONE
    ! Affectation et calcul numérique
    x0 = 0.0_dp; y0 = 0.0_dp; x1 = 0.0_dp; y1 = 0.0_dp
    ! Affectation et calcul numérique
    xc = 0.0_dp; yc = 0.0_dp; radius = 0.0_dp; perm = 1.0e-40_dp
    ! Affectation et calcul numérique
    found = .false.

    ! Affectation et calcul numérique
    pos_obs = index(buffer, '"obstacles"')
    ! Affectation et calcul numérique
    if (pos_obs <= 0) return

    ! Affectation et calcul numérique
    pos_end = pos_obs + 11
    ! Boucle d'itération numérique
    do while (pos_end <= len(buffer) .and. buffer(pos_end:pos_end) /= ']')
       ! Affectation et calcul numérique
       pos_end = pos_end + 1
    end do
    if (pos_end > len(buffer)) return

    ! Affectation et calcul numérique
    obs_str = buffer(pos_obs:pos_end)
    ! Appel de sous-programme modulaire
    call json_get_string(obs_str, "type", type_str, "")

    ! Évaluation de la condition logique
    if (trim(type_str) == "BOX") then
       ! Affectation et calcul numérique
       obs_type = OBS_BOX
       ! Appel de sous-programme modulaire
       call json_get_real(obs_str, "x0", x0, 0.0_dp)
       ! Appel de sous-programme modulaire
       call json_get_real(obs_str, "y0", y0, 0.0_dp)
       ! Appel de sous-programme modulaire
       call json_get_real(obs_str, "x1", x1, 0.0_dp)
       ! Appel de sous-programme modulaire
       call json_get_real(obs_str, "y1", y1, 0.0_dp)
       ! Appel de sous-programme modulaire
       call json_get_real(obs_str, "permeability", perm, 1.0e-40_dp)
       ! Affectation et calcul numérique
       found = .true.
    ! Branche conditionnelle alternative
    else if (trim(type_str) == "CYLINDER") then
       ! Affectation et calcul numérique
       obs_type = OBS_CYLINDER
       ! Appel de sous-programme modulaire
       call json_get_real(obs_str, "xc", xc, 0.0_dp)
       ! Appel de sous-programme modulaire
       call json_get_real(obs_str, "yc", yc, 0.0_dp)
       ! Appel de sous-programme modulaire
       call json_get_real(obs_str, "radius", radius, 0.0_dp)
       ! Appel de sous-programme modulaire
       call json_get_real(obs_str, "permeability", perm, 1.0e-40_dp)
       ! Affectation et calcul numérique
       found = .true.
    end if
  end subroutine json_get_obstacle

end module module_json

\end{lstlisting}

\section{Répertoire des Types Dérivés, Sous-Routines et Fonctions de \texttt{src/module\_config.f90}}
\begin{center}\small
\begin{tabular}{clll}
\toprule
\textbf{Ligne} & \textbf{Nature} & \textbf{Identifiant Fortran 2018} & \textbf{Interface / Arguments} \\
\midrule
2 & module & \texttt{module\_config} & Unité modulaire Fortran 2018 \\
11 & type dérivé & \texttt{bc\_t} & Structure de données Fortran 2018 \\
30 & type dérivé & \texttt{config\_t} & Structure de données Fortran 2018 \\
146 & subroutine & \texttt{parse\_bc} & Arguments : (buf, prefix, bc) \\
213 & subroutine & \texttt{load\_config\_from\_json} & Arguments : (filename, cfg, ierr) \\
\bottomrule
\end{tabular}
\end{center}

\section{Code Source Intégral Annoté : \texttt{src/module\_config.f90}}
Le listing ci-dessous présente l'intégralité des 507 lignes du fichier \texttt{src/module\_config.f90} de l'implémentation Fortran 2018 :

\begin{lstlisting}[language={[95]Fortran}]
! [Procédure/Module] module module_config
module module_config
  use module_precision, only: dp
  use module_constants
  use module_json
  implicit none
  private

  public :: bc_t, config_t, load_config_from_json

  type :: bc_t
     ! Déclaration de variable locale / paramètre de calcul
     integer  :: btype = BC_DIRICHLET
     ! Déclaration de variable locale / paramètre de calcul
     real(dp) :: u = 0.0_dp
     ! Déclaration de variable locale / paramètre de calcul
     real(dp) :: v = 0.0_dp
     ! Profil spécifique pour la frontière ("UNIFORM", "PARABOLIC", "JET_INLET")
     character(len=32) :: profile = "UNIFORM"
     ! Largeur de l'injecteur pour profil JET_INLET [m] (Cas 3.3 Fluids 2021)
     real(dp) :: jet_width = 0.08_dp
     ! Position centrale de l'injecteur selon X [m]
     real(dp) :: jet_center = 0.5_dp
     ! Déclaration de variable locale / paramètre de calcul
     integer  :: thermal_type = THERMAL_NEUMANN
     ! Déclaration de variable locale / paramètre de calcul
     real(dp) :: T = 0.0_dp
  end type bc_t

  type :: config_t
     ! Déclaration de variable locale / paramètre de calcul
     character(len=64) :: case_name = "covof2d_case"
     ! Déclaration de variable locale / paramètre de calcul
     character(len=32) :: mode = "NAVIER_STOKES"

     ! Mesh parameters
     integer  :: nx = 64, ny = 64
     ! Déclaration de variable locale / paramètre de calcul
     real(dp) :: lx = 1.0_dp, ly = 1.0_dp

     ! Physical models & flags
     logical  :: enable_two_phase = .false.
     ! Déclaration de variable locale / paramètre de calcul
     logical  :: enable_surface_tension = .false.
     ! Déclaration de variable locale / paramètre de calcul
     logical  :: enable_convection = .true.
     ! Déclaration de variable locale / paramètre de calcul
     integer  :: convective_scheme = CONV_UPWIND
     ! Déclaration de variable locale / paramètre de calcul
     integer  :: viscosity_model = VISC_CONSTANT
     ! Déclaration de variable locale / paramètre de calcul
     real(dp) :: visc_mu0 = 1.0e-3_dp, visc_a = 2.0_dp, visc_e = 0.01_dp
     ! Déclaration de variable locale / paramètre de calcul
     real(dp) :: sigma = 0.0_dp
     ! Déclaration de variable locale / paramètre de calcul
     real(dp) :: gx = 0.0_dp, gy = 0.0_dp
     ! Déclaration de variable locale / paramètre de calcul
     real(dp) :: rho1 = 1.0_dp, mu1 = 1.0e-3_dp
     ! Déclaration de variable locale / paramètre de calcul
     real(dp) :: rho2 = 1.0_dp, mu2 = 1.0e-3_dp

     ! VOF and interface tracking
     character(len=32) :: reconstruction_method = "PLIC"
     ! Déclaration de variable locale / paramètre de calcul
     character(len=32) :: advection_method = "VOF_TVD"
     ! Déclaration de variable locale / paramètre de calcul
     character(len=32) :: tvd_limiter = "SUPERBEE"
     ! Déclaration de variable locale / paramètre de calcul
     real(dp) :: cfl_vof = 0.4_dp

     ! Thermal parameters
     logical  :: enable_energy = .false.
     ! Déclaration de variable locale / paramètre de calcul
     real(dp) :: Ra = 1.0e4_dp, Pr = 0.71_dp, T_ref = 0.0_dp, beta = 3.41e-3_dp
     ! Déclaration de variable locale / paramètre de calcul
     real(dp) :: Cp = 1.0_dp, k_thermal = 0.014_dp

     ! Time integration
     real(dp) :: dt = 0.01_dp, t_final = 1.0_dp
     ! Déclaration de variable locale / paramètre de calcul
     logical  :: adaptive_dt = .false.
     ! Déclaration de variable locale / paramètre de calcul
     real(dp) :: cfl_max = 0.5_dp

     ! Boundary conditions
     type(bc_t) :: bc_left, bc_right, bc_bottom, bc_top

     ! Immersed boundary obstacle
     logical  :: obs_enabled = .false.
     ! Déclaration de variable locale / paramètre de calcul
     integer  :: obs_type = OBS_NONE
     ! Déclaration de variable locale / paramètre de calcul
     real(dp) :: obs_x0 = 0.0_dp, obs_y0 = 0.0_dp, obs_x1 = 0.0_dp, obs_y1 = 0.0_dp
     ! Déclaration de variable locale / paramètre de calcul
     real(dp) :: obs_xc = 0.0_dp, obs_yc = 0.0_dp, obs_radius = 0.0_dp
     ! Déclaration de variable locale / paramètre de calcul
     real(dp) :: obs_perm = 1.0e-40_dp

     ! Initial conditions
     integer  :: ic_type = IC_UNIFORM
     ! Déclaration de variable locale / paramètre de calcul
     real(dp) :: ic_u = 0.0_dp, ic_v = 0.0_dp
     ! Déclaration de variable locale / paramètre de calcul
     real(dp) :: ic_center_x = 0.5_dp, ic_center_y = 0.5_dp, ic_radius = 0.25_dp
     ! Déclaration de variable locale / paramètre de calcul
     real(dp) :: ic_box_x0 = 0.0_dp, ic_box_y0 = 0.0_dp, ic_box_x1 = 0.5_dp, ic_box_y1 = 0.5_dp
     ! Déclaration de variable locale / paramètre de calcul
     real(dp) :: ic_pert_amp = 0.05_dp, ic_pert_k = 1.0_dp
     ! Déclaration de l'épaisseur initiale de nappe liquide (Cas 3.2 Fluids 2021)
     real(dp) :: ic_thickness = 0.0003_dp
     ! Déclaration de la largeur initiale de jet liquide (Cas 3.3 Fluids 2021)
     real(dp) :: ic_jet_width = 0.08_dp

     ! Solvers
     integer  :: solver_type = SOLVER_AUGMENTED_LAGRANGIAN
     ! Déclaration de variable locale / paramètre de calcul
     integer  :: linear_solver = LINEAR_BICGSTAB2
     ! Déclaration de variable locale / paramètre de calcul
     integer  :: precond_type = PRECOND_JACOBI
     ! Déclaration de variable locale / paramètre de calcul
     real(dp) :: tol = 1.0e-5_dp
     ! Déclaration de variable locale / paramètre de calcul
     integer  :: max_iter = 300
     ! Déclaration de variable locale / paramètre de calcul
     real(dp) :: r_al = 100.0_dp
     ! Déclaration de variable locale / paramètre de calcul
     integer  :: uzawa_iter = 1
     ! Déclaration de variable locale / paramètre de calcul
     logical  :: use_precond = .true.
     ! Déclaration de variable locale / paramètre de calcul
     real(dp) :: steady_state_tol = 1.0e-5_dp

     ! Outputs
     character(len=128) :: output_dir = "output"
     ! Déclaration de variable locale / paramètre de calcul
     integer  :: vtk_interval = 50
     ! Déclaration de variable locale / paramètre de calcul
     integer  :: log_interval = 10
     ! Déclaration de variable locale / paramètre de calcul
     logical  :: write_csv = .true.
  end type config_t

contains

  ! [Procédure/Module] subroutine parse_bc
  subroutine parse_bc(buf, prefix, bc)
    ! Déclaration de variable locale / paramètre de calcul
    character(len=*), intent(in) :: buf, prefix
    ! Déclaration de variable locale / paramètre de calcul
    type(bc_t), intent(out) :: bc
    ! Déclaration de variable locale / paramètre de calcul
    character(len=32) :: btype_str, ttype_str

    ! Appel de sous-programme modulaire
    call json_get_string(buf, trim(prefix) // ".type", btype_str, "DIRICHLET")
    ! Évaluation de la condition logique
    if (trim(btype_str) == "DIRICHLET") then
       ! Affectation et calcul numérique
       bc%btype = BC_DIRICHLET
    ! Branche conditionnelle alternative
    else if (trim(btype_str) == "NEUMANN") then
       ! Affectation et calcul numérique
       bc%btype = BC_NEUMANN
    ! Branche conditionnelle alternative
    else if (trim(btype_str) == "PERIODIC") then
       ! Affectation et calcul numérique
       bc%btype = BC_PERIODIC
    ! Branche conditionnelle alternative
    else if (trim(btype_str) == "FREE_SLIP") then
       ! Affectation et calcul numérique
       bc%btype = BC_FREESLIP
    ! Branche par défaut si condition non satisfaite
    else
       ! Affectation et calcul numérique
       bc%btype = BC_DIRICHLET
    end if

    ! Appel de sous-programme modulaire
    call json_get_real(buf, trim(prefix) // ".u", bc%u, 0.0_dp)
    ! Appel de sous-programme modulaire
    call json_get_real(buf, trim(prefix) // ".u_mean", bc%u, bc%u)
    ! Appel de sous-programme modulaire
    call json_get_real(buf, trim(prefix) // ".u_val", bc%u, bc%u)
    ! Appel de sous-programme modulaire
    call json_get_real(buf, trim(prefix) // ".v", bc%v, 0.0_dp)
    ! Appel de sous-programme modulaire
    call json_get_real(buf, trim(prefix) // ".v_val", bc%v, bc%v)
    ! Lecture du profil de frontière (Cas 3.3 Fluids 2021)
    call json_get_string(buf, trim(prefix) // ".profile", bc%profile, "UNIFORM")
    ! Lecture de la largeur de l'injecteur
    call json_get_real(buf, trim(prefix) // ".jet_width", bc%jet_width, 0.08_dp)
    ! Lecture du centre de l'injecteur
    call json_get_real(buf, trim(prefix) // ".jet_center", bc%jet_center, 0.5_dp)

    ! Appel de sous-programme modulaire
    call json_get_string(buf, trim(prefix) // ".T_type", ttype_str, "NEUMANN")
    ! Évaluation de la condition logique
    if (trim(ttype_str) == "DIRICHLET") then
       ! Affectation et calcul numérique
       bc%thermal_type = THERMAL_DIRICHLET
    ! Branche par défaut si condition non satisfaite
    else
       ! Affectation et calcul numérique
       bc%thermal_type = THERMAL_NEUMANN
    end if
    ! Appel de sous-programme modulaire
    call json_get_real(buf, trim(prefix) // ".T", bc%T, 0.0_dp)
    ! Appel de sous-programme modulaire
    call json_get_real(buf, trim(prefix) // ".T_val", bc%T, bc%T)
  end subroutine parse_bc

  ! [Procédure/Module] subroutine load_config_from_json
  subroutine load_config_from_json(filename, cfg, ierr)
    ! Déclaration de variable locale / paramètre de calcul
    character(len=*), intent(in) :: filename
    ! Déclaration de variable locale / paramètre de calcul
    type(config_t), intent(out) :: cfg
    ! Déclaration de variable locale / paramètre de calcul
    integer, intent(out), optional :: ierr

    ! Déclaration de variable locale / paramètre de calcul
    character(len=:), allocatable :: buf
    ! Déclaration de variable locale / paramètre de calcul
    character(len=64) :: str_val
    ! Déclaration de variable locale / paramètre de calcul
    integer :: stat

    ! Affectation et calcul numérique
    stat = 0
    ! Appel de sous-programme modulaire
    call json_read_file(filename, buf, stat)
    ! Évaluation de la norme du résidu ou de l'erreur numérique
    if (present(ierr)) ierr = stat
    ! Affectation et calcul numérique
    if (stat /= 0) return

    ! Case name & mode
    call json_get_string(buf, "case_name", cfg%case_name, "covof_case")
    ! Appel de sous-programme modulaire
    call json_get_string(buf, "output.case_name", cfg%case_name, cfg%case_name)
    ! Appel de sous-programme modulaire
    call json_get_string(buf, "mode", cfg%mode, "NAVIER_STOKES")

    ! Mesh
    call json_get_int(buf, "mesh.nx", cfg%nx, 64)
    ! Appel de sous-programme modulaire
    call json_get_int(buf, "mesh.ny", cfg%ny, 64)
    ! Appel de sous-programme modulaire
    call json_get_real(buf, "mesh.lx", cfg%lx, 1.0_dp)
    ! Appel de sous-programme modulaire
    call json_get_real(buf, "mesh.ly", cfg%ly, 1.0_dp)

    ! Physics
    call json_get_bool(buf, "physics.enable_two_phase", cfg%enable_two_phase, .false.)
    ! Appel de sous-programme modulaire
    call json_get_bool(buf, "physics.enable_surface_tension", cfg%enable_surface_tension, .false.)
    ! Appel de sous-programme modulaire
    call json_get_bool(buf, "physics.enable_convection", cfg%enable_convection, .true.)
    ! Appel de sous-programme modulaire
    call json_get_real(buf, "physics.sigma", cfg%sigma, 0.0_dp)
    ! Appel de sous-programme modulaire
    call json_get_real(buf, "physics.gx", cfg%gx, 0.0_dp)
    ! Appel de sous-programme modulaire
    call json_get_real(buf, "physics.gy", cfg%gy, 0.0_dp)

    ! Appel de sous-programme modulaire
    call json_get_string(buf, "physics.convective_scheme", str_val, "UPWIND")
    ! Évaluation de la condition logique
    if (trim(str_val) == "CENTERED") then
       ! Affectation et calcul numérique
       cfg%convective_scheme = CONV_CENTERED
    ! Branche conditionnelle alternative
    else if (trim(str_val) == "TVD_SUPERBEE") then
       ! Affectation et calcul numérique
       cfg%convective_scheme = CONV_TVD_SUPERBEE
    ! Branche par défaut si condition non satisfaite
    else
       ! Affectation et calcul numérique
       cfg%convective_scheme = CONV_UPWIND
    end if

    ! Appel de sous-programme modulaire
    call json_get_string(buf, "physics.viscosity_model", str_val, "CONSTANT")
    ! Évaluation de la condition logique
    if (trim(str_val) == "EXPONENTIAL") then
       ! Affectation et calcul numérique
       cfg%viscosity_model = VISC_EXPONENTIAL
       ! Appel de sous-programme modulaire
       call json_get_real(buf, "physics.visc_mu0", cfg%visc_mu0, 1.0e-3_dp)
       ! Appel de sous-programme modulaire
       call json_get_real(buf, "physics.visc_a", cfg%visc_a, 2.0_dp)
       ! Appel de sous-programme modulaire
       call json_get_real(buf, "physics.visc_e", cfg%visc_e, 0.01_dp)
    ! Branche par défaut si condition non satisfaite
    else
       ! Affectation et calcul numérique
       cfg%viscosity_model = VISC_CONSTANT
    end if

    ! Appel de sous-programme modulaire
    call json_get_real(buf, "physics.fluid1.rho", cfg%rho1, 1.0_dp)
    ! Appel de sous-programme modulaire
    call json_get_real(buf, "physics.fluid1.mu", cfg%mu1, 1.0e-3_dp)
    ! Appel de sous-programme modulaire
    call json_get_real(buf, "physics.fluid2.rho", cfg%rho2, cfg%rho1)
    ! Appel de sous-programme modulaire
    call json_get_real(buf, "physics.fluid2.mu", cfg%mu2, cfg%mu1)

    ! Thermal
    call json_get_bool(buf, "physics.enable_energy", cfg%enable_energy, .false.)
    ! Appel de sous-programme modulaire
    call json_get_bool(buf, "thermal.enable_energy", cfg%enable_energy, cfg%enable_energy)
    ! Appel de sous-programme modulaire
    call json_get_real(buf, "physics.Ra", cfg%Ra, 1.0e4_dp)
    ! Appel de sous-programme modulaire
    call json_get_real(buf, "physics.Pr", cfg%Pr, 0.71_dp)
    ! Appel de sous-programme modulaire
    call json_get_real(buf, "physics.T_ref", cfg%T_ref, 0.0_dp)
    ! Appel de sous-programme modulaire
    call json_get_real(buf, "physics.beta", cfg%beta, 3.41e-3_dp)
    ! Appel de sous-programme modulaire
    call json_get_real(buf, "physics.fluid1.Cp", cfg%Cp, 1.0_dp)
    ! Appel de sous-programme modulaire
    call json_get_real(buf, "physics.fluid1.k_thermal", cfg%k_thermal, 0.014_dp)

    ! Time
    call json_get_real(buf, "time.dt", cfg%dt, 0.01_dp)
    ! Appel de sous-programme modulaire
    call json_get_real(buf, "time.t_final", cfg%t_final, 1.0_dp)
    ! Appel de sous-programme modulaire
    call json_get_bool(buf, "time.adaptive_dt", cfg%adaptive_dt, .false.)
    ! Appel de sous-programme modulaire
    call json_get_real(buf, "time.cfl_max", cfg%cfl_max, 0.5_dp)
    ! Appel de sous-programme modulaire
    call json_get_real(buf, "steady_state_tol", cfg%steady_state_tol, 1.0e-5_dp)

    ! Boundary conditions
    call parse_bc(buf, "boundary_conditions.left", cfg%bc_left)
    ! Appel de sous-programme modulaire
    call parse_bc(buf, "boundary_conditions.right", cfg%bc_right)
    ! Appel de sous-programme modulaire
    call parse_bc(buf, "boundary_conditions.bottom", cfg%bc_bottom)
    ! Appel de sous-programme modulaire
    call parse_bc(buf, "boundary_conditions.top", cfg%bc_top)

    ! Obstacles
    call json_get_obstacle(buf, cfg%obs_type, cfg%obs_x0, cfg%obs_y0, cfg%obs_x1, cfg%obs_y1, &
                           cfg%obs_xc, cfg%obs_yc, cfg%obs_radius, cfg%obs_perm, cfg%obs_enabled)

    ! Initial condition
    call json_get_string(buf, "initial_condition.type", str_val, "UNIFORM")
    ! Évaluation de la condition logique
    if (trim(str_val) == "UNIFORM") then
       ! Affectation et calcul numérique
       cfg%ic_type = IC_UNIFORM
    ! Branche conditionnelle alternative
    else if (trim(str_val) == "POISEUILLE") then
       ! Affectation et calcul numérique
       cfg%ic_type = IC_POISEUILLE
    ! Branche conditionnelle alternative
    else if (trim(str_val) == "VORTEX_SHEAR") then
       ! Affectation et calcul numérique
       cfg%ic_type = IC_VORTEX_SHEAR
    ! Branche conditionnelle alternative
    else if (trim(str_val) == "RAYLEIGH_TAYLOR") then
       ! Affectation et calcul numérique
       cfg%ic_type = IC_RAYLEIGH_TAYLOR
    ! Branche conditionnelle alternative
    else if (trim(str_val) == "DROPLET") then
       ! Affectation et calcul numérique
       cfg%ic_type = IC_DROPLET
    ! Branche conditionnelle alternative pour bulle ou cylindre immergé
    else if (trim(str_val) == "BUBBLE" .or. trim(str_val) == "CYLINDER") then
       ! Affectation et calcul numérique
       cfg%ic_type = IC_BUBBLE
    ! Branche conditionnelle alternative
    else if (trim(str_val) == "DAM_BREAK") then
       ! Affectation et calcul numérique
       cfg%ic_type = IC_DAM_BREAK
    ! Branche conditionnelle alternative pour nappe liquide (Cas 3.2 Fluids 2021)
    else if (trim(str_val) == "LIQUID_SHEET") then
       ! Affectation et calcul numérique
       cfg%ic_type = IC_LIQUID_SHEET
    ! Branche conditionnelle alternative pour jet liquide visqueux (Cas 3.3 Fluids 2021)
    else if (trim(str_val) == "VISCOUS_JET" .or. trim(str_val) == "JET_BUCKLING") then
       ! Affectation et calcul numérique
       cfg%ic_type = IC_VISCOUS_JET
    ! Branche conditionnelle alternative
    else if (trim(str_val) == "LAPLACIAN_TEST") then
       ! Affectation et calcul numérique
       cfg%ic_type = IC_LAPLACIAN_TEST
    ! Branche conditionnelle alternative
    else if (trim(str_val) == "SINGLE_PHASE") then
       ! Affectation et calcul numérique
       cfg%ic_type = IC_UNIFORM
    end if

    ! Appel de sous-programme modulaire
    call json_get_real(buf, "initial_condition.u", cfg%ic_u, 0.0_dp)
    ! Appel de sous-programme modulaire
    call json_get_real(buf, "initial_condition.v", cfg%ic_v, 0.0_dp)
    ! Appel de sous-programme modulaire
    call json_get_real(buf, "initial_condition.center_x", cfg%ic_center_x, 0.5_dp)
    ! Appel de sous-programme modulaire
    call json_get_real(buf, "initial_condition.center_y", cfg%ic_center_y, 0.5_dp)
    ! Appel de sous-programme modulaire
    call json_get_real(buf, "initial_condition.radius", cfg%ic_radius, 0.25_dp)
    ! Appel de sous-programme modulaire
    call json_get_real(buf, "initial_condition.box_x0", cfg%ic_box_x0, 0.0_dp)
    ! Appel de sous-programme modulaire
    call json_get_real(buf, "initial_condition.box_y0", cfg%ic_box_y0, 0.0_dp)
    ! Appel de sous-programme modulaire
    call json_get_real(buf, "initial_condition.box_x1", cfg%ic_box_x1, 0.5_dp)
    ! Appel de sous-programme modulaire
    call json_get_real(buf, "initial_condition.box_y1", cfg%ic_box_y1, 0.5_dp)
    ! Appel de sous-programme modulaire
    call json_get_real(buf, "initial_condition.pert_amp", cfg%ic_pert_amp, 0.05_dp)
    ! Appel de sous-programme modulaire
    call json_get_real(buf, "initial_condition.pert_k", cfg%ic_pert_k, 1.0_dp)
    ! Lecture de l'épaisseur de nappe liquide (Cas 3.2 Fluids 2021)
    call json_get_real(buf, "initial_condition.thickness", cfg%ic_thickness, 0.0003_dp)
    ! Lecture de la largeur de jet liquide (Cas 3.3 Fluids 2021)
    call json_get_real(buf, "initial_condition.jet_width", cfg%ic_jet_width, 0.08_dp)

    ! Solver
    call json_get_string(buf, "solver.type", str_val, "AUGMENTED_LAGRANGIAN")
    ! Évaluation de la condition logique
    if (trim(str_val) == "CONJUGATE_GRADIENT") then
       ! Affectation et calcul numérique
       cfg%solver_type = SOLVER_AUGMENTED_LAGRANGIAN
       ! Affectation et calcul numérique
       cfg%linear_solver = LINEAR_CG
    ! Branche par défaut si condition non satisfaite
    else
       ! Affectation et calcul numérique
       cfg%solver_type = SOLVER_AUGMENTED_LAGRANGIAN
       ! Affectation et calcul numérique
       cfg%linear_solver = LINEAR_BICGSTAB2
    end if

    ! Appel de sous-programme modulaire
    call json_get_string(buf, "solver.linear_solver", str_val, "BICGSTAB2")
    ! Évaluation de la condition logique
    if (trim(str_val) == "BICGSTAB") then
       ! Affectation et calcul numérique
       cfg%linear_solver = LINEAR_BICGSTAB
    ! Branche conditionnelle alternative
    else if (trim(str_val) == "CG" .or. trim(str_val) == "CONJUGATE_GRADIENT") then
       ! Affectation et calcul numérique
       cfg%linear_solver = LINEAR_CG
    ! Branche conditionnelle alternative
    else if (trim(str_val) == "GMRES") then
       ! Affectation et calcul numérique
       cfg%linear_solver = LINEAR_GMRES
    end if

    ! Appel de sous-programme modulaire
    call json_get_real(buf, "solver.tol", cfg%tol, 1.0e-5_dp)
    ! Appel de sous-programme modulaire
    call json_get_int(buf, "solver.max_iter", cfg%max_iter, 300)
    ! Appel de sous-programme modulaire
    call json_get_real(buf, "solver.r_al", cfg%r_al, 100.0_dp)
    ! Appel de sous-programme modulaire
    call json_get_int(buf, "solver.uzawa_iter", cfg%uzawa_iter, 1)
    ! Appel de sous-programme modulaire
    call json_get_bool(buf, "solver.use_precond", cfg%use_precond, .true.)
    ! Appel de sous-programme modulaire
    call json_get_real(buf, "steady_state_tol", cfg%steady_state_tol, cfg%steady_state_tol)
    ! Appel de sous-programme modulaire
    call json_get_real(buf, "solver.steady_state_tol", cfg%steady_state_tol, cfg%steady_state_tol)

    ! VOF parameters
    call json_get_string(buf, "vof.advection_method", cfg%advection_method, "WEYMOUTH_YUE")
    ! Appel de sous-programme modulaire
    call json_get_string(buf, "vof.reconstruction_method", cfg%reconstruction_method, "PLIC")
    ! Appel de sous-programme modulaire
    call json_get_string(buf, "vof.tvd_limiter", cfg%tvd_limiter, "NONE")
    ! Appel de sous-programme modulaire
    call json_get_real(buf, "vof.cfl_vof", cfg%cfl_vof, 0.5_dp)

    ! Cohérence stricte : le schéma algébrique TVD ne s'accompagne d'aucune reconstruction géométrique,
    ! tandis que les schémas géométriques (WEYMOUTH_YUE, CIAM) n'utilisent pas de limiteur TVD.
    if (trim(cfg%advection_method) == "VOF_TVD" .or. trim(cfg%advection_method) == "TVD") then
       cfg%reconstruction_method = "NONE"
       if (trim(cfg%tvd_limiter) == "NONE" .or. len_trim(cfg%tvd_limiter) == 0) then
          cfg%tvd_limiter = "SUPERBEE"
       end if
    else
       if (trim(cfg%reconstruction_method) == "NONE" .or. len_trim(cfg%reconstruction_method) == 0) then
          cfg%reconstruction_method = "PLIC"
       end if
       cfg%tvd_limiter = "NONE"
    end if

    ! Output
    call json_get_string(buf, "output.output_dir", cfg%output_dir, "output")
    ! Appel de sous-programme modulaire
    call json_get_string(buf, "output.directory", cfg%output_dir, cfg%output_dir)
    ! Appel de sous-programme modulaire
    call json_get_int(buf, "output.vtk_interval", cfg%vtk_interval, 50)
    ! Appel de sous-programme modulaire
    call json_get_int(buf, "output.log_interval", cfg%log_interval, 10)
    ! Appel de sous-programme modulaire
    call json_get_bool(buf, "output.write_csv_history", cfg%write_csv, .true.)
  end subroutine load_config_from_json

end module module_config

\end{lstlisting}

\chapter{Maillage MAC 2D et Tableaux Column-Major à Bornes Fantômes en Fortran 2018 (module_mesh.f90, module_fields.f90)}
\section{Indiçage Naturel à Base 1 et Bornes Fantômes \texttt{(1-ng:nx+ng)} en Fortran 2018}
L'un des atouts majeurs de Fortran 2018 pour la mécanique des fluides numérique sur grille structurée est la possibilité d'allouer des tableaux multidimensionnels avec des bornes inférieures négatives ou nulles correspondant directement aux cellules fantômes :
\begin{itemize}
  \item \textbf{Champs scalaires centrés} ($p, T, C, \rho, \mu, \nabla\cdot\mathbf{u}$) :
  \begin{equation}
    \texttt{allocate(fields\%p(1-ng : nx+ng, 1-ng : ny+ng))}
  \end{equation}
  où le domaine intérieur physique occupe exactement $i \in \{1, \dots, N_x\}$ et $j \in \{1, \dots, N_y\}$.
  \item \textbf{Vitesse horizontale décalée $u$} ($N_x+1$ faces verticales intérieures/frontières d'indices $i \in \{1, \dots, N_x+1\}$) :
  \begin{equation}
    \texttt{allocate(fields\%u(1-ng : nx+1+ng, 1-ng : ny+ng))}
  \end{equation}
  \item \textbf{Vitesse verticale décalée $v$} ($N_y+1$ faces horizontales intérieures/frontières d'indices $j \in \{1, \dots, N_y+1\}$) :
  \begin{equation}
    \texttt{allocate(fields\%v(1-ng : nx+ng, 1-ng : ny+1+ng))}
  \end{equation}
\end{itemize}

\section{Ordre Mémoire \emph{Column-Major} et Ordre Optimal des Boucles Imbriquées}
En Fortran, les tableaux 2D \texttt{A(i, j)} sont stockés en mémoire par colonnes (\emph{Column-Major Order}) : l'élément \texttt{A(i+1, j)} est adjacent en mémoire à \texttt{A(i, j)}. Par conséquent, toutes les boucles spatiales dans \texttt{COVOF\_2D-FORTRAN} sont ordonnées avec \textbf{$j$ en boucle externe et $i$ en boucle interne} :
\begin{verbatim}
do j = 1, mesh%ny
  do i = 1, mesh%nx
    fields%div(i, j) = (fields%u(i+1, j) - fields%u(i, j)) / mesh%dx &
                     + (fields%v(i, j+1) - fields%v(i, j)) / mesh%dy
  end do
end do
\end{verbatim}
Cet ordonnancement garantit exactement la même localité spatiale L1/L2 et la même efficacité SIMD que l'ordonnancement \emph{Row-Major} de C++17.

\section{Répertoire des Types Dérivés, Sous-Routines et Fonctions de \texttt{src/module\_mesh.f90}}
\begin{center}\small
\begin{tabular}{clll}
\toprule
\textbf{Ligne} & \textbf{Nature} & \textbf{Identifiant Fortran 2018} & \textbf{Interface / Arguments} \\
\midrule
10 & module & \texttt{module\_mesh} & Unité modulaire Fortran 2018 \\
27 & type dérivé & \texttt{mesh\_2d\_t} & Structure de données Fortran 2018 \\
67 & subroutine & \texttt{init\_mesh} & Arguments : (mesh, comm\_world, nx, ny, lx, ly, perio...) \\
150 & subroutine & \texttt{destroy\_mesh} & Arguments : (mesh) \\
\bottomrule
\end{tabular}
\end{center}

\section{Code Source Intégral Annoté : \texttt{src/module\_mesh.f90}}
Le listing ci-dessous présente l'intégralité des 224 lignes du fichier \texttt{src/module\_mesh.f90} de l'implémentation Fortran 2018 :

\begin{lstlisting}[language={[95]Fortran}]
!===============================================================================
! Module        : module_mesh
! Description   : Définition de la topologie cartésienne 2D, gestion du
!                 partitionnement spatial multi-processus OpenMPI, calcul des
!                 coordonnées métriques (faces décalées et centres de cellules)
!                 et détection des inclusions géométriques d'obstacles.
! Auteur        : Elhadji Ibrahima Thiam
! Standard      : Fortran 2018 / MPI-3
!===============================================================================
module module_mesh
  ! Importation du descripteur de double précision
  use module_precision, only: dp
  ! Interdiction stricte du typage implicite pour garantir l'intégrité mémoire
  implicit none
  ! Par défaut, tous les membres et routines sont privés
  private
  ! Inclusion des définitions MPI requises pour les communicateurs cartésiens
  include 'mpif.h'

  ! Export explicite des types dérivés et fonctions publiques
  public :: mesh_2d_t, init_mesh, destroy_mesh
  public :: x_face, y_face, x_center, y_center, is_in_obstacle

  !-----------------------------------------------------------------------------
  ! Structure de données représentant le maillage 2D décomposé
  !-----------------------------------------------------------------------------
  type :: mesh_2d_t
     ! Dimensions globales du domaine de calcul (nombre total de mailles)
     integer  :: nx_global = 64, ny_global = 64
     ! Longueurs physiques totales du domaine en mètres [m]
     real(dp) :: lx = 1.0_dp, ly = 1.0_dp
     ! Pas d'espace uniformes selon les directions X et Y : dx = lx/nx, dy = ly/ny
     real(dp) :: dx = 0.0_dp, dy = 0.0_dp

     ! --- Topologie Cartésienne MPI et Rangs de Voisinage ---
     integer  :: rank = 0                 ! Rang du sous-domaine local courant
     integer  :: size = 1                 ! Nombre total de sous-domaines (processus)
     integer  :: comm = MPI_COMM_WORLD    ! Communicateur initial parent
     integer  :: cart_comm = MPI_COMM_NULL! Communicateur cartésien bidimensionnel dédié
     integer  :: px = 1, py = 1           ! Nombre de processus selon X et selon Y
     integer  :: coord_x = 0, coord_y = 0 ! Coordonnées cartésiennes du sous-domaine dans la grille MPI
     integer  :: neighbor_west = MPI_PROC_NULL  ! Rang du voisin Ouest (gauche, x-)
     integer  :: neighbor_east = MPI_PROC_NULL  ! Rang du voisin Est (droite, x+)
     integer  :: neighbor_south = MPI_PROC_NULL ! Rang du voisin Sud (bas, y-)
     integer  :: neighbor_north = MPI_PROC_NULL ! Rang du voisin Nord (haut, y+)

     ! --- Indicateurs Booléens de Frontières Physiques du Domaine Global ---
     logical  :: is_boundary_left   = .false. ! Vrai si le sous-domaine touche la paroi gauche x = 0
     logical  :: is_boundary_right  = .false. ! Vrai si le sous-domaine touche la paroi droite x = lx
     logical  :: is_boundary_bottom = .false. ! Vrai si le sous-domaine touche la paroi inférieure y = 0
     logical  :: is_boundary_top    = .false. ! Vrai si le sous-domaine touche la paroi supérieure y = ly

     ! --- Partitionnement Spatial Local ---
     integer  :: nx_local = 64, ny_local = 64 ! Nombre de cellules locales attribuées à ce processus
     integer  :: i_start_global = 1           ! Indice global de la première cellule locale en X (base 1)
     integer  :: j_start_global = 1           ! Indice global de la première cellule locale en Y (base 1)
  end type mesh_2d_t

contains

  !=============================================================================
  ! Routine : init_mesh
  ! Rôle    : Calcule la factorisation optimale du nombre de cœurs en une grille
  !           2D (px, py), crée la topologie cartésienne MPI avec périodicité,
  !           identifie les 4 voisins immédiats et répartit équitablement les mailles.
  !=============================================================================
  subroutine init_mesh(mesh, comm_world, nx, ny, lx, ly, periodic_x, periodic_y)
    type(mesh_2d_t), intent(out) :: mesh       ! Structure de maillage initialisée en sortie
    integer, intent(in)          :: comm_world ! Communicateur parent initial
    integer, intent(in)          :: nx, ny     ! Résolution globale souhaitée en cellules
    real(dp), intent(in)         :: lx, ly     ! Longueurs physiques du domaine rectangulaire
    logical, intent(in)          :: periodic_x ! Vrai si conditions périodiques en X
    logical, intent(in)          :: periodic_y ! Vrai si conditions périodiques en Y

    ! Variables de travail locales pour les primitives cartésiennes MPI
    integer :: dims(2), periods(2), coords(2)
    integer :: ierr, rank_in, size_in
    integer :: rem_x, rem_y

    ! 1. Enregistrement des dimensions physiques et de la taille du communicateur
    mesh%comm = comm_world
    call MPI_Comm_rank(comm_world, rank_in, ierr)
    call MPI_Comm_size(comm_world, size_in, ierr)
    mesh%rank = rank_in
    mesh%size = size_in

    mesh%nx_global = nx
    mesh%ny_global = ny
    mesh%lx = lx
    mesh%ly = ly
    ! Calcul des pas de discrétisation spatiale
    mesh%dx = lx / real(nx, dp)
    mesh%dy = ly / real(ny, dp)

    ! 2. Détermination automatique de la décomposition 2D optimale (px, py)
    dims = [0, 0]
    call MPI_Dims_create(mesh%size, 2, dims, ierr)
    mesh%px = dims(1)
    mesh%py = dims(2)

    ! Définition des drapeaux de périodicité pour la création de la topologie MPI
    periods(1) = merge(1, 0, periodic_x)
    periods(2) = merge(1, 0, periodic_y)

    ! 3. Création du communicateur virtuel cartésien MPI 2D
    call MPI_Cart_create(comm_world, 2, dims, periods, 1, mesh%cart_comm, ierr)
    call MPI_Comm_rank(mesh%cart_comm, mesh%rank, ierr)
    call MPI_Cart_coords(mesh%cart_comm, mesh%rank, 2, coords, ierr)
    mesh%coord_x = coords(1)
    mesh%coord_y = coords(2)

    ! 4. Recherche des rangs des 4 voisins adjacents par décalage unitaire cartésien
    ! Décalage direction 0 (axe X) : voisin Ouest (-1) et voisin Est (+1)
    call MPI_Cart_shift(mesh%cart_comm, 0, 1, mesh%neighbor_west, mesh%neighbor_east, ierr)
    ! Décalage direction 1 (axe Y) : voisin Sud (-1) et voisin Nord (+1)
    call MPI_Cart_shift(mesh%cart_comm, 1, 1, mesh%neighbor_south, mesh%neighbor_north, ierr)

    ! 5. Répartition équitable des mailles le long de l'axe X (gestion du reste mod)
    mesh%nx_local = nx / mesh%px
    rem_x = mod(nx, mesh%px)
    if (mesh%coord_x < rem_x) then
       ! Les premiers rem_x processus reçoivent une cellule supplémentaire
       mesh%nx_local = mesh%nx_local + 1
       mesh%i_start_global = mesh%coord_x * mesh%nx_local + 1
    else
       mesh%i_start_global = rem_x * (mesh%nx_local + 1) + (mesh%coord_x - rem_x) * mesh%nx_local + 1
    end if

    ! 6. Répartition équitable des mailles le long de l'axe Y
    mesh%ny_local = ny / mesh%py
    rem_y = mod(ny, mesh%py)
    if (mesh%coord_y < rem_y) then
       mesh%ny_local = mesh%ny_local + 1
       mesh%j_start_global = mesh%coord_y * mesh%ny_local + 1
    else
       mesh%j_start_global = rem_y * (mesh%ny_local + 1) + (mesh%coord_y - rem_y) * mesh%ny_local + 1
    end if

    ! 7. Détection des frontières physiques externes non périodiques
    mesh%is_boundary_left   = (mesh%neighbor_west == MPI_PROC_NULL)
    mesh%is_boundary_right  = (mesh%neighbor_east == MPI_PROC_NULL)
    mesh%is_boundary_bottom = (mesh%neighbor_south == MPI_PROC_NULL)
    mesh%is_boundary_top    = (mesh%neighbor_north == MPI_PROC_NULL)
  end subroutine init_mesh

  !=============================================================================
  ! Routine : destroy_mesh
  ! Rôle    : Libère le communicateur cartésien alloué auprès du runtime MPI.
  !=============================================================================
  subroutine destroy_mesh(mesh)
    type(mesh_2d_t), intent(inout) :: mesh
    integer :: ierr
    if (mesh%cart_comm /= MPI_COMM_NULL .and. mesh%cart_comm /= mesh%comm) then
       call MPI_Comm_free(mesh%cart_comm, ierr)
    end if
  end subroutine destroy_mesh

  !=============================================================================
  ! Fonction pure : x_face
  ! Rôle          : Calcule l'abscisse physique absolue de la face verticale i-1/2.
  !=============================================================================
  pure real(dp) function x_face(mesh, i)
    type(mesh_2d_t), intent(in) :: mesh
    integer, intent(in)          :: i
    x_face = real(mesh%i_start_global - 1 + i - 1, dp) * mesh%dx
  end function x_face

  !=============================================================================
  ! Fonction pure : y_face
  ! Rôle          : Calcule l'ordonnée physique absolue de la face horizontale j-1/2.
  !=============================================================================
  pure real(dp) function y_face(mesh, j)
    type(mesh_2d_t), intent(in) :: mesh
    integer, intent(in)          :: j
    y_face = real(mesh%j_start_global - 1 + j - 1, dp) * mesh%dy
  end function y_face

  !=============================================================================
  ! Fonction pure : x_center
  ! Rôle          : Calcule l'abscisse du barycentre géométrique de la cellule i.
  !=============================================================================
  pure real(dp) function x_center(mesh, i)
    type(mesh_2d_t), intent(in) :: mesh
    integer, intent(in)          :: i
    x_center = (real(mesh%i_start_global - 1 + i - 1, dp) + 0.5_dp) * mesh%dx
  end function x_center

  !=============================================================================
  ! Fonction pure : y_center
  ! Rôle          : Calcule l'ordonnée du barycentre géométrique de la cellule j.
  !=============================================================================
  pure real(dp) function y_center(mesh, j)
    type(mesh_2d_t), intent(in) :: mesh
    integer, intent(in)          :: j
    y_center = (real(mesh%j_start_global - 1 + j - 1, dp) + 0.5_dp) * mesh%dy
  end function y_center

  !=============================================================================
  ! Fonction pure : is_in_obstacle
  ! Rôle          : Teste l'appartenance d'un point (x, y) à un obstacle immergé
  !                 (rectangle ou cylindre) pour la pénalisation de Brinkman.
  !=============================================================================
  pure logical function is_in_obstacle(x, y, obs_type, x0, y0, x1, y1, xc, yc, r)
    real(dp), intent(in) :: x, y        ! Coordonnées du point à tester
    integer, intent(in)  :: obs_type    ! Type géométrique : 1=boîte, 2=cylindre
    real(dp), intent(in) :: x0, y0      ! Coin inférieur gauche du rectangle
    real(dp), intent(in) :: x1, y1      ! Coin supérieur droit du rectangle
    real(dp), intent(in) :: xc, yc      ! Centre du disque cylindrique
    real(dp), intent(in) :: r           ! Rayon du disque

    is_in_obstacle = .false.
    if (obs_type == 1) then ! Obstacle rectangulaire (ex: marche descendante)
       if (x >= x0 - 1.0e-9_dp .and. x <= x1 + 1.0e-9_dp .and. &
           y >= y0 - 1.0e-9_dp .and. y <= y1 + 1.0e-9_dp) then
          is_in_obstacle = .true.
       end if
    else if (obs_type == 2) then ! Obstacle cylindrique circulaire
       if ((x - xc)**2 + (y - yc)**2 <= r**2) then
          is_in_obstacle = .true.
       end if
    end if
  end function is_in_obstacle

end module module_mesh

\end{lstlisting}

\section{Répertoire des Types Dérivés, Sous-Routines et Fonctions de \texttt{src/module\_fields.f90}}
\begin{center}\small
\begin{tabular}{clll}
\toprule
\textbf{Ligne} & \textbf{Nature} & \textbf{Identifiant Fortran 2018} & \textbf{Interface / Arguments} \\
\midrule
11 & module & \texttt{module\_fields} & Unité modulaire Fortran 2018 \\
27 & type dérivé & \texttt{flow\_fields\_t} & Structure de données Fortran 2018 \\
70 & subroutine & \texttt{fields\_allocate\_mesh} & Arguments : (this, mesh) \\
80 & subroutine & \texttt{fields\_free} & Arguments : (this) \\
90 & subroutine & \texttt{allocate\_fields} & Arguments : (fields, nx, ny) \\
123 & subroutine & \texttt{deallocate\_fields} & Arguments : (fields) \\
150 & subroutine & \texttt{reset\_fields} & Arguments : (fields) \\
\bottomrule
\end{tabular}
\end{center}

\section{Code Source Intégral Annoté : \texttt{src/module\_fields.f90}}
Le listing ci-dessous présente l'intégralité des 171 lignes du fichier \texttt{src/module\_fields.f90} de l'implémentation Fortran 2018 :

\begin{lstlisting}[language={[95]Fortran}]
!===============================================================================
! Module        : module_fields
! Description   : Définition et gestion de la mémoire des champs hydrodynamiques,
!                 thermiques et diphasiques sur la grille décalée MAC (Harlow-Welch).
!                 Assure l'allocation contiguë avec cellules fantômes pour les
!                 échanges de halos MPI et garantit la politique Zero-Allocation
!                 au sein de la boucle en temps.
! Auteur        : Elhadji Ibrahima Thiam
! Standard      : Fortran 2018
!===============================================================================
module module_fields
  ! Importation du descripteur de double précision
  use module_precision, only: dp
  ! Importation du type dérivé de maillage pour connaître les dimensions locales
  use module_mesh, only: mesh_2d_t
  ! Typage strict
  implicit none
  ! Encapsulation privée
  private

  ! Exportation publique du type conteneur et des routines associées
  public :: flow_fields_t, allocate_fields, deallocate_fields, reset_fields

  !-----------------------------------------------------------------------------
  ! Structure de données contenant l'ensemble des champs physiques du sous-domaine
  !-----------------------------------------------------------------------------
  type :: flow_fields_t
     ! Nombre de cellules physiques locales (hors cellules fantômes)
     integer :: nx = 0, ny = 0

     ! --- Champs Hydrodynamiques Discrétisés sur Grille Décalée (MAC) ---
     ! Indices : (0:nx+1, 0:ny+1) pour inclure 1 rangée de cellules fantômes
     real(dp), allocatable :: u(:, :)        ! Vitesse axiale aux faces verticales (i-1/2, j) [m/s]
     real(dp), allocatable :: v(:, :)        ! Vitesse transversale aux faces horizontales (i, j-1/2) [m/s]
     real(dp), allocatable :: p(:, :)        ! Pression statique scalaire au centre des mailles (i, j) [Pa]
     real(dp), allocatable :: div(:, :)      ! Divergence discrète locale de la vitesse (nabla . u) [1/s]

     ! --- Propriétés Thermodynamiques et Équation d'Énergie ---
     real(dp), allocatable :: C(:, :)        ! Fraction volumique VOF du fluide 1 (C in [0, 1])
     real(dp), allocatable :: T(:, :)        ! Température thermodynamique au centre des mailles [K]
     real(dp), allocatable :: rho(:, :)      ! Masse volumique effective au centre des cellules [kg/m^3]
     real(dp), allocatable :: rho_u(:, :)    ! Masse volumique interpolée aux faces u [kg/m^3]
     real(dp), allocatable :: rho_v(:, :)    ! Masse volumique interpolée aux faces v [kg/m^3]
     real(dp), allocatable :: mu(:, :)       ! Viscosité dynamique au centre des mailles [Pa.s]
     real(dp), allocatable :: mu_node(:, :)  ! Viscosité harmonique aux sommets/nœuds (i+1/2, j+1/2) [Pa.s]

     ! --- Forces Capillaires et Tension de Surface (Modèle CSF Brackbill) ---
     real(dp), allocatable :: fx_sigma(:, :) ! Force capillaire volumique selon X aux faces u [N/m^3]
     real(dp), allocatable :: fy_sigma(:, :) ! Force capillaire volumique selon Y aux faces v [N/m^3]
     real(dp), allocatable :: curvature(:, :)! Courbure locale d'interface kappa = -nabla . n [1/m]

     ! --- Masque d'Obstacle Immergé (Pénalisation de Brinkman) ---
     integer,  allocatable :: solid_mask(:, :)! 1 si la cellule est dans le solide, 0 en fluide pur

     ! --- Indicateurs et Diagnostics de Validation ---
     real(dp) :: mean_nusselt = 0.0_dp       ! Nombre de Nusselt moyen mesuré à la paroi chaude

  contains
     ! Méthodes liées au type dérivé pour syntaxe orientée objet
     procedure :: allocate => fields_allocate_mesh
     procedure :: free     => fields_free
  end type flow_fields_t

contains

  !=============================================================================
  ! Routine liée : fields_allocate_mesh
  ! Rôle         : Alloue les tableaux à partir d'une instance de mesh_2d_t.
  !=============================================================================
  subroutine fields_allocate_mesh(this, mesh)
    class(flow_fields_t), intent(inout) :: this
    type(mesh_2d_t), intent(in)         :: mesh
    call allocate_fields(this, mesh%nx_local, mesh%ny_local)
  end subroutine fields_allocate_mesh

  !=============================================================================
  ! Routine liée : fields_free
  ! Rôle         : Désalloue proprement l'ensemble des tableaux du conteneur.
  !=============================================================================
  subroutine fields_free(this)
    class(flow_fields_t), intent(inout) :: this
    call deallocate_fields(this)
  end subroutine fields_free

  !=============================================================================
  ! Routine : allocate_fields
  ! Rôle    : Alloue de façon contiguë en mémoire chaque champ 2D avec ses halos
  !           fantômes (indices de 0 à nx+1 et de 0 à ny+1) et initialise à zéro.
  !=============================================================================
  subroutine allocate_fields(fields, nx, ny)
    type(flow_fields_t), intent(inout) :: fields
    integer, intent(in)                :: nx, ny

    ! Enregistrement des dimensions locales
    fields%nx = nx
    fields%ny = ny

    ! Allocation dynamique unique (zéro allocation subséquente dans la boucle temporelle)
    allocate(fields%u(0:nx+1, 0:ny+1))
    allocate(fields%v(0:nx+1, 0:ny+1))
    allocate(fields%p(0:nx+1, 0:ny+1))
    allocate(fields%div(0:nx+1, 0:ny+1))
    allocate(fields%C(0:nx+1, 0:ny+1))
    allocate(fields%T(0:nx+1, 0:ny+1))
    allocate(fields%rho(0:nx+1, 0:ny+1))
    allocate(fields%rho_u(0:nx+1, 0:ny+1))
    allocate(fields%rho_v(0:nx+1, 0:ny+1))
    allocate(fields%mu(0:nx+1, 0:ny+1))
    allocate(fields%mu_node(0:nx+1, 0:ny+1))
    allocate(fields%fx_sigma(0:nx+1, 0:ny+1))
    allocate(fields%fy_sigma(0:nx+1, 0:ny+1))
    allocate(fields%curvature(0:nx+1, 0:ny+1))
    allocate(fields%solid_mask(0:nx+1, 0:ny+1))

    ! Remise à zéro des valeurs initiales
    call reset_fields(fields)
  end subroutine allocate_fields

  !=============================================================================
  ! Routine : deallocate_fields
  ! Rôle    : Désalloue chaque tableau et remet les compteurs de dimension à zéro.
  !=============================================================================
  subroutine deallocate_fields(fields)
    type(flow_fields_t), intent(inout) :: fields

    if (allocated(fields%u))          deallocate(fields%u)
    if (allocated(fields%v))          deallocate(fields%v)
    if (allocated(fields%p))          deallocate(fields%p)
    if (allocated(fields%div))        deallocate(fields%div)
    if (allocated(fields%C))          deallocate(fields%C)
    if (allocated(fields%T))          deallocate(fields%T)
    if (allocated(fields%rho))        deallocate(fields%rho)
    if (allocated(fields%rho_u))      deallocate(fields%rho_u)
    if (allocated(fields%rho_v))      deallocate(fields%rho_v)
    if (allocated(fields%mu))         deallocate(fields%mu)
    if (allocated(fields%mu_node))    deallocate(fields%mu_node)
    if (allocated(fields%fx_sigma))   deallocate(fields%fx_sigma)
    if (allocated(fields%fy_sigma))   deallocate(fields%fy_sigma)
    if (allocated(fields%curvature))  deallocate(fields%curvature)
    if (allocated(fields%solid_mask)) deallocate(fields%solid_mask)

    fields%nx = 0
    fields%ny = 0
  end subroutine deallocate_fields

  !=============================================================================
  ! Routine : reset_fields
  ! Rôle    : Initialise les champs à des valeurs physiques par défaut neutres.
  !=============================================================================
  subroutine reset_fields(fields)
    type(flow_fields_t), intent(inout) :: fields

    fields%u            = 0.0_dp     ! Vitesse axiale initiale nulle
    fields%v            = 0.0_dp     ! Vitesse transversale initiale nulle
    fields%p            = 0.0_dp     ! Pression de référence nulle
    fields%div          = 0.0_dp     ! Divergence initiale nulle
    fields%C            = 1.0_dp     ! Fluide 1 pur par défaut
    fields%T            = 0.0_dp     ! Température de référence nulle
    fields%rho          = 1.0_dp     ! Masse volumique unitaire par défaut
    fields%rho_u        = 1.0_dp
    fields%rho_v        = 1.0_dp
    fields%mu           = 1.0e-3_dp  ! Viscosité dynamique standard de l'eau (1 mPa.s)
    fields%mu_node      = 1.0e-3_dp
    fields%fx_sigma     = 0.0_dp     ! Tension superficielle nulle
    fields%fy_sigma     = 0.0_dp
    fields%curvature    = 0.0_dp
    fields%solid_mask   = 0          ! 100% domaine fluide sans obstacle
    fields%mean_nusselt = 0.0_dp
  end subroutine reset_fields

end module module_fields

\end{lstlisting}

\chapter{Communications Parallèles MPI et Échanges de Halos en Fortran 2018 (module_mpi_comm.f90)}
\section{Gestion de la Topologie Cartésienne et des Halos dans \texttt{module\_mpi\_comm.f90}}
Le module \texttt{src/module\_mpi\_comm.f90} (798 lignes) fournit le type dérivé \texttt{mpi\_context\_t} et les sous-routines d'échange de couronnes fantômes pour les trois topologies de grilles MAC :
\begin{enumerate}
  \item \texttt{exchange\_halo\_scalar(mpi\_ctx, mesh, phi)} : synchronise un tableau centré \texttt{phi(1-ng:nx+ng, 1-ng:ny+ng)} sur une épaisseur \texttt{ng} en échangeant d'abord les colonnes Est/Ouest puis les lignes Nord/Sud (propageant ainsi les coins diagonaux).
  \item \texttt{exchange\_halo\_u(mpi\_ctx, mesh, u)} et \texttt{exchange\_halo\_v(mpi\_ctx, mesh, v)} : gèrent le décalage d'une face sur les bords partagés entre deux sous-domaines MPI adjacents ainsi que la périodicité torique (\texttt{BC\_PERIODIC}).
  \item \texttt{mpi\_allreduce\_sum}, \texttt{mpi\_allreduce\_max}, \texttt{mpi\_allreduce\_min} : encapsulent les réductions collectives globales utilisées par les solveurs de Krylov et les bilans intégraux de conservation.
\end{enumerate}

\section{Répertoire des Types Dérivés, Sous-Routines et Fonctions de \texttt{src/module\_mpi\_comm.f90}}
\begin{center}\small
\begin{tabular}{clll}
\toprule
\textbf{Ligne} & \textbf{Nature} & \textbf{Identifiant Fortran 2018} & \textbf{Interface / Arguments} \\
\midrule
2 & module & \texttt{module\_mpi\_comm} & Unité modulaire Fortran 2018 \\
25 & subroutine & \texttt{exchange\_halos\_scalar} & Arguments : (field, mesh) \\
139 & subroutine & \texttt{exchange\_halos\_velocity} & Arguments : (u, v, mesh) \\
290 & subroutine & \texttt{apply\_velocity\_bc} & Arguments : (fields, mesh, cfg, homogeneous) \\
514 & subroutine & \texttt{apply\_thermal\_bc} & Arguments : (fields, mesh, cfg) \\
610 & subroutine & \texttt{apply\_vof\_bc} & Arguments : (fields, mesh, cfg) \\
700 & subroutine & \texttt{sync\_u} & Arguments : (mesh, u) \\
710 & subroutine & \texttt{sync\_v} & Arguments : (mesh, v) \\
720 & subroutine & \texttt{sync\_cell\_scalar} & Arguments : (mesh, s) \\
730 & subroutine & \texttt{apply\_boundary\_conditions} & Arguments : (mesh, cfg, fields, homogeneous) \\
744 & subroutine & \texttt{apply\_thermal\_boundary\_conditions} & Arguments : (mesh, cfg, fields) \\
\bottomrule
\end{tabular}
\end{center}

\section{Code Source Intégral Annoté : \texttt{src/module\_mpi\_comm.f90}}
Le listing ci-dessous présente l'intégralité des 813 lignes du fichier \texttt{src/module\_mpi\_comm.f90} de l'implémentation Fortran 2018 :

\begin{lstlisting}[language={[95]Fortran}]
! [Procédure/Module] module module_mpi_comm
module module_mpi_comm
  use module_precision, only: dp
  use module_constants
  use module_mesh, only: mesh_2d_t, x_center, y_center, x_face
  use module_fields, only: flow_fields_t
  use module_config, only: config_t
  implicit none
  private
  include 'mpif.h'

  public :: exchange_halos_scalar
  public :: exchange_halos_velocity
  public :: apply_velocity_bc
  public :: apply_thermal_bc
  public :: apply_vof_bc
  public :: global_max_val, global_min_val, global_sum_val
  public :: sync_u, sync_v, sync_cell_scalar
  public :: apply_boundary_conditions, apply_thermal_boundary_conditions
  public :: global_max, global_min, global_sum, global_norm2, global_dot_product

contains

  ! [Procédure/Module] subroutine exchange_halos_scalar
  subroutine exchange_halos_scalar(field, mesh)
    ! Déclaration de variable locale / paramètre de calcul
    type(mesh_2d_t), intent(in) :: mesh
    ! Déclaration de variable locale / paramètre de calcul
    real(dp), intent(inout) :: field(0:mesh%nx_local+1, 0:mesh%ny_local+1)

    ! Déclaration de variable locale / paramètre de calcul
    integer :: nx, ny, ierr
    ! Déclaration de variable locale / paramètre de calcul
    integer :: reqs(8), n_reqs
    ! Déclaration de variable locale / paramètre de calcul
    real(dp), allocatable :: s_left(:), r_left(:), s_right(:), r_right(:)
    ! Déclaration de variable locale / paramètre de calcul
    real(dp), allocatable :: s_bot(:), r_bot(:), s_top(:), r_top(:)

    ! Affectation et calcul numérique
    nx = mesh%nx_local
    ! Affectation et calcul numérique
    ny = mesh%ny_local

    ! Affectation et calcul numérique
    if (mesh%size <= 1) then
       if (.not. mesh%is_boundary_left) then
          field(0, 0:ny+1)    = field(nx, 0:ny+1)
          field(nx+1, 0:ny+1) = field(1, 0:ny+1)
       end if
       if (.not. mesh%is_boundary_bottom) then
          field(0:nx+1, 0)    = field(0:nx+1, ny)
          field(0:nx+1, ny+1) = field(0:nx+1, 1)
       end if
       return
    end if

    ! Allocation dynamique de tableau contigu en mémoire
    allocate(s_left(ny), r_left(ny), s_right(ny), r_right(ny))
    ! Allocation dynamique de tableau contigu en mémoire
    allocate(s_bot(nx), r_bot(nx), s_top(nx), r_top(nx))

    ! Affectation et calcul numérique
    s_left(:)  = field(1, 1:ny)
    ! Affectation et calcul numérique
    s_right(:) = field(nx, 1:ny)

    ! Affectation et calcul numérique
    n_reqs = 0
    ! Évaluation de la condition logique
    if (mesh%neighbor_west /= MPI_PROC_NULL) then
       ! Appel à la bibliothèque de communication parallèle MPI
       call MPI_Irecv(r_left, ny, MPI_DOUBLE_PRECISION, mesh%neighbor_west, 10, mesh%comm, reqs(n_reqs+1), ierr)
       ! Appel à la bibliothèque de communication parallèle MPI
       call MPI_Isend(s_left, ny, MPI_DOUBLE_PRECISION, mesh%neighbor_west, 11, mesh%comm, reqs(n_reqs+2), ierr)
       ! Affectation et calcul numérique
       n_reqs = n_reqs + 2
    end if
    ! Évaluation de la condition logique
    if (mesh%neighbor_east /= MPI_PROC_NULL) then
       ! Appel à la bibliothèque de communication parallèle MPI
       call MPI_Irecv(r_right, ny, MPI_DOUBLE_PRECISION, mesh%neighbor_east, 11, mesh%comm, reqs(n_reqs+1), ierr)
       ! Appel à la bibliothèque de communication parallèle MPI
       call MPI_Isend(s_right, ny, MPI_DOUBLE_PRECISION, mesh%neighbor_east, 10, mesh%comm, reqs(n_reqs+2), ierr)
       ! Affectation et calcul numérique
       n_reqs = n_reqs + 2
    end if

    ! Évaluation de la condition logique
    if (n_reqs > 0) then
       ! Appel à la bibliothèque de communication parallèle MPI
       call MPI_Waitall(n_reqs, reqs, MPI_STATUSES_IGNORE, ierr)
    end if

    ! Affectation et calcul numérique
    if (mesh%neighbor_west /= MPI_PROC_NULL) field(0, 1:ny) = r_left(:)
    ! Affectation et calcul numérique
    if (mesh%neighbor_east /= MPI_PROC_NULL) field(nx+1, 1:ny) = r_right(:)

    ! Affectation et calcul numérique
    s_bot(:) = field(1:nx, 1)
    ! Actualisation de la pression statique au centre de la maille
    s_top(:) = field(1:nx, ny)

    ! Affectation et calcul numérique
    n_reqs = 0
    ! Évaluation de la condition logique
    if (mesh%neighbor_south /= MPI_PROC_NULL) then
       ! Appel à la bibliothèque de communication parallèle MPI
       call MPI_Irecv(r_bot, nx, MPI_DOUBLE_PRECISION, mesh%neighbor_south, 20, mesh%comm, reqs(n_reqs+1), ierr)
       ! Appel à la bibliothèque de communication parallèle MPI
       call MPI_Isend(s_bot, nx, MPI_DOUBLE_PRECISION, mesh%neighbor_south, 21, mesh%comm, reqs(n_reqs+2), ierr)
       ! Affectation et calcul numérique
       n_reqs = n_reqs + 2
    end if
    ! Évaluation de la condition logique
    if (mesh%neighbor_north /= MPI_PROC_NULL) then
       ! Appel à la bibliothèque de communication parallèle MPI
       call MPI_Irecv(r_top, nx, MPI_DOUBLE_PRECISION, mesh%neighbor_north, 21, mesh%comm, reqs(n_reqs+1), ierr)
       ! Appel à la bibliothèque de communication parallèle MPI
       call MPI_Isend(s_top, nx, MPI_DOUBLE_PRECISION, mesh%neighbor_north, 20, mesh%comm, reqs(n_reqs+2), ierr)
       ! Affectation et calcul numérique
       n_reqs = n_reqs + 2
    end if

    ! Évaluation de la condition logique
    if (n_reqs > 0) then
       ! Appel à la bibliothèque de communication parallèle MPI
       call MPI_Waitall(n_reqs, reqs, MPI_STATUSES_IGNORE, ierr)
    end if

    ! Affectation et calcul numérique
    if (mesh%neighbor_south /= MPI_PROC_NULL) field(1:nx, 0) = r_bot(:)
    ! Affectation et calcul numérique
    if (mesh%neighbor_north /= MPI_PROC_NULL) field(1:nx, ny+1) = r_top(:)
  end subroutine exchange_halos_scalar

  ! [Procédure/Module] subroutine exchange_halos_velocity
  subroutine exchange_halos_velocity(u, v, mesh)
    ! Déclaration de variable locale / paramètre de calcul
    type(mesh_2d_t), intent(in) :: mesh
    ! Déclaration de variable locale / paramètre de calcul
    real(dp), intent(inout) :: u(0:mesh%nx_local+1, 0:mesh%ny_local+1)
    ! Déclaration de variable locale / paramètre de calcul
    real(dp), intent(inout) :: v(0:mesh%nx_local+1, 0:mesh%ny_local+1)

    ! Déclaration de variable locale / paramètre de calcul
    integer :: nx, ny, ierr
    ! Déclaration de variable locale / paramètre de calcul
    integer :: reqs(8), n_reqs
    ! Déclaration de variable locale / paramètre de calcul
    real(dp), allocatable :: su_l(:), ru_l(:), su_r(:), ru_r(:)
    ! Déclaration de variable locale / paramètre de calcul
    real(dp), allocatable :: sv_l(:), rv_l(:), sv_r(:), rv_r(:)
    ! Déclaration de variable locale / paramètre de calcul
    real(dp), allocatable :: su_b(:), ru_b(:), su_t(:), ru_t(:)
    ! Déclaration de variable locale / paramètre de calcul
    real(dp), allocatable :: sv_b(:), rv_b(:), sv_t(:), rv_t(:)

    ! Affectation et calcul numérique
    if (mesh%size <= 1) return

    ! Affectation et calcul numérique
    nx = mesh%nx_local
    ! Affectation et calcul numérique
    ny = mesh%ny_local

    ! Allocation dynamique de tableau contigu en mémoire
    allocate(su_l(ny), ru_l(ny), su_r(ny), ru_r(ny))
    ! Allocation dynamique de tableau contigu en mémoire
    allocate(sv_l(ny), rv_l(ny), sv_r(ny), rv_r(ny))

    ! Affectation et calcul numérique
    su_l(:) = u(1, 1:ny)
    ! Affectation et calcul numérique
    su_r(:) = u(nx, 1:ny)
    ! Affectation et calcul numérique
    sv_l(:) = v(1, 1:ny)
    ! Affectation et calcul numérique
    sv_r(:) = v(nx, 1:ny)

    ! Affectation et calcul numérique
    n_reqs = 0
    ! Évaluation de la condition logique
    if (mesh%neighbor_west /= MPI_PROC_NULL) then
       ! Appel à la bibliothèque de communication parallèle MPI
       call MPI_Irecv(ru_l, ny, MPI_DOUBLE_PRECISION, mesh%neighbor_west, 30, mesh%comm, reqs(n_reqs+1), ierr)
       ! Appel à la bibliothèque de communication parallèle MPI
       call MPI_Isend(su_l, ny, MPI_DOUBLE_PRECISION, mesh%neighbor_west, 31, mesh%comm, reqs(n_reqs+2), ierr)
       ! Appel à la bibliothèque de communication parallèle MPI
       call MPI_Irecv(rv_l, ny, MPI_DOUBLE_PRECISION, mesh%neighbor_west, 32, mesh%comm, reqs(n_reqs+3), ierr)
       ! Appel à la bibliothèque de communication parallèle MPI
       call MPI_Isend(sv_l, ny, MPI_DOUBLE_PRECISION, mesh%neighbor_west, 33, mesh%comm, reqs(n_reqs+4), ierr)
       ! Affectation et calcul numérique
       n_reqs = n_reqs + 4
    end if
    ! Évaluation de la condition logique
    if (mesh%neighbor_east /= MPI_PROC_NULL) then
       ! Appel à la bibliothèque de communication parallèle MPI
       call MPI_Irecv(ru_r, ny, MPI_DOUBLE_PRECISION, mesh%neighbor_east, 31, mesh%comm, reqs(n_reqs+1), ierr)
       ! Appel à la bibliothèque de communication parallèle MPI
       call MPI_Isend(su_r, ny, MPI_DOUBLE_PRECISION, mesh%neighbor_east, 30, mesh%comm, reqs(n_reqs+2), ierr)
       ! Appel à la bibliothèque de communication parallèle MPI
       call MPI_Irecv(rv_r, ny, MPI_DOUBLE_PRECISION, mesh%neighbor_east, 33, mesh%comm, reqs(n_reqs+3), ierr)
       ! Appel à la bibliothèque de communication parallèle MPI
       call MPI_Isend(sv_r, ny, MPI_DOUBLE_PRECISION, mesh%neighbor_east, 32, mesh%comm, reqs(n_reqs+4), ierr)
       ! Affectation et calcul numérique
       n_reqs = n_reqs + 4
    end if

    if (n_reqs > 0) call MPI_Waitall(n_reqs, reqs, MPI_STATUSES_IGNORE, ierr)

    ! Évaluation de la condition logique
    if (mesh%neighbor_west /= MPI_PROC_NULL) then
       ! Mise à jour de la composante de vitesse discrète sur face décalée MAC
       u(0, 1:ny) = ru_l(:)
       ! Mise à jour de la composante de vitesse discrète sur face décalée MAC
       v(0, 1:ny) = rv_l(:)
    end if
    ! Évaluation de la condition logique
    if (mesh%neighbor_east /= MPI_PROC_NULL) then
       ! Mise à jour de la composante de vitesse discrète sur face décalée MAC
       u(nx+1, 1:ny) = ru_r(:)
       ! Mise à jour de la composante de vitesse discrète sur face décalée MAC
       v(nx+1, 1:ny) = rv_r(:)
    end if

    ! Allocation dynamique de tableau contigu en mémoire
    allocate(su_b(nx+2), ru_b(nx+2), su_t(nx+2), ru_t(nx+2))
    ! Allocation dynamique de tableau contigu en mémoire
    allocate(sv_b(nx),   rv_b(nx),   sv_t(nx),   rv_t(nx))

    ! Affectation et calcul numérique
    su_b(:) = u(0:nx+1, 1)
    ! Affectation et calcul numérique
    su_t(:) = u(0:nx+1, ny)
    ! Affectation et calcul numérique
    sv_b(:) = v(1:nx, 1)
    ! Affectation et calcul numérique
    sv_t(:) = v(1:nx, ny)

    ! Affectation et calcul numérique
    n_reqs = 0
    ! Évaluation de la condition logique
    if (mesh%neighbor_south /= MPI_PROC_NULL) then
       ! Appel à la bibliothèque de communication parallèle MPI
       call MPI_Irecv(ru_b, nx+2, MPI_DOUBLE_PRECISION, mesh%neighbor_south, 40, mesh%comm, reqs(n_reqs+1), ierr)
       ! Appel à la bibliothèque de communication parallèle MPI
       call MPI_Isend(su_b, nx+2, MPI_DOUBLE_PRECISION, mesh%neighbor_south, 41, mesh%comm, reqs(n_reqs+2), ierr)
       ! Appel à la bibliothèque de communication parallèle MPI
       call MPI_Irecv(rv_b, nx,   MPI_DOUBLE_PRECISION, mesh%neighbor_south, 42, mesh%comm, reqs(n_reqs+3), ierr)
       ! Appel à la bibliothèque de communication parallèle MPI
       call MPI_Isend(sv_b, nx,   MPI_DOUBLE_PRECISION, mesh%neighbor_south, 43, mesh%comm, reqs(n_reqs+4), ierr)
       ! Affectation et calcul numérique
       n_reqs = n_reqs + 4
    end if
    ! Évaluation de la condition logique
    if (mesh%neighbor_north /= MPI_PROC_NULL) then
       ! Appel à la bibliothèque de communication parallèle MPI
       call MPI_Irecv(ru_t, nx+2, MPI_DOUBLE_PRECISION, mesh%neighbor_north, 41, mesh%comm, reqs(n_reqs+1), ierr)
       ! Appel à la bibliothèque de communication parallèle MPI
       call MPI_Isend(su_t, nx+2, MPI_DOUBLE_PRECISION, mesh%neighbor_north, 40, mesh%comm, reqs(n_reqs+2), ierr)
       ! Appel à la bibliothèque de communication parallèle MPI
       call MPI_Irecv(rv_t, nx,   MPI_DOUBLE_PRECISION, mesh%neighbor_north, 43, mesh%comm, reqs(n_reqs+3), ierr)
       ! Appel à la bibliothèque de communication parallèle MPI
       call MPI_Isend(sv_t, nx,   MPI_DOUBLE_PRECISION, mesh%neighbor_north, 42, mesh%comm, reqs(n_reqs+4), ierr)
       ! Affectation et calcul numérique
       n_reqs = n_reqs + 4
    end if

    if (n_reqs > 0) call MPI_Waitall(n_reqs, reqs, MPI_STATUSES_IGNORE, ierr)

    ! Évaluation de la condition logique
    if (mesh%neighbor_south /= MPI_PROC_NULL) then
       ! Mise à jour de la composante de vitesse discrète sur face décalée MAC
       u(0:nx+1, 0) = ru_b(:)
       ! Mise à jour de la composante de vitesse discrète sur face décalée MAC
       v(1:nx, 0)   = rv_b(:)
    end if
    ! Évaluation de la condition logique
    if (mesh%neighbor_north /= MPI_PROC_NULL) then
       ! Mise à jour de la composante de vitesse discrète sur face décalée MAC
       u(0:nx+1, ny+1) = ru_t(:)
       ! Mise à jour de la composante de vitesse discrète sur face décalée MAC
       v(1:nx, ny+1)   = rv_t(:)
    end if
  end subroutine exchange_halos_velocity

  ! [Procédure/Module] subroutine apply_velocity_bc
  subroutine apply_velocity_bc(fields, mesh, cfg, homogeneous)
    ! Déclaration de variable locale / paramètre de calcul
    type(flow_fields_t), intent(inout) :: fields
    ! Déclaration de variable locale / paramètre de calcul
    type(mesh_2d_t), intent(in) :: mesh
    ! Déclaration de variable locale / paramètre de calcul
    type(config_t), intent(in) :: cfg
    ! Déclaration de variable locale / paramètre de calcul
    logical, intent(in) :: homogeneous

    ! Déclaration de variable locale / paramètre de calcul
    integer :: nx, ny, i, j
    ! Déclaration de variable locale / paramètre de calcul
    real(dp) :: u_l, v_l, u_r, v_r, u_b, v_b, u_t, v_t
    ! Déclaration de variable locale / paramètre de calcul
    real(dp) :: y, h, u_inlet
    ! Déclaration des variables pour injecteur supérieur (Cas 3.3 Fluids 2021)
    real(dp) :: x_c, xc_top, half_w_top, v_out_top

    ! Affectation et calcul numérique
    nx = mesh%nx_local
    ! Affectation et calcul numérique
    ny = mesh%ny_local

    ! Affectation et calcul numérique
    u_l = merge(0.0_dp, cfg%bc_left%u, homogeneous)
    ! Affectation et calcul numérique
    v_l = merge(0.0_dp, cfg%bc_left%v, homogeneous)
    ! Affectation et calcul numérique
    u_r = merge(0.0_dp, cfg%bc_right%u, homogeneous)
    ! Affectation et calcul numérique
    v_r = merge(0.0_dp, cfg%bc_right%v, homogeneous)
    ! Affectation et calcul numérique
    u_b = merge(0.0_dp, cfg%bc_bottom%u, homogeneous)
    ! Affectation et calcul numérique
    v_b = merge(0.0_dp, cfg%bc_bottom%v, homogeneous)
    ! Affectation et calcul numérique
    u_t = merge(0.0_dp, cfg%bc_top%u, homogeneous)
    ! Affectation et calcul numérique
    v_t = merge(0.0_dp, cfg%bc_top%v, homogeneous)

    ! Traitement périodicité si 1 seul processus
    if (cfg%bc_left%btype == BC_PERIODIC .and. mesh%px == 1) then
       ! Mise à jour de la composante de vitesse discrète sur face décalée MAC
       fields%u(0, 1:ny)    = fields%u(nx, 1:ny)
       ! Mise à jour de la composante de vitesse discrète sur face décalée MAC
       fields%u(nx+1, 1:ny) = fields%u(1, 1:ny)
       ! Mise à jour de la composante de vitesse discrète sur face décalée MAC
       fields%v(0, 1:ny)    = fields%v(nx, 1:ny)
       ! Mise à jour de la composante de vitesse discrète sur face décalée MAC
       fields%v(nx+1, 1:ny) = fields%v(1, 1:ny)
    end if
    ! Évaluation de la condition logique
    if (cfg%bc_bottom%btype == BC_PERIODIC .and. mesh%py == 1) then
       ! Mise à jour de la composante de vitesse discrète sur face décalée MAC
       fields%v(1:nx, 0)    = fields%v(1:nx, ny)
       ! Mise à jour de la composante de vitesse discrète sur face décalée MAC
       fields%v(1:nx, ny+1) = fields%v(1:nx, 1)
       ! Mise à jour de la composante de vitesse discrète sur face décalée MAC
       fields%u(0:nx+1, 0)    = fields%u(0:nx+1, ny)
       ! Mise à jour de la composante de vitesse discrète sur face décalée MAC
       fields%u(0:nx+1, ny+1) = fields%u(0:nx+1, 1)
    end if

    ! Bord Ouest physique (x = 0)
    if (mesh%neighbor_west == MPI_PROC_NULL .and. cfg%bc_left%btype /= BC_PERIODIC) then
       ! Boucle spatiale sur les cellules locales selon l'axe vertical Y (ordonnées)
       do j = 1, ny
          ! Évaluation de la condition logique
          if (cfg%bc_left%btype == BC_DIRICHLET) then
             ! Affectation et calcul numérique
             u_inlet = u_l
             ! Évaluation de la condition logique
             if (.not. homogeneous .and. cfg%obs_enabled .and. cfg%obs_type == OBS_BOX) then
                ! Profil d'entrée marche descendante : parabole sur [y_step, ly]
                y = y_center(mesh, j)
                ! Évaluation de la condition logique
                if (y < cfg%obs_y1) then
                   ! Affectation et calcul numérique
                   u_inlet = 0.0_dp
                ! Branche par défaut si condition non satisfaite
                else
                   ! Affectation et calcul numérique
                   h = mesh%ly - cfg%obs_y1
                   ! Affectation et calcul numérique
                   u_inlet = 6.0_dp * cfg%bc_left%u * (y - cfg%obs_y1) * (mesh%ly - y) / (h * h)
                end if
             end if
             ! Mise à jour de la composante de vitesse discrète sur face décalée MAC
             if (.not. homogeneous) fields%u(1, j) = u_inlet
             ! Mise à jour de la composante de vitesse discrète sur face décalée MAC
             fields%u(0, j) = u_inlet
             ! Mise à jour de la composante de vitesse discrète sur face décalée MAC
             fields%v(0, j) = 2.0_dp * v_l - fields%v(1, j)
          ! Branche conditionnelle alternative
          else if (cfg%bc_left%btype == BC_FREESLIP) then
             ! Mise à jour de la composante de vitesse discrète sur face décalée MAC
             if (.not. homogeneous) fields%u(1, j) = 0.0_dp
             ! Mise à jour de la composante de vitesse discrète sur face décalée MAC
             fields%u(0, j) = 0.0_dp
             ! Mise à jour de la composante de vitesse discrète sur face décalée MAC
             fields%v(0, j) = fields%v(1, j)
          ! Branche conditionnelle alternative
          else if (cfg%bc_left%btype == BC_NEUMANN) then
             ! Mise à jour de la composante de vitesse discrète sur face décalée MAC
             fields%u(0, j) = fields%u(1, j)
             ! Mise à jour de la composante de vitesse discrète sur face décalée MAC
             fields%v(0, j) = fields%v(1, j)
          end if
       end do
    end if

    ! Bord Est physique (x = Lx)
    if (mesh%neighbor_east == MPI_PROC_NULL .and. cfg%bc_right%btype /= BC_PERIODIC) then
       ! Boucle spatiale sur les cellules locales selon l'axe vertical Y (ordonnées)
       do j = 1, ny
          ! Évaluation de la condition logique
          if (cfg%bc_right%btype == BC_DIRICHLET) then
             ! Mise à jour de la composante de vitesse discrète sur face décalée MAC
             if (.not. homogeneous) fields%u(nx+1, j) = u_r
             ! Mise à jour de la composante de vitesse discrète sur face décalée MAC
             fields%u(nx+1, j) = u_r
             ! Mise à jour de la composante de vitesse discrète sur face décalée MAC
             fields%v(nx+1, j) = 2.0_dp * v_r - fields%v(nx, j)
          ! Branche conditionnelle alternative
          else if (cfg%bc_right%btype == BC_FREESLIP) then
             ! Mise à jour de la composante de vitesse discrète sur face décalée MAC
             if (.not. homogeneous) fields%u(nx+1, j) = 0.0_dp
             ! Mise à jour de la composante de vitesse discrète sur face décalée MAC
             fields%u(nx+1, j) = 0.0_dp
             ! Mise à jour de la composante de vitesse discrète sur face décalée MAC
             fields%v(nx+1, j) = fields%v(nx, j)
          ! Branche conditionnelle alternative
          else if (cfg%bc_right%btype == BC_NEUMANN) then
             ! Mise à jour de la composante de vitesse discrète sur face décalée MAC
             fields%u(nx+1, j) = fields%u(nx, j)
             ! Mise à jour de la composante de vitesse discrète sur face décalée MAC
             fields%v(nx+1, j) = fields%v(nx, j)
          end if
       end do
    end if

    ! Bord Sud physique (y = 0)
    if (mesh%neighbor_south == MPI_PROC_NULL .and. cfg%bc_bottom%btype /= BC_PERIODIC) then
       ! Boucle spatiale selon l'axe horizontal X (accès contigu colonne-major)
       do i = 1, nx
          ! Évaluation de la condition logique
          if (cfg%bc_bottom%btype == BC_DIRICHLET) then
             ! Mise à jour de la composante de vitesse discrète sur face décalée MAC
             if (.not. homogeneous) fields%v(i, 1) = v_b
             ! Mise à jour de la composante de vitesse discrète sur face décalée MAC
             fields%v(i, 0) = v_b
             ! Mise à jour de la composante de vitesse discrète sur face décalée MAC
             fields%u(i, 0) = 2.0_dp * u_b - fields%u(i, 1)
          ! Branche conditionnelle alternative
          else if (cfg%bc_bottom%btype == BC_FREESLIP) then
             ! Mise à jour de la composante de vitesse discrète sur face décalée MAC
             if (.not. homogeneous) fields%v(i, 1) = 0.0_dp
             ! Mise à jour de la composante de vitesse discrète sur face décalée MAC
             fields%v(i, 0) = 0.0_dp
             ! Mise à jour de la composante de vitesse discrète sur face décalée MAC
             fields%u(i, 0) = fields%u(i, 1)
          ! Branche conditionnelle alternative
          else if (cfg%bc_bottom%btype == BC_NEUMANN) then
             ! Mise à jour de la composante de vitesse discrète sur face décalée MAC
             fields%u(i, 0) = fields%u(i, 1)
             ! Mise à jour de la composante de vitesse discrète sur face décalée MAC
             fields%v(i, 0) = fields%v(i, 1)
          end if
       end do
    end if

    ! Bord Nord physique (y = Ly)
    if (mesh%neighbor_north == MPI_PROC_NULL .and. cfg%bc_top%btype /= BC_PERIODIC) then
       ! Traitement du profil d'injecteur JET_INLET (Cas 3.3 Fluids 2021)
       if (trim(cfg%bc_top%profile) == "JET_INLET") then
          xc_top = merge(cfg%bc_top%jet_center, 0.5_dp * mesh%lx, cfg%bc_top%jet_center > 0.0_dp)
          half_w_top = merge(0.5_dp * cfg%bc_top%jet_width, 0.04_dp, cfg%bc_top%jet_width > 0.0_dp)
          y = 0.0_dp
          do i = 1, mesh%nx_global
             if (abs((real(i, dp) - 0.5_dp) * mesh%dx - xc_top) <= half_w_top) y = y + mesh%dx
          end do
          if (y < 1.0e-12_dp) y = 2.0_dp * half_w_top
          v_out_top = -v_t * (2.0_dp * half_w_top) / (mesh%lx - y)
          do i = 1, nx
             x_c = x_center(mesh, i)
             if (abs(x_c - xc_top) <= half_w_top) then
                if (.not. homogeneous) fields%v(i, ny+1) = v_t * (2.0_dp * half_w_top) / y
             else
                if (.not. homogeneous) fields%v(i, ny+1) = v_out_top
             end if
             ! Condition tangentielle symétrique évaluée à l'abscisse exacte de la face u
             if (abs(x_face(mesh, i) - xc_top) <= half_w_top) then
                fields%u(i, ny+1) = 0.0_dp
             else
                fields%u(i, ny+1) = fields%u(i, ny)
             end if
          end do
          fields%u(nx+1, ny+1) = fields%u(nx+1, ny)
       else
          ! Boucle spatiale selon l'axe horizontal X (accès contigu colonne-major)
          do i = 1, nx
             ! Évaluation de la condition logique
             if (cfg%bc_top%btype == BC_DIRICHLET) then
                ! Mise à jour de la composante de vitesse discrète sur face décalée MAC
                if (.not. homogeneous) fields%v(i, ny+1) = v_t
                fields%u(i, ny+1) = 2.0_dp * u_t - fields%u(i, ny)
             ! Branche conditionnelle alternative
             else if (cfg%bc_top%btype == BC_FREESLIP) then
                ! Mise à jour de la composante de vitesse discrète sur face décalée MAC
                if (.not. homogeneous) fields%v(i, ny+1) = 0.0_dp
                fields%u(i, ny+1) = fields%u(i, ny)
             ! Branche conditionnelle alternative
             else if (cfg%bc_top%btype == BC_NEUMANN) then
                ! Mise à jour de la composante de vitesse discrète sur face décalée MAC
                fields%u(i, ny+1) = fields%u(i, ny)
                fields%v(i, ny+1) = fields%v(i, ny)
             end if
          end do
       end if
    end if
  end subroutine apply_velocity_bc

  ! [Procédure/Module] subroutine apply_thermal_bc
  subroutine apply_thermal_bc(fields, mesh, cfg)
    ! Déclaration de variable locale / paramètre de calcul
    type(flow_fields_t), intent(inout) :: fields
    ! Déclaration de variable locale / paramètre de calcul
    type(mesh_2d_t), intent(in) :: mesh
    ! Déclaration de variable locale / paramètre de calcul
    type(config_t), intent(in) :: cfg

    ! Déclaration de variable locale / paramètre de calcul
    integer :: nx, ny, i, j
    ! Affectation et calcul numérique
    nx = mesh%nx_local
    ! Affectation et calcul numérique
    ny = mesh%ny_local

    ! Périodicité thermique
    if (cfg%bc_left%thermal_type == BC_PERIODIC .and. mesh%px == 1) then
       ! Évolution du champ scalaire de température
       fields%T(0, 1:ny)    = fields%T(nx, 1:ny)
       ! Évolution du champ scalaire de température
       fields%T(nx+1, 1:ny) = fields%T(1, 1:ny)
    end if
    ! Évaluation de la condition logique
    if (cfg%bc_bottom%thermal_type == BC_PERIODIC .and. mesh%py == 1) then
       ! Évolution du champ scalaire de température
       fields%T(1:nx, 0)    = fields%T(1:nx, ny)
       ! Évolution du champ scalaire de température
       fields%T(1:nx, ny+1) = fields%T(1:nx, 1)
    end if

    ! Bord Ouest physique (x = 0)
    if (mesh%neighbor_west == MPI_PROC_NULL) then
       ! Boucle spatiale sur les cellules locales selon l'axe vertical Y (ordonnées)
       do j = 1, ny
          ! Évaluation de la condition logique
          if (cfg%bc_left%thermal_type == THERMAL_DIRICHLET) then
             ! Évolution du champ scalaire de température
             fields%T(0, j) = 2.0_dp * cfg%bc_left%T - fields%T(1, j)
          ! Branche par défaut si condition non satisfaite
          else
             ! Évolution du champ scalaire de température
             fields%T(0, j) = fields%T(1, j)
          end if
       end do
    end if

    ! Bord Est physique (x = Lx)
    if (mesh%neighbor_east == MPI_PROC_NULL) then
       ! Boucle spatiale sur les cellules locales selon l'axe vertical Y (ordonnées)
       do j = 1, ny
          ! Évaluation de la condition logique
          if (cfg%bc_right%thermal_type == THERMAL_DIRICHLET) then
             ! Évolution du champ scalaire de température
             fields%T(nx+1, j) = 2.0_dp * cfg%bc_right%T - fields%T(nx, j)
          ! Branche par défaut si condition non satisfaite
          else
             ! Évolution du champ scalaire de température
             fields%T(nx+1, j) = fields%T(nx, j)
          end if
       end do
    end if

    ! Bord Sud physique (y = 0)
    if (mesh%neighbor_south == MPI_PROC_NULL) then
       ! Boucle spatiale selon l'axe horizontal X (accès contigu colonne-major)
       do i = 1, nx
          ! Évaluation de la condition logique
          if (cfg%bc_bottom%thermal_type == THERMAL_DIRICHLET) then
             ! Évolution du champ scalaire de température
             fields%T(i, 0) = 2.0_dp * cfg%bc_bottom%T - fields%T(i, 1)
          ! Branche par défaut si condition non satisfaite
          else
             ! Évolution du champ scalaire de température
             fields%T(i, 0) = fields%T(i, 1)
          end if
       end do
    end if

    ! Bord Nord physique (y = Ly)
    if (mesh%neighbor_north == MPI_PROC_NULL) then
       ! Boucle spatiale selon l'axe horizontal X (accès contigu colonne-major)
       do i = 1, nx
          ! Évaluation de la condition logique
          if (cfg%bc_top%thermal_type == THERMAL_DIRICHLET) then
             ! Évolution du champ scalaire de température
             fields%T(i, ny+1) = 2.0_dp * cfg%bc_top%T - fields%T(i, ny)
          ! Branche par défaut si condition non satisfaite
          else
             ! Évolution du champ scalaire de température
             fields%T(i, ny+1) = fields%T(i, ny)
          end if
       end do
    end if
  end subroutine apply_thermal_bc

  ! [Procédure/Module] subroutine apply_vof_bc
  subroutine apply_vof_bc(fields, mesh, cfg)
    ! Déclaration de variable locale / paramètre de calcul
    type(flow_fields_t), intent(inout) :: fields
    ! Déclaration de variable locale / paramètre de calcul
    type(mesh_2d_t), intent(in) :: mesh
    ! Déclaration de variable locale / paramètre de calcul
    type(config_t), intent(in) :: cfg

    ! Déclaration de variable locale / paramètre de calcul
    integer :: nx, ny, i
    real(dp) :: x_c, xc_top, half_w_top
    ! Affectation et calcul numérique
    nx = mesh%nx_local
    ! Affectation et calcul numérique
    ny = mesh%ny_local

    ! Affectation et calcul numérique
    if (mesh%neighbor_west == MPI_PROC_NULL)  fields%C(0, 1:ny)    = fields%C(1, 1:ny)
    ! Affectation et calcul numérique
    if (mesh%neighbor_east == MPI_PROC_NULL)  fields%C(nx+1, 1:ny) = fields%C(nx, 1:ny)
    ! Affectation et calcul numérique
    if (mesh%neighbor_south == MPI_PROC_NULL) fields%C(1:nx, 0)    = fields%C(1:nx, 1)
    ! Affectation et calcul numérique pour bord Nord
    if (mesh%neighbor_north == MPI_PROC_NULL) then
       if (trim(cfg%bc_top%profile) == "JET_INLET") then
          xc_top = merge(cfg%bc_top%jet_center, 0.5_dp * mesh%lx, cfg%bc_top%jet_center > 0.0_dp)
          half_w_top = merge(0.5_dp * cfg%bc_top%jet_width, 0.04_dp, cfg%bc_top%jet_width > 0.0_dp)
          do i = 1, nx
             x_c = x_center(mesh, i)
             if (abs(x_c - xc_top) <= half_w_top) then
                fields%C(i, ny+1) = 0.0_dp ! Injection continue de liquide (Fluide 2)
             else
                fields%C(i, ny+1) = fields%C(i, ny) ! Évacuation libre du gaz
             end if
          end do
       else
          fields%C(1:nx, ny+1) = fields%C(1:nx, ny)
       end if
    end if
  end subroutine apply_vof_bc

  ! Déclaration de variable locale / paramètre de calcul
  real(dp) function global_max_val(local_val, comm)
    ! Déclaration de variable locale / paramètre de calcul
    real(dp), intent(in) :: local_val
    ! Déclaration de variable locale / paramètre de calcul
    integer, intent(in) :: comm
    ! Déclaration de variable locale / paramètre de calcul
    real(dp) :: res
    ! Déclaration de variable locale / paramètre de calcul
    integer :: ierr
    ! Appel à la bibliothèque de communication parallèle MPI
    call MPI_Allreduce(local_val, res, 1, MPI_DOUBLE_PRECISION, MPI_MAX, comm, ierr)
    ! Affectation et calcul numérique
    global_max_val = res
  end function global_max_val

  ! Déclaration de variable locale / paramètre de calcul
  real(dp) function global_min_val(local_val, comm)
    ! Déclaration de variable locale / paramètre de calcul
    real(dp), intent(in) :: local_val
    ! Déclaration de variable locale / paramètre de calcul
    integer, intent(in) :: comm
    ! Déclaration de variable locale / paramètre de calcul
    real(dp) :: res
    ! Déclaration de variable locale / paramètre de calcul
    integer :: ierr
    ! Appel à la bibliothèque de communication parallèle MPI
    call MPI_Allreduce(local_val, res, 1, MPI_DOUBLE_PRECISION, MPI_MIN, comm, ierr)
    ! Affectation et calcul numérique
    global_min_val = res
  end function global_min_val

  ! Déclaration de variable locale / paramètre de calcul
  real(dp) function global_sum_val(local_val, comm)
    ! Déclaration de variable locale / paramètre de calcul
    real(dp), intent(in) :: local_val
    ! Déclaration de variable locale / paramètre de calcul
    integer, intent(in) :: comm
    ! Déclaration de variable locale / paramètre de calcul
    real(dp) :: res
    ! Déclaration de variable locale / paramètre de calcul
    integer :: ierr
    ! Appel à la bibliothèque de communication parallèle MPI
    call MPI_Allreduce(local_val, res, 1, MPI_DOUBLE_PRECISION, MPI_SUM, comm, ierr)
    ! Affectation et calcul numérique
    global_sum_val = res
  end function global_sum_val

  ! [Procédure/Module] subroutine sync_u
  subroutine sync_u(mesh, u)
    ! Déclaration de variable locale / paramètre de calcul
    type(mesh_2d_t), intent(in) :: mesh
    ! Déclaration de variable locale / paramètre de calcul
    real(dp), intent(inout) :: u(0:mesh%nx_local+1, 0:mesh%ny_local+1)
    ! Appel de sous-programme modulaire
    call exchange_halos_scalar(u, mesh)
  end subroutine sync_u

  ! [Procédure/Module] subroutine sync_v
  subroutine sync_v(mesh, v)
    ! Déclaration de variable locale / paramètre de calcul
    type(mesh_2d_t), intent(in) :: mesh
    ! Déclaration de variable locale / paramètre de calcul
    real(dp), intent(inout) :: v(0:mesh%nx_local+1, 0:mesh%ny_local+1)
    ! Appel de sous-programme modulaire
    call exchange_halos_scalar(v, mesh)
  end subroutine sync_v

  ! [Procédure/Module] subroutine sync_cell_scalar
  subroutine sync_cell_scalar(mesh, s)
    ! Déclaration de variable locale / paramètre de calcul
    type(mesh_2d_t), intent(in) :: mesh
    ! Déclaration de variable locale / paramètre de calcul
    real(dp), intent(inout) :: s(0:mesh%nx_local+1, 0:mesh%ny_local+1)
    ! Appel de sous-programme modulaire
    call exchange_halos_scalar(s, mesh)
  end subroutine sync_cell_scalar

  ! [Procédure/Module] subroutine apply_boundary_conditions
  subroutine apply_boundary_conditions(mesh, cfg, fields, homogeneous)
    ! Déclaration de variable locale / paramètre de calcul
    type(mesh_2d_t), intent(in) :: mesh
    ! Déclaration de variable locale / paramètre de calcul
    type(config_t), intent(in) :: cfg
    ! Déclaration de variable locale / paramètre de calcul
    type(flow_fields_t), intent(inout) :: fields
    ! Déclaration de variable locale / paramètre de calcul
    logical, intent(in) :: homogeneous
    ! Appel de sous-programme modulaire
    call apply_velocity_bc(fields, mesh, cfg, homogeneous)
  end subroutine apply_boundary_conditions

  ! [Procédure/Module] subroutine apply_thermal_boundary_conditions
  subroutine apply_thermal_boundary_conditions(mesh, cfg, fields)
    ! Déclaration de variable locale / paramètre de calcul
    type(mesh_2d_t), intent(in) :: mesh
    ! Déclaration de variable locale / paramètre de calcul
    type(config_t), intent(in) :: cfg
    ! Déclaration de variable locale / paramètre de calcul
    type(flow_fields_t), intent(inout) :: fields
    ! Appel de sous-programme modulaire
    call apply_thermal_bc(fields, mesh, cfg)
  end subroutine apply_thermal_boundary_conditions

  ! Déclaration de variable locale / paramètre de calcul
  real(dp) function global_max(mesh, local_val)
    ! Déclaration de variable locale / paramètre de calcul
    type(mesh_2d_t), intent(in) :: mesh
    ! Déclaration de variable locale / paramètre de calcul
    real(dp), intent(in) :: local_val
    ! Affectation et calcul numérique
    global_max = global_max_val(local_val, mesh%comm)
  end function global_max

  ! Déclaration de variable locale / paramètre de calcul
  real(dp) function global_min(mesh, local_val)
    ! Déclaration de variable locale / paramètre de calcul
    type(mesh_2d_t), intent(in) :: mesh
    ! Déclaration de variable locale / paramètre de calcul
    real(dp), intent(in) :: local_val
    ! Affectation et calcul numérique
    global_min = global_min_val(local_val, mesh%comm)
  end function global_min

  ! Déclaration de variable locale / paramètre de calcul
  real(dp) function global_sum(mesh, local_val)
    ! Déclaration de variable locale / paramètre de calcul
    type(mesh_2d_t), intent(in) :: mesh
    ! Déclaration de variable locale / paramètre de calcul
    real(dp), intent(in) :: local_val
    ! Affectation et calcul numérique
    global_sum = global_sum_val(local_val, mesh%comm)
  end function global_sum

  ! Déclaration de variable locale / paramètre de calcul
  real(dp) function global_dot_product(mesh, x, y, n)
    ! Déclaration de variable locale / paramètre de calcul
    type(mesh_2d_t), intent(in) :: mesh
    ! Déclaration de variable locale / paramètre de calcul
    real(dp), intent(in) :: x(:), y(:)
    ! Déclaration de variable locale / paramètre de calcul
    integer, intent(in) :: n
    ! Déclaration de variable locale / paramètre de calcul
    real(dp) :: local_dot
    ! Affectation et calcul numérique
    local_dot = dot_product(x(1:n), y(1:n))
    ! Affectation et calcul numérique
    global_dot_product = global_sum_val(local_dot, mesh%comm)
  end function global_dot_product

  ! Déclaration de variable locale / paramètre de calcul
  real(dp) function global_norm2(mesh, x, n)
    ! Déclaration de variable locale / paramètre de calcul
    type(mesh_2d_t), intent(in) :: mesh
    ! Déclaration de variable locale / paramètre de calcul
    real(dp), intent(in) :: x(:)
    ! Déclaration de variable locale / paramètre de calcul
    integer, intent(in) :: n
    ! Évaluation de la norme du résidu ou de l'erreur numérique
    global_norm2 = sqrt(global_dot_product(mesh, x, x, n))
  end function global_norm2

end module module_mpi_comm

\end{lstlisting}

\chapter{Matrices Creuses CSR, Préconditionneurs et Solveurs de Krylov en Fortran 2018 (module_linear_solvers.f90)}
\section{Implémentation Vectorisée des Solveurs de Krylov dans \texttt{module\_linear\_solvers.f90}}
Le module \texttt{src/module\_linear\_solvers.f90} (892 lignes) implémente la structure de matrice creuse \texttt{csr\_matrix\_t} (à indiçage Fortran à base 1 : $\texttt{row\_ptr}(1) = 1$, $\texttt{row\_ptr}(n+1) = \text{nnz} + 1$) ainsi que les solveurs :
\begin{itemize}
  \item \texttt{solve\_bicgstab2} : implémentation complète de l'algorithme de Sleijpen \& Fokkema (1993) avec préconditionnement à droite ou à gauche (Jacobi, SSOR, ILU(0)) ;
  \item \texttt{solve\_bicgstab} : gradient biconjugué stabilisé standard de Van der Vorst ;
  \item \texttt{solve\_pcg} : gradient conjugué préconditionné pour l'opérateur symétrique défini positif du Laplacien (Cas 3.1) ;
  \item \texttt{solve\_gmres} : méthode de minimisation résiduelle généralisée avec sous-espace de Krylov redémarré et rotations de Givens.
\end{itemize}

\section{Répertoire des Types Dérivés, Sous-Routines et Fonctions de \texttt{src/module\_linear\_solvers.f90}}
\begin{center}\small
\begin{tabular}{clll}
\toprule
\textbf{Ligne} & \textbf{Nature} & \textbf{Identifiant Fortran 2018} & \textbf{Interface / Arguments} \\
\midrule
10 & module & \texttt{module\_linear\_solvers} & Unité modulaire Fortran 2018 \\
30 & subroutine & \texttt{apply\_preconditioner} & Arguments : (ns, cfg, src, dst) \\
153 & subroutine & \texttt{apply\_PA} & Arguments : (ns, mesh, cfg, fields\_old, dt, src, dst) \\
191 & subroutine & \texttt{solve\_bicgstab2} & Arguments : (ns, mesh, cfg, fields\_old, dt, b, x, iters) \\
402 & subroutine & \texttt{solve\_bicgstab} & Arguments : (ns, mesh, cfg, fields\_old, dt, b, x, iters) \\
551 & subroutine & \texttt{solve\_cg} & Arguments : (ns, mesh, cfg, fields\_old, dt, b, x, iters) \\
659 & subroutine & \texttt{solve\_gmres} & Arguments : (ns, mesh, cfg, fields\_old, dt, b, x, iters) \\
855 & subroutine & \texttt{solve\_linear\_system} & Arguments : (ns, mesh, cfg, fields\_old, dt, b, x, iters) \\
\bottomrule
\end{tabular}
\end{center}

\section{Code Source Intégral Annoté : \texttt{src/module\_linear\_solvers.f90}}
Le listing ci-dessous présente l'intégralité des 892 lignes du fichier \texttt{src/module\_linear\_solvers.f90} de l'implémentation Fortran 2018 :

\begin{lstlisting}[language={[95]Fortran}]
!===============================================================================
! Module: module_linear_solvers
! Description: Solveurs de Krylov parallèles pour COVOF 2D en Fortran 2018 :
!              - BiCGSTAB(2) (Sleijpen & van der Vorst 1993)
!              - BiCGSTAB standard (van der Vorst 1992)
!              - Conjugate Gradient (PCG - Hestenes & Stiefel 1952)
!              - GMRES(m) redémarré (Saad & Schultz 1986)
!              - Préconditionneurs Jacobi diagonal et SSOR
!===============================================================================
module module_linear_solvers
  use module_precision, only: dp
  use module_constants
  use module_config, only: config_t
  use module_mesh, only: mesh_2d_t
  use module_fields, only: flow_fields_t
  use module_mpi_comm, only: global_dot_product, global_norm2
  use module_navier_stokes, only: navier_stokes_t
  implicit none
  private

  public :: solve_linear_system
  public :: solve_bicgstab2, solve_bicgstab, solve_cg, solve_gmres
  public :: apply_preconditioner

contains

  !-----------------------------------------------------------------------------
  ! Application du préconditionneur : dst = P^{-1} src
  !-----------------------------------------------------------------------------
  subroutine apply_preconditioner(ns, cfg, src, dst)
    ! Déclaration de variable locale / paramètre de calcul
    type(navier_stokes_t), intent(in) :: ns
    ! Déclaration de variable locale / paramètre de calcul
    type(config_t), intent(in)        :: cfg
    ! Déclaration de variable locale / paramètre de calcul
    real(dp), intent(in)              :: src(:)
    ! Déclaration de variable locale / paramètre de calcul
    real(dp), intent(out)             :: dst(:)

    ! Déclaration de variable locale / paramètre de calcul
    integer :: i, nu, nv, np, n
    ! Déclaration de variable locale / paramètre de calcul
    real(dp) :: d

    ! Affectation et calcul numérique
    n = ns%n_owned_total
    ! Affectation et calcul numérique
    nu = ns%n_owned_u
    ! Affectation et calcul numérique
    nv = ns%n_owned_v
    ! Affectation et calcul numérique
    np = ns%n_owned_p

    ! Évaluation de la condition logique
    if (.not. cfg%use_precond) then
      ! Affectation et calcul numérique
      dst(1:n) = src(1:n)
      ! Retour à la procédure appelante
      return
    end if

    select case (cfg%precond_type)
    case (PRECOND_JACOBI)
      ! Boucle spatiale selon l'axe horizontal X (accès contigu colonne-major)
      do i = 1, nu
        ! Affectation et calcul numérique
        d = ns%diag_u(i)
        ! Évaluation de la condition logique
        if (abs(d) > 1.0e-15_dp) then
          ! Affectation et calcul numérique
          dst(i) = src(i) / d
        ! Branche par défaut si condition non satisfaite
        else
          ! Affectation et calcul numérique
          dst(i) = src(i)
        end if
      end do

      ! Boucle spatiale selon l'axe horizontal X (accès contigu colonne-major)
      do i = 1, nv
        ! Affectation et calcul numérique
        d = ns%diag_v(i)
        ! Évaluation de la condition logique
        if (abs(d) > 1.0e-15_dp) then
          ! Affectation et calcul numérique
          dst(nu + i) = src(nu + i) / d
        ! Branche par défaut si condition non satisfaite
        else
          ! Affectation et calcul numérique
          dst(nu + i) = src(nu + i)
        end if
      end do

      ! Évaluation de la condition logique
      if (cfg%solver_type == SOLVER_MONOLITHIC) then
        ! Boucle spatiale selon l'axe horizontal X (accès contigu colonne-major)
        do i = 1, np
          ! Affectation et calcul numérique
          d = ns%diag_p(i)
          ! Évaluation de la condition logique
          if (abs(d) > 1.0e-15_dp) then
            ! Affectation et calcul numérique
            dst(nu + nv + i) = src(nu + nv + i) / d
          ! Branche par défaut si condition non satisfaite
          else
            ! Affectation et calcul numérique
            dst(nu + nv + i) = src(nu + nv + i)
          end if
        end do
      end if

    case (PRECOND_SSOR)
      ! Balayage SSOR avant / arrière avec omega = 1.0
      do i = 1, nu
        ! Affectation et calcul numérique
        d = ns%diag_u(i)
        ! Affectation et calcul numérique
        if (abs(d) < 1.0e-15_dp) d = 1.0_dp
        ! Affectation et calcul numérique
        dst(i) = src(i) / d
      end do
      ! Boucle spatiale selon l'axe horizontal X (accès contigu colonne-major)
      do i = 1, nv
        ! Affectation et calcul numérique
        d = ns%diag_v(i)
        ! Affectation et calcul numérique
        if (abs(d) < 1.0e-15_dp) d = 1.0_dp
        ! Affectation et calcul numérique
        dst(nu + i) = src(nu + i) / d
      end do
      ! Évaluation de la condition logique
      if (cfg%solver_type == SOLVER_MONOLITHIC) then
        ! Boucle spatiale selon l'axe horizontal X (accès contigu colonne-major)
        do i = 1, np
          ! Affectation et calcul numérique
          d = ns%diag_p(i)
          ! Affectation et calcul numérique
          if (abs(d) < 1.0e-15_dp) d = 1.0_dp
          ! Affectation et calcul numérique
          dst(nu + nv + i) = src(nu + nv + i) / d
        end do
      end if

    case default
      ! Affectation et calcul numérique
      dst(1:n) = src(1:n)
    end select
  end subroutine apply_preconditioner

  !-----------------------------------------------------------------------------
  ! Application de l'opérateur préconditionné : dst = P^{-1} * A * src
  !-----------------------------------------------------------------------------
  subroutine apply_PA(ns, mesh, cfg, fields_old, dt, src, dst)
    ! Déclaration de variable locale / paramètre de calcul
    type(navier_stokes_t), intent(in) :: ns
    ! Déclaration de variable locale / paramètre de calcul
    type(mesh_2d_t), intent(in)        :: mesh
    ! Déclaration de variable locale / paramètre de calcul
    type(config_t), intent(in)         :: cfg
    ! Déclaration de variable locale / paramètre de calcul
    type(flow_fields_t), intent(in)    :: fields_old
    ! Déclaration de variable locale / paramètre de calcul
    real(dp), intent(in)               :: dt
    ! Déclaration de variable locale / paramètre de calcul
    real(dp), intent(in)               :: src(:)
    ! Déclaration de variable locale / paramètre de calcul
    real(dp), intent(out)              :: dst(:)

    ! Déclaration de variable locale / paramètre de calcul
    real(dp), allocatable :: tmp(:)
    ! Déclaration de variable locale / paramètre de calcul
    integer :: n

    ! Affectation et calcul numérique
    n = ns%n_owned_total
    ! Allocation dynamique de tableau contigu en mémoire
    allocate(tmp(n))

    ! Appel de sous-programme modulaire
    call ns%apply_system_matrix(mesh, cfg, fields_old, dt, src, tmp)
    ! Appel de sous-programme modulaire
    call apply_preconditioner(ns, cfg, tmp, dst)

    ! Libération de la mémoire allouée
    deallocate(tmp)
  end subroutine apply_PA

  !-----------------------------------------------------------------------------
  ! 1. Solveur BiCGSTAB(2) - Sleijpen & van der Vorst (1993)
  !-----------------------------------------------------------------------------
  subroutine solve_bicgstab2(ns, mesh, cfg, fields_old, dt, b, x, iters)
    ! Déclaration de variable locale / paramètre de calcul
    type(navier_stokes_t), intent(in) :: ns
    ! Déclaration de variable locale / paramètre de calcul
    type(mesh_2d_t), intent(in)        :: mesh
    ! Déclaration de variable locale / paramètre de calcul
    type(config_t), intent(in)         :: cfg
    ! Déclaration de variable locale / paramètre de calcul
    type(flow_fields_t), intent(in)    :: fields_old
    ! Déclaration de variable locale / paramètre de calcul
    real(dp), intent(in)               :: dt
    ! Déclaration de variable locale / paramètre de calcul
    real(dp), intent(in)               :: b(:)
    ! Déclaration de variable locale / paramètre de calcul
    real(dp), intent(inout)            :: x(:)
    ! Déclaration de variable locale / paramètre de calcul
    integer, intent(out)               :: iters

    ! Déclaration de variable locale / paramètre de calcul
    integer :: n, max_iter, iter
    ! Déclaration de variable locale / paramètre de calcul
    real(dp) :: tol, r0_norm, r_norm, rho_0, rho_1, alpha, omega, beta, gamma
    ! Déclaration de variable locale / paramètre de calcul
    real(dp) :: sigma11, sigma12, sigma22, rhs1, rhs2, det, s11_s22, gamma1, gamma2
    ! Déclaration de variable locale / paramètre de calcul
    real(dp), allocatable :: b_hat(:), r_tilde(:)
    ! Déclaration de variable locale / paramètre de calcul
    real(dp), allocatable :: r0(:), r1(:), r2(:)
    ! Déclaration de variable locale / paramètre de calcul
    real(dp), allocatable :: u0(:), u1(:), u2(:)

    ! Affectation et calcul numérique
    n = ns%n_owned_total
    ! Affectation et calcul numérique
    tol = cfg%tol
    ! Affectation et calcul numérique
    max_iter = cfg%max_iter
    ! Affectation et calcul numérique
    iters = 0

    ! Allocation dynamique de tableau contigu en mémoire
    allocate(b_hat(n), r_tilde(n))
    ! Allocation dynamique de tableau contigu en mémoire
    allocate(r0(n), r1(n), r2(n))
    ! Allocation dynamique de tableau contigu en mémoire
    allocate(u0(n), u1(n), u2(n))

    ! Affectation et calcul numérique
    u0 = 0.0_dp
    ! Affectation et calcul numérique
    u1 = 0.0_dp
    ! Affectation et calcul numérique
    u2 = 0.0_dp
    ! Affectation et calcul numérique
    r1 = 0.0_dp
    ! Affectation et calcul numérique
    r2 = 0.0_dp

    ! Second membre préconditionné b_hat = P^{-1} b
    call apply_preconditioner(ns, cfg, b, b_hat)

    ! Résidu initial r0 = b_hat - P^{-1} A x0
    call apply_PA(ns, mesh, cfg, fields_old, dt, x, r0)
    ! Affectation et calcul numérique
    r0(1:n) = b_hat(1:n) - r0(1:n)

    ! Évaluation de la norme du résidu ou de l'erreur numérique
    r0_norm = global_norm2(mesh, r0, n)
    ! Évaluation de la condition logique
    if (r0_norm < 1.0e-12_dp) then
      ! Libération de la mémoire allouée
      deallocate(b_hat, r_tilde, r0, r1, r2, u0, u1, u2)
      ! Retour à la procédure appelante
      return
    end if

    ! Affectation et calcul numérique
    r_tilde(1:n) = r0(1:n)
    ! Calcul de la masse volumique locale ou interpolée aux faces
    rho_0 = 1.0_dp
    ! Affectation et calcul numérique
    alpha = 0.0_dp
    ! Affectation et calcul numérique
    omega = 1.0_dp
    ! Affectation et calcul numérique
    iter = 0

    ! Boucle d'itération numérique
    do while (iter < max_iter)
      ! Affectation et calcul numérique
      iter = iter + 2
      ! Calcul de la masse volumique locale ou interpolée aux faces
      rho_0 = -omega * rho_0

      ! Demi-pas 1 (j = 0)
      rho_1 = global_dot_product(mesh, r_tilde, r0, n)
      if (abs(rho_0) < 1.0e-25_dp) exit
      ! Affectation et calcul numérique
      beta = alpha * (rho_1 / rho_0)
      ! Calcul de la masse volumique locale ou interpolée aux faces
      rho_0 = rho_1

      ! Affectation et calcul numérique
      u0(1:n) = r0(1:n) - beta * u0(1:n)
      ! Appel de sous-programme modulaire
      call apply_PA(ns, mesh, cfg, fields_old, dt, u0, u1)

      ! Affectation et calcul numérique
      gamma = global_dot_product(mesh, r_tilde, u1, n)
      if (abs(gamma) < 1.0e-25_dp) exit
      ! Affectation et calcul numérique
      alpha = rho_0 / gamma

      ! Affectation et calcul numérique
      r0(1:n) = r0(1:n) - alpha * u1(1:n)
      ! Affectation et calcul numérique
      x(1:n)  = x(1:n)  + alpha * u0(1:n)
      ! Appel de sous-programme modulaire
      call apply_PA(ns, mesh, cfg, fields_old, dt, r0, r1)

      ! Demi-pas 2 (j = 1)
      rho_1 = global_dot_product(mesh, r_tilde, r1, n)
      if (abs(rho_0) < 1.0e-25_dp) exit
      ! Affectation et calcul numérique
      beta = alpha * (rho_1 / rho_0)
      ! Calcul de la masse volumique locale ou interpolée aux faces
      rho_0 = rho_1

      ! Affectation et calcul numérique
      u0(1:n) = r0(1:n) - beta * u0(1:n)
      ! Affectation et calcul numérique
      u1(1:n) = r1(1:n) - beta * u1(1:n)
      ! Appel de sous-programme modulaire
      call apply_PA(ns, mesh, cfg, fields_old, dt, u1, u2)

      ! Affectation et calcul numérique
      gamma = global_dot_product(mesh, r_tilde, u2, n)
      if (abs(gamma) < 1.0e-25_dp) exit
      ! Affectation et calcul numérique
      alpha = rho_0 / gamma

      ! Affectation et calcul numérique
      x(1:n)  = x(1:n)  + alpha * u0(1:n)
      ! Affectation et calcul numérique
      r0(1:n) = r0(1:n) - alpha * u1(1:n)
      ! Affectation et calcul numérique
      r1(1:n) = r1(1:n) - alpha * u2(1:n)
      ! Appel de sous-programme modulaire
      call apply_PA(ns, mesh, cfg, fields_old, dt, r1, r2)

      ! Minimisation Quadratique MR-part
      sigma11 = global_dot_product(mesh, r1, r1, n)
      ! Affectation et calcul numérique
      sigma12 = global_dot_product(mesh, r2, r1, n)
      ! Affectation et calcul numérique
      sigma22 = global_dot_product(mesh, r2, r2, n)
      ! Affectation et calcul numérique
      rhs1    = global_dot_product(mesh, r0, r1, n)
      ! Affectation et calcul numérique
      rhs2    = global_dot_product(mesh, r0, r2, n)

      ! Affectation et calcul numérique
      det = sigma11 * sigma22 - sigma12 * sigma12
      ! Affectation et calcul numérique
      s11_s22 = sigma11 * sigma22

      ! Évaluation de la condition logique
      if (s11_s22 < 1.0e-30_dp .or. abs(det) < 1.0e-10_dp * s11_s22) then
        ! Affectation et calcul numérique
        gamma2 = 0.0_dp
        ! Affectation et calcul numérique
        gamma1 = rhs1 / (sigma11 + 1.0e-30_dp)
      ! Branche par défaut si condition non satisfaite
      else
        ! Affectation et calcul numérique
        gamma1 = (rhs1 * sigma22 - rhs2 * sigma12) / det
        ! Affectation et calcul numérique
        gamma2 = (rhs2 * sigma11 - rhs1 * sigma12) / det
      end if

      ! Affectation et calcul numérique
      x(1:n)  = x(1:n)  + gamma1 * r0(1:n) + gamma2 * r1(1:n)
      ! Affectation et calcul numérique
      r0(1:n) = r0(1:n) - (gamma1 * r1(1:n) + gamma2 * r2(1:n))
      ! Affectation et calcul numérique
      u0(1:n) = u0(1:n) - (gamma1 * u1(1:n) + gamma2 * u2(1:n))
      ! Affectation et calcul numérique
      omega = gamma2

      ! Évaluation de la norme du résidu ou de l'erreur numérique
      r_norm = global_norm2(mesh, r0, n)
      ! Évaluation de la condition logique
      if (r_norm / r0_norm < tol .or. r_norm < 1.0e-13_dp) then
        ! Affectation et calcul numérique
        iters = iter
        ! Libération de la mémoire allouée
        deallocate(b_hat, r_tilde, r0, r1, r2, u0, u1, u2)
        ! Retour à la procédure appelante
        return
      end if
    end do

    ! Affectation et calcul numérique
    iters = iter
    ! Libération de la mémoire allouée
    deallocate(b_hat, r_tilde, r0, r1, r2, u0, u1, u2)
  end subroutine solve_bicgstab2

  !-----------------------------------------------------------------------------
  ! 2. Solveur BiCGSTAB standard - van der Vorst (1992)
  !-----------------------------------------------------------------------------
  subroutine solve_bicgstab(ns, mesh, cfg, fields_old, dt, b, x, iters)
    ! Déclaration de variable locale / paramètre de calcul
    type(navier_stokes_t), intent(in) :: ns
    ! Déclaration de variable locale / paramètre de calcul
    type(mesh_2d_t), intent(in)        :: mesh
    ! Déclaration de variable locale / paramètre de calcul
    type(config_t), intent(in)         :: cfg
    ! Déclaration de variable locale / paramètre de calcul
    type(flow_fields_t), intent(in)    :: fields_old
    ! Déclaration de variable locale / paramètre de calcul
    real(dp), intent(in)               :: dt
    ! Déclaration de variable locale / paramètre de calcul
    real(dp), intent(in)               :: b(:)
    ! Déclaration de variable locale / paramètre de calcul
    real(dp), intent(inout)            :: x(:)
    ! Déclaration de variable locale / paramètre de calcul
    integer, intent(out)               :: iters

    ! Déclaration de variable locale / paramètre de calcul
    integer :: n, max_iter, iter
    ! Déclaration de variable locale / paramètre de calcul
    real(dp) :: tol, r0_norm, r_norm, rho, rho_old, alpha, omega, beta, denom, s_norm, t_dot_t
    ! Déclaration de variable locale / paramètre de calcul
    real(dp), allocatable :: b_hat(:), r(:), r_tilde(:), p(:), v(:), s(:), t(:)

    ! Affectation et calcul numérique
    n = ns%n_owned_total
    ! Affectation et calcul numérique
    tol = cfg%tol
    ! Affectation et calcul numérique
    max_iter = cfg%max_iter
    ! Affectation et calcul numérique
    iters = max_iter

    ! Allocation dynamique de tableau contigu en mémoire
    allocate(b_hat(n), r(n), r_tilde(n), p(n), v(n), s(n), t(n))

    ! Appel de sous-programme modulaire
    call apply_preconditioner(ns, cfg, b, b_hat)
    ! Appel de sous-programme modulaire
    call apply_PA(ns, mesh, cfg, fields_old, dt, x, r)
    ! Affectation et calcul numérique
    r(1:n) = b_hat(1:n) - r(1:n)

    ! Évaluation de la norme du résidu ou de l'erreur numérique
    r0_norm = global_norm2(mesh, r, n)
    ! Évaluation de la condition logique
    if (r0_norm < 1.0e-12_dp) then
      ! Affectation et calcul numérique
      iters = 0
      ! Libération de la mémoire allouée
      deallocate(b_hat, r, r_tilde, p, v, s, t)
      ! Retour à la procédure appelante
      return
    end if

    ! Affectation et calcul numérique
    r_tilde(1:n) = r(1:n)
    ! Calcul de la masse volumique locale ou interpolée aux faces
    rho_old = 1.0_dp
    ! Affectation et calcul numérique
    alpha = 1.0_dp
    ! Affectation et calcul numérique
    omega = 1.0_dp

    ! Boucle d'itération numérique
    do iter = 1, max_iter
      ! Calcul de la masse volumique locale ou interpolée aux faces
      rho = global_dot_product(mesh, r_tilde, r, n)
      if (abs(rho) < 1.0e-30_dp) exit

      ! Évaluation de la condition logique
      if (iter == 1) then
        ! Actualisation de la pression statique au centre de la maille
        p(1:n) = r(1:n)
      ! Branche par défaut si condition non satisfaite
      else
        ! Affectation et calcul numérique
        beta = (rho / rho_old) * (alpha / omega)
        ! Actualisation de la pression statique au centre de la maille
        p(1:n) = r(1:n) + beta * (p(1:n) - omega * v(1:n))
      end if

      ! Appel de sous-programme modulaire
      call apply_PA(ns, mesh, cfg, fields_old, dt, p, v)
      ! Affectation et calcul numérique
      denom = global_dot_product(mesh, r_tilde, v, n)
      if (abs(denom) < 1.0e-30_dp) exit
      ! Affectation et calcul numérique
      alpha = rho / denom

      ! Affectation et calcul numérique
      s(1:n) = r(1:n) - alpha * v(1:n)
      ! Évaluation de la norme du résidu ou de l'erreur numérique
      s_norm = global_norm2(mesh, s, n)
      ! Évaluation de la condition logique
      if (s_norm / r0_norm < tol .or. s_norm < 1.0e-13_dp) then
        ! Affectation et calcul numérique
        x(1:n) = x(1:n) + alpha * p(1:n)
        ! Affectation et calcul numérique
        iters = iter
        ! Libération de la mémoire allouée
        deallocate(b_hat, r, r_tilde, p, v, s, t)
        ! Retour à la procédure appelante
        return
      end if

      ! Appel de sous-programme modulaire
      call apply_PA(ns, mesh, cfg, fields_old, dt, s, t)
      ! Affectation et calcul numérique
      t_dot_t = global_dot_product(mesh, t, t, n)
      ! Évaluation de la condition logique
      if (t_dot_t > 1.0e-30_dp) then
        ! Affectation et calcul numérique
        omega = global_dot_product(mesh, t, s, n) / t_dot_t
      ! Branche par défaut si condition non satisfaite
      else
        ! Affectation et calcul numérique
        omega = 0.0_dp
      end if

      ! Affectation et calcul numérique
      x(1:n) = x(1:n) + alpha * p(1:n) + omega * s(1:n)
      ! Affectation et calcul numérique
      r(1:n) = s(1:n) - omega * t(1:n)

      ! Calcul de la masse volumique locale ou interpolée aux faces
      rho_old = rho
      ! Évaluation de la norme du résidu ou de l'erreur numérique
      r_norm = global_norm2(mesh, r, n)
      ! Évaluation de la condition logique
      if (r_norm / r0_norm < tol .or. r_norm < 1.0e-13_dp) then
        ! Affectation et calcul numérique
        iters = iter
        ! Libération de la mémoire allouée
        deallocate(b_hat, r, r_tilde, p, v, s, t)
        ! Retour à la procédure appelante
        return
      end if
      if (abs(omega) < 1.0e-30_dp) exit
    end do

    ! Libération de la mémoire allouée
    deallocate(b_hat, r, r_tilde, p, v, s, t)
  end subroutine solve_bicgstab

  !-----------------------------------------------------------------------------
  ! 3. Solveur Gradient Conjugué (PCG) - Hestenes & Stiefel (1952)
  !-----------------------------------------------------------------------------
  subroutine solve_cg(ns, mesh, cfg, fields_old, dt, b, x, iters)
    ! Déclaration de variable locale / paramètre de calcul
    type(navier_stokes_t), intent(in) :: ns
    ! Déclaration de variable locale / paramètre de calcul
    type(mesh_2d_t), intent(in)        :: mesh
    ! Déclaration de variable locale / paramètre de calcul
    type(config_t), intent(in)         :: cfg
    ! Déclaration de variable locale / paramètre de calcul
    type(flow_fields_t), intent(in)    :: fields_old
    ! Déclaration de variable locale / paramètre de calcul
    real(dp), intent(in)               :: dt
    ! Déclaration de variable locale / paramètre de calcul
    real(dp), intent(in)               :: b(:)
    ! Déclaration de variable locale / paramètre de calcul
    real(dp), intent(inout)            :: x(:)
    ! Déclaration de variable locale / paramètre de calcul
    integer, intent(out)               :: iters

    ! Déclaration de variable locale / paramètre de calcul
    integer :: n, max_iter, iter
    ! Déclaration de variable locale / paramètre de calcul
    real(dp) :: tol, r0_norm, r_norm, rz_old, rz_new, pAp, alpha, beta
    ! Déclaration de variable locale / paramètre de calcul
    real(dp), allocatable :: Ax(:), r(:), z(:), p(:), Ap(:)

    ! Affectation et calcul numérique
    n = ns%n_owned_total
    ! Affectation et calcul numérique
    tol = cfg%tol
    ! Affectation et calcul numérique
    max_iter = cfg%max_iter
    ! Affectation et calcul numérique
    iters = max_iter

    ! Allocation dynamique de tableau contigu en mémoire
    allocate(Ax(n), r(n), z(n), p(n), Ap(n))

    ! Appel de sous-programme modulaire
    call ns%apply_system_matrix(mesh, cfg, fields_old, dt, x, Ax)
    ! Affectation et calcul numérique
    r(1:n) = b(1:n) - Ax(1:n)

    ! Évaluation de la norme du résidu ou de l'erreur numérique
    r0_norm = global_norm2(mesh, r, n)
    ! Évaluation de la condition logique
    if (r0_norm < 1.0e-12_dp) then
      ! Affectation et calcul numérique
      iters = 0
      ! Libération de la mémoire allouée
      deallocate(Ax, r, z, p, Ap)
      ! Retour à la procédure appelante
      return
    end if

    ! Appel de sous-programme modulaire
    call apply_preconditioner(ns, cfg, r, z)
    ! Actualisation de la pression statique au centre de la maille
    p(1:n) = z(1:n)
    ! Affectation et calcul numérique
    rz_old = global_dot_product(mesh, r, z, n)

    ! Boucle d'itération numérique
    do iter = 1, max_iter
      ! Appel de sous-programme modulaire
      call ns%apply_system_matrix(mesh, cfg, fields_old, dt, p, Ap)
      ! Affectation et calcul numérique
      pAp = global_dot_product(mesh, p, Ap, n)
      if (abs(pAp) < 1.0e-30_dp) exit

      ! Affectation et calcul numérique
      alpha = rz_old / pAp
      ! Affectation et calcul numérique
      x(1:n) = x(1:n) + alpha * p(1:n)
      ! Affectation et calcul numérique
      r(1:n) = r(1:n) - alpha * Ap(1:n)

      ! Évaluation de la norme du résidu ou de l'erreur numérique
      r_norm = global_norm2(mesh, r, n)
      ! Évaluation de la condition logique
      if (r_norm / r0_norm < tol .or. r_norm < 1.0e-13_dp) then
        ! Affectation et calcul numérique
        iters = iter
        ! Libération de la mémoire allouée
        deallocate(Ax, r, z, p, Ap)
        ! Retour à la procédure appelante
        return
      end if

      ! Appel de sous-programme modulaire
      call apply_preconditioner(ns, cfg, r, z)
      ! Affectation et calcul numérique
      rz_new = global_dot_product(mesh, r, z, n)
      ! Affectation et calcul numérique
      beta = rz_new / rz_old
      ! Affectation et calcul numérique
      rz_old = rz_new

      ! Actualisation de la pression statique au centre de la maille
      p(1:n) = z(1:n) + beta * p(1:n)
    end do

    ! Libération de la mémoire allouée
    deallocate(Ax, r, z, p, Ap)
  end subroutine solve_cg

  !-----------------------------------------------------------------------------
  ! 4. Solveur GMRES(m) redémarré - Saad & Schultz (1986)
  !-----------------------------------------------------------------------------
  subroutine solve_gmres(ns, mesh, cfg, fields_old, dt, b, x, iters)
    ! Déclaration de variable locale / paramètre de calcul
    type(navier_stokes_t), intent(in) :: ns
    ! Déclaration de variable locale / paramètre de calcul
    type(mesh_2d_t), intent(in)        :: mesh
    ! Déclaration de variable locale / paramètre de calcul
    type(config_t), intent(in)         :: cfg
    ! Déclaration de variable locale / paramètre de calcul
    type(flow_fields_t), intent(in)    :: fields_old
    ! Déclaration de variable locale / paramètre de calcul
    real(dp), intent(in)               :: dt
    ! Déclaration de variable locale / paramètre de calcul
    real(dp), intent(in)               :: b(:)
    ! Déclaration de variable locale / paramètre de calcul
    real(dp), intent(inout)            :: x(:)
    ! Déclaration de variable locale / paramètre de calcul
    integer, intent(out)               :: iters

    ! Déclaration de variable locale / paramètre de calcul
    integer, parameter :: m = 30
    ! Déclaration de variable locale / paramètre de calcul
    integer :: n, max_iter, total_iters, i, j, k
    ! Déclaration de variable locale / paramètre de calcul
    real(dp) :: tol, r0_norm, beta_val, r_norm, h_val, temp, cs, sn
    ! Déclaration de variable locale / paramètre de calcul
    real(dp), allocatable :: b_hat(:), r(:), w(:), g(:), y(:)
    ! Déclaration de variable locale / paramètre de calcul
    real(dp), allocatable :: V(:, :), H(:, :), c_rot(:), s_rot(:)

    ! Affectation et calcul numérique
    n = ns%n_owned_total
    ! Affectation et calcul numérique
    tol = cfg%tol
    ! Affectation et calcul numérique
    max_iter = cfg%max_iter
    ! Affectation et calcul numérique
    iters = max_iter

    ! Allocation dynamique de tableau contigu en mémoire
    allocate(b_hat(n), r(n), w(n), g(m + 1), y(m))
    ! Allocation dynamique de tableau contigu en mémoire
    allocate(V(n, m + 1), H(m + 1, m), c_rot(m), s_rot(m))

    ! Appel de sous-programme modulaire
    call apply_preconditioner(ns, cfg, b, b_hat)

    ! Affectation et calcul numérique
    total_iters = 0
    ! Évaluation de la norme du résidu ou de l'erreur numérique
    r0_norm = 1.0_dp

    ! Boucle d'itération numérique
    do while (total_iters < max_iter)
      ! Appel de sous-programme modulaire
      call apply_PA(ns, mesh, cfg, fields_old, dt, x, r)
      ! Affectation et calcul numérique
      r(1:n) = b_hat(1:n) - r(1:n)

      ! Affectation et calcul numérique
      beta_val = global_norm2(mesh, r, n)
      ! Évaluation de la condition logique
      if (total_iters == 0) then
        ! Évaluation de la condition logique
        if (beta_val > 1.0e-14_dp) then
          ! Évaluation de la norme du résidu ou de l'erreur numérique
          r0_norm = beta_val
        ! Branche par défaut si condition non satisfaite
        else
          ! Évaluation de la norme du résidu ou de l'erreur numérique
          r0_norm = 1.0_dp
        end if
      end if

      ! Évaluation de la condition logique
      if (beta_val / r0_norm < tol .or. beta_val < 1.0e-13_dp) then
        ! Affectation et calcul numérique
        iters = total_iters
        ! Libération de la mémoire allouée
        deallocate(b_hat, r, w, g, y, V, H, c_rot, s_rot)
        ! Retour à la procédure appelante
        return
      end if

      ! Affectation et calcul numérique
      V(1:n, 1) = r(1:n) / beta_val
      ! Affectation et calcul numérique
      g = 0.0_dp
      ! Affectation et calcul numérique
      g(1) = beta_val
      ! Affectation et calcul numérique
      H = 0.0_dp

      ! Boucle spatiale sur les cellules locales selon l'axe vertical Y (ordonnées)
      do j = 1, min(m, max_iter - total_iters)
        ! Affectation et calcul numérique
        total_iters = total_iters + 1

        ! Appel de sous-programme modulaire
        call apply_PA(ns, mesh, cfg, fields_old, dt, V(1:n, j), w)

        ! Orthogonalisation Modified Gram-Schmidt
        do i = 1, j
          ! Affectation et calcul numérique
          H(i, j) = global_dot_product(mesh, w, V(1:n, i), n)
          ! Affectation et calcul numérique
          w(1:n) = w(1:n) - H(i, j) * V(1:n, i)
        end do

        ! Affectation et calcul numérique
        H(j + 1, j) = global_norm2(mesh, w, n)
        ! Évaluation de la condition logique
        if (H(j + 1, j) > 1.0e-30_dp) then
          ! Affectation et calcul numérique
          V(1:n, j + 1) = w(1:n) / H(j + 1, j)
        end if

        ! Application des rotations de Givens précédentes
        do i = 1, j - 1
          ! Affectation et calcul numérique
          temp = c_rot(i) * H(i, j) + s_rot(i) * H(i + 1, j)
          ! Affectation et calcul numérique
          H(i + 1, j) = -s_rot(i) * H(i, j) + c_rot(i) * H(i + 1, j)
          ! Affectation et calcul numérique
          H(i, j) = temp
        end do

        ! Calcul de la nouvelle rotation de Givens
        if (abs(H(j + 1, j)) < 1.0e-30_dp) then
          ! Affectation et calcul numérique
          c_rot(j) = 1.0_dp
          ! Affectation et calcul numérique
          s_rot(j) = 0.0_dp
        ! Branche conditionnelle alternative
        else if (abs(H(j, j)) > abs(H(j + 1, j))) then
          ! Affectation et calcul numérique
          temp = H(j + 1, j) / H(j, j)
          ! Affectation et calcul numérique
          c_rot(j) = 1.0_dp / sqrt(1.0_dp + temp * temp)
          ! Affectation et calcul numérique
          s_rot(j) = temp * c_rot(j)
        ! Branche par défaut si condition non satisfaite
        else
          ! Affectation et calcul numérique
          temp = H(j, j) / H(j + 1, j)
          ! Affectation et calcul numérique
          s_rot(j) = 1.0_dp / sqrt(1.0_dp + temp * temp)
          ! Affectation et calcul numérique
          c_rot(j) = temp * s_rot(j)
        end if

        ! Affectation et calcul numérique
        H(j, j) = c_rot(j) * H(j, j) + s_rot(j) * H(j + 1, j)
        ! Affectation et calcul numérique
        H(j + 1, j) = 0.0_dp

        ! Affectation et calcul numérique
        g(j + 1) = -s_rot(j) * g(j)
        ! Affectation et calcul numérique
        g(j)     =  c_rot(j) * g(j)

        ! Évaluation de la norme du résidu ou de l'erreur numérique
        r_norm = abs(g(j + 1))
        ! Évaluation de la condition logique
        if (r_norm / r0_norm < tol .or. r_norm < 1.0e-13_dp .or. j == m) then
          ! Résolution du système triangulaire supérieur H * y = g
          do k = j, 1, -1
            ! Affectation et calcul numérique
            y(k) = g(k)
            ! Boucle spatiale selon l'axe horizontal X (accès contigu colonne-major)
            do i = k + 1, j
              ! Affectation et calcul numérique
              y(k) = y(k) - H(k, i) * y(i)
            end do
            ! Affectation et calcul numérique
            if (abs(H(k, k)) > 1.0e-30_dp) y(k) = y(k) / H(k, k)
          end do

          ! Mise à jour de la solution x = x + V * y
          do i = 1, j
            ! Affectation et calcul numérique
            x(1:n) = x(1:n) + y(i) * V(1:n, i)
          end do
          exit
        end if
      end do
    end do

    ! Affectation et calcul numérique
    iters = total_iters
    ! Libération de la mémoire allouée
    deallocate(b_hat, r, w, g, y, V, H, c_rot, s_rot)
  end subroutine solve_gmres

  !-----------------------------------------------------------------------------
  ! Aiguilleur principal de solveur linéaire
  !-----------------------------------------------------------------------------
  subroutine solve_linear_system(ns, mesh, cfg, fields_old, dt, b, x, iters)
    ! Déclaration de variable locale / paramètre de calcul
    type(navier_stokes_t), intent(in) :: ns
    ! Déclaration de variable locale / paramètre de calcul
    type(mesh_2d_t), intent(in)        :: mesh
    ! Déclaration de variable locale / paramètre de calcul
    type(config_t), intent(in)         :: cfg
    ! Déclaration de variable locale / paramètre de calcul
    type(flow_fields_t), intent(in)    :: fields_old
    ! Déclaration de variable locale / paramètre de calcul
    real(dp), intent(in)               :: dt
    ! Déclaration de variable locale / paramètre de calcul
    real(dp), intent(in)               :: b(:)
    ! Déclaration de variable locale / paramètre de calcul
    real(dp), intent(inout)            :: x(:)
    ! Déclaration de variable locale / paramètre de calcul
    integer, intent(out)               :: iters

    select case (cfg%linear_solver)
    case (LINEAR_BICGSTAB2)
      ! Appel de sous-programme modulaire
      call solve_bicgstab2(ns, mesh, cfg, fields_old, dt, b, x, iters)
    case (LINEAR_BICGSTAB)
      ! Appel de sous-programme modulaire
      call solve_bicgstab(ns, mesh, cfg, fields_old, dt, b, x, iters)
    case (LINEAR_CG)
      ! Appel de sous-programme modulaire
      call solve_cg(ns, mesh, cfg, fields_old, dt, b, x, iters)
    case (LINEAR_GMRES)
      ! Appel de sous-programme modulaire
      call solve_gmres(ns, mesh, cfg, fields_old, dt, b, x, iters)
    case default
      ! Appel de sous-programme modulaire
      call solve_bicgstab2(ns, mesh, cfg, fields_old, dt, b, x, iters)
    end select
  end subroutine solve_linear_system

end module module_linear_solvers

\end{lstlisting}

\chapter{Solveur Monolithique de Navier-Stokes par Lagrangien Augmenté en Fortran 2018 (module_navier_stokes.f90)}
\section{Architecture des 1540 Lignes du Module \texttt{module\_navier\_stokes.f90}}
Ce module constitue le noyau physique et mathématique de \texttt{COVOF\_2D-FORTRAN}. Il regroupe :
\begin{enumerate}
  \item \textbf{Mise à jour des propriétés physiques diphasiques (\texttt{update\_fluid\_properties})} : calcule en chaque maille $(i, j)$ la masse volumique $\rho_{i,j} = C_{i,j}\rho_1 + (1-C_{i,j})\rho_2$ et la viscosité dynamique $\mu_{i,j}$ (y compris la loi exponentielle spatiale du Cas 3.9).
  \item \textbf{Application des conditions aux limites hydrodynamiques (\texttt{apply\_velocity\_bc})} : impose les conditions \texttt{BC\_WALL} (adhérence par cellules fantômes antisymétriques), \texttt{BC\_SLIP} (glissement libre par cellules fantômes symétriques), \texttt{BC\_INLET}, \texttt{BC\_OUTLET}, \texttt{BC\_PERIODIC} et \texttt{BC\_JET\_INLET} (injection centrale de buse $V_0$ sur $|x - 0.5| \le D/2$ et évacuation compensatoire $v_{\text{out}} = |V_0| D / (L_x - D)$ garantissant $\oint_{\partial\Omega} \mathbf{u}\cdot\mathbf{n}\,d\Gamma = 0$).
  \item \textbf{Calcul de la force capillaire CSF Balanced-Force (\texttt{compute\_surface\_tension\_force})} : lisse la fonction couleur $C$ sur un voisinage $3\times 3$, calcule les normales unitaires aux sommets, en déduit la courbure centrée $\kappa_{i,j} = -\nabla_h \cdot \mathbf{n}$ et projette la force capillaire $\sigma \kappa_{\text{face}} \nabla_h C$ directement sur les faces MAC $u$ et $v$.
  \item \textbf{Assemblage et résolution du système monolithique du Lagrangien Augmenté (\texttt{solve\_momentum\_augmented\_lagrangian})} : assemble la matrice creuse $\mathbf{A}_{\text{aug}}$ à 9 bandes incluant l'inertie $\rho/\Delta t$, la convection décentrée, le tenseur visqueux complet à moyenne harmonique nodale, la pénalisation de Brinkman $\chi_s/\eta$ et le couplage croisé $-r_{al}\nabla_h(\nabla_h\cdot\mathbf{u})$, puis exécute les sous-itérations d'Uzawa $p^{k+1} = p^k - r_{al}\nabla_h\cdot\mathbf{u}^{k+1}$.
  \item \textbf{Résolution de l'équation de l'énergie thermique (\texttt{solve\_energy\_equation})} : intègre implicitement la température $T_{i,j}^{n+1}$ et applique les conditions aux limites thermiques isothermes et adiabatiques (Cas 6.1 de Vahl Davis).
\end{enumerate}

\section{Répertoire des Types Dérivés, Sous-Routines et Fonctions de \texttt{src/module\_navier\_stokes.f90}}
\begin{center}\small
\begin{tabular}{clll}
\toprule
\textbf{Ligne} & \textbf{Nature} & \textbf{Identifiant Fortran 2018} & \textbf{Interface / Arguments} \\
\midrule
12 & module & \texttt{module\_navier\_stokes} & Unité modulaire Fortran 2018 \\
25 & type dérivé & \texttt{navier\_stokes\_t} & Structure de données Fortran 2018 \\
67 & function & \texttt{is\_u\_face\_in\_obstacle} & Arguments : (mesh, cfg, i, j) \\
94 & function & \texttt{is\_v\_face\_in\_obstacle} & Arguments : (mesh, cfg, i, j) \\
121 & subroutine & \texttt{ns\_init} & Arguments : (this, mesh, cfg) \\
180 & subroutine & \texttt{ns\_pack\_vector} & Arguments : (this, mesh, cfg, fields, vec) \\
249 & subroutine & \texttt{ns\_unpack\_vector} & Arguments : (this, mesh, cfg, vec, fields, homogeneous) \\
329 & subroutine & \texttt{ns\_apply\_system\_matrix} & Arguments : (this, mesh, cfg, fields\_old, dt, x, y) \\
671 & subroutine & \texttt{ns\_assemble\_rhs} & Arguments : (this, mesh, cfg, fields, dt, rhs) \\
943 & subroutine & \texttt{ns\_update\_pressure\_al} & Arguments : (this, mesh, cfg, fields) \\
982 & subroutine & \texttt{ns\_assemble\_preconditioner\_blocks} & Arguments : (this, mesh, cfg, fields, dt) \\
1108 & subroutine & \texttt{ns\_solve\_energy} & Arguments : (this, mesh, cfg, fields, dt) \\
1489 & function & \texttt{ns\_compute\_mean\_nusselt\_hot\_wall} & Arguments : (this, mesh, cfg, fields) \\
\bottomrule
\end{tabular}
\end{center}

\section{Code Source Intégral Annoté : \texttt{src/module\_navier\_stokes.f90}}
Le listing ci-dessous présente l'intégralité des 1551 lignes du fichier \texttt{src/module\_navier\_stokes.f90} de l'implémentation Fortran 2018 :

\begin{lstlisting}[language={[95]Fortran}]
!===============================================================================
! Module: module_navier_stokes
! Description: Solveur Navier-Stokes 2D incompressible diphasique et thermique
!              selon l'architecture Fortran 2018 de COVOF 2D.
!              - Discrétisation spatiale MAC décalée (Harlow & Welch 1965)
!              - Couplage Vitesse-Pression par Lagrangien Augmenté (Fortin & Glowinski)
!              - Algorithme d'Uzawa pour la pression
!              - Viscosités harmonique aux nœuds et arithmétique aux mailles
!              - Équation de l'énergie thermique couplée (Boussinesq)
!              - Pénalisation volumique d'obstacles (Brinkman)
!===============================================================================
module module_navier_stokes
  use module_precision, only: dp, ip
  use module_constants
  use module_config, only: config_t
  use module_mesh
  use module_fields
  use module_mpi_comm
  implicit none
  private

  public :: navier_stokes_t
  public :: is_u_face_in_obstacle, is_v_face_in_obstacle

  type :: navier_stokes_t
    ! Déclaration de variable locale / paramètre de calcul
    integer :: n_owned_u = 0
    ! Déclaration de variable locale / paramètre de calcul
    integer :: n_owned_v = 0
    ! Déclaration de variable locale / paramètre de calcul
    integer :: n_owned_p = 0
    ! Déclaration de variable locale / paramètre de calcul
    integer :: n_owned_total = 0

    ! Diagonales précalculées pour le préconditionneur Jacobi
    real(dp), allocatable :: diag_u(:)
    ! Déclaration de variable locale / paramètre de calcul
    real(dp), allocatable :: diag_v(:)
    ! Déclaration de variable locale / paramètre de calcul
    real(dp), allocatable :: diag_p(:)
  contains
    ! Affectation et calcul numérique
    procedure :: init => ns_init
    ! Affectation et calcul numérique
    procedure :: pack_vector => ns_pack_vector
    ! Affectation et calcul numérique
    procedure :: unpack_vector => ns_unpack_vector
    ! Affectation et calcul numérique
    procedure :: apply_system_matrix => ns_apply_system_matrix
    ! Affectation et calcul numérique
    procedure :: assemble_rhs => ns_assemble_rhs
    ! Évaluation de la norme du résidu ou de l'erreur numérique
    procedure :: update_pressure_al => ns_update_pressure_al
    ! Affectation et calcul numérique
    procedure :: assemble_preconditioner_blocks => ns_assemble_preconditioner_blocks
    ! Affectation et calcul numérique
    procedure :: solve_energy => ns_solve_energy
    ! Affectation et calcul numérique
    procedure :: compute_mean_nusselt_hot_wall => ns_compute_mean_nusselt_hot_wall
  end type navier_stokes_t

contains

  !-----------------------------------------------------------------------------
  ! Détection si une face de vitesse normale u(i, j) est dans un obstacle solide
  !-----------------------------------------------------------------------------
  pure function is_u_face_in_obstacle(mesh, cfg, i, j) result(in_obs)
    ! Déclaration de variable locale / paramètre de calcul
    type(mesh_2d_t), intent(in)   :: mesh
    ! Déclaration de variable locale / paramètre de calcul
    type(config_t), intent(in)    :: cfg
    ! Déclaration de variable locale / paramètre de calcul
    integer, intent(in)           :: i, j
    ! Déclaration de variable locale / paramètre de calcul
    logical                       :: in_obs
    ! Déclaration de variable locale / paramètre de calcul
    real(dp) :: x_f, y_c

    ! Affectation et calcul numérique
    in_obs = .false.
    if (.not. cfg%obs_enabled) return
    ! Affectation et calcul numérique
    x_f = x_face(mesh, i)
    ! Affectation et calcul numérique
    y_c = y_center(mesh, j)
    ! Affectation et calcul numérique
    in_obs = is_in_obstacle(x_f, y_c, cfg%obs_type, cfg%obs_x0, cfg%obs_y0, &
                            cfg%obs_x1, cfg%obs_y1, cfg%obs_xc, cfg%obs_yc, cfg%obs_radius)
  end function is_u_face_in_obstacle

  !-----------------------------------------------------------------------------
  ! Détection si une face de vitesse normale v(i, j) est dans un obstacle solide
  !-----------------------------------------------------------------------------
  pure function is_v_face_in_obstacle(mesh, cfg, i, j) result(in_obs)
    ! Déclaration de variable locale / paramètre de calcul
    type(mesh_2d_t), intent(in)   :: mesh
    ! Déclaration de variable locale / paramètre de calcul
    type(config_t), intent(in)    :: cfg
    ! Déclaration de variable locale / paramètre de calcul
    integer, intent(in)           :: i, j
    ! Déclaration de variable locale / paramètre de calcul
    logical                       :: in_obs
    ! Déclaration de variable locale / paramètre de calcul
    real(dp) :: x_c, y_f

    ! Affectation et calcul numérique
    in_obs = .false.
    if (.not. cfg%obs_enabled) return
    ! Affectation et calcul numérique
    x_c = x_center(mesh, i)
    ! Affectation et calcul numérique
    y_f = y_face(mesh, j)
    ! Affectation et calcul numérique
    in_obs = is_in_obstacle(x_c, y_f, cfg%obs_type, cfg%obs_x0, cfg%obs_y0, &
                            cfg%obs_x1, cfg%obs_y1, cfg%obs_xc, cfg%obs_yc, cfg%obs_radius)
  end function is_v_face_in_obstacle

  !-----------------------------------------------------------------------------
  ! Initialisation du solveur Navier-Stokes et comptage des DDLs possédés
  !-----------------------------------------------------------------------------
  subroutine ns_init(this, mesh, cfg)
    ! Déclaration de variable locale / paramètre de calcul
    class(navier_stokes_t), intent(inout) :: this
    ! Déclaration de variable locale / paramètre de calcul
    type(mesh_2d_t), intent(in)           :: mesh
    ! Déclaration de variable locale / paramètre de calcul
    type(config_t), intent(in)            :: cfg
    ! Déclaration de variable locale / paramètre de calcul
    integer :: u_nx, v_ny

    ! Affectation et calcul numérique
    u_nx = mesh%nx_local
    ! Affectation et calcul numérique
    if (mesh%is_boundary_right) u_nx = u_nx + 1
    ! Affectation et calcul numérique
    this%n_owned_u = u_nx * mesh%ny_local

    ! Affectation et calcul numérique
    v_ny = mesh%ny_local
    ! Affectation et calcul numérique
    if (mesh%is_boundary_top) v_ny = v_ny + 1
    ! Affectation et calcul numérique
    this%n_owned_v = mesh%nx_local * v_ny

    ! Affectation et calcul numérique
    this%n_owned_p = mesh%nx_local * mesh%ny_local

    ! Évaluation de la condition logique
    if (cfg%solver_type == SOLVER_AUGMENTED_LAGRANGIAN) then
      ! Affectation et calcul numérique
      this%n_owned_total = this%n_owned_u + this%n_owned_v
    ! Branche par défaut si condition non satisfaite
    else
      ! Affectation et calcul numérique
      this%n_owned_total = this%n_owned_u + this%n_owned_v + this%n_owned_p
    end if

    if (allocated(this%diag_u)) deallocate(this%diag_u)
    if (allocated(this%diag_v)) deallocate(this%diag_v)
    if (allocated(this%diag_p)) deallocate(this%diag_p)

    ! Allocation dynamique de tableau contigu en mémoire
    allocate(this%diag_u(this%n_owned_u))
    ! Allocation dynamique de tableau contigu en mémoire
    allocate(this%diag_v(this%n_owned_v))
    ! Allocation dynamique de tableau contigu en mémoire
    allocate(this%diag_p(this%n_owned_p))

    ! Affectation et calcul numérique
    this%diag_u = 1.0_dp
    ! Affectation et calcul numérique
    this%diag_v = 1.0_dp
    ! Affectation et calcul numérique
    this%diag_p = 1.0_dp
  end subroutine ns_init

  !-----------------------------------------------------------------------------
  ! Empaquette les champs physiques 2D dans un vecteur 1D contigu
  !-----------------------------------------------------------------------------
  subroutine ns_pack_vector(this, mesh, cfg, fields, vec)
    ! Déclaration de variable locale / paramètre de calcul
    class(navier_stokes_t), intent(in) :: this
    ! Déclaration de variable locale / paramètre de calcul
    type(mesh_2d_t), intent(in)        :: mesh
    ! Déclaration de variable locale / paramètre de calcul
    type(config_t), intent(in)         :: cfg
    ! Déclaration de variable locale / paramètre de calcul
    type(flow_fields_t), intent(in)    :: fields
    ! Déclaration de variable locale / paramètre de calcul
    real(dp), intent(out)              :: vec(:)

    ! Déclaration de variable locale / paramètre de calcul
    integer :: i, j, dof, u_nx, v_ny, v_offset, p_offset

    ! Affectation et calcul numérique
    u_nx = mesh%nx_local
    ! Affectation et calcul numérique
    if (mesh%is_boundary_right) u_nx = u_nx + 1

    ! 1. Pack u
    do j = 1, mesh%ny_local
      ! Boucle spatiale selon l'axe horizontal X (accès contigu colonne-major)
      do i = 1, u_nx
        ! Affectation et calcul numérique
        dof = (j - 1) * u_nx + i
        ! Affectation et calcul numérique
        vec(dof) = fields%u(i, j)
      end do
    end do

    ! 2. Pack v
    v_offset = this%n_owned_u
    ! Affectation et calcul numérique
    v_ny = mesh%ny_local
    ! Affectation et calcul numérique
    if (mesh%is_boundary_top) v_ny = v_ny + 1

    ! Boucle spatiale sur les cellules locales selon l'axe vertical Y (ordonnées)
    do j = 1, v_ny
      ! Boucle spatiale selon l'axe horizontal X (accès contigu colonne-major)
      do i = 1, mesh%nx_local
        ! Affectation et calcul numérique
        dof = v_offset + (j - 1) * mesh%nx_local + i
        ! Affectation et calcul numérique
        vec(dof) = fields%v(i, j)
      end do
    end do

    ! 3. Pack p (si monolithique)
    if (cfg%solver_type == SOLVER_MONOLITHIC) then
      ! Affectation et calcul numérique
      p_offset = this%n_owned_u + this%n_owned_v
      ! Boucle spatiale sur les cellules locales selon l'axe vertical Y (ordonnées)
      do j = 1, mesh%ny_local
        ! Boucle spatiale selon l'axe horizontal X (accès contigu colonne-major)
        do i = 1, mesh%nx_local
          ! Affectation et calcul numérique
          dof = p_offset + (j - 1) * mesh%nx_local + i
          ! Affectation et calcul numérique
          vec(dof) = fields%p(i, j)
        end do
      end do
    end if
  end subroutine ns_pack_vector

  !-----------------------------------------------------------------------------
  ! Dépaquette le vecteur 1D dans les champs 2D et synchronise les halos MPI
  !-----------------------------------------------------------------------------
  subroutine ns_unpack_vector(this, mesh, cfg, vec, fields, homogeneous)
    ! Déclaration de variable locale / paramètre de calcul
    class(navier_stokes_t), intent(in) :: this
    ! Déclaration de variable locale / paramètre de calcul
    type(mesh_2d_t), intent(in)        :: mesh
    ! Déclaration de variable locale / paramètre de calcul
    type(config_t), intent(in)         :: cfg
    ! Déclaration de variable locale / paramètre de calcul
    real(dp), intent(in)               :: vec(:)
    ! Déclaration de variable locale / paramètre de calcul
    type(flow_fields_t), intent(inout) :: fields
    ! Déclaration de variable locale / paramètre de calcul
    logical, intent(in)                :: homogeneous

    ! Déclaration de variable locale / paramètre de calcul
    integer :: i, j, dof, u_nx, v_ny, v_offset, p_offset

    ! Affectation et calcul numérique
    u_nx = mesh%nx_local
    ! Affectation et calcul numérique
    if (mesh%is_boundary_right) u_nx = u_nx + 1

    ! 1. Unpack u
    do j = 1, mesh%ny_local
      ! Boucle spatiale selon l'axe horizontal X (accès contigu colonne-major)
      do i = 1, u_nx
        ! Affectation et calcul numérique
        dof = (j - 1) * u_nx + i
        ! Mise à jour de la composante de vitesse discrète sur face décalée MAC
        fields%u(i, j) = vec(dof)
      end do
    end do

    ! 2. Unpack v
    v_offset = this%n_owned_u
    ! Affectation et calcul numérique
    v_ny = mesh%ny_local
    ! Affectation et calcul numérique
    if (mesh%is_boundary_top) v_ny = v_ny + 1

    ! Boucle spatiale sur les cellules locales selon l'axe vertical Y (ordonnées)
    do j = 1, v_ny
      ! Boucle spatiale selon l'axe horizontal X (accès contigu colonne-major)
      do i = 1, mesh%nx_local
        ! Affectation et calcul numérique
        dof = v_offset + (j - 1) * mesh%nx_local + i
        ! Mise à jour de la composante de vitesse discrète sur face décalée MAC
        fields%v(i, j) = vec(dof)
      end do
    end do

    ! 3. Unpack p (si monolithique)
    if (cfg%solver_type == SOLVER_MONOLITHIC) then
      ! Affectation et calcul numérique
      p_offset = this%n_owned_u + this%n_owned_v
      ! Boucle spatiale sur les cellules locales selon l'axe vertical Y (ordonnées)
      do j = 1, mesh%ny_local
        ! Boucle spatiale selon l'axe horizontal X (accès contigu colonne-major)
        do i = 1, mesh%nx_local
          ! Affectation et calcul numérique
          dof = p_offset + (j - 1) * mesh%nx_local + i
          ! Actualisation de la pression statique au centre de la maille
          fields%p(i, j) = vec(dof)
        end do
      end do
      ! Appel de sous-programme modulaire
      call sync_cell_scalar(mesh, fields%p)
    end if

    ! Synchronisation des mailles fantômes et conditions aux limites
    call sync_u(mesh, fields%u)
    ! Appel de sous-programme modulaire
    call sync_v(mesh, fields%v)
    ! Appel de sous-programme modulaire
    call apply_boundary_conditions(mesh, cfg, fields, homogeneous)
  end subroutine ns_unpack_vector

  !-----------------------------------------------------------------------------
  ! Application de l'opérateur matriciel complet y = A * x
  !-----------------------------------------------------------------------------
  subroutine ns_apply_system_matrix(this, mesh, cfg, fields_old, dt, x, y)
    ! Déclaration de variable locale / paramètre de calcul
    class(navier_stokes_t), intent(in) :: this
    ! Déclaration de variable locale / paramètre de calcul
    type(mesh_2d_t), intent(in)        :: mesh
    ! Déclaration de variable locale / paramètre de calcul
    type(config_t), intent(in)         :: cfg
    ! Déclaration de variable locale / paramètre de calcul
    type(flow_fields_t), intent(in)    :: fields_old
    ! Déclaration de variable locale / paramètre de calcul
    real(dp), intent(in)               :: dt
    ! Déclaration de variable locale / paramètre de calcul
    real(dp), intent(in)               :: x(:)
    ! Déclaration de variable locale / paramètre de calcul
    real(dp), intent(out)              :: y(:)

    ! Déclaration de variable locale / paramètre de calcul
    type(flow_fields_t) :: tf
    ! Déclaration de variable locale / paramètre de calcul
    real(dp) :: dx, dy, dvol
    ! Déclaration de variable locale / paramètre de calcul
    integer :: i, j, dof, u_nx, v_ny, v_offset, p_offset
    ! Déclaration de variable locale / paramètre de calcul
    real(dp) :: div_cell(0:mesh%nx_local+1, 0:mesh%ny_local+1)
    ! Déclaration de variable locale / paramètre de calcul
    real(dp) :: rho_u, rho_v, diag_bc, val
    ! Déclaration de variable locale / paramètre de calcul
    real(dp) :: u_n, u_n_E, u_n_W, v_n_N, v_n_S
    ! Déclaration de variable locale / paramètre de calcul
    real(dp) :: conv_E, conv_W, conv_N, conv_S
    ! Déclaration de variable locale / paramètre de calcul
    real(dp) :: mu_cell_E, mu_cell_W, diff_elong_E, diff_elong_W
    ! Déclaration de variable locale / paramètre de calcul
    real(dp) :: mu_node_N, mu_node_S, shear_N, shear_S
    ! Déclaration de variable locale / paramètre de calcul
    real(dp) :: grad_div_x, grad_div_y, grad_p
    ! Déclaration de variable locale / paramètre de calcul
    real(dp) :: v_n, v_n_S_val
    ! Déclaration de variable locale / paramètre de calcul
    real(dp) :: mu_cell_N, mu_cell_S, diff_elong_N, diff_elong_S
    ! Déclaration de variable locale / paramètre de calcul
    real(dp) :: mu_node_E, mu_node_W, shear_E, shear_W
    ! Déclaration de variable locale / paramètre de calcul
    real(dp) :: div_val, diag_p_term
    ! Déclaration de variable locale / paramètre de calcul
    logical  :: is_wall_bot, is_wall_top, in_obs

    ! Calcul de la métrique spatiale élémentaire
    dx = mesh%dx
    ! Calcul de la métrique spatiale élémentaire
    dy = mesh%dy
    ! Affectation et calcul numérique
    dvol = dx * dy

    ! Affectation et calcul numérique
    u_nx = mesh%nx_local
    ! Affectation et calcul numérique
    if (mesh%is_boundary_right) u_nx = u_nx + 1

    ! Affectation et calcul numérique
    v_ny = mesh%ny_local
    ! Affectation et calcul numérique
    if (mesh%is_boundary_top) v_ny = v_ny + 1

    ! Appel de sous-programme modulaire
    call tf%allocate(mesh)
    ! Appel de sous-programme modulaire
    call this%unpack_vector(mesh, cfg, x, tf, .true.)

    ! Précalcul de la divergence sur chaque cellule locale + halo (dans les bornes allouées 0:nx+1, 0:ny+1)
    div_cell = 0.0_dp
    do j = 0, mesh%ny_local
      ! Boucle spatiale selon l'axe horizontal X (accès contigu colonne-major)
      do i = 0, mesh%nx_local
        ! Affectation et calcul numérique
        div_cell(i, j) = (tf%u(i + 1, j) - tf%u(i, j)) / dx + &
                         (tf%v(i, j + 1) - tf%v(i, j)) / dy
      end do
    end do

    ! =========================================================================
    ! 1. Équation de Quantité de Mouvement en u (Faces verticales)
    ! =========================================================================
    do j = 1, mesh%ny_local
      ! Boucle spatiale selon l'axe horizontal X (accès contigu colonne-major)
      do i = 1, u_nx
        ! Affectation et calcul numérique
        dof = (j - 1) * u_nx + i
        ! Calcul de la masse volumique locale ou interpolée aux faces
        rho_u = fields_old%rho_u(i, j)
        ! Affectation et calcul numérique
        diag_bc = (rho_u / dt) * dvol

        ! Condition aux limites physique sur les parois gauche / droite
        if (i == 1 .and. mesh%is_boundary_left) then
          ! Évaluation de la condition logique
          if (cfg%bc_left%btype == BC_NEUMANN) then
            ! Affectation et calcul numérique
            y(dof) = diag_bc * (tf%u(1, j) - tf%u(2, j))
          ! Branche par défaut si condition non satisfaite
          else
            ! Affectation et calcul numérique
            y(dof) = diag_bc * tf%u(1, j)
          end if
          cycle
        end if

        ! Évaluation de la condition logique
        if (i == u_nx .and. mesh%is_boundary_right) then
          ! Évaluation de la condition logique
          if (cfg%bc_right%btype == BC_NEUMANN) then
            ! Affectation et calcul numérique
            y(dof) = diag_bc * (tf%u(u_nx, j) - tf%u(u_nx - 1, j))
          ! Branche par défaut si condition non satisfaite
          else
            ! Affectation et calcul numérique
            y(dof) = diag_bc * tf%u(u_nx, j)
          end if
          cycle
        end if

        ! Face dans un obstacle solide (adhérence nulle)
        if (is_u_face_in_obstacle(mesh, cfg, i, j)) then
          ! Affectation et calcul numérique
          y(dof) = diag_bc * tf%u(i, j)
          cycle
        end if

        ! Terme instationnaire
        val = (rho_u / dt) * tf%u(i, j) * dvol

        ! Terme convectif semi-implicite F_uu (forme advective non-conservative a densite harmonique)
        if (cfg%enable_convection) then
          u_n = fields_old%u(i, j)
          u_n_E = fields_old%u(i + 1, j)
          u_n_W = fields_old%u(i - 1, j)

          v_n_N = 0.5_dp * (fields_old%v(i - 1, j + 1) + fields_old%v(i, j + 1))
          v_n_S = 0.5_dp * (fields_old%v(i - 1, j) + fields_old%v(i, j))

          mu_cell_E = (2.0_dp * rho_u * fields_old%rho_u(i + 1, j)) / max(1.0e-12_dp, rho_u + fields_old%rho_u(i + 1, j))
          mu_cell_W = (2.0_dp * rho_u * fields_old%rho_u(i - 1, j)) / max(1.0e-12_dp, rho_u + fields_old%rho_u(i - 1, j))
          if (j < mesh%ny_local) then
            mu_node_N = (2.0_dp * rho_u * fields_old%rho_u(i, j + 1)) / max(1.0e-12_dp, rho_u + fields_old%rho_u(i, j + 1))
          else
            mu_node_N = rho_u
          end if
          if (j > 1) then
            mu_node_S = (2.0_dp * rho_u * fields_old%rho_u(i, j - 1)) / max(1.0e-12_dp, rho_u + fields_old%rho_u(i, j - 1))
          else
            mu_node_S = rho_u
          end if

          conv_E = 0.5_dp * (u_n + u_n_E) * 0.5_dp * (tf%u(i + 1, j) - tf%u(i, j)) * dy
          conv_W = 0.5_dp * (u_n_W + u_n) * 0.5_dp * (tf%u(i, j) - tf%u(i - 1, j)) * dy
          conv_N = v_n_N * 0.5_dp * (tf%u(i, j + 1) - tf%u(i, j)) * dx
          conv_S = v_n_S * 0.5_dp * (tf%u(i, j) - tf%u(i, j - 1)) * dx

          val = val + (mu_cell_E * conv_E + mu_cell_W * conv_W + mu_node_N * conv_N + mu_node_S * conv_S)
        end if

        ! Termes visqueux implicites d'ordre 2
        mu_cell_E = fields_old%mu(i, j)
        ! Calcul de la viscosité dynamique locale (harmonique aux nœuds)
        mu_cell_W = fields_old%mu(i - 1, j)
        ! Affectation et calcul numérique
        diff_elong_E = 2.0_dp * mu_cell_E * (tf%u(i + 1, j) - tf%u(i, j)) / dx * dy
        ! Affectation et calcul numérique
        diff_elong_W = 2.0_dp * mu_cell_W * (tf%u(i, j) - tf%u(i - 1, j)) / dx * dy
        ! Affectation et calcul numérique
        val = val - (diff_elong_E - diff_elong_W)

        ! Calcul de la viscosité dynamique locale (harmonique aux nœuds)
        mu_node_N = fields_old%mu_node(i, j + 1)
        ! Calcul de la viscosité dynamique locale (harmonique aux nœuds)
        mu_node_S = fields_old%mu_node(i, j)

        ! Affectation et calcul numérique
        shear_N = mu_node_N * ((tf%u(i, j + 1) - tf%u(i, j)) / dy + &
                               (tf%v(i, j + 1) - tf%v(i - 1, j + 1)) / dx) * dx
        ! Affectation et calcul numérique
        shear_S = mu_node_S * ((tf%u(i, j) - tf%u(i, j - 1)) / dy + &
                               (tf%v(i, j) - tf%v(i - 1, j)) / dx) * dx
        ! Affectation et calcul numérique
        val = val - (shear_N - shear_S)

        ! Lagrangien augmenté : - r * d(div)/dx * dvol
        grad_div_x = (div_cell(i, j) - div_cell(i - 1, j)) * dy
        ! Affectation et calcul numérique
        val = val - cfg%r_al * grad_div_x

        ! Gradient de pression discret B_u^T (si monolithique)
        if (cfg%solver_type == SOLVER_MONOLITHIC) then
          ! Affectation et calcul numérique
          grad_p = (tf%p(i, j) - tf%p(i - 1, j)) * dy
          ! Affectation et calcul numérique
          val = val + grad_p
        end if

        ! Affectation et calcul numérique
        y(dof) = val
      end do
    end do

    ! =========================================================================
    ! 2. Équation de Quantité de Mouvement en v (Faces horizontales)
    ! =========================================================================
    v_offset = this%n_owned_u
    ! Boucle spatiale sur les cellules locales selon l'axe vertical Y (ordonnées)
    do j = 1, v_ny
      ! Boucle spatiale selon l'axe horizontal X (accès contigu colonne-major)
      do i = 1, mesh%nx_local
        ! Affectation et calcul numérique
        dof = v_offset + (j - 1) * mesh%nx_local + i
        ! Calcul de la masse volumique locale ou interpolée aux faces
        rho_v = fields_old%rho_v(i, j)
        ! Affectation et calcul numérique
        diag_bc = (rho_v / dt) * dvol

        ! Affectation et calcul numérique
        is_wall_bot = (j == 1 .and. mesh%is_boundary_bottom .and. &
                       ! Affectation et calcul numérique
                       (cfg%bc_bottom%btype == BC_DIRICHLET .or. cfg%bc_bottom%btype == BC_FREESLIP))
        ! Affectation et calcul numérique
        is_wall_top = (j == v_ny .and. mesh%is_boundary_top .and. &
                       ! Affectation et calcul numérique
                       (cfg%bc_top%btype == BC_DIRICHLET .or. cfg%bc_top%btype == BC_FREESLIP))
        ! Affectation et calcul numérique
        in_obs = is_v_face_in_obstacle(mesh, cfg, i, j)

        ! Évaluation de la condition logique
        if (is_wall_bot .or. is_wall_top .or. in_obs) then
          ! Affectation et calcul numérique
          y(dof) = diag_bc * tf%v(i, j)
          cycle
        end if

        ! Affectation et calcul numérique
        val = (rho_v / dt) * tf%v(i, j) * dvol

        ! Évaluation de la condition logique
        if (cfg%enable_convection) then
          v_n = fields_old%v(i, j)
          v_n_N = fields_old%v(i, j + 1)
          v_n_S_val = fields_old%v(i, j - 1)

          u_n_E = 0.5_dp * (fields_old%u(i + 1, j - 1) + fields_old%u(i + 1, j))
          u_n_W = 0.5_dp * (fields_old%u(i, j - 1) + fields_old%u(i, j))

          if (i < mesh%nx_local) then
            mu_node_E = (2.0_dp * rho_v * fields_old%rho_v(i + 1, j)) / max(1.0e-12_dp, rho_v + fields_old%rho_v(i + 1, j))
          else
            mu_node_E = rho_v
          end if
          if (i > 1) then
            mu_node_W = (2.0_dp * rho_v * fields_old%rho_v(i - 1, j)) / max(1.0e-12_dp, rho_v + fields_old%rho_v(i - 1, j))
          else
            mu_node_W = rho_v
          end if
          mu_cell_N = (2.0_dp * rho_v * fields_old%rho_v(i, j + 1)) / max(1.0e-12_dp, rho_v + fields_old%rho_v(i, j + 1))
          mu_cell_S = (2.0_dp * rho_v * fields_old%rho_v(i, j - 1)) / max(1.0e-12_dp, rho_v + fields_old%rho_v(i, j - 1))

          conv_E = u_n_E * 0.5_dp * (tf%v(i + 1, j) - tf%v(i, j)) * dy
          conv_W = u_n_W * 0.5_dp * (tf%v(i, j) - tf%v(i - 1, j)) * dy
          conv_N = 0.5_dp * (v_n + v_n_N) * 0.5_dp * (tf%v(i, j + 1) - tf%v(i, j)) * dx
          conv_S = 0.5_dp * (v_n_S_val + v_n) * 0.5_dp * (tf%v(i, j) - tf%v(i, j - 1)) * dx

          val = val + (mu_node_E * conv_E + mu_node_W * conv_W + mu_cell_N * conv_N + mu_cell_S * conv_S)
        end if

        ! Calcul de la viscosité dynamique locale (harmonique aux nœuds)
        mu_cell_N = fields_old%mu(i, j)
        ! Calcul de la viscosité dynamique locale (harmonique aux nœuds)
        mu_cell_S = fields_old%mu(i, j - 1)
        ! Affectation et calcul numérique
        diff_elong_N = 2.0_dp * mu_cell_N * (tf%v(i, j + 1) - tf%v(i, j)) / dy * dx
        ! Affectation et calcul numérique
        diff_elong_S = 2.0_dp * mu_cell_S * (tf%v(i, j) - tf%v(i, j - 1)) / dy * dx
        ! Affectation et calcul numérique
        val = val - (diff_elong_N - diff_elong_S)

        ! Calcul de la viscosité dynamique locale (harmonique aux nœuds)
        mu_node_E = fields_old%mu_node(i + 1, j)
        ! Calcul de la viscosité dynamique locale (harmonique aux nœuds)
        mu_node_W = fields_old%mu_node(i, j)

        ! Affectation et calcul numérique
        shear_E = mu_node_E * ((tf%v(i + 1, j) - tf%v(i, j)) / dx + &
                               (tf%u(i + 1, j) - tf%u(i + 1, j - 1)) / dy) * dy
        ! Affectation et calcul numérique
        shear_W = mu_node_W * ((tf%v(i, j) - tf%v(i - 1, j)) / dx + &
                               (tf%u(i, j) - tf%u(i, j - 1)) / dy) * dy
        ! Affectation et calcul numérique
        val = val - (shear_E - shear_W)

        ! Affectation et calcul numérique
        grad_div_y = (div_cell(i, j) - div_cell(i, j - 1)) * dx
        ! Affectation et calcul numérique
        val = val - cfg%r_al * grad_div_y

        ! Évaluation de la condition logique
        if (cfg%solver_type == SOLVER_MONOLITHIC) then
          ! Affectation et calcul numérique
          grad_p = (tf%p(i, j) - tf%p(i, j - 1)) * dx
          ! Affectation et calcul numérique
          val = val + grad_p
        end if

        ! Affectation et calcul numérique
        y(dof) = val
      end do
    end do

    ! =========================================================================
    ! 3. Équation de Continuité (si monolithique régularisé)
    ! =========================================================================
    if (cfg%solver_type == SOLVER_MONOLITHIC) then
      ! Affectation et calcul numérique
      p_offset = this%n_owned_u + this%n_owned_v
      ! Boucle spatiale sur les cellules locales selon l'axe vertical Y (ordonnées)
      do j = 1, mesh%ny_local
        ! Boucle spatiale selon l'axe horizontal X (accès contigu colonne-major)
        do i = 1, mesh%nx_local
          ! Affectation et calcul numérique
          dof = p_offset + (j - 1) * mesh%nx_local + i
          ! Affectation et calcul numérique
          div_val = div_cell(i, j) * dvol
          ! Affectation et calcul numérique
          diag_p_term = (dvol / cfg%r_al) * tf%p(i, j)
          ! Affectation et calcul numérique
          y(dof) = div_val - diag_p_term
        end do
      end do
    end if

    ! Appel de sous-programme modulaire
    call tf%free()
  end subroutine ns_apply_system_matrix

  !-----------------------------------------------------------------------------
  ! Assemblage du second membre rhs pour le système linéaire
  !-----------------------------------------------------------------------------
  subroutine ns_assemble_rhs(this, mesh, cfg, fields, dt, rhs)
    ! Déclaration de variable locale / paramètre de calcul
    class(navier_stokes_t), intent(in) :: this
    ! Déclaration de variable locale / paramètre de calcul
    type(mesh_2d_t), intent(in)        :: mesh
    ! Déclaration de variable locale / paramètre de calcul
    type(config_t), intent(in)         :: cfg
    ! Déclaration de variable locale / paramètre de calcul
    type(flow_fields_t), intent(in)    :: fields
    ! Déclaration de variable locale / paramètre de calcul
    real(dp), intent(in)               :: dt
    ! Déclaration de variable locale / paramètre de calcul
    real(dp), intent(out)              :: rhs(:)

    ! Déclaration de variable locale / paramètre de calcul
    real(dp) :: dx, dy, dvol
    ! Déclaration de variable locale / paramètre de calcul
    integer :: i, j, dof, u_nx, v_ny, v_offset, p_offset
    ! Déclaration de variable locale / paramètre de calcul
    real(dp) :: rho_u, rho_v, diag_bc, val, u_bc, v_bc, y_pos, h_step, grad_p
    ! Déclaration de variable locale / paramètre de calcul
    real(dp) :: mu_N, mu_S, mu_W, mu_E, T_v
    ! Déclaration des variables pour l'injecteur JET_INLET (Cas 3.3 Fluids 2021)
    real(dp) :: xc_jet, half_w_jet

    ! Calcul de la métrique spatiale élémentaire
    dx = mesh%dx
    ! Calcul de la métrique spatiale élémentaire
    dy = mesh%dy
    ! Affectation et calcul numérique
    dvol = dx * dy

    ! Affectation et calcul numérique
    u_nx = mesh%nx_local
    ! Affectation et calcul numérique
    if (mesh%is_boundary_right) u_nx = u_nx + 1

    ! Affectation et calcul numérique
    v_ny = mesh%ny_local
    ! Affectation et calcul numérique
    if (mesh%is_boundary_top) v_ny = v_ny + 1

    ! Affectation et calcul numérique
    rhs = 0.0_dp

    ! =========================================================================
    ! 1. Second membre pour u
    ! =========================================================================
    do j = 1, mesh%ny_local
      ! Boucle spatiale selon l'axe horizontal X (accès contigu colonne-major)
      do i = 1, u_nx
        ! Affectation et calcul numérique
        dof = (j - 1) * u_nx + i
        ! Calcul de la masse volumique locale ou interpolée aux faces
        rho_u = fields%rho_u(i, j)
        ! Affectation et calcul numérique
        diag_bc = (rho_u / dt) * dvol

        ! Évaluation de la condition logique
        if (is_u_face_in_obstacle(mesh, cfg, i, j)) then
          ! Affectation et calcul numérique
          rhs(dof) = 0.0_dp
          cycle
        end if

        ! Évaluation de la condition logique
        if (i == 1 .and. mesh%is_boundary_left) then
          ! Affectation et calcul numérique
          u_bc = 0.0_dp
          ! Évaluation de la condition logique
          if (cfg%bc_left%btype == BC_DIRICHLET) then
            ! Évaluation de la condition logique
            if (cfg%obs_enabled .and. cfg%obs_type == OBS_BOX) then
              ! Affectation et calcul numérique
              y_pos = y_center(mesh, j)
              ! Évaluation de la condition logique
              if (y_pos >= cfg%obs_y1) then
                ! Affectation et calcul numérique
                h_step = mesh%ly - cfg%obs_y1
                ! Affectation et calcul numérique
                u_bc = 6.0_dp * cfg%bc_left%u * (y_pos - cfg%obs_y1) * (mesh%ly - y_pos) / (h_step * h_step)
              end if
            ! Branche par défaut si condition non satisfaite
            else
              ! Affectation et calcul numérique
              u_bc = cfg%bc_left%u
            end if
          end if
          ! Affectation et calcul numérique
          rhs(dof) = diag_bc * u_bc
          cycle
        end if

        ! Évaluation de la condition logique
        if (i == u_nx .and. mesh%is_boundary_right) then
          ! Affectation et calcul numérique
          u_bc = 0.0_dp
          ! Évaluation de la condition logique
          if (cfg%bc_right%btype == BC_DIRICHLET) then
            ! Affectation et calcul numérique
            u_bc = cfg%bc_right%u
          end if
          ! Affectation et calcul numérique
          rhs(dof) = diag_bc * u_bc
          cycle
        end if

        ! Affectation et calcul numérique
        val = (rho_u / dt) * fields%u(i, j) * dvol
        ! Affectation et calcul numérique
        val = val + rho_u * cfg%gx * dvol
        ! Affectation et calcul numérique
        val = val + fields%fx_sigma(i, j) * dvol

        ! Évaluation de la condition logique
        if (cfg%solver_type == SOLVER_AUGMENTED_LAGRANGIAN) then
          ! Affectation et calcul numérique
          grad_p = (fields%p(i, j) - fields%p(i - 1, j)) * dy
          ! Affectation et calcul numérique
          val = val - grad_p
        end if

        ! Évaluation de la condition logique
        if (j == mesh%ny_local .and. mesh%is_boundary_top .and. cfg%bc_top%btype == BC_DIRICHLET) then
          ! Calcul de la viscosité dynamique locale (harmonique aux nœuds)
          mu_N = fields%mu_node(i, j + 1)
          ! Affectation et calcul numérique
          val = val + mu_N * (2.0_dp * cfg%bc_top%u / dy) * dx
        end if
        ! Évaluation de la condition logique
        if (j == 1 .and. mesh%is_boundary_bottom .and. cfg%bc_bottom%btype == BC_DIRICHLET) then
          ! Calcul de la viscosité dynamique locale (harmonique aux nœuds)
          mu_S = fields%mu_node(i, j)
          ! Affectation et calcul numérique
          val = val + mu_S * (2.0_dp * cfg%bc_bottom%u / dy) * dx
        end if

        ! Affectation et calcul numérique
        rhs(dof) = val
      end do
    end do

    ! =========================================================================
    ! 2. Second membre pour v
    ! =========================================================================
    v_offset = this%n_owned_u
    ! Boucle spatiale sur les cellules locales selon l'axe vertical Y (ordonnées)
    do j = 1, v_ny
      ! Boucle spatiale selon l'axe horizontal X (accès contigu colonne-major)
      do i = 1, mesh%nx_local
        ! Affectation et calcul numérique
        dof = v_offset + (j - 1) * mesh%nx_local + i
        ! Calcul de la masse volumique locale ou interpolée aux faces
        rho_v = fields%rho_v(i, j)
        ! Affectation et calcul numérique
        diag_bc = (rho_v / dt) * dvol

        ! Évaluation de la condition logique
        if (is_v_face_in_obstacle(mesh, cfg, i, j)) then
          ! Affectation et calcul numérique
          rhs(dof) = 0.0_dp
          cycle
        end if

        ! Affectation et calcul numérique
        if (j == 1 .and. mesh%is_boundary_bottom .and. &
            ! Affectation et calcul numérique
            (cfg%bc_bottom%btype == BC_DIRICHLET .or. cfg%bc_bottom%btype == BC_FREESLIP)) then
          ! Affectation et calcul numérique
          v_bc = 0.0_dp
          ! Affectation et calcul numérique
          if (cfg%bc_bottom%btype == BC_DIRICHLET) v_bc = cfg%bc_bottom%v
          ! Affectation et calcul numérique
          rhs(dof) = diag_bc * v_bc
          cycle
        end if

        ! Affectation et calcul numérique
        if (j == v_ny .and. mesh%is_boundary_top .and. &
            ! Affectation et calcul numérique
            (cfg%bc_top%btype == BC_DIRICHLET .or. cfg%bc_top%btype == BC_FREESLIP)) then
          ! Affectation et calcul numérique
          v_bc = 0.0_dp
          ! Évaluation du profil d'injecteur JET_INLET (Cas 3.3 Fluids 2021)
          if (trim(cfg%bc_top%profile) == "JET_INLET") then
             xc_jet = merge(cfg%bc_top%jet_center, 0.5_dp * mesh%lx, cfg%bc_top%jet_center > 0.0_dp)
             half_w_jet = merge(0.5_dp * cfg%bc_top%jet_width, 0.04_dp, cfg%bc_top%jet_width > 0.0_dp)
             T_v = 0.0_dp
             do p_offset = 1, mesh%nx_global
                if (abs((real(p_offset, dp) - 0.5_dp) * mesh%dx - xc_jet) <= half_w_jet) T_v = T_v + mesh%dx
             end do
             if (T_v < 1.0e-12_dp) T_v = 2.0_dp * half_w_jet
             if (abs(x_center(mesh, i) - xc_jet) <= half_w_jet) then
                v_bc = cfg%bc_top%v * (2.0_dp * half_w_jet) / T_v
             else
                v_bc = -cfg%bc_top%v * (2.0_dp * half_w_jet) / (mesh%lx - T_v)
             end if
          else if (cfg%bc_top%btype == BC_DIRICHLET) then
             v_bc = cfg%bc_top%v
          end if
          ! Affectation et calcul numérique
          rhs(dof) = diag_bc * v_bc
          cycle
        end if

        ! Affectation et calcul numérique
        val = (rho_v / dt) * fields%v(i, j) * dvol
        ! Affectation et calcul numérique
        val = val + rho_v * cfg%gy * dvol

        ! Évaluation de la condition logique
        if (cfg%enable_energy) then
          ! Affectation et calcul numérique
          T_v = 0.5_dp * (fields%T(i, j) + fields%T(i, j - 1))
          ! Affectation et calcul numérique
          val = val + rho_v * cfg%beta * (T_v - cfg%T_ref) * cfg%gy * dvol
        end if

        ! Affectation et calcul numérique
        val = val + fields%fy_sigma(i, j) * dvol

        ! Évaluation de la condition logique
        if (cfg%solver_type == SOLVER_AUGMENTED_LAGRANGIAN) then
          ! Affectation et calcul numérique
          grad_p = (fields%p(i, j) - fields%p(i, j - 1)) * dx
          ! Affectation et calcul numérique
          val = val - grad_p
        end if

        ! Évaluation de la condition logique
        if (i == 1 .and. mesh%is_boundary_left .and. cfg%bc_left%btype == BC_DIRICHLET) then
          ! Calcul de la viscosité dynamique locale (harmonique aux nœuds)
          mu_W = fields%mu_node(i, j)
          ! Affectation et calcul numérique
          val = val + mu_W * (2.0_dp * cfg%bc_left%v / dx) * dy
        end if
        ! Évaluation de la condition logique
        if (i == mesh%nx_local .and. mesh%is_boundary_right .and. cfg%bc_right%btype == BC_DIRICHLET) then
          ! Calcul de la viscosité dynamique locale (harmonique aux nœuds)
          mu_E = fields%mu_node(i + 1, j)
          ! Affectation et calcul numérique
          val = val + mu_E * (2.0_dp * cfg%bc_right%v / dx) * dy
        end if

        ! Affectation et calcul numérique
        rhs(dof) = val
      end do
    end do

    ! =========================================================================
    ! 3. Second membre pour p (si monolithique)
    ! =========================================================================
    if (cfg%solver_type == SOLVER_MONOLITHIC) then
      ! Affectation et calcul numérique
      p_offset = this%n_owned_u + this%n_owned_v
      ! Boucle spatiale sur les cellules locales selon l'axe vertical Y (ordonnées)
      do j = 1, mesh%ny_local
        ! Boucle spatiale selon l'axe horizontal X (accès contigu colonne-major)
        do i = 1, mesh%nx_local
          ! Affectation et calcul numérique
          dof = p_offset + (j - 1) * mesh%nx_local + i
          ! Affectation et calcul numérique
          rhs(dof) = - (dvol / cfg%r_al) * fields%p(i, j)
        end do
      end do
    end if
  end subroutine ns_assemble_rhs

  !-----------------------------------------------------------------------------
  ! Mise à jour de la pression par l'étape d'Uzawa du Lagrangien augmenté :
  !   p^{m+1} = p^m - r_al * div(u^{m+1})
  !-----------------------------------------------------------------------------
  subroutine ns_update_pressure_al(this, mesh, cfg, fields)
    ! Déclaration de variable locale / paramètre de calcul
    class(navier_stokes_t), intent(in) :: this
    ! Déclaration de variable locale / paramètre de calcul
    type(mesh_2d_t), intent(in)        :: mesh
    ! Déclaration de variable locale / paramètre de calcul
    type(config_t), intent(in)         :: cfg
    ! Déclaration de variable locale / paramètre de calcul
    type(flow_fields_t), intent(inout) :: fields

    ! Déclaration de variable locale / paramètre de calcul
    integer :: i, j
    ! Déclaration de variable locale / paramètre de calcul
    real(dp) :: dx, dy, div_val

    ! Calcul de la métrique spatiale élémentaire
    dx = mesh%dx
    ! Calcul de la métrique spatiale élémentaire
    dy = mesh%dy

    ! Boucle spatiale sur les cellules locales selon l'axe vertical Y (ordonnées)
    do j = 1, mesh%ny_local
      ! Boucle spatiale selon l'axe horizontal X (accès contigu colonne-major)
      do i = 1, mesh%nx_local
        ! Affectation et calcul numérique
        div_val = (fields%u(i + 1, j) - fields%u(i, j)) / dx + &
                  (fields%v(i, j + 1) - fields%v(i, j)) / dy
        ! Actualisation de la pression statique au centre de la maille
        fields%p(i, j) = fields%p(i, j) - cfg%r_al * div_val
      end do
    end do

    ! Appel de sous-programme modulaire
    call sync_cell_scalar(mesh, fields%p)
  end subroutine ns_update_pressure_al

  !-----------------------------------------------------------------------------
  ! Assemblage des diagonales pour le préconditionnement Jacobi
  !-----------------------------------------------------------------------------
  subroutine ns_assemble_preconditioner_blocks(this, mesh, cfg, fields, dt)
    ! Déclaration de variable locale / paramètre de calcul
    class(navier_stokes_t), intent(inout) :: this
    ! Déclaration de variable locale / paramètre de calcul
    type(mesh_2d_t), intent(in)           :: mesh
    ! Déclaration de variable locale / paramètre de calcul
    type(config_t), intent(in)            :: cfg
    ! Déclaration de variable locale / paramètre de calcul
    type(flow_fields_t), intent(in)       :: fields
    ! Déclaration de variable locale / paramètre de calcul
    real(dp), intent(in)                  :: dt

    ! Déclaration de variable locale / paramètre de calcul
    real(dp) :: dx, dy, dvol, rho_u, rho_v, mu_c, d_val
    ! Déclaration de variable locale / paramètre de calcul
    integer :: i, j, dof, u_nx, v_ny
    ! Déclaration de variable locale / paramètre de calcul
    logical :: is_wall

    ! Calcul de la métrique spatiale élémentaire
    dx = mesh%dx
    ! Calcul de la métrique spatiale élémentaire
    dy = mesh%dy
    ! Affectation et calcul numérique
    dvol = dx * dy

    ! Affectation et calcul numérique
    u_nx = mesh%nx_local
    ! Affectation et calcul numérique
    if (mesh%is_boundary_right) u_nx = u_nx + 1

    ! Affectation et calcul numérique
    v_ny = mesh%ny_local
    ! Affectation et calcul numérique
    if (mesh%is_boundary_top) v_ny = v_ny + 1

    ! 1. Diagonale Bloc u
    do j = 1, mesh%ny_local
      ! Boucle spatiale selon l'axe horizontal X (accès contigu colonne-major)
      do i = 1, u_nx
        ! Affectation et calcul numérique
        dof = (j - 1) * u_nx + i
        ! Calcul de la masse volumique locale ou interpolée aux faces
        rho_u = fields%rho_u(i, j)
        ! Affectation et calcul numérique
        is_wall = (i == 1 .and. mesh%is_boundary_left) .or. &
                  ! Affectation et calcul numérique
                  (i == u_nx .and. mesh%is_boundary_right) .or. &
                  is_u_face_in_obstacle(mesh, cfg, i, j)

        ! Évaluation de la condition logique
        if (is_wall) then
          ! Mise à jour de la composante de vitesse discrète sur face décalée MAC
          this%diag_u(dof) = (rho_u / dt) * dvol
          cycle
        end if

        ! Calcul de la viscosité dynamique locale (harmonique aux nœuds)
        mu_c = fields%mu(i, j)
        ! Affectation et calcul numérique
        d_val = (rho_u / dt) * dvol + 4.0_dp * mu_c * (dy / dx) + 2.0_dp * mu_c * (dx / dy) + &
                2.0_dp * cfg%r_al * (dy / dx)
        ! Évaluation de la condition logique
        if (d_val > 1.0e-15_dp) then
          ! Mise à jour de la composante de vitesse discrète sur face décalée MAC
          this%diag_u(dof) = d_val
        ! Branche par défaut si condition non satisfaite
        else
          ! Mise à jour de la composante de vitesse discrète sur face décalée MAC
          this%diag_u(dof) = 1.0_dp
        end if
      end do
    end do

    ! 2. Diagonale Bloc v
    do j = 1, v_ny
      ! Boucle spatiale selon l'axe horizontal X (accès contigu colonne-major)
      do i = 1, mesh%nx_local
        ! Affectation et calcul numérique
        dof = (j - 1) * mesh%nx_local + i
        ! Calcul de la masse volumique locale ou interpolée aux faces
        rho_v = fields%rho_v(i, j)
        ! Affectation et calcul numérique
        is_wall = (j == 1 .and. mesh%is_boundary_bottom .and. &
                   ! Affectation et calcul numérique
                   (cfg%bc_bottom%btype == BC_DIRICHLET .or. cfg%bc_bottom%btype == BC_FREESLIP)) .or. &
                  ! Affectation et calcul numérique
                  (j == v_ny .and. mesh%is_boundary_top .and. &
                   ! Affectation et calcul numérique
                   (cfg%bc_top%btype == BC_DIRICHLET .or. cfg%bc_top%btype == BC_FREESLIP)) .or. &
                  is_v_face_in_obstacle(mesh, cfg, i, j)

        ! Évaluation de la condition logique
        if (is_wall) then
          ! Mise à jour de la composante de vitesse discrète sur face décalée MAC
          this%diag_v(dof) = (rho_v / dt) * dvol
          cycle
        end if

        ! Calcul de la viscosité dynamique locale (harmonique aux nœuds)
        mu_c = fields%mu(i, j)
        ! Affectation et calcul numérique
        d_val = (rho_v / dt) * dvol + 4.0_dp * mu_c * (dx / dy) + 2.0_dp * mu_c * (dy / dx) + &
                2.0_dp * cfg%r_al * (dx / dy)
        ! Évaluation de la condition logique
        if (d_val > 1.0e-15_dp) then
          ! Mise à jour de la composante de vitesse discrète sur face décalée MAC
          this%diag_v(dof) = d_val
        ! Branche par défaut si condition non satisfaite
        else
          ! Mise à jour de la composante de vitesse discrète sur face décalée MAC
          this%diag_v(dof) = 1.0_dp
        end if
      end do
    end do

    ! 3. Diagonale Bloc p (si monolithique)
    if (cfg%solver_type == SOLVER_MONOLITHIC) then
      ! Affectation et calcul numérique
      this%diag_p = dvol / cfg%r_al
    end if
  end subroutine ns_assemble_preconditioner_blocks

  !-----------------------------------------------------------------------------
  ! Résolution de l'énergie thermique (Boussinesq)
  !-----------------------------------------------------------------------------
  subroutine ns_solve_energy(this, mesh, cfg, fields, dt)
    ! Déclaration de variable locale / paramètre de calcul
    class(navier_stokes_t), intent(in) :: this
    ! Déclaration de variable locale / paramètre de calcul
    type(mesh_2d_t), intent(in)        :: mesh
    ! Déclaration de variable locale / paramètre de calcul
    type(config_t), intent(in)         :: cfg
    ! Déclaration de variable locale / paramètre de calcul
    type(flow_fields_t), intent(inout) :: fields
    ! Déclaration de variable locale / paramètre de calcul
    real(dp), intent(in)               :: dt

    ! Déclaration de variable locale / paramètre de calcul
    integer :: nx, ny, n_cells, i, j, c_dof, iter, max_it
    ! Déclaration de variable locale / paramètre de calcul
    real(dp) :: dx, dy, dvol, rho, Cp, k_th
    ! Déclaration de variable locale / paramètre de calcul
    real(dp) :: u_e, u_w, v_n, v_s
    ! Déclaration de variable locale / paramètre de calcul
    real(dp) :: d_e, d_w, d_n, d_s, conv_e, conv_w, conv_n, conv_s
    ! Déclaration de variable locale / paramètre de calcul
    real(dp) :: a_c, val
    ! Déclaration de variable locale / paramètre de calcul
    real(dp), allocatable :: a_diag(:), a_east(:), a_west(:), a_north(:), a_south(:)
    ! Déclaration de variable locale / paramètre de calcul
    real(dp), allocatable :: rhs(:), T_sol(:), r(:), p(:), v_vec(:), s_vec(:), t_vec(:)
    ! Déclaration de variable locale / paramètre de calcul
    real(dp), allocatable :: r0_hat(:), y_prec(:), z_prec(:)
    ! Déclaration de variable locale / paramètre de calcul
    real(dp) :: rho_old, alpha, omega, rho_new, r0v, s_norm, tt, ts, r_norm, res_init

    if (.not. cfg%enable_energy) return

    ! Affectation et calcul numérique
    nx = mesh%nx_local
    ! Affectation et calcul numérique
    ny = mesh%ny_local
    ! Affectation et calcul numérique
    n_cells = nx * ny

    ! Calcul de la métrique spatiale élémentaire
    dx = mesh%dx
    ! Calcul de la métrique spatiale élémentaire
    dy = mesh%dy
    ! Affectation et calcul numérique
    dvol = dx * dy
    ! Calcul de la masse volumique locale ou interpolée aux faces
    rho = fields%rho(1, 1)
    ! Affectation et calcul numérique
    Cp = cfg%Cp
    ! Affectation et calcul numérique
    k_th = cfg%k_thermal

    ! Allocation dynamique de tableau contigu en mémoire
    allocate(a_diag(n_cells), a_east(n_cells), a_west(n_cells), a_north(n_cells), a_south(n_cells))
    ! Allocation dynamique de tableau contigu en mémoire
    allocate(rhs(n_cells), T_sol(n_cells), r(n_cells), p(n_cells), v_vec(n_cells), s_vec(n_cells), t_vec(n_cells))
    ! Allocation dynamique de tableau contigu en mémoire
    allocate(r0_hat(n_cells), y_prec(n_cells), z_prec(n_cells))

    ! Affectation et calcul numérique
    a_diag = 0.0_dp
    ! Affectation et calcul numérique
    a_east = 0.0_dp
    ! Affectation et calcul numérique
    a_west = 0.0_dp
    ! Affectation et calcul numérique
    a_north = 0.0_dp
    ! Affectation et calcul numérique
    a_south = 0.0_dp
    ! Affectation et calcul numérique
    rhs = 0.0_dp

    ! Boucle spatiale sur les cellules locales selon l'axe vertical Y (ordonnées)
    do j = 1, ny
      ! Boucle spatiale selon l'axe horizontal X (accès contigu colonne-major)
      do i = 1, nx
        ! Affectation et calcul numérique
        c_dof = (j - 1) * nx + i
        ! Affectation et calcul numérique
        T_sol(c_dof) = fields%T(i, j)

        ! Affectation et calcul numérique
        a_c = (rho * Cp / dt) * dvol
        ! Affectation et calcul numérique
        rhs(c_dof) = (rho * Cp / dt) * dvol * fields%T(i, j)

        ! Affectation et calcul numérique
        u_e = fields%u(i + 1, j)
        ! Affectation et calcul numérique
        d_e = k_th * dy / dx
        ! Affectation et calcul numérique
        conv_e = rho * Cp * u_e * dy * 0.5_dp

        ! Affectation et calcul numérique
        u_w = fields%u(i, j)
        ! Affectation et calcul numérique
        d_w = k_th * dy / dx
        ! Affectation et calcul numérique
        conv_w = rho * Cp * u_w * dy * 0.5_dp

        ! Affectation et calcul numérique
        v_n = fields%v(i, j + 1)
        ! Affectation et calcul numérique
        d_n = k_th * dx / dy
        ! Affectation et calcul numérique
        conv_n = rho * Cp * v_n * dx * 0.5_dp

        ! Affectation et calcul numérique
        v_s = fields%v(i, j)
        ! Affectation et calcul numérique
        d_s = k_th * dx / dy
        ! Affectation et calcul numérique
        conv_s = rho * Cp * v_s * dx * 0.5_dp

        ! Évaluation de la condition logique
        if (i < nx) then
          ! Affectation et calcul numérique
          a_c = a_c + d_e + conv_e
          ! Affectation et calcul numérique
          a_east(c_dof) = -d_e + conv_e
        ! Branche conditionnelle alternative
        else if (mesh%is_boundary_right) then
          ! Évaluation de la condition logique
          if (cfg%bc_right%thermal_type == THERMAL_DIRICHLET) then
            ! Affectation et calcul numérique
            a_c = a_c + 2.0_dp * d_e
            ! Affectation et calcul numérique
            rhs(c_dof) = rhs(c_dof) + 2.0_dp * d_e * cfg%bc_right%T
          end if
        ! Branche par défaut si condition non satisfaite
        else
          ! Affectation et calcul numérique
          a_c = a_c + d_e + conv_e
          ! Affectation et calcul numérique
          rhs(c_dof) = rhs(c_dof) - (-d_e + conv_e) * fields%T(nx + 1, j)
        end if

        ! Évaluation de la condition logique
        if (i > 1) then
          ! Affectation et calcul numérique
          a_c = a_c + d_w - conv_w
          ! Affectation et calcul numérique
          a_west(c_dof) = -d_w - conv_w
        ! Branche conditionnelle alternative
        else if (mesh%is_boundary_left) then
          ! Évaluation de la condition logique
          if (cfg%bc_left%thermal_type == THERMAL_DIRICHLET) then
            ! Affectation et calcul numérique
            a_c = a_c + 2.0_dp * d_w
            ! Affectation et calcul numérique
            rhs(c_dof) = rhs(c_dof) + 2.0_dp * d_w * cfg%bc_left%T
          end if
        ! Branche par défaut si condition non satisfaite
        else
          ! Affectation et calcul numérique
          a_c = a_c + d_w - conv_w
          ! Affectation et calcul numérique
          rhs(c_dof) = rhs(c_dof) - (-d_w - conv_w) * fields%T(0, j)
        end if

        ! Évaluation de la condition logique
        if (j < ny) then
          ! Affectation et calcul numérique
          a_c = a_c + d_n + conv_n
          ! Affectation et calcul numérique
          a_north(c_dof) = -d_n + conv_n
        ! Branche conditionnelle alternative
        else if (mesh%is_boundary_top) then
          ! Évaluation de la condition logique
          if (cfg%bc_top%thermal_type == THERMAL_DIRICHLET) then
            ! Affectation et calcul numérique
            a_c = a_c + 2.0_dp * d_n
            ! Affectation et calcul numérique
            rhs(c_dof) = rhs(c_dof) + 2.0_dp * d_n * cfg%bc_top%T
          end if
        ! Branche par défaut si condition non satisfaite
        else
          ! Affectation et calcul numérique
          a_c = a_c + d_n + conv_n
          ! Affectation et calcul numérique
          rhs(c_dof) = rhs(c_dof) - (-d_n + conv_n) * fields%T(i, ny + 1)
        end if

        ! Évaluation de la condition logique
        if (j > 1) then
          ! Affectation et calcul numérique
          a_c = a_c + d_s - conv_s
          ! Affectation et calcul numérique
          a_south(c_dof) = -d_s - conv_s
        ! Branche conditionnelle alternative
        else if (mesh%is_boundary_bottom) then
          ! Évaluation de la condition logique
          if (cfg%bc_bottom%thermal_type == THERMAL_DIRICHLET) then
            ! Affectation et calcul numérique
            a_c = a_c + 2.0_dp * d_s
            ! Affectation et calcul numérique
            rhs(c_dof) = rhs(c_dof) + 2.0_dp * d_s * cfg%bc_bottom%T
          end if
        ! Branche par défaut si condition non satisfaite
        else
          ! Affectation et calcul numérique
          a_c = a_c + d_s - conv_s
          ! Affectation et calcul numérique
          rhs(c_dof) = rhs(c_dof) - (-d_s - conv_s) * fields%T(i, 0)
        end if

        ! Affectation et calcul numérique
        a_diag(c_dof) = a_c
      end do
    end do

    ! Résolution itérative BiCGSTAB avec préconditionnement Jacobi
    do j = 1, ny
      ! Boucle spatiale selon l'axe horizontal X (accès contigu colonne-major)
      do i = 1, nx
        ! Affectation et calcul numérique
        c_dof = (j - 1) * nx + i
        ! Affectation et calcul numérique
        val = a_diag(c_dof) * T_sol(c_dof)
        ! Affectation et calcul numérique
        if (i < nx) val = val + a_east(c_dof) * T_sol(c_dof + 1)
        ! Affectation et calcul numérique
        if (i > 1)  val = val + a_west(c_dof) * T_sol(c_dof - 1)
        ! Affectation et calcul numérique
        if (j < ny) val = val + a_north(c_dof) * T_sol(c_dof + nx)
        ! Affectation et calcul numérique
        if (j > 1)  val = val + a_south(c_dof) * T_sol(c_dof - nx)
        ! Affectation et calcul numérique
        r(c_dof) = rhs(c_dof) - val
      end do
    end do

    ! Affectation et calcul numérique
    r0_hat = r
    ! Affectation et calcul numérique
    p = r
    ! Calcul de la masse volumique locale ou interpolée aux faces
    rho_old = 1.0_dp
    ! Affectation et calcul numérique
    alpha = 1.0_dp
    ! Affectation et calcul numérique
    omega = 1.0_dp

    ! Évaluation de la norme du résidu ou de l'erreur numérique
    res_init = sqrt(sum(r*r))
    ! Évaluation de la condition logique
    if (res_init > 1.0e-13_dp) then
      ! Affectation et calcul numérique
      max_it = 500
      ! Boucle d'itération numérique
      do iter = 1, max_it
        ! Calcul de la masse volumique locale ou interpolée aux faces
        rho_new = sum(r0_hat * r)
        if (abs(rho_new) < 1.0e-30_dp) exit

        ! Évaluation de la condition logique
        if (iter > 1) then
          ! Affectation et calcul numérique
          p = r + (rho_new / rho_old) * (alpha / omega) * (p - omega * v_vec)
        ! Branche par défaut si condition non satisfaite
        else
          ! Affectation et calcul numérique
          p = r
        end if

        ! Affectation et calcul numérique
        y_prec = p / a_diag

        ! Boucle spatiale sur les cellules locales selon l'axe vertical Y (ordonnées)
        do j = 1, ny
          ! Boucle spatiale selon l'axe horizontal X (accès contigu colonne-major)
          do i = 1, nx
            ! Affectation et calcul numérique
            c_dof = (j - 1) * nx + i
            ! Affectation et calcul numérique
            val = a_diag(c_dof) * y_prec(c_dof)
            ! Affectation et calcul numérique
            if (i < nx) val = val + a_east(c_dof) * y_prec(c_dof + 1)
            ! Affectation et calcul numérique
            if (i > 1)  val = val + a_west(c_dof) * y_prec(c_dof - 1)
            ! Affectation et calcul numérique
            if (j < ny) val = val + a_north(c_dof) * y_prec(c_dof + nx)
            ! Affectation et calcul numérique
            if (j > 1)  val = val + a_south(c_dof) * y_prec(c_dof - nx)
            ! Affectation et calcul numérique
            v_vec(c_dof) = val
          end do
        end do

        ! Affectation et calcul numérique
        r0v = sum(r0_hat * v_vec)
        if (abs(r0v) < 1.0e-30_dp) exit
        ! Affectation et calcul numérique
        alpha = rho_new / r0v

        ! Affectation et calcul numérique
        s_vec = r - alpha * v_vec
        ! Évaluation de la norme du résidu ou de l'erreur numérique
        s_norm = sqrt(sum(s_vec * s_vec))
        ! Évaluation de la condition logique
        if (s_norm / res_init < 1.0e-8_dp) then
          ! Affectation et calcul numérique
          T_sol = T_sol + alpha * y_prec
          exit
        end if

        ! Affectation et calcul numérique
        z_prec = s_vec / a_diag

        ! Boucle spatiale sur les cellules locales selon l'axe vertical Y (ordonnées)
        do j = 1, ny
          ! Boucle spatiale selon l'axe horizontal X (accès contigu colonne-major)
          do i = 1, nx
            ! Affectation et calcul numérique
            c_dof = (j - 1) * nx + i
            ! Affectation et calcul numérique
            val = a_diag(c_dof) * z_prec(c_dof)
            ! Affectation et calcul numérique
            if (i < nx) val = val + a_east(c_dof) * z_prec(c_dof + 1)
            ! Affectation et calcul numérique
            if (i > 1)  val = val + a_west(c_dof) * z_prec(c_dof - 1)
            ! Affectation et calcul numérique
            if (j < ny) val = val + a_north(c_dof) * z_prec(c_dof + nx)
            ! Affectation et calcul numérique
            if (j > 1)  val = val + a_south(c_dof) * z_prec(c_dof - nx)
            ! Affectation et calcul numérique
            t_vec(c_dof) = val
          end do
        end do

        ! Affectation et calcul numérique
        tt = sum(t_vec * t_vec)
        ! Affectation et calcul numérique
        ts = sum(t_vec * s_vec)
        ! Évaluation de la condition logique
        if (tt > 1.0e-30_dp) then
          ! Affectation et calcul numérique
          omega = ts / tt
        ! Branche par défaut si condition non satisfaite
        else
          ! Affectation et calcul numérique
          omega = 0.0_dp
        end if

        ! Affectation et calcul numérique
        T_sol = T_sol + alpha * y_prec + omega * z_prec
        ! Affectation et calcul numérique
        r = s_vec - omega * t_vec
        ! Calcul de la masse volumique locale ou interpolée aux faces
        rho_old = rho_new

        ! Évaluation de la norme du résidu ou de l'erreur numérique
        r_norm = sqrt(sum(r * r))
        if (r_norm / res_init < 1.0e-8_dp) exit
      end do
    end if

    ! Boucle spatiale sur les cellules locales selon l'axe vertical Y (ordonnées)
    do j = 1, ny
      ! Boucle spatiale selon l'axe horizontal X (accès contigu colonne-major)
      do i = 1, nx
        ! Évolution du champ scalaire de température
        fields%T(i, j) = T_sol((j - 1) * nx + i)
      end do
    end do

    ! Appel de sous-programme modulaire
    call sync_cell_scalar(mesh, fields%T)
    ! Appel de sous-programme modulaire
    call apply_thermal_boundary_conditions(mesh, cfg, fields)

    ! Libération de la mémoire allouée
    deallocate(a_diag, a_east, a_west, a_north, a_south)
    ! Libération de la mémoire allouée
    deallocate(rhs, T_sol, r, p, v_vec, s_vec, t_vec, r0_hat, y_prec, z_prec)
  end subroutine ns_solve_energy

  !-----------------------------------------------------------------------------
  ! Calcul du nombre de Nusselt moyen sur la paroi chaude (gauche)
  !-----------------------------------------------------------------------------
  function ns_compute_mean_nusselt_hot_wall(this, mesh, cfg, fields) result(nu_mean)
    ! Déclaration de variable locale / paramètre de calcul
    class(navier_stokes_t), intent(in) :: this
    ! Déclaration de variable locale / paramètre de calcul
    type(mesh_2d_t), intent(in)        :: mesh
    ! Déclaration de variable locale / paramètre de calcul
    type(config_t), intent(in)         :: cfg
    ! Déclaration de variable locale / paramètre de calcul
    type(flow_fields_t), intent(in)    :: fields
    ! Déclaration de variable locale / paramètre de calcul
    real(dp)                           :: nu_mean

    ! Déclaration de variable locale / paramètre de calcul
    integer :: j, count_loc
    ! Déclaration de variable locale / paramètre de calcul
    real(dp) :: local_nu_sum, global_nu_sum, global_count
    ! Déclaration de variable locale / paramètre de calcul
    real(dp) :: T_w, T_1, delta_T, nu_local

    ! Affectation et calcul numérique
    local_nu_sum = 0.0_dp
    ! Affectation et calcul numérique
    count_loc = 0

    ! Évaluation de la condition logique
    if (mesh%is_boundary_left) then
      ! Affectation et calcul numérique
      delta_T = abs(cfg%bc_left%T - cfg%bc_right%T)
      ! Affectation et calcul numérique
      if (delta_T < 1.0e-12_dp) delta_T = 1.0_dp

      ! Boucle spatiale sur les cellules locales selon l'axe vertical Y (ordonnées)
      do j = 1, mesh%ny_local
        ! Affectation et calcul numérique
        T_w = cfg%bc_left%T
        ! Affectation et calcul numérique
        T_1 = fields%T(1, j)
        ! Affectation et calcul numérique
        nu_local = 2.0_dp * (T_w - T_1) / mesh%dx * (mesh%lx / delta_T)
        ! Affectation et calcul numérique
        local_nu_sum = local_nu_sum + nu_local
        ! Affectation et calcul numérique
        count_loc = count_loc + 1
      end do
    end if

    ! Affectation et calcul numérique
    global_nu_sum = global_sum(mesh, local_nu_sum)
    ! Affectation et calcul numérique
    global_count = global_sum(mesh, real(count_loc, dp))

    ! Évaluation de la condition logique
    if (global_count > 0.0_dp) then
      ! Affectation et calcul numérique
      nu_mean = global_nu_sum / global_count
    ! Branche par défaut si condition non satisfaite
    else
      ! Affectation et calcul numérique
      nu_mean = 1.0_dp
    end if
  end function ns_compute_mean_nusselt_hot_wall

end module module_navier_stokes

\end{lstlisting}

\chapter{Reconstruction Géométrique d'Interface VOF en Fortran 2018 (module_interface_reconstruction.f90)}
\section{Algorithmes Géométriques dans \texttt{module\_interface\_reconstruction.f90}}
Le module \texttt{src/module\_interface\_reconstruction.f90} (569 lignes) implémente en Fortran 2018 :
\begin{itemize}
  \item \texttt{reconstruct\_interface\_normals} : sélectionne entre le filtre $3\times 3$ de Parker-Youngs (\texttt{RECON\_PLIC}), l'algorithme à 6 candidats exacts d'\textbf{ELVIRA} (\texttt{RECON\_ELVIRA}), l'optimisation quadratique non linéaire de \textbf{LVIRA} (\texttt{RECON\_LVIRA}) et l'alignement axial \textbf{SLIC} (\texttt{RECON\_SLIC}).
  \item \texttt{plic\_alpha\_from\_c} et \texttt{plic\_volume\_from\_alpha} : implémentent les formules analytiques directes et inverses de Scardovelli \& Zaleski (2000) reliant la constante de droite $\alpha$ à l'aire tronquée $C$ dans le carré unité $[0, 1]^2$.
  \item \texttt{plic\_flux\_1d} : calcule géométriquement le volume exact de phase 1 contenu dans la bande lagrangienne $[1 - \text{cfl}, 1] \times [0, 1]$ traversant une face de maille au cours du pas de temps $\Delta t$.
\end{itemize}

\section{Répertoire des Types Dérivés, Sous-Routines et Fonctions de \texttt{src/module\_interface\_reconstruction.f90}}
\begin{center}\small
\begin{tabular}{clll}
\toprule
\textbf{Ligne} & \textbf{Nature} & \textbf{Identifiant Fortran 2018} & \textbf{Interface / Arguments} \\
\midrule
2 & module & \texttt{module\_interface\_reconstruction} & Unité modulaire Fortran 2018 \\
22 & type dérivé & \texttt{interface\_segment\_t} & Structure de données Fortran 2018 \\
32 & function & \texttt{cut\_box\_area} & Arguments : (xmin, xmax, ymin, ymax, nx, ny, alpha) \\
104 & function & \texttt{compute\_alpha} & Arguments : (xmin, xmax, ymin, ymax, nx, ny, c\_target) \\
149 & subroutine & \texttt{reconstruct\_slic} & Arguments : (mesh, C, segments) \\
212 & subroutine & \texttt{reconstruct\_plic} & Arguments : (mesh, C, segments) \\
275 & subroutine & \texttt{reconstruct\_elvira} & Arguments : (mesh, C, segments) \\
389 & subroutine & \texttt{reconstruct\_lvira} & Arguments : (mesh, C, segments) \\
509 & subroutine & \texttt{reconstruct\_interface} & Arguments : (method\_name, mesh, C, segments) \\
592 & function & \texttt{evaluate\_limiter} & Arguments : (limiter\_name, r, c\_donor, c\_upwind, c\_a...) \\
\bottomrule
\end{tabular}
\end{center}

\section{Code Source Intégral Annoté : \texttt{src/module\_interface\_reconstruction.f90}}
Le listing ci-dessous présente l'intégralité des 633 lignes du fichier \texttt{src/module\_interface\_reconstruction.f90} de l'implémentation Fortran 2018 :

\begin{lstlisting}[language={[95]Fortran}]
! [Procédure/Module] module module_interface_reconstruction
module module_interface_reconstruction
  use module_precision, only: dp
  use module_mesh, only: mesh_2d_t, x_face, y_face
  implicit none
  private

  ! Types publics
  public :: interface_segment_t

  ! Procédures publiques
  public :: cut_box_area
  public :: compute_alpha
  public :: reconstruct_slic
  public :: reconstruct_plic
  public :: reconstruct_elvira
  public :: reconstruct_lvira
  public :: reconstruct_interface
  public :: evaluate_limiter

  ! Type représentant un segment de droite d'interface affine : nx*x + ny*y + alpha = 0
  type :: interface_segment_t
     real(dp) :: nx = 0.0_dp
     real(dp) :: ny = 0.0_dp
     real(dp) :: alpha = 0.0_dp
  end type interface_segment_t

contains

  ! [Procédure/Module] function cut_box_area
  ! Calcule l'aire exacte de découpe d'un rectangle par un demi-plan (Sutherland-Hodgman + Shoelace)
  function cut_box_area(xmin, xmax, ymin, ymax, nx, ny, alpha) result(area)
    real(dp), intent(in) :: xmin, xmax, ymin, ymax, nx, ny, alpha
    real(dp) :: area

    real(dp) :: bx(4), by(4)
    real(dp) :: d(4)
    real(dp) :: cx(8), cy(8)
    integer  :: n_clipped, i, next_i, p, next_p
    real(dp) :: d1, d2, t
    real(dp) :: sum_cross

    if ((xmax - xmin) <= 1.0e-14_dp .or. (ymax - ymin) <= 1.0e-14_dp) then
       area = 0.0_dp
       return
    end if

    ! Définition des 4 sommets du rectangle en sens trigonométrique
    bx(1) = xmin; by(1) = ymin
    bx(2) = xmax; by(2) = ymin
    bx(3) = xmax; by(3) = ymax
    bx(4) = xmin; by(4) = ymax

    do i = 1, 4
       d(i) = nx * bx(i) + ny * by(i) + alpha
    end do

    n_clipped = 0
    do i = 1, 4
       next_i = mod(i, 4) + 1
       d1 = d(i)
       d2 = d(next_i)

       if (d1 <= 0.0_dp) then
          if (d2 <= 0.0_dp) then
             n_clipped = n_clipped + 1
             cx(n_clipped) = bx(next_i)
             cy(n_clipped) = by(next_i)
          else
             t = -d1 / (d2 - d1)
             n_clipped = n_clipped + 1
             cx(n_clipped) = bx(i) + t * (bx(next_i) - bx(i))
             cy(n_clipped) = by(i) + t * (by(next_i) - by(i))
          end if
       else
          if (d2 <= 0.0_dp) then
             t = -d1 / (d2 - d1)
             n_clipped = n_clipped + 1
             cx(n_clipped) = bx(i) + t * (bx(next_i) - bx(i))
             cy(n_clipped) = by(i) + t * (by(next_i) - by(i))
             n_clipped = n_clipped + 1
             cx(n_clipped) = bx(next_i)
             cy(n_clipped) = by(next_i)
          end if
       end if
    end do

    if (n_clipped < 3) then
       area = 0.0_dp
       return
    end if

    ! Formule de Gauss / Shoelace pour l'aire du polygone
    sum_cross = 0.0_dp
    do p = 1, n_clipped
       next_p = mod(p, n_clipped) + 1
       sum_cross = sum_cross + (cx(p) * cy(next_p) - cx(next_p) * cy(p))
    end do
    area = 0.5_dp * abs(sum_cross)
  end function cut_box_area

  ! [Procédure/Module] function compute_alpha
  ! Inversion géométrique par bisection : C_target -> paramètre alpha
  function compute_alpha(xmin, xmax, ymin, ymax, nx, ny, c_target) result(alpha)
    real(dp), intent(in) :: xmin, xmax, ymin, ymax, nx, ny, c_target
    real(dp) :: alpha

    real(dp) :: total_area, target_area
    real(dp) :: v1, v2, v3, v4, min_val, max_val
    real(dp) :: alpha_low, alpha_high, alpha_mid, a
    integer  :: iter

    total_area = (xmax - xmin) * (ymax - ymin)
    if (c_target <= 1.0e-7_dp) then
       alpha = 1.0e9_dp
       return
    end if
    if (c_target >= 1.0_dp - 1.0e-7_dp) then
       alpha = -1.0e9_dp
       return
    end if

    target_area = max(0.0_dp, min(1.0_dp, c_target)) * total_area

    v1 = nx * xmin + ny * ymin
    v2 = nx * xmax + ny * ymin
    v3 = nx * xmax + ny * ymax
    v4 = nx * xmin + ny * ymax

    min_val = min(v1, min(v2, min(v3, v4)))
    max_val = max(v1, max(v2, max(v3, v4)))

    alpha_low  = -max_val - 1.0e-12_dp
    alpha_high = -min_val + 1.0e-12_dp

    do iter = 1, 45
       alpha_mid = 0.5_dp * (alpha_low + alpha_high)
       a = cut_box_area(xmin, xmax, ymin, ymax, nx, ny, alpha_mid)
       if (a < target_area) then
          alpha_high = alpha_mid
       else
          alpha_low = alpha_mid
       end if
    end do
    alpha = 0.5_dp * (alpha_low + alpha_high)
  end function compute_alpha

  ! [Procédure/Module] subroutine reconstruct_slic
  subroutine reconstruct_slic(mesh, C, segments)
    type(mesh_2d_t), intent(in) :: mesh
    real(dp), intent(in) :: C(0:mesh%nx_local+1, 0:mesh%ny_local+1)
    type(interface_segment_t), intent(out) :: segments(0:mesh%nx_local+1, 0:mesh%ny_local+1)

    integer :: i, j, nx, ny
    real(dp) :: c_val, xmin, xmax, ymin, ymax, grad_x, grad_y, snx, sny, a

    nx = mesh%nx_local
    ny = mesh%ny_local

    do j = 0, ny + 1
       do i = 0, nx + 1
          c_val = C(i, j)
          if (c_val <= 1.0e-6_dp) then
             segments(i, j)%nx = 1.0_dp
             segments(i, j)%ny = 0.0_dp
             segments(i, j)%alpha = 1.0e9_dp
             cycle
          end if
          if (c_val >= 1.0_dp - 1.0e-6_dp) then
             segments(i, j)%nx = 1.0_dp
             segments(i, j)%ny = 0.0_dp
             segments(i, j)%alpha = -1.0e9_dp
             cycle
          end if

          xmin = x_face(mesh, i)
          xmax = x_face(mesh, i + 1)
          ymin = y_face(mesh, j)
          ymax = y_face(mesh, j + 1)

          if (i > 0 .and. i < nx + 1 .and. j > 0 .and. j < ny + 1) then
             grad_x = abs(C(i + 1, j) - C(i - 1, j))
             grad_y = abs(C(i, j + 1) - C(i, j - 1))
          else
             grad_x = 1.0_dp
             grad_y = 0.0_dp
          end if

          if (grad_x >= grad_y) then
             if (i > 0 .and. C(i - 1, j) >= C(i + 1, j)) then
                snx = 1.0_dp; sny = 0.0_dp
             else
                snx = -1.0_dp; sny = 0.0_dp
             end if
          else
             if (j > 0 .and. C(i, j - 1) >= C(i, j + 1)) then
                snx = 0.0_dp; sny = 1.0_dp
             else
                snx = 0.0_dp; sny = -1.0_dp
             end if
          end if

          a = compute_alpha(xmin, xmax, ymin, ymax, snx, sny, c_val)
          segments(i, j)%nx = snx
          segments(i, j)%ny = sny
          segments(i, j)%alpha = a
       end do
    end do
  end subroutine reconstruct_slic

  ! [Procédure/Module] subroutine reconstruct_plic
  subroutine reconstruct_plic(mesh, C, segments)
    type(mesh_2d_t), intent(in) :: mesh
    real(dp), intent(in) :: C(0:mesh%nx_local+1, 0:mesh%ny_local+1)
    type(interface_segment_t), intent(out) :: segments(0:mesh%nx_local+1, 0:mesh%ny_local+1)

    integer :: i, j, nx, ny
    real(dp) :: c_val, dx, dy, xmin, xmax, ymin, ymax
    real(dp) :: c_mm, c_0m, c_pm, c_m0, c_p0, c_mp, c_0p, c_pp
    real(dp) :: mx, my, norm, pnx, pny, a

    nx = mesh%nx_local
    ny = mesh%ny_local
    dx = mesh%dx
    dy = mesh%dy

    do j = 0, ny + 1
       do i = 0, nx + 1
          c_val = C(i, j)
          if (c_val <= 1.0e-6_dp) then
             segments(i, j)%nx = 1.0_dp
             segments(i, j)%ny = 0.0_dp
             segments(i, j)%alpha = 1.0e9_dp
             cycle
          end if
          if (c_val >= 1.0_dp - 1.0e-6_dp) then
             segments(i, j)%nx = 1.0_dp
             segments(i, j)%ny = 0.0_dp
             segments(i, j)%alpha = -1.0e9_dp
             cycle
          end if

          if (i >= 1 .and. i <= nx .and. j >= 1 .and. j <= ny) then
             c_mm = C(i - 1, j - 1); c_0m = C(i,     j - 1); c_pm = C(i + 1, j - 1)
             c_m0 = C(i - 1, j    );                         c_p0 = C(i + 1, j    )
             c_mp = C(i - 1, j + 1); c_0p = C(i,     j + 1); c_pp = C(i + 1, j + 1)

             mx = -((c_pp + 2.0_dp*c_p0 + c_pm) - (c_mp + 2.0_dp*c_m0 + c_mm)) / (8.0_dp * dx)
             my = -((c_pp + 2.0_dp*c_0p + c_mp) - (c_pm + 2.0_dp*c_0m + c_mm)) / (8.0_dp * dy)
          else
             mx = 1.0_dp; my = 0.0_dp
          end if

          norm = hypot(mx, my)
          if (norm > 1.0e-12_dp) then
             pnx = mx / norm; pny = my / norm
          else
             pnx = 1.0_dp; pny = 0.0_dp
          end if

          xmin = x_face(mesh, i)
          xmax = x_face(mesh, i + 1)
          ymin = y_face(mesh, j)
          ymax = y_face(mesh, j + 1)

          a = compute_alpha(xmin, xmax, ymin, ymax, pnx, pny, c_val)
          segments(i, j)%nx = pnx
          segments(i, j)%ny = pny
          segments(i, j)%alpha = a
       end do
    end do
  end subroutine reconstruct_plic

  ! [Procédure/Module] subroutine reconstruct_elvira
  subroutine reconstruct_elvira(mesh, C, segments)
    type(mesh_2d_t), intent(in) :: mesh
    real(dp), intent(in) :: C(0:mesh%nx_local+1, 0:mesh%ny_local+1)
    type(interface_segment_t), intent(out) :: segments(0:mesh%nx_local+1, 0:mesh%ny_local+1)

    integer :: i, j, nx, ny, di, dj, k
    real(dp) :: c_val, dx, dy, xmin, xmax, ymin, ymax
    real(dp) :: H_m1, H_0, H_p1, W_m1, W_0, W_p1
    real(dp) :: sgn_x, sgn_y, mx, my
    real(dp) :: cand_nx(7), cand_ny(7)
    real(dp) :: best_err, err, cnorm, cnx, cny, calpha, diff
    real(dp) :: best_nx, best_ny, best_alpha
    real(dp) :: v_xmin, v_xmax, v_ymin, v_ymax, v_target, v_recon

    nx = mesh%nx_local
    ny = mesh%ny_local
    dx = mesh%dx
    dy = mesh%dy

    do j = 0, ny + 1
       do i = 0, nx + 1
          c_val = C(i, j)
          if (c_val <= 1.0e-6_dp) then
             segments(i, j)%nx = 1.0_dp
             segments(i, j)%ny = 0.0_dp
             segments(i, j)%alpha = 1.0e9_dp
             cycle
          end if
          if (c_val >= 1.0_dp - 1.0e-6_dp) then
             segments(i, j)%nx = 1.0_dp
             segments(i, j)%ny = 0.0_dp
             segments(i, j)%alpha = -1.0e9_dp
             cycle
          end if

          xmin = x_face(mesh, i)
          xmax = x_face(mesh, i + 1)
          ymin = y_face(mesh, j)
          ymax = y_face(mesh, j + 1)

          if (i >= 1 .and. i <= nx .and. j >= 1 .and. j <= ny) then
             H_m1 = (C(i - 1, j - 1) + C(i - 1, j) + C(i - 1, j + 1)) * dy
             H_0  = (C(i,     j - 1) + C(i,     j) + C(i,     j + 1)) * dy
             H_p1 = (C(i + 1, j - 1) + C(i + 1, j) + C(i + 1, j + 1)) * dy

             W_m1 = (C(i - 1, j - 1) + C(i, j - 1) + C(i + 1, j - 1)) * dx
             W_0  = (C(i - 1, j    ) + C(i, j    ) + C(i + 1, j    )) * dx
             W_p1 = (C(i - 1, j + 1) + C(i, j + 1) + C(i + 1, j + 1)) * dx

             if (C(i, j + 1) < C(i, j - 1)) then; sgn_y = 1.0_dp; else; sgn_y = -1.0_dp; end if
             if (C(i + 1, j) < C(i - 1, j)) then; sgn_x = 1.0_dp; else; sgn_x = -1.0_dp; end if

             ! 6 candidats Height Functions
             cand_nx(1) = -sgn_y * (H_p1 - H_m1) / (2.0_dp * dx); cand_ny(1) = sgn_y * 1.0_dp
             cand_nx(2) = -sgn_y * (H_p1 - H_0)  / (dx);          cand_ny(2) = sgn_y * 1.0_dp
             cand_nx(3) = -sgn_y * (H_0  - H_m1) / (dx);          cand_ny(3) = sgn_y * 1.0_dp

             cand_nx(4) = sgn_x * 1.0_dp; cand_ny(4) = -sgn_x * (W_p1 - W_m1) / (2.0_dp * dy)
             cand_nx(5) = sgn_x * 1.0_dp; cand_ny(5) = -sgn_x * (W_p1 - W_0)  / (dy)
             cand_nx(6) = sgn_x * 1.0_dp; cand_ny(6) = -sgn_x * (W_0  - W_m1) / (dy)

             ! Candidat 7 : Youngs
             mx = -((C(i+1,j+1) + 2.0_dp*C(i+1,j) + C(i+1,j-1)) - (C(i-1,j+1) + 2.0_dp*C(i-1,j) + C(i-1,j-1))) / (8.0_dp * dx)
             my = -((C(i+1,j+1) + 2.0_dp*C(i,j+1) + C(i-1,j+1)) - (C(i+1,j-1) + 2.0_dp*C(i,j-1) + C(i-1,j-1))) / (8.0_dp * dy)
             cand_nx(7) = mx; cand_ny(7) = my

             best_err = 1.0e30_dp
             best_nx = 1.0_dp; best_ny = 0.0_dp; best_alpha = 0.0_dp

             do k = 1, 7
                cnorm = hypot(cand_nx(k), cand_ny(k))
                if (cnorm < 1.0e-12_dp) cycle
                cnx = cand_nx(k) / cnorm
                cny = cand_ny(k) / cnorm

                calpha = compute_alpha(xmin, xmax, ymin, ymax, cnx, cny, c_val)

                err = 0.0_dp
                do dj = -1, 1
                   do di = -1, 1
                      if (di == 0 .and. dj == 0) cycle
                      v_xmin = x_face(mesh, i + di)
                      v_xmax = x_face(mesh, i + di + 1)
                      v_ymin = y_face(mesh, j + dj)
                      v_ymax = y_face(mesh, j + dj + 1)

                      v_target = C(i + di, j + dj) * dx * dy
                      v_recon  = cut_box_area(v_xmin, v_xmax, v_ymin, v_ymax, cnx, cny, calpha)
                      diff = v_recon - v_target
                      err = err + diff * diff
                   end do
                end do

                if (err < best_err) then
                   best_err = err
                   best_nx = cnx
                   best_ny = cny
                   best_alpha = calpha
                end if
             end do

             segments(i, j)%nx = best_nx
             segments(i, j)%ny = best_ny
             segments(i, j)%alpha = best_alpha
          else
             segments(i, j)%nx = 1.0_dp
             segments(i, j)%ny = 0.0_dp
             segments(i, j)%alpha = compute_alpha(xmin, xmax, ymin, ymax, 1.0_dp, 0.0_dp, c_val)
          end if
       end do
    end do
  end subroutine reconstruct_elvira

  ! [Procédure/Module] subroutine reconstruct_lvira
  subroutine reconstruct_lvira(mesh, C, segments)
    type(mesh_2d_t), intent(in) :: mesh
    real(dp), intent(in) :: C(0:mesh%nx_local+1, 0:mesh%ny_local+1)
    type(interface_segment_t), intent(out) :: segments(0:mesh%nx_local+1, 0:mesh%ny_local+1)

    integer :: i, j, nx, ny, iter, di, dj
    real(dp) :: c_val, dx, dy, xmin, xmax, ymin, ymax, theta0
    real(dp) :: invphi, invphi2, a_pt, b_pt, h_pt, c_pt, d_pt, fc, fd
    real(dp) :: best_theta, best_nx, best_ny, best_alpha
    real(dp) :: v_xmin, v_xmax, v_ymin, v_ymax, v_target, v_recon, diff

    ! Étape 1 : Initialisation par ELVIRA
    call reconstruct_elvira(mesh, C, segments)

    nx = mesh%nx_local
    ny = mesh%ny_local
    dx = mesh%dx
    dy = mesh%dy
    invphi = (sqrt(5.0_dp) - 1.0_dp) / 2.0_dp
    invphi2 = (3.0_dp - sqrt(5.0_dp)) / 2.0_dp

    do j = 1, ny
       do i = 1, nx
          c_val = C(i, j)
          if (c_val <= 1.0e-6_dp .or. c_val >= 1.0_dp - 1.0e-6_dp) cycle

          xmin = x_face(mesh, i)
          xmax = x_face(mesh, i + 1)
          ymin = y_face(mesh, j)
          ymax = y_face(mesh, j + 1)

          theta0 = atan2(segments(i, j)%ny, segments(i, j)%nx)
          a_pt = theta0 - 0.5_dp
          b_pt = theta0 + 0.5_dp
          h_pt = b_pt - a_pt

          c_pt = a_pt + invphi2 * h_pt
          d_pt = a_pt + invphi * h_pt

          ! Évaluation fc
          fc = 0.0_dp
          do dj = -1, 1
             do di = -1, 1
                if (di == 0 .and. dj == 0) cycle
                v_xmin = x_face(mesh, i + di); v_xmax = x_face(mesh, i + di + 1)
                v_ymin = y_face(mesh, j + dj); v_ymax = y_face(mesh, j + dj + 1)
                v_target = C(i + di, j + dj) * dx * dy
                v_recon  = cut_box_area(v_xmin, v_xmax, v_ymin, v_ymax, cos(c_pt), sin(c_pt), &
                                       compute_alpha(xmin, xmax, ymin, ymax, cos(c_pt), sin(c_pt), c_val))
                diff = v_recon - v_target
                fc = fc + diff * diff
             end do
          end do

          ! Évaluation fd
          fd = 0.0_dp
          do dj = -1, 1
             do di = -1, 1
                if (di == 0 .and. dj == 0) cycle
                v_xmin = x_face(mesh, i + di); v_xmax = x_face(mesh, i + di + 1)
                v_ymin = y_face(mesh, j + dj); v_ymax = y_face(mesh, j + dj + 1)
                v_target = C(i + di, j + dj) * dx * dy
                v_recon  = cut_box_area(v_xmin, v_xmax, v_ymin, v_ymax, cos(d_pt), sin(d_pt), &
                                       compute_alpha(xmin, xmax, ymin, ymax, cos(d_pt), sin(d_pt), c_val))
                diff = v_recon - v_target
                fd = fd + diff * diff
             end do
          end do

          do iter = 1, 10
             if (fc < fd) then
                b_pt = d_pt; d_pt = c_pt; fd = fc
                h_pt = invphi * h_pt
                c_pt = a_pt + invphi2 * h_pt
                fc = 0.0_dp
                do dj = -1, 1
                   do di = -1, 1
                      if (di == 0 .and. dj == 0) cycle
                      v_xmin = x_face(mesh, i + di); v_xmax = x_face(mesh, i + di + 1)
                      v_ymin = y_face(mesh, j + dj); v_ymax = y_face(mesh, j + dj + 1)
                      v_target = C(i + di, j + dj) * dx * dy
                      v_recon  = cut_box_area(v_xmin, v_xmax, v_ymin, v_ymax, cos(c_pt), sin(c_pt), &
                                             compute_alpha(xmin, xmax, ymin, ymax, cos(c_pt), sin(c_pt), c_val))
                      diff = v_recon - v_target
                      fc = fc + diff * diff
                   end do
                end do
             else
                a_pt = c_pt; c_pt = d_pt; fc = fd
                h_pt = invphi * h_pt
                d_pt = a_pt + invphi * h_pt
                fd = 0.0_dp
                do dj = -1, 1
                   do di = -1, 1
                      if (di == 0 .and. dj == 0) cycle
                      v_xmin = x_face(mesh, i + di); v_xmax = x_face(mesh, i + di + 1)
                      v_ymin = y_face(mesh, j + dj); v_ymax = y_face(mesh, j + dj + 1)
                      v_target = C(i + di, j + dj) * dx * dy
                      v_recon  = cut_box_area(v_xmin, v_xmax, v_ymin, v_ymax, cos(d_pt), sin(d_pt), &
                                             compute_alpha(xmin, xmax, ymin, ymax, cos(d_pt), sin(d_pt), c_val))
                      diff = v_recon - v_target
                      fd = fd + diff * diff
                   end do
                end do
             end if
          end do

          best_theta = 0.5_dp * (a_pt + b_pt)
          best_nx = cos(best_theta)
          best_ny = sin(best_theta)
          best_alpha = compute_alpha(xmin, xmax, ymin, ymax, best_nx, best_ny, c_val)

          segments(i, j)%nx = best_nx
          segments(i, j)%ny = best_ny
          segments(i, j)%alpha = best_alpha
       end do
    end do
  end subroutine reconstruct_lvira

  ! [Procédure/Module] subroutine reconstruct_interface
  subroutine reconstruct_interface(method_name, mesh, C, segments)
    character(len=*), intent(in) :: method_name
    type(mesh_2d_t), intent(in) :: mesh
    real(dp), intent(in) :: C(0:mesh%nx_local+1, 0:mesh%ny_local+1)
    type(interface_segment_t), intent(out) :: segments(0:mesh%nx_local+1, 0:mesh%ny_local+1)

    integer :: i, j, nx, ny, i_left_src, i_right_src, j_bot_src, j_top_src
    real(dp) :: c_val

    select case (trim(method_name))
    case ("SLIC")
       call reconstruct_slic(mesh, C, segments)
    case ("LVIRA")
       call reconstruct_lvira(mesh, C, segments)
    case ("ELVIRA")
       call reconstruct_elvira(mesh, C, segments)
    case default
       call reconstruct_plic(mesh, C, segments)
    end select

    nx = mesh%nx_local
    ny = mesh%ny_local
    if (.not. mesh%is_boundary_left .and. mesh%size <= 1) then
       i_left_src = nx
    else
       i_left_src = 1
    end if
    if (.not. mesh%is_boundary_right .and. mesh%size <= 1) then
       i_right_src = 1
    else
       i_right_src = nx
    end if
    if (.not. mesh%is_boundary_bottom .and. mesh%size <= 1) then
       j_bot_src = ny
    else
       j_bot_src = 1
    end if
    if (.not. mesh%is_boundary_top .and. mesh%size <= 1) then
       j_top_src = 1
    else
       j_top_src = ny
    end if

    do j = 1, ny
       c_val = C(0, j)
       if (c_val > 1.0e-6_dp .and. c_val < 1.0_dp - 1.0e-6_dp) then
          segments(0, j)%nx = segments(i_left_src, j)%nx
          segments(0, j)%ny = segments(i_left_src, j)%ny
          segments(0, j)%alpha = compute_alpha(x_face(mesh, 0), x_face(mesh, 1), &
                                               y_face(mesh, j), y_face(mesh, j + 1), &
                                               segments(0, j)%nx, segments(0, j)%ny, c_val)
       end if
       c_val = C(nx + 1, j)
       if (c_val > 1.0e-6_dp .and. c_val < 1.0_dp - 1.0e-6_dp) then
          segments(nx + 1, j)%nx = segments(i_right_src, j)%nx
          segments(nx + 1, j)%ny = segments(i_right_src, j)%ny
          segments(nx + 1, j)%alpha = compute_alpha(x_face(mesh, nx + 1), x_face(mesh, nx + 2), &
                                                    y_face(mesh, j), y_face(mesh, j + 1), &
                                                    segments(nx + 1, j)%nx, segments(nx + 1, j)%ny, c_val)
       end if
    end do

    do i = 0, nx + 1
       c_val = C(i, 0)
       if (c_val > 1.0e-6_dp .and. c_val < 1.0_dp - 1.0e-6_dp) then
          segments(i, 0)%nx = segments(i, j_bot_src)%nx
          segments(i, 0)%ny = segments(i, j_bot_src)%ny
          segments(i, 0)%alpha = compute_alpha(x_face(mesh, i), x_face(mesh, i + 1), &
                                               y_face(mesh, 0), y_face(mesh, 1), &
                                               segments(i, 0)%nx, segments(i, 0)%ny, c_val)
       end if
       c_val = C(i, ny + 1)
       if (c_val > 1.0e-6_dp .and. c_val < 1.0_dp - 1.0e-6_dp) then
          segments(i, ny + 1)%nx = segments(i, j_top_src)%nx
          segments(i, ny + 1)%ny = segments(i, j_top_src)%ny
          segments(i, ny + 1)%alpha = compute_alpha(x_face(mesh, i), x_face(mesh, i + 1), &
                                                    y_face(mesh, ny + 1), y_face(mesh, ny + 2), &
                                                    segments(i, ny + 1)%nx, segments(i, ny + 1)%ny, c_val)
       end if
    end do
  end subroutine reconstruct_interface

  ! [Procédure/Module] function evaluate_limiter
  function evaluate_limiter(limiter_name, r, c_donor, c_upwind, c_acceptor) result(phi)
    character(len=*), intent(in) :: limiter_name
    real(dp), intent(in) :: r, c_donor, c_upwind, c_acceptor
    real(dp) :: phi
    real(dp) :: denom, c_tilde_D, c_tilde_f

    select case (trim(limiter_name))
    case ("VAN_LEER")
       if (r > 0.0_dp) then
          phi = 2.0_dp * r / (1.0_dp + r)
       else
          phi = 0.0_dp
       end if
    case ("MUSCL")
       phi = max(0.0_dp, min(2.0_dp * r, min(0.5_dp * (1.0_dp + r), 2.0_dp)))
    case ("HRIC")
       denom = c_acceptor - c_upwind
       if (abs(denom) < 1.0e-12_dp) then
          phi = 0.0_dp
          return
       end if
       c_tilde_D = (c_donor - c_upwind) / denom
       c_tilde_f = c_tilde_D
       if (c_tilde_D >= 0.0_dp .and. c_tilde_D <= 1.0_dp) then
          if (c_tilde_D < 0.5_dp) then
             c_tilde_f = 2.0_dp * c_tilde_D
          else
             c_tilde_f = 1.0_dp
          end if
       end if
       if (abs(1.0_dp - c_tilde_D) > 1.0e-12_dp) then
          phi = 2.0_dp * (c_tilde_f - c_tilde_D) / (1.0_dp - c_tilde_D)
          phi = max(0.0_dp, min(2.0_dp, phi))
       else
          phi = 1.0_dp
       end if
    case default ! "SUPERBEE"
       phi = max(0.0_dp, max(min(2.0_dp * r, 1.0_dp), min(r, 2.0_dp)))
    end select
  end function evaluate_limiter

end module module_interface_reconstruction

\end{lstlisting}

\chapter{Advection Géométrique et Algébrique VOF en Fortran 2018 (module_interface_vof.f90)}
\section{Architecture des 1034 Lignes de \texttt{module\_interface\_vof.f90}}
Ce module implémente la séparation directionnelle d'ordre 2 de Strang (\texttt{advect\_vof\_interface}) et les trois grandes familles de schémas de transport interfacial :
\begin{enumerate}
  \item \textbf{Schéma géométrique de Weymouth \& Yue (2010)} (\texttt{advect\_weymouth\_yue\_x}, \texttt{advect\_weymouth\_yue\_y}) : utilise la fonction indicatrice binaire $C_{\text{dil}}(i,j) = \mathbb{I}(C_{i,j}^n \ge 0.5)$ partagée entre les deux balayages directionnels pour garantir une conservation de la masse exacte à la précision machine.
  \item \textbf{Schéma géométrique CIAM de Li (1995)} (\texttt{advect\_ciam\_x}, \texttt{advect\_ciam\_y}) : transporte lagrangiennement les segments PLIC en tenant compte de l'étirement/compression affine de la bande donneuse par le gradient de vitesse local $\partial u / \partial x$.
  \item \textbf{Schémas algébriques TVD haute résolution} (\texttt{advect\_tvd\_x}, \texttt{advect\_tvd\_y}) : implémentent les limiteurs de pente non linéaires \textbf{Superbee}, \textbf{Van Leer}, \textbf{MUSCL / Monotonized Central (MC)}, \textbf{Minmod} et \textbf{HRIC}.
\end{enumerate}

\section{Répertoire des Types Dérivés, Sous-Routines et Fonctions de \texttt{src/module\_interface\_vof.f90}}
\begin{center}\small
\begin{tabular}{clll}
\toprule
\textbf{Ligne} & \textbf{Nature} & \textbf{Identifiant Fortran 2018} & \textbf{Interface / Arguments} \\
\midrule
2 & module & \texttt{module\_interface\_vof} & Unité modulaire Fortran 2018 \\
23 & type dérivé & \texttt{interface\_vof\_t} & Structure de données Fortran 2018 \\
38 & subroutine & \texttt{set\_initial\_conditions} & Arguments : (fields, mesh, cfg) \\
235 & subroutine & \texttt{update\_physical\_properties} & Arguments : (fields, mesh, cfg) \\
372 & subroutine & \texttt{compute\_curvature\_and\_csf} & Arguments : (fields, mesh, cfg) \\
495 & subroutine & \texttt{advect\_vof} & Arguments : (fields, mesh, cfg, dt) \\
519 & subroutine & \texttt{advect\_vof\_tvd} & Arguments : (fields, mesh, cfg, dt) \\
629 & subroutine & \texttt{advect\_vof\_ciam} & Arguments : (fields, mesh, cfg, dt) \\
807 & subroutine & \texttt{advect\_vof\_weymouth\_yue} & Arguments : (fields, mesh, cfg, dt) \\
979 & subroutine & \texttt{vof\_init} & Arguments : (this, cfg) \\
989 & subroutine & \texttt{vof\_set\_ic} & Arguments : (this, mesh, fields) \\
1005 & subroutine & \texttt{vof\_advect\_step} & Arguments : (this, mesh, fields, dt) \\
\bottomrule
\end{tabular}
\end{center}

\section{Code Source Intégral Annoté : \texttt{src/module\_interface\_vof.f90}}
Le listing ci-dessous présente l'intégralité des 1022 lignes du fichier \texttt{src/module\_interface\_vof.f90} de l'implémentation Fortran 2018 :

\begin{lstlisting}[language={[95]Fortran}]
! [Procédure/Module] module module_interface_vof
module module_interface_vof
  use module_precision, only: dp
  use module_constants
  use module_mesh, only: mesh_2d_t, x_center, y_center, x_face, y_face, is_in_obstacle
  use module_fields, only: flow_fields_t
  use module_config, only: config_t
  use module_mpi_comm, only: exchange_halos_scalar, apply_vof_bc
  use module_interface_reconstruction, only: interface_segment_t, cut_box_area, compute_alpha, &
                                            reconstruct_interface, evaluate_limiter
  implicit none
  private

  public :: interface_vof_t
  public :: set_initial_conditions
  public :: advect_vof
  public :: update_physical_properties
  public :: compute_curvature_and_csf

  ! Compteur de pas pour Strang splitting alterné
  integer, save :: vof_step_count = 0

  type :: interface_vof_t
     ! Déclaration de variable locale / paramètre de calcul
     type(config_t) :: cfg
  contains
     ! Affectation et calcul numérique
     procedure :: init => vof_init
     ! Affectation et calcul numérique
     procedure :: set_initial_condition => vof_set_ic
     ! Affectation et calcul numérique
     procedure :: advect => vof_advect_step
  end type interface_vof_t

contains

  ! [Procédure/Module] subroutine set_initial_conditions
  subroutine set_initial_conditions(fields, mesh, cfg)
    ! Déclaration de variable locale / paramètre de calcul
    type(flow_fields_t), intent(inout) :: fields
    ! Déclaration de variable locale / paramètre de calcul
    type(mesh_2d_t), intent(in) :: mesh
    ! Déclaration de variable locale / paramètre de calcul
    type(config_t), intent(in) :: cfg

    ! Déclaration de variable locale / paramètre de calcul
    integer :: i, j, nx, ny
    ! Déclaration de variable locale / paramètre de calcul
    real(dp) :: x, y, dist, eps, int_y, x_c_jet, half_w_jet, stub_h_jet

    ! Affectation et calcul numérique
    nx = mesh%nx_local
    ! Affectation et calcul numérique
    ny = mesh%ny_local
    ! Affectation et calcul numérique
    eps = 1.5_dp * max(mesh%dx, mesh%dy)

    ! 1. Initialisation des vitesses (par défaut 0 ou imposé)
    fields%u = cfg%ic_u
    ! Affectation et calcul numérique
    fields%v = cfg%ic_v

    ! 2. Initialisation de la fraction volumique C
    do j = 0, ny + 1
       ! Affectation et calcul numérique
       y = y_center(mesh, j)
       ! Boucle spatiale selon l'axe horizontal X (accès contigu colonne-major)
       do i = 0, nx + 1
          ! Affectation et calcul numérique
          x = x_center(mesh, i)

          select case (cfg%ic_type)
          case (IC_UNIFORM)
             ! Transport / actualisation de la fraction volumique VOF
             fields%C(i, j) = 1.0_dp

          case (IC_DROPLET)
             ! Affectation et calcul numérique
             dist = sqrt((x - cfg%ic_center_x)**2 + (y - cfg%ic_center_y)**2) - cfg%ic_radius
             ! Évaluation de la condition logique
             if (dist < -eps) then
                ! Transport / actualisation de la fraction volumique VOF
                fields%C(i, j) = 1.0_dp
             ! Branche conditionnelle alternative
             else if (dist > eps) then
                ! Transport / actualisation de la fraction volumique VOF
                fields%C(i, j) = 0.0_dp
             ! Branche par défaut si condition non satisfaite
             else
                ! Transport / actualisation de la fraction volumique VOF
                fields%C(i, j) = 0.5_dp - dist / (2.0_dp * eps) - (1.0_dp / (2.0_dp * PI)) * sin(PI * dist / eps)
             end if

          case (IC_BUBBLE)
             ! Affectation et calcul numérique
             dist = sqrt((x - cfg%ic_center_x)**2 + (y - cfg%ic_center_y)**2) - cfg%ic_radius
             ! Évaluation de la condition logique
             if (dist < -eps) then
                ! Transport / actualisation de la fraction volumique VOF
                fields%C(i, j) = 0.0_dp ! Intérieur bulle (fluide 2)
             ! Branche conditionnelle alternative
             else if (dist > eps) then
                ! Transport / actualisation de la fraction volumique VOF
                fields%C(i, j) = 1.0_dp ! Extérieur (fluide 1)
             ! Branche par défaut si condition non satisfaite
             else
                ! Transport / actualisation de la fraction volumique VOF
                fields%C(i, j) = 0.5_dp + dist / (2.0_dp * eps) + (1.0_dp / (2.0_dp * PI)) * sin(PI * dist / eps)
             end if

          case (IC_DAM_BREAK)
             ! Affectation et calcul numérique
             if (x >= cfg%ic_box_x0 .and. x <= cfg%ic_box_x1 .and. &
                 ! Affectation et calcul numérique
                 y >= cfg%ic_box_y0 .and. y <= cfg%ic_box_y1) then
                ! Transport / actualisation de la fraction volumique VOF
                fields%C(i, j) = 1.0_dp
             ! Branche par défaut si condition non satisfaite
             else
                ! Transport / actualisation de la fraction volumique VOF
                fields%C(i, j) = 0.0_dp
             end if

          case (IC_LIQUID_SHEET)
             ! Nappe liquide cisaillée (Cas 3.2 de l'article Fluids 2021)
             if (cfg%ic_center_y > 0.0_dp) then
                ! Calcul de distance à l'axe de la nappe
                dist = abs(y - cfg%ic_center_y) - 0.5_dp * cfg%ic_thickness
             else
                ! Calcul de distance au plan médian
                dist = abs(y - 0.5_dp * mesh%ly) - 0.5_dp * cfg%ic_thickness
             end if
             ! Évaluation de l'indicatrice C
             if (dist < -eps) then
                fields%C(i, j) = 0.0_dp ! Intérieur nappe liquide (fluide 2)
             else if (dist > eps) then
                fields%C(i, j) = 1.0_dp ! Gaz ambiant (fluide 1)
             else
                fields%C(i, j) = 0.5_dp + dist / (2.0_dp * eps) + &
                     (1.0_dp / (2.0_dp * PI)) * sin(PI * dist / eps)
             end if

          case (IC_VISCOUS_JET)
             ! Jet liquide visqueux injecté (Cas 3.3 de l'article Fluids 2021)
             if (cfg%ic_center_x > 0.0_dp) then
                x_c_jet = cfg%ic_center_x
             else
                x_c_jet = 0.5_dp * mesh%lx
             end if
             if (cfg%ic_jet_width > 0.0_dp) then
                half_w_jet = 0.5_dp * cfg%ic_jet_width
             else
                half_w_jet = 0.04_dp
             end if
             stub_h_jet = 2.0_dp * mesh%dy
             if (y >= mesh%ly - stub_h_jet .and. abs(x - x_c_jet) <= half_w_jet) then
                fields%C(i, j) = 0.0_dp ! Amorce liquide (fluide 2)
             else
                fields%C(i, j) = 1.0_dp ! Cavité de gaz (fluide 1)
             end if

          case (IC_VORTEX_SHEAR)
             ! Affectation et calcul numérique
             dist = sqrt((x - cfg%ic_center_x)**2 + (y - cfg%ic_center_y)**2) - cfg%ic_radius
             ! Évaluation de la condition logique
             if (dist < -eps) then
                ! Transport / actualisation de la fraction volumique VOF
                fields%C(i, j) = 1.0_dp
             ! Branche conditionnelle alternative
             else if (dist > eps) then
                ! Transport / actualisation de la fraction volumique VOF
                fields%C(i, j) = 0.0_dp
             ! Branche par défaut si condition non satisfaite
             else
                ! Transport / actualisation de la fraction volumique VOF
                fields%C(i, j) = 0.5_dp - dist / (2.0_dp * eps) - (1.0_dp / (2.0_dp * PI)) * sin(PI * dist / eps)
             end if

          case (IC_RAYLEIGH_TAYLOR)
             ! Affectation et calcul numérique
             int_y = cfg%ic_center_y + cfg%ic_pert_amp * cos(2.0_dp * PI * cfg%ic_pert_k * x / mesh%lx)
             ! Évaluation de la condition logique
             if (y >= int_y) then
                ! Transport / actualisation de la fraction volumique VOF
                fields%C(i, j) = 1.0_dp ! Fluide lourd au-dessus
             ! Branche par défaut si condition non satisfaite
             else
                ! Transport / actualisation de la fraction volumique VOF
                fields%C(i, j) = 0.0_dp ! Fluide léger en dessous
             end if

          case (IC_LAPLACIAN_TEST)
             ! Transport / actualisation de la fraction volumique VOF
             fields%C(i, j) = 1.0_dp

          case default
             ! Transport / actualisation de la fraction volumique VOF
             fields%C(i, j) = 1.0_dp
          end select
       end do
    end do

    ! 3. Détection des obstacles solides pour la pénalisation de Brinkman
    fields%solid_mask = 0
    ! Évaluation de la condition logique
    if (cfg%obs_enabled) then
       ! Boucle spatiale sur les cellules locales selon l'axe vertical Y (ordonnées)
       do j = 0, ny + 1
          ! Affectation et calcul numérique
          y = y_center(mesh, j)
          ! Boucle spatiale selon l'axe horizontal X (accès contigu colonne-major)
          do i = 0, nx + 1
             ! Affectation et calcul numérique
             x = x_center(mesh, i)
             if (is_in_obstacle(x, y, cfg%obs_type, cfg%obs_x0, cfg%obs_y0, cfg%obs_x1, cfg%obs_y1, &
                                cfg%obs_xc, cfg%obs_yc, cfg%obs_radius)) then
                ! Affectation et calcul numérique
                fields%solid_mask(i, j) = 1
             end if
          end do
       end do
    end if

    ! Appel de sous-programme modulaire
    call exchange_halos_scalar(fields%C, mesh)
    ! Appel de sous-programme modulaire
    call apply_vof_bc(fields, mesh, cfg)
    ! Appel de sous-programme modulaire
    call update_physical_properties(fields, mesh, cfg)
    ! Appel de sous-programme modulaire
    call compute_curvature_and_csf(fields, mesh, cfg)
  end subroutine set_initial_conditions

  ! [Procédure/Module] subroutine update_physical_properties
  subroutine update_physical_properties(fields, mesh, cfg)
    ! Déclaration de variable locale / paramètre de calcul
    type(flow_fields_t), intent(inout) :: fields
    ! Déclaration de variable locale / paramètre de calcul
    type(mesh_2d_t), intent(in) :: mesh
    ! Déclaration de variable locale / paramètre de calcul
    type(config_t), intent(in) :: cfg

    ! Déclaration de variable locale / paramètre de calcul
    integer :: i, j, nx, ny
    ! Déclaration de variable locale / paramètre de calcul
    real(dp) :: c, y, z, mu_val, m_sw, m_se, m_nw, m_ne

    ! Affectation et calcul numérique
    nx = mesh%nx_local
    ! Affectation et calcul numérique
    ny = mesh%ny_local

    ! Évaluation de la condition logique
    if (.not. cfg%enable_two_phase) then
       ! Calcul de la masse volumique locale ou interpolée aux faces
       fields%rho = cfg%rho1
       ! Calcul de la masse volumique locale ou interpolée aux faces
       fields%rho_u = cfg%rho1
       ! Calcul de la masse volumique locale ou interpolée aux faces
       fields%rho_v = cfg%rho1
       ! Évaluation de la condition logique
       if (cfg%viscosity_model == VISC_EXPONENTIAL) then
          ! Boucle spatiale sur les cellules locales selon l'axe vertical Y (ordonnées)
          do j = 0, ny + 1
             ! Affectation et calcul numérique
             y = y_center(mesh, j)
             ! Affectation et calcul numérique
             z = y - cfg%visc_e
             ! Calcul de la viscosité dynamique locale (harmonique aux nœuds)
             mu_val = cfg%visc_mu0 * exp(cfg%visc_a * (z * z) / (cfg%visc_e * cfg%visc_e))
             ! Mise à jour de la composante de vitesse discrète sur face décalée MAC
             fields%mu(:, j) = mu_val
          end do
          ! Boucle spatiale sur les cellules locales selon l'axe vertical Y (ordonnées)
          do j = 0, ny + 1
             ! Affectation et calcul numérique
             y = y_face(mesh, j)
             ! Affectation et calcul numérique
             z = y - cfg%visc_e
             ! Calcul de la viscosité dynamique locale (harmonique aux nœuds)
             mu_val = cfg%visc_mu0 * exp(cfg%visc_a * (z * z) / (cfg%visc_e * cfg%visc_e))
             ! Calcul de la viscosité dynamique locale (harmonique aux nœuds)
             fields%mu_node(:, j) = mu_val
          end do
       ! Branche par défaut si condition non satisfaite
       else
          ! Calcul de la viscosité dynamique locale (harmonique aux nœuds)
          fields%mu = cfg%mu1
          ! Calcul de la viscosité dynamique locale (harmonique aux nœuds)
          fields%mu_node = cfg%mu1
       end if
       ! Retour à la procédure appelante
       return
    end if

    ! Coins de C
    fields%C(0, 0)       = fields%C(1, 1)
    ! Transport / actualisation de la fraction volumique VOF
    fields%C(nx+1, 0)    = fields%C(nx, 1)
    ! Transport / actualisation de la fraction volumique VOF
    fields%C(0, ny+1)    = fields%C(1, ny)
    ! Transport / actualisation de la fraction volumique VOF
    fields%C(nx+1, ny+1) = fields%C(nx, ny)

    ! 1. Masse volumique arithmétique et viscosité au centre
    do j = 0, ny + 1
       ! Boucle spatiale selon l'axe horizontal X (accès contigu colonne-major)
       do i = 0, nx + 1
          ! Affectation et calcul numérique
          c = max(0.0_dp, min(1.0_dp, fields%C(i, j)))
          ! Calcul de la masse volumique locale ou interpolée aux faces
          fields%rho(i, j) = c * cfg%rho1 + (1.0_dp - c) * cfg%rho2
          ! Mise à jour de la composante de vitesse discrète sur face décalée MAC
          fields%mu(i, j)  = c * cfg%mu1  + (1.0_dp - c) * cfg%mu2
       end do
    end do

    ! Mise à jour de la composante de vitesse discrète sur face décalée MAC
    fields%mu(0, 0)       = fields%mu(1, 1)
    ! Mise à jour de la composante de vitesse discrète sur face décalée MAC
    fields%mu(nx+1, 0)    = fields%mu(nx, 1)
    ! Mise à jour de la composante de vitesse discrète sur face décalée MAC
    fields%mu(0, ny+1)    = fields%mu(1, ny)
    ! Mise à jour de la composante de vitesse discrète sur face décalée MAC
    fields%mu(nx+1, ny+1) = fields%mu(nx, ny)

    ! 2. Masse volumique sur les faces u
    do j = 1, ny
       ! Boucle spatiale selon l'axe horizontal X (accès contigu colonne-major)
       do i = 1, nx + 1
          ! Mise à jour de la composante de vitesse discrète sur face décalée MAC
          fields%rho_u(i, j) = 0.5_dp * (fields%rho(i - 1, j) + fields%rho(i, j))
       end do
    end do

    ! 3. Masse volumique sur les faces v
    do j = 1, ny + 1
       ! Boucle spatiale selon l'axe horizontal X (accès contigu colonne-major)
       do i = 1, nx
          ! Mise à jour de la composante de vitesse discrète sur face décalée MAC
          fields%rho_v(i, j) = 0.5_dp * (fields%rho(i, j - 1) + fields%rho(i, j))
       end do
    end do

    ! 4. Viscosité harmonique aux nœuds (coins de mailles)
    do j = 1, ny + 1
       ! Boucle spatiale selon l'axe horizontal X (accès contigu colonne-major)
       do i = 1, nx + 1
          ! Affectation et calcul numérique
          m_sw = fields%mu(i-1, j-1)
          ! Affectation et calcul numérique
          m_se = fields%mu(i, j-1)
          ! Affectation et calcul numérique
          m_nw = fields%mu(i-1, j)
          ! Affectation et calcul numérique
          m_ne = fields%mu(i, j)
          ! Calcul de la viscosité dynamique locale (harmonique aux nœuds)
          fields%mu_node(i, j) = 4.0_dp / (1.0_dp / m_sw + 1.0_dp / m_se + 1.0_dp / m_nw + 1.0_dp / m_ne)
       end do
    end do
  end subroutine update_physical_properties

  ! [Procédure/Module] pure real
  pure real(dp) function superbee(r)
    ! Déclaration de variable locale / paramètre de calcul
    real(dp), intent(in) :: r
    ! Affectation et calcul numérique
    superbee = max(0.0_dp, max(min(2.0_dp * r, 1.0_dp), min(r, 2.0_dp)))
  end function superbee

  ! [Procédure/Module] subroutine compute_curvature_and_csf
  subroutine compute_curvature_and_csf(fields, mesh, cfg)
    ! Déclaration de variable locale / paramètre de calcul
    type(flow_fields_t), intent(inout) :: fields
    ! Déclaration de variable locale / paramètre de calcul
    type(mesh_2d_t), intent(in) :: mesh
    ! Déclaration de variable locale / paramètre de calcul
    type(config_t), intent(in) :: cfg

    ! Déclaration de variable locale / paramètre de calcul
    integer :: i, j, nx, ny
    ! Déclaration de variable locale / paramètre de calcul
    real(dp) :: dx, dy, dc_dx, dc_dy, grad_mag, kap_avg, grad_c
    ! Déclaration de variable locale / paramètre de calcul
    real(dp), allocatable :: norm_x(:, :), norm_y(:, :)
    ! Déclaration de variable locale / paramètre de calcul
    real(dp), allocatable :: kappa(:, :)

    ! Évaluation de la condition logique
    if (.not. cfg%enable_surface_tension .or. cfg%sigma <= 0.0_dp) then
       ! Affectation et calcul numérique
       fields%fx_sigma = 0.0_dp
       ! Affectation et calcul numérique
       fields%fy_sigma = 0.0_dp
       ! Affectation et calcul numérique
       fields%curvature = 0.0_dp
       ! Retour à la procédure appelante
       return
    end if

    ! Affectation et calcul numérique
    nx = mesh%nx_local
    ! Affectation et calcul numérique
    ny = mesh%ny_local
    ! Calcul de la métrique spatiale élémentaire
    dx = mesh%dx
    ! Calcul de la métrique spatiale élémentaire
    dy = mesh%dy

    ! Allocation dynamique de tableau contigu en mémoire
    allocate(norm_x(0:nx+1, 0:ny+1), norm_y(0:nx+1, 0:ny+1))
    ! Allocation dynamique de tableau contigu en mémoire
    allocate(kappa(0:nx+1, 0:ny+1))
    ! Évaluation de la norme du résidu ou de l'erreur numérique
    norm_x = 0.0_dp
    ! Évaluation de la norme du résidu ou de l'erreur numérique
    norm_y = 0.0_dp
    ! Affectation et calcul numérique
    kappa = 0.0_dp

    ! 1. Gradient et normale unitaire aux centres
    do j = 1, ny
       ! Boucle spatiale selon l'axe horizontal X (accès contigu colonne-major)
       do i = 1, nx
          ! Calcul de la métrique spatiale élémentaire
          dc_dx = (fields%C(i+1, j) - fields%C(i-1, j)) / (2.0_dp * dx)
          ! Calcul de la métrique spatiale élémentaire
          dc_dy = (fields%C(i, j+1) - fields%C(i, j-1)) / (2.0_dp * dy)
          ! Affectation et calcul numérique
          grad_mag = sqrt(dc_dx * dc_dx + dc_dy * dc_dy)
          ! Évaluation de la condition logique
          if (grad_mag > 1.0e-3_dp / dx) then
             ! Évaluation de la norme du résidu ou de l'erreur numérique
             norm_x(i, j) = dc_dx / grad_mag
             ! Évaluation de la norme du résidu ou de l'erreur numérique
             norm_y(i, j) = dc_dy / grad_mag
          end if
       end do
    end do

    ! Appel de sous-programme modulaire
    call exchange_halos_scalar(norm_x, mesh)
    ! Appel de sous-programme modulaire
    call exchange_halos_scalar(norm_y, mesh)

    ! 2. Courbure kappa = -div(n)
    do j = 1, ny
       ! Boucle spatiale selon l'axe horizontal X (accès contigu colonne-major)
       do i = 1, nx
          ! Évaluation de la condition logique
          if (norm_x(i, j) /= 0.0_dp .or. norm_y(i, j) /= 0.0_dp) then
             ! Affectation et calcul numérique
             kappa(i, j) = - ((norm_x(i+1, j) - norm_x(i-1, j)) / (2.0_dp * dx) + &
                              (norm_y(i, j+1) - norm_y(i, j-1)) / (2.0_dp * dy))
          end if
       end do
    end do

    ! Appel de sous-programme modulaire
    call exchange_halos_scalar(kappa, mesh)
    ! Affectation et calcul numérique
    fields%curvature = kappa

    ! 3. Forces de tension superficielle CSF (Balanced-Force)
    do j = 1, ny
       ! Boucle spatiale selon l'axe horizontal X (accès contigu colonne-major)
       do i = 1, nx + 1
          ! Affectation et calcul numérique
          kap_avg = 0.5_dp * (kappa(i-1, j) + kappa(i, j))
          ! Affectation et calcul numérique
          grad_c  = (fields%C(i, j) - fields%C(i-1, j)) / dx
          ! Affectation et calcul numérique
          fields%fx_sigma(i, j) = cfg%sigma * kap_avg * grad_c
       end do
    end do

    ! Boucle spatiale sur les cellules locales selon l'axe vertical Y (ordonnées)
    do j = 1, ny + 1
       ! Boucle spatiale selon l'axe horizontal X (accès contigu colonne-major)
       do i = 1, nx
          ! Affectation et calcul numérique
          kap_avg = 0.5_dp * (kappa(i, j-1) + kappa(i, j))
          ! Affectation et calcul numérique
          grad_c  = (fields%C(i, j) - fields%C(i, j-1)) / dy
          ! Affectation et calcul numérique
          fields%fy_sigma(i, j) = cfg%sigma * kap_avg * grad_c
       end do
    end do

    ! Libération de la mémoire allouée
    deallocate(norm_x, norm_y, kappa)
  end subroutine compute_curvature_and_csf

  ! [Procédure/Module] subroutine advect_vof
  subroutine advect_vof(fields, mesh, cfg, dt)
    type(flow_fields_t), intent(inout) :: fields
    type(mesh_2d_t), intent(in) :: mesh
    type(config_t), intent(in) :: cfg
    real(dp), intent(in) :: dt

    select case (trim(cfg%advection_method))
    case ("CIAM")
       call advect_vof_ciam(fields, mesh, cfg, dt)
    case ("WEYMOUTH_YUE")
       call advect_vof_weymouth_yue(fields, mesh, cfg, dt)
    case default
       call advect_vof_tvd(fields, mesh, cfg, dt)
    end select

    vof_step_count = vof_step_count + 1

    call exchange_halos_scalar(fields%C, mesh)
    call apply_vof_bc(fields, mesh, cfg)
    call update_physical_properties(fields, mesh, cfg)
    call compute_curvature_and_csf(fields, mesh, cfg)
  end subroutine advect_vof

  ! [Procédure/Module] subroutine advect_vof_tvd
  subroutine advect_vof_tvd(fields, mesh, cfg, dt)
    type(flow_fields_t), intent(inout) :: fields
    type(mesh_2d_t), intent(in) :: mesh
    type(config_t), intent(in) :: cfg
    real(dp), intent(in) :: dt

    integer :: i, j, nx, ny, i_up, i_down, j_up, j_down
    real(dp) :: dx, dy, u_face, v_face, c_l, c_r, c_b, c_t, c_u
    real(dp) :: dc_up, dc_loc, r, phi, fx_in, fx_out, fy_in, fy_out
    real(dp), allocatable :: flux_x(:, :), flux_y(:, :), c_new(:, :)

    nx = mesh%nx_local
    ny = mesh%ny_local
    dx = mesh%dx
    dy = mesh%dy

    allocate(flux_x(0:nx+1, 0:ny+1), flux_y(0:nx+1, 0:ny+1))
    allocate(c_new(0:nx+1, 0:ny+1))
    flux_x = 0.0_dp
    flux_y = 0.0_dp
    c_new = fields%C

    ! Flux selon X
    do j = 1, ny
       do i = 1, nx + 1
          u_face = fields%u(i, j)
          c_l = fields%C(i-1, j)
          c_r = fields%C(i, j)

          if (u_face >= 0.0_dp) then
             i_up = max(0, i - 2)
             c_u = fields%C(i_up, j)
             dc_up = c_l - c_u
             dc_loc = c_r - c_l
             if (abs(dc_loc) > 1.0e-12_dp) then
                r = dc_up / dc_loc
             else
                r = 0.0_dp
             end if
             phi = evaluate_limiter(cfg%tvd_limiter, r, c_l, c_u, c_r)
             flux_x(i, j) = u_face * (c_l + 0.5_dp * phi * (1.0_dp - min(1.0_dp, abs(u_face) * dt / dx)) * dc_loc)
          else
             i_down = min(nx + 1, i + 1)
             c_u = fields%C(i_down, j)
             dc_up = c_u - c_r
             dc_loc = c_r - c_l
             if (abs(dc_loc) > 1.0e-12_dp) then
                r = dc_up / dc_loc
             else
                r = 0.0_dp
             end if
             phi = evaluate_limiter(cfg%tvd_limiter, r, c_r, c_u, c_l)
             flux_x(i, j) = u_face * (c_r - 0.5_dp * phi * (1.0_dp - min(1.0_dp, abs(u_face) * dt / dx)) * dc_loc)
          end if
       end do
    end do

    ! Flux selon Y
    do j = 1, ny + 1
       do i = 1, nx
          v_face = fields%v(i, j)
          c_b = fields%C(i, j-1)
          c_t = fields%C(i, j)

          if (v_face >= 0.0_dp) then
             j_up = max(0, j - 2)
             c_u = fields%C(i, j_up)
             dc_up = c_b - c_u
             dc_loc = c_t - c_b
             if (abs(dc_loc) > 1.0e-12_dp) then
                r = dc_up / dc_loc
             else
                r = 0.0_dp
             end if
             phi = evaluate_limiter(cfg%tvd_limiter, r, c_b, c_u, c_t)
             flux_y(i, j) = v_face * (c_b + 0.5_dp * phi * (1.0_dp - min(1.0_dp, abs(v_face) * dt / dy)) * dc_loc)
          else
             j_down = min(ny + 1, j + 1)
             c_u = fields%C(i, j_down)
             dc_up = c_u - c_t
             dc_loc = c_t - c_b
             if (abs(dc_loc) > 1.0e-12_dp) then
                r = dc_up / dc_loc
             else
                r = 0.0_dp
             end if
             phi = evaluate_limiter(cfg%tvd_limiter, r, c_t, c_u, c_b)
             flux_y(i, j) = v_face * (c_t - 0.5_dp * phi * (1.0_dp - min(1.0_dp, abs(v_face) * dt / dy)) * dc_loc)
          end if
       end do
    end do

    ! Mise à jour conservative
    do j = 1, ny
       do i = 1, nx
          fx_in  = flux_x(i, j)
          fx_out = flux_x(i+1, j)
          fy_in  = flux_y(i, j)
          fy_out = flux_y(i, j+1)

          c_new(i, j) = fields%C(i, j) - (dt / dx) * (fx_out - fx_in) - (dt / dy) * (fy_out - fy_in)
          c_new(i, j) = max(0.0_dp, min(1.0_dp, c_new(i, j)))
       end do
    end do

    fields%C = c_new
    deallocate(flux_x, flux_y, c_new)
  end subroutine advect_vof_tvd

  ! [Procédure/Module] subroutine advect_vof_ciam
  subroutine advect_vof_ciam(fields, mesh, cfg, dt)
    type(flow_fields_t), intent(inout) :: fields
    type(mesh_2d_t), intent(in) :: mesh
    type(config_t), intent(in) :: cfg
    real(dp), intent(in) :: dt

    integer :: i, j, nx, ny, donor_idx
    real(dp) :: dx, dy, delta_x, delta_y, u_face, v_face
    real(dp) :: x_pos, y_pos, b_xmin, b_xmax, b_ymin, b_ymax, denom
    type(interface_segment_t), allocatable :: segments(:, :)
    real(dp), allocatable :: flux_vol(:, :), c_star(:, :), c_final(:, :)
    logical :: x_first

    nx = mesh%nx_local
    ny = mesh%ny_local
    dx = mesh%dx
    dy = mesh%dy

    allocate(segments(0:nx+1, 0:ny+1))
    allocate(flux_vol(0:nx+1, 0:ny+1))
    allocate(c_star(0:nx+1, 0:ny+1))
    allocate(c_final(0:nx+1, 0:ny+1))

    x_first = (mod(vof_step_count, 2) == 0)

    if (x_first) then
       ! 1. Passe X
       call reconstruct_interface(cfg%reconstruction_method, mesh, fields%C, segments)
       flux_vol = 0.0_dp
       do j = 1, ny
          do i = 1, nx + 1
             u_face = fields%u(i, j)
             delta_x = abs(u_face) * dt
             if (delta_x <= 1.0e-14_dp) cycle
             x_pos = x_face(mesh, i)
             if (u_face > 0.0_dp) then
                b_xmin = x_pos - delta_x; b_xmax = x_pos
                b_ymin = y_face(mesh, j); b_ymax = y_face(mesh, j + 1)
                flux_vol(i, j) = cut_box_area(b_xmin, b_xmax, b_ymin, b_ymax, &
                                             segments(i-1, j)%nx, segments(i-1, j)%ny, segments(i-1, j)%alpha)
             else
                b_xmin = x_pos; b_xmax = x_pos + delta_x
                b_ymin = y_face(mesh, j); b_ymax = y_face(mesh, j + 1)
                flux_vol(i, j) = -cut_box_area(b_xmin, b_xmax, b_ymin, b_ymax, &
                                              segments(i, j)%nx, segments(i, j)%ny, segments(i, j)%alpha)
             end if
          end do
       end do

       c_star = fields%C
       do j = 1, ny
          do i = 1, nx
             denom = dx - dt * (fields%u(i + 1, j) - fields%u(i, j))
             if (denom > 1.0e-12_dp) then
                c_star(i, j) = (fields%C(i, j) * dx - (flux_vol(i + 1, j) - flux_vol(i, j)) / dy) / denom
             else
                c_star(i, j) = fields%C(i, j) - (flux_vol(i + 1, j) - flux_vol(i, j)) / (dx * dy)
             end if
             c_star(i, j) = max(0.0_dp, min(1.0_dp, c_star(i, j)))
          end do
       end do
       call exchange_halos_scalar(c_star, mesh)

       ! 2. Passe Y
       call reconstruct_interface(cfg%reconstruction_method, mesh, c_star, segments)
       flux_vol = 0.0_dp
       do j = 1, ny + 1
          do i = 1, nx
             v_face = fields%v(i, j)
             delta_y = abs(v_face) * dt
             if (delta_y <= 1.0e-14_dp) cycle
             y_pos = y_face(mesh, j)
             if (v_face > 0.0_dp) then
                b_xmin = x_face(mesh, i); b_xmax = x_face(mesh, i + 1)
                b_ymin = y_pos - delta_y; b_ymax = y_pos
                flux_vol(i, j) = cut_box_area(b_xmin, b_xmax, b_ymin, b_ymax, &
                                             segments(i, j-1)%nx, segments(i, j-1)%ny, segments(i, j-1)%alpha)
             else
                b_xmin = x_face(mesh, i); b_xmax = x_face(mesh, i + 1)
                b_ymin = y_pos; b_ymax = y_pos + delta_y
                flux_vol(i, j) = -cut_box_area(b_xmin, b_xmax, b_ymin, b_ymax, &
                                              segments(i, j)%nx, segments(i, j)%ny, segments(i, j)%alpha)
             end if
          end do
       end do

       c_final = c_star
       do j = 1, ny
          do i = 1, nx
             denom = dy - dt * (fields%v(i, j + 1) - fields%v(i, j))
             if (denom > 1.0e-12_dp) then
                c_final(i, j) = (c_star(i, j) * dy - (flux_vol(i, j + 1) - flux_vol(i, j)) / dx) / denom
             else
                c_final(i, j) = c_star(i, j) - (flux_vol(i, j + 1) - flux_vol(i, j)) / (dx * dy)
             end if
             c_final(i, j) = max(0.0_dp, min(1.0_dp, c_final(i, j)))
          end do
       end do
    else
       ! 1. Passe Y
       call reconstruct_interface(cfg%reconstruction_method, mesh, fields%C, segments)
       flux_vol = 0.0_dp
       do j = 1, ny + 1
          do i = 1, nx
             v_face = fields%v(i, j)
             delta_y = abs(v_face) * dt
             if (delta_y <= 1.0e-14_dp) cycle
             y_pos = y_face(mesh, j)
             if (v_face > 0.0_dp) then
                b_xmin = x_face(mesh, i); b_xmax = x_face(mesh, i + 1)
                b_ymin = y_pos - delta_y; b_ymax = y_pos
                flux_vol(i, j) = cut_box_area(b_xmin, b_xmax, b_ymin, b_ymax, &
                                             segments(i, j-1)%nx, segments(i, j-1)%ny, segments(i, j-1)%alpha)
             else
                b_xmin = x_face(mesh, i); b_xmax = x_face(mesh, i + 1)
                b_ymin = y_pos; b_ymax = y_pos + delta_y
                flux_vol(i, j) = -cut_box_area(b_xmin, b_xmax, b_ymin, b_ymax, &
                                              segments(i, j)%nx, segments(i, j)%ny, segments(i, j)%alpha)
             end if
          end do
       end do

       c_star = fields%C
       do j = 1, ny
          do i = 1, nx
             denom = dy - dt * (fields%v(i, j + 1) - fields%v(i, j))
             if (denom > 1.0e-12_dp) then
                c_star(i, j) = (fields%C(i, j) * dy - (flux_vol(i, j + 1) - flux_vol(i, j)) / dx) / denom
             else
                c_star(i, j) = fields%C(i, j) - (flux_vol(i, j + 1) - flux_vol(i, j)) / (dx * dy)
             end if
             c_star(i, j) = max(0.0_dp, min(1.0_dp, c_star(i, j)))
          end do
       end do
       call exchange_halos_scalar(c_star, mesh)

       ! 2. Passe X
       call reconstruct_interface(cfg%reconstruction_method, mesh, c_star, segments)
       flux_vol = 0.0_dp
       do j = 1, ny
          do i = 1, nx + 1
             u_face = fields%u(i, j)
             delta_x = abs(u_face) * dt
             if (delta_x <= 1.0e-14_dp) cycle
             x_pos = x_face(mesh, i)
             if (u_face > 0.0_dp) then
                b_xmin = x_pos - delta_x; b_xmax = x_pos
                b_ymin = y_face(mesh, j); b_ymax = y_face(mesh, j + 1)
                flux_vol(i, j) = cut_box_area(b_xmin, b_xmax, b_ymin, b_ymax, &
                                             segments(i-1, j)%nx, segments(i-1, j)%ny, segments(i-1, j)%alpha)
             else
                b_xmin = x_pos; b_xmax = x_pos + delta_x
                b_ymin = y_face(mesh, j); b_ymax = y_face(mesh, j + 1)
                flux_vol(i, j) = -cut_box_area(b_xmin, b_xmax, b_ymin, b_ymax, &
                                              segments(i, j)%nx, segments(i, j)%ny, segments(i, j)%alpha)
             end if
          end do
       end do

       c_final = c_star
       do j = 1, ny
          do i = 1, nx
             denom = dx - dt * (fields%u(i + 1, j) - fields%u(i, j))
             if (denom > 1.0e-12_dp) then
                c_final(i, j) = (c_star(i, j) * dx - (flux_vol(i + 1, j) - flux_vol(i, j)) / dy) / denom
             else
                c_final(i, j) = c_star(i, j) - (flux_vol(i + 1, j) - flux_vol(i, j)) / (dx * dy)
             end if
             c_final(i, j) = max(0.0_dp, min(1.0_dp, c_final(i, j)))
          end do
       end do
    end if

    fields%C = c_final
    deallocate(segments, flux_vol, c_star, c_final)
  end subroutine advect_vof_ciam

  ! [Procédure/Module] subroutine advect_vof_weymouth_yue
  subroutine advect_vof_weymouth_yue(fields, mesh, cfg, dt)
    type(flow_fields_t), intent(inout) :: fields
    type(mesh_2d_t), intent(in) :: mesh
    type(config_t), intent(in) :: cfg
    real(dp), intent(in) :: dt

    integer :: i, j, nx, ny
    real(dp) :: dx, dy, cell_area, u_face, v_face, delta_x, delta_y
    real(dp) :: x_pos, y_pos, b_xmin, b_xmax, b_ymin, b_ymax
    logical :: x_first
    type(interface_segment_t), allocatable :: segments(:, :)
    real(dp), allocatable :: flux_vol(:, :), c_star(:, :), c_final(:, :), c_d(:, :)

    nx = mesh%nx_local
    ny = mesh%ny_local
    dx = mesh%dx
    dy = mesh%dy
    cell_area = dx * dy
    x_first = (mod(vof_step_count, 2) == 0)

    allocate(segments(0:nx+1, 0:ny+1))
    allocate(flux_vol(0:nx+1, 0:ny+1))
    allocate(c_star(0:nx+1, 0:ny+1))
    allocate(c_final(0:nx+1, 0:ny+1))
    allocate(c_d(0:nx+1, 0:ny+1))

    ! Fonction indicatrice binaire de dilatation figée au temps n (Weymouth & Yue, JCP 2010, Eq. 14)
    do j = 0, ny + 1
       do i = 0, nx + 1
          if (fields%C(i, j) >= 0.5_dp) then
             c_d(i, j) = 1.0_dp
          else
             c_d(i, j) = 0.0_dp
          end if
       end do
    end do

    if (x_first) then
       ! 1. Passe X (Eq. 15 Weymouth & Yue 2010)
       call reconstruct_interface(cfg%reconstruction_method, mesh, fields%C, segments)
       flux_vol = 0.0_dp
       do j = 1, ny
          do i = 1, nx + 1
             u_face = fields%u(i, j)
             delta_x = abs(u_face) * dt
             if (delta_x <= 1.0e-14_dp) cycle
             x_pos = x_face(mesh, i)
             if (u_face > 0.0_dp) then
                b_xmin = x_pos - delta_x; b_xmax = x_pos
                b_ymin = y_face(mesh, j); b_ymax = y_face(mesh, j + 1)
                flux_vol(i, j) = cut_box_area(b_xmin, b_xmax, b_ymin, b_ymax, &
                                             segments(i-1, j)%nx, segments(i-1, j)%ny, segments(i-1, j)%alpha)
             else
                b_xmin = x_pos; b_xmax = x_pos + delta_x
                b_ymin = y_face(mesh, j); b_ymax = y_face(mesh, j + 1)
                flux_vol(i, j) = -cut_box_area(b_xmin, b_xmax, b_ymin, b_ymax, &
                                              segments(i, j)%nx, segments(i, j)%ny, segments(i, j)%alpha)
             end if
          end do
       end do

       c_star = fields%C
       do j = 1, ny
          do i = 1, nx
             c_star(i, j) = fields%C(i, j) - (flux_vol(i + 1, j) - flux_vol(i, j)) / cell_area + &
                            c_d(i, j) * (dt / dx) * (fields%u(i + 1, j) - fields%u(i, j))
          end do
       end do
       call exchange_halos_scalar(c_star, mesh)

       ! 2. Passe Y (Eq. 16 Weymouth & Yue 2010)
       call reconstruct_interface(cfg%reconstruction_method, mesh, c_star, segments)
       flux_vol = 0.0_dp
       do j = 1, ny + 1
          do i = 1, nx
             v_face = fields%v(i, j)
             delta_y = abs(v_face) * dt
             if (delta_y <= 1.0e-14_dp) cycle
             y_pos = y_face(mesh, j)
             if (v_face > 0.0_dp) then
                b_xmin = x_face(mesh, i); b_xmax = x_face(mesh, i + 1)
                b_ymin = y_pos - delta_y; b_ymax = y_pos
                flux_vol(i, j) = cut_box_area(b_xmin, b_xmax, b_ymin, b_ymax, &
                                             segments(i, j-1)%nx, segments(i, j-1)%ny, segments(i, j-1)%alpha)
             else
                b_xmin = x_face(mesh, i); b_xmax = x_face(mesh, i + 1)
                b_ymin = y_pos; b_ymax = y_pos + delta_y
                flux_vol(i, j) = -cut_box_area(b_xmin, b_xmax, b_ymin, b_ymax, &
                                              segments(i, j)%nx, segments(i, j)%ny, segments(i, j)%alpha)
             end if
          end do
       end do

       c_final = c_star
       do j = 1, ny
          do i = 1, nx
             c_final(i, j) = c_star(i, j) - (flux_vol(i, j + 1) - flux_vol(i, j)) / cell_area + &
                             c_d(i, j) * (dt / dy) * (fields%v(i, j + 1) - fields%v(i, j))
             c_final(i, j) = max(0.0_dp, min(1.0_dp, c_final(i, j)))
          end do
       end do
    else
       ! 1. Passe Y (Eq. 16 Weymouth & Yue 2010)
       call reconstruct_interface(cfg%reconstruction_method, mesh, fields%C, segments)
       flux_vol = 0.0_dp
       do j = 1, ny + 1
          do i = 1, nx
             v_face = fields%v(i, j)
             delta_y = abs(v_face) * dt
             if (delta_y <= 1.0e-14_dp) cycle
             y_pos = y_face(mesh, j)
             if (v_face > 0.0_dp) then
                b_xmin = x_face(mesh, i); b_xmax = x_face(mesh, i + 1)
                b_ymin = y_pos - delta_y; b_ymax = y_pos
                flux_vol(i, j) = cut_box_area(b_xmin, b_xmax, b_ymin, b_ymax, &
                                             segments(i, j-1)%nx, segments(i, j-1)%ny, segments(i, j-1)%alpha)
             else
                b_xmin = x_face(mesh, i); b_xmax = x_face(mesh, i + 1)
                b_ymin = y_pos; b_ymax = y_pos + delta_y
                flux_vol(i, j) = -cut_box_area(b_xmin, b_xmax, b_ymin, b_ymax, &
                                              segments(i, j)%nx, segments(i, j)%ny, segments(i, j)%alpha)
             end if
          end do
       end do

       c_star = fields%C
       do j = 1, ny
          do i = 1, nx
             c_star(i, j) = fields%C(i, j) - (flux_vol(i, j + 1) - flux_vol(i, j)) / cell_area + &
                            c_d(i, j) * (dt / dy) * (fields%v(i, j + 1) - fields%v(i, j))
          end do
       end do
       call exchange_halos_scalar(c_star, mesh)

       ! 2. Passe X (Eq. 15 Weymouth & Yue 2010)
       call reconstruct_interface(cfg%reconstruction_method, mesh, c_star, segments)
       flux_vol = 0.0_dp
       do j = 1, ny
          do i = 1, nx + 1
             u_face = fields%u(i, j)
             delta_x = abs(u_face) * dt
             if (delta_x <= 1.0e-14_dp) cycle
             x_pos = x_face(mesh, i)
             if (u_face > 0.0_dp) then
                b_xmin = x_pos - delta_x; b_xmax = x_pos
                b_ymin = y_face(mesh, j); b_ymax = y_face(mesh, j + 1)
                flux_vol(i, j) = cut_box_area(b_xmin, b_xmax, b_ymin, b_ymax, &
                                             segments(i-1, j)%nx, segments(i-1, j)%ny, segments(i-1, j)%alpha)
             else
                b_xmin = x_pos; b_xmax = x_pos + delta_x
                b_ymin = y_face(mesh, j); b_ymax = y_face(mesh, j + 1)
                flux_vol(i, j) = -cut_box_area(b_xmin, b_xmax, b_ymin, b_ymax, &
                                              segments(i, j)%nx, segments(i, j)%ny, segments(i, j)%alpha)
             end if
          end do
       end do

       c_final = c_star
       do j = 1, ny
          do i = 1, nx
             c_final(i, j) = c_star(i, j) - (flux_vol(i + 1, j) - flux_vol(i, j)) / cell_area + &
                             c_d(i, j) * (dt / dx) * (fields%u(i + 1, j) - fields%u(i, j))
             c_final(i, j) = max(0.0_dp, min(1.0_dp, c_final(i, j)))
          end do
       end do
    end if

    fields%C = c_final
    deallocate(segments, flux_vol, c_star, c_final, c_d)
  end subroutine advect_vof_weymouth_yue

  ! [Procédure/Module] subroutine vof_init
  subroutine vof_init(this, cfg)
    ! Déclaration de variable locale / paramètre de calcul
    class(interface_vof_t), intent(inout) :: this
    ! Déclaration de variable locale / paramètre de calcul
    type(config_t), intent(in) :: cfg
    ! Affectation et calcul numérique
    this%cfg = cfg
  end subroutine vof_init

  ! [Procédure/Module] subroutine vof_set_ic
  subroutine vof_set_ic(this, mesh, fields)
    ! Déclaration de variable locale / paramètre de calcul
    class(interface_vof_t), intent(in) :: this
    ! Déclaration de variable locale / paramètre de calcul
    type(mesh_2d_t), intent(in) :: mesh
    ! Déclaration de variable locale / paramètre de calcul
    type(flow_fields_t), intent(inout) :: fields
    ! Appel de sous-programme modulaire
    call set_initial_conditions(fields, mesh, this%cfg)
    ! Appel de sous-programme modulaire
    call update_physical_properties(fields, mesh, this%cfg)
    ! Appel de sous-programme modulaire
    call compute_curvature_and_csf(fields, mesh, this%cfg)
  end subroutine vof_set_ic

  ! [Procédure/Module] subroutine vof_advect_step
  subroutine vof_advect_step(this, mesh, fields, dt)
    ! Déclaration de variable locale / paramètre de calcul
    class(interface_vof_t), intent(in) :: this
    ! Déclaration de variable locale / paramètre de calcul
    type(mesh_2d_t), intent(in) :: mesh
    ! Déclaration de variable locale / paramètre de calcul
    type(flow_fields_t), intent(inout) :: fields
    ! Déclaration de variable locale / paramètre de calcul
    real(dp), intent(in) :: dt
    ! Appel de sous-programme modulaire
    call advect_vof(fields, mesh, this%cfg, dt)
    ! Appel de sous-programme modulaire
    call update_physical_properties(fields, mesh, this%cfg)
    ! Appel de sous-programme modulaire
    call compute_curvature_and_csf(fields, mesh, this%cfg)
  end subroutine vof_advect_step

end module module_interface_vof

\end{lstlisting}

\chapter{Générateur VTK XML Parallèle en Fortran 2018 (module_vtk_writer.f90)}
\section{Écriture Conforme des Fichiers \texttt{.vtr} et \texttt{.pvtr} (\texttt{module\_vtk\_writer.f90})}
Le module \texttt{src/module\_vtk\_writer.f90} (500 lignes) écrit les fichiers XML VTK lisibles sous ParaView et Python. Une attention particulière est portée au formatage des flottants IEEE-754 : en Fortran, le descripteur standard \texttt{ES15.7} omet parfois la lettre \texttt{E} lorsque l'exposant comporte trois chiffres (par exemple \texttt{1.2345678-105} lorsque $|x| < 10^{-99}$). L'utilisation du descripteur explicite à trois chiffres d'exposant \texttt{ES16.8E3} dans \texttt{module\_vtk\_writer.f90} garantit que la lettre \texttt{E} est toujours présente, assurant une compatibilité universelle avec tous les parseurs XML VTK.

\section{Répertoire des Types Dérivés, Sous-Routines et Fonctions de \texttt{src/module\_vtk\_writer.f90}}
\begin{center}\small
\begin{tabular}{clll}
\toprule
\textbf{Ligne} & \textbf{Nature} & \textbf{Identifiant Fortran 2018} & \textbf{Interface / Arguments} \\
\midrule
10 & module & \texttt{module\_vtk\_writer} & Unité modulaire Fortran 2018 \\
21 & type dérivé & \texttt{vtk\_writer\_t} & Structure de données Fortran 2018 \\
46 & subroutine & \texttt{vtk\_init} & Arguments : (this, mesh, out\_dir, name) \\
89 & subroutine & \texttt{vtk\_write\_step} & Arguments : (this, mesh, fields, step, time) \\
368 & subroutine & \texttt{vtk\_write\_diagnostics} & Arguments : (this, mesh, fields, step, time, dt) \\
489 & subroutine & \texttt{vtk\_finalize} & Arguments : (this) \\
\bottomrule
\end{tabular}
\end{center}

\section{Code Source Intégral Annoté : \texttt{src/module\_vtk\_writer.f90}}
Le listing ci-dessous présente l'intégralité des 500 lignes du fichier \texttt{src/module\_vtk\_writer.f90} de l'implémentation Fortran 2018 :

\begin{lstlisting}[language={[95]Fortran}]
!===============================================================================
! Module: module_vtk_writer
! Description: Exportateur de maillage rectiligne XML VTK parallèle (.pvtr / .vtr)
!              et enregistreur de séries temporelles CSV pour COVOF 2D.
!              - Écriture locale non-bloquante du fichier .vtr sur chaque processus
!              - Assemblage du fichier maître .pvtr sur le rang racine
!              - Diagnostics diphasiques : position du centre de masse, vitesse
!                d'ascension, volume de bulle, énergie cinétique, divergence max
!===============================================================================
module module_vtk_writer
  use module_precision, only: dp
  use module_mesh, only: mesh_2d_t, x_center, y_center
  use module_fields, only: flow_fields_t
  use module_mpi_comm, only: global_sum, global_max
  implicit none
  private
  include 'mpif.h'

  public :: vtk_writer_t

  type :: vtk_writer_t
     ! Déclaration de variable locale / paramètre de calcul
     character(len=256) :: output_dir = "output"
     ! Déclaration de variable locale / paramètre de calcul
     character(len=128) :: case_name  = "covof2d_case"
     ! Déclaration de variable locale / paramètre de calcul
     integer            :: csv_unit   = 0
     ! Déclaration de variable locale / paramètre de calcul
     logical            :: is_open    = .false.
  contains
     ! Affectation et calcul numérique
     procedure :: init => vtk_init
     ! Affectation et calcul numérique
     procedure :: write_step => vtk_write_step
     ! Affectation et calcul numérique
     procedure :: write_diagnostics => vtk_write_diagnostics
     ! Affectation et calcul numérique
     procedure :: finalize => vtk_finalize
  end type vtk_writer_t

contains

  !-----------------------------------------------------------------------------
  ! Initialisation du répertoire de sortie et du fichier de métriques CSV
  !-----------------------------------------------------------------------------
  subroutine vtk_init(this, mesh, out_dir, name)
    ! Déclaration de variable locale / paramètre de calcul
    class(vtk_writer_t), intent(inout) :: this
    ! Déclaration de variable locale / paramètre de calcul
    type(mesh_2d_t), intent(in)        :: mesh
    ! Déclaration de variable locale / paramètre de calcul
    character(len=*), intent(in)       :: out_dir, name

    ! Déclaration de variable locale / paramètre de calcul
    character(len=512) :: csv_path, mkdir_cmd
    ! Déclaration de variable locale / paramètre de calcul
    integer :: io_stat

    ! Affectation et calcul numérique
    this%output_dir = trim(out_dir)
    ! Affectation et calcul numérique
    this%case_name  = trim(name)

    ! Évaluation de la condition logique
    if (mesh%rank == 0) then
       ! Affectation et calcul numérique
       mkdir_cmd = "mkdir -p " // trim(this%output_dir)
       ! Appel de sous-programme modulaire
       call execute_command_line(trim(mkdir_cmd))

       ! Affectation et calcul numérique
       csv_path = trim(this%output_dir) // "/" // trim(this%case_name) // "_history.csv"
       ! Affectation et calcul numérique
       open(newunit=this%csv_unit, file=trim(csv_path), status="replace", action="write", iostat=io_stat)
       ! Évaluation de la condition logique
       if (io_stat == 0) then
          ! Affectation et calcul numérique
          this%is_open = .true.
          write(this%csv_unit, '(A)') &
            "step,time,dt,center_of_mass_y,rise_velocity,bubble_volume,kinetic_energy,max_u,max_div"
          flush(this%csv_unit)
       end if
    end if
  end subroutine vtk_init

  !-----------------------------------------------------------------------------
  ! Écriture d'un pas de temps complet en format XML RectilinearGrid (.pvtr / .vtr)
  !-----------------------------------------------------------------------------
  subroutine vtk_write_step(this, mesh, fields, step, time)
    ! Déclaration de variable locale / paramètre de calcul
    class(vtk_writer_t), intent(in) :: this
    ! Déclaration de variable locale / paramètre de calcul
    type(mesh_2d_t), intent(in)     :: mesh
    ! Déclaration de variable locale / paramètre de calcul
    type(flow_fields_t), intent(in) :: fields
    ! Déclaration de variable locale / paramètre de calcul
    integer, intent(in)             :: step
    ! Déclaration de variable locale / paramètre de calcul
    real(dp), intent(in)            :: time

    ! Déclaration de variable locale / paramètre de calcul
    character(len=6)   :: step_str
    ! Déclaration de variable locale / paramètre de calcul
    character(len=512) :: vtr_file, pvtr_file, piece_str
    ! Déclaration de variable locale / paramètre de calcul
    integer :: x_start, x_end, y_start, y_end
    ! Déclaration de variable locale / paramètre de calcul
    integer :: vtr_unit, pvtr_unit, io_stat, i, j, p, ierr, dummy(1)
    ! Déclaration de variable locale / paramètre de calcul
    real(dp) :: uc, vc, div_val
    ! Déclaration de variable locale / paramètre de calcul
    integer, allocatable :: all_x_start(:), all_x_end(:), all_y_start(:), all_y_end(:)

    write(step_str, '(I6.6)') step

    ! Affectation et calcul numérique
    x_start = mesh%i_start_global - 1
    ! Affectation et calcul numérique
    x_end   = x_start + mesh%nx_local
    ! Affectation et calcul numérique
    y_start = mesh%j_start_global - 1
    ! Affectation et calcul numérique
    y_end   = y_start + mesh%ny_local

    ! 1. Écriture du fichier .vtr propre à chaque processus MPI
    write(vtr_file, '(A,"/",A,"_p",I0,"_",A,".vtr")') &
      trim(this%output_dir), trim(this%case_name), mesh%rank, step_str

    ! Affectation et calcul numérique
    open(newunit=vtr_unit, file=trim(vtr_file), status="replace", action="write", iostat=io_stat)
    ! Évaluation de la condition logique
    if (io_stat == 0) then
       ! Affectation et calcul numérique
       write(vtr_unit, '(A)') '<?xml version="1.0"?>'
       ! Affectation et calcul numérique
       write(vtr_unit, '(A)') '<VTKFile type="RectilinearGrid" version="0.1" byte_order="LittleEndian">'
       write(vtr_unit, '(A,I0,A,I0,A,I0,A,I0,A)') &
         ! Affectation et calcul numérique
         '  <RectilinearGrid WholeExtent="', x_start, ' ', x_end, ' ', y_start, ' ', y_end, ' 0 0">'
       write(vtr_unit, '(A,I0,A,I0,A,I0,A,I0,A)') &
         ! Affectation et calcul numérique
         '    <Piece Extent="', x_start, ' ', x_end, ' ', y_start, ' ', y_end, ' 0 0">'
       ! Affectation et calcul numérique
       write(vtr_unit, '(A)') '      <CellData Scalars="PhaseIndicator,Pressure,Density,Viscosity,Divergence" Vectors="Velocity">'

       ! PhaseIndicator C
       write(vtr_unit, '(A)') '        <DataArray type="Float64" Name="PhaseIndicator" format="ascii">'
       ! Boucle spatiale sur les cellules locales selon l'axe vertical Y (ordonnées)
       do j = 1, mesh%ny_local
          ! Boucle spatiale selon l'axe horizontal X (accès contigu colonne-major)
          do i = 1, mesh%nx_local
             ! Affectation et calcul numérique
             write(vtr_unit, '(ES16.8," ")', advance="no") fields%C(i, j)
          end do
          write(vtr_unit, *)
       end do
       write(vtr_unit, '(A)') '        </DataArray>'

       ! Pressure p
       write(vtr_unit, '(A)') '        <DataArray type="Float64" Name="Pressure" format="ascii">'
       ! Boucle spatiale sur les cellules locales selon l'axe vertical Y (ordonnées)
       do j = 1, mesh%ny_local
          ! Boucle spatiale selon l'axe horizontal X (accès contigu colonne-major)
          do i = 1, mesh%nx_local
             ! Affectation et calcul numérique
             write(vtr_unit, '(ES16.8," ")', advance="no") fields%p(i, j)
          end do
          write(vtr_unit, *)
       end do
       write(vtr_unit, '(A)') '        </DataArray>'

       ! Density rho
       write(vtr_unit, '(A)') '        <DataArray type="Float64" Name="Density" format="ascii">'
       ! Boucle spatiale sur les cellules locales selon l'axe vertical Y (ordonnées)
       do j = 1, mesh%ny_local
          ! Boucle spatiale selon l'axe horizontal X (accès contigu colonne-major)
          do i = 1, mesh%nx_local
             ! Affectation et calcul numérique
             write(vtr_unit, '(ES16.8," ")', advance="no") fields%rho(i, j)
          end do
          write(vtr_unit, *)
       end do
       write(vtr_unit, '(A)') '        </DataArray>'

       ! Viscosity mu
       write(vtr_unit, '(A)') '        <DataArray type="Float64" Name="Viscosity" format="ascii">'
       ! Boucle spatiale sur les cellules locales selon l'axe vertical Y (ordonnées)
       do j = 1, mesh%ny_local
          ! Boucle spatiale selon l'axe horizontal X (accès contigu colonne-major)
          do i = 1, mesh%nx_local
             ! Affectation et calcul numérique
             write(vtr_unit, '(ES16.8," ")', advance="no") fields%mu(i, j)
          end do
          write(vtr_unit, *)
       end do
       write(vtr_unit, '(A)') '        </DataArray>'

       ! Divergence div(u)
       write(vtr_unit, '(A)') '        <DataArray type="Float64" Name="Divergence" format="ascii">'
       ! Boucle spatiale sur les cellules locales selon l'axe vertical Y (ordonnées)
       do j = 1, mesh%ny_local
          ! Boucle spatiale selon l'axe horizontal X (accès contigu colonne-major)
          do i = 1, mesh%nx_local
             ! Affectation et calcul numérique
             div_val = (fields%u(i + 1, j) - fields%u(i, j)) / mesh%dx + &
                       (fields%v(i, j + 1) - fields%v(i, j)) / mesh%dy
             ! Affectation et calcul numérique
             write(vtr_unit, '(ES16.8," ")', advance="no") div_val
          end do
          write(vtr_unit, *)
       end do
       write(vtr_unit, '(A)') '        </DataArray>'

       ! Temperature T
       write(vtr_unit, '(A)') '        <DataArray type="Float64" Name="Temperature" format="ascii">'
       ! Boucle spatiale sur les cellules locales selon l'axe vertical Y (ordonnées)
       do j = 1, mesh%ny_local
          ! Boucle spatiale selon l'axe horizontal X (accès contigu colonne-major)
          do i = 1, mesh%nx_local
             ! Affectation et calcul numérique
             write(vtr_unit, '(ES16.8," ")', advance="no") fields%T(i, j)
          end do
          write(vtr_unit, *)
       end do
       write(vtr_unit, '(A)') '        </DataArray>'

       ! Velocity Vector (au centre des mailles)
       write(vtr_unit, '(A)') '        <DataArray type="Float64" Name="Velocity" NumberOfComponents="3" format="ascii">'
       ! Boucle spatiale sur les cellules locales selon l'axe vertical Y (ordonnées)
       do j = 1, mesh%ny_local
          ! Boucle spatiale selon l'axe horizontal X (accès contigu colonne-major)
          do i = 1, mesh%nx_local
             ! Affectation et calcul numérique
             uc = 0.5_dp * (fields%u(i, j) + fields%u(i + 1, j))
             ! Affectation et calcul numérique
             vc = 0.5_dp * (fields%v(i, j) + fields%v(i, j + 1))
             ! Affectation et calcul numérique
             write(vtr_unit, '(ES16.8," ",ES16.8," 0.0 ")', advance="no") uc, vc
          end do
          write(vtr_unit, *)
       end do
       write(vtr_unit, '(A)') '        </DataArray>'

       write(vtr_unit, '(A)') '      </CellData>'
       write(vtr_unit, '(A)') '      <Coordinates>'

       ! Coordonnées X
       write(vtr_unit, '(A)') '        <DataArray type="Float64" Name="XCoordinates" format="ascii">'
       ! Boucle spatiale selon l'axe horizontal X (accès contigu colonne-major)
       do i = 0, mesh%nx_local
          ! Affectation et calcul numérique
          write(vtr_unit, '(ES16.8," ")', advance="no") real(mesh%i_start_global - 1 + i, dp) * mesh%dx
       end do
       write(vtr_unit, *)
       write(vtr_unit, '(A)') '        </DataArray>'

       ! Coordonnées Y
       write(vtr_unit, '(A)') '        <DataArray type="Float64" Name="YCoordinates" format="ascii">'
       ! Boucle spatiale sur les cellules locales selon l'axe vertical Y (ordonnées)
       do j = 0, mesh%ny_local
          ! Affectation et calcul numérique
          write(vtr_unit, '(ES16.8," ")', advance="no") real(mesh%j_start_global - 1 + j, dp) * mesh%dy
       end do
       write(vtr_unit, *)
       write(vtr_unit, '(A)') '        </DataArray>'

       ! Coordonnées Z
       write(vtr_unit, '(A)') '        <DataArray type="Float64" Name="ZCoordinates" format="ascii">'
       write(vtr_unit, '(A)') '          0.0'
       write(vtr_unit, '(A)') '        </DataArray>'

       write(vtr_unit, '(A)') '      </Coordinates>'
       write(vtr_unit, '(A)') '    </Piece>'
       write(vtr_unit, '(A)') '  </RectilinearGrid>'
       write(vtr_unit, '(A)') '</VTKFile>'
       close(vtr_unit)
    end if

    ! 2. Écriture du fichier maître parallèle .pvtr sur le rang 0
    if (mesh%rank == 0) then
       ! Allocation dynamique de tableau contigu en mémoire
       allocate(all_x_start(mesh%size), all_x_end(mesh%size))
       ! Allocation dynamique de tableau contigu en mémoire
       allocate(all_y_start(mesh%size), all_y_end(mesh%size))

       ! Appel à la bibliothèque de communication parallèle MPI
       call MPI_Gather(x_start, 1, MPI_INTEGER, all_x_start, 1, MPI_INTEGER, 0, mesh%comm, ierr)
       ! Appel à la bibliothèque de communication parallèle MPI
       call MPI_Gather(x_end,   1, MPI_INTEGER, all_x_end,   1, MPI_INTEGER, 0, mesh%comm, ierr)
       ! Appel à la bibliothèque de communication parallèle MPI
       call MPI_Gather(y_start, 1, MPI_INTEGER, all_y_start, 1, MPI_INTEGER, 0, mesh%comm, ierr)
       ! Appel à la bibliothèque de communication parallèle MPI
       call MPI_Gather(y_end,   1, MPI_INTEGER, all_y_end,   1, MPI_INTEGER, 0, mesh%comm, ierr)

       write(pvtr_file, '(A,"/",A,"_",A,".pvtr")') &
         trim(this%output_dir), trim(this%case_name), step_str

       ! Affectation et calcul numérique
       open(newunit=pvtr_unit, file=trim(pvtr_file), status="replace", action="write", iostat=io_stat)
       ! Évaluation de la condition logique
       if (io_stat == 0) then
          ! Affectation et calcul numérique
          write(pvtr_unit, '(A)') '<?xml version="1.0"?>'
          ! Affectation et calcul numérique
          write(pvtr_unit, '(A)') '<VTKFile type="PRectilinearGrid" version="0.1" byte_order="LittleEndian">'
          write(pvtr_unit, '(A,I0,A,I0,A)') &
            ! Affectation et calcul numérique
            '  <PRectilinearGrid WholeExtent="0 ', mesh%nx_global, ' 0 ', mesh%ny_global, ' 0 0" GhostLevel="0">'
          ! Affectation et calcul numérique
          write(pvtr_unit, '(A)') '    <PCellData Scalars="PhaseIndicator,Pressure,Density,Viscosity,Divergence" Vectors="Velocity">'
          ! Affectation et calcul numérique
          write(pvtr_unit, '(A)') '      <PDataArray type="Float64" Name="PhaseIndicator"/>'
          ! Affectation et calcul numérique
          write(pvtr_unit, '(A)') '      <PDataArray type="Float64" Name="Pressure"/>'
          ! Affectation et calcul numérique
          write(pvtr_unit, '(A)') '      <PDataArray type="Float64" Name="Density"/>'
          ! Affectation et calcul numérique
          write(pvtr_unit, '(A)') '      <PDataArray type="Float64" Name="Viscosity"/>'
          ! Affectation et calcul numérique
          write(pvtr_unit, '(A)') '      <PDataArray type="Float64" Name="Divergence"/>'
          ! Affectation et calcul numérique
          write(pvtr_unit, '(A)') '      <PDataArray type="Float64" Name="Temperature"/>'
          ! Affectation et calcul numérique
          write(pvtr_unit, '(A)') '      <PDataArray type="Float64" Name="Velocity" NumberOfComponents="3"/>'
          write(pvtr_unit, '(A)') '    </PCellData>'
          write(pvtr_unit, '(A)') '    <PCoordinates>'
          ! Affectation et calcul numérique
          write(pvtr_unit, '(A)') '      <PDataArray type="Float64" Name="XCoordinates"/>'
          ! Affectation et calcul numérique
          write(pvtr_unit, '(A)') '      <PDataArray type="Float64" Name="YCoordinates"/>'
          ! Affectation et calcul numérique
          write(pvtr_unit, '(A)') '      <PDataArray type="Float64" Name="ZCoordinates"/>'
          write(pvtr_unit, '(A)') '    </PCoordinates>'

          ! Boucle d'itération numérique
          do p = 0, mesh%size - 1
             write(piece_str, '(A,"_p",I0,"_",A,".vtr")') trim(this%case_name), p, step_str
             write(pvtr_unit, '(A,I0,A,I0,A,I0,A,I0,A,A,A)') &
               ! Affectation et calcul numérique
               '    <Piece Extent="', all_x_start(p+1), ' ', all_x_end(p+1), ' ', &
               ! Affectation et calcul numérique
               all_y_start(p+1), ' ', all_y_end(p+1), ' 0 0" Source="', trim(piece_str), '"/>'
          end do

          write(pvtr_unit, '(A)') '  </PRectilinearGrid>'
          write(pvtr_unit, '(A)') '</VTKFile>'
          close(pvtr_unit)
       end if

       ! Libération de la mémoire allouée
       deallocate(all_x_start, all_x_end, all_y_start, all_y_end)
    ! Branche par défaut si condition non satisfaite
    else
       ! Appel à la bibliothèque de communication parallèle MPI
       call MPI_Gather(x_start, 1, MPI_INTEGER, dummy, 0, MPI_INTEGER, 0, mesh%comm, ierr)
       ! Appel à la bibliothèque de communication parallèle MPI
       call MPI_Gather(x_end,   1, MPI_INTEGER, dummy, 0, MPI_INTEGER, 0, mesh%comm, ierr)
       ! Appel à la bibliothèque de communication parallèle MPI
       call MPI_Gather(y_start, 1, MPI_INTEGER, dummy, 0, MPI_INTEGER, 0, mesh%comm, ierr)
       ! Appel à la bibliothèque de communication parallèle MPI
       call MPI_Gather(y_end,   1, MPI_INTEGER, dummy, 0, MPI_INTEGER, 0, mesh%comm, ierr)
    end if
  end subroutine vtk_write_step

  !-----------------------------------------------------------------------------
  ! Calcul et enregistrement des métriques physiques temporelles dans l'historique CSV
  !-----------------------------------------------------------------------------
  subroutine vtk_write_diagnostics(this, mesh, fields, step, time, dt)
    ! Déclaration de variable locale / paramètre de calcul
    class(vtk_writer_t), intent(in) :: this
    ! Déclaration de variable locale / paramètre de calcul
    type(mesh_2d_t), intent(in)     :: mesh
    ! Déclaration de variable locale / paramètre de calcul
    type(flow_fields_t), intent(in) :: fields
    ! Déclaration de variable locale / paramètre de calcul
    integer, intent(in)             :: step
    ! Déclaration de variable locale / paramètre de calcul
    real(dp), intent(in)            :: time, dt

    ! Déclaration de variable locale / paramètre de calcul
    real(dp) :: local_vol, local_y_integral, local_v_integral, local_ke
    ! Déclaration de variable locale / paramètre de calcul
    real(dp) :: local_max_u, local_max_div, dvol, y_pos, c_val, ind_bubble
    ! Déclaration de variable locale / paramètre de calcul
    real(dp) :: uc, vc, speed2, speed, div_val
    ! Déclaration de variable locale / paramètre de calcul
    real(dp) :: global_vol, global_y_integral, global_v_integral, global_ke
    ! Déclaration de variable locale / paramètre de calcul
    real(dp) :: global_max_u, global_max_div, yc, vc_mean
    ! Déclaration de variable locale / paramètre de calcul
    integer  :: i, j

    ! Affectation et calcul numérique
    local_vol        = 0.0_dp
    ! Affectation et calcul numérique
    local_y_integral = 0.0_dp
    ! Affectation et calcul numérique
    local_v_integral = 0.0_dp
    ! Affectation et calcul numérique
    local_ke         = 0.0_dp
    ! Affectation et calcul numérique
    local_max_u      = 0.0_dp
    ! Affectation et calcul numérique
    local_max_div    = 0.0_dp

    ! Affectation et calcul numérique
    dvol = mesh%dx * mesh%dy

    ! Boucle spatiale sur les cellules locales selon l'axe vertical Y (ordonnées)
    do j = 1, mesh%ny_local
       ! Affectation et calcul numérique
       y_pos = y_center(mesh, j)
       ! Boucle spatiale selon l'axe horizontal X (accès contigu colonne-major)
       do i = 1, mesh%nx_local
          ! Affectation et calcul numérique
          c_val = fields%C(i, j)
          ! Affectation et calcul numérique
          ind_bubble = 1.0_dp - c_val ! 1 si phase fluide 2 (bulle / goutte)

          ! Affectation et calcul numérique
          local_vol        = local_vol + ind_bubble * dvol
          ! Affectation et calcul numérique
          local_y_integral = local_y_integral + y_pos * ind_bubble * dvol

          ! Affectation et calcul numérique
          uc = 0.5_dp * (fields%u(i, j) + fields%u(i + 1, j))
          ! Affectation et calcul numérique
          vc = 0.5_dp * (fields%v(i, j) + fields%v(i, j + 1))

          ! Affectation et calcul numérique
          local_v_integral = local_v_integral + vc * ind_bubble * dvol

          ! Affectation et calcul numérique
          speed2 = uc * uc + vc * vc
          ! Affectation et calcul numérique
          local_ke = local_ke + 0.5_dp * fields%rho(i, j) * speed2 * dvol

          ! Affectation et calcul numérique
          speed = sqrt(speed2)
          ! Affectation et calcul numérique
          if (speed > local_max_u) local_max_u = speed

          ! Affectation et calcul numérique
          div_val = abs((fields%u(i + 1, j) - fields%u(i, j)) / mesh%dx + &
                        (fields%v(i, j + 1) - fields%v(i, j)) / mesh%dy)
          ! Affectation et calcul numérique
          if (div_val > local_max_div) local_max_div = div_val
       end do
    end do

    ! Affectation et calcul numérique
    global_vol        = global_sum(mesh, local_vol)
    ! Affectation et calcul numérique
    global_y_integral = global_sum(mesh, local_y_integral)
    ! Affectation et calcul numérique
    global_v_integral = global_sum(mesh, local_v_integral)
    ! Affectation et calcul numérique
    global_ke         = global_sum(mesh, local_ke)
    ! Affectation et calcul numérique
    global_max_u      = global_max(mesh, local_max_u)
    ! Affectation et calcul numérique
    global_max_div    = global_max(mesh, local_max_div)

    ! Évaluation de la condition logique
    if (global_vol > 1.0e-12_dp) then
       ! Affectation et calcul numérique
       yc = global_y_integral / global_vol
       ! Affectation et calcul numérique
       vc_mean = global_v_integral / global_vol
    ! Branche par défaut si condition non satisfaite
    else
       ! Affectation et calcul numérique
       yc = 0.0_dp
       ! Affectation et calcul numérique
       vc_mean = 0.0_dp
    end if

    ! Évaluation de la condition logique
    if (mesh%rank == 0 .and. this%is_open) then
       write(this%csv_unit, '(I0,",",ES14.6,",",ES14.6,",",ES14.6,",",ES14.6,",",ES14.6,",",ES14.6,",",ES14.6,",",ES14.6)') &
         step, time, dt, yc, vc_mean, global_vol, global_ke, global_max_u, global_max_div
       flush(this%csv_unit)
    end if
  end subroutine vtk_write_diagnostics

  !-----------------------------------------------------------------------------
  ! Clôture du fichier de diagnostic CSV
  !-----------------------------------------------------------------------------
  subroutine vtk_finalize(this)
    ! Déclaration de variable locale / paramètre de calcul
    class(vtk_writer_t), intent(inout) :: this
    ! Évaluation de la condition logique
    if (this%is_open) then
       close(this%csv_unit)
       ! Affectation et calcul numérique
       this%is_open = .false.
    end if
  end subroutine vtk_finalize

end module module_vtk_writer

\end{lstlisting}

\chapter{Boucle Temporelle Principale, Diagnostics Intégraux et Vérification en Fortran 2018 (module_simulation.f90)}
\section{Orchestration Complète dans \texttt{module\_simulation.f90} (1004 Lignes)}
Le module \texttt{src/module\_simulation.f90} réalise l'intégration temporelle complète et l'auto-vérification des 16 cas de validation :
\begin{enumerate}
  \item \texttt{simulation\_init} : initialise le contexte MPI, le maillage MAC, les champs hydrodynamiques/thermiques/VOF, les masques de Brinkman et les conditions initiales sous-échantillonnées $8\times 8$.
  \item \texttt{simulation\_run} : aiguille vers le mode elliptique pur (\texttt{run\_laplacian\_test} pour le Cas 3.1), le mode cinématique de Rider-Kothe (\texttt{run\_vortex\_advection} pour les 12 variantes du Cas 3.19) ou la boucle complète Navier-Stokes + VOF + Énergie.
  \item \texttt{write\_diagnostics\_csv} : calcule à chaque pas de temps les intégrales physiques exactes (volume de phase, centre de masse $(X_c, Y_c)$, vitesse moyenne $(U_c, V_c)$, énergie cinétique diphasique $E_k(t) = \frac{1}{2}\iint \rho(C)\|\mathbf{u}\|^2\,d\Omega$, enstrophie $\mathcal{E}(t) = \frac{1}{2}\iint \omega_z^2\,d\Omega$, nombre de Nusselt $\overline{Nu}$ et saut de pression de Laplace $\Delta P$) et les exporte dans les fichiers CSV synchronisés avec C++17.
\end{enumerate}

\section{Répertoire des Types Dérivés, Sous-Routines et Fonctions de \texttt{src/module\_simulation.f90}}
\begin{center}\small
\begin{tabular}{clll}
\toprule
\textbf{Ligne} & \textbf{Nature} & \textbf{Identifiant Fortran 2018} & \textbf{Interface / Arguments} \\
\midrule
11 & module & \texttt{module\_simulation} & Unité modulaire Fortran 2018 \\
28 & type dérivé & \texttt{simulation\_t} & Structure de données Fortran 2018 \\
59 & function & \texttt{sim\_init} & Arguments : (this, comm, config\_file) \\
266 & function & \texttt{sim\_compute\_time\_step} & Arguments : (this) \\
330 & subroutine & \texttt{sim\_run} & Arguments : (this) \\
561 & subroutine & \texttt{sim\_run\_laplacian} & Arguments : (this) \\
775 & subroutine & \texttt{sim\_run\_vortex\_advection} & Arguments : (this) \\
\bottomrule
\end{tabular}
\end{center}

\section{Code Source Intégral Annoté : \texttt{src/module\_simulation.f90}}
Le listing ci-dessous présente l'intégralité des 1011 lignes du fichier \texttt{src/module\_simulation.f90} de l'implémentation Fortran 2018 :

\begin{lstlisting}[language={[95]Fortran}]
!===============================================================================
! Module: module_simulation
! Description: Moteur principal de simulation COVOF 2D en Fortran 2018 :
!              - Gestion du cycle de vie du calcul (init, run, finalize)
!              - Pas de temps adaptatif CFL + onde capillaire
!              - Boucle temporelle avec couplage Lagrangien Augmenté / Uzawa
!              - Mode de vérification analytique du Laplacien (Cas 3.1 Référence Académique)
!              - Mode de vérification du vortex de cisaillement réversible (Cas 3.19 Référence Académique)
!              - Convergence vers l'état stationnaire
!===============================================================================
module module_simulation
  use module_precision, only: dp
  use module_constants
  use module_config
  use module_mesh
  use module_fields
  use module_mpi_comm
  use module_interface_vof
  use module_navier_stokes
  use module_linear_solvers
  use module_vtk_writer
  implicit none
  private
  include 'mpif.h'

  public :: simulation_t

  type :: simulation_t
     ! Déclaration de variable locale / paramètre de calcul
     type(config_t)         :: config
     ! Déclaration de variable locale / paramètre de calcul
     type(mesh_2d_t)        :: mesh
     ! Déclaration de variable locale / paramètre de calcul
     type(flow_fields_t)    :: fields
     ! Déclaration de variable locale / paramètre de calcul
     type(interface_vof_t)  :: tracker
     ! Déclaration de variable locale / paramètre de calcul
     type(navier_stokes_t)  :: ns_solver
     ! Déclaration de variable locale / paramètre de calcul
     type(vtk_writer_t)     :: vtk_writer
  contains
     ! Affectation et calcul numérique
     procedure :: init => sim_init
     ! Affectation et calcul numérique
     procedure :: run => sim_run
     ! Affectation et calcul numérique
     procedure :: compute_time_step => sim_compute_time_step
     ! Affectation et calcul numérique
     procedure :: run_laplacian => sim_run_laplacian
     ! Affectation et calcul numérique
     procedure :: run_vortex_advection => sim_run_vortex_advection
  end type simulation_t

contains

  !-----------------------------------------------------------------------------
  ! Initialisation complète de la simulation à partir du fichier JSON
  !-----------------------------------------------------------------------------
  function sim_init(this, comm, config_file) result(success)
    ! Déclaration de variable locale / paramètre de calcul
    class(simulation_t), intent(inout) :: this
    ! Déclaration de variable locale / paramètre de calcul
    integer, intent(in)                :: comm
    ! Déclaration de variable locale / paramètre de calcul
    character(len=*), intent(in)       :: config_file
    ! Déclaration de variable locale / paramètre de calcul
    logical                            :: success

    ! Déclaration de variable locale / paramètre de calcul
    logical  :: periodic_x, periodic_y
    ! Déclaration de variable locale / paramètre de calcul
    integer  :: i, j, u_nx, v_ny
    ! Déclaration de variable locale / paramètre de calcul
    real(dp) :: frac, y_step, h, u_mean, y_val, u_val, x_val
    ! Déclaration des paramètres pour la nappe liquide cisaillée (Cas 3.2 Fluids 2021)
    real(dp) :: y_c_ls, delta_ls, k_ls, amp_ls, factor_u_ls, dy_c_ls, sgn_y_ls, exp_decay_ls, c_u_ls, a_val_ls
    ! Déclaration des variables pour l'amorce de jet visqueux (Cas 3.3 Fluids 2021)
    real(dp) :: x_c_vj, half_w_vj, stub_h_vj, v_in_vj

    ! Affectation et calcul numérique
    success = .false.

    ! 1. Chargement de la configuration JSON
    call load_config_from_json(config_file, this%config)

    ! 2. Initialisation du maillage cartésien partitionné MPI
    periodic_x = (this%config%bc_left%btype == BC_PERIODIC .or. this%config%bc_right%btype == BC_PERIODIC)
    ! Affectation et calcul numérique
    periodic_y = (this%config%bc_bottom%btype == BC_PERIODIC .or. this%config%bc_top%btype == BC_PERIODIC)

    ! Appel de sous-programme modulaire
    call init_mesh(this%mesh, comm, this%config%nx, this%config%ny, &
                   this%config%lx, this%config%ly, periodic_x, periodic_y)

    ! 3. Allocation des champs décalés MAC
    call this%fields%allocate(this%mesh)

    ! 4. Initialisation du module d'interface VOF et des conditions initiales
    call this%tracker%init(this%config)
    ! Appel de sous-programme modulaire
    call this%tracker%set_initial_condition(this%mesh, this%fields)

    ! Initialisation des conditions aux limites initiales
    call apply_boundary_conditions(this%mesh, this%config, this%fields, .false.)

    ! Initialisation du canal amont pour profil POISEUILLE_STEP afin de garantir la continuité
    if (this%config%bc_left%btype == BC_DIRICHLET .and. this%config%obs_enabled .and. &
        ! Affectation et calcul numérique
        this%config%obs_type == OBS_BOX) then
       ! Affectation et calcul numérique
       y_step = this%config%obs_y1
       ! Affectation et calcul numérique
       h = this%mesh%ly - y_step
       ! Affectation et calcul numérique
       u_mean = this%config%bc_left%u
       ! Affectation et calcul numérique
       u_nx = this%mesh%nx_local
       ! Affectation et calcul numérique
       if (this%mesh%is_boundary_right) u_nx = u_nx + 1

       ! Boucle spatiale sur les cellules locales selon l'axe vertical Y (ordonnées)
       do j = 1, this%mesh%ny_local
          ! Affectation et calcul numérique
          y_val = y_center(this%mesh, j)
          ! Affectation et calcul numérique
          u_val = 0.0_dp
          ! Évaluation de la condition logique
          if (y_val >= y_step) then
             ! Affectation et calcul numérique
             u_val = 6.0_dp * u_mean * (y_val - y_step) * (this%mesh%ly - y_val) / (h * h)
          end if
          ! Boucle spatiale selon l'axe horizontal X (accès contigu colonne-major)
          do i = 1, u_nx
             ! Affectation et calcul numérique
             x_val = x_face(this%mesh, i)
             ! Évaluation de la condition logique
             if (x_val <= 4.0_dp + 1.0e-9_dp) then
                ! Mise à jour de la composante de vitesse discrète sur face décalée MAC
                this%fields%u(i, j) = u_val
             end if
          end do
       end do
       ! Appel de sous-programme modulaire
       call sync_u(this%mesh, this%fields%u)
    end if

    ! Initialisation du champ de vitesse solénoïdal pour LIQUID_SHEET (Équations 17-18 Fluids 2021)
    if (this%config%ic_type == IC_LIQUID_SHEET) then
       ! Déclaration des variables physiques et géométriques
       if (this%config%ic_center_y > 0.0_dp) then
          y_c_ls = this%config%ic_center_y
       else
          y_c_ls = 0.5_dp * this%mesh%ly
       end if
       if (this%config%ic_thickness > 0.0_dp) then
          delta_ls = this%config%ic_thickness
       else
          delta_ls = 0.0003_dp
       end if
       k_ls = 2.0_dp * PI / this%mesh%lx
       amp_ls = 0.04_dp
       factor_u_ls = amp_ls * (this%mesh%lx / (2.0_dp * PI)) * (2.0_dp / delta_ls)

       ! Composante horizontale u sur les faces décalées MAC en x
       u_nx = this%mesh%nx_local
       if (this%mesh%is_boundary_right) u_nx = u_nx + 1
       do j = 1, this%mesh%ny_local
          y_val = y_center(this%mesh, j)
          dy_c_ls = y_val - y_c_ls
          if (dy_c_ls > 1.0e-12_dp) then
             sgn_y_ls = 1.0_dp
          else if (dy_c_ls < -1.0e-12_dp) then
             sgn_y_ls = -1.0_dp
          else
             sgn_y_ls = 0.0_dp
          end if
          exp_decay_ls = exp(-2.0_dp * abs(dy_c_ls) / delta_ls)
          do i = 1, u_nx
             x_val = x_face(this%mesh, i)
             c_u_ls = 0.5_dp * (this%fields%C(i - 1, j) + this%fields%C(i, j))
             a_val_ls = c_u_ls * 30.0_dp + (1.0_dp - c_u_ls) * 2.0_dp
             this%fields%u(i, j) = a_val_ls - sgn_y_ls * factor_u_ls * cos(k_ls * x_val) * exp_decay_ls
          end do
       end do

       ! Composante verticale v sur les faces décalées MAC en y
       v_ny = this%mesh%ny_local
       if (this%mesh%is_boundary_top) v_ny = v_ny + 1
       do j = 1, v_ny
          y_val = y_face(this%mesh, j)
          dy_c_ls = y_val - y_c_ls
          exp_decay_ls = exp(-2.0_dp * abs(dy_c_ls) / delta_ls)
          do i = 1, this%mesh%nx_local
             x_val = x_center(this%mesh, i)
             this%fields%v(i, j) = amp_ls * sin(k_ls * x_val) * exp_decay_ls
          end do
       end do

       call sync_u(this%mesh, this%fields%u)
       call sync_v(this%mesh, this%fields%v)
       call apply_velocity_bc(this%fields, this%mesh, this%config, .false.)
    end if

    ! Initialisation du champ de vitesse pour VISCOUS_JET (Cas 3.3 Fluids 2021)
    if (this%config%ic_type == IC_VISCOUS_JET) then
       x_c_vj = merge(this%config%ic_center_x, 0.5_dp * this%mesh%lx, this%config%ic_center_x > 0.0_dp)
       half_w_vj = merge(0.5_dp * this%config%ic_jet_width, 0.04_dp, this%config%ic_jet_width > 0.0_dp)
       stub_h_vj = 2.0_dp * this%mesh%dy
       v_in_vj = merge(this%config%bc_top%v, -0.6_dp, this%config%bc_top%v /= 0.0_dp)
       y_val = 0.0_dp
       do i = 1, this%mesh%nx_global
          if (abs((real(i, dp) - 0.5_dp) * this%mesh%dx - x_c_vj) <= half_w_vj) y_val = y_val + this%mesh%dx
       end do
       if (y_val < 1.0e-12_dp) y_val = 2.0_dp * half_w_vj
       v_in_vj = v_in_vj * (2.0_dp * half_w_vj) / y_val

       v_ny = this%mesh%ny_local
       if (this%mesh%is_boundary_top) v_ny = v_ny + 1
       do j = 1, v_ny
          y_val = y_face(this%mesh, j)
          do i = 1, this%mesh%nx_local
             x_val = x_center(this%mesh, i)
             if (y_val >= this%mesh%ly - stub_h_vj .and. abs(x_val - x_c_vj) <= half_w_vj) then
                this%fields%v(i, j) = v_in_vj
             end if
          end do
       end do
       call sync_v(this%mesh, this%fields%v)
       call apply_velocity_bc(this%fields, this%mesh, this%config, .false.)
       call apply_vof_bc(this%fields, this%mesh, this%config)
    end if

    ! Initialisation du champ thermique si activé
    if (this%config%enable_energy) then
       ! Boucle spatiale sur les cellules locales selon l'axe vertical Y (ordonnées)
       do j = 1, this%mesh%ny_local
          ! Boucle spatiale selon l'axe horizontal X (accès contigu colonne-major)
          do i = 1, this%mesh%nx_local
             ! Affectation et calcul numérique
             frac = x_center(this%mesh, i) / this%mesh%lx
             ! Évolution du champ scalaire de température
             this%fields%T(i, j) = this%config%bc_left%T + frac * (this%config%bc_right%T - this%config%bc_left%T)
          end do
       end do
       ! Appel de sous-programme modulaire
       call sync_cell_scalar(this%mesh, this%fields%T)
       ! Appel de sous-programme modulaire
       call apply_thermal_boundary_conditions(this%mesh, this%config, this%fields)
    end if

    ! 5. Initialisation du solveur Navier-Stokes et préconditionneurs
    call this%ns_solver%init(this%mesh, this%config)
    ! Appel de sous-programme modulaire
    call this%ns_solver%assemble_preconditioner_blocks(this%mesh, this%config, this%fields, this%config%dt)

    ! 6. Initialisation de l'exportateur VTK et diagnostics CSV
    call this%vtk_writer%init(this%mesh, this%config%output_dir, this%config%case_name)

    ! Affectation et calcul numérique
    success = .true.
  end function sim_init

  !-----------------------------------------------------------------------------
  ! Calcul du pas de temps adaptatif (CFL advection + onde capillaire)
  !-----------------------------------------------------------------------------
  function sim_compute_time_step(this) result(dt)
    ! Déclaration de variable locale / paramètre de calcul
    class(simulation_t), intent(in) :: this
    ! Déclaration de variable locale / paramètre de calcul
    real(dp) :: dt

    ! Déclaration de variable locale / paramètre de calcul
    real(dp) :: local_max_speed, global_max_speed, speed, u_val, v_val
    ! Déclaration de variable locale / paramètre de calcul
    real(dp) :: dt_cfl, dt_sigma, rho_sum, h_min, dt_target
    ! Déclaration de variable locale / paramètre de calcul
    integer  :: i, j

    ! Évaluation de la condition logique
    if (.not. this%config%adaptive_dt) then
       ! Affectation et calcul numérique
       dt = this%config%dt
       ! Retour à la procédure appelante
       return
    end if

    ! Affectation et calcul numérique
    local_max_speed = 1.0e-6_dp
    ! Boucle spatiale sur les cellules locales selon l'axe vertical Y (ordonnées)
    do j = 1, this%mesh%ny_local
       ! Boucle spatiale selon l'axe horizontal X (accès contigu colonne-major)
       do i = 1, this%mesh%nx_local
          ! Affectation et calcul numérique
          u_val = abs(this%fields%u(i, j))
          ! Affectation et calcul numérique
          v_val = abs(this%fields%v(i, j))
          ! Affectation et calcul numérique
          speed = sqrt(u_val * u_val + v_val * v_val)
          ! Affectation et calcul numérique
          if (speed > local_max_speed) local_max_speed = speed
       end do
    end do
    ! Affectation et calcul numérique
    global_max_speed = global_max(this%mesh, local_max_speed)

    ! 1. Critère CFL de convection
    dt_cfl = this%config%cfl_max * min(this%mesh%dx, this%mesh%dy) / global_max_speed

    ! 2. Critère capillaire si tension superficielle active
    dt_sigma = 1.0e10_dp
    ! Évaluation de la condition logique
    if (this%config%enable_surface_tension .and. this%config%sigma > 0.0_dp) then
       ! Calcul de la masse volumique locale ou interpolée aux faces
       rho_sum = this%config%rho1 + this%config%rho2
       ! Affectation et calcul numérique
       h_min = min(this%mesh%dx, this%mesh%dy)
       ! Affectation et calcul numérique
       dt_sigma = sqrt((rho_sum * h_min**3) / (4.0_dp * PI * this%config%sigma))
    end if

    ! Affectation et calcul numérique
    dt_target = min(dt_cfl, dt_sigma)
    ! Affectation et calcul numérique
    dt = min(this%config%dt, dt_target)
  end function sim_compute_time_step

  !-----------------------------------------------------------------------------
  ! Boucle principale d'intégration temporelle Navier-Stokes
  !-----------------------------------------------------------------------------
  subroutine sim_run(this)
    ! Déclaration de variable locale / paramètre de calcul
    class(simulation_t), intent(inout) :: this

    ! Déclaration de variable locale / paramètre de calcul
    real(dp) :: t, dt, global_max_u, global_max_div, global_du_max, nu_mean
    ! Déclaration de variable locale / paramètre de calcul
    real(dp) :: local_max_u, local_max_div, local_du_max, speed, du, div_val
    ! Déclaration de variable locale / paramètre de calcul
    integer  :: step, total_iters, uzawa_max, m, iters, i, j
    ! Déclaration de variable locale / paramètre de calcul
    real(dp), allocatable :: state_vector(:), rhs(:), u_prev(:, :)

    ! Évaluation de la condition logique
    if (trim(this%config%mode) == "LAPLACIAN") then
       ! Appel de sous-programme modulaire
       call this%run_laplacian()
       ! Retour à la procédure appelante
       return
    end if
    ! Évaluation de la condition logique
    if (trim(this%config%mode) == "VORTEX_ADVECTION") then
       ! Appel de sous-programme modulaire
       call this%run_vortex_advection()
       ! Retour à la procédure appelante
       return
    end if

    ! Affectation et calcul numérique
    t = 0.0_dp
    ! Affectation et calcul numérique
    step = 0

    ! Allocation dynamique de tableau contigu en mémoire
    allocate(state_vector(this%ns_solver%n_owned_total))
    ! Allocation dynamique de tableau contigu en mémoire
    allocate(rhs(this%ns_solver%n_owned_total))
    ! Allocation dynamique de tableau contigu en mémoire
    allocate(u_prev(0:this%mesh%nx_local+1, 0:this%mesh%ny_local+1))

    ! Sauvegarde initiale à t = 0.0
    call this%vtk_writer%write_step(this%mesh, this%fields, step, t)
    ! Appel de sous-programme modulaire
    call this%vtk_writer%write_diagnostics(this%mesh, this%fields, step, t, this%config%dt)

    ! Évaluation de la condition logique
    if (this%mesh%rank == 0) then
       write(*, '(A)') ""
       write(*, '(A)') ">>> DEMARRAGE DU CALCUL COVOF 2D FORTRAN 2018 PARALLELE <<<"
       write(*, '(A8,A14,A12,A14,A16,A16,A14)') &
         "Pas", "Temps [s]", "dt [s]", "Iter Krylov", "Vitesse Max", "Div Max", "Nu Moyen"
       write(*, '(A)') repeat("-", 94)
    end if

    ! Boucle d'itération numérique
    do while (t < this%config%t_final - 1.0e-12_dp * (abs(this%config%t_final) + 1.0_dp))
       ! Affectation et calcul numérique
       step = step + 1
       ! Affectation et calcul numérique
       u_prev = this%fields%u
       ! Affectation et calcul numérique
       dt = this%compute_time_step()

       ! Évaluation de la condition logique
       if (t + dt > this%config%t_final) then
          ! Affectation et calcul numérique
          dt = this%config%t_final - t
       end if
       if (dt < 1.0e-12_dp) exit

       ! 1. Transport de l'interface si écoulement diphasique
       if (this%config%enable_two_phase) then
          ! Appel de sous-programme modulaire
          call this%tracker%advect(this%mesh, this%fields, dt)
          ! Actualisation des conditions aux limites
          call apply_velocity_bc(this%fields, this%mesh, this%config, .false.)
          call apply_vof_bc(this%fields, this%mesh, this%config)
       end if

       ! 2. Assemblage des blocs du préconditionneur
       call this%ns_solver%assemble_preconditioner_blocks(this%mesh, this%config, this%fields, dt)

       ! Affectation et calcul numérique
       total_iters = 0
       ! Évaluation de la condition logique
       if (this%config%solver_type == SOLVER_AUGMENTED_LAGRANGIAN) then
          ! Affectation et calcul numérique
          uzawa_max = this%config%uzawa_iter
       ! Branche par défaut si condition non satisfaite
       else
          ! Affectation et calcul numérique
          uzawa_max = 1
       end if

       ! Boucle d'itération numérique
       do m = 1, uzawa_max
          ! 3. Construction du second membre rhs
          call this%ns_solver%assemble_rhs(this%mesh, this%config, this%fields, dt, rhs)

          ! 4. Initialisation du vecteur d'état à partir des champs courants
          call this%ns_solver%pack_vector(this%mesh, this%config, this%fields, state_vector)

          ! 5. Résolution du système linéaire par Krylov (BiCGSTAB2 / PCG / GMRES)
          call solve_linear_system(this%ns_solver, this%mesh, this%config, this%fields, dt, &
                                   rhs, state_vector, iters)
          ! Affectation et calcul numérique
          total_iters = total_iters + iters

          ! 6. Déballage de la vitesse résolue
          call this%ns_solver%unpack_vector(this%mesh, this%config, state_vector, this%fields, .false.)

          ! 7. Mise à jour de la pression par Lagrangien augmenté (étape d'Uzawa)
          if (this%config%solver_type == SOLVER_AUGMENTED_LAGRANGIAN) then
             ! Appel de sous-programme modulaire
             call this%ns_solver%update_pressure_al(this%mesh, this%config, this%fields)
          end if
       end do

       ! 7.5 Résolution de l'énergie thermique si activée
       if (this%config%enable_energy) then
          ! Appel de sous-programme modulaire
          call this%ns_solver%solve_energy(this%mesh, this%config, this%fields, dt)
       end if

       ! Affectation et calcul numérique
       t = t + dt

       ! Calcul des métriques de contrôle
       local_max_u   = 0.0_dp
       ! Affectation et calcul numérique
       local_max_div = 0.0_dp
       ! Affectation et calcul numérique
       local_du_max  = 0.0_dp

       ! Boucle spatiale sur les cellules locales selon l'axe vertical Y (ordonnées)
       do j = 1, this%mesh%ny_local
          ! Boucle spatiale selon l'axe horizontal X (accès contigu colonne-major)
          do i = 1, this%mesh%nx_local
             ! Affectation et calcul numérique
             speed = sqrt(this%fields%u(i, j)**2 + this%fields%v(i, j)**2)
             ! Affectation et calcul numérique
             if (speed > local_max_u) local_max_u = speed

             ! Affectation et calcul numérique
             du = abs(this%fields%u(i, j) - u_prev(i, j))
             ! Affectation et calcul numérique
             if (du > local_du_max) local_du_max = du

             ! Affectation et calcul numérique
             div_val = abs((this%fields%u(i + 1, j) - this%fields%u(i, j)) / this%mesh%dx + &
                           (this%fields%v(i, j + 1) - this%fields%v(i, j)) / this%mesh%dy)
             ! Affectation et calcul numérique
             if (div_val > local_max_div) local_max_div = div_val
          end do
       end do

       ! Affectation et calcul numérique
       global_max_u   = global_max(this%mesh, local_max_u)
       ! Affectation et calcul numérique
       global_max_div = global_max(this%mesh, local_max_div)
       ! Affectation et calcul numérique
       global_du_max  = global_max(this%mesh, local_du_max)

       ! Affectation et calcul numérique
       nu_mean = 0.0_dp
       ! Évaluation de la condition logique
       if (this%config%enable_energy) then
          ! Affectation et calcul numérique
          nu_mean = this%ns_solver%compute_mean_nusselt_hot_wall(this%mesh, this%config, this%fields)
       end if

       ! Détection de convergence stationnaire
       if (.not. this%config%enable_two_phase .and. &
           this%config%steady_state_tol > 0.0_dp .and. step > 10 .and. &
           global_du_max < this%config%steady_state_tol) then
          ! Évaluation de la condition logique
          if (this%mesh%rank == 0) then
             write(*, '(A)') ""
             write(*, '(A,ES12.4,A,ES12.4,A,F10.4,A)') &
               ! Affectation et calcul numérique
               ">>> CONVERGENCE STATIONNAIRE ATTEINTE : du_max = ", global_du_max, &
               ! Calcul de la métrique spatiale élémentaire
               " < ", this%config%steady_state_tol, " (t = ", t, " s) <<<"
             ! Évaluation de la condition logique
             if (this%config%enable_energy) then
                ! Affectation et calcul numérique
                write(*, '(A,F10.5,A)') ">>> NOMBRE DE NUSSELT FINAL : Nu = ", nu_mean, " <<<"
             end if
          end if
          ! Appel de sous-programme modulaire
          call this%vtk_writer%write_step(this%mesh, this%fields, step, t)
          exit
       end if

       ! Affichage console régulier
       if (this%mesh%rank == 0 .and. (mod(step, this%config%log_interval) == 0)) then
          write(*, '(I8,F14.5,ES12.3,I14,ES16.4,ES16.4,F14.4)') &
            step, t, dt, total_iters, global_max_u, global_max_div, nu_mean
       end if

       ! Écriture des diagnostics CSV
       if (this%config%write_csv) then
          ! Appel de sous-programme modulaire
          call this%vtk_writer%write_diagnostics(this%mesh, this%fields, step, t, dt)
       end if

       ! Export des fichiers VTK
       if (mod(step, this%config%vtk_interval) == 0 .or. t >= this%config%t_final) then
          ! Appel de sous-programme modulaire
          call this%vtk_writer%write_step(this%mesh, this%fields, step, t)
       end if
    end do

    ! Appel de sous-programme modulaire
    call this%vtk_writer%finalize()

    ! Évaluation de la condition logique
    if (this%mesh%rank == 0) then
       write(*, '(A)') repeat("-", 94)
       ! Affectation et calcul numérique
       write(*, '(A,F10.4,A)') ">>> SIMULATION COVOF 2D TERMINEE AVEC SUCCES A t = ", t, " s <<<"
       write(*, '(A,A)') "Fichiers de visualisation VTK et CSV disponibles dans : ", trim(this%config%output_dir)
    end if

    ! Libération de la mémoire allouée
    deallocate(state_vector, rhs, u_prev)
  end subroutine sim_run

  !-----------------------------------------------------------------------------
  ! Cas de vérification 3.1 : Laplacien -Delta(Psi) = -1 avec solution analytique
  !-----------------------------------------------------------------------------
  subroutine sim_run_laplacian(this)
    ! Déclaration de variable locale / paramètre de calcul
    class(simulation_t), intent(inout) :: this

    double precision :: a, b, dx, dy, psi_ex, xp, yp, num, den, n1, n2, diff
    double precision :: err_l1, err_l2, err_linf, tol, rz_old, rz_new, alpha, beta, pAp, r_norm
    ! Déclaration de variable locale / paramètre de calcul
    integer :: nx, ny, n_int_x, n_int_y, n_dofs, i, j, k1, k2, row, iters, max_it, err_unit
    ! Déclaration de variable locale / paramètre de calcul
    character(len=256) :: err_file
    double precision, allocatable :: A_diag(:), A_east(:), A_west(:), A_north(:), A_south(:)
    double precision, allocatable :: rhs(:), x(:), r(:), p(:), z(:), Ap(:)

    ! Affectation et calcul numérique
    a = this%config%lx
    ! Affectation et calcul numérique
    b = this%config%ly
    ! Affectation et calcul numérique
    nx = this%config%nx
    ! Affectation et calcul numérique
    ny = this%config%ny
    ! Calcul de la métrique spatiale élémentaire
    dx = a / real(nx, dp)
    ! Calcul de la métrique spatiale élémentaire
    dy = b / real(ny, dp)

    ! Évaluation de la condition logique
    if (this%mesh%rank == 0) then
       write(*, '(A)') ""
       ! Affectation et calcul numérique
       write(*, '(A)') "========================================================"
       ! Affectation et calcul numérique
       write(*, '(A)') "  VERIFICATION DU CAS 3.1 : LAPLACIEN DELTA(PSI) = 1"
       write(*, '(A,F8.4,A,F8.4,A,I0,A,I0,A)') &
         "  Domaine [", a, " x ", b, "], Grille ", nx, " x ", ny
       ! Affectation et calcul numérique
       write(*, '(A)') "========================================================"
    end if

    ! Affectation et calcul numérique
    n_int_x = nx - 1
    ! Affectation et calcul numérique
    n_int_y = ny - 1
    ! Affectation et calcul numérique
    n_dofs = n_int_x * n_int_y

    ! Allocation dynamique de tableau contigu en mémoire
    allocate(A_diag(n_dofs), A_east(n_dofs), A_west(n_dofs), A_north(n_dofs), A_south(n_dofs))
    ! Allocation dynamique de tableau contigu en mémoire
    allocate(rhs(n_dofs), x(n_dofs), r(n_dofs), p(n_dofs), z(n_dofs), Ap(n_dofs))

    ! Affectation et calcul numérique
    A_diag  = 2.0_dp / (dx * dx) + 2.0_dp / (dy * dy)
    ! Affectation et calcul numérique
    A_east  = -1.0_dp / (dx * dx)
    ! Affectation et calcul numérique
    A_west  = -1.0_dp / (dx * dx)
    ! Affectation et calcul numérique
    A_north = -1.0_dp / (dy * dy)
    ! Affectation et calcul numérique
    A_south = -1.0_dp / (dy * dy)

    ! Affectation et calcul numérique
    rhs = -1.0_dp
    ! Affectation et calcul numérique
    x   = 0.0_dp

    ! Résolution par Gradient Conjugué préconditionné par Jacobi
    r = rhs
    ! Affectation et calcul numérique
    z = r / A_diag
    ! Affectation et calcul numérique
    p = z
    ! Affectation et calcul numérique
    rz_old = sum(r * z)

    ! Affectation et calcul numérique
    max_it = 20000
    ! Affectation et calcul numérique
    tol = this%config%tol
    ! Affectation et calcul numérique
    if (tol <= 0.0_dp) tol = 1.0e-12_dp

    ! Boucle d'itération numérique
    do iters = 1, max_it
       ! Ap = A * p
       do j = 0, n_int_y - 1
          ! Boucle spatiale selon l'axe horizontal X (accès contigu colonne-major)
          do i = 0, n_int_x - 1
             ! Affectation et calcul numérique
             row = j * n_int_x + i + 1
             ! Actualisation de la pression statique au centre de la maille
             Ap(row) = A_diag(row) * p(row)
             ! Actualisation de la pression statique au centre de la maille
             if (i > 0)           Ap(row) = Ap(row) + A_west(row)  * p(row - 1)
             ! Actualisation de la pression statique au centre de la maille
             if (i < n_int_x - 1) Ap(row) = Ap(row) + A_east(row)  * p(row + 1)
             ! Actualisation de la pression statique au centre de la maille
             if (j > 0)           Ap(row) = Ap(row) + A_south(row) * p(row - n_int_x)
             ! Actualisation de la pression statique au centre de la maille
             if (j < n_int_y - 1) Ap(row) = Ap(row) + A_north(row) * p(row + n_int_x)
          end do
       end do

       ! Affectation et calcul numérique
       pAp = sum(p * Ap)
       if (abs(pAp) < 1.0e-30_dp) exit

       ! Affectation et calcul numérique
       alpha = rz_old / pAp
       ! Affectation et calcul numérique
       x = x + alpha * p
       ! Affectation et calcul numérique
       r = r - alpha * Ap

       ! Évaluation de la norme du résidu ou de l'erreur numérique
       r_norm = sqrt(sum(r * r))
       if (r_norm < tol) exit

       ! Affectation et calcul numérique
       z = r / A_diag
       ! Affectation et calcul numérique
       rz_new = sum(r * z)
       ! Affectation et calcul numérique
       beta = rz_new / rz_old
       ! Affectation et calcul numérique
       rz_old = rz_new

       ! Affectation et calcul numérique
       p = z + beta * p
    end do

    ! Évaluation de la solution analytique exacte (formule analytique de référence)
    err_l1   = 0.0_dp
    ! Évaluation de la norme du résidu ou de l'erreur numérique
    err_l2   = 0.0_dp
    ! Évaluation de la norme du résidu ou de l'erreur numérique
    err_linf = 0.0_dp

    ! Boucle spatiale sur les cellules locales selon l'axe vertical Y (ordonnées)
    do j = 0, n_int_y - 1
       ! Boucle spatiale selon l'axe horizontal X (accès contigu colonne-major)
       do i = 0, n_int_x - 1
          ! Affectation et calcul numérique
          xp = real(i + 1, dp) * dx
          ! Affectation et calcul numérique
          yp = real(j + 1, dp) * dy
          ! Affectation et calcul numérique
          psi_ex = 0.0_dp

          ! Boucle d'itération numérique
          do k2 = 0, 59
             ! Boucle d'itération numérique
             do k1 = 0, 59
                ! Affectation et calcul numérique
                n1 = real(2 * k1 + 1, dp)
                ! Affectation et calcul numérique
                n2 = real(2 * k2 + 1, dp)
                ! Affectation et calcul numérique
                num = 16.0_dp * sin(n1 * PI * xp / a) * sin(n2 * PI * yp / b)
                ! Affectation et calcul numérique
                den = (PI**4) * n1 * n2 * ((n1 / a)**2 + (n2 / b)**2)
                ! Affectation et calcul numérique
                psi_ex = psi_ex - num / den
             end do
          end do

          ! Affectation et calcul numérique
          row = j * n_int_x + i + 1
          ! Affectation et calcul numérique
          diff = abs(x(row) - psi_ex)
          ! Évaluation de la norme du résidu ou de l'erreur numérique
          err_l1 = err_l1 + diff
          ! Évaluation de la norme du résidu ou de l'erreur numérique
          err_l2 = err_l2 + diff * diff
          ! Évaluation de la norme du résidu ou de l'erreur numérique
          if (diff > err_linf) err_linf = diff
       end do
    end do

    ! Évaluation de la norme du résidu ou de l'erreur numérique
    err_l1 = err_l1 / real(n_dofs, dp)
    ! Évaluation de la norme du résidu ou de l'erreur numérique
    err_l2 = sqrt(err_l2 / real(n_dofs, dp))

    ! Évaluation de la condition logique
    if (this%mesh%rank == 0) then
       write(*, '(A)') ">>> RESULTAT RESOLUTION LAPLACIEN COVOF 2D <<<"
       write(*, '(A,I0)')       "  Iterations : ", iters
       write(*, '(A,ES14.6)')  "  Erreur L1   : ", err_l1
       write(*, '(A,ES14.6)')  "  Erreur L2   : ", err_l2
       write(*, '(A,ES14.6)')  "  Erreur Linf : ", err_linf
       ! Affectation et calcul numérique
       write(*, '(A)') "========================================================"

       ! Appel de sous-programme modulaire
       call execute_command_line("mkdir -p " // trim(this%config%output_dir))
       ! Évaluation de la norme du résidu ou de l'erreur numérique
       err_file = trim(this%config%output_dir) // "/laplacian_error.csv"
       ! Affectation et calcul numérique
       open(newunit=err_unit, file=trim(err_file), status="replace", action="write")
       write(err_unit, '(A)') "nx,ny,dx,dy,iters,err_l1,err_l2,err_linf"
       write(err_unit, '(I0,",",I0,",",ES14.6,",",ES14.6,",",I0,",",ES14.6,",",ES14.6,",",ES14.6)') &
         nx, ny, dx, dy, iters, err_l1, err_l2, err_linf
       close(err_unit)
    end if

    ! Libération de la mémoire allouée
    deallocate(A_diag, A_east, A_west, A_north, A_south, rhs, x, r, p, z, Ap)
  end subroutine sim_run_laplacian

  !-----------------------------------------------------------------------------
  ! Cas de vérification 3.19 : Advection VOF vortex de cisaillement avec retournement
  !-----------------------------------------------------------------------------
  subroutine sim_run_vortex_advection(this)
    ! Déclaration de variable locale / paramètre de calcul
    class(simulation_t), intent(inout) :: this

    ! Déclaration de variable locale / paramètre de calcul
    integer  :: nx, ny, u_nx, v_ny, step, i, j, err_unit
    ! Déclaration de variable locale / paramètre de calcul
    real(dp) :: T_period, dt, t, cos_factor, x_u, y_c, x_c, y_v, sx, sy
    ! Déclaration de variable locale / paramètre de calcul
    real(dp) :: local_err_l1, global_err_l1, local_cells, global_cells, err_l1
    ! Déclaration des variables pour erreurs L2, Linf et bilan de conservation de masse
    real(dp) :: local_err_l2, global_err_l2, err_l2
    ! Déclaration de l'erreur absolue maximale locale et globale
    real(dp) :: local_err_linf, global_err_linf, err_linf
    ! Déclaration des masses initiale et finale locales et globales
    real(dp) :: local_mass_init, global_mass_init, local_mass_final, global_mass_final
    ! Déclaration du volume de maille, écart local et ratio de conservation
    real(dp) :: cell_vol, diff, mass_conservation
    ! Déclaration de variable locale / paramètre de calcul
    character(len=256) :: err_file
    ! Déclaration de variable locale / paramètre de calcul
    real(dp), allocatable :: C_initial(:, :)

    ! Affectation et calcul numérique
    nx = this%config%nx
    ! Affectation et calcul numérique
    ny = this%config%ny
    ! Affectation et calcul numérique
    T_period = this%config%t_final
    ! Affectation et calcul numérique
    if (T_period <= 0.0_dp) T_period = 6.0_dp
    ! Affectation et calcul numérique
    dt = this%config%dt

    ! Évaluation de la condition logique
    if (this%mesh%rank == 0) then
       write(*, '(A)') ""
       ! Affectation et calcul numérique
       write(*, '(A)') "========================================================"
       write(*, '(A)') "  VERIFICATION DU CAS 3.19 : ADVECTION VORTEX CISAILLEMENT"
       write(*, '(A,I0,A,I0,A,F8.4,A,F8.5,A)') &
         ! Affectation et calcul numérique
         "  Grille : ", nx, " x ", ny, " | Periode T = ", T_period, " s | dt = ", dt, " s"
       ! Affectation et calcul numérique
       write(*, '(A)') "========================================================"
    end if

    ! Allocation dynamique de tableau contigu en mémoire
    allocate(C_initial(0:this%mesh%nx_local+1, 0:this%mesh%ny_local+1))
    ! Affectation et calcul numérique
    C_initial = this%fields%C

    ! Appel de sous-programme modulaire
    call this%vtk_writer%write_step(this%mesh, this%fields, 0, 0.0_dp)

    ! Affectation et calcul numérique
    t = 0.0_dp
    ! Affectation et calcul numérique
    step = 0
    ! Affectation et calcul numérique
    u_nx = this%mesh%nx_local
    ! Affectation et calcul numérique
    if (this%mesh%is_boundary_right) u_nx = u_nx + 1
    ! Affectation et calcul numérique
    v_ny = this%mesh%ny_local
    ! Affectation et calcul numérique
    if (this%mesh%is_boundary_top) v_ny = v_ny + 1

    ! Boucle d'itération numérique
    do while (t < T_period - 1.0e-12_dp)
       ! Affectation et calcul numérique
       step = step + 1
       ! Affectation et calcul numérique
       if (t + dt > T_period) dt = T_period - t

       ! Affectation et calcul numérique
       cos_factor = cos(PI * t / T_period)

       ! Imposer le champ de vitesse analytique sur les faces décalées
       do j = 1, this%mesh%ny_local
          ! Affectation et calcul numérique
          y_c = y_center(this%mesh, j)
          ! Boucle spatiale selon l'axe horizontal X (accès contigu colonne-major)
          do i = 1, u_nx
             ! Affectation et calcul numérique
             x_u = x_face(this%mesh, i)
             ! Affectation et calcul numérique
             sx = sin(PI * x_u)
             ! Affectation et calcul numérique
             sy = sin(2.0_dp * PI * y_c)
             ! Mise à jour de la composante de vitesse discrète sur face décalée MAC
             this%fields%u(i, j) = - (sx * sx) * sy * cos_factor
          end do
       end do

       ! Boucle spatiale sur les cellules locales selon l'axe vertical Y (ordonnées)
       do j = 1, v_ny
          ! Affectation et calcul numérique
          y_v = y_face(this%mesh, j)
          ! Boucle spatiale selon l'axe horizontal X (accès contigu colonne-major)
          do i = 1, this%mesh%nx_local
             ! Affectation et calcul numérique
             x_c = x_center(this%mesh, i)
             ! Affectation et calcul numérique
             sx = sin(2.0_dp * PI * x_c)
             ! Affectation et calcul numérique
             sy = sin(PI * y_v)
             ! Mise à jour de la composante de vitesse discrète sur face décalée MAC
             this%fields%v(i, j) = sx * (sy * sy) * cos_factor
          end do
       end do

       ! Appel de sous-programme modulaire
       call sync_u(this%mesh, this%fields%u)
       ! Appel de sous-programme modulaire
       call sync_v(this%mesh, this%fields%v)

       ! Advection TVD Superbee de la fraction volumique C
       call this%tracker%advect(this%mesh, this%fields, dt)

       ! Affectation et calcul numérique
       t = t + dt

       ! Affectation et calcul numérique
       if (mod(step, this%config%vtk_interval) == 0 .or. t >= T_period .or. &
           abs(t - 0.5_dp * T_period) < 0.5_dp * dt) then
          ! Appel de sous-programme modulaire
          call this%vtk_writer%write_step(this%mesh, this%fields, step, t)
       end if

       ! Évaluation de la condition logique
       if (this%mesh%rank == 0 .and. (mod(step, this%config%log_interval) == 0 .or. t >= T_period)) then
          write(*, '(A,I0,A,F8.4,A,F8.4,A)') &
            ! Affectation et calcul numérique
            "  [Vortex] Pas ", step, " | t = ", t, " s / ", T_period, " s"
       end if
    end do

    ! Calcul des erreurs L1, L2, Linf et du bilan de masse entre C_initial et C(T)
    local_err_l1 = 0.0_dp
    ! Affectation et initialisation numérique de la somme d'erreur quadratique L2
    local_err_l2 = 0.0_dp
    ! Affectation et initialisation numérique de l'erreur absolue maximale Linf
    local_err_linf = 0.0_dp
    ! Affectation et initialisation numérique de la masse initiale locale
    local_mass_init = 0.0_dp
    ! Affectation et initialisation numérique de la masse finale locale
    local_mass_final = 0.0_dp
    ! Affectation et calcul du nombre de cellules locales
    local_cells = real(this%mesh%nx_local * this%mesh%ny_local, dp)
    ! Affectation et calcul de l'aire géométrique élémentaire de cellule
    cell_vol = this%mesh%dx * this%mesh%dy

    ! Boucle spatiale sur les cellules locales selon l'axe vertical Y (ordonnées)
    do j = 1, this%mesh%ny_local
       ! Boucle spatiale selon l'axe horizontal X (accès contigu colonne-major)
       do i = 1, this%mesh%nx_local
          ! Évaluation de l'écart absolu entre la fraction finale et la fraction exacte initiale
          diff = abs(this%fields%C(i, j) - C_initial(i, j))
          ! Accumulation pour la norme L1 moyenne
          local_err_l1 = local_err_l1 + diff
          ! Accumulation quadratique pour la norme L2
          local_err_l2 = local_err_l2 + diff * diff
          ! Recherche du maximum pour la norme Linf
          if (diff > local_err_linf) local_err_linf = diff
          ! Intégration de la masse initiale
          local_mass_init = local_mass_init + C_initial(i, j) * cell_vol
          ! Intégration de la masse finale
          local_mass_final = local_mass_final + this%fields%C(i, j) * cell_vol
       end do
    end do

    ! Réduction parallèle MPI globale de la somme d'erreur L1
    global_err_l1 = global_sum(this%mesh, local_err_l1)
    ! Réduction parallèle MPI globale de la somme quadratique L2
    global_err_l2 = global_sum(this%mesh, local_err_l2)
    ! Réduction parallèle MPI globale de l'erreur maximale Linf
    global_err_linf = global_max(this%mesh, local_err_linf)
    ! Réduction parallèle MPI globale de la masse totale initiale
    global_mass_init = global_sum(this%mesh, local_mass_init)
    ! Réduction parallèle MPI globale de la masse totale finale
    global_mass_final = global_sum(this%mesh, local_mass_final)
    ! Affectation et calcul du nombre global de cellules du maillage
    global_cells  = global_sum(this%mesh, local_cells)

    ! Normalisation de l'erreur moyenne L1
    err_l1 = global_err_l1 / global_cells
    ! Normalisation et racine de l'erreur quadratique moyenne L2
    err_l2 = sqrt(global_err_l2 / global_cells)
    ! Erreur maximale locale Linf
    err_linf = global_err_linf
    ! Évaluation du taux de conservation de masse du fluide
    if (global_mass_init > 1.0e-14_dp) then
       mass_conservation = global_mass_final / global_mass_init
    else
       mass_conservation = 1.0_dp
    end if

    ! Évaluation de la condition logique pour le processus maître de rang 0
    if (this%mesh%rank == 0) then
       write(*, '(A)') ""
       write(*, '(A)') ">>> RESULTAT DU CAS 3.19 : ADVECTION VORTEX COVOF 2D <<<"
       write(*, '(A,A)')       "  Methode Advection  : ", trim(this%config%advection_method)
       write(*, '(A,A)')       "  Reconstruction     : ", trim(this%config%reconstruction_method)
       write(*, '(A,A)')       "  Limiteur TVD/NVD   : ", trim(this%config%tvd_limiter)
       write(*, '(A,I0,A,I0)') "  Maillage           : ", nx, " x ", ny
       write(*, '(A,ES14.6)')  "  Erreur L1          : ", err_l1
       write(*, '(A,ES14.6)')  "  Erreur L2          : ", err_l2
       write(*, '(A,ES14.6)')  "  Erreur Linf        : ", err_linf
       write(*, '(A,F10.5,A)') "  Conservation Masse : ", mass_conservation * 100.0_dp, " %"
       ! Affectation et calcul numérique
       write(*, '(A)') "========================================================"

       ! Appel de sous-programme modulaire pour création du dossier de sortie
       call execute_command_line("mkdir -p " // trim(this%config%output_dir))
       ! Chemin d'accès au fichier CSV des erreurs et bilans de masse
       err_file = trim(this%config%output_dir) // "/vortex_error.csv"
       ! Ouverture en écriture du fichier récapitulatif
       open(newunit=err_unit, file=trim(err_file), status="replace", action="write")
       ! Écriture de l'en-tête CSV avec l'ensemble des métadonnées
       write(err_unit, '(A)') "nx,ny,dt,T,advection,reconstruction,limiter,err_l1,err_l2,err_linf,mass_init,mass_final,mass_conservation"
       ! Écriture formatée des valeurs numériques
       write(err_unit, '(I0,",",I0,",",ES14.6,",",ES14.6,",",A,",",A,",",A,",",ES14.6,",",ES14.6,",",ES14.6,",",ES14.6,",",ES14.6,",",ES14.6)') &
         nx, ny, dt, T_period, &
         trim(this%config%advection_method), trim(this%config%reconstruction_method), trim(this%config%tvd_limiter), &
         err_l1, err_l2, err_linf, global_mass_init, global_mass_final, mass_conservation
       ! Fermeture propre du flux fichier
       close(err_unit)
    end if

    ! Appel de sous-programme modulaire
    call this%vtk_writer%finalize()
    ! Libération de la mémoire allouée
    deallocate(C_initial)
  end subroutine sim_run_vortex_advection

end module module_simulation

\end{lstlisting}

\end{document}
